\documentclass[a4paper,11pt]{ltjsarticle}
\usepackage[margin=23mm,headheight=16pt]{geometry}
\usepackage{luatexja-fontspec}
\setmainfont{Latin Modern Roman}
\setsansfont{Latin Modern Sans}
\setmonofont{Latin Modern Mono}[Scale=MatchLowercase]
\setmainjfont{HaranoAjiMincho-Regular.otf}[BoldFont=HaranoAjiGothic-Medium.otf]
\setsansjfont{HaranoAjiGothic-Medium.otf}
\usepackage{amsmath,amssymb,amsthm,mathtools}
\newcommand{\mathscr}{\mathcal}
\usepackage{booktabs,array,enumitem,xcolor,fancyhdr,needspace}
\usepackage[unicode,colorlinks=true,linkcolor=blue!45!black,urlcolor=blue!50!black,citecolor=blue!45!black]{hyperref}
\hypersetup{pdftitle={非可換 cubic threefold：形式変形空間・残余圏・貼り合わせと点加群},pdfauthor={Research note}}
\setlength{\emergencystretch}{2em}
\setlength{\parskip}{0.28em}
\setlist{nosep,leftmargin=1.7em}
\allowdisplaybreaks[2]
\pagestyle{fancy}
\fancyhf{}
\fancyhead[L]{\small 非可換 cubic threefold}
\fancyhead[R]{\small 研究ノート・第18版}
\fancyfoot[C]{\thepage}
\renewcommand{\headrulewidth}{0.3pt}
\theoremstyle{definition}
\newtheorem{theorem}{命題}[section]
\newtheorem{lemma}[theorem]{補題}
\newtheorem{remark}[theorem]{注意}
\newtheorem{proposition}[theorem]{命題}
\newtheorem{corollary}[theorem]{系}
\newcommand{\CC}{\mathbb C}
\newcommand{\OO}{\mathcal O}
\newcommand{\PP}{\mathbb P}
\newcommand{\Aff}{\mathbb A}
\newcommand{\Ku}{\operatorname{Ku}}
\newcommand{\HH}{\operatorname{HH}}
\newcommand{\HP}{\operatorname{HP}}
\newcommand{\Gr}{\operatorname{Gr}}
\newcommand{\Spf}{\operatorname{Spf}}
\newcommand{\Def}{\operatorname{Def}}
\newcommand{\Sym}{\operatorname{Sym}}
\newcommand{\Ext}{\operatorname{Ext}}
\newcommand{\Hom}{\operatorname{Hom}}
\newcommand{\Perf}{\operatorname{Perf}}
\newcommand{\qgr}{\operatorname{qgr}}
\newcommand{\Dsg}{D_{\mathrm{sg}}}
\newcommand{\Art}{\operatorname{Art}_{\CC}}
\newcommand{\NC}{\operatorname{NC}}
\newcommand{\Pois}{\operatorname{Pois}}
\newcommand{\RHom}{\operatorname{RHom}}
\newcommand{\id}{\operatorname{id}}
\newcommand{\im}{\operatorname{im}}
\newcommand{\rk}{\operatorname{rank}}
\newcommand{\frakm}{\mathfrak m}
\newcommand{\MF}{\operatorname{MF}}
\newcommand{\CatDef}{\operatorname{CatDef}}
\newcommand{\Pic}{\operatorname{Pic}}
\newcommand{\Ind}{\operatorname{Ind}}
\title{\bfseries 非可換 cubic threefold\\[0.35em]
\large 形式変形空間・残余圏・貼り合わせと点加群}
\author{}
\date{2026年10月30日・第18版}

\begin{document}
\maketitle
\vspace{-1.5em}
\begin{abstract}
\noindent\textbf{第18版の主な進展.}
実際のBrill--Noether境界の横断三次項について、Hessian側の係数が全次数で零であることを証明した。
固定平面のPicard族に沿うnormal coneが同じ抽象楕円曲線を与えることと、
Hesse pencilの非定数な $j$ 不変量を用いる。
有限個のTate係数の零性からの推測ではない。
また横断係数の全形式接触変換則を求め、この消滅が普遍形式正規化で保存されることを示した。

第17版の全通常商の分裂
\[
 \bigoplus_n\mathcal C_n/J_n
 =\mathbb C[\overline U_4,\overline J_6]\cdot[\Theta_8]
       \oplus\mathcal K^{\mathrm{nor}}
\]
において、実際の境界の全係数は有限次元の後半部分に入る。
さらに45変数の斉次Gr\"obner次数上界から
\[
 \mathcal K_n^{\mathrm{nor}}=0\quad(n>N_*),
 \qquad N_*=20\cdot18^{\,2^{43}}-10
\]
という明示上界を得た。
この数は極めて大きいが、normal消滅を満たす境界はその $N_*$-jet に普遍形式右同値となる。
従って実際の境界には、$P+H/96$ に有限個の頂点成分を加えた多項式の形式代表が存在する。
この無限合成の収束や代数化は、ここでは主張しない。

Hessianについて全座標 $a_{ij}$ がその類を殺す明示恒等式を得た。
一方、$\partial_{q_j}H$ も全頂点局所コホモロジーの非零元となるため、
全頂点加群と、その指定された同変成分 $\mathcal K^{\mathrm{nor}}$ は区別しなければならない。

二十六次では $\dim\mathcal C_{26}=6377$、$\dim\mathcal U_{20}=3712$ などの
全指標を厳密に求め、二入力の法三次Tate再帰と独立保存監査でnormal係数の零性を確認した。
$\mathcal K_{26}^{\mathrm{nor}}=0$、$E_{\mathrm{vert}}=t^{10}$、
および全次数の $P+H/96$ 正規形は未確定である。
前版の26次開始の代数的・収束的偶正規形はそのまま維持する。

本稿は非可換cubic threefoldの形式変形空間、残余圏、貼り合わせと点加群の継続研究ノートである。
前版までの議論を履歴として残し、今回の新しい結果を第66節から第72節にまとめた。
証明、有限計算、保存値の監査、新規再帰、未実行の探索を区別して記録する。

\end{abstract}

\section{対象と、今回確定する主張}

以下、$X=V(F)\subset\PP^4_{\CC}$ は滑らかなcubic threefoldとする。
「非可換cubic threefold」には複数の定式化があり、変形対象を明示する必要がある。

\begin{center}
\small
\begin{tabular}{p{.18\textwidth}p{.36\textwidth}p{.36\textwidth}}
\toprule
水準 & 変形する対象 & 今回の結論\\
\midrule
構造層 &
$\OO_X$ のmarkedな平坦結合代数層 &
古典的形式基底を5本の二次式で記述\\
次数付き代数 &
$\CC[x_0,\ldots,x_4]/(F)$ の次数を保つ変形 &
一般の形式構成と、Fermatの代数的1パラメータ族\\
残余圏 &
$\Ku(X)=\langle\OO_X,\OO_X(1)\rangle^\perp$ のdg変形 &
古典的Morita変形の10次元基底、固定ファイバー、Fermat族の二次類\\
\bottomrule
\end{tabular}
\end{center}

基底環は可換な局所Artin $\CC$-代数とし、特殊ファイバーとの同一視を固定する。
従って、ここでの「非可換」は変形後の構造代数についての語である。
非可換Artin環をパラメータとするLaudal型変形理論を扱うものではない。

構造層の変形はYekutieliの定義\cite[Definitions 2.7--2.8]{Yek}に従う。
$A\in\Art$ に対し、$\NC_X(A)$ を、位相空間 $X$ 上の平坦な結合 $A$-代数層
$\mathscr B$ と同一視 $\mathscr B/\frakm_A\mathscr B\simeq\OO_X$ の組の、
同一視を保つ代数層同型による同値類とする。
$\Pois_X(A)$ は、可換積そのものも変形でき、特殊ファイバーで括弧が零となる
Poisson代数層の対応する関手である。

\medskip
\noindent\textbf{結論の位置づけ。}
第2節の形式環と第5節の代数族は第1版の内容である。
第7節では、第1版で条件付きだった圏論的比較を、古典的Artin基底上の
compactly generated Morita変形の意味で完成する。
第8節以下のHesse比較では、単なる点曲線の一致で済ませず、
次数シフトと三次周期性を含むdg模型を比較する。
第\ref{sec:gluing-objects}--\ref{sec:points}節は第3版の追加であり、
貼り合わせの局所変形問題と、明示的な点加群スキームを扱う。
第\ref{sec:derived-gluing}--\ref{sec:reconstruction}節が第4版の追加であり、
導来構造、代数的モジュライでの実現、楕円曲線の復元を扱う。
第\ref{sec:general-tangent-comparison}--\ref{sec:five-generator}節が第5版の追加である。
一般の接錐の比較とBrill--Noether境界の六次接錐、
明示的族の評価核による記述、五生成元代数の内在的分類を扱う。
第\ref{sec:picard-compatibility}--\ref{sec:boundary-torelli}節が第6版の追加である。
相対Picardによる適合部分空間の全次数の記述、一般評価核点の境界の形式局所構造、
および六次境界によるTorelliを扱う。
第\ref{sec:global-bn-singular}--\ref{sec:boundary-higher-jets}節が第7版の追加である。
境界の全特異部分、埋め込み解消と対数標準性、普遍同変方程式の9次までの斉次性、
および10次の候補を扱う。
第\ref{sec:tate-boundary-tenth}節が第8版の追加である。
Tate行列による実際の対象・incidence・境界の構成、係数 $1/96$ の厳密計算、
全ての滑らかな三次式でのJacobian障害、10次jetでの非同型、および
incidenceの形式法束模型の非存在を証明する。
第\ref{sec:tate-algebraization-twelfth}節が第9版の追加である。
Tate再帰の有限代数化、収束評価、実際のétale変形族、および12次項と座標変更の作用を扱う。
第\ref{sec:contact-twelve-thirteen}節が第10版の追加である。
12次商の次元・生成元・係数汎関数を決定し、実際の12次類の消滅を証明する。
また13次の接触商の消滅により、実際の境界を13次まで正規化する。
第\ref{sec:contact-fourteen}節が第11版の追加である。
14次の接触階数・商・係数汎関数を決定し、低次安定化群の作用を含めて、
実際の14次項の消去を証明する。対角評価の限界と非対角微分の役割も明示する。
第\ref{sec:all-order-parity}--\ref{sec:contact-fifteen}節が第12版の追加である。
Clifford反転により全奇数次を同時に消去し、12次・14次の正規化との両立を証明する。
またPfaffian微分の全次数の全射性、15次の接触商の消滅、
実際の15次係数の厳密計算、およびKoszulホモロジーの非消滅を扱う。

第\ref{sec:algebraic-parity}--\ref{sec:contact-sixteen}節が第13版の追加である。
全次数偶化の代数性と収束性、16次接触商と低次安定化群の作用、
一般テンソルによる厳密な階数証明、および実際の16次係数の判定を扱う。




\section{構造層の形式変形空間：17次元のPfaffian特異点}

\subsection{主張}

$0\leq i<j\leq4$ に対して10個の変数 $a_{ij}$ をとり、
\[
 q_{ijkl}=a_{ij}a_{kl}-a_{ik}a_{jl}+a_{il}a_{jk}
 \qquad(0\leq i<j<k<l\leq4)
\]
とおく。これは5個のPlücker二次式である。

\begin{theorem}[markedな古典的形式変形基底]\label{thm:base}
非標準的な関手の同型
\[
 \NC_X(A)\simeq\operatorname{Hom}_{\mathrm{cont,loc},\CC}(R_{\NC},A)
\]
が存在し、
\begin{equation}\label{eq:base}
 R_{\NC}=
 \frac{\CC[[t_1,\ldots,t_{10},a_{01},a_{02},a_{03},a_{04},
                         a_{12},a_{13},a_{14},a_{23},a_{24},a_{34}]]}
 {(q_{0123},q_{0124},q_{0134},q_{0234},q_{1234})}.
\end{equation}
従って古典的形式変形空間は
\[
 \widehat{\Aff}^{10}\times
 \widehat{C(\Gr(2,5))}_{\,0}
\]
である。
\end{theorem}

これはYekutieliの量子化定理\cite[Corollary 11.2]{Yek}と、
MayanskiyのSchouten括弧計算\cite[Section 2.2]{May}の相対版から導く帰結である。
以下では非被約な基底上の議論を含めて証明する。

\subsection{証明：Poisson変形への移行}

三次超曲面の標準的なコホモロジー計算から
\[
 H^1(X,\OO_X)=H^2(X,\OO_X)=0.
\]
YekutieliのCorollary 11.2は
\[
 \Pois_X(A)\xrightarrow{\;\sim\;}\NC_X(A)
\]
を与え、パラメータ環の準同型との可換性を明記している。
従ってこれはArtin変形関手の同型である。
この定理は、固定した $\OO_X[[h]]$ 上の大域star積だけに限定されず、
可換な構造層の変形も含むことが重要である。

Poisson変形の可換代数層は通常のスキーム変形を与える。
実際、滑らかなアフィン開集合上では、同論文のGeometrization Theorem 3.6により、
基礎の可換代数層は $A\otimes\OO_U$ に局所的に同型である。
これらは可換代数層の同型で貼り合わされるので、通常の $A$-スキーム $X_A$ を与える。
従って、$X$ の通常の変形と、その上の相対Poisson括弧を分類すればよい。

\subsection{証明：通常の変形の普遍形式族}

Euler列と法束列から
\[
 H^0(T_X)=H^2(T_X)=0,\qquad h^1(T_X)=10.
\]
特に通常のmarked変形には無限小自己同型がなく、障害もない。
$W_3=H^0(\PP^4,\OO(3))$ は35次元であり、
\[
 \mathfrak{gl}_5\longrightarrow W_3,\qquad M\longmapsto M(F)
\]
は単射である。実際、核の元は $X$ のベクトル場を与え、その消滅からスカラー行列となるが、
スカラー行列 $\lambda I$ は $F$ を $3\lambda F$ に写す。
従って核は零である。

$\mathfrak{gl}_5\cdot F$ の10次元の補空間をとり、その基底を $G_1,\ldots,G_{10}$ とする。
\[
 S=\CC[[t_1,\ldots,t_{10}]],\qquad
 F_t=F+\sum_{\alpha=1}^{10}t_\alpha G_\alpha
\]
は形式的なsliceを与える。これは $GL_5$ の作用の微分と補空間の直和が同型であること、
および形式的逆関数定理による。

$\OO_X(1)$ の延長の障害・曖昧さはそれぞれ $H^2(\OO_X)$、$H^1(\OO_X)$ に入り、
ともに零である。その切断も延長し、変形は再び $\PP^4$ 内の三次超曲面となる。
従ってこのsliceは抽象的な通常のmarked変形をも表す。
特殊ファイバーで恒等となる自己同型は、Artin環の長さに関する帰納法と
$H^0(T_X)=0$ から恒等である。

\subsection{証明：相対括弧と冪零基底}

$p:X_S\to\Spf S$ を上の普遍族とする。
線形形式の自由加群 $V_S=p_*\OO_{X_S}(1)$ と随伴の同一視を選ぶ。
基底からの可逆層の因子が現れても、局所形式基底では自明化できる。
相対Euler列・余法束列から、任意のArtin基底変換と可換な同型
\[
 p_*\Lambda^2T_{X_S/S}\simeq\Lambda^2 V_S,\qquad
 p_*\Lambda^3T_{X_S/S}\simeq\Sym^2V_S
\]
を得る。階数はそれぞれ10、15である。
例えば $\Lambda^2T\simeq\Omega^1(2)$ と
\[
 0\longrightarrow\OO_{X_S}(-1)
 \longrightarrow\Omega^1_{\PP^4_S/S}(2)|_{X_S}
 \longrightarrow\Omega^1_{X_S/S}(2)\longrightarrow0
\]
を用いる。左端は特殊ファイバーで全コホモロジーが零であり、Artin基底上でも同じである。

$\epsilon_{ij}=x_j\,dx_i-x_i\,dx_j$ をこの同型を通じて双ベクトルとみなし、
$\pi=\sum_{i<j}a_{ij}\epsilon_{ij}$ と書く。
Mayanskiyの計算は、規約をこのように固定すると
\begin{align}
 \frac14[\pi,\pi]
 ={}&q_{1234}(F_t)_0-q_{0234}(F_t)_1+q_{0134}(F_t)_2 \notag\\
    &-q_{0124}(F_t)_3+q_{0123}(F_t)_4, \label{eq:bracket}
\end{align}
を与える。ここで $(F_t)_i=\partial F_t/\partial x_i$ である。
この式は微分と斉次Euler公式から得る多項式恒等式なので、相対的にも成立する。

滑らかさから $F_0,\ldots,F_4$ は線形独立である。
定数係数の非自明な線形関係があれば、座標変換後に $F$ が一変数を含まず、
$X$ が特異な射影錐になるからである。
15個の二次単項式で勾配を表した係数行列の、ある $5\times5$ 小行列式は特殊ファイバーで非零である。
従って形式基底上で可逆であり、
\[
 S^5\longrightarrow\Sym^2V_S,\qquad (b_i)\longmapsto\sum_i b_i(F_t)_i
\]
は分裂単射となる。この単射性は任意のArtin基底変換の後も保たれる。

よって \eqref{eq:bracket} が零であることは、非被約基底上でも
\[
 q_{0123}=q_{0124}=q_{0134}=q_{0234}=q_{1234}=0
\]
と同値である。量子化による関手同型と合わせて命題\ref{thm:base}を得る。
\hfill$\square$

\subsection{次元、特異点、および主張の限界}

$\Gr(2,5)$ の次元は6なので、そのアフィン錐は7次元である。
従って
\[
 \dim R_{\NC}=17,\qquad
 \dim_\CC\frakm_{R_{\NC}}/\frakm_{R_{\NC}}^2=20.
\]
5本の二次式は最小生成系だが、余次元は3である。
従って「20変数から5を引いて15次元」と数えてはならず、この環は完全交叉ではない。
例えば $a_{01}\ne0$ のチャートでは、3本の式から $a_{23},a_{24},a_{34}$ を解け、
残りの2本は従属する。

錐は頂点以外で滑らかであり、被約な特異点集合は $a=0$、すなわち10次元の可換変形の部分である。
Jacobianイデアルが定める特異部分スキームの厚みまでは、この言い方では主張していない。

HKR分解と上のコホモロジー計算は
\[
 \HH^2(X)\simeq H^1(T_X)\oplus H^0(\Lambda^2T_X)\simeq\CC^{10}\oplus\CC^{10},
 \qquad \HH^3(X)\simeq\CC^{15}
\]
を与える\cite{May}。$\HH^3$ は障害の入り得る空間であり、最小方程式の本数ではない。
また、\eqref{eq:base} は\textbf{古典的集合値関手}の環であり、
derived formal moduliの全ての高次構造をこの通常の環に同定する主張ではない。
markingを忘れれば $\operatorname{Aut}(X)$ の有限群作用も考慮する必要がある。
積分解はsliceと自明化の選択に依存する。

\section{同時に動かせる非可換方向は最大4次元}

17次元の形式空間の特異性には、具体的な意味がある。
固定したcubic上のPoissonテンソルは、$\Lambda^2V$ の分解可能な元、
すなわちGrassmannian錐の点に対応する\cite[Theorem 1]{May}。
任意の二つのPoisson方向を独立なパラメータで同時に量子化できるとは限らない。

\begin{lemma}\label{lem:linear}
$V$ を5次元とし、線形部分空間 $L\subset\Lambda^2V$ の全ての元が分解可能とする。
このとき、ある直線 $\ell\subset V$ または3次元部分空間 $W\subset V$ に対して
\[
 L\subset\ell\wedge V
 \quad\text{または}\quad
 L\subset\Lambda^2W.
\]
特に $\dim L\le4$ である。
\end{lemma}

\begin{proof}
非比例な $a,b\in L$ に対し $(a+b)^2=0$ なので $a\wedge b=0$。
対応する二つの2平面は直線で交わり、基底を選んで
$\langle e_0,e_1\rangle$、$\langle e_0,e_2\rangle$ と書ける。
双方と交わる別の2平面は、$\langle e_0\rangle$ を含むか、
$W=\langle e_0,e_1,e_2\rangle$ に含まれる。
後者で $e_0$ を含まないものが一つでもあれば、他の全ての2平面も $W$ に含まれる。
そうでなければ全てが $\langle e_0\rangle$ を含む。
\end{proof}

従って、可換点を通る滑らかな形式基底
$\Spf\CC[[s_1,\ldots,s_r]]$ 上の族の非可換接方向の像は最大4次元である。
実際、$a_{ij}=a_{ij}(s)$ をPlücker式に代入した次数2の項は、
一次項の像が補題の $L$ になることを意味する。
幾何学的方向も含めた接写像の階数は最大14となる。
これは\textbf{基底の次元や一般点での階数の上限ではなく、可換点での接写像の上限}である。
$\ell$ を固定して $\ell\wedge V$ を動かすと、4方向の上限は達成される。

もう一つの帰結として、非零Poissonテンソル $\pi$ における
$[\pi,-]:H^0(\Lambda^2T_X)\to H^0(\Lambda^3T_X)$ の階数は3である。
これはPlücker式の微分と分裂単射な勾配写像の合成だからである。
MayanskiyのPoissonコホモロジー公式\cite[Theorem 2]{May}を代入すると
\[
 \dim\HP^0=1,\quad \dim\HP^1=0,\quad
 \dim\HP^2=17,\quad \dim\HP^3=12.
\]
従って障害群が非零でも、Poisson対の通常のモジュライはこの点で滑らかであり得る。

\section{一般のcubic：次数を保つ形式量子化}

\subsection{線形束に対応するJacobian括弧}

座標を $(u,v,x,y,z)$ とする。三次式 $F=F(u,v,x,y,z)$ に対して
\begin{equation}\label{eq:jac}
 \{x,y\}=F_z,\qquad \{y,z\}=F_x,\qquad \{z,x\}=F_y,
 \qquad \{u,-\}=\{v,-\}=0
\end{equation}
とおく。これは $\CC[u,v]$ 上の三変数Jacobian Poisson括弧であり、
$F$ はCasimirである。括弧の係数は二次式なので
$\deg\{f,g\}=\deg f+\deg g$ となる。

この括弧は $F=0$ に降り、次数0の局所化を通して $X$ 上のPoisson構造を与える。
対応する分解可能な2形式は $u\,dv-v\,du$ に比例する。
任意の非零分解可能な方向は線形座標変換でこの形にできる。
ただしArtin基底上の普遍Plücker解が常に二つの線形形式の外積として選べる、
という主張ではない。

\subsection{中央の三次式を同時に量子化する}

\begin{theorem}[各Poisson方向の形式的な次数付きモデル]\label{thm:formalgraded}
\eqref{eq:jac} に対して、次数を保つ結合形式変形 $B_h$ と次数3の中心元
$\widehat F=F+O(h)\in B_h$ が存在する。
$u,v$ は厳密に中心的であり、
\[
 A_h^{\mathrm{form}}=B_h/(\widehat F)
\]
は各次数で有限自由な $\CC[[h]]$-加群となる。そのHilbert級数は
\[
 H_{A_h^{\mathrm{form}}}(t)=\frac{1-t^3}{(1-t)^5}
\]
であり、一次交換子は \eqref{eq:jac} を与える。
\end{theorem}

\begin{proof}
アフィン5空間上のKontsevichの多項式formality写像 $\mathcal U$ を使う
\cite[Sections 4.6, 6, 8]{Kon}。
対応する正規化star積
\[
 f*g=fg+\frac h2\{f,g\}+O(h^2)
\]
の各グラフでは、$n$ 個の双ベクトル頂点が係数の次数を $2n$ 増し、
$2n$ 本の辺が微分で同じ次数を減らす。
従って積は次数を保つ。双ベクトルには $u,v$ 方向の微分がないため、
各双微分作用素は $x,y,z$ のみを微分する。
単位の正規化から $u*f=uf=f*u$、$v*f=vf=f*v$ が成り立つ。

twistした $L_\infty$ 写像の接写像は、Poisson複体からHochschild複体への鎖写像である。
$[\pi,F]=0$ なので
\[
 \widehat F=\mathcal U^{h\pi}_1(F)
 =\sum_{n\ge0}\frac{h^n}{n!}\mathcal U_{n+1}(\pi,\ldots,\pi,F)
\]
はstar積について中心的である。係数は三次斉次式である。
実際、$F$ 頂点を加えたグラフにも辺は $2n$ 本あり、
入力の多項式次数は $2n+3$ だから、出力は次数3となる。

最後に、各次数 $d$ の乗法写像
\[
 B_{h,d-3}\xrightarrow{\ *\widehat F\ }B_{h,d}
\]
は有限自由 $\CC[[h]]$-加群の間の写像であり、$h=0$ で非零多項式 $F$ による単射となる。
従って最大階数のある小行列式が単元であり、分裂単射で余核は自由である。
余核の階数は通常の三次超曲面と一致する。
\end{proof}

ここで完備化は各次数で行う。命題は形式級数としての構成であり、
$h\in\CC^\times$ への収束・代数化を与えてはいない。
また、商で使うのは一般に補正された $\widehat F$ である。
可換な方程式 $F$ をそのまま量子化代数で割ればよい、とは限らない。
次節は、この補正を有限な式で全て書ける例である。

\section{Fermat cubicの明示的な代数的非可換族}

\subsection{定義と中心元}

$h$ を次数0の中心パラメータとし、
\begin{equation}\label{eq:skly}
 \mathscr S=
 \frac{\CC[h]\langle x,y,z\rangle}
 {(xy-yx-hz^2,\ yz-zy-hx^2,\ zx-xz-hy^2)}
\end{equation}
とおく。$x,y,z$ の次数は1とする。
\[
 D=x^3+y^3+z^3,\qquad
 T=xyz+xzy+yxz+yzx+zxy+zyx
\]
と書く。

\begin{theorem}[次数3の中心元の明示式]\label{thm:center}
$\mathscr S$ において
\begin{equation}\label{eq:center}
 G_h=D-\frac{h^2}{3}T
\end{equation}
は中心的である。等価な表示は
\begin{equation}\label{eq:centerordered}
 G_h=(1-h^3)(x^3+z^3)+(1+h^3)y^3-2h^2zyx
\end{equation}
である。この恒等式は全ての $h$ に対して成り立つ。
\end{theorem}

Sklyanin型代数の中心と三次ポテンシャルについては
Etingof--Ginzburg\cite{EG}の体系的な研究がある。
ここでは規約の符号に依存しないよう、上の表示から直接証明する。

\begin{proof}
$C=y^3-z^3$ とおく。\eqref{eq:skly} から
\[
 [yz,x]=[zy,x]=hC
\]
であり、さらに
\begin{align*}
 [D,x]
 &=h(zy^2z-yz^2y)\\
 &=h[zy,yz]
 =h^2[yz,x^2]
 =h^3(Cx+xC).
\end{align*}
6個の単語を $zyx$ の順に整理すると
\[
 T=6zyx+3h(x^3-C).
\]
従って
\[
 [T,x]=6hCx-3h[C,x]=3h(Cx+xC).
\]
ゆえに $[D-h^2T/3,x]=0$。
$D,T$ と関係式は巡回置換で不変なので、$y,z$ との交換子も零である。
上の $T$ の表示を代入すれば \eqref{eq:centerordered} を得る。
$h$ で割る操作は用いていない。
\end{proof}

この計算は、元の三次式 $D$ をそのまま使う問題点も示している。
$[D,x]=h^3(Cx+xC)$ であり、一般に $D$ 自身は中心元ではない。
中心性を保証する $-h^2T/3$ の補正が本質的である。

\subsection{非可換cubic代数族と平坦性}

次数1の中心変数 $u,v$ を加え、
\begin{equation}\label{eq:family}
 \boxed{\quad
 \mathscr A=\mathscr S[u,v]/(G_h+u^3+v^3).
 \quad}
\end{equation}
$S_\lambda=\mathscr S/(h-\lambda)$、
$A_\lambda=\mathscr A/(h-\lambda)$ と書く。

\begin{theorem}[全パラメータで平坦な次数付き族]\label{thm:family}
$\mathscr A$ は各次数で有限自由な $\CC[h]$-加群であり、
\[
 A_0\simeq
 \CC[x,y,z,u,v]/(x^3+y^3+z^3+u^3+v^3).
\]
全ての $\lambda\in\CC$ に対して
\[
 H_{A_\lambda}(t)=\frac{1-t^3}{(1-t)^5}.
\]
また $\lambda\ne0$ なら $A_\lambda$ は非可換である。
\end{theorem}

\begin{proof}
$S_\lambda$ のSklyaninパラメータは
\[
 (p,q,r)=(1,-1,-\lambda)
\]
である。少なくとも二つのパラメータは非零で、
$p^3=1$ と $q^3=-1$ は異なる。
従ってArtin--Tate--Van den Berghの定理、または
Iyudu--Shkarinの代数的証明\cite{IS}により、
全ての $\lambda$ で $S_\lambda$ は3次元AS正則・Koszulで、
$H_{S_\lambda}(t)=(1-t)^{-3}$ である。

$\mathscr S_n$ は有限生成 $\CC[h]$-加群で、その各閉ファイバーの次元は一定である。
PID上の有限生成加群の構造定理から、ねじれ部分は零となり、$\mathscr S_n$ は自由である。
次に $u^3+G_h+v^3$ は中心変数 $u$ についてmonicであるため、割り算により
\[
 \mathscr A\simeq
 \mathscr S[v]\oplus\mathscr S[v]u\oplus\mathscr S[v]u^2
\]
という加群の直和分解を得る。従って次数ごとの自由性と
\[
 H_{A_\lambda}(t)
 =\frac{1+t+t^2}{(1-t)^4}
 =\frac{1-t^3}{(1-t)^5}
\]
が従う。

$\lambda\ne0$ では $[x,y]=\lambda z^2$ である。
$S_\lambda$ の次数2で $z^2\ne0$ は関係式から直接分かり、
次数3の斉次元で割っても次数2の成分は変わらない。
従って $A_\lambda$ は非可換である。
\end{proof}

$B_\lambda=S_\lambda[u,v]$ はAS正則次元5・Gorensteinパラメータ5である。
$G_\lambda+u^3+v^3$ は中心的な非零因子なので、change-of-ringsの議論
\cite[Lemma 1.2]{HeYe}から
$A_\lambda$ はAS Gorenstein次元4・パラメータ2となる：
\[
 \Ext^i_{A_\lambda}(\CC,A_\lambda)=
 \begin{cases}
 \CC(2)&i=4,\\
 0&i\ne4,
 \end{cases}
\]
右加群についても同様である。
この性質とHilbert級数は、次数付き超曲面モデルとしてcubic threefoldに適合する。
一方、$\qgr A_\lambda$ の有限大域次元やdg滑らかさには別の議論が必要である。
前者は次節で開集合 $\lambda^6\ne-27$ 上について証明する。

\subsection{半古典極限と点加群の三次曲線}

一次交換子は
\[
 \{x,y\}=z^2,\qquad\{y,z\}=x^2,\qquad\{z,x\}=y^2,
 \qquad \{u,-\}=\{v,-\}=0
\]
を与える。これはFermat式の、$u,v$ に相対的なJacobian括弧の $1/3$ 倍である。
射影チャート $u\ne0$ では $v/u$ がCasimirとなり、
線形束 $[u:v]$ がPoisson葉の幾何を支配する。

$S_\lambda$ の関係式を多重線形化すると、
点 $(x:y:z)$ と次の点 $(x':y':z')$ の間に
\[
 M_\lambda
 \begin{pmatrix}x'\\y'\\z'\end{pmatrix}=0,\qquad
 M_\lambda=
 \begin{pmatrix}
 -y&x&-\lambda z\\
 -\lambda x&-z&y\\
 z&-\lambda y&-x
 \end{pmatrix}
\]
を得る。直接計算により
\[
 \det M_\lambda=-\lambda\bigl(x^3+y^3+z^3+\lambda^2xyz\bigr).
\]
$\lambda\ne0$ で現れるHesse曲線は
\begin{equation}\label{eq:elliptic}
 E_\lambda:\quad x^3+y^3+z^3+\lambda^2xyz=0.
\end{equation}
その三つの偏微分に共通射影零点があるための条件は
$\lambda^6=-27$ である。従って
$\lambda\ne0,\ \lambda^6\ne-27$ でこの曲線は滑らかである。
$\lambda=0$ では $S_0$ 自体の点スキームは $\PP^2$ なので、
極限の三次曲線と点スキームを同一視してはいけない。

曲線の式に最初の変化が $\lambda^2$ で現れることは、
一次情報だけでは捉えられない幾何の一例である。
ただし、これだけで $\Ku(A_\lambda)$ の変形を計算したことにはならない。

\subsection{標準中心元との一致}

Walton--Wang--Yakimov\cite[Lemma 2.19]{WWY}に記載された標準的なSklyanin中心元に
$(a,b,c)=(1,-1,-\lambda)$ を代入すると、
\[
 g=(\lambda^4-\lambda)y^3+(1+\lambda^3)yxz
       +(\lambda^3-1)xyz-(\lambda+\lambda^4)x^3.
\]
関係式を用いて整理すると厳密に $g=-\lambda G_\lambda$ となる。
従って $\lambda\ne0,\ \lambda^6\ne-27$ で
\[
 S_\lambda/(G_\lambda)\simeq B(E_\lambda,\mathcal L,\sigma),
 \qquad
 \qgr(S_\lambda/(G_\lambda))\simeq\operatorname{coh}E_\lambda.
\]
$B$ はtwisted homogeneous coordinate ringである。
$\lambda=0$ では通常のFermat平面三次曲線の斉次座標環となる。
従って $\lambda^6\ne-27$ の全域で、中心商のtailsは滑らかな楕円曲線である。

\section{Fermat族の射影的正則性：大域次元3}

\begin{theorem}\label{thm:smooth}
$\lambda^6\ne-27$ とする。このとき
\[
 \operatorname{gldim}(\qgr A_\lambda)=3.
\]
ここで大域次元はcoherent対象間のExtの消滅で測る。
従って $A_\lambda$ は、tailsの有限大域次元を定義とする意味で
非可換graded isolated singularityである。
\end{theorem}

以下の証明は、matrix factorizationの定理と非可換Serre双対を組み合わせるものである。
この結論を、そのままperfect diagonalの意味でのdg滑らかさと言い換えない。

\begin{proof}
$S=S_\lambda,\ g=G_\lambda,\ A=A_\lambda,\ C=S/(g)$ と略記する。
$S$ は左右noetherianな3次元AS正則代数である。
滑らかな楕円曲線の場合のこの性質と中心商の記述はSklyaninの標準的結果
\cite[Section 2.3]{WWY}であり、$\lambda=0$ は可換多項式環である。

\textbf{双対性の準備。}
三次tensor
\[
 w=xyz+yzx+zxy-(xzy+zyx+yxz)-\lambda(x^3+y^3+z^3)
\]
は巡回不変で、その左偏微分は $S$ の三つの関係式である。
AS正則性とsuperpotentialの判定\cite[Lemma 2.15]{IM}から $S$ は3-Calabi--Yauである。
中心多項式拡大 $B=S[u,v]$ と中心正則元 $W=g+u^3+v^3$ による商を考えると、
$A$ のbalanced dualizing complexは
\[
 A(-2)[4]
\]
となる。従って $\chi$、有限局所コホモロジー次元、および
perfect対象上のSerre双対が使える\cite[Theorem 4.2, Corollary 4.3]{YZ}。
ここでは有限大域次元を先に仮定していない。

\textbf{特異圏の消滅。}
$\pi_A:\operatorname{gr}A\to\qgr A$ を商関手とし、
\[
 \mathcal P=\operatorname{thick}\{\pi_AA(n)\mid n\in\mathbb Z\},
 \qquad \mathcal Q=D^b(\qgr A)/\mathcal P
\]
とおく。$\mathcal Q=0$ を示す。

任意のcoherent対象 $M=\pi_A N$ の有限生成graded lift $N$ を選ぶ。
$B$ はnoetherian AS正則であり、$W$ は中心正則元なので、
Cassidy--Conner--Kirkman--Moore\cite[Theorems 4.7, 5.5]{CCKM}により
$\Dsg^{\mathrm{gr}}(A)$ の対象を $W$ のmatrix factorizationで表せる。
その奇微分を $d$ と書くと $d^2=W$ であり、
中心変数による形式微分から
\[
 d(\partial_ud)+(\partial_ud)d=3u^2,\qquad
 d(\partial_vd)+(\partial_vd)d=3v^2.
\]
従って特異圏では $u^2,v^2$ の作用が零となる。
$D^b(\operatorname{gr}A)\to\mathcal Q$ はgraded perfect複体を零にするので、
この零化は $\mathcal Q$ にも降りる。
導来商関手の本質的全射性を仮定する必要はなく、選んだ実際のlift $N$ に適用すればよい。

$N$ に対する二重Koszul複体 $K(u^2,v^2;N)$ のコホモロジーは
$(u^2,v^2)$ に殺される。
$I=(u,v)$ とすれば $I^3\subset(u^2,v^2)$ なので、
各コホモロジーは長さ高々3の $I$-filtrationを持ち、その商は $C$-加群となる。

一方、$D^b(\qgr C)=D^b(\operatorname{coh}E_\lambda)$ は $\pi_CC(n)$ で厚く生成される。
これはample sequenceの生成性からも分かるが、
大域次元1と二段の有限表示を使えば直接に示せる。
二項複体の負次数コホモロジーから零次コホモロジーへの連結写像は
$\Ext^2$ に入り零なので、零次コホモロジーがその二項複体の直和因子となるからである。

$u,v$ は $A$ 上の中心正則列である。
$u$ の非零因子性はmonic表示から、$A/(u)$ 上の $v$ の非零因子性は $g$ の正則性から従う。
従ってKoszul分解により $\pi_A C(n)\in\mathcal P$。
前段のfiltrationとexactな制限関手から
\[
 \pi_AK(u^2,v^2;N)\in\mathcal P
\]
を得る。

しかし $\mathcal Q$ では $u^2,v^2$ の作用は零なので、この二重coneは
$M$ の四つのshift・twistの直和になる。
その直和が零である以上、$M=0$ である。
全てのcoherent対象について成立するので $D^b(\qgr A)=\mathcal P$。

\textbf{一様な大域次元の上界。}
双対性を全てのperfect対象へ拡張する正確な結果として、
Chen\cite[Theorem 3.10]{Chen}のgeneralized Serre functor
$-\otimes_A^{\mathbf L}R_A[-1]$ を用いる。
$\pi_A A(n)$ はコホモロジー次元3から有限射影次元を持つので適用でき、
有限coneと直和因子を通して $\mathcal P$ 全体に拡張する。
いま全てのcoherent対象が $\mathcal P$ に入るので
\[
 \Ext^i(M,N)^*\simeq\Ext^{3-i}(N,M(-2)).
\]
$i>3$ なら右辺は零である。
さらに $\Ext^3(\OO,\OO(-2))\simeq\CC$ なので、大域次元はちょうど3となる。
\end{proof}

$\lambda^6=-27$ の6点については、この証明は適用しない。
第\ref{sec:exceptional}節では非perfectな行列因子化を構成して、
これらのファイバーの大域次元が無限であることを示す。

\section{Kuznetsov成分の形式Morita変形と固定ファイバー}

\subsection{Hochschild変形のhomotopy abelian性}

滑らかなcubic threefoldでは
\[
 D^b(X)=\langle\Ku(X),\OO_X,\OO_X(1)\rangle.
\]
Kuznetsovの計算\cite[Theorem 8.9]{Kuz}は
\[
 \HH^i(\Ku(X))=
 \begin{cases}
 \CC&i=0,\\
 \CC^{10}&i=2,\\
 0&i\ne0,2
 \end{cases}
\]
を与える。また自然な幾何学的写像
\[
 H^1(T_X)\longrightarrow\HH^2(\Ku(X))
\]
は同型である\cite[Theorem 4.4(2)]{JLLZ}。

このコホモロジーの次数だけで、Hochschild dg Lie代数
$L=CC^\bullet(\Ku(X))[1]$ のhomotopy abelian性が分かる。
その最小 $L_\infty$ モデルの台は次数 $-1$ と $1$ のみである。
$n$ 項演算 $\ell_n$ の次数は $2-n$ なので、
次数 $-1$ の入力が $m$ 個あると、出力次数は
\[
 (n-2m)+(2-n)=2-2m
\]
となる。これは偶数であり、台には存在しない。
従って全ての高次括弧が零である。

よってHochschild/Maurer--Cartan形式変形問題の古典的基底は
$\Spf\CC[[s_1,\ldots,s_{10}]]$ である。
$\HH^0$ に由来する高次自己同型は、集合値の基底には表れていない。
これだけでuncurved Morita変形との比較が済むわけではない。
Keller--Lowen\cite{KL}の曲率に関する注意を踏まえ、以下では実在する
幾何学的変形を使って比較を証明する。

\subsection{実在するMorita変形の形式基底}

$\mathcal C=\Ind(\Ku(X))$ とし、$\CatDef^c_{\mathcal C}$ を
compactly generatedな線形圏の標識付き変形空間とする。
以後は可換局所Artin $\CC$-代数に制限し、$\pi_0$ で同値類を取る。

\begin{theorem}[古典的形式categorical Torelli]\label{thm:formal-torelli}
自然な写像
\[
 \Def_X\longrightarrow\pi_0\CatDef^c_{\mathcal C}
\]
は同型である。従って右辺は滑らかな10次元の形式空間で表される。
\end{theorem}

\begin{proof}
$\mathcal C$ はproperなcompact generatorを持つのでtamely compactly generatedである。
Blanc--Katzarkov--Pandit\cite[Section 3.3]{BKP}が記述するLurieの比較定理により
\[
 \theta:\CatDef^c_{\mathcal C}\longrightarrow\CatDef^\wedge_{\mathcal C}
\]
は $(-1)$-truncatedである。特に各Artin基底で $\pi_0\theta$ は単射であり、
右辺の形式完備化はHochschild複体で制御される。
前小節から $\pi_0\CatDef^\wedge_{\mathcal C}$ は滑らかな10次元形式空間である。

通常の三次三様体の変形に相対例外対の直交補を対応させる写像を考える。
その合成
\[
 \Def_X\longrightarrow\pi_0\CatDef^c_{\mathcal C}
 \longrightarrow\pi_0\CatDef^\wedge_{\mathcal C}
\]
の微分は $H^1(T_X)\xrightarrow{\sim}\HH^2(\Ku(X))$ である。
両端は滑らかな10次元の表現可能な形式変形関手なので、形式逆関数定理から合成は同型。
従って $\pi_0\theta$ は全射でもあり、結論を得る。
\end{proof}

ここで使う文献の入力は、全Artin基底に対する $\theta$ の単射性である。
\cite[Theorem 4.27]{BKP}の一変数形式族に対する全射性だけから、
上の多変数の主張を結論しているわけではない。
また、圏の変形の高次空間まで離散的な $\Def_X$ と同一視する主張ではない。

\subsection{非可換層から残余圏への自然な写像}

構造層の結合代数変形から、Lowen\cite[Theorem 3.26]{LowenPre}により
$\operatorname{Qch}(X)$ の平坦なGrothendieck圏の変形 $\mathcal D_A$ を得る。
\[
 \Ext^1(\OO_X(i),\OO_X(i))=
 \Ext^2(\OO_X(i),\OO_X(i))=0
\]
なので各 $\OO_X(i)$ は一意な同型類として持ち上がる。持ち上げを $E_i$ とする。
$G=\bigoplus_{i=0}^4\OO_X(i)$ は $\PP^4$ のBeilinson生成子の制限であり、
随伴と閉埋め込みの押し出しの保守性によりcompact generatorである。
De Deken--Lowen\cite[Propositions 5.8--5.10]{DDL}から
$G_A=\bigoplus_{i=0}^4E_i$ もcompact generatorとなる。

Derived base changeとNakayamaから $(E_0,E_1)$ は相対例外対である。
その右直交成分は $P_A(G_A)$ で生成される。ただし
$P_A=L_{E_0}L_{E_1}$ である。
\cite[Lemma 5.14]{DDL}からmutationは基底変換と可換なので、
\[
 \RHom_{\mathcal D_A}(P_A(G_A),P_A(G_A))
\]
は残余圏の実在するMorita変形を与える。生成子の選択を変えてもMorita同値類は変わらない。
従ってArtin基底上の自然な写像
\[
 \Phi:\NC_X\longrightarrow\pi_0\CatDef^c_{\mathcal C}
\]
が構成される。この議論で $\OO_X(i)$ がfull helixであるとは仮定しない。

\begin{lemma}\label{lem:bivector-zero}
接写像 $\HH^2(X)\to\HH^2(\Ku(X))$ は、通常のHKR分解の
$H^0(\Lambda^2T_X)$ を零に写す。
\end{lemma}
\begin{proof}
JLLZ\cite[Theorems 3.10, 4.4(2)]{JLLZ}のTodd補正付きHKRの加群構造と
$\gamma:\HH^2(\Ku(X))\to\operatorname{Hom}(H^{2,1},H^{1,2})$ の単射性を使う。
通常のHKRとの差も確認する。局所体積形式 $\mu_i$ を選び、
$\iota_{D_i}\mu_i=d(\iota_\pi\mu_i)$ によりベクトル場 $D_i$ を定める。
$\mu_j=f_{ij}\mu_i$ なら $D_j-D_i=\pm\pi^\sharp d\log f_{ij}$ となる。
従って $\pi\lrcorner c_1(K_X)$ は $H^1(T_X)$ で零である。
残りのTodd補正は $H^2(\OO_X)=0$ に入るので、双ベクトルに対して両HKRは一致する。
その $H^{2,1}=H^1(\Omega_X^2)$ への作用は
$H^1(\Omega_X^2)\to H^1(\OO_X)=0$ であり、$\gamma$ の単射性から結論する。
この局所計算にPoisson条件は不要である。
\end{proof}

\subsection{残余圏を固定した形式ファイバー}

\begin{theorem}[固定した残余圏の形式ファイバー]\label{thm:fibre}
上で構成した $\Phi$ の、未変形の残余圏の上の古典的形式ファイバーは非標準的に
\[
 \Spf\CC[[a_{ij}]]/(q_{ijkl})
\]
と同型である。特に次元は7である。
\end{theorem}

\begin{proof}
環準同型を $s_i\mapsto f_i(t,a)$ と書く。
$\det(\partial f_i/\partial t_j)(0)\ne0$ なので、
$t'_i=f_i(t,a),\ a'_{ij}=a_{ij}$ は形式的な座標変換を与える。
$a$ を固定するためPlückerイデアルも保存する。
新しい座標で $\Phi$ は $t'$ への射影となり、結論を得る。
\end{proof}

$\Phi$ の接写像が双ベクトル方向を殺すことは、
任意の純Poisson族で残余圏が全次数で一定であることを意味しない。
固定残余圏の条件は一般に幾何学的パラメータ $t$ の高次補正を必要とし得る。
幾何学的sliceの座標を用いると、この補正は $t(a)=O(a^2)$ から始まる。

\subsection{明示的モデルの残余圏}

$\lambda^6\ne-27$ とする。\eqref{eq:family} のAS Gorenstein代数に
Orlovの定理\cite[Theorem 2.5]{Orlov}を適用すると半直交分解
\[
 D^b(\qgr A_\lambda)
 =\langle\mathcal K_\lambda,\OO,\OO(1)\rangle,\qquad
 \mathcal K_\lambda\simeq\Dsg^{\mathrm{gr}}(A_\lambda)
\]
を得る。埋め込みの選択によりshift・mutationの規約は変わる。
第6節の孤立特異点性と $\omega_A=A(-2)$ に、
Ueyamaのstable Cohen--Macaulay圏のSerre公式\cite{Ueyama}を適用すると、
残余圏のSerre関手は
\[
 S_{\mathcal K_\lambda}=(-2)[3]
\]
となる。ここで括弧内は特異圏側のgrade shiftを表す。
中心的な三次斉次元のmatrix factorizationの周期性 $(3)\simeq[2]$ から
\[
 S_{\mathcal K_\lambda}=(1)[1],\qquad
 \boxed{S_{\mathcal K_\lambda}^{\,3}=[5]}
\]
を得る。従って通常のcubic threefoldと同じ分数Calabi--Yau次元 $5/3$ を持つ。
grade shiftは $S_{\mathcal K_\lambda}[-1]$ で既に決まるので、
これだけを追加しても、全体圏への埋め込みやgluingの情報を補えない。
また、この関手等式だけから $\mathcal K_\lambda$ がある可換cubicの
Kuznetsov成分と同値であることは従わない。次節では別の具体的な比較により同定する。

\section{Fermat族の残余圏とHesse族の比較}

\subsection{主張と、二次類の規約}

\[
 Y_s=V(x^3+y^3+z^3+u^3+v^3+sxyz)\subset\PP^4
\]
とおく。偏微分の計算から、$Y_s$ は $s^3\ne-27$ のとき滑らかである。
以下のdg圏は常にpretriangulated hullとidempotent completionを取った
Morita圏として扱う。

\begin{theorem}[Hesse比較]\label{thm:hesse}
$\lambda^6\ne-27$ なら、dg Morita同値
\[
 \mathcal K_\lambda\simeq\Ku(Y_{\lambda^2})
\]
が存在する。また $R=\CC[[h]]$ 上では、この同値を
$h=0$ の通常のOrlov同一視に還元する同値として構成できる。
従って、$X_0=Y_0$ のJacobian環による同一視
\[
 \HH^2(\Ku(X_0))\simeq
 \left(\frac{\CC[x,y,z,u,v]}{(x^2,y^2,z^2,u^2,v^2)}\right)_3
\]
において、最初の非零の変形係数は
\[
 \boxed{\kappa_2=[xyz]\ne0.}
\]
ここで $\kappa_2$ は $h^2$ の係数を意味する。
二階微分そのものを使う規約では $2[xyz]$ となる。
\end{theorem}

証明は楕円曲線の導来圏だけでなく、grade shiftと $(3)\simeq[2]$ を比較する。
点曲線の式だけからこの命題を結論することはできない。

\subsection{楕円曲線のblockと、次数シフト}

まず体上で $C=S_\lambda/(G_\lambda)$ とする。
前述のtwisted homogeneous coordinate ringの同一視と、
Etingof--Ginzburg\cite[Theorem 3.6.2]{EG}、Orlovの定理から
\[
 \MF^{\mathbb Z}(S_\lambda,G_\lambda)
 \simeq\Dsg^{\mathrm{gr}}(C)\simeq D^b(E_\lambda)
\]
を得る。$\lambda=0$ では通常の平面三次曲線のOrlov同値である。
この同値は以下のwindow構成によりdg関手で実現する。

一般に $C=B(E,L,\sigma)$、$\deg L=3$、$\sigma=t_p$ を楕円曲線の平行移動とする。
右加群の規約
\[
 \Gamma_*(F)_n=H^0(E,F\otimes L_n),\qquad
 L_n=L\otimes\sigma^*L\otimes\cdots\otimes(\sigma^{n-1})^*L
\]
を採用すると、tailsの次数シフトは
\[
 H(F)=\sigma_*(F\otimes L)
\]
である。特異圏側の次数シフトを $\rho$ と書く。
適切なOrlov windowでは
\begin{equation}\label{eq:rotation}
 \rho=T_{\OO_E}\circ H.
\end{equation}

この公式は非可換の場合にもwindowから直接に従う。
$q:D^b(\operatorname{gr}C)\to\Dsg^{\mathrm{gr}}(C)$ とし、
$\Psi_i=qR\Gamma_{\ge i}$ とおく。Gorenstein parameterは零なので各 $\Psi_i$ は同値である。
内部次数の切断から $(1)\Psi_1=\Psi_0H$、
$q(R)=\Psi_0(\OO_E)[1]$、および三角
\[
 \Psi_1F\longrightarrow\Psi_0F\longrightarrow
 \RHom(\OO_E,F)\otimes q(R)
\]
を得る。楕円曲線のSerre双対を使うと
$\Psi_1=\Psi_0T_{\OO_E}^{-1}$ となり、\eqref{eq:rotation}を得る。
体の場合は $R$ をその体と読む。

\subsection{平行移動の消去を明示する}

$\rho_0=T_{\OO_E}\circ(-\otimes L)$ を可換な平面三次曲線のgrade shiftとする。
平行移動は $\OO_E$ を保存するので、\eqref{eq:rotation}から
\[
 \rho_{\NC}=t_{-p}^*\rho_0.
\]
原点を $o\in E$ とし $P_q=\OO_E(q-o)$、$A_q=(-\otimes P_q)$ とおく。
楕円曲線の対角はCartier因子なので、spherical twistのFourier--Mukai核は
\[
 K_{T_{\OO_E}}=I_\Delta[1],\qquad
 K_{T_{P_q}}=(P_q^{-1}\boxtimes P_q)\otimes I_\Delta[1].
\]
Theorem of the square、または両射影の各ファイバーへの制限から
\[
 (P_q^{-1}\boxtimes P_q)\otimes I_\Delta
 \simeq I_{\{y=x-q\}}
\]
となる。従って $T_{P_q}=t_q^*T_{\OO_E}$ であり、
\begin{equation}\label{eq:conjugacy}
 A_q\rho_0A_q^{-1}=t_q^*\rho_0,
 \qquad \rho_{\NC}=A_{-p}\rho_0A_{-p}^{-1}.
\end{equation}
従って平行移動のパラメータは、次数零の直線束による共役で取り除ける。

この比較では三次周期性 $\rho^3\simeq[2]$ も保存する。
二つの周期性同型の差は $\HH^0(D^b(E))=R$ の単元である。
体 $\CC$ 上ではその三乗根、形式基底上では特殊ファイバーで1となる三乗根を取り、
共役の同型を正規化する。高次のcoherenceの差も、各小拡大の核 $I$ に対する
$H^{>0}(\mathbb Z/3,I)=0$ により帰納的に消去できる。

\subsection{正則DVR上の比較とdg基底変換}

形式的な二次類を読むには、各複素ファイバーでの同値だけでは不足する。
そこで $R=\CC[[h]]$ 上で構成する。ここでは
\[
 E_R=V(x^3+y^3+z^3+h^2xyz)\subset\PP^2_R
\]
は滑らかな楕円曲線である。$S_R$ の関係行列は $E_R$ 上で階数2となる。
その核は $\sigma_R:E_R\to E_R$ を定め、転置行列から逆写像を得る。
$\sigma_0=\id$ であり、$\sigma_R$ は平行移動である。
例えばSklyaninの標準的規約では $p_h=[1:-1:-h]$、原点は $[1:-1:0]$ である。
符号の逆の規約を用いても \eqref{eq:conjugacy} の結論は同じである。

関係行列から $S_R\to B(E_R,L,\sigma_R)$ を得る。
$G_h$ は一般ファイバーで零に写るので、各次数の $R$-torsion freenessから全体で零に写る。
特殊ファイバーは通常の三次曲線の斉次座標環である。
次数 $n>0$ の両辺は階数 $3n$ の有限自由 $R$-加群なので、Nakayamaにより
\begin{equation}\label{eq:relative-thcr}
 C_R=S_R/(G_h)\simeq B(E_R,L,\sigma_R).
\end{equation}

相対Koszul複体とその双対は、各内部次数で有限自由 $R$-複体である。
$h=0$ で多項式環のKoszul複体になるので完全であり、
\[
 \RHom_{S_R}(R,S_R)=R(3)[-3],\qquad
 \RHom_{C_R}(R,C_R)=R[-2]
\]
を得る。$\operatorname{gldim}R=1$ と相対bimodule分解から
$S_R$ のgraded global dimensionは4以下、$C_R$ の両側graded injective dimensionは有限となる。

Noetherian性も相対的に確認できる。$\operatorname{gr}_hS_R=\CC[x,y,z][h]$ の
斉次イデアルの有限な生成元を持ち上げる。固定した内部次数での係数の修正は
有限自由 $R$-加群内で収束するので、任意の斉次イデアルは $S_R$ 自身で有限生成である。
非斉次イデアルには最高内部次数のイデアルを使い、次数を下げる有限回の帰納法を適用する。
これで通常の両側Noetherian性を得る。環全体が $h$-進完備であることは仮定しない。

従ってBurke--Stevenson\cite[Definition 5.1, Lemma 6.2, Theorem 6.4(3)]{BS}を
正則DVR $R$ 上で適用でき、$a=0$ のOrlov windowを得る。
非可換MFとの比較は次のように確認する。
$C_R$ のGorenstein-projective加群 $M$ についてchange of ringsから
$\Ext^{>1}_{S_R}(M,S_R)=0$ となる。
graded localな $S_R$ の有限最小分解を双対化すると
$\operatorname{pd}_{S_R}M\le1$ なので、
\[
 0\to P_1\xrightarrow\varphi P_0\to M\to0
\]
がある。$G_hM=0$ より $G_h:P_0\to P_0(3)$ を $\psi:P_0\to P_1(3)$ に持ち上げれば
$\varphi\psi=G_h$、$\psi\varphi=G_h$ となる。
分解の比較は安定射とMFのホモトピー射を同一視する。
これは\cite{CCKM}の議論をgraded localなDVR上で行ったものである。

dg構造と基底変換についても構成を固定する。
複体のdg圏、perfect複体によるdg quotient、complete resolutionを用いれば、
上のcokernel・window関手はdg準関手となる。
windowの逆関手 $qR\Gamma_{\ge i}$ は、内部次数の切断とderived global sectionからなる。
\eqref{eq:relative-thcr}の下で各成分は
$R\Gamma(E_R,F\otimes L_n)$ であり、$E_R\to\operatorname{Spec}R$ がsmooth properで
$F$ がperfectなので、任意のderived base changeと可換である。
また $C_R$ は $R$ 上flatで $\operatorname{gldim}R=1$ なので、
$-\otimes_R^{\mathbf L}\mathbb C$ は $D^b(\operatorname{gr}C_R)$ を
$D^b(\operatorname{gr}C_0)$ に送り、perfect複体をperfect複体に送る。
従って特異圏の間の特殊化関手が定まる。
内部次数の切断は成分ごとの操作であり、この関手によってwindowを比較できる。
両windowが逆であることは三角と随伴射で実現され、その射も同時に基底変換する。
従って
\[
 \MF^{\mathbb Z}(S_R,G_h)\simeq\Perf(E_R)
\]
は $h=0$ のOrlov同値と両立するdg Morita同値である。
Artin商へはこのdg圏をderived base changeする。
特異なArtin基底の\emph{絶対}特異圏を直接使うことはしない。
より具体的には、$\Perf(E_R)$ の有限個のample twistから生成子を選び、
そのwindow像の自己準同型dg代数を取ればよい。
その特殊化は $\Perf(E_0)$ の生成子の自己準同型代数となるので、
基底変換後のMFのMorita圏に生成子が欠落することもない。

相対版の \eqref{eq:conjugacy} では $p_0=o$ なので $A_{-p_0}=\id$。
従って標識を保存して、$\MF(S_R,G_h)$ を
$\MF(R[x,y,z],x^3+y^3+z^3+h^2xyz)$ と、次数シフト込みで比較できる。

\subsection{中心変数を加えるThom--Sebastiani型補題}

次に、必要な非可換テンソル積の主張を証明する。
可換の場合の対応する定理はBallard--Favero--Katzarkov\cite[Theorem 5.15, Corollary 5.18]{BFK}
だが、その定理の仮定をそのまま非可換代数へ拡張して引用することはしない。

$S$ を上の代数、$g$ を中心三次元とし、$u,v$ を中心変数とする。
\[
 \Gamma=\mathbb Z^3/\langle3e_1-3e_2,3e_1-3e_3\rangle,
 \qquad w=3e_1=3e_2=3e_3
\]
で次数を付け、$x,y,z$ の次数を $e_1$、$u$ を $e_2$、$v$ を $e_3$ とする。

\begin{lemma}\label{lem:ts}
外部テンソル積はdg Morita同値
\[
 \MF^{\mathbb Z}(S,g)\otimes_R^{\mathbf L}
 \MF^{\mathbb Z}(R[u],u^3)\otimes_R^{\mathbf L}
 \MF^{\mathbb Z}(R[v],v^3)
 \simeq\MF^\Gamma(S[u,v],g+u^3+v^3)
\]
を与える。$R=\CC$ または上の正則DVRとする。
\end{lemma}

\begin{proof}
有限自由MFの奇微分の外部テンソル和は、Koszul符号を付けると
$(d_1+d_2+d_3)^2=g+u^3+v^3$ を満たす。
完全忠実性はHom複体の直接計算である。
次数 $n$ のHomは
\[
 \bigoplus_{i+j+k=n}
 \operatorname{Hom}^i(P,P')\otimes_R
 \operatorname{Hom}^j(Q,Q')\otimes_R
 \operatorname{Hom}^k(T,T')
\]
に分解し、微分もテンソル積微分になる。
三つの $[2]$ と三次内部シフトを同一視するのが $\Gamma$ の関係式である。
有限自由MFと非負次数の代数では、固定した $n$ に寄与する分解は有限なので
無限積・完備化の問題は生じない。
$u^3,v^3$ の因子は $\Perf(RA_2)$ という有限自由なtilting模型を持つので、
相対的なテンソル積はこの模型でderived tensorとして計算できる。

生成性を示す。$A=S[u,v]/(g+u^3+v^3)$、$C=S/(g)=A/(u,v)$ とする。
MFの微分を $u,v$ で微分すると $u^2,v^2$ の作用は特異圏で零となる。
従って任意の対象 $N$ は $K(u^2,v^2;N)$ のshift・twistの直和因子である。
このKoszul複体のコホモロジーは $(u^2,v^2)$ に殺され、
$(u,v)^3\subset(u^2,v^2)$ より有限filtrationの商は $C$-加群となる。

$u,v$ は $A$ の中心正則列なので、$C$-free加群の制限は $A$ 上perfectである。
従って制限は $\Dsg(C)\to\Dsg(A)$ に降りる。
$C$-加群 $M$ の安定化を $P$ とすると、その制限の安定化は
\begin{equation}\label{eq:stabilization}
 P\boxtimes(u,u^2)\boxtimes(v,v^2)
\end{equation}
となる。これは次の分解から直接に分かる。
$M$ の長さ1の $S$-分解を $S[u,v]$ に延長し、$u,v$ のKoszul分解とテンソルする。
その上で $g$ の零ホモトピーと、$u^3=u\cdot u^2$、$v^3=v\cdot v^2$ の
Koszul零ホモトピーを足して周期化すると、\eqref{eq:stabilization}の微分そのものになる。
異なる因子の奇作用は反可換なので追加の高次ホモトピーは不要である。
これは制限加群の自由分解の周期的末尾を与える。

一変数について行列で確認しておく。
$P=(P_1\xrightarrow\varphi P_0\xrightarrow\psi P_1(3))$、$M=\operatorname{coker}\varphi$
とすると、$uM=0$ とした $M$ の $S[u]$-分解は
\[
 P_1(-1)\xrightarrow{\binom{-u}{\varphi}}
 P_1\oplus P_0(-1)\xrightarrow{(\varphi\ \ u)}P_0.
\]
$g+u^3$ の零ホモトピーは $s_0=(\psi,u^2)^t$、$s_1=(-u^2,\psi)$ で、
$ds+sd=g+u^3$、$s^2=0$ を満たす。周期化した末尾の行列は
\[
 \alpha=\begin{pmatrix}\varphi&u\\-u^2&\psi\end{pmatrix},
 \qquad
 \beta=\begin{pmatrix}\psi&-u\\u^2&\varphi\end{pmatrix},
 \qquad \alpha\beta=\beta\alpha=(g+u^3)I.
\]
この周期化が $M$ を分解することも確認できる。
元の分解を $A=S[u]/(g+u^3)$ にテンソルした複体のホモロジーは
次数0の $M$ と次数1の $M(-3)$ のみである。
周期化を周期の個数で濾過すると、$s$ が隣接する二つの $M$ を同型で結ぶ。
$ds=(g+u^3)$ によるconnecting mapでこの同型が分かる。
従ってそれらは対ごとに消え、次数0の $M$ だけが残る。
各全次数で濾過は有限なので、このspectral sequenceは収束する。
これで\eqref{eq:stabilization}は単なる因子化の恒等式でなく、制限の安定化を表す。

以上から任意の対象が外部積の像で厚く生成される。
$\Gamma$-次数の $C$-加群は $\Gamma/\mathbb Ze_1$ の9個の剰余類に分かれるが、
第二・第三因子の次数シフトがこれらをすべて与える。
直和因子を加えれば主張を得る。
\end{proof}

\subsection{通常の整数次数への降下}

$\Gamma\to\mathbb Z$、$e_i\mapsto1$ の核は
$T\simeq(\mathbb Z/3)^2$ である。
$\Gamma$-次数付きMFは、通常の次数付きMFに $H=\operatorname{Hom}(T,\CC^*)$ の
同変構造を付けた圏である。具体的には $H$ は $u,v$ を独立に三乗根倍する。
双対群の作用で再び同変化すると元の圏を回復する
\cite[Theorem 4.2]{Elagin}。これは有限群の平均化によるdg構成でも成立し、
$|H|$ が可逆な $R$ 上でも同じ証明でよい。

補題\ref{lem:ts}の左辺では、この双対作用はgrade shiftの差
\[
 \rho_2\rho_1^{-1},\qquad\rho_3\rho_1^{-1}
\]
で与えられる。三乗が恒等になる同型は各因子の三次周期性から得る。
\eqref{eq:conjugacy}は第一因子の $\rho_1$ と周期性を保存するので、
テンソル積の同値はこの群作用を保存し、通常の整数次数のMF圏へ降りる。
これで
\[
 \MF^{\mathbb Z}(S_R[u,v],G_h+u^3+v^3)
 \simeq
 \MF^{\mathbb Z}(R[x,y,z,u,v],x^3+y^3+z^3+h^2xyz+u^3+v^3)
\]
を得る。双方にGorenstein parameter 2のOrlov分解を適用すれば
命題\ref{thm:hesse}の圏同値を得る。形式的には全構成が $h=0$ の同一視を保つ。
従って二次係数は $[xyz]$ であり、この単項式はFermatのJacobianイデアルに属さない。

\subsection{dg smoothnessへの帰結}

\begin{theorem}\label{thm:dg-smooth}
$\lambda^6\ne-27$ なら $D^b(\qgr A_\lambda)$ の標準的dg enhancementは
smooth properである。
\end{theorem}
\begin{proof}
命題\ref{thm:hesse}から $\mathcal K_\lambda$ はsmooth properであり、
他の二成分も $\Perf(\CC)$ なのでsmooth properである。
$\qgr A_\lambda$ は $\chi$、有限コホモロジー次元と第6節の大域次元3を持つので、
coherent対象間のtotal Extは有限次元である。
従って半直交分解のgluing bimoduleはpointwise properである。

smoothな圏上のpointwise proper加群はperfectである。
実際、perfectな対角bimoduleを有限個のrepresentable bimoduleから構成し、
その構成に $M\otimes^{\mathbf L}-$ を作用させればよい。
従って各gluing bimoduleはperfectであり、Orlovのgluing定理
\cite[Proposition 3.8, Theorem 3.25]{OrlovGlue}から全体がsmoothとなる。
propernessは既に確認した。
\end{proof}

これは全体圏が $D^b(Y_{\lambda^2})$ と同値であるという主張ではない。
上で同定したのは残余成分であり、例外対とのgluingは別のデータである。

\section{二次類の解釈と、残余圏を固定する明示的補正}

\subsection{対称性による独立の検査}

$\omega^3=1$、$\omega\ne1$ とする。
この族は $u\mapsto\omega u$ と $v\mapsto\omega v$ を持ち、
半線形な対称性 $(h,x)\mapsto(\omega h,\omega x)$、
$(h,x,y)\mapsto(-h,y,x)$ も持つ。
Fermat Jacobian環の次数3は5変数のsquare-freeな三次単項式10個を基底とする。
最初の二対称性の不変部分は $\CC[xyz]$ のみである。

幾何学的形式変形基底との標識付き同一視を取ると、対称性を保つこの族の像は
$Y_{s(h)}$ と書け、半線形対称性から
\[
 s(\omega h)=\omega^2s(h),\qquad s(-h)=s(h)
\]
となる。従って $s(h)=h^2f(h^6)$ であり、許される次数は $2,8,14,\ldots$ である。
対称性だけでは $f(0)$ を決定できない。
命題\ref{thm:hesse}が、通常のmarked Hesse座標で $s(h)=h^2$ を与える。

\subsection{中心式の可換化が二次類を与えない理由}

対称化した中心式は
\[
 G_h=x^3+y^3+z^3-\frac{h^2}{3}\sum_{\tau\in S_3}\tau(xyz)
\]
なので、文字を可換として並べれば二次項は $-2h^2xyz$ になる。
しかしこの操作は周囲の代数の乗法の変形を捨ててしまう。
残余圏の変形はその乗法も含むため、正しい係数は命題\ref{thm:hesse}の
$+h^2xyz$ である。

補助的な楕円曲線の $j$-不変量も計算できる。
\[
 j(E_s)=\frac{s^3(216-s^3)^3}{(s^3+27)^3},\qquad
 j(E_{h^2})=512h^6+O(h^{12}).
\]
これは $s=0$ の楕円曲線の自己同型による粗モジュライ写像の分岐を反映する。
$j$ の最初の変化が6次であることは、marked変形の2次類が零であることを意味しない。
添付コードはHesse式をWeierstrass式へ変換してこの有理式を厳密検証する。

\subsection{固定した残余圏を持つ形式族}

$t=(1-h^3)^{1/3}\in\CC[[h]]$ を $t(0)=1$ の分枝とし、
\[
 S'_h=\frac{\CC[[h]]\langle x,y,z\rangle}
 {(xy-tyx-hz^2,\ yz-tzy-hx^2,\ zx-txz-hy^2)}
\]
とする。そのSklyaninパラメータは $(a,b,c)=(1,-t,-h)$ であり、
\[
 a^3+b^3+c^3=1-t^3-h^3=0.
\]
従って一般ファイバーの点曲線は固定したFermat曲線
$E_0:x^3+y^3+z^3=0$ であり、中心切断に対応する相対楕円曲線も
$E_0\times\operatorname{Spf}\mathbb C[[h]]$ と取れる。
ただし $h=0$ で $S'_0$ 自身の点スキームは $\PP^2$ である。

標準中心元\cite[Lemma 2.19]{WWY}を正規化すると
\[
 G'_h=(1+h^3)(x^3+z^3)+(1-2h^3)y^3-3h^2xyz
\]
を得る。具体的には同文献の中心元 $g$ に上のパラメータを代入し、
$yxz=t^{-1}(xyz-hz^3)$ を用いると $g=-hG'_h$ となる。
$S'_h$ は $h$-torsion freeなので $G'_h$ は中心元である。
その平坦性は、一般ファイバーのSklyanin Hilbert級数と特殊ファイバーの多項式環から、
有限生成 $\CC[[h]]$-加群の構造定理を各次数に適用して従う。

\[
 A'_h=S'_h[u,v]/(G'_h+u^3+v^3)
\]
とおく。monicな $u$ の式からこれも次数ごとに平坦であり、
前節と同じ次数シフト・テンソル積の比較により
\[
 \mathcal K(A'_h)\simeq\Ku(X_0)\otimes\CC[[h]]
\]
を得る。一方、$t=1-h^3/3+O(h^6)$ なので、元の族と半古典括弧は同じである。
元の中心元を $xyz$ 順序で書くと
\[
 G_h=(1+h^3)(x^3+z^3)+(1-h^3)y^3-2h^2xyz.
\]
従って両族の関係式を $h^3$ で比較して
\[
 G'_h-G_h\equiv-h^2xyz\pmod{h^3}
\]
となる。これは $\kappa_2=[xyz]$ を消す幾何学的補正を明示する。
高次には中心性を保つため乗法も補正しており、
元の代数のまま中心でない $h^2xyz$ だけを引いているのではない。

\section{6個の例外ファイバーは射影的に特異である}\label{sec:exceptional}

\begin{theorem}\label{thm:exceptional}
全ての $\lambda\in\CC$ について
\[
 \operatorname{gldim}(\qgr A_\lambda)=
 \begin{cases}
 3&\lambda^6\ne-27,\\
 \infty&\lambda^6=-27.
 \end{cases}
\]
ここで大域次元はcoherent対象間のExtによるホモロジー次元である。
\end{theorem}

\subsection{Fourier変換と中心切断の三角形}

$\omega^2+\omega+1=0$、$\lambda_0=\omega-\omega^2$ とすると $\lambda_0^2=-3$ である。
前節の対称性 $h\mapsto\omega h$、$h\mapsto-h$ は6個の例外ファイバーを同型にする。
従って $\lambda_0$ だけを考えればよい。
\[
 P=x+y+z,\quad Q=x+\omega y+\omega^2z,\quad R=x+\omega^2y+\omega z
\]
とおく。ここで $R$ はこの節の第三生成元を表す。
直接代入により
\[
 S_{\lambda_0}\simeq
 \CC\langle P,Q,R\rangle/(P^2+2RQ,\ Q^2+2PR,\ R^2+2QP).
\]
二次関係から $P^3=Q^3=R^3$。
一方、自由代数内のFourier展開から
\[
 P^3+Q^3+R^3
 =3\left(x^3+y^3+z^3+\sum_{\tau\in S_3}\tau(xyz)\right)
 =3G_{\lambda_0}
\]
なので、厳密に $G_{\lambda_0}=P^3$ である。

量子多項式環
\[
 T=\CC\langle P,Q,R\rangle/(PR+2RP,\ QP+2PQ,\ RQ+2QR)
\]
と位数3の自己同型 $\tau(P)=R,\tau(Q)=P,\tau(R)=Q$ をとる。
Zhang積 $a*b=a\tau^{\deg a}(b)$ で計算すると
$S_{\lambda_0}=T^\tau$、$P*P*P=-2PQR$ であり、
\[
 C:=S_{\lambda_0}/(G_{\lambda_0})\simeq(T/(PQR))^\tau.
\]
Zhang同値\cite[Theorem 1.1]{Zhang}によりgraded加群圏とtailsは保存される。

$T/(PQR)$ は三角形 $E=V(PQR)\subset\PP^2$ のtwisted homogeneous coordinate ringである。
$q=-2$ とすると、その自己同型は各成分で
\[
 \begin{array}{c|c}
 P=0 &[Q:R]\mapsto[qQ:R]\\
 Q=0 &[P:R]\mapsto[P:qR]\\
 R=0 &[P:Q]\mapsto[qP:Q]
 \end{array}
\]
となり、節点では一致する。そのgraph方程式が $T$ の多重線形関係である。
各次数の切断は、三つの頂点を検出する純冪と各辺の混合単項式で生成され、
次数 $n\ge1$ の次元は両辺とも $3n$ となる。これでTHCRとの同型を直接確認できる。
ample sequenceのmultidegreeは $(n,n,n)$ であり、
特異な射影スキームにも適用できるTHCRの定理\cite{AV,SVdB}から
\[
 \qgr C\simeq\operatorname{coh}E
\]
を得る。

$C/PC$ は辺 $L=V(P)$ の構造層に対応する。節点での完備局所模型は
\[
 B=\CC[[a,b]]/(ab),\qquad \OO_L=B/(a).
\]
この加群の自由分解は $a,b$ の掛け算が交互に現れる無限周期分解であり、
剰余体とテンソルすると全ての正次数のTorが非零となる。
従って $\OO_L$ はperfectでない。
$\pi_C C(n)$ は直線束なので、$\pi_C(C/PC)$ は
\[
 D^b(\qgr C)/\operatorname{thick}\{\pi_C C(n)\}
\]
で非零である。この段階で滑らかな楕円曲線に対するOrlov同値は用いていない。

\subsection{特異性を検出する階数4の行列因子化}

$S=S_{\lambda_0}$ 上の階数1の因子化 $(P,P^2)$ を使い、
\[
 Q=(P,P^2)\boxtimes(u,u^2)\boxtimes(v,v^2)
\]
を作る。これは $S[u,v]$ 上の $W=P^3+u^3+v^3$ のgraded MFで、各parityの階数は4である。
外積作用 $\varepsilon_i$ と縮約 $\iota_i$ により、全微分は
\[
 d=P\varepsilon_1+P^2\iota_1+
 u\varepsilon_2+u^2\iota_2+v\varepsilon_3+v^2\iota_3
\]
と書ける。中心性から異なる因子は反可換し、$d^2=W I$ となる。
内部次数シフトは通常の三次MFのテンソル積に従う。
この等式は添付コードで $8\times8$ 行列の全64成分について検証する。

$N=\operatorname{coker}Q$ とする。$(u,v)$ は $A=A_{\lambda_0}$ の中心正則列なので、
導来制限
\[
 -\otimes_A^{\mathbf L}C:D^b(\qgr A)\to D^b(\qgr C)
\]
はTor振幅2を持ち、$\pi_AA(n)$ を $\pi_CC(n)$ に送る。
graded torsion加群のTorもtorsionなので、tailsへの降下も定義される。
従って両側のline objectが生成する厚い部分圏で割った商への関手を得る。

$Q$ を $u=v=0$ に特殊化すると、$(P,P^2)$ のshift・twistを変えた4個の直和となる。
これは導来制限の像の周期的自由分解の末尾を表す。
実際、$N$ の自由分解を $C$ とテンソルすれば、Tor振幅2により十分高次の末尾は
この特殊化と一致し、有限の初期部分は上の商圏ではperfectな差にしかならない。
前小節の非perfect性からその像は非零であり、従って
\[
 \pi_A N\notin\operatorname{thick}\{\pi_AA(n)\}.
\]

もし $\qgr A$ の大域次元が有限の $d$ なら、任意のcoherent対象 $M$ の
line objectの有限和による表示を $d+1$ 段まで取り、最後の核を $K$ とする。
対応する有限複体の連結類は $\Ext^{d+1}(M,K)=0$ に入る。
従って $M$ はその有限複体の直和因子であり、上の厚い部分圏に入る。
これは $N$ と矛盾する。以上で命題\ref{thm:exceptional}を得る。

Zhang twistは中心切断 $C$ の解析にだけ用いた。
中心変数 $u,v$ の付加がそのtwistと可換するとは仮定していない。

\section{二段階の貼り合わせと単一対象の変形}\label{sec:gluing-objects}

\subsection{二つのgluing object}

$V=H^0(X,\OO_X(1))$、$\dim V=5$ とし、評価列
\[
 0\longrightarrow Q\longrightarrow V\otimes\OO_X
 \longrightarrow\OO_X(1)\longrightarrow0
\]
で $Q=\Omega^1_{\PP^4}(1)|_X$ を定める。
$Q=L_{\OO_X}\OO_X(1)[-1]$ は例外対象であり、mutationにより
\[
 \mathcal D=\Perf(X)=\langle\mathcal K,Q,\OO_X\rangle,\qquad
 \mathcal B=\OO_X^\perp=\langle\mathcal K,Q\rangle
\]
を得る。$\mathcal K=\Ku(X)$、包含を $i:\mathcal K\to\mathcal D$、
$j:\mathcal B\to\mathcal D$ と書き、
\[
 E=i^!Q\in\mathcal K,\qquad G=j^!\OO_X\in\mathcal B
\]
とおく。$E$ は $\mathcal K$ に $Q$ を加えるgluing object、
$G$ は $\mathcal B$ に最後の $\OO_X$ を加えるgluing objectである。
$E$ の構成と次の三角はFeyzbakhsh--Liu--Zhang
\cite[Section 3.2, equations (9)--(12)]{FLZ}にも記載される。
\[
 \OO_X(-1)[1]\longrightarrow E\longrightarrow Q
 \longrightarrow\OO_X(-1)[2]. \tag{E}\label{eq:E-triangle}
\]
以下、随伴の同一視および必要な1次元のtrace線の自明化を固定する。
例えば上の三角は、選択なしには
$\OO_X(-1)[1]\otimes\Ext^2(Q,\OO_X(-1))^*$ を左端に持つ。
従って以下の裸の $\Lambda^2V$ との同一視を、選択なしの
$GL(V)$-同変な同一視と解釈する必要はない。

\begin{theorem}[二つの対象の自己Ext]\label{thm:gluing-ext}
上の規約の下で
\[
 \Ext^r_{\mathcal K}(E,E)=
 \begin{cases}\CC&r=0,\\ \Lambda^2V&r=1,\\0&\text{その他},\end{cases}
 \qquad
 \Ext^r_{\mathcal B}(G,G)=
 \begin{cases}\CC&r=0,\\ \Sym^2V&r=2,\\0&\text{その他}.\end{cases}
\]
従って $E$ の古典的な標識付き対象変形の同値類は滑らかな10次元の形式空間をなし、
$G$ の対象変形は剛的で、障害の入り得る空間は15次元である。
\end{theorem}

\begin{proof}
随伴から $\RHom_{\mathcal K}(E,E)=\RHom_{\mathcal D}(E,Q)$ である。
評価列を $\OO_X(1)$ でtwistすると
\[
 \RHom(\OO_X(-1),Q)=R\Gamma(Q(1))
 =\ker(V\otimes V\to\Sym^2V)=\Lambda^2V
\]
を次数零に得る。$Q$ の例外性と\eqref{eq:E-triangle}に
$\RHom(-,Q)$ を適用すれば三角
\[
 \CC\longrightarrow\RHom(E,E)\longrightarrow
 \Lambda^2V[-1]\longrightarrow\CC[1]
\]
を得るので、最初の計算が従う。

一方、Serre関手 $S_{\mathcal D}=(-2)[3]$ により
\[
 \mathcal D=\langle\OO_X(-2)[3],\mathcal B\rangle
\]
と回転できる。
$\RHom(\OO_X,\OO_X(-2)[3])=\CC$ なので、右mutationは
\[
 \OO_X(-2)[2]\longrightarrow G\longrightarrow\OO_X
 \longrightarrow\OO_X(-2)[3] \tag{G}\label{eq:G-triangle}
\]
を与える。随伴と $\RHom(-,\OO_X)$ により
\[
 \CC\longrightarrow\RHom(G,G)\longrightarrow
 H^0(\OO_X(2))[-2]\longrightarrow\CC[1]
\]
を得て、二番目の計算も従う。

対象の変形は自己準同型dg代数のMaurer--Cartan理論で制御される
\cite[Proposition 6.1]{ELOI}\cite[Theorems 7.2, 11.8]{ELOII}。
$E$ のstrictly unitalな最小 $A_\infty$ 模型は次数0の単位と次数1しか持たない。
正次数入力だけを持つ $m_n$ の出力は次数2なので全て零であり、
次数零のscalar gaugeも可換Artin基底上では自明に作用する。
従って
\[
 \Def_E\simeq\Spf\CC[[z_1,\ldots,z_{10}]].
\]
これは対象の同値類の形式基底であり、scalar自己同型を含むstack全体ではない。
\end{proof}

\subsection{非可換接方向がこの対象を動かすことの証明}

半直交分解 $\mathcal T=\langle\mathcal A,F\rangle$ に対し、
$F$ が例外対象で $E_F=i^!F$ なら、Hochschild gluing triangleは
\begin{equation}\label{eq:gluing-triangle}
 CC^\bullet(\mathcal T)\longrightarrow
 CC^\bullet(\mathcal A)\oplus CC^\bullet(\CC)
 \longrightarrow\RHom(E_F,E_F)
 \longrightarrow CC^\bullet(\mathcal T)[1].
\end{equation}
これは\cite[Theorem 7.7, Remark 7.8]{Kuz}の完全列のdg模型である。
次数零での第二の射は、二つのscalarの差の作用である。

$\mathcal B=\langle\mathcal K,Q\rangle$ に適用すると、命題
\ref{thm:gluing-ext}と第7節の計算から
\[
 \HH^1(\mathcal B)=0,\qquad
 0\longrightarrow\Ext^1(E,E)\longrightarrow\HH^2(\mathcal B)
 \longrightarrow\HH^2(\mathcal K)\longrightarrow0,\qquad
 \HH^3(\mathcal B)=0.
\]
特に $\dim\HH^2(\mathcal B)=20$ である。
次に $\mathcal D=\langle\mathcal B,\OO_X\rangle$ に適用すると、
$\Ext^1(G,G)=0$ から
$\HH^2(\mathcal D)\to\HH^2(\mathcal B)$ は単射であり、
両辺が20次元なので同型となる。また
\[
 \Ext^2(G,G)\xrightarrow{\sim}\HH^3(X).
\]
二つの制限射の合成 $\HH^2(\mathcal D)\to\HH^2(\mathcal K)$ の核は、
補題\ref{lem:bivector-zero}から $H^0(\Lambda^2T_X)$ である。従って
\begin{equation}\label{eq:gluing-tangent}
 H^0(X,\Lambda^2T_X)\xrightarrow{\sim}
 \Ext^1_{\mathcal K}(E,E).
\end{equation}
これは単なる次元比較ではなく、制限射の核とgluing完全列による同型である。
非可換層変形から $E$ の変形への接写像は、上三角dg模型における
混合Hochschild cochainの制限により、この同型と一致する。

\begin{theorem}[滑らかな対象変形空間内の錐]\label{thm:object-cone}
命題\ref{thm:fibre}の固定残余圏の形式ファイバー $\mathfrak F_X$ から、
gluing object $E$ の変形空間への自然な写像は閉埋め込みである。
形式座標を選べば
\[
 \mathfrak F_X=
 \Spf\CC[[a_{ij}]]/(q_{ijkl})
 \lhook\joinrel\longrightarrow
 \Spf\CC[[a_{ij}]]=\Def_E.
\]
この7次元の部分空間は、原点では10次元全体を接空間に持つ。
\end{theorem}

\begin{proof}
構造層の変形で持ち上げた $\OO_X,\OO_X(1)$ から、相対mutationで $Q_A$ を作り、
右随伴で $E_A=i_A^!Q_A$ を得る。perfectな射影bimoduleを使えば、
これらはArtin基底変換と可換である。
固定残余圏の上で標識を保つ同値 $\mathcal K_A\simeq\mathcal K\otimes A$ を選び、
$E_A$ を移す。異なる選択の差は恒等に特殊化する自己同値である。
$\HH^1(\mathcal K)=0$ と第7節の形式比較により、その同型類の形式群は自明なので、
対象の同値類はこの選択に依存しない。

従って $\mathfrak F_X\to\Def_E$ は定義され、接写像は
\eqref{eq:gluing-tangent}の同型である。
対応する完備局所環の射
$\CC[[z_1,\ldots,z_{10}]]\to\CC[[a_{ij}]]/(q)$
は余接空間上同型なので全射である。
$z_i$ の像を $\CC[[a_{ij}]]$ に持ち上げると、その一次項は可逆であり、
形式座標変換を定める。この座標で核はちょうどPlückerイデアルである。
\end{proof}

\subsection{元の例外対との対応}

元の分解で使う二つの表現対象も書ける。
$\pi=i^*:\mathcal D\to\mathcal K$ とすれば、Serre双対性から
\[
 P_r:=i^!\OO_X(r)=S_{\mathcal K}\pi(\OO_X(r+2))[-3],
 \qquad r=0,1.
\]
従って $F\in\mathcal K$ に対し
$\RHom_{\mathcal D}(F,\OO_X(r))=\RHom_{\mathcal K}(F,P_r)$ である。
評価三角に $i^!$ を適用すると
\[
 E\longrightarrow V\otimes P_0\longrightarrow P_1\longrightarrow E[1]
\]
を得る。さらにrotation $\rho$ に対して
$\rho(P_0)\simeq E$ である\cite[Remark 7.2]{FLZ}。
可換の場合には二回のEuler mutationから
\[
 \pi(\OO_X(2))=\Omega^2_{\PP^4}(2)|_X[2],\qquad
 P_0=S_{\mathcal K}\bigl(\Omega^2_{\PP^4}(2)|_X\bigr)[-1]
\]
とも書ける。単一対象 $E$ の自己Extと、二対象間の射を含む全体の貼り合わせは、
変形問題として区別する必要がある。

特に第9節の固定残余圏族 $A'_h$ の一次Poisson類は非零なので、
\eqref{eq:gluing-tangent}から、固定した $\Ku(X_0)$ 内のgluing objectは
既に $h$ の一次から非自明に動く。
この方向は第4節の線形束 $\langle u,v\rangle$ に対応する分解可能双ベクトルの直線である。
残余圏が一定であることは、この追加対象まで一定であることを意味しない。

\section{gluing bimoduleの形式変形基底}\label{sec:gluing-moduli}

\subsection{対象と自己Ext}

例外ブロックを
\[
 \mathcal P=\langle\OO_X,\OO_X(1)\rangle
 \simeq\Perf(\CC Q_5)
\]
と書く。$Q_5$ は二頂点間に5本の平行な矢を持つKronecker quiverである。
$\mathcal D=\langle\mathcal K,\mathcal P\rangle$ のgluing bimoduleを
\[
 M(F,P)=\RHom_{\mathcal D}(F,P),\qquad
 F\in\mathcal K,\ P\in\mathcal P
\]
とする。これは $\mathcal K^{\mathrm{op}}\otimes\mathcal P$ 上のperfectな
双関手として扱う。以下 $\Def_M$ は、二つのdg圏を固定して、このperfect
bimoduleを可換局所Artin環上で標識付きに変形する群oidである。

$\mathcal P$ の相対Hochschild複体は、頂点の $\CC^2$ と
矢空間の $\operatorname{End}(V)$ からなり、微分は
$(a,b)\mapsto(b-a)\id_V$ である。長さ2の矢道は存在しないので
\[
 \HH^0(\mathcal P)=\CC,\qquad
 \HH^1(\mathcal P)=\mathfrak{pgl}(V),\qquad
 \HH^{\ge2}(\mathcal P)=0.
\]
gluing triangle
\begin{equation}\label{eq:full-gluing-triangle}
 CC^\bullet(\mathcal D)\longrightarrow
 CC^\bullet(\mathcal K)\oplus CC^\bullet(\mathcal P)
 \longrightarrow\RHom(M,M)\longrightarrow CC^\bullet(\mathcal D)[1]
\end{equation}
と、第7節およびHKRの計算から次を得る。

\begin{theorem}[全体の貼り合わせの接空間と障害]\label{thm:bimodule-ext}
\[
 \dim\Ext^r(M,M)=
 \begin{cases}1&r=0,\\34&r=1,\\15&r=2,\\0&\text{その他}.\end{cases}
\]
さらに完全列と同型
\begin{gather}
 0\longrightarrow\mathfrak{pgl}(V)\longrightarrow\Ext^1(M,M)
 \xrightarrow{p}H^0(\Lambda^2T_X)\longrightarrow0,
 \label{eq:bimodule-tangent}\\
 \partial:\Ext^2(M,M)\xrightarrow{\sim}\HH^3(X)
 \simeq H^0(\OO_X(2))
 \label{eq:bimodule-obstruction}
\end{gather}
を持つ。24次元の部分空間は、Kroneckerブロックの矢の座標変更の微分である。
\end{theorem}

\subsection{二次障害は五つのPlücker式である}

方向付きdg模型では、二つのブロックをまたぐHochschild cochainは、
入力にちょうど一つの $M$、出力に $M$ を持つ。
両端に $\mathcal K,\mathcal P$ のbar入力を持つこの複体は、
シフト後にはbimoduleのbar自己準同型複体となる。
混合入力へのbrace挿入はbimodule premorphismの合成なので、
\[
 \RHom_{\mathrm{bar}}(M,M)
 \longrightarrow CC^\bullet(\mathcal D)[1]
\]
はdg Lie射である。そのコホモロジー射は
\eqref{eq:full-gluing-triangle}の連結射である。

$\xi\in\Ext^1(M,M)$、$\pi=p(\xi)$ とすると、二次延長の障害は
Yoneda平方 $\xi^2$ であり、上のdg Lie射によって
\begin{equation}\label{eq:yoneda-schouten}
 \partial(\xi^2)=\frac12[\pi,\pi]\in H^0(\Lambda^3T_X)
\end{equation}
となる。ここでは大域双ベクトルの反対称化した双微分cochainとHKRを用いれば、
Hochschildコホモロジー上でGerstenhaber括弧がSchouten括弧に対応することを確認できる。
cochainそのものの厳密な等式を用いているわけではない。

$\pi=\pi(a)$ を第2節の座標で書くと、\eqref{eq:bracket}から
障害写像の15成分が張る二次式の空間は
\[
 \langle q_{0123},q_{0124},q_{0134},q_{0234},q_{1234}\rangle
\]
に一致する。$F_0,\ldots,F_4$ の線形独立性がここでも必要である。
従って独立な二次式は5本であり、その零点の余次元は3である。
また\eqref{eq:bimodule-obstruction}が同型なので、
24個の矢座標方向を含む交差項も含め、二次障害は $\pi$ だけで決まる。

\subsection{全次数の古典的形式基底}

\begin{theorem}[固定した両ブロック間の全形式gluing変形]\label{thm:full-gluing-base}
古典的な同値類の関手 $\pi_0\Def_M$ は、非標準的に
\[
 \Spf
 \frac{\CC[[b_1,\ldots,b_{24},a_{ij}\ (0\le i<j\le4)]]}
 {(q_{ijkl})}
 \simeq
 \widehat{\operatorname{PGL}(V)}_{\id}
 \times\widehat{C(\Gr(2,5))}_0
\]
で表される。次元は31、接空間次元は34である。
恒等に特殊化する矢の座標変更で割ると、
命題\ref{thm:fibre}の7次元の固定残余圏ファイバーを得る。
\end{theorem}

\begin{proof}
\textbf{Maurer--Cartanによる局所環。}
$\Ext^{-1}(M,M)=0$ で自己Extは上に有界なので、
\cite[Theorem 11.8, Proposition 11.10]{ELOII}により導来変形を
h-projective模型の変形に移せる。
\cite[Proposition 6.1]{ELOI}と\cite[Theorem 7.2]{ELOII}により、
strictly unitalな最小 $A_\infty$ 模型 $H=\Ext^*(M,M)$ を使える。
$H^0=\CC\cdot1,\ H^1=\CC^{34},\ H^2=\CC^{15}$ であり、
可換Artin基底の極大イデアル $\mathfrak m$ 上の方程式は
\[
 \kappa(z)=\sum_{n\ge2}m_n(z,\ldots,z)=0,\qquad
 z\in H^1\otimes\mathfrak m.
\]
次数零で恒等に特殊化する可逆射は $(1+c)1$ に限られる。
strict unitにより二つのMC元 $z,z'$ の間でこの射が閉じる条件は
$(1+c)(z-z')=0$ なので、$z=z'$ である。
従って同値類の関手の環は
\[
 R_M=P/I,\qquad P=\CC[[z_1,\ldots,z_{34}]],
\]
ただし $I$ は $\kappa$ の15成分が生成するイデアルである。
その二次項は前小節のYoneda平方である。

この比較で対象のperfect性も保たれる。
固定した対象対でbimoduleを評価するか、compact generatorのMorita模型を選んで考える。
Artin環 $A$ は有限長の剰余体filtrationを持つため、特殊ファイバーが
有界有限次元の複体なら、その評価の導来変形も基底上で有界coherentとなる。
有限Tor振幅の判定から $A$-perfectである。
さらに $\mathcal K^{\mathrm{op}}\otimes\mathcal P$ のperfectな対角分解を
テンソルすれば、圏上でもperfectとなる。
従ってここでのMC関手は、最初に定義したperfect bimoduleの関手を計算している。

\textbf{既存の非可換族から34方向の族を作る。}
$R_{\mathfrak F}=\CC[[a_{ij}]]/(q)$ 上の固定残余圏族を取る。
各 $R_{\mathfrak F}/\mathfrak m^n$ で残余圏の標識付き自明化を選べる。
隣接する自明化の差は恒等に特殊化する自己同値であり、
$\HH^1(\mathcal K)=0$ と第7節の比較から、その同型類は一意である。
従って各隣接段階で自然同型を選び、逐次的に整合する自明化を取れる。
$\HH^0$ は高次のscalarの選択を残すが、この逐次的な選択の存在を妨げない。
また矢空間は階数5の自由加群なので、その基底も帰納的に整合して選べる。
これにより、固定した $\mathcal K,\mathcal P$ 上の形式bimodule $M_a$ を得る。

その $\mathcal P$-作用を矢の自己同型で変える。
左作用の規約は $(g\cdot M)(F,P)=M(F,g^{-1}P)$ とする。
$GL(V)$ のscalar部分はKronecker代数の内部自己同型なので、
同値類上の有効な作用は $\operatorname{PGL}(V)$ である。こうして
\begin{equation}\label{eq:framed-family}
 \widehat{\operatorname{PGL}(V)}_{\id}\times\Spf R_{\mathfrak F}
 \longrightarrow\pi_0\Def_M
\end{equation}
を構成する。
\eqref{eq:bimodule-tangent}によって接写像は同型である。

\textbf{高次の追加方程式がないこと。}
形式座標を直して $P=\CC[[a_{ij},b_1,\ldots,b_{24}]]$、
$J=(q_{ijkl})$ とすれば、\eqref{eq:framed-family}の環射は
$P/I\to P/J$ となり、$I\subset J$ である。
前小節の計算から $I$ の二次項は5本のPlücker式を張る。
これらは $J/\mathfrak m_PJ$ の基底であり、
\[
 J\cap\mathfrak m_P^3=\mathfrak m_PJ
\]
である。実際、$f=\sum c_iq_i$ の二次項が零なら、
$q_i$ の線形独立性から全ての $c_i(0)$ が零になる。
従って $I$ 内に $f_i=\sum_j A_{ij}q_j$ で $A(0)$ 可逆となる元を選べる。
$A$ は $P$ 上可逆なので、全ての $q_j$ は $I$ に入る。よって $I=J$。

この結果、\eqref{eq:framed-family}は同型である。
構成から形式 $\operatorname{PGL}(V)$ の左移動に同変なので、
この作用で割った同値類は $\Spf R_{\mathfrak F}$ となる。
\end{proof}

\begin{remark}[古典的基底と高次構造]
命題\ref{thm:full-gluing-base}は可換局所Artin基底上の標識付き同値類についての主張である。
$M$ のscalar自己同型や導来モジュライの構造は、この通常の完備局所環には含まれない。
古典的イデアルが二次式だけで生成されることから、
自己準同型 $A_\infty$ 代数全体のformalityまでは結論していない。
\end{remark}

\subsection{最後の例外対象を付加できる正確な条件}

単一対象 $E$ の変形 $\widetilde E$ から、
一つの例外対象 $Q$ を加えた圏
$\mathcal B_{\widetilde E}=\langle\mathcal K,Q\rangle_{\widetilde E}$
を作れる。ここで $G\in\mathcal B$ が
$\mathcal B_{\widetilde E}$ のperfect対象に持ち上がる条件を考える。

$\Ext^1(G,G)=0$ なので、持ち上げが存在すれば同型類は一意である。
さらに特殊ファイバーで
$\RHom(Q,G)=\RHom(Q,\OO_X)=V^*$ は次数零だけなので、
導来Nakayamaから、任意のperfectな持ち上げでもこのHomは
次数零の階数5自由加群である。
$G$ をgluing objectとして $\OO_X$ を付加すると、
二つの例外対象のブロックは再びKronecker型となり、
mutation $\OO_X(1)=R_{\OO_X}(Q)[1]$ により元の
$\langle\OO_X,\OO_X(1)\rangle$ を回復する。
逆に、$\Def_M$ の任意の変形から、この二段階の分解と $G$ の持ち上げが得られる。
矢の基底変更は評価射のmutationで定まる $Q$ の同型類を保存するので、
$\Def_M\to\Def_E$ は形式 $\operatorname{PGL}(V)$ の商を経由する。
矢の基底を忘れれば、以上の二つの構成は同型類上で互いに逆である。

従って命題\ref{thm:full-gluing-base}と命題\ref{thm:object-cone}により、
\[
 \left\{\widetilde E\in\Def_E:
 G\text{ が }\mathcal B_{\widetilde E}\text{ のperfect対象に持ち上がる}\right\}
 =
 \widehat{C(\Gr(2,5))}_0\ \subset\ \widehat{\Aff}^{10}.
\]
すなわち、$E$ 単独の変形は障害を持たないが、
最後の例外対象を付加する条件が5本のPlücker式を課す。
これが固定残余圏の7次元特異空間のgluingによる解釈である。

\section{点加群スキームと長さ8の非被約点}\label{sec:points}

\subsection{点加群と中心変数による分離}

体 $k$ 上の次数1生成代数 $A$ に対し、可換 $k$-代数 $T$ 上のpoint moduleとは、
cyclicな次数付き右 $A\otimes_kT$-加群 $N=\bigoplus_{n\ge0}N_n$ で、
$N_0=T$、全ての $N_n$ が階数1の局所自由 $T$-加群となるものをいう。
その同型類の関手を表すスキームを $\mathscr P(A)$ と書く。
定義とtruncationによる構成については\cite[Section 4]{RZ}を参照。
以下では実際に有限段階で安定することを示すので、
strongly noetherian性からの一般的な表現可能性定理には依存しない。

$S=k\langle x,y,z\rangle/(\text{quadratic relations})$ の
二次関係の多重線形化が、$E\subset\PP^2$ 上の自己同型 $\sigma$ のgraphであるとする。
また $\sigma$ は幾何学的固定点を持たず、三次の中心元 $g$ は
このgraphに対応する普遍point moduleを零化すると仮定する。
$u,v$ を中心変数とし、
\[
 B=S[u,v],\qquad A=B/(g+u^3+v^3)
\]
とおく。$S_{\mathrm{ab}},A_{\mathrm{ab}}$ は可換化を表す。

\begin{lemma}[中心変数を付加したpoint scheme]\label{lem:point-splitting}
上の仮定の下で、スキームとして
\[
 \mathscr P(A)=E\ \sqcup\ \operatorname{Proj}A_{\mathrm{ab}}
\]
である。point moduleのshiftは、第一成分で $\sigma$、第二成分で恒等となる。
\end{lemma}

\begin{proof}
連続する二つの座標点を
$p=(x,y,z,u,v)$、$p'=(x',y',z',u',v')$ と書く。
中心性から $x_iu'-ux_i'=0$、$x_iv'-vx_i'=0$ などを得る。
$u\ne0$ のチャートでは $u'$ も局所的に可逆である。
実際、ある幾何学的点で $u'=0$ なら中心性から $p'$ の全座標が零となる。
従って両側を規格化すると $p=p'$ であり、残る二次式はちょうど
$S_{\mathrm{ab}}$ の関係式である。$v\ne0$ でも同じである。

最初の幾何学的点で $u=v=0$ なら、少なくとも一つの $x_i$ が非零なので、
中心性から $u'=v'=0$ も従う。
このような幾何学的点では、
$(x':y':z')=\sigma(x:y:z)$ である。
固定点がないため、ある $2\times2$ 小行列式
$\Delta=x_ix_j'-x_jx_i'$ は非零であり、局所近傍で可逆となる。
二本の中心性の式から $u=u'=0$ がスキームとして従い、
同様に $v=v'=0$ である。従ってこの成分の法方向に冪零構造も残らない。
これらの近傍は全ての幾何学的点を覆うので、二次関係スキームは
\[
 \operatorname{Graph}(\sigma)\ \sqcup\
 \Delta_{\operatorname{Proj}(S_{\mathrm{ab}}[u,v])}
\]
である。両者の台は交わらない。

$X_n(B)\subset(\PP^4)^n$ を $N_0,\ldots,N_n$ のtruncationスキームとする。
上の二次関係がgraphなので、全ての $n\ge2$ で
$X_n(B)$ は最初の座標だけから決まる。
三次式は $X_3(B)$ 上で、$E$ 成分では仮定から零、diagonal成分では
可換化した $g+u^3+v^3$ となる。その零点はshiftで保たれる。
従って $X_n(A)\simeq X_3(A)$ が全ての $n\ge3$ で成立し、
このスキームが無限point moduleの関手を表す。
議論は任意の冪零を含む試験環上で成り立つ。
\end{proof}

\subsection{元のFermat族の全ての非零ファイバー}

\begin{theorem}[明示的族のpoint scheme]\label{thm:point-scheme}
$\lambda\in\CC^\times$ とし、
\[
 J_0=\CC[x,y,z]/(x^2,y^2,z^2)
\]
とおく。このとき
\[
 \boxed{
 \mathscr P(A_\lambda)
 =E_{\lambda^2}\ \sqcup\
 \coprod_{\zeta^3=-1}\operatorname{Spec}J_0
 },\qquad
 E_s=V(x^3+y^3+z^3+sxyz).
\]
各非被約点の長さは8、Hilbert関数は $(1,3,3,1)$ である。
この式は $\lambda^6=-27$ の三角形のファイバーにも成立する。
$\lambda=0$ では $\mathscr P(A_0)=X_0$ である。
\end{theorem}

\begin{proof}
第5節の多重線形化行列 $M_\lambda$ の階数は、
$E_{\lambda^2}$ 上で常に2である。
実際、階数が1以下なら、二つの小行列式
$xy-\lambda z^2$ と $xy+\lambda z^2$ がともに零となり $z=0$。
巡回した式から $x=y=0$ となるので矛盾する。
逆向きの関係行列も同じ性質を持つため、関係スキームは自己同型のgraphである。
diagonalに代入した二次式は
$\lambda x^2=\lambda y^2=\lambda z^2=0$ なので、固定点はない。

$G_\lambda$ が普遍point moduleを零化することは、滑らかな一般ファイバーで
$S_\lambda/(G_\lambda)=B(E_{\lambda^2},L,\sigma)$ から従う。
非零パラメータ上の相対Hesse曲線は、一般ファイバーが幾何学的整で
基底上flatな整スキームである。
$G$ の多重線形化はその上の可逆層の切断なので、一般ファイバーでの零化から
全体で零となる。これにより三角形の6ファイバーも含めて
補題\ref{lem:point-splitting}を適用できる。

可換化は
\[
 (A_\lambda)_{\mathrm{ab}}=
 \frac{\CC[x,y,z,u,v]}
 {(x^2,y^2,z^2,\ u^3+v^3-2\lambda^2xyz)}
\]
である。そのProjの台は $[0:0:0:1:\zeta]$、$\zeta^3=-1$ の3点である。
$u=1$ とし、$T=xyz\in J_0$ とおくと $T^2=0$ である。
$\zeta$ の枝で
\[
 v=\zeta+\frac{2\lambda^2}{3\zeta^2}T
\]
と置けば、三次式は厳密に消える。
従って各局所成分は $\operatorname{Spec}J_0$ である。
$J_0$ の基底は
\[
 1,x,y,z,xy,xz,yz,xyz
\]
であり、長さとHilbert関数も従う。
\end{proof}

第一座標の $\PP^4$ への埋め込みで、その斉次イデアルは
\[
 (u,v,D+\lambda^2xyz)\ \cap\
 (x^2,y^2,z^2,u^3+v^3-2\lambda^2xyz)
\]
である。Hilbert多項式は $3n+24$ となる。
従って、楕円曲線と三つの被約点だけを記述すると、
零次元部分の長さ24の情報を失う。

\subsection{残余圏を固定する族でも非被約構造は動く}

第9節の族 $A'_h$ を $K=\CC((h))$ 上で考え、
\[
 r=(1-h^3)^{1/3},\qquad
 a=\frac{1-r}{h}=\frac{h^2}{3}+\frac{h^5}{9}+O(h^8)
\]
とおく。この $r$ は第9節の $t$ と同じ元である。
可換化した二次関係は
\[
 x^2=a yz,\qquad y^2=a zx,\qquad z^2=a xy
\]
である。また
\[
 a^3=\frac{(1-r)^2}{1+r+r^2},\qquad
 a^3-1=-\frac{3r}{1+r+r^2}\ne0.
\]
これらの二次式は $\PP^2_K$ に共通零点を持たないので、
対応するSklyaninの自己同型は幾何学的固定点を持たない。

\begin{theorem}[固定残余圏族の非被約点]\label{thm:fixed-points}
\[
 J_a=
 \frac{K[x,y,z]}
 {(x^2-a yz,\ y^2-a zx,\ z^2-a xy)}
\]
とおくと、
\[
 \mathscr P(A'_h\otimes K)
 =(E_0)_K\ \sqcup\
 \coprod_{\zeta^3=-1}\operatorname{Spec}J_a.
\]
$J_a$ は長さ8でHilbert関数 $(1,3,3,1)$ を持つ。
その局所代数に内在する三次形式は
\[
 \Phi_a(X,Y,Z)=a(X^3+Y^3+Z^3)+6XYZ.
\]
これは滑らかな平面三次曲線を定め、その $j$-不変量は
\begin{equation}\label{eq:fat-j}
 j_{\mathrm{fat}}(a)
 =\frac{110592(a^3-1)^3}{a^3(a^3+8)^3}
 =-5832h^{-6}\bigl(1+O(h^3)\bigr)
\end{equation}
である。従って残余圏と被約な楕円曲線成分が一定でも、
非被約点の局所代数は一定ではない。
\end{theorem}

\begin{proof}
二次関係から $(1-a^3)x^2y=0$ とその置換を得る。
従って指数型 $(2,1,0)$ の三次単項式は零であり、
\[
 x^3=y^3=z^3=a\,xyz,\qquad
 x\,xyz=y\,xyz=z\,xyz=0.
\]
次数4以上は零となり、$1,x,y,z,xy,xz,yz,xyz$ が張る。
三つの二次式の射影共通零点がないので、多項式環のCohen--Macaulay性から
それらは正則列であり、Hilbert級数は $(1+t)^3$ である。
従ってこの8元は独立でもあり、基底となる。
添付コードはこの積表の512個の基底三つ組について結合律を厳密に確認する。

$T=xyz$ とおくと、中心元の可換化は
\[
 (G'_h)_{\mathrm{ab}}
 =\bigl(2(1+h^3)a+(1-2h^3)a-3h^2\bigr)T
 =3(a-h^2)T.
\]
$u=1$、$\zeta^3=-1$ の枝で
\[
 v=\zeta-\frac{a-h^2}{\zeta^2}T
\]
と置けば三次式が消える。補題\ref{lem:point-splitting}から最初の分解を得る。

極大イデアルを $\mathfrak n$ と書くと、
$\mathfrak n/\mathfrak n^2$ は $x,y,z$ の3次元空間、
$\mathfrak n^3=K\cdot T$ は1次元で、これらと三重積は局所代数に内在する。
直接展開により
\[
 (Xx+Yy+Zz)^3
 =\bigl(a(X^3+Y^3+Z^3)+6XYZ\bigr)T.
\]
従って局所代数の同型は、$\Phi_a$ を線形座標変換と非零scalar倍の範囲で保存する。
$a\ne0$、$a^3\ne-8$ なのでその零点は滑らかなHesse曲線である。
第9節で検証した $j(E_s)$ に $s=6/a$ を代入すると
\eqref{eq:fat-j}を得る。
これは局所代数の同型不変量であり、$h$ に関して非定数である。
\end{proof}

この三次形式はMacaulay dual cubicとも解釈できるが、
ここでは極大イデアルの三重積から直接定義した。
$a=0$ では $\Phi_0=6XYZ$ は三角形であり、一般の $a$ における滑らかな三次曲線と異なる。
固定残余圏族の第一座標のpoint schemeの斉次イデアルは
\[
 (u,v,D)\ \cap\
 \bigl(x^2-a yz,y^2-a zx,z^2-a xy,\,
 u^3+v^3+3(a-h^2)xyz\bigr)
\]
である。

\subsection{平坦な代数族と点加群の特殊化}

$a\in R=\CC[[h]]$ として同じ $J_a$ を定義すると、
$1-a^3$ は単元なので、これは $R$ 上階数8の自由代数である。
上の $v$ の式と固定楕円曲線のgraph族は、相対truncationスキーム内で
一般ファイバーの3個の非被約点と曲線のflatな水平閉包を与える。
その特殊ファイバーは
\[
 E_0\ \sqcup\ 3\operatorname{Spec}J_0
\]
であり、Fermat三次元多様体 $X_0$ の真の閉部分スキームである。

一方、$N_0,\ldots,N_3$ を扱う相対length-4 truncationスキームは
$(\PP^4_R)^3$ の閉部分スキームとして射影的であり、
一般ファイバーは上の1次元スキーム、特殊ファイバーは対角上の $X_0$ である。
Hilbert多項式の次数が異なるため、この相対truncationスキームはflatでない。
これは $A'_h$ 自体の次数ごとのflat性と両立する。
基底全体で無限point module関手の一様な有限表現可能性を主張する必要はない。

point schemeは次数付き代数、従って偏極を持つtailsの三つ組
\[
 (\qgr A,\pi A,(1))
\]
の不変量である。実際、ここではAS Gorensteinのdepthから
$H^0_{\mathfrak m}(A)=H^1_{\mathfrak m}(A)=0$ なので、
局所コホモロジー完全列により
$A\simeq\bigoplus_{n\ge0}\operatorname{Hom}_{\qgr A}(\pi A,\pi A(n))$
と再構成できる。任意の偏極を忘れたdg Morita同値に対する不変量としては用いていない。

\section{貼り合わせの導来変形と形式弧上のtorsion}\label{sec:derived-gluing}

本節では命題\ref{thm:full-gluing-base}の古典的基底に残る導来構造を決定する。
固定する圏は $\mathcal K=\Ku(X)$ と
$\mathcal P=\langle\OO_X,\OO_X(1)\rangle$ であり、変形対象はその間の
perfect gluing bimodule $M$ である。
導来基底には非正次数の可換dg Artin $\CC$-代数を用いる。
この問題は、自己準同型の結合的な $A_\infty$ 構造全体を決定する問題とは区別する。

\subsection{scalarを除いたChevalley--Eilenberg模型}

\begin{theorem}[導来形式基底の非標準的な正規形]\label{thm:derived-gluing-normal}
$M$ のscalar自己同型を除いた可換導来形式変形は、次の完備可換dg代数で表される。
\begin{gather}
 \mathscr R_M\simeq
 \CC[[b_1,\ldots,b_{24},a_{ij}]]
 \otimes\Lambda(\theta_0,\ldots,\theta_4,\rho_1,\ldots,\rho_{10}),
 \label{eq:derived-gluing-normal}\\
 |a_{ij}|=|b_i|=0,\qquad |\theta_i|=|\rho_j|=-1,\qquad
 d\theta_i=q_i(a),\quad d\rho_j=0,\quad da_{ij}=db_i=0.
 \nonumber
\end{gather}
ここで $q_i$ は5本のPlücker式であり、符号と順序は任意に固定する。
scalarを戻した導来形式変形は、この形式空間と
$B\widehat{\mathbb G}_{m,1}$ との積である。

従って $\RHom(M,M)$ のcommutator dg Lie代数は
$L_\infty$ 同値の意味でformalである。
この結論は結合的な $A_\infty$ 代数のformalityを主張しない。
\end{theorem}

\begin{proof}
\textbf{次数とstrict unit。}
命題\ref{thm:bimodule-ext}と
\cite[Theorem 11.8, Proposition 11.10, Theorem 7.2]{ELOII}により、
strictly unitalな最小 $A_\infty$ 模型
\[
 H^0=\CC\cdot1,\qquad H^1=\CC^{34},\qquad H^2=\CC^{15}
\]
を取る。その反対称化を $L$ とし、括弧 $\ell_n$ の次数を $2-n$ とする。
strict unitはcommutatorについて中心的であり、
$n\ne2$ の括弧にunitを入力すると零になる。
残りの入力が全て次数1なら出力は次数2であり、
一つでも次数2の入力を持てば出力は次数3以上となって零になる。
従ってscalar以外の全ての演算は
\[
 \ell_n:\Sym^nH^1\longrightarrow H^2\qquad(n\ge2)
\]
に限られる。次数1の入力についてgraded反対称性は通常の対称性になる。
規約を揃えてMaurer--Cartan方程式を
\[
 \kappa(z)=\sum_{n\ge2}\frac1{n!}\ell_n(z,\ldots,z)=0
\]
と書く。この $\kappa$ は前節の結合的なMC方程式と同じである。

よって $L_+=H^1\oplus H^2$ の完備Chevalley--Eilenberg代数は
\begin{equation}\label{eq:derived-gluing-ce}
 \CC[[z_1,\ldots,z_{34}]]\otimes\Lambda(\eta_1,\ldots,\eta_{15}),
 \qquad |\eta_j|=-1,\quad d\eta_j=\kappa_j(z),\quad dz_i=0.
\end{equation}
実際、CE代数の生成元は $(L_+[1])^\vee$ であり、
次数1の自己Extは次数零の座標、次数2の自己Extは次数 $-1$ の座標を与える。
この代数から可換dg Artin環への射はMC解そのものであり、
simplex上の微分形式を用いれば同じ比較が写像空間とhomotopyにも成り立つ
\cite[Proposition 1.1, Section 4]{Getzler}。
上の次数の議論により、$\eta$ を含む未記載の微分項は存在しない。
中心的な $H^0=\CC\cdot1$ は $L_\infty$ 直和因子であり、
その形式moduliが $B\widehat{\mathbb G}_{m,1}$ である。

\textbf{15個の方程式の行変換。}
命題\ref{thm:full-gluing-base}の証明で得た形式座標を用いて
\[
 P=\CC[[a_{ij},b_1,\ldots,b_{24}]],\qquad
 (\kappa_1,\ldots,\kappa_{15})=(q_0,\ldots,q_4)\subset P
\]
とする。二次項の階数は5なので、$\eta$ の定数係数での基底変換により、
最初の5成分の二次項が独立であるとしてよい。
すると
\[
 \begin{pmatrix}\kappa_1\\ \vdots\\ \kappa_5\end{pmatrix}=Aq,
 \qquad
 \begin{pmatrix}\kappa_6\\ \vdots\\ \kappa_{15}\end{pmatrix}=Bq,
 \qquad A\in\operatorname{GL}_5(P),\quad B\in\operatorname{Mat}_{10\times5}(P).
\]
$A(0)$ が可逆なのは二次項の独立性による。
そこで15個の負次数生成元を行列
\[
 \begin{pmatrix}A^{-1}&0\\-BA^{-1}&I_{10}\end{pmatrix}
\]
で変換すると、その微分は $(q_0,\ldots,q_4,0,\ldots,0)^t$ となる。
行列の係数は $P$ に属し $dP=0$ なので、係数を微分する追加項はない。
これが\eqref{eq:derived-gluing-normal}である。

新しいCE微分は純粋に二次式である。
双対の完備coalgebraに戻せば、これは高次括弧 $\ell_{n\ge3}$ が零である
最小 $L_\infty$ 模型を与える。
残る二項括弧は、基底変更を除いて元のcohomology bracketであるから、
commutator dg Lie代数のformalityを得る。
\end{proof}

\begin{remark}[この正規形が保存する情報]
次数1の元 $x,y$ に対してcommutatorが記録するのは
$m_2(x,y)+m_2(y,x)$ であり、結合積の反対称部分は記録しない。
高次演算でも同様に、対称化されない $A_\infty$ データは残り得る。
従って上の証明から結合的formalityは従わない。
また形式座標と負次数生成元の変換は非標準的であり、
導来正規形の $\operatorname{PGL}(V)$-同変性も主張しない。
命題\ref{thm:full-gluing-base}で証明した古典的な同変性は引き続き用いられる。
\end{remark}

\subsection{形式弧を実際のperfect族として扱う}

$R=\CC[[h]]$、$F_\eta=\CC((h))$ とし、命題\ref{thm:full-gluing-base}の基底に
\begin{equation}\label{eq:gluing-formal-arc}
 b(h)\in hR^{24},\qquad a(h)\in hR^{10},\qquad q(a(h))=0
\end{equation}
という形式弧を取る。
$\Spf$ の原点完備化には、この弧の「非零の閉点」が存在するわけではない。
以下のgeneric fibreは、実際の $R$-対象を作ってから $h$ を可逆化したものである。

\begin{lemma}[形式弧のeffectivity]\label{lem:gluing-arc-effective}
\eqref{eq:gluing-formal-arc}が定める整合的なArtin変形は、
固定した双圏 $(\mathcal K^{\mathrm{op}}\otimes\mathcal P)\otimes R$ 上の
perfect bimodule $M_R$ として実現できる。
その自己準同型複体は $R$-perfectであり、$R/(h^n)$ と $F_\eta$ への
導来基底変換と可換である。
\end{lemma}

\begin{proof}
周囲の圏 $\Perf(\mathcal K^{\mathrm{op}}\otimes\mathcal P)$ はsmooth properである。
Toën--Vaquié\cite[Theorem 3.6, Corollaries 3.15, 3.21--3.22]{TV}により、
負の自己Extが零である $M$ の近傍の対象モジュライはArtin 1-stackである。
弧から得られる $M_n$ はこの同じ開部分に属する。
affineなsmooth atlas $U$ と $M$ 上の点 $u_0$ を選ぶ。
smooth性により、$M_{n+1}$ の上で $u_n$ を $u_{n+1}$ へ逐次持ち上げられる。
剰余体は代数閉体 $\CC$ なので、各持ち上げは局所Artin環自身の上で選べる。
$R=\varprojlim R/(h^n)$ と $U$ のaffine性により、
整合する $u_n$ は $\operatorname{Spec}R\to U$ を定める。
普遍対象の引戻しが求める $M_R$ である。

このmoduliでperfect性は定義に含まれる。
Morita模型で基底上のperfect性を先に確認しても、smoothな対角を用いれば
圏上のperfect性を得る\cite[Lemma 2.8]{TV}。
proper性とperfect性から自己準同型複体の有限性と基底変換が従う。
この議論は完備環 $R$ 上の実現であり、弧が有限型の複素代数曲線に由来するとは主張しない。
\end{proof}

\subsection{Jacobian複体とSmith標準形}

\begin{theorem}[形式弧上の自己Extとtorsion]\label{thm:gluing-arc-ext}
\eqref{eq:gluing-formal-arc}で $a(h)\ne0$ とし、
\[
 \nu=\min_{i<j}\operatorname{ord}_h a_{ij}(h)\ge1,\qquad
 T_\nu=(R/(h^{\nu}))^{\oplus3}
\]
とおく。固定した双圏上の自己Extは
\begin{equation}\label{eq:gluing-arc-ext}
 \Ext^0(M_R,M_R)=R,\qquad
 \Ext^1(M_R,M_R)=R^{31},\qquad
 \Ext^2(M_R,M_R)=R^{12}\oplus T_\nu
\end{equation}
であり、その他の次数では零である。
従って $M_\eta=M_R\otimes_R F_\eta$ について
\[
 \bigl(\dim_{F_\eta}\Ext^0(M_\eta,M_\eta),
       \dim_{F_\eta}\Ext^1(M_\eta,M_\eta),
       \dim_{F_\eta}\Ext^2(M_\eta,M_\eta)\bigr)=(1,31,12)
\]
となる。この計算は全ての非零の分解可能形式弧に成り立つ。
\end{theorem}

\begin{proof}
\textbf{MC方程式の微分。}
最小模型を弧のMC元 $z(h)$ でtwistする。
自己準同型複体の次数1から次数2への微分は
\[
 x\longmapsto\sum_{n\ge1}\frac1{n!}\ell_{n+1}(z,\ldots,z,x)
       =D\kappa_z(x)
\]
である\cite[Proposition 4.4]{Getzler}。
unitの微分は零であり、次数2から出る微分も零である。
従って $\RHom(M_R,M_R)$ は、次数 $0,1,2$ に置かれた
\begin{equation}\label{eq:gluing-jacobian-complex}
 [\,R\xrightarrow{0}R^{34}\xrightarrow{D\kappa_z}R^{15}\,]
\end{equation}
と同値である。
この比較は各 $R/(h^n)$ 上で最小模型の準同型をtwistして作れ、
その式は $h$-進に整合する。
左辺のperfect性と右辺の有限自由性により、逆極限でも同じ比較が成り立つ。

形式座標と方程式の行変換は、この複体上でも可逆変換になる。
実際 $\kappa(z(y))=C(y)q(y)$ なら、MC軌跡 $q=0$ の上で
\[
 D\kappa_z\,Dz_y=C(y)\,Dq_y
\]
であり、$DC$ に由来する項は $q=0$ により消える。
可逆な定数項を持つ全ての変換行列は $R$ 上でも可逆である。
よって\eqref{eq:gluing-jacobian-complex}の非零となり得る微分は、
適当な基底で
\begin{equation}\label{eq:gluing-wedge-differential}
 R^{24}\oplus\Lambda^2R^5
 \longrightarrow\Lambda^4R^5\oplus R^{10},\qquad
 (b,v)\longmapsto(a(h)\wedge v,0)
\end{equation}
となる。ここでは $q(a)=\tfrac12a\wedge a$ という規格化を用いた。

\textbf{DVR上の階数2交代行列。}
$a=h^{\nu} a'$ と書くと、$a'$ はprimitiveであり、$R$ は整域なので $q(a')=0$ である。
基底を置換して $a'_{01}$ がunitとしてよい。
添字 $0,1$ を含む三つのPfaffian式により、残る
$a'_{23},a'_{24},a'_{34}$ を $a'_{01},a'_{0i},a'_{1i}$ で表せる。
従って $a'$ は $R$ 上の独立なprimitiveな二ベクトルのwedgeであり、
$\operatorname{GL}_5(R)$ による変換で $e_0\wedge e_1$ になる。

$R^5=U\oplus W$、$U=\langle e_0,e_1\rangle$、$\rk W=3$ とする。
$(e_0\wedge e_1)\wedge-$ は
$\Lambda^2U\oplus(U\otimes W)$ 上で零であり、
$\Lambda^2W$ を $\Lambda^2U\otimes\Lambda^2W$ に同型に写す。
後者は階数3の直和因子なので、
\eqref{eq:gluing-wedge-differential}の非零のSmith因子は
\[
                         h^{\nu},\quad h^{\nu},\quad h^{\nu}
\]
に尽きる。kernelの階数は $24+7=31$、cokernelは $R^{12}\oplus T_\nu$ であり、
主張を得る。
\end{proof}

\begin{remark}[特殊ファイバーで戻る三つの接方向]
中心の $(1,34,15)$ と弧上の計算は、導来基底変換によって整合する。
実際、完全列
\[
 0\longrightarrow\Ext^1(M_R,M_R)\otimes_R\CC
 \longrightarrow\Ext^1(M,M)
 \longrightarrow\operatorname{Tor}_1^R(\Ext^2(M_R,M_R),\CC)
 \longrightarrow0
\]
の次元は $31,34,3$ である。
また $\Ext^2(M_R,M_R)\otimes_R\CC$ の次元は $12+3=15$ である。
従って障害加群のtorsionが、原点で増える三つの接方向を正確に記録している。
$a(h)$ が恒等的に零なら微分も零であり、自己Extの階数は $(1,34,15)$ のままである。
\end{remark}

\subsection{全体圏のHochschild cohomology}

\begin{theorem}[形式弧上のHochschild加群]\label{thm:gluing-arc-hh}
命題\ref{thm:gluing-arc-ext}の仮定のもとで
$\mathcal D_R=\langle\mathcal K_R,\mathcal P_R\rangle_{M_R}$ とする。
以下のHochschild cohomologyは $R$ 線形のものである。
\begin{equation}\label{eq:gluing-arc-hh}
 \HH^i(\mathcal D_R)=
 \begin{cases}
 R&i=0,\\
 R^{17}&i=2,\\
 R^{12}\oplus(R/(h^{\nu}))^{\oplus3}&i=3,\\
 0&\text{その他}.
 \end{cases}
\end{equation}
従ってgeneric fibreのHochschild次元は、次数 $0,2,3$ でそれぞれ $1,17,12$ である。
\end{theorem}

\begin{proof}
perfect bimoduleによるtriangular gluingはsmooth proper性を保つ。
\eqref{eq:full-gluing-triangle}は相対bar模型でも成り立ち
\cite[Theorem 7.7, Remark 7.8]{Kuz}、全ての複体は $R$-perfectである。
固定したブロックについて
\[
 \HH^*(\mathcal K_R)=R\text{ in degree }0\;\oplus\;R^{10}\text{ in degree }2,
 \qquad
 \HH^*(\mathcal P_R)=R\text{ in degree }0\;\oplus\;R^{24}\text{ in degree }1.
\]
次数零の射 $R^2\to\operatorname{End}(M_R)=R$ はscalarの差を取る全射である。
また
\[
 \HH^1(\mathcal P_R)=R^{24}\longrightarrow\Ext^1(M_R,M_R)=R^{31}
\]
は直和因子への単射であり、そのcokernelは $R^7$ である。
これは命題\ref{thm:full-gluing-base}の古典的な同変性から分かる。
左移動の微分は $b$ 因子上の可逆行列であり、$a$ 成分を持たない。
ここで導来正規形の同変性を仮定する必要はない。

残る問題はcharacteristic map
\begin{equation}\label{eq:gluing-characteristic-arc}
 \chi_R:\HH^2(\mathcal K_R)=R^{10}\longrightarrow\Ext^2(M_R,M_R)
\end{equation}
である。この射は次元の比較だけでは決まらない。
命題\ref{thm:fibre}の証明で、全NC族の形式基底を
\[
 \CC[[s_1,\ldots,s_{10},a_{ij}]]/(q)
\]
とし、残余圏への写像を $s$ への射影に直した。
$a=a(h)$ と矢のframingを固定し、
\[
 s_j=\epsilon c_j(h),\qquad\epsilon^2=0,\qquad c_j(h)\in R
\]
と引き戻す。この族は $\HH^2(\mathcal K_R)$ の任意の基底方向を
全体圏の一次変形へ持ち上げる。従って、その障害である $\chi_R$ は零である。

入力のNC族を形式族のまま扱う場合にも、この議論は厳密に行える。
各 $R/(h^n)[\epsilon]$ 上で整合するKodaira--Spencer類を取り、逆極限に移ればよい。
ここでは $\mathcal D_R\otimes_RR/(h^n)$ を、
$(s=0,a(h),b(h))$ に沿う既存のframed NC族と整合的に同一視し、
その同一視によって類を運ぶ。
$CC^\bullet(\mathcal D_R)$ は完備環上のperfect複体であり、
有限段階のcohomologyは有限長なのでMittag--Leffler条件を満たす。
Milnor完全列の $\varprojlim{}^1$ は零になり、
\[
 \HH^i(\mathcal D_R)\simeq
 \varprojlim_n\HH^i(\mathcal D_R\otimes_RR/(h^n))
\]
を得る。従って変動する圏について別の代数化定理を仮定せずに、
必要な $R$ 上の類と $\chi_R=0$ を得られる。

gluingの長完全列は
\[
 0\longrightarrow R^7\longrightarrow\HH^2(\mathcal D_R)
 \longrightarrow R^{10}\longrightarrow0,
 \qquad
 \HH^3(\mathcal D_R)\simeq\Ext^2(M_R,M_R)
\]
を与える。最初の列は $R$-加群として分裂する。
次数零の差写像と矢座標写像の単射性から $\HH^0=R,\HH^1=0$ も従い、
全ての主張を得る。
\end{proof}

特殊ファイバーに戻るときには
\[
 0\longrightarrow\HH^2(\mathcal D_R)\otimes_R\CC
 \longrightarrow\HH^2(X)
 \longrightarrow\operatorname{Tor}_1^R(\HH^3(\mathcal D_R),\CC)
 \longrightarrow0
\]
があり、$17+3=20$ を再現する。同様に次数3は $12+3=15$ となる。
このように自己ExtとHochschild cohomologyの両方で、
同じ三つの $h^{\nu}$-torsionが原点での跳びを記録する。

\subsection{滑らかな古典的moduliに残る12個の導来方向}

$F_\eta$ 上では $a(h)$ は非零の分解可能双ベクトルである。
$a_{01}\ne0$ のchartでは、三つのPfaffian式が正則列となり、
残る二つは関係式 $Aq=0$ によってその線形結合で表される。
三つの方程式を局所座標とし、対応するKoszul生成元との組を除くと、
$M_\eta$ におけるscalarを除いた導来局所模型は
\begin{equation}\label{eq:generic-gluing-derived}
 F_\eta[[u_1,\ldots,u_{31}]]\otimes\Lambda(v_1,\ldots,v_{12}),
 \qquad |v_i|=-1,\quad d=0
\end{equation}
となる。従って古典的moduliは31次元で滑らかであり、
commutator $L_\infty$ 代数はこの点でhomotopy abelianである。
12個の方向のうち10個は\eqref{eq:derived-gluing-normal}の $\rho$、
残り2個は5本のPfaffian式が余次元3を定めることによる。

これは、非零の $\Ext^2$ が必ず古典的な障害方程式を生むわけではないことを示す。
scalarを除いたvirtual dimensionは $34-15=19$ のままであり、
generic fibreでは古典的次元31との差が12個の導来方向として残る。
本節では全体圏のHochschild dg Lie代数のformalityや、
導来moduliの $\operatorname{PGL}(V)$ 商の積表示までは結論していない。


\section{安定対象モジュライ上の適合条件とBrill--Noether接錐}\label{sec:bn-algebraic}

\subsection{既知の10次元モジュライの中のgluing object}

第\ref{sec:gluing-objects}節の記号
$\mathcal K=\Ku(X)$、$Q=\Omega^1_{\PP^4}(1)|_X$、
$E=i^!Q$、$G=j^!\OO_X$ を引き続き用いる。
$\ell\subset X$ を直線とし、$v=[I_\ell]$ とおくと
\[
 \operatorname{ch}(v)=\left(1,0,-\frac{H^2}{3},0\right),\qquad
 \operatorname{ch}(E)=3\operatorname{ch}(v)=(3,0,-H^2,0).
\]
Serre-invariantな安定性条件 $\sigma$ を選び、対象の位相を固定する。
Feyzbakhsh--Liu--Zhang\cite[Section 3.2, Proposition 5.5]{FLZ}によれば、
$E$ は $\sigma$-stableであり、class $3v$ の他のstable対象は、
shiftを除いてGieseker-stableな層である。

\begin{theorem}[安定対象モジュライでの位置づけ]\label{thm:stable-ambient}
半安定対象の粗モジュライ空間
$M^{\mathrm{ss}}_\sigma(\mathcal K,3v)$ は既約かつ射影的であり、
そのstable部分
\[
 M^{\mathrm s}=M^{\mathrm s}_\sigma(\mathcal K,3v)
\]
は滑らかな10次元空間である。
$e=[E]$ はこのstable部分の点で、その完備局所環は
$\CC[[z_1,\ldots,z_{10}]]$ である。
階数3のnormalized Ulrich束は $M^{\mathrm{ss}}_\sigma(\mathcal K,3v)$ の
非空な稠密開部分をなし、$M^{\mathrm s}$ との交わりも非空かつ稠密である。
$E$ 自身はこのUlrich束の部分には属さない。
\end{theorem}

\begin{proof}
既約性、射影性とUlrich束の開埋め込みは
\cite[Section 6, Corollary 6.5]{FP}による。
stable部分の滑らかさは\cite[Theorem 1.2]{PY}であり、
その次元は
\[
 1-\chi(3v,3v)=1-9\chi(v,v)=10
\]
である。$e$ での完備局所環は命題\ref{thm:gluing-ext}からも直接得られる。
非空なUlrichの開部分と非空なstable部分は、既約な
$M^{\mathrm{ss}}_\sigma(\mathcal K,3v)$ 内で交わる。
一方、\eqref{eq:E-triangle}から
$\mathcal H^{-1}(E)=\OO_X(-1)$、$\mathcal H^0(E)=Q$ なので、
$E$ は層ではなく二項複体である。
\end{proof}

ここでnormalized Ulrich束とは、次数付き切断加群の生成元を次数1に置く
\cite{FP}の規約であり、Chern characterは $(3,0,-H^2,0)$ である。
生成元を次数0に置く規約とは $\OO_X(1)$ によるtwistだけ異なる。
また、射影的と述べたのは半安定対象の粗モジュライであり、
stable部分そのもののpropernessを主張してはいない。

\subsection{相対対象モジュライによる適合条件の代数化}

第\ref{sec:gluing-moduli}節では、$E$ の変形 $\widetilde E$ に対し、
$G$ が $\mathcal B_{\widetilde E}$ のperfect対象へ持ち上がる条件を
閉形式部分空間 $\mathfrak F_X\subset\Def_E$ として決定した。
この条件を実際の代数的な局所部分空間として実現する。

\begin{theorem}[代数的な適合部分空間]\label{thm:algebraic-compatibility}
$e\in M^{\mathrm s}$ の適当な\'etale近傍 $(U,u)$ と、
有限型の閉部分スキーム $(C,c)\subset(U,u)$ を選べて、次を満たす。
$U$ は滑らかな10次元スキームで、$\mathcal K$ 内のuniversal perfect対象
$\mathcal E_U$ を持ち、$\mathcal E_u\simeq E$ である。
$C$ は、相対圏 $\mathcal B_{\mathcal E_U}$ のsimple perfect対象の
モジュライのうち、$g=[G]$ を通る一つの\'etale局所的な枝である。
completed closed immersionは
\[
 \widehat C_c\longrightarrow\widehat U_u
 \ \simeq\
 \mathfrak F_X\longrightarrow\Def_E
\]
と一致する。従って適当な形式座標の下で
\[
 \widehat{\OO}_{C,c}
  \simeq\CC[[a_{ij}]]/(q_{ijkl}),\qquad
 \widehat{\OO}_{U,u}\simeq\CC[[a_{ij}]].
\]
特に $C$ は $c$ で7次元、余次元3であり、$T_cC=T_uU$ は10次元である。
\end{theorem}

\begin{proof}
\textbf{相対圏の構成。}
To\"en--Vaqui\'e\cite[Corollary 3.22]{TV}によれば、
negative自己Extを持たないsimple対象のstackは、
局所有限表示なalgebraic space上の、\'etale局所自明な
$\mathbb G_m$-gerbeである。
stable部分の\'etaleスキームchartを選び、このgerbeを自明化すれば、
universal対象 $\mathcal E_U$ を持つ $U$ が得られる。
安定性はここでは開条件であり、\cite[Proposition 6.3]{FP}の
Clifford加群のGiesekerモジュライによる記述を使っても確認できる。

$U$ をaffineに縮小する。
$\mathcal K$ のcompact generatorの自己準同型dg代数を $A$ とすれば、
$A_U=A\otimes_\CC\OO(U)$ は $\OO(U)$ 上smooth properである。
$\mathcal E_U$ に対応するperfectな $A_U$-加群を使い、
$\mathcal K_U$ と一つの例外対象 $Q$ を上三角に貼り合わせる。具体的には
\[
 \RHom_{\mathcal B_{\mathcal E_U}}(F,Q)
   =\RHom_{\mathcal K_U}(F,\mathcal E_U),\qquad
 \RHom_{\mathcal B_{\mathcal E_U}}(Q,F)=0,\qquad
 \RHom(Q,Q)=\OO_U
\]
と定め、pretriangulatedかつidempotent completeな包を取る。
この相対圏はsmooth properである。
propernessは各Hom複体の $\OO(U)$-perfect性による。
smoothnessは、上三角dg代数の対角bimoduleを、
二つの対角成分のperfectな対角bimoduleとperfectな混合加群から構成できる
ことによる。この構成は体に限らず可換基底環上で成立する。

\textbf{相対simple対象のモジュライ。}
\cite[Corollaries 2.13, 3.15, 3.17, 3.22]{TV}を基底環 $\OO(U)$ 上で適用する。
$\mathcal B_{\mathcal E_U}$ のsimple perfect対象のstackからscalar自己同型を
除いたalgebraic spaceを $\mathcal H$ と書くと、
$f:\mathcal H\to U$ は局所有限表示である。
$g=[G]$ における\emph{相対}接空間は
\[
 T_g(\mathcal H/U)=\Ext^1_{\mathcal B}(G,G)=0.
\]
残差体は $\CC$ なので、局所有限表示とこの消滅から、
$f$ は $g$ でunramifiedである。

$\mathcal H$ の $g$ を通る\'etaleスキームchartを取り、
unramifiedな射の\'etale局所構造
\cite[Lemma 41.17.1, Tag 04HH]{StacksUnramified}を適用する。
$U$ をさらに\'etale局所化すると、指定点を通る開いた枝 $C$ が
$U$ へ閉埋め込みになる。必要に応じて有限型の近傍へ縮小する。

\textbf{完備化の同定。}
この\'etale局所化は、指定した $\CC$-点での形式変形関手を変えない。
$g$ での $\mathcal H$ の完備化は、
$\mathcal E_A$ に沿って $G$ を
$\mathcal B_{\mathcal E_A}$ のsimple perfect対象へ持ち上げる変形関手である。
special fiberでnegative自己Extがなく $\operatorname{End}(G)=\CC$ なので、
局所Artin基底上のperfectな持ち上げはこのsimple locusに入る。
第\ref{sec:gluing-moduli}節によれば、その持ち上げは存在すれば同型類が一意で、
存在条件はちょうど $\mathfrak F_X$ である。
従って完成された閉埋め込みは
$\mathfrak F_X\subset\Def_E$ と一致する。
\end{proof}

$C$ 上のuniversalな $G$ 自身が必要なら、
simple対象の $\mathbb G_m$-gerbeをさらに\'etale局所自明化すればよい。
この命題は指定対象 $G$ を通る局所的な枝を構成している。
相対モジュライの任意の有限型部分の像はconstructibleだが、
その全体のpropernessや大域的な像の閉性はここでは結論しない。
また、一般の形式イデアルの代数化を仮定したのではなく、
実際の相対対象モジュライを先に構成した点がこの証明の要点である。

\subsection{Brill--Noether incidenceの局所構造}

$p\in X$ に対し $P_p=i^*\OO_p\in\mathcal K$ とおく。
相対射影を使えば、これらは $X$ を基底とするperfectな族になる。
$F\in\mathcal K$ に対し、随伴とSerre双対性から
\[
 \Ext^3_{\mathcal K}(P_p,F)=\Ext^3_X(\OO_p,F)
 \simeq\operatorname{Hom}_X(F,\OO_p\otimes\omega_X)^*.
\]
従って
\[
 b_p(F):=\dim_\CC\Ext^3_{\mathcal K}(P_p,F)
       =\dim_\CC\operatorname{Hom}_X(F,\OO_p).
\]
$E$ の場合は、任意の $p$ について
\[
 \dim\Ext^2_{\mathcal K}(P_p,E)=1,\qquad b_p(E)=4.
\]
一方、$F$ が階数3のtorsion-free層なら
\[
 b_p(F)\ge4
 \quad\Longleftrightarrow\quad
 F\text{ は }p\text{ で局所自由でない}.
 \tag{BN}\label{eq:bn-nonfree}
\]
実際、$b_p(F)$ は局所加群 $F_p$ の最小生成元数である。
3個の生成元による $\OO_{X,p}^{\,3}\twoheadrightarrow F_p$ の核は、
torsion-free加群の階数零の部分加群なので零となる。

\begin{theorem}[滑らかなincidenceとその法束]\label{thm:bn-incidence}
命題\ref{thm:algebraic-compatibility}の $U$ を縮小し、
$X\times U$ 上のcohomology-jump部分スキーム
\[
 \mathcal I
   =\{(p,t)\in X\times U:b_p(\mathcal E_t)\ge4\}
\]
をdeterminantalなスキーム構造で定める。
すると $\mathcal I$ は $X\times\{u\}$ の近傍で滑らかな9次元スキームであり、
$\mathcal I_u=X$ はスキームとして成立する。
さらに、固定したtrace線の自明化の下で
\[
 \mathcal N_{X/\mathcal I}
 \simeq
 \ker\bigl(\Ext^1(E,E)\otimes\OO_X\longrightarrow Q(1)\bigr)
 \simeq\Lambda^2Q.
\]
\end{theorem}

\begin{proof}
$E$ の $p$ における導来ファイバーは、次数 $-1$ に1次元、次数0に4次元を持つ。
perfect複体の局所表示から可逆な部分を消去すると、
$(p,u)$ の近傍でuniversal対象は
\[
 [\,\OO\xrightarrow{\,d\,}\OO^{\,4}\,]
 \qquad\text{（次数 }-1,0\text{）}
\]
と表示できる。この表示で $\mathcal I$ は $d$ の4成分の零点である。
special fiberではこれらの成分が全て零なので、$\mathcal I_u=X$ となる。

基底方向の微分はcohomologyのKodaira--Spencer写像
\[
 \Ext^1(E,E)\longrightarrow
 \operatorname{Hom}(\mathcal H^{-1}(E),\mathcal H^0(E))
 =H^0(Q(1))
\]
を $p$ で評価したものである。
この写像は命題\ref{thm:gluing-ext}で計算した同型と一致する。
実際、\eqref{eq:E-triangle}の射を
$\iota:\OO_X(-1)[1]\to E$、$\beta:E\to Q$ と書くと、
その写像は
\[
 \alpha\longmapsto
 \beta[1]\circ\alpha\circ\iota:
 \OO_X(-1)[1]\longrightarrow Q[1].
 \tag{KS}\label{eq:bn-edge-map}
\]
随伴 $\RHom(E,E)\simeq\RHom(E,Q)$ は $\beta$ との合成で与えられ、
三角へ $\RHom(-,Q)$ を適用する次の写像は $\iota$ との合成なので、
これは既に用いた自己Extの同型そのものである。

評価列の外積から
\[
 0\longrightarrow\Lambda^2Q
 \longrightarrow\Lambda^2V\otimes\OO_X
 \longrightarrow Q(1)\longrightarrow0
\]
を得る。従って $H^0(Q(1))=\Lambda^2V$ の評価は各点で全射である。
$d$ の4成分の微分が全射となるので、
$\mathcal I$ は13次元の $X\times U$ 内で滑らかな余次元4の部分スキームである。
法束の計算もこの接空間列と外積の評価列から従う。
$X$ のpropernessにより、必要な縮小を $U$ 上で一斉に行える。
\end{proof}

\subsection{Yoneda積による一次条件とFermatでの接錐}

$\alpha\in\Ext^1(E,E)$ に対して、Yoneda積
\[
 m_{\alpha,p}:\Ext^2_{\mathcal K}(P_p,E)
       \longrightarrow\Ext^3_{\mathcal K}(P_p,E),
 \qquad b\longmapsto\alpha[2]\circ b
\]
を考え、その零点を
$Z_\alpha=\{p\in X:m_{\alpha,p}=0\}$ とする。
これは $\mathcal K$ 内の対象と射だけで書かれた一次Brill--Noether条件である。
局所双対性と導来ファイバーの関手性により、
$m_{\alpha,p}$ は\eqref{eq:bn-edge-map}の $p$ における評価に一致する。

\begin{lemma}[一次条件の階数による記述]\label{lem:bn-rank}
\eqref{eq:bn-edge-map}によって $\alpha$ を
$a\in H^0(Q(1))=\Lambda^2V$ と同一視し、対応する切断を $s_a$ と書く。
$X\subset\PP(V^*)$ の規約の下で
\[
 Z_\alpha=Z(s_a)
   =X\cap\PP\bigl(\ker(a:V^*\to V)\bigr).
\]
$a\ne0$ の階数が2ならこの零点は平面三次曲線であり、
階数が4なら高々1点である。従って集合として
\[
 \{a\in\Lambda^2V:\dim Z(s_a)\ge1\}
   =C(\Gr(2,V)).
 \tag{*}\label{eq:first-bn-cone}
\]
ここでは $a=0$ も右辺に含める。
\end{lemma}

\begin{proof}
$p\in\PP(V^*)$ を評価functionalと見れば、$s_a(p)$ は contraction
$\iota_p(a)$ である。5次元の交代形式の非零の階数は2または4である。
階数2ならkernelの射影化は平面であり、階数4なら1点である。
滑らかなcubic threefoldは平面を含まないので、前者の交叉は次元1の平面三次曲線となる。
例えば平面の理想を $(l_1,l_2)$ として $F=l_1q_1+l_2q_2$ と書けたなら、
平面上の $q_1,q_2$ の共通零点で全ての偏微分が零になり、滑らかさに反する。
最後に階数2以下の交代形式は分解可能な外積であり、
その錐はPl\"ucker埋め込みのGrassmannian錐である。
\end{proof}

一般の $X$ に対して、上の $\Lambda^2V$ との同一視と、
\eqref{eq:gluing-tangent}を通じてPoisson空間から得る同一視の比較には、
次元の一致以上の議論が必要である。
ここではまずFermatの場合の群作用による証明を与える。
全ての滑らかな $X$ に対する比較は、第\ref{sec:general-tangent-comparison}節で証明する。

\begin{theorem}[FermatでのBrill--Noether接錐]\label{thm:fermat-bn-cone}
$X_0=V(x^3+y^3+z^3+u^3+v^3)$ とする。
命題\ref{thm:algebraic-compatibility}の適合部分空間 $C$ の接錐は、
\eqref{eq:bn-edge-map}の座標でPl\"ucker錐である。
その台は圏論的に
\[
 \left|\operatorname{TC}_cC\right|
 =\{\alpha\in\Ext^1(E,E):\dim Z_\alpha\ge1\}
\]
と特徴付けられる。これは10次元の接空間全体ではなく、7次元の接錐についての主張である。
\end{theorem}

\begin{proof}
被覆群
$\widetilde\Gamma=\mu_3^5\rtimes S_5$ の線形作用で計算し、
最後に $\Gamma=(\mu_3^5/\mu_{3,\mathrm{diag}})\rtimes S_5$ へ降ろす。
$V$ の座標を $x_0,\ldots,x_4$ とし、
$g^*x_i=\zeta_i x_i$ とする。Fermat三次式はこの作用で不変なので、
同変なadjunctionは
\[
 \omega_{X_0}\simeq\OO_{X_0}(-2)\otimes\det V
\]
を与える。よってPoisson側の自然な表現は
\[
 H^0(\Lambda^2T_{X_0})
 \simeq H^0(\Omega^1_{X_0}(2))\otimes(\det V)^*
 \simeq\Lambda^2V\otimes(\det V)^*.
\]
最後の同型は、conormal列
$0\to\OO_{X_0}(-1)\to Q(1)\to\Omega^1_{X_0}(2)\to0$
と $\OO_{X_0}(-1)$ のコホモロジーの消滅から従う。

対象側ではtrace線
\[
 L=\Ext^2(Q,\OO_{X_0}(-1))
   \simeq H^3(\OO_{X_0}(-2))
   \simeq(\det V)^*
\]
を保持する。\eqref{eq:E-triangle}の左端は正確には
$\OO_{X_0}(-1)[1]\otimes L^*$ なので、
\[
 \Ext^1(E,E)
 \simeq H^0(Q(1))\otimes L
 \simeq\Lambda^2V\otimes(\det V)^*.
\]
従って二つの側のadjunctionおよびtraceのcharacterは一致する。
各 $(i,j)$ のweight lineのcharacterは
\[
 \chi_{ij}(\zeta)
  =\frac{\zeta_i\zeta_j}{\zeta_0\zeta_1\zeta_2\zeta_3\zeta_4}.
\]
対角 $\mu_3$ 上では $\zeta^{\,2-5}=1$ なので、この表現は $\Gamma$ に降りる。
10個のweightは相異なり、$S_5$ はその直線を推移的に置換する。
従ってこの群作用と可換する自己準同型はscalarだけである。

\eqref{eq:gluing-tangent}は自然な同型なので、
上の二つの同一視の下では非零scalar倍となる。
命題\ref{thm:object-cone}および命題\ref{thm:algebraic-compatibility}の
接錐が、この座標でもPl\"ucker錐であることが従う。
その台の記述は補題\ref{lem:bn-rank}である。
\end{proof}

右辺のHilbert多項式の跳躍条件から別途定めたスキーム構造まで
Pl\"uckerイデアルと一致するとは主張していない。
また、$C$ 自体の全ての点を\eqref{eq:bn-nonfree}の条件で特徴付けたわけではない。
一般の $X$ についてHochschild写像と\eqref{eq:bn-edge-map}を比較する証明は、
命題\ref{thm:general-tangent-comparison}で与える。

\subsection{明示的な固定残余圏族から得る層の変形}

\begin{theorem}[非定数な階数3の層の族と特殊化の制限]\label{thm:gluing-sheaf-arc}
第9節の固定残余圏族 $A'_h$ は、$R=\CC[[h]]$ 上で、
固定した $X_0$ の $\mathcal K$ 内のperfect対象の非定数な族
$\mathcal E_R$ を与える。special fiberは $E$ であり、
その幾何学的generic fiberはGieseker-stableなtorsion-free層で、
\[
 \rk=3,\qquad c_1=0,\qquad c_2=H^2,\qquad c_3=0
\]
を持つ。
generic fiberの非局所自由な点の集合の水平閉包を
$S\subset X_0\times\operatorname{Spec}R$ とすれば、集合として
\[
 S_0\subset
 X_0\cap\{u=v=0\}
 =V(x^3+y^3+z^3)\subset\PP^2.
\]
\end{theorem}

\begin{proof}
既存の形式族は
$\Def_E\simeq\widehat U_u$ への形式弧を与える。
$U$ をaffineに取れば、対応する連続準同型
$\widehat{\OO}_{U,u}\to R$ を $\OO(U)$ に制限して
$\operatorname{Spec}R\to U$ を得る。
universal perfect対象を引き戻した $\mathcal E_R$ は、
全てのArtin切断で元の形式族と一致する実際の $R$ 上の族である。

$A'_h$ の一次Poisson類は非零であり、
\eqref{eq:gluing-tangent}によりこの弧の接ベクトルも非零である。
従ってgeneric pointの像は $e$ ではない。
$U$ をstable locus内に取り、universal対象の局所Tor-amplitudeを
$[-1,0]$ に縮小しておく。
さらに $U$ を縮小し、$e$ の上の点を $u$ だけにしておく。
\cite[Proposition 5.5]{FLZ}によれば、$U\setminus\{u\}$ の閉点は
Gieseker-stableな層を表す。このamplitudeと正の階数によりshiftは零である。
層であることとGieseker-stabilityの開性、および有限型 $\CC$-スキームの
Jacobson性から、この結論は $U\setminus\{u\}$ の全ての点に成り立つ。
従って非定数な弧のgeneric fiberについても同じ結論が従う。
purityと正の階数からtorsion-freeであり、Chern類はclass $3v$ から得られる。

命題\ref{thm:fermat-bn-cone}の座標比較により、この弧の接方向は、
非零scalar倍を除いて $a=u\wedge v$ である。
命題\ref{thm:bn-incidence}の局所二項表示では、微分は
$h\,s_a+O(h^2)$ と書ける。
generic fiberの非局所自由な部分はその微分の4成分で定義される。
水平閉包のイデアルはこれらの成分を $h$ で割ったものを含むので、
$h=0$ で $s_a$ は零になる。従って $S_0\subset Z(s_a)$。
最後に
\[
 \iota_p(u\wedge v)=u(p)v-v(p)u
\]
から $Z(s_a)=X_0\cap\{u=v=0\}$ を得る。
\end{proof}

この段階では非局所自由部分の特殊化の包含までを得た。
第\ref{sec:evaluation-gluing}節では、この族のgeneric fiberを評価核として記述する。
それにより、非局所自由部分が固定された平面三次曲線そのものであり、
局所自由でもUlrichでもないことを証明する。

\section{非被約点からの平行移動パラメータの復元}\label{sec:reconstruction}

第\ref{sec:points}節の局所代数 $J_a$ が記録する情報を、
Sklyaninの平行移動パラメータと比較する。
まず代数閉体 $k$、$\operatorname{char}k=0$ 上の幾何学的ファイバーを扱い、
最後に $\CC[[h]]$ 上の形式族へ戻る。
以下の一般ファイバーについての個数は幾何学的な個数であり、
$\CC((h))$ 上の非自明なformまで $j$ だけで分類するという主張ではない。
本節では $t=a^3$ と書き、第9節の三乗根は $r$ と書く。

\subsection{Jacobian環のassociated formとしての三次形式}

\[
 f_a=x^3+y^3+z^3-3axyz,\qquad
 J_a=k[x,y,z]/(x^2-ayz,y^2-azx,z^2-axy)
\]
とおく。$a^3\ne1$ のとき $f_a$ は滑らかな平面三次曲線を定め、
$J_a$ はそのJacobian環である。

\begin{theorem}[associated formと逆数関係]\label{thm:fat-associated}
$a^3\ne1$ とする。局所代数 $J_a$ のsocle三次形式
\[
 \Phi_a(X,Y,Z)=a(X^3+Y^3+Z^3)+6XYZ
\]
は、$f_a$ のMacaulay inverse system、従ってそのassociated formである。
さらに $a\ne0$、$a^3\ne-8$ のとき、通常の楕円曲線の $j$ の規約で
\begin{equation}\label{eq:associated-reciprocal-j}
 j(V(f_a))=\frac{27a^3(a^3+8)^3}{(a^3-1)^3},\qquad
 j(V(f_a))\,j_{\mathrm{fat}}(a)=1728^2.
\end{equation}
\end{theorem}

\begin{proof}
$\partial_x f_a=3(x^2-ayz)$ とその巡回置換からJacobian環の表示を得る。
第\ref{sec:points}節の積表を使い、$T=xyz$ とおくと
\[
 (Xx+Yy+Zz)^3=\Phi_a(X,Y,Z)T.
\]
また $(\partial_X^2-a\partial_Y\partial_Z)\Phi_a=0$ とその巡回置換が成り立つ。
$\Phi_a$ のapolar algebraのHilbert関数は $(1,3,3,1)$ であり、
$J_a$ も同じ長さ8を持つので、この三つの二次作用素がそのannihilatorを生成する。
これがinverse systemとしての同一視を与える。
$j$ の式はHesse三次曲線の公式と\eqref{eq:fat-j}から従う。
これはAlper--Isaev--Kruzhilin\cite[Section 3.2, (3.11), (3.18)]{AIK}の
associated-form involutionの具体例である。
同論文の規格化 $J=j/1728$ では、この関係は $J\mapsto1/J$ と書かれる。
\end{proof}

associated formはHessianとは区別される\cite[Remark 3.1]{AIK}。
実際、直接微分すると
\begin{equation}\label{eq:fat-hessian-distinction}
 \operatorname{Hess}(f_a)
 =-54a^2(x^3+y^3+z^3)+54(4-a^3)xyz
\end{equation}
となり、$\Phi_a$ とは異なる。
この逆数関係を、固定曲線 $E_0$ と $V(\Phi_a)$ の間の
3-isogenyと解釈することはできない。
後で構成する3-isogenyは $V(\Phi_a)$ 自身を始域とする。

\subsection{固定Fermat曲線の次数2の商}

Sklyaninの平行移動の原点を
\[
 E_0=V(x^3+y^3+z^3),\qquad O=[1:-1:0]
\]
とする。第9節のSklyaninパラメータは
\[
 p(h)=[1:-r:-h],\qquad r^3+h^3=1,
 \qquad a=\frac{1-r}{h}=-\frac{x+y}{z}
\]
であり、点曲線の自己同型は $p(h)$ による平行移動である
\cite[Section 2.3]{WWY}。

\begin{theorem}[markedな係数からの復元]\label{thm:translation-double-cover}
有理関数 $a=-(x+y)/z$ は次数2の射
$E_0\to\PP^1_a$ に延長し、$E_0/\{\pm1\}$ を与える。
その分岐値は $a=0$ および $a^3=4$ である。
一般の $a$ から平行移動点を復元する式は
\begin{equation}\label{eq:translation-quadratic}
 (1-a^3)h^2+3a^2h-3a=0,\qquad r=1-ah.
\end{equation}
従って $a$ は平行移動点を符号まで決定する。
\end{theorem}

\begin{proof}
$\omega$ を原始三乗根とし、$Q_\pm=[1:-\omega^{\pm1}:0]$ とおく。
直線 $x+y=0$ は $O$ における変曲接線であり、
$z=0$ の交点は $O,Q_+,Q_-$ なので
\[
 \operatorname{div}(a)=2O-Q_+-Q_-.
\]
従って $a$ は次数2である。
$[x:y:z]\mapsto[y:x:z]$ は $a$ と $O$ を固定する非自明なinvolutionなので、
楕円曲線の $[-1]$ に等しい。
その固定点は $O$ と $x=y$ 上の三点であり、後者では $a^3=4$ となる。
また $r=1-ah$ を $r^3+h^3=1$ に代入し、一般点で $h$ を消去すると
\eqref{eq:translation-quadratic}を得る。その判別式は
\[
 9a^4+12a(1-a^3)=3a(4-a^3).
\]
二根の交換は $(r,h)\mapsto(1/r,-h/r)$、すなわち上の $[-1]$ である。
\end{proof}

ここでの分岐値 $a^3=4$ と、socle三次曲線が特異になる $a^3=-8$ は異なる。
前者は平行移動点の符号を忘れることに由来し、後者は局所代数の三次形式に由来する。

\subsection{粗い局所代数からは四つのSklyanin同型類が残る}

\begin{theorem}[幾何学的一般点における有限多価の復元]\label{thm:four-sklyanin-classes}
$t=a^3$ とおくと、$E_0\to\PP^1_t$ は
$\operatorname{Aut}(E_0,O)\simeq\mu_6$ による商である。
socle三次曲線の不変量は
\begin{equation}\label{eq:fat-j-degree-four}
 \mathcal J(t)=\frac{110592(t-1)^3}{t(t+8)^3}
\end{equation}
という次数4の写像であり、合成 $E_0\to\PP^1_j$ の次数は24である。

同じ埋め込まれた曲線 $E_0$ に付随するSklyanin代数を $S(p)$ と書く。
$p,q\in E_0\setminus E_0[3]$ に対して
\begin{equation}\label{eq:fixed-sklyanin-isom}
 S(p)\simeq_{\mathrm{gr}}S(q)
 \quad\Longleftrightarrow\quad
 q\in\operatorname{Aut}(E_0,O)\cdot p.
\end{equation}
従って、一般の非被約点の局所代数の同型類からは、
平行移動点として24個、Sklyanin代数の次数付き同型類として4個の候補が残る。
\end{theorem}

\begin{proof}
$[x:y:z]\mapsto[y:x:z]$ と $[x:y:z]\mapsto[x:y:\omega z]$ は
$\operatorname{Aut}(E_0,O)$ を生成し、$a$ に対してそれぞれ
$a\mapsto a$、$a\mapsto\omega^{-1}a$ と作用する。
$t=a^3$ は不変で次数6なので、この群の商座標である。
\eqref{eq:fat-j-degree-four}の分子と分母は互いに素であり、その最大次数は4である。

幾何代数の同型判定\cite[Theorem 2.8]{IMRelations}によれば、
次数付き同型は点曲線の平行移動自己同型の射影的共役に対応する。
$\OO_{E_0}(1)=\OO_{E_0}(3O)$ なので
\[
 \operatorname{Aut}(\PP^2,E_0)
 =E_0[3]\rtimes\operatorname{Aut}(E_0,O)
\]
である\cite[Theorem 4.12]{IMRelations}。
平行移動を $T_p$ と書くと、$T_s\alpha$ による共役は
\[
 (T_s\alpha)T_p(T_s\alpha)^{-1}=T_{\alpha(p)}
\]
だから\eqref{eq:fixed-sklyanin-isom}を得る。
これは\cite[Theorem 4.16 and Lemma 4.17]{IMRelations}の $i=0$ の場合でもある。

最後に、$\Phi_a$ が滑らかな一般点では、$J_a$ の局所代数としての同型類と
$\Phi_a$ の射影同値類は同値である。
一方向は極大イデアルの次数付き環と三重積を取り、逆方向はinverse systemを使う。
滑らかな平面三次曲線の射影同値類は、代数閉体上で $j$ により決まる。
実際、同じ $j$ を持つ楕円曲線間の同型を平行移動で調整すれば、
二つの次数3の線形束を一致させられる。
これは $[3]:E\to E$ の全射性による。
従って一般の局所代数の同型類は $\mathcal J$ の一般ファイバーを与え、
そこには四つの $t$、それぞれに六つの平行移動点がある。
\end{proof}

候補は具体的に
\[
 j\,t(t+8)^3-110592(t-1)^3=0
\]
を解いて求められる。さらに一つの $t$ を選ぶと
\[
 (t-1)r^2+(t+2)r+(t-1)=0,\qquad h^3=1-r^3
\]
から六つの点を復元できる。
最初の二次式の根は互いに逆数で、判別式は $3t(4-t)$ である。
ここまでの分類は三生成元のSklyanin代数についてのものである。
第\ref{sec:five-generator}節では、これを五生成元の三次元代数 $A'_h$ の
次数付き同型分類へ拡張する。

\subsection{socle三次曲線の3-isogenyと \texorpdfstring{$X_0(3)$}{X0(3)}}

$a\ne0$ として $b=-2/a$ と置けば、$V(\Phi_a)$ は
\[
 C_b:\ x^3+y^3+z^3-3bxyz=0
\]
と書ける。以下 $b^3\ne1$ とする。

\begin{theorem}[明示的な巡回3-isogeny]\label{thm:fat-three-isogeny}
$C_b$ の自己同型
\[
 \delta[x:y:z]=[x:\omega y:\omega^2z]
\]
は非零の3-torsion点による平行移動であり、その商は
\begin{equation}\label{eq:fat-isogeny-quotient}
 [x:y:z]\longmapsto[X:Y:Z:T]=[x^3:y^3:z^3:xyz],\qquad
 X+Y+Z=3bT,\quad XYZ=T^3
\end{equation}
で与えられる。商曲線の $j$ は
\begin{equation}\label{eq:fat-isogeny-j}
 j_3(t)=-\frac{110592(t+9)^3}{t^3(t+8)}.
\end{equation}
幾何学的一般点で $t$ は、楕円曲線と選ばれた位数3の巡回部分群の組の
粗モジュライ $X_0(3)$ の座標となる。
$s=-27t/(t+8)$ と書けば、二つの $j$-写像は
\begin{equation}\label{eq:x03-j-pair}
 j_{\mathrm{fat}}=\frac{(s+3)^3(s+27)}s,\qquad
 j_3=\frac{(s+27)(s+243)^3}{s^3}.
\end{equation}
双対isogenyへ移るFricke involutionは
\begin{equation}\label{eq:fat-fricke}
 s\longmapsto\frac{729}s,
 \qquad\text{すなわち}\qquad t\longmapsto-t-8
\end{equation}
である。
\end{theorem}

\begin{proof}
$\delta$ の射影空間内の固定点は三つの座標頂点だけで、どれも $C_b$ 上にはない。
楕円曲線上の有限位数の自己同型が平行移動でなければ、
$1-\alpha$ の全射性から固定点を持つ。従って $\delta$ は位数3の平行移動である。
\eqref{eq:fat-isogeny-quotient}の一般ファイバーは $\delta$ の三点軌道だから、
これはその巡回部分群による商である。
$Z$ を消去すると $XY(3bT-X-Y)=T^3$ を得る。
$x_W=-T$、$y_W=X$、$z_W=Y$ とすればWeierstrass形は
\[
 y_W^2z_W+3b x_Wy_Wz_W+y_Wz_W^2=x_W^3
\]
となり、その不変量は
\[
 c_4=9b(9b^3-8),\qquad\Delta=27(b^3-1).
\]
従って $j=c_4^3/\Delta$ に $b^3=-8/t$ を代入して
\eqref{eq:fat-isogeny-j}を得る。
商曲線の原点を始域の原点の像に選べば、この次数3の射はisogenyである。
Weierstrass表示で選んだ変曲点を原点に変更しても $j$ は変わらない。

$a\mapsto\omega a$ に対して $b\mapsto\omega^2b$ となり、
$[x:y:z]\mapsto[\omega x:y:z]$ が曲線と巡回対称性を同一視する。
従ってこの組の粗モジュライは $t=a^3$ に降下する。
一般の楕円曲線の $E[3]\simeq(\mathbb Z/3)^2$ には位数3の部分群が四つあり、
一般の自己同型 $\pm1$ はそれぞれを保つ。
従って $X_0(3)\to\PP^1_j$ の次数は4である。
$t$ の写像も次数4なので、$\PP^1_t$ からこの粗モジュライへの写像は次数1となる。

\eqref{eq:x03-j-pair}は直接代入で従う。
また $s\mapsto729/s$、すなわち $t\mapsto-t-8$ は二つの $j$ を交換する。
幾何学的一般点では異なる位数3の部分群から得られる商の $j$ は異なるので、
これが双対isogenyに対応する。
この最後の事実は、一般の楕円曲線の自己準同型環が $\mathbb Z$ であることからも分かる。
二つの異なる商が同型なら、一方のisogenyと他方の双対を合成して得る次数9の
自己準同型は $\pm[3]$ であり、双対を取ることで二つの核が一致してしまう。
\end{proof}

Hesse pencilの共通する3-torsionとそのsymplectic構造については
\cite[Theorem 1.1 and Sections 2--4]{ATT}を参照されたい。
ここで $X_0(3)$ が記録する巡回部分群は、socle三次曲線上に選んだ対称性である。
固定Sklyanin曲線 $E_0$ は、平行移動パラメータ $p$ を動かすための別の曲線である。

\subsection{形式近傍と、一般のSklyaninパラメータ}

$O$ の形式近傍では
\[
 t=a(h)^3=\frac{h^6}{27}+\frac{h^9}{27}+O(h^{12}),\qquad
 \frac1{\mathcal J(t)}=-\frac{t}{216}+O(t^2).
\]
後者は $t$ に関して逆形式級数を持つ。
従って $O$ に近づく枝を指定すれば、$1/j_{\mathrm{fat}}$ から $t$ は一意に復元される。
大域的な四つの $t$ の曖昧さはこの形式近傍では消え、形式商の座標環は
\[
 \widehat{\OO}_{E_0,O}^{\mu_6}=\CC[[t]]
\]
となる。整域上の非零の形式弧を復元するときに残る曖昧さは $\mu_6$ の作用である。
ただし粗い商のArtin値点は、群作用による軌道の関手とは異なる。
例えば $\CC[\epsilon]/(\epsilon^2)$ 上の $h=0$ と $h=\epsilon$ は
同じ $t=0$ を与えるが、同じ $\mu_6$ 軌道には属さない。

一般のSklyaninパラメータ $[A:B:C]$ でも、$C=-1$ と規格化して
\[
 d=A+B,\qquad w=AB,\qquad
 \tau=\frac{A^3+B^3+C^3}{ABC}
\]
とおけば、$A^3+B^3=d^3-3wd$ より
\begin{equation}\label{eq:general-sklyanin-recovery}
 w=\frac{1-d^3}{\tau-3d},\qquad
 A,B:\ U^2-dU+\frac{1-d^3}{\tau-3d}=0
\end{equation}
を得る。点曲線は $x^3+y^3+z^3-\tau xyz=0$、
対角上の二次関係は $J_d$ のものである。
従ってmarkedな $(\tau,d)$ は、一般点で $A\leftrightarrow B$ までパラメータを復元する。
この二次式の判別式は
\[
 \frac{d^3+\tau d^2-4}{\tau-3d}
\]
である。元のFermat族では $d=0$、$\tau=-\lambda^2$、
固定残余圏族では $\tau=0$、$d=a(h)$ となる。
両族で異なる不変量が変化することを、この座標でまとめて記述できる。

\subsection{三つの検出次数：貼り合わせ、局所環、粗モジュライ}

第9節の固定残余圏族を $R=\CC[[h]]$ 上で考える。
第13節の水平な非被約成分の局所環は、$a=a(h)$ とした $J_a$ である。
この成分のflat性と、相対point-module truncation全体の非flat性は区別する。

\begin{theorem}[固定残余圏族で最初に見える次数]\label{thm:orders-126}
この族に沿って、次の三つの次数が得られる。
\begin{enumerate}
 \item gluing object $E_h$ の変形の一次類は非零である。
 \item $J_{a(h)}$ のmarkedな可換局所代数としての変形は
 $\bmod\,h^2$ では定数であるが、$\bmod\,h^3$ では定数でない。
 \item socle三次形式の粗モジュライの局所座標 $\tau=1/j_{\mathrm{fat}}$ は
 \[
   \tau=-\frac{h^6}{5832}+O(h^9)
 \]
 であり、最初の非零項は6次である。
\end{enumerate}
\end{theorem}

\begin{proof}
第一の主張は第11節の接写像の同型と、非零の一次Poisson類による。
第二を確認する。まず $a=\epsilon$、$\epsilon^2=0$ とした
$J_\epsilon$ のKodaira--Spencer類は、complete intersection
$J_0=\CC[x,y,z]/(x^2,y^2,z^2)$ の余接複体によって
\[
 T^1_{J_0}=\operatorname{coker}
 \bigl[J_0^3\xrightarrow{\operatorname{diag}(2x,2y,2z)}J_0^3\bigr]
\]
における $(-yz,-zx,-xy)$ で表される。
$yz\notin xJ_0$ なのでこれは非零である。

$a(h)=h^2/3+O(h^5)$ であるから、$\bmod\,h^2$ では定数である。
さらに、$\bmod\,h^3$ で定数族へのmarkedな同型があったと仮定する。
定数族の生成元 $x$ の像は
$x+h f+h^2g$、$f,g\in J_0$ と書ける。
その平方の $h$ の係数は $2xf$ なので、$f\in\operatorname{ann}_{J_0}(x)=xJ_0$。
従って $f^2=0$ である。$h^2$ の係数は
\[
 \frac13yz+2xg
\]
となり、$yz\notin xJ_0$ に反する。よって定数化できない。
第三は\eqref{eq:fat-j}の逆数と $a(h)$ の展開から従う。
\end{proof}

従って、同じ非可換族に対しても、貼り合わせ対象、水平な非被約成分の局所環、
そのsocle三次形式の粗モジュライは、それぞれ1次・2次・6次から変化を検出する。
これは普遍的な写像 $\Def_E\to\operatorname{Def}_{J_0}$ を構成したという意味ではない。
ここで比較したのは、次数付き代数模型も指定された一本の形式族である。

\section{一般のcubic threefoldにおける二つの接写像の比較}
\label{sec:general-tangent-comparison}

第\ref{sec:bn-algebraic}節でFermatの場合に証明した接錐の同定を、
全ての滑らかな三次三次元多様体へ拡張する。
ここで比較するのは、Hochschild gluingから得る
\[
 \kappa_F:H^0(X,\Lambda^2T_X)\xrightarrow{\sim}\Ext^1(E,E)
\]
と、cohomologyのKodaira--Spencer写像
\[
 e_F:\Ext^1(E,E)\xrightarrow{\sim}H^0(Q(1))\otimes L_F,
 \qquad L_F=\Ext^2(Q,\OO_X(-1))
\]
である。後者は\eqref{eq:bn-edge-map}のtrace線を保持した形である。
conormal列とadjunctionが与える同型を
\[
 \lambda_F:H^0(Q(1))\otimes L_F
       \xrightarrow{\sim}H^0(X,\Lambda^2T_X)
\]
と書く。$L_F\simeq H^3(X,\OO_X(-2))$ の規約を一貫して用いる。

\begin{theorem}[一般の $X$ における接座標の比較]
\label{thm:general-tangent-comparison}
HKR、連結射、およびadjunctionの符号・正規化を固定すると、
$F$ に依存しない定数 $c\in\CC^\times$ が存在して
\[
               \lambda_F\circ e_F\circ\kappa_F
                    =c\,\id_{H^0(\Lambda^2T_X)}
\]
となる。従って任意の滑らかなcubic threefoldに対して、
命題\ref{thm:algebraic-compatibility}の適合部分空間 $C$ の接錐は、
edge写像 $e_F$ の座標でもPl\"ucker錐である。
特に、その台は
\[
 |\operatorname{TC}_cC|
 =\{\alpha\in\Ext^1(E,E):\dim Z_\alpha\ge1\}
\]
である。
\end{theorem}

証明では、滑らかな三次式の開集合だけで同変性を使うことはしない。
その開集合上には判別式を分母に持つ正則関数が存在するためである。
まず、特異なものも含む全ての非零三次式の上に比較写像を延長する。

\subsection{全非零三次式上の相対gluing}

$S=\Sym^3V\setminus\{0\}$ とし、
$p:\mathcal X\to S$ を $\PP(V^*)\times S$ 内の普遍三次超曲面とする。
この射はflat、proper、Gorensteinである。
以下のHochschild複体と自己Extは全て基底 $S$ に相対的なものとする。
以下の $\Perf(\mathcal X/S)$ は、$\mathcal X$ 上でperfectな複体の圏を
$S$-線形dg圏として見るという意味であり、単に $S$-perfectな複体全体を意味しない。

\begin{lemma}[特異三次式を含む普遍的な持ち上げ]
\label{lem:universal-tangent-lift}
$S$ 上で相対圏
\[
 \mathcal D=\Perf(\mathcal X/S)
       =\langle\mathcal K,Q,\OO_{\mathcal X}\rangle,
 \qquad\mathcal B=\langle\mathcal K,Q\rangle
\]
と $E=i^!Q$、$G=j^!\OO_{\mathcal X}$ を構成できる。
$L=R^3p_*\OO_{\mathcal X}(-2)$ とおくと、相対自己Ext層は
\[
 \mathcal E xt^{0,1}(E,E)=\OO_S,\ \Lambda^2V\otimes L,
 \qquad
 \mathcal E xt^{0,2}(G,G)=\OO_S,\ \Sym^2V\otimes L
\]
であり、記載した以外の次数では零である。
gluingの連結射
\[
 \delta:\Lambda^2V\otimes L
      \longrightarrow\mathcal H^2CC^\bullet(\mathcal B/S)
\]
は、一意に
\[
 \widetilde\delta:\Lambda^2V\otimes L
      \longrightarrow\mathcal H^2CC^\bullet(\mathcal D/S)
\]
へ持ち上がる。
\end{lemma}

\begin{proof}
$\OO,\OO(1),Q$ の例外性は、射影空間のEuler列と
三次超曲面の制限列から従い、超曲面の滑らかさを必要としない。
また $\mathcal X/S$ はGorensteinなので、perfect対象に対する相対Serre双対性がある。
従って第\ref{sec:gluing-objects}節と同じmutationが定義され、
trace線を保持すると
\[
 \OO(-1)[1]\otimes L^*\longrightarrow E\longrightarrow Q,
 \qquad
 \OO(-2)[2]\otimes L^*\longrightarrow G\longrightarrow\OO
\]
という三角を得る。Euler列と低次数の制限列の計算は相対基底上でも成立し、
これらの三角から主張した自己Ext層を得る。
それらは局所自由であり、任意の基底変換と可換である。

Hochschild gluing triangleは、上三角dg圏のbar複体において
混合cochainを分けることで得られる。
この構成は任意の可換基底環上で成立し、smoothnessを仮定しない
\cite[Section 4.5, (4.5.1)--(4.5.2)]{KellerHoch}。
$\mathcal D=\langle\mathcal B,\OO\rangle$ に適用し、
$\mathcal E xt^1(G,G)=0$ を用いると、完全列
\[
 0\longrightarrow\mathcal H^2CC^\bullet(\mathcal D/S)
 \longrightarrow\mathcal H^2CC^\bullet(\mathcal B/S)
 \longrightarrow\Sym^2V\otimes L
\]
を得る。ここで $\Perf(S)$ の相対Hochschild複体は $\OO_S$ のみである。

$\delta$ と最後の射との合成は、局所自由層の間の射
\[
 \chi:\Lambda^2V\otimes L\longrightarrow\Sym^2V\otimes L
\]
である。滑らかな三次式の開集合上では、
第\ref{sec:gluing-objects}節の二段階のgluing計算により
$\HH^2(X)\to\HH^2(\mathcal B)$ は同型であるから、$\chi$ は零である。
$S$ は整かつ被約であり、この開集合は稠密なので、$\chi=0$ は $S$ 全体で成立する。
完全列から $\widetilde\delta$ の存在と一意性を得る。

相対複体の扱いを明記しておく。
$CC^\bullet(\mathcal D/S)$ は対角の相対自己Extと比較できる
\cite[Section 8.1, Theorem 8.9]{ToenMorita}。
普遍族はflatかつseparatedであり、対角の構造層はpseudo-coherentである。
従って各次数の相対Hochschild cohomologyはcoherentであり、flatな局所化と可換である。
$\mathcal B,\mathcal K$ についてもgluing triangleを使って局所化できる。
上の完全列は各affine基底上で構成して層化したものである。
特異ファイバーでHochschild cohomologyの任意の基底変換が同型になるとは仮定しない。
必要なのは自然な基底変換射と、上で計算した二つの自己Extの基底変換だけである。
\end{proof}

\subsection{反対称化とambient双ベクトル}

特異な可換スキームに対しても、Hochschildの二次類には
一次交換子を取る自然な射
\[
 \mathcal H^2CC^\bullet(\mathcal D/S)
 \longrightarrow
 p_*\mathcal H om_{\OO_{\mathcal X}}
   (\Lambda^2\Omega^1_{\mathcal X/S},\OO_{\mathcal X})
\]
がある。局所Hochschild 2-cocycle $\mu$ に対して
$\mu(f,g)-\mu(g,f)$ を取る射である。
正確には、対角のExtに対する局所大域スペクトル系列のedge射で
局所の次数2のExt層の大域切断へ送り、その上で反対称化する。
結合性の一次式からこれはbiderivationになり、
1-cochainのcoboundaryは可換代数上で対称なので、この射はcohomology上で定義される。
重なりでのcocycleの差もcoboundaryなので、局所のbiderivationは貼り合わさる。

この相対記述は通常のbar模型で確認できる。
$S$ を $F(v)\ne0$ の開集合で覆い、固定した線形座標変換で $v=e_0$ とする。
この開集合上の $\mathcal X$ は $x_j\ne0$、$1\le j\le4$ のaffine開集合で覆われる。
各座標環とその交叉は、$x_0$ についてmonicな三次式の商なので基底環上freeである。
従ってrelative bar/\v Cech模型で上の反対称化を行える。
ここでは特異スキームについてのHKR同型を用いていない。

$\Omega^1_{\PP^4/S}|_{\mathcal X}\twoheadrightarrow\Omega^1_{\mathcal X/S}$ から、
biderivationをambient双ベクトルへ単射で送れる。
さらに
\[
 \Lambda^2T_{\PP^4}\simeq\Omega^2_{\PP^4}(5)\otimes(\det V)^*,
 \qquad H^*(\PP^4,\Omega^2_{\PP^4}(2))=0
\]
なので、制限写像は
\[
 W\otimes\OO_S\xrightarrow{\sim}
 p_*(\Lambda^2T_{\PP^4}|_{\mathcal X}),
 \quad W=H^0(\PP^4,\Omega^2_{\PP^4}(5))\otimes(\det V)^*
\]
という同型である。消滅はEuler列の外積の大域切断が
$0\to\Lambda^2V\to V\otimes V\to\Sym^2V\to0$ となることから直接確認できる。

補題\ref{lem:universal-tangent-lift}の持ち上げと合成して、
全ての非零三次式上で正則な、$GL(V)$-同変射
\[
 \Psi:\Lambda^2V\otimes L\longrightarrow W\otimes\OO_S
\]
を得る。滑らかなファイバーで、ambientへの単射を $j_F$ と書けば、
符号とHKRの一定の正規化因子を除いて
\[
                         \Psi_F=j_F\circ\kappa_F^{-1}\circ e_F^{-1}
\]
である。すなわち、これはgluingの逆写像のHKR双ベクトルをambient空間へ送る写像である。

\subsection{三次式への一次性と表現の重複度1}

\begin{proof}[命題\ref{thm:general-tangent-comparison}の証明]
普遍超曲面の完全列
\[
 0\longrightarrow\OO_{\PP^4}(-5)
 \xrightarrow{\ F\ }\OO_{\PP^4}(-2)
 \longrightarrow\OO_{X_F}(-2)\longrightarrow0
\]
の境界により、
\[
 \partial_F:L_F\xrightarrow{\sim}
 H^4(\PP^4,\OO(-5))\simeq(\det V)^*
\]
と自明化する。この境界は
\[
                         \partial_{\lambda F}=\lambda^{-1}\partial_F
\]
を満たす。実際、上の完全列で中央と右端を恒等に保つ比較射は、
左端に $\lambda^{-1}$ を掛ける射である。

$F$ と $\lambda F$ は同じ超曲面と同じ内在的gluingを与えるので、
この自明化で $\Psi_F$ を
$\Lambda^2V\otimes(\det V)^*\to W$ と書けば
\[
                            \Psi_{\lambda F}=\lambda\Psi_F
\]
となる。
$\Psi$ は $\Sym^3V\setminus\{0\}$ 全体で正則であり、
この35次元affine空間から除いた点の余次元は2以上である。
従って各行列成分は原点へ正則に延長する。
上のhomogeneityから各成分は $F$ の係数について一次式である。
よって $\Psi$ は $GL(V)$-同変な線形射
\[
 \Sym^3V\otimes\Lambda^2V\otimes(\det V)^*
                   \longrightarrow W
\]
に対応する。

Euler--Koszul列とPieri則によって
\[
 H^0(\PP^4,\Omega^2(5))\simeq\mathbb S_{(3,1,1)}V,
 \qquad
 \Sym^3V\otimes\Lambda^2V
  \simeq\mathbb S_{(4,1)}V\oplus\mathbb S_{(3,1,1)}V.
\]
特に目標の既約表現の重複度は1である。
ここで最初の表現は、完全列
\[
 0\longrightarrow H^0(\Omega^2(5))
 \longrightarrow\Lambda^2V\otimes\Sym^3V
 \longrightarrow V\otimes\Sym^4V
 \longrightarrow\Sym^5V\longrightarrow0
\]
の核としても得られる。
従って、この形の非零の同変射はscalar倍を除いて一意である。

標準的な同変射を具体的に書く。
$a=v\wedge w$ に対して $\eta_a=v\,dw-w\,dv$ とし、線形に延長する。
$\eta_a$ は $\Omega^1_{\PP^4}(2)$ の切断である。
affine錐のEulerベクトル場を $\mathsf E$ とすると
\[
 \iota_{\mathsf E}\eta_a=0,\quad
 \iota_{\mathsf E}d\eta_a=2\eta_a,\quad
 \iota_{\mathsf E}dF=3F.
\]
従って
\[
                 \Omega(F,a)=3F\,d\eta_a-2\,dF\wedge\eta_a
\]
は水平でweight 5の二形式であり、$H^0(\Omega^2_{\PP^4}(5))$ に属する。
これは $F,a$ に双線形かつ $GL(V)$-同変である。
$X_F$ 上では
\[
                         \Omega(F,a)|_{X_F}=-2\,dF\wedge\eta_a
\]
となる。conormal列の外積とadjunctionにより、これは
$\lambda_F(a)$ のambient双ベクトルへの像の一定非零scalar倍である。
滑らかな $F$ では $\lambda_F$ は同型なので、この同変射は非零である。

重複度1により $\Psi$ はこの射の一定scalar倍となる。
また滑らかなファイバーで $\kappa_F^{-1}$ は同型であり、
HKR反対称化とambientへの単射も双ベクトル上で単射なので、$\Psi$ は非零である。
従って二つの写像を結ぶscalarは非零であり、
$\lambda_F e_F\kappa_F=c\,\id$ を得る。

命題\ref{thm:object-cone}の形式座標の一次部分は $\kappa_F$ である。
従ってedge座標への変更もPl\"ucker錐をscalar倍で保ち、
接錐はスキームとして同じPl\"ucker錐になる。
その台のBrill--Noetherによる記述は補題\ref{lem:bn-rank}から従う。
\end{proof}

この証明は、一般の $X$ について二つの接座標の比較を与える。
一方、適合部分空間 $C$ の高次の方程式や大域的な閉包を、
Brill--Noether条件だけで特徴付けたわけではない。
また、零点集合の次元跳躍から別途作ったスキーム構造とPl\"uckerイデアルの
一致までは主張していない。
第\ref{sec:picard-compatibility}節では別の全次数の持ち上げを用い、
適合部分空間そのものを相対Picardの形式像として記述する。

\subsection{任意の適合形式弧への特殊化の帰結}

\begin{theorem}[一般の $X$ と任意の非定数な適合形式弧]
\label{thm:general-compatible-arc-bn}
$R=\CC[[h]]$ とし、$C$ を通る非定数な形式弧を取る。
Pl\"ucker形式座標でその最初の非零項を
\[
 a(h)=h^\nu a_\nu+O(h^{\nu+1}),\qquad a_\nu\ne0
\]
と書く。すると $a_\nu$ は分解可能である。
対応する実際の $R$-perfect対象 $\mathcal E_R$ の幾何学的一般ファイバーは、
class $3[I_\ell]$ のGieseker-stableなtorsion-free層である。
その非局所自由な部分をFittingイデアルで定め、
そのスキーム論的水平閉包を $S_R\subset X\times\operatorname{Spec}R$ とすると、
\[
 (S_R)_0\ \subset\
 C_0:=X\cap\PP\bigl(\ker(a_\nu:V^*\to V)\bigr)
\]
は閉部分スキームの包含である。$C_0$ は平面三次曲線である。

さらに $C_0$ が滑らかで、一般ファイバーの非局所自由な部分が1次元なら、
$(S_R)_0=C_0$ であり、$S_R/R$ は平面三次曲線の滑らかな族になる。
\end{theorem}

\begin{proof}
$q(a(h))=0$ の最初の非零となり得る係数から $q(a_\nu)=0$ を得る。
従って $a_\nu$ は分解可能である。
命題\ref{thm:algebraic-compatibility}のuniversal対象を弧へ引き戻せば、
実際の $R$-perfect対象が得られる。
非定数性と\cite[Proposition 5.5]{FLZ}から、
命題\ref{thm:gluing-sheaf-arc}と同じ議論で一般ファイバーは主張した層になる。

局所二項表示 $[\OO\xrightarrow d\OO^4]$ において、
命題\ref{thm:general-tangent-comparison}から、
非零scalarと局所自明化を除けば
\[
                         d=h^\nu s_{a_\nu}+O(h^{\nu+1})
\]
である。一般ファイバーの非局所自由部分のFittingイデアルは
$d$ の四つの成分で生成される。
水平閉包はこのイデアルの $h$-saturationで定まるので、
$d$ の各成分を $h^\nu$ で割ったものを含む。
特殊化により $S_0\subset Z(s_{a_\nu})=C_0$ をスキームとして得る。

$h$-saturationにより $\OO_{S_R}$ は $R$-torsion-free、従って $R$-flatである。
$S_R$ はproperで、1次元の一般ファイバーを持つなら特殊ファイバーも1次元である。
滑らかな平面三次曲線 $C_0$ は整であり、その1次元の閉部分スキームは
$C_0$ 全体に等しい。従って $S_0=C_0$。
flatnessと特殊ファイバーの滑らかさから $S_R\to\operatorname{Spec}R$ は
$S_0$ の近傍で滑らかであり、propernessから全体で滑らかになる。

最後に $H^1(C_0,\OO_{C_0}(1))=0$ と $h^0(\OO_{C_0}(1))=3$ より、
$H^0(S_R,\OO_{S_R}(1))$ は階数3の自由 $R$-加群であり、
$V\otimes R$ からの制限は全射である。
その核は階数2の直和因子で、$S_R$ を含む相対平面を定める。
$X$ とこの平面との交叉は $R$-flatな三次曲線族で、特殊ファイバーは $C_0$。
その構造層から $\OO_{S_R}$ への全射の核は、flatnessとNakayamaにより
特殊ファイバーの近傍で零であり、propernessにより全体で零である。
よって $S_R$ はこの平面切断の族に一致する。
\end{proof}

最後の結論では、一般ファイバーの非局所自由な部分が1次元であることを仮定した。
その部分が空または0次元である可能性を、接錐の計算だけで除いてはいない。
命題\ref{thm:all-compatible-arcs}では相対Picardによる記述を使い、
全ての非定数な適合形式弧で、その水平閉包が平面三次曲線の族に等しいことを示す。

\subsection{Brill--Noether境界の六次接錐}\label{sec:sextic-boundary}

本節は命題\ref{thm:bn-incidence}のincidenceから得られる、
局所自由でない対象の境界についての計算である。
適合部分空間 $C$ との包含や同一視は仮定しない。
$\mathcal I\subset X\times U$ を同命題のincidenceとし、
$D\subset U$ を射 $\mathcal I\to U$ のscheme-theoretic imageとする。

$W=\Lambda^2V$ とし、$a\in W$ の5個の符号付きPfaffianを
$q(a)=(q_0(a),\ldots,q_4(a))$ と書く。
符号は、階数4の場合に $[q(a)]\in\PP(V^*)$ が $a$ のkernel点になるように取る。
$X=V(F)$ の方程式を用いて
\[
 H_X=V\bigl(F(q_0(a),\ldots,q_4(a))\bigr)\subset\PP(W)
\]
と定める。

\begin{theorem}[境界の重複度と接錐]\label{thm:sextic-boundary}
$U$ を $e=[E]$ の周りで縮小すると、$D$ は既約かつ被約な因子である。
cohomologyの接座標 $T_eU=\Lambda^2V$ に関して
\[
 \operatorname{TC}_eD
 =\operatorname{Spec}
 \frac{\CC[a_{ij}]}{\bigl(F(q_0(a),\ldots,q_4(a))\bigr)}
\]
がスキームとして成立する。特に $\operatorname{mult}_eD=6$ である。
$H_X$ は正規な既約六次超曲面で、特異部分の台は $\Gr(2,V)$ である。
また
\[
 \PP_X(\Lambda^2Q)\longrightarrow H_X
\]
は解消であり、階数2の $[a]$ 上のファイバーは
$X\cap\PP(\ker a)$ という平面三次曲線である。
\end{theorem}

\begin{proof}
\textbf{像とblowup。}
$X$ は射影的なので $\mathcal I\to U$ はproperである。
$\mathcal I$ はcentral fiber $X$ の近傍で滑らかである。
非滑らかな点の閉集合のproperな像をまず除き、
次に $X$ に会わない成分のproperな像を除けば、滑らかな既約9次元スキームとしてよい。
そのscheme-theoretic imageはintegralである。
実際、affine開集合 $\operatorname{Spec}A\subset U$ 上で、そのイデアルは
\[
 A\longrightarrow\Gamma(\mathcal I\times_U\operatorname{Spec}A,\OO)
\]
の核であり、右辺は空でなければ整域である。
proper性によりこの核はcoherentなイデアルを与える。

$\mathcal I_e=X$ がスキームとして成立するので
$\mathfrak m_e\OO_{\mathcal I}=\mathcal I_{X/\mathcal I}$ である。
従ってblowupの普遍性からproperな射
\[
 f:B=\operatorname{Bl}_X\mathcal I\longrightarrow\operatorname{Bl}_eU
\]
を得る。sourceの例外因子を $J=\PP_X(\Lambda^2Q)$、targetの例外因子を
$E_U=\PP(T_eU)$ とすると、Cartier因子として
\begin{equation}\label{eq:sextic-exceptional-order}
                         f^*E_U=J
\end{equation}
であり、係数は1である。
$B\setminus J$ は $\mathcal I\setminus X$ と同型なので、
$f(B)$ は $D$ のstrict transform $\widetilde D$ である。

\textbf{例外因子上の射。}
normal bundleの包含から、$f|_J$ は
\[
 J=\{(p,[a])\in X\times\PP(W):\iota_p(a)=0\}
   \longrightarrow\PP(W)
\]
という射である。
階数4の部分では $[a]$ のkernelは一点であり、その点は $[q(a)]$ である。
kernel点への写像
\[
 \PP(W)_{\rk4}\longrightarrow\PP(V^*)
\]
は、ファイバーが $\PP(\Lambda^2\ker p)$ 内の開集合となる滑らかな射である。
従って $X$ の逆像は滑らかである。
階数2の場合のファイバーは平面三次曲線であり、必ず非空である。
ここで滑らかな三次三次元多様体は平面を含まない。
実際、平面のイデアルを $(u,v)$ として $F=uQ_1+vQ_2$ と書ければ、
平面上の二つの二次式 $Q_1,Q_2$ の共通零点で $X$ が特異になってしまう。
よって $J$ の像は $H_X$ であり、階数4の部分では $J\to H_X$ は同型である。

$J$ は $X$ 上の $\PP^5$-束なので既約である。
Pfaffianの共通零点 $\Gr(2,V)$ は余次元3であり、そこに別の超曲面成分は存在し得ない。
階数4の部分での滑らかさと合わせて、$F(q)$ は既約かつ被約な六次式である。
また $\widetilde D\cap E_U$ の台は8次元の $H_X$ なので、$D$ は9次元の因子となる。

\textbf{接錐の重複度。}
$H_X$ の階数4の部分の閉点 $z$ を取る。
$f^{-1}(z)$ は $J$ 上の一点 $y$ だけであり、
$J$ の接方向の写像は8次元の同型である。
\eqref{eq:sextic-exceptional-order}により、$J$ の法方向も
$E_U$ の法方向へ非零に写る。従って $df_y$ の階数は9であり、
$f$ は $y$ でunramifiedである。
proper性とファイバーの一意性から、$z$ の近傍で $f$ をfiniteにできる。
そのファイバーは被約な一点 $\operatorname{Spec}\CC$ なので、
finite加群に対するNakayamaにより局所的な構造環の射は全射である。
従って $f$ はそこで閉埋め込みとなり、$\widetilde D$ は滑らかで、
$E_U$ と横断的に交わる。

$U$ は正則なので $D$ の局所方程式を $d$ と書ける。
その最低次斉次項の台は既約な $H_X$ だから
$\operatorname{in}(d)=c\,F(q)^m$、$m\ge1$ である。
上の横断性は例外因子上での重複度が1であることを意味し、$m=1$ を得る。
これがスキームとしての接錐と重複度6を証明する。

\textbf{特異部分と解消。}
階数4の部分は既に滑らかである。
階数2では $q=0$ であり、$F$ は三次式なので $F(q)$ の全ての一次偏微分が零になる。
従って特異部分の台は $\Gr(2,V)$ である。
これは $H_X$ 内で余次元2なので、超曲面の $S_2$ 性と合わせて正規性を得る。
$J$ は滑らかであり、$J\to H_X$ はproperかつ双有理なので解消である。
\end{proof}

\begin{remark}[横断的な三次錐]
階数2の点の近傍で $a_{01}\ne0$ とし、交代行列をblock表示する。
可逆な $2\times2$ blockに対するSchur complementを取ると、
残る $3\times3$ 交代行列の三成分がGrassmannianに横断する座標になる。
この座標でPfaffianのkernel vectorは三成分に関して正確に線形である。
従って $H_X$ の横断sliceは、対応する平面切断
$X\cap\PP(\ker a)$ を底とする三次錐になる。
平面切断が滑らかな場合には、横断sliceの解消の例外曲線はその楕円曲線である。
本節は $D$ の接錐を決定しており、$D$ 自身をその錐へ形式同型にする主張ではない。
\end{remark}

命題\ref{thm:general-tangent-comparison}と合わせると、同じ接空間 $T_eU$ の中で
\[
 \bigl(\operatorname{Sing}(\operatorname{TC}_eD)\bigr)_{\mathrm{red}}
                  =\operatorname{TC}_eC
\]
が成立する。すなわち、7次元の適合条件の接錐は、
非局所自由な境界の六次接錐の特異部分を被約化したものに一致する。
この接錐の等式だけでは、$C\subset D$ や
$C=\operatorname{Sing}(D)_{\mathrm{red}}$ は導かれない。
命題\ref{thm:compatibility-in-bn-boundary}では相対Picardの全次数の記述から
形式包含 $\widehat C_c\subset\widehat D_e$ を別に証明する。
後者の特異部分全体との等式は、第7版の命題\ref{thm:global-bn-singular}で証明する。


\section{固定残余圏族の貼り合わせ対象：楕円曲線上の評価核}\label{sec:evaluation-gluing}

第\ref{sec:bn-algebraic}節では、明示的な固定残余圏族のgluing objectが、
generic fiberで階数3の安定torsion-free層になることを示した。
本節ではFermatの対称性と対象の変形を用い、この層を具体的に記述する。

\subsection{平面三次曲線上の次数3の線形束}

$X\subset\PP^4$ は滑らかなcubic threefold、
$C=X\cap\PP^2$ は滑らかな平面三次曲線とし、$j:C\hookrightarrow X$ を包含とする。
$L\in\Pic^3(C)$ に対して
\begin{equation}\label{eq:evaluation-kernel}
 0\longrightarrow F_L\longrightarrow H^0(C,L)\otimes\OO_X
 \xrightarrow{\mathrm{ev}}j_*L\longrightarrow0
\end{equation}
で $F_L$ を定める。次数3なので $h^0(C,L)=3$、$h^1(C,L)=0$ であり、
$L$ は大域切断で生成される。

\begin{theorem}[評価核の安定性と特異部分]\label{thm:evaluation-sheaf}
$F_L$ はGieseker-stableな階数3のtorsion-free層であり、
\[
 \operatorname{ch}(F_L)=(3,0,-H^2,0),\qquad
 F_L^{**}=H^0(C,L)\otimes\OO_X,\qquad F_L^{**}/F_L=j_*L.
\]
非局所自由な部分は、Fittingイデアルによるスキーム構造を含めて $C$ である。
さらに $L\ne\OO_C(1)$ なら $F_L\in\Ku(X)$ であり、
\[
 h^2(X,F_L(-2))=3.
\]
従ってこの層はUlrichではない。
\end{theorem}

\begin{proof}
$C$ は余次元2の局所完全交叉である。
$C$ 上で $L$ を局所自明化し、評価射の一成分を単元に直すと、
\[
 F_L\simeq\mathcal I_C\oplus\OO_X^{\oplus2}
\]
という局所表示を得る。
従ってtorsion-free性、二重双対とその商、
$\operatorname{Fitt}_3(F_L)=\mathcal I_C$ が従う。
$\chi(C,L)=3$ と $H\cdot C=3$ から
$\operatorname{ch}(j_*L)=(0,0,H^2,0)$ であり、Chern characterも得られる。

安定性を示す。$A\subset F_L$ を階数 $r=1,2$ の非零部分層とする。
$F_L\subset\OO_X^3$ なので $\mu_H(A)\le0$ である。
$\mu_H(A)<0$ なら既にreduced Hilbert polynomialは $p_A<p_{F_L}$ となる。
$\mu_H(A)=0$ とする。$\OO_X^3$ の中での飽和を $\overline A$ とすると、
飽和で生じるeffectiveな因子の次数は $\OO_X^3$ の半安定性により零である。
従って $c_1(\overline A)=0$。
$\det(\overline A)^{**}=\OO_X$ のPl\"ucker座標は定数なので、
そのgenericな $r$ 次元部分空間は固定される。
飽和性から $\overline A=W\otimes\OO_X$、$\dim W=r$ と書ける。

$K_W=F_L\cap(W\otimes\OO_X)$ とおく。
$0\ne W\subset H^0(C,L)$ の評価像は、整な曲線 $C$ 上でgeneric rank 1を持つので、
$\operatorname{ch}_2(K_W)=-[C]$ である。
$A\subset K_W$ の商は余次元2以上のtorsion層だから
\[
 \frac{H\cdot\operatorname{ch}_2(A)}r
 \le-\frac3r<-1
 =\frac{H\cdot\operatorname{ch}_2(F_L)}3.
\]
$c_1=0$ のもとでHilbert多項式の一次係数を比較すれば、
再び $p_A<p_{F_L}$ を得る。

\eqref{eq:evaluation-kernel}の大域切断の射は同型なので $H^*(X,F_L)=0$。
$L\ne\OO_C(1)$ なら $L(-1)$ は非自明な次数0の線形束であり、
$H^*(C,L(-1))=0$ である。
$\OO_X(-1)$ もacyclicなので $H^*(X,F_L(-1))=0$、従って $F_L\in\Ku(X)$。
最後に同じ完全列を $\OO_X(-2)$ でtwistすると
\[
 H^2(X,F_L(-2))\simeq H^1(C,L(-2))\simeq\CC^3
\]
となる。これはnormalized UlrichのACM条件\cite[Section 6]{FP}を満たさない。
\end{proof}

\subsection{Picardの原点をgluing objectへ送る射影}

$V=H^0(X,\OO_X(1))$、$U\subset V$ を $C$ の平面を切り出す2次元部分空間とし、
$\pi=i^*=L_{\OO_X}L_{\OO_X(1)}$ と書く。
$L$ の族に対して
\[
 \mathcal P_L=\pi(j_*L)[-1]
\]
を考える。$L\ne\OO_C(1)$ では命題\ref{thm:evaluation-sheaf}と
\eqref{eq:evaluation-kernel}から $\mathcal P_L\simeq F_L$ である。

trace線を
\[
 T=\Ext^2(Q,\OO_X(-1))\simeq H^3(\OO_X(-2))
\]
とする。$C$ のKoszul列をtwistした複体
\[
 [\,\OO_X(-1)\otimes\det U\longrightarrow
       \OO_X\otimes U\longrightarrow\OO_X(1)\,]
 \simeq j_*\OO_C(1)
\]
は次数 $-2,-1,0$ に置かれる。
$\pi$ は右の二項を殺し、\eqref{eq:E-triangle}から
$\pi(\OO_X(-1))[1]\simeq T\otimes E$ なので
\begin{equation}\label{eq:picard-special-object}
 \mathcal P_{\OO_C(1)}\simeq\det U\otimes T\otimes E.
\end{equation}
従って全ての平面についてPicardのこの点は、粗モジュライの同じ点 $[E]$ へ写る。
固定した $C$ では1次元因子を自明化して $\mathcal P_{\OO_C(1)}\simeq E$ とする。

\subsection{二対象の最小模型とPicard方向の接写像}

\begin{lemma}[射影の局所的な積の公式]\label{lem:mutation-product}
$A=\OO_X(-1)$、$B=Q$ とし、trace線を自明化する。
$0\ne a\in W=\Hom(A,B)=\Lambda^2V$ として
$F_a=\operatorname{Cone}(a:A\to B)$ とおく。
$F_a$ の形式変形は、$a$ の一座標を1に規格化した $f\equiv a$ と、
一つのパラメータ $t$ を用いて記述できる。
$E$ の形式変形の座標を $w\in W$ とすれば、
射影 $\pi$ が誘導する写像は適切なこの座標で
\[
                       (f,t)\longmapsto w=tf
\]
である。特に原点での接写像は
$(\dot f\bmod\CC a,\dot t)\mapsto\dot t\,a$ である。
\end{lemma}

\begin{proof}
Euler列とSerre双対性から、二対象 $A,B$ のfull dg部分圏のcohomologyは
\[
 \begin{aligned}
 \Hom^*(A,A)=\CC\id_A,\quad \Hom^*(B,B)=\CC\id_B,
 &\qquad\text{in degree }0,\\
 \Hom^*(A,B)=W\text{ in degree }0,
 &\qquad\Hom^*(B,A)=\CC g\text{ in degree }2.
 \end{aligned}
\]
である。strictly unitalな最小 $A_\infty$ 模型を取る。
非単位の合成可能な入力は $W$ と $g$ が交互に並ぶものに限られる。
$m_n$ の次数 $2-n$ から、偶数個の入力の出力は対角Homの次数2、
奇数個の入力の出力は $A\to B$ の次数1または $B\to A$ の次数3となる。
いずれも零なので、radical同士の積と全ての高次演算は零である。

$F_a$ は $A[1]\oplus B$ に微分 $a$ を入れたもので、
その自己Extは次数0で $\CC$、次数1で $(W/\CC a)\oplus\CC g$、他で零である。
$E$ の次数1の $W$ 成分は、三角の射 $\iota,\beta$ による
$\alpha\mapsto\beta[1]\alpha\iota$ そのものなので、
第\ref{sec:bn-algebraic}節のcohomologyの接座標と一致する。
従って最小模型によるArtin変形\cite{ELOI,ELOII}を
\[
 F_{f,t}:\ \delta=f+t g
\]
と書ける。一方 $E$ は微分 $g$ の対象であり、その変形を
$E_w:\delta=g+w$ と書ける。
これらの記述は可換Artin基底およびその形式逆極限上で用いる。
非triangularな行列を、任意の通常のパラメータで自動的にperfectな対象と
みなすことはしていない。

$A,B\in\OO_X^\perp$ なので、ここで $\pi$ は $L_B$ と一致する。
$\RHom(B,F_{f,t})$ は次数 $0,1$ の二項複体 $[R\xrightarrow{t}R]$ である。
left mutationの評価射のconeを消去すると $E_{tf}$ が残る。
より具体的には、$A[1]\oplus B$ 上の
\[
 \operatorname{diag}(1,t):F_{f,t}\longrightarrow E_{tf}
\]
は閉射であり、そのconeから $A$ 上の恒等写像を消去すると
$\operatorname{Cone}(t:B\to B)$ が残る。
$\pi$ はこのconeを殺し、$E_{tf}$ は $B$ と $\OO_X$ の右直交にある。
従って $\pi(F_{f,t})\simeq E_{tf}$ となる。
\end{proof}

\begin{lemma}[Picard方向はPl\"ucker直線を与える]\label{lem:picard-gluing-ks}
$\Pic^3(C)$ の $\OO_C(1)$ における形式近傍から $\Def_E$ への写像
$L\mapsto\mathcal P_L$ の接写像は非零である。
$\Ext^1(E,E)=\Lambda^2V\otimes T$ の同一視の下で、その像は
$\det U\otimes T$ というPl\"ucker直線である。
\end{lemma}

\begin{proof}
Euler列とKoszul列を比較すると
\[
 0\longrightarrow\OO_X(-1)\otimes\det U
   \xrightarrow{a}Q\longrightarrow F_{\OO_C(1)}\longrightarrow0
\]
を得る。ここで $a$ は $\Lambda^2U$ の非零元に対応する。
従って補題\ref{lem:mutation-product}を適用できる。

評価核の完全列から
\[
 \RHom(\OO_X(1),F_L)\simeq R\Gamma(C,L(-1))[-1].
\]
$L=\OO_C(1)$ では右辺は次数1と2に1次元ずつを持つ。
$0\ne\xi\in H^1(C,\OO_C)$ による線形束の一次変形は、その二つの間の微分を
\[
 H^0(C,\OO_C)\xrightarrow{\,\cup\xi\,}H^1(C,\OO_C)
\]
という非零写像に変える。
補題の模型では $f$ 方向はこの微分を変えず、$g$ 方向が非零に変える。
後者は $\delta=g$ の対象 $E$ が $\Ku(X)$ に属することからも確認できる。
従ってPicard方向の $\dot t$ は零でなく、射影後の接方向は $\dot t\,a$ である。
traceを戻すと、$H^1(\OO_C)\simeq\det U\otimes T$ から
$\Lambda^2V\otimes T$ へのtautologicalな直線の包含となる。
\end{proof}

\subsection{Fermatの固定部分空間と明示的族の同定}

\begin{theorem}[明示的族のgeneric gluing sheaf]\label{thm:explicit-gluing-sheaf}
第9節の固定残余圏族 $A'_h$ を $R=\CC[[h]]$ 上で考え、
\[
 C_0=X_0\cap\{u=v=0\}=V(x^3+y^3+z^3)\subset\PP^2
\]
とする。$L_0=\OO_{C_0}(1)$ を特殊値とする次数3の線形束の族 $L_R$ が存在し、
$\Pic^3(C_0)$ へのその形式弧の一次項は非零である。
適切な基底上の1次元因子の自明化の下で
\[
 \mathcal E_R\simeq i^*(j_*L_R)[-1].
\]
特に $L_\eta\ne\OO_{C_0}(1)$ であり、generic fiberは
\begin{equation}\label{eq:explicit-gluing-evaluation}
 0\longrightarrow\mathcal E_\eta\longrightarrow
 H^0(C_0,L_\eta)\otimes\OO_{X_0}
 \longrightarrow j_*L_\eta\longrightarrow0
\end{equation}
で表される。従って
\[
 \mathcal E_\eta^{**}\simeq\OO_{X_0}^{\oplus3},\qquad
 \mathcal E_\eta^{**}/\mathcal E_\eta\simeq j_*L_\eta,
 \qquad\operatorname{Fitt}_3(\mathcal E_\eta)=\mathcal I_{C_0}.
\]
この層はGieseker-stableであり、局所自由でなく、Ulrichでもない。
\end{theorem}

\begin{proof}
$\Gamma=\mu_3(u)\times\mu_3(v)$ とする。
$\Gamma$ は $C_0$ を点ごとに固定するので、$\Pic^3(C_0)$ と
$L\mapsto i^*(j_*L)[-1]$ の粗モジュライの像を固定する。
\eqref{eq:picard-special-object}の1次元因子のcharacterは
\[
 \det\langle u,v\rangle\otimes T:
 (\zeta_u\zeta_v)(\zeta_u\zeta_v)^{-1}=1
\]
であり、特殊値も $\Gamma$-同変に $E$ と同一視できる。

命題\ref{thm:fermat-bn-cone}のtraceを保持した表現を用いると
\[
 \Ext^1(E,E)=\Lambda^2V\otimes(\det V)^*,\qquad
 \dim\Ext^1(E,E)^\Gamma=1.
\]
固定部分は $u\wedge v$ の直線である。
有限群の作用は、characteristic zeroでReynolds作用素を使って形式的に線形化できる。
従って滑らかな $\Def_E$ の固定部分 $(\Def_E)^\Gamma$ は滑らかな形式曲線である。
補題\ref{lem:picard-gluing-ks}と形式逆関数定理から
\[
 \widehat{\Pic^3(C_0)}_{\,\OO_{C_0}(1)}
       \xrightarrow{\sim}(\Def_E)^\Gamma.
\]

$A'_h$ の関係式は $\Gamma$-作用を保ち、例外対象の持ち上げとmutationも自然である。
残余圏を固定した $\mathcal K_R$ と同一視した後、輸送された $\Gamma$ の作用と
一定の $\Gamma$ の作用との差は、特殊ファイバーで恒等となる自己同値である。
第7節の形式比較と $\HH^1(\mathcal K)=0$ により、その同型類は自明である。
従って $\mathcal E_R$ の標識付き変形の同値類は $\Gamma$ で固定される。
ここでは古典的な同値類上の固定性だけを使い、coherentな群作用の厳密化は仮定しない。

ゆえにこの弧は、上のPicard形式曲線の一意な弧に対応する。
既存の非零の一次gluing類と補題\ref{lem:picard-gluing-ks}から、
Picardへの弧も一次項が非零である。
Picard曲線のaffine chartに完備局所環の射を制限し、Poincar\'e束を引き戻せば、
実際の $R$ 上の $L_R$ を得る。
simple対象の普遍\'etale chartで比較すれば、必要な基底上の線形束の因子を除いて
$\mathcal E_R\simeq i^*(j_*L_R)[-1]$ となる。この因子は局所環 $R$ 上では自由である。

非定数性により $L_\eta\ne\OO_{C_0}(1)$ なので、generic fiberでは射影が不要となる。
\eqref{eq:explicit-gluing-evaluation}と残りの結論は
命題\ref{thm:evaluation-sheaf}から従う。
\end{proof}

従って命題\ref{thm:gluing-sheaf-arc}の特殊化の包含は、
この明示的な固定残余圏族については、固定曲線 $C_0$ そのものへ強化される。
一方、線形束 $L_R$ のPicard座標を、Sklyaninの平行移動点 $p(h)$ と直接結ぶ
明示公式はここでは得ていない。
この段階の証明は一本のFermat族を扱うものである。
命題\ref{thm:all-compatible-arcs}では、任意の滑らかな $X$ の
全ての非定数な適合形式弧へ、この評価核表示を拡張する。

\section{五生成元代数の内在的復元と次数付き分類}\label{sec:five-generator}

本節では、第\ref{sec:reconstruction}節の三生成元Sklyanin代数の復元を、
中心変数を二つ加えた五生成元代数の分類へ拡張する。
まず $k$ は標数零の代数閉体とする。
\[
 S=S(\alpha,\beta,\gamma)
 =\frac{k\langle x,y,z\rangle}
 {(\alpha yz+\beta zy+\gamma x^2,\
   \alpha zx+\beta xz+\gamma y^2,\
   \alpha xy+\beta yx+\gamma z^2)}
\]
の点曲線 $E$ が滑らかな楕円曲線であると仮定する。
具体的には $\alpha\beta\gamma\ne0$ および
\[
 (\alpha^3+\beta^3+\gamma^3)^3\ne27\alpha^3\beta^3\gamma^3
\]
を満たす通常の楕円Sklyaninの場合を扱う。
$L=\OO_E(1)$、$\sigma=T_p$ とし、
標準的な点加群の商の核を生成する次数3の中心元を $g_S$ と書く。
\cite[Definition--Lemma 2.13, Lemma 2.19]{WWY}により
\[
 S/(g_S)\simeq B(E,L,\sigma)
\]
であり、$g_S$ は非零scalar倍の選択を除いて定まる。
以後
\[
 A(S,g_S)=S[u,v]/(g_S+u^3+v^3),\qquad \deg u=\deg v=1
\]
とおく。$u,v$ は中心変数である。

\subsection{次数1の中心とquadratic cover}

\begin{theorem}[中心平面からのSklyanin代数の復元]\label{thm:intrinsic-sklyanin}
$A=A(S,g_S)$ の次数1の中心は
\[
 U_A:=Z(A)\cap A_1=ku\oplus kv
\]
である。従って $A$ から内在的に
\[
 B_A=A/(U_A)=S/(g_S)
\]
を回復できる。さらに、次数2までの積から作るquadratic cover
\begin{equation}\label{eq:intrinsic-quadratic-cover}
 \mathcal Q(B_A):=
 \frac{T_k((B_A)_1)}
 {\left(\ker\bigl[(B_A)_1\otimes(B_A)_1\longrightarrow(B_A)_2\bigr]\right)}
 \simeq S
\end{equation}
により、元の三生成元代数そのものを回復できる。
この構成は次数付き同型に関して関手的である。
\end{theorem}

\begin{proof}
$W=S_1$ とし、三つの二次関係が張る空間を $R_S\subset W\otimes W$ とする。
各関係はそれぞれ $\gamma x^2,\gamma y^2,\gamma z^2$ を含み、
$\gamma\ne0$ である。従って
\[
 R_S\cap\Lambda^2W=0.
\]
もし $w\in W$ が中心的なら、全ての $w'\in W$ について
$w\otimes w'-w'\otimes w\in R_S\cap\Lambda^2W$ となる。
$\dim W=3$ なので $w=0$ である。よって $Z(S)\cap S_1=0$。

三次関係による商は次数2までを変えないので、
\[
 A_1=W\oplus(ku\oplus kv),\qquad
 A_2=S_2\oplus(W\otimes(ku\oplus kv))\oplus\Sym^2(ku\oplus kv)
\]
である。$w+\lambda u+\mu v$ が中心的であるための条件は、
$w$ が $S$ の次数1の中心に入ることに等しい。
従って $U_A=ku\oplus kv$ を得る。
これで $B_A$ の表示も従う。

$g_S$ は次数3なので、$S$ と $B_A$ の次数1、次数2およびその間の積は同じである。
$S$ は二次関係で定義されるため、\eqref{eq:intrinsic-quadratic-cover}を得る。
中心、商、積のkernelはいずれも内在的な構成なので、関手性も従う。
\end{proof}

特に、この復元は中心全体の生成元を決定する必要がなく、
$\sigma$ が有限位数か無限位数かにも依存しない。

\subsection{次数3の中心元の線と全ての同型の形}

\begin{lemma}[次数3の中心元の線の不変性]\label{lem:canonical-central-line}
上のクラスのSklyanin代数の次数付き同型
$\phi:S\xrightarrow{\sim}S'$ は
\[
 \phi(g_S)=\lambda g_{S'}\qquad(\lambda\in k^\times)
\]
を満たす。
\end{lemma}

\begin{proof}
$S$ の次数3の元で、全てのpoint module $N$ に対して
$N_0$ を零化するものの空間を $I_3(S)$ とする。
$S\twoheadrightarrow B(E,L,\sigma)$ の核は $(g_S)$ である
\cite[Lemma 2.19]{WWY}。
商の次数3の元は $L_3=L\otimes\sigma^*L\otimes(\sigma^2)^*L$ の切断であり、
各点のpoint moduleへの作用はその評価である。
非零切断はある点で非零なので
\[
 I_3(S)=k g_S.
\]
次数付き同型はpoint moduleとその次数零の部分を保存するから、
$I_3(S)$ を $I_3(S')$ へ写す。
従って次数3の中心元の線は保存される。
\end{proof}

\begin{theorem}[五生成元代数の次数付き同型判定]\label{thm:five-generator-isom}
上のクラスの $S,S'$ と次数3の中心元の選択について、
\begin{equation}\label{eq:five-generator-isom}
 A(S,g_S)\simeq_{\mathrm{gr}}A(S',g_{S'})
 \quad\Longleftrightarrow\quad
 S\simeq_{\mathrm{gr}}S'
\end{equation}
が成り立つ。
さらに全ての次数付き同型は、次数1の分解
\[
 S_1\oplus(ku\oplus kv),\qquad S'_1\oplus(ku'\oplus kv')
\]
についてblock diagonalである。
すなわち、同型の全体は
\[
 \phi:S\xrightarrow{\sim}S',\qquad
 \tau:ku\oplus kv\xrightarrow{\sim}ku'\oplus kv'
\]
であって、ある $\lambda\in k^\times$ に対し
\begin{equation}\label{eq:five-generator-iso-pairs}
 \phi(g_S)=\lambda g_{S'},\qquad
 \Sym^3(\tau)(u^3+v^3)=\lambda(u'^3+v'^3)
\end{equation}
を満たす組と一致する。
\end{theorem}

\begin{proof}
\textbf{必要性と混合項の消滅。}
$\Phi:A(S,g_S)\to A(S',g_{S'})$ を次数付き同型とする。
命題\ref{thm:intrinsic-sklyanin}により中心平面を保存し、商とquadratic cover上で
$\phi:S\xrightarrow{\sim}S'$ を誘導する。
$X=\phi(x),Y=\phi(y),Z=\phi(z)$ とおけば
\[
 \Phi(x)=X+\ell_x,\quad \Phi(y)=Y+\ell_y,\quad
 \Phi(z)=Z+\ell_z,\qquad \ell_i\in ku'\oplus kv'
\]
と書ける。
$\gamma x^2+\alpha yz+\beta zy=0$ の像を取り、
$S'_1\otimes(ku'\oplus kv')$ 成分だけを見ると
\[
 2\gamma X\ell_x+(\alpha+\beta)(Y\ell_z+Z\ell_y)=0.
\]
$X,Y,Z$ は独立であり、$2\gamma\ne0$ なので $\ell_x=0$。
巡回した二つの関係から $\ell_y=\ell_z=0$ も従う。
よって混合項は存在しない。
なお $S\hookrightarrow A$ は、$u$ のmonicな三次関係による
$S[v]$-加群分解から単射である。

補題\ref{lem:canonical-central-line}により $\phi(g_S)=\lambda g_{S'}$ である。
$S'[u',v']$ の次数3で商の核は
$k(g_{S'}+u'^3+v'^3)$ なので、三次関係の像が零になる条件は
\eqref{eq:five-generator-iso-pairs}に一致する。

\textbf{十分性。}
$\phi:S\xrightarrow{\sim}S'$ を取り、
$\phi(g_S)=\lambda g_{S'}$ とする。
$k$ は代数閉体なので $\delta^3=\lambda$ を満たす $\delta\in k^\times$ を選べる。
$u\mapsto\delta u'$、$v\mapsto\delta v'$ とすれば、三次関係の像は
$\lambda(g_{S'}+u'^3+v'^3)$ であり、$\phi$ は求める同型へ延長する。
一般の\eqref{eq:five-generator-iso-pairs}についても同じ議論が成り立つ。
\end{proof}

同じ証明は、$u^3+v^3$ を相異なる三つの零点を持つ二変数三次式に置き換えても成り立つ。
代数閉体上では任意の相異なる三点の組は射影同値であり、
三次式のscalarも吸収できるため、この二変数三次式から追加の同型類のmodulusは生じない。

\subsection{固定Fermat曲線と四つの非可換三fold}

第\ref{sec:reconstruction}節と同様に
$E_0=V(x^3+y^3+z^3)$、$O=[1:-1:0]$ とする。
$p\in E_0\setminus E_0[3]$ に対して $S(p)=A(E_0,T_p)$ と書き、
その標準次数3の中心元の任意の非零倍を $g_p$ として
\[
 \mathcal A(p)=S(p)[u,v]/(g_p+u^3+v^3)
\]
とおく。その次数付き同型類は $g_p$ の非零scalar倍に依存しない。
第9節の $A'_h$ の幾何学的一般ファイバーは、この形に属する。

\begin{theorem}[固定楕円曲線上での完全な次数付き分類]\label{thm:four-five-generator}
$p,q\in E_0\setminus E_0[3]$ に対して
\[
 \mathcal A(p)\simeq_{\mathrm{gr}}\mathcal A(q)
 \quad\Longleftrightarrow\quad q\in\mu_6\cdot p.
\]
従って $t=a^3$ は、この構成に属する五生成元代数の幾何学的な次数付き同型類を
分離する座標である。
一般の非被約点の局所代数 $J_a$ の同型類に対して、
対応する五生成元代数の次数付き同型類はちょうど四つである。
\end{theorem}

\begin{proof}
命題\ref{thm:five-generator-isom}と
\eqref{eq:fixed-sklyanin-isom}から最初の同値が従う。
Sklyaninの同型判定そのものは、幾何対の射影共役による
\cite[Theorem 2.8, Theorem 4.16]{IMRelations}である。
$t=a^3$ は $E_0/\mu_6$ の座標であり、
\eqref{eq:fat-j-degree-four}は次数4であるから、最後の個数も従う。
三乗根の選択から新たな同型類が増えないことは、
命題\ref{thm:five-generator-isom}の十分性が保証している。
\end{proof}

この結果はpoint schemeの抽象スキームだけでは五生成元代数を分類できないことも示す。
一般の同じ $j_{\mathrm{fat}}$ に対応する四つの代数のpoint schemeは、
shiftを忘れれば全て $E_0\sqcup3\operatorname{Spec}J_a$ と同型である。
その楕円曲線成分上のpoint-module shiftは異なる平行移動を記録している。

また、第\ref{sec:points}節の再構成
\[
 A\simeq\bigoplus_{n\ge0}
 \operatorname{Hom}_{\qgr A}(\pi A,\pi A(n))
\]
により、構造対象と次数シフトを保つtailsの三つ組
$(\qgr A,\pi A,(1))$ の同値についても同じ同型判定が成り立つ。
ここで $S[u,v]$ はAS-regularであり、中心的な正則三次式による商 $A$ は
AS-Gorenstein、depth 4なので、この再構成に必要な低次数の局所コホモロジーは消滅する。
次数シフトの自然同型は積との整合性を含めて保つものとする。
これは、構造対象や偏極を指定しない圏同値の分類ではない。

\subsection{自己同型群の大きさと二重化する点}

$\operatorname{PAut}_{\mathrm{gr}}(A)
=\operatorname{Aut}_{\mathrm{gr}}(A)/k^\times$ と書く。
分母は全ての次数1生成元を同じscalar倍にする作用である。
$G_p=\{\alpha\in\operatorname{Aut}(E,O):\alpha(p)=p\}$ とおく。

\begin{theorem}[射影的な次数付き自己同型群]\label{thm:graded-automorphisms}
滑らかな平面三次曲線 $E$、$\OO_E(1)=\OO_E(3O)$ と
$p\in E\setminus E[3]$ に対して、上の代数 $A$ は完全列
\[
 1\longrightarrow\mu_3
 \longrightarrow\operatorname{PAut}_{\mathrm{gr}}(A)
 \longrightarrow(E[3]\rtimes G_p)\times S_3
 \longrightarrow1
\]
を持つ。特に
\[
 |\operatorname{PAut}_{\mathrm{gr}}(A)|=162\,|G_p|.
\]
固定Fermat曲線の場合、この位数は $t\ne4$ で162、$t=4$ で324である。
後者では $j_{\mathrm{fat}}=432$ となる。
\end{theorem}

\begin{proof}
命題\ref{thm:five-generator-isom}により全ての自己同型はblock diagonalである。
$S$ の射影的自己同型は、$(E,T_p)$ の射影的自己同型であり、
\cite[Theorem 4.12]{IMRelations}から
\[
 \operatorname{PAut}_{\mathrm{gr}}(S)=E[3]\rtimes G_p
\]
となる。二変数三次式の三つの零点の射影的自己同型群は $S_3$ である。
これらの任意の組の線形持ち上げを選ぶと、
$g$ と $u^3+v^3$ のscalar倍率は、第二の持ち上げを三乗根倍することで一致させられる。
従って右側への射は全射である。

両方の射影的作用が恒等なら、線形作用は
$S_1$ 上で $\lambda\id$、$\langle u,v\rangle$ 上で $\mu\id$ である。
三次関係を保つ条件は $\lambda^3=\mu^3$。
共通scalarを除けば $\mu/\lambda\in\mu_3$ だけが残るので、kernelを得る。

Fermatの場合は $\operatorname{Aut}(E_0,O)=\mu_6$ である。
位数3の自己同型で固定される点は3-torsionに入り、位数6の自己同型の固定点は $O$ だけである。
従って $p\notin E_0[3]$ で非自明なstabilizerを持つのは
$p\in E_0[2]\setminus\{O\}$ の場合だけで、そのstabilizerは $\{\pm1\}$ である。
命題\ref{thm:translation-double-cover}から、この三点の共通する $t$ は4である。
最後に $\mathcal J(4)=432$ を代入して得る。
\end{proof}

従って $j_{\mathrm{fat}}=432$ の四つの次数付き同型類のうち、
$t=4$ に対応する一つは射影的自己同型群が324元であり、残る三つは162元である。
実際、対応する四次式は
\[
 432\,t(t+8)^3-110592(t-1)^3
 =432(t-4)(t^3-228t^2+48t-64)
\]
と因数分解する。右側の三次式の判別式は $-2^{14}3^{11}\ne0$ であり、
$t=4$ での値は $-3456\ne0$ なので、四つの $t$ は相異なる。
同じ非被約局所環を持つ代数でも、自己同型群の大きさは異なり得る。

\subsection{DVR上の形式族への適用と基底体の制限}

$R=\CC[[h]]$ 上の次数ごとに有限自由な二つのSklyanin型族 $S_R,S'_R$ を考える。
ここでは、これらは次数1で生成され、次数2の関係のみで定義される $R$-代数とする。
その幾何学的一般ファイバーの点曲線は滑らかとし、
次数3の中心元 $G,G'$ は一般ファイバーで標準次数3の中心元の線を張るとする。
さらに $G,G'$ は次数3の自由加群のprimitiveな元であると仮定する。
この条件は第9節の $G'_h$ で成立する。

このとき、相対次数付き代数
\[
 A_R=S_R[u,v]/(G+u^3+v^3),\qquad
 A'_R=S'_R[u',v']/(G'+u'^3+v'^3)
\]
についても
\[
 A_R\simeq_{R,\mathrm{gr}}A'_R
 \quad\Longleftrightarrow\quad
 S_R\simeq_{R,\mathrm{gr}}S'_R
\]
が成り立つ。
実際、一般ファイバーの中心の計算と $R$-torsion freenessから
$Z(A_R)_1=Ru\oplus Rv$ となり、相対quadratic coverで必要性を得る。
逆に $\phi:S_R\to S'_R$ が与えられると、一般ファイバーで
$\phi(G)=\lambda G'$、$\lambda\in\CC((h))^\times$ である。
両辺のprimitive性を $S'_{R,3}$ の基底で比較すれば
$\lambda\in R^\times$ と分かる。
全ての $R$ の単元は三乗根を持つので、中心変数の共通スケールで十分性を得る。
特殊ファイバーでの同一視を保つ場合には、$\lambda(0)=1$ に対して
定数項1の三乗根を選べばよい。

このDVR上の主張を、任意の非被約Artin基底上の同型関手の同一視へ拡張してはいない。
例えば第9節の族を $R_N=R/(h^N)$、$N\ge2$ へ制限すると、
交換子が $h$ で割り切れるため
\[
 h^{N-1}x,\qquad h^{N-1}y,\qquad h^{N-1}z
\]
も次数1の中心に入る。従って中心の形成はこの非平坦な基底変換と可換ではない。
また、任意の基底体では必要な三乗根が存在するとは限らず、
\eqref{eq:five-generator-iso-pairs}のscalar条件を保持しなければならない。
一般の $\CC((h))$ 上のformを、幾何学的な $t$ または $j$ だけで分類する主張ではない。

\section{適合部分空間の相対Picard記述}\label{sec:picard-compatibility}

本節では、7次元の適合部分空間を、平面三次曲線上の線形束の変形から
高次まで特徴付ける。第\ref{sec:evaluation-gluing}節の一本のFermat族についての記述を、
任意の滑らかなcubic threefoldと任意の非定数な適合形式弧へ拡張する。

\subsection{全ての平面切断にわたる形式Picard}

$\mathcal G=\Gr(2,V)$、$\mathcal S\subset V\otimes\OO_{\mathcal G}$ をtautological束とし、
\[
 T=H^3(X,\OO_X(-2)),\qquad
 W=\Ext^1(E,E)=\Lambda^2V\otimes T,\qquad
 \mathcal D=\det\mathcal S\otimes T
\]
とおく。$U\in\mathcal G$ の定める平面三次曲線を
$C_U=X\cap\PP((V/U)^*)$ と書く。
$X$ は平面を含まないので、全ての $C_U$ は平面三次曲線である。
本節では、その被約性、既約性、滑らかさを仮定しない。

Universal曲線 $p:\mathcal C\to\mathcal G$ はflatで、
$\mathcal C\subset X\times\mathcal G$ は余次元2の相対正則埋め込みである。
Koszul列とcohomology and base changeから
\[
 p_*\OO_{\mathcal C}=\OO_{\mathcal G},\qquad
 R^1p_*\OO_{\mathcal C}=\mathcal D,\qquad
 R^{\ge2}p_*\OO_{\mathcal C}=0
\]
を得る。相対Picardの $L=\OO_{C_U}(1)$ というsectionに沿う形式完備化を
$\widehat{\mathcal P}$ とする。
このsectionでtranslateすると、$\widehat{\mathcal P}$ は相対次元1の可換形式群である。
曲線上の線形束の障害群 $H^2$ は零であり、Picardの滑らかさは
\cite[Tag 0DPJ]{StacksPicardSmooth}によっても得られる。
標数0の形式対数\cite[Appendix A.2, Proposition A.2.2]{PicardFormal}により、
接写像が恒等となる一意の形式群同型を得る：
\begin{equation}\label{eq:relative-picard-log}
 \widehat{\mathcal P}
 \simeq\widehat{\operatorname{Tot}(\mathcal D)}_{\mathcal G}
\end{equation}

線形束 $L$ の各Artin変形に対して
\[
 \mathcal P_L=i^*(j_*L)[-1]
\]
を取る。$j$ は正則埋め込みなので $j_*L$ はrelative perfectであり、
相対mutationはperfect性と基底変換を保つ。
\eqref{eq:picard-special-object}は、特異または非被約な $C_U$ に対しても
Koszul列からそのまま成立する。従ってzero section全体が $e=[E]$ に写り、
\[
 \Phi:\widehat{\mathcal P}\longrightarrow\Def_E
\]
を得る。全体上のPoincar\'e束の存在は仮定しない。
局所的な線形束の選択の差は基底からの線形束であり、粗対象モジュライへの射は降下する。

\begin{theorem}[相対Picard像とそのblowup]\label{thm:picard-formal-image}
$\Phi$ のscheme-theoretic imageは、$\Def_E$ の中の閉形式部分スキーム
$P$ である。$P$ は整な7次元空間で、抽象的にはGrassmannian錐の頂点完備化に同型である。
また
\[
 \widehat{\mathcal P}\simeq
 \widehat{\operatorname{Bl}_eP}_{\mathcal G}
\]
となり、$\Phi$ はそのblowdownと $P\hookrightarrow\Def_E$ の合成である。
cohomologyの接座標 $W=\Lambda^2V\otimes T$ では、
$\operatorname{TC}_eP$ はPl\"ucker錐である。
\end{theorem}

\begin{proof}
補題\ref{lem:picard-gluing-ks}の接写像の計算は、全ての平面三次曲線に適用できる。
実際、$h^0(\OO_{C_U})=h^1(\OO_{C_U})=1$ であり、
$0\ne\xi\in H^1(\OO_{C_U})$ とのcup積は $1\mapsto\xi$ という非零写像である。
この点を二対象模型で確認する。
Euler列とKoszul列から
$F_0=\operatorname{Cone}(\OO_X(-1)\otimes\det U\to Q)$ であり、
線束のArtin変形の評価核 $F_L$ は、このperfect対象の変形となる。
full dg部分圏 $\langle\OO_X(-1),Q\rangle$ とambient圏の自己Hom dg代数は同じなので、
一次Kodaira--Spencer類は、同じ最小模型の $(\dot f,\dot t)$ として計算できる。
ここで必要なのは $\epsilon^2=0$ 上の一次計算だけである。
$F_L\in\OO_X^\perp$ なので、Euler列と評価列から自然に
\[
 \RHom(Q,F_L)
  \simeq\RHom(\OO_X(1),F_L)[1]
  \simeq R\Gamma(C_U,L(-1))
\]
となる。左辺は二対象模型で $[R\xrightarrow{t}R]$ であり、$f$ は微分に寄与しない。
右辺の一次微分は上のcup積なので、$t$ の一次項は非零である。
$\det U$ を保持すると、この方向は $\det U\otimes T$ に値を取る。
二対象のmutationの公式から、zero sectionの法方向の写像は
\[
 \mathcal D=\det\mathcal S\otimes T
    \lhook\joinrel\longrightarrow W\otimes\OO_{\mathcal G}
\]
というtautologicalな包含になる。

\eqref{eq:relative-picard-log}により
\[
 R_P:=\Gamma(\widehat{\mathcal P},\OO)
   =\prod_{n\ge0}H^0(\mathcal G,\mathcal D^{-n})
\]
は、Pl\"ucker section ringの頂点完備化である。
GrassmannianのPl\"ucker埋め込みのprojective normalityと次数1での生成から、
$R_P\simeq\CC[[a_{ij}]]/(q_{ijkl})$ は正規な整域であり、
$\dim R_P=7$、$\mathfrak m_{R_P}/\mathfrak m_{R_P}^2=W^*$ である。

$\Def_E=\Spf A$、$A\simeq\CC[[z_1,\ldots,z_{10}]]$ と書く。
$\Phi$ が誘導する連続環準同型 $A\to R_P$ は、上の接写像から余接空間上同型になる。
完備局所環のNakayamaによりこの射は全射である。
従って $P=\Spf R_P\hookrightarrow\Def_E$ が閉形式像となる。

Pl\"ucker錐の頂点のblowupはtautological線束の全空間である。
trace線を含めれば、そのzero sectionに沿う完備化は
$\widehat{\operatorname{Tot}(\mathcal D)}_{\mathcal G}$ となる。
affine形式空間への射は大域切断の環準同型で決まるので、
元の $\Phi$ はこの標準的な縮約を通じて $P$ へ写る。
これによりblowupの主張も得られる。
\end{proof}

\subsection{Picard族に沿う最後のgluing対象の持ち上げ}

ここでは接錐の比較を用いず、Picard族そのものが適合条件を満たすことを証明する。
最後のgluing対象の障害を、特異な三次式も含む普遍族上で考えることが要点である。

\begin{theorem}[Picard族の全次数の適合性]\label{thm:picard-compatible}
$X\subset\PP(V^*)$ を滑らかなcubic threefold、
$U\subset V$ を2次元部分空間、$C=X\cap\PP((V/U)^*)$ を平面切断とする。
$j:C\hookrightarrow X$、$\pi=i^*$、
\[
 T=H^3(X,\OO_X(-2)),\qquad D_U=\det U\otimes T
\]
とおく。可換局所Artin $\CC$-代数 $R$ 上で、
$L_R$ を $\OO_C(1)$ の線形束としての変形とする。
このとき
\[
 E_{L_R}=D_U^{-1}\otimes\pi(j_*L_R)[-1]\in\mathcal K_R
\]
に沿って、$G$ は $\mathcal B_{E_{L_R}}$ のperfect対象に持ち上がる。
従ってPicard形式近傍の射は、閉形式部分空間
$\widehat C_c\subset\Def_E$ を経由する。
この主張は平面切断 $C$ が特異な場合にも成立する。
\end{theorem}

\begin{proof}
\textbf{普遍族を取る基底。}
$V=U\oplus Z$、$\dim Z=3$ と分解し、
\[
 \Sigma_U=\{F\in\Sym^3V:F|_Z\ne0\}
 \simeq\Aff^{25}\times(\Aff^{10}\setminus\{0\}),
 \qquad H_U=GL(U)\times GL(Z)
\]
とおく。$H_U$ はcharacteristic zeroでreductiveである。
除いた部分は余次元10の線形部分空間なので、局所cohomologyの計算から
\begin{equation}\label{eq:picard-universal-base-cohom}
 \Gamma(\Sigma_U,\OO)=\CC[\Sym^3V],\qquad
 H^r(\Sigma_U,\OO)=0\quad(1\le r\le8)
\end{equation}
を得る。特に係数を判別式で割ることは許していない。

$\mathcal X\to\Sigma_U$ を普遍三次超曲面、
$q:\mathcal C\to\Sigma_U$ をその $\PP(Z^*)$ による切断とする。
両者はflatかつproperであり、$\mathcal C\hookrightarrow\mathcal X$ は
余次元2の正則埋め込みである。
Euler列、制限列、Gorenstein双対性による第\ref{sec:gluing-objects}節と
補題\ref{lem:universal-tangent-lift}の計算は、
特異・非被約なファイバーについてもperfect対象に対して成立する。
trace線を保持すると、相対自己Extは
\[
 \mathcal E xt^r_{\mathcal B/\Sigma_U}(G,G)=
 \begin{cases}
  \OO_{\Sigma_U}&r=0,\\
  \Sym^2V\otimes T&r=2,\\
  0&\text{その他}
 \end{cases}
\]
であり、任意の基底変換と可換である。
普遍制限列の境界写像から、同変な同一視
\[
 T\simeq(\det V)^*\otimes\OO_{\Sigma_U},\qquad
 R^1q_*\OO_{\mathcal C}\simeq D_U
       =\det U\otimes(\det V)^*\simeq(\det Z)^*
\]
を得る。また $q_*\OO_{\mathcal C}=\OO_{\Sigma_U}$、
$R^rq_*\OO_{\mathcal C}=0$ ($r\ge2$) である。

\textbf{同変なPicardパラメータ。}
$R_n=\bigoplus_{r=0}^nD_U^{-r}$ を次数 $n$ で切った対称代数とし、
$\Sigma_n=\operatorname{Spec}_{\Sigma_U}R_n$ とする。
$R_n\twoheadrightarrow R_{n-1}$ の核は $I_n=D_U^{-n}$ で、
正次数のイデアルにより零化される。
$\OO_{\mathcal C}(1)$ の線形束としての持ち上げの障害は
\[
 H^2(\mathcal C,q^*I_n)^{H_U}=0
\]
に入る。消滅はLeray列と\eqref{eq:picard-universal-base-cohom}による。
同変な場合も、不変部分を取る関手が完全なので同じ結論を得る。
一次変形の類を $D_U\otimes D_U^*$ の恒等元に取れば、
$H_U$-同変な線形束の整合的な族 $L_n$ を逐次構成できる。
各ファイバーでは $h^1(\OO_C)=1$、$H^2(\OO_C)=0$ なので、
この族はPicard形式近傍の座標を与える。
ここでは相対Picardスキーム上の大域的なPoincar\'e束の存在を仮定していない。

\eqref{eq:picard-special-object}によって
\[
 \mathcal E_n=D_U^{-1}\otimes\pi(j_*L_n)[-1]
\]
は $E$ の同変な変形である。
$j$ が正則埋め込みなので $j_*L_n$ はperfectであり、相対mutationもperfect性を保つ。
これを使って $\mathcal B_n=\langle\mathcal K_{\Sigma_n},Q\rangle_{\mathcal E_n}$
を構成する。

\textbf{変化するdg圏での持ち上げ障害。}
square-zero拡大に沿うdg加群の持ち上げを、局所的なsemifree模型で計算する。
次数付きの持ち上げに微分と加群作用を持ち上げると、その平方の失敗は
$I_n$ を係数とする次数2の自己準同型cocycleである。
持ち上げを変えるとcoboundaryだけ変わり、その類が零なら次数1の補正で
平方を零にできる。
これは微分の持ち上げの障害計算\cite[Theorem 3.8]{LowenObstruction}、
およびMaurer--Cartan模型\cite[Proposition 7.1]{ELOI}のcurvature計算である。
ここでは既に指定された $\mathcal B_n$ 上のbar模型に適用するので、
圏が一定であることを仮定していない。

この構成の還元が導来還元であることも確認する。
affine基底上で $\mathcal B_n$ のcofibrant模型を取り、前段階の加群をsemifreeに取る。
その各評価は基底上h-projectiveかつperfectなので、有界な有限projective複体と
homotopy同値である。
微分を補正した後も次数付き評価はprojectiveであり、
crude lifting lemma\cite[Corollary 3.11]{LowenObstruction}により、
このhomotopy同値を有界な有限projective複体の持ち上げへ輸送できる。
従って新しい評価はK-flatで、通常の還元と導来還元が一致する。
圏のcofibrant模型の基底変換と導来テンソル積の結合則から、
加群としての導来還元も前段階の加群になる。
無限複体の次数付きflat性だけからこの結論を出しているわけではない。

\v Cech descentも含めると、$G_{n-1}$ の持ち上げ障害は
\[
 \mathbb H^2\bigl(\Sigma_U,
       \RHom_{\mathcal B/\Sigma_U}(G,G)\otimes I_n\bigr)^{H_U}
\]
に入り、存在すれば同型類の差は同じ複体の次数1に入る。
相対自己Extと\eqref{eq:picard-universal-base-cohom}から、次数1は零、次数2は
\begin{equation}\label{eq:picard-obstruction-equivariants}
 \left[\CC[\Sym^3V]\otimes\Sym^2V\otimes(\det V)^*
                         \otimes D_U^{-n}\right]^{H_U}
\end{equation}
となる。$H^3(\Sigma_U,\OO)=0$ により、local-to-global列の
次数2からのtransgressionもない。
各段階を同変な加群の圏で構成しているので、同型類の一意性だけから
coherentな群作用の存在を推測する必要はない。

\textbf{全ての次数で障害が消えること。}
\eqref{eq:picard-obstruction-equivariants}の元を、多項式同変写像
\[
 b:\Sym^3V\longrightarrow
       \Sym^2V\otimes(\det V)^*\otimes D_U^{-n}
\]
と見る。$a_s=(s\id_U,\id_Z)$ の作用について、
$F$ の $\Sym^3Z$、$U\otimes\Sym^2Z$、$\Sym^2U\otimes Z$、$\Sym^3U$
成分のweightは、それぞれ $0,1,2,3$ である。
一方、$D_U$ のweightは0であり、目標の
\[
 \Sym^2U\otimes(\det V)^*\otimes D_U^{-n},\quad
 U\otimes Z\otimes(\det V)^*\otimes D_U^{-n},\quad
 \Sym^2Z\otimes(\det V)^*\otimes D_U^{-n}
\]
のweightは $0,-1,-2$ である。
従って多項式写像 $b$ の負weight成分は零であり、残る成分は
$F|_Z$ のみを用いて $\Sym^2U$ の成分に値を取る。
ところが $SL(U)$ は $F|_Z$、$D_U$、$(\det V)^*$ に自明に作用し、
$\Sym^2U$ は $SL(U)$-不変ベクトルを持たない。従って $b=0$ である。
この議論は全ての $n\ge1$ に成立する。

\textbf{perfect性とArtin基底への特殊化。}
以上から $G_0=G$ を出発点として、整合的な同変dg加群 $G_n$ を得る。
特異な三次式を含む普遍基底上では、各評価が基底上perfectであるという
pseudo-perfect性だけを用いれば十分である。
評価複体の導来還元は $G$ の評価でありperfectなので、
冪零イデアルに関するperfect性の判定\cite[Lemma 15.80.4, Tag 07LU]{StacksNilPerfect}
により、この性質は各段階で保たれる。

最後に滑らかな $X_F$ に制限する。
$\mathcal K_F$ はsmooth properで、$\mathcal E_n$ はperfectなので、
$\mathcal B_n|_F$ もsmooth properである。
このgluingのsmoothnessはArtin基底上でも成立する。
実際、$A$ を $\mathcal K_F\otimes R_n$ のdg模型、$M$ を $\mathcal E_n$ の
perfect加群とし、上三角代数を $T_{\mathrm{tri}}$、その頂点冪等元を $e,f$ と書く。
対角bimoduleの三角
\[
 \begin{split}
 T_{\mathrm{tri}}e\otimes_A^{\mathbf L}M\otimes_{R_n}^{\mathbf L}fT_{\mathrm{tri}}
 &\longrightarrow
 (T_{\mathrm{tri}}e\otimes_A^{\mathbf L}eT_{\mathrm{tri}})
 \oplus(T_{\mathrm{tri}}f\otimes_{R_n}^{\mathbf L}fT_{\mathrm{tri}})\\
 &\longrightarrow\Delta_{T_{\mathrm{tri}}}
 \end{split}
\]
の左項と中項は、$A$ のsmoothnessと $M$ のperfect性によりperfectである。
従って対角もperfectであり、体上という仮定は必要ない。
perfectな対角bimoduleの有限分解とテンソルすれば、
pseudo-perfectな $G_n|_F$ はperfectな $\mathcal B_n|_F$-加群になる。
任意のArtin線形束変形は、上のPicardパラメータ族のある有限段階の引き戻しなので、
その引き戻しが所要のperfectな $G$ の持ち上げを与える。
\end{proof}

\subsection{適合条件との高次までの一致}

\begin{theorem}[適合部分空間の幾何学的記述]\label{thm:picard-compatibility-equality}
命題\ref{thm:algebraic-compatibility}の適合部分空間の完備化 $\widehat C_c$ は、
命題\ref{thm:picard-formal-image}の $P$ に、
$\Def_E$ 内の閉形式部分スキームとして一致する。
従って
\[
 \widehat{\mathcal P}
 \simeq\widehat{\operatorname{Bl}_c C}_{\mathcal G}.
\]
\end{theorem}

\begin{proof}
命題\ref{thm:picard-compatible}により、固定した各 $U\in\mathcal G$ のPicard形式曲線は
適合条件を満たす。
適合イデアルの任意の元を $\widehat{\mathcal P}$ に引き戻し、形式対数の次数ごとに展開する。
その $n$ 次係数は $H^0(\mathcal G,\mathcal D^{-n})$ の切断であり、
全ての閉点 $U$ で零になる。$\mathcal G$ は被約なので、この係数自体が零である。
従って $\Phi$ 全体が $\widehat C_c$ を経由し、$P\subset\widehat C_c$ を得る。

既存の形式変形の計算により $\widehat C_c$ は整な7次元空間であり、
命題\ref{thm:picard-formal-image}により $P$ も整な7次元空間である。
正則完備局所環 $A$ 内の対応する二つの素イデアルは、包含関係を持ち高さが等しいので一致する。
blowupの主張は命題\ref{thm:picard-formal-image}から従う。
この証明は、全次数での包含と等次元性を使っており、接錐の一致だけからは結論していない。
\end{proof}

\subsection{相対cohomologyの単純零点}

\begin{lemma}[Picardのidentity sectionの重複度]\label{lem:relative-picard-theta}
$L=\OO_{C_U}(1)$ のsectionの形式近傍で、$Rp_*L(-1)$ は局所的に
\[
 [\,\mathcal L_0\xrightarrow{\,d\,}\mathcal L_1\,]
 \qquad\text{（次数 }0,1\text{）}
\]
という二つの線束の間の複体で表せる。
zero sectionの局所形式パラメータを $t$ とすると、線束を自明化した表示で
$d=t\cdot u$、$u$ は単元である。
この主張は特異または非被約な平面三次曲線を含む。
\end{lemma}

\begin{proof}
$p$ はproper flatな曲線の族であり、$L(-1)$ は線束なので、
$Rp_*L(-1)$ はperfectかつTor amplitude $[0,1]$ を持つ。
zero sectionの各点では、導来ファイバーは $R\Gamma(\OO_{C_U})$ であり、
次数0と1にそれぞれ1次元のcohomologyを持つ。
局所的なperfect表示から可逆な微分成分を消去すると、上の二項表示になる。

zero sectionに沿ってcohomologyは $\OO_{\mathcal G}$ と $\mathcal D$ であり、
基底変換と可換なので $d$ はそこで恒等的に零になる。従って $d=tu$ と書ける。
$d$ の法方向の一次項はPicardのKodaira--Spencer類とのcup積
\[
 H^0(\OO_{C_U})\longrightarrow H^1(\OO_{C_U}),\qquad1\longmapsto\xi
\]
である。これは非零の法方向に対して同型なので、$u$ はzero section上で単元であり、
形式近傍全体でも単元となる。
\end{proof}

この補題は、被約でない曲線上の任意の非自明な次数0の線束がacyclicであると
仮定するものではない。必要な消滅は、identity sectionの近傍の相対複体から直接得られる。

\subsection{全ての非定数な適合形式弧の一般ファイバー}

\begin{theorem}[一般gluing sheafと非自由支持の決定]\label{thm:all-compatible-arcs}
$X$ を任意の滑らかなcubic threefold、$R=\CC[[h]]$、
$\gamma:\Spf R\to\widehat C_c$ を非定数の適合形式弧とする。
対応する貼り合わせ対象を $E_R$ とすると、次を満たす平面の族 $U_R$ と
線形束 $L_R$ が得られる。
\[
 C_R=X_R\cap\PP((V_R/U_R)^*),\qquad
 L_R|_{C_0}=\OO_{C_0}(1).
\]
一般ファイバーでは
\begin{equation}\label{eq:all-arcs-evaluation}
 0\longrightarrow E_\eta
 \longrightarrow H^0(C_\eta,L_\eta)\otimes\OO_{X_\eta}
 \longrightarrow j_*L_\eta\longrightarrow0
\end{equation}
となる。$E_\eta$ はGieseker-stableな階数3のtorsion-free層であり、
\[
 \operatorname{Fitt}_3(E_\eta)=\mathcal I_{C_\eta},\qquad
 E_\eta^{**}\simeq\OO_{X_\eta}^{\oplus3},\qquad
 E_\eta^{**}/E_\eta\simeq j_*L_\eta.
\]
従って非局所自由部分の水平閉包は、スキームとして $C_R$ に等しい。

さらに、cohomologyの接座標で $\gamma$ の最初の非零係数を $h^m a_m$ とすると、
$a_m$ は分解可能であり、対応する $U_0\in\Gr(2,V)$ に対して
\[
 C_0=X\cap\PP(\operatorname{ann}U_0)=Z(s_{a_m}).
\]
一般の層は局所自由でもACMでもなく、normalized Ulrichではない。
\end{theorem}

\begin{proof}
$\gamma$ が非定数なので、$\mathfrak m_c\OO_R=(h^m)$ は非零の可逆イデアルである。
blowupの普遍性により $\gamma$ は一意に $\operatorname{Bl}_cC$ へ持ち上がる。
命題\ref{thm:picard-compatibility-equality}により、この持ち上げは相対Picardの形式点となり、
$U_R$ と $L_R$ の形式族を与える。
Grassmannianの射影性と、properな $C_R$ 上のGrothendieck existenceにより、
これらは実際の $R$ 上の平面族と線束に実現する。
基底からの線束の曖昧さは局所環 $R$ 上では自明である。

Picardの形式パラメータ $t_R$ は非零である。
補題\ref{lem:relative-picard-theta}から
\[
 R\Gamma(C_R,L_R(-1))\simeq[R\xrightarrow{\,t_Ru\,}R]
\]
となり、一般ファイバーではacyclicである。
また $h^0(\OO_{C_0}(1))=3$、$h^1(\OO_{C_0}(1))=0$ なので、
$p_*L_R$ は階数3の自由加群で、$R^1p_*L_R=0$。
特殊ファイバーでの大域生成性、Nakayamaとproper性により評価射は全射となる。
その核を $F_R$ とすれば
\[
 0\longrightarrow F_R\longrightarrow p_*L_R\otimes\OO_{X_R}
 \longrightarrow j_*L_R\longrightarrow0
\]
は $R$-flatな層の完全列であり、$\pi F_R\simeq E_R$ である。
一般ファイバーで $F_\eta$ と $F_\eta(-1)$ はacyclicなので
$F_\eta\in\Ku(X_\eta)$、従って $E_\eta=F_\eta$ を得る。
安定性は、既存のstable対象モジュライのchartと
\cite[Proposition 5.5]{FLZ}から従う。

$C_R\subset X_R$ は余次元2の正則埋め込みである。
線束を局所自明化し、評価の一成分を単元に直せば
$F_R\simeq\mathcal I_{C_R}\oplus\OO_{X_R}^2$ となる。
これがFittingイデアル、二重双対とその商を与える。
$C_R$ は $R$-flatなので、$C_\eta$ のスキームとしての閉包は $C_R$ である。
blowupの例外点は $[a_m]$ であり、法方向のtautologicalな接写像から
特殊平面は $U_0$、特殊曲線は表示の $C_0$ となる。

最後に、特殊ファイバーの $L_0(-2)=\OO_{C_0}(-1)$ は大域切断を持たない。
これは非被約な平面三次曲線についても、$\PP^2$ 上の完全列
$0\to\OO(-4)\to\OO(-1)\to\OO_{C_0}(-1)\to0$ から従う。
半連続性により $H^0(C_\eta,L_\eta(-2))=0$ であり、
$\chi(C_\eta,L_\eta(-2))=-3$ だから $h^1(C_\eta,L_\eta(-2))=3$。
\eqref{eq:all-arcs-evaluation}と $H^1(\OO_X(-2))=H^2(\OO_X(-2))=0$ から
\[
 h^2(X_\eta,E_\eta(-2))=h^1(C_\eta,L_\eta(-2))=3
\]
となり、ACM消滅を満たさない。
\end{proof}

\begin{theorem}[適合部分空間は非局所自由境界に含まれる]
\label{thm:compatibility-in-bn-boundary}
第\ref{sec:sextic-boundary}節のBrill--Noether境界 $D$ に対して、
$\widehat C_c\subset\widehat D_e$ が閉形式部分スキームとして成立する。
\end{theorem}

\begin{proof}
固定した各 $U\in\mathcal G$ のPicard形式曲線 $\CC[[t]]$ を取る。
補題\ref{lem:relative-picard-theta}と評価核の表示により、
その一般対象の非局所自由部分は非空な平面三次曲線 $C_U$ である。
従って一般点はBN incidenceの像 $D$ に属する。
$D$ の局所方程式 $d$ は $\CC((t))$ で零になるので、整域 $\CC[[t]]$ でも零である。
$\Phi^*d$ の形式対数の各次数の係数は、$\mathcal G$ 上の切断であり、
全ての閉点 $U$ で零になる。従って各係数は零であり、
$d$ は $P=\widehat C_c$ 上で零となる。
この包含は全次数のPicard記述を使う帰結であり、接錐の包含だけから推論したものではない。
\end{proof}

\begin{remark}
全 $R$ 上の等式は $E_R\simeq\pi F_R$ である。
特殊ファイバーでは $F_0$ は層、$E$ は二項複体なので、
一般ファイバーでの $E_\eta=F_\eta$ と区別する。
本節は適合部分空間そのものの相対Picard記述を与えるが、
Brill--Noether境界の特異部分全体が $C$ に一致することまでは主張しない。
\end{remark}

\section{評価核の点における境界：相対三次錐と解消}
\label{sec:local-bn-cubic}

前節は $e=[E]$ における適合部分空間を形式的に記述した。
本節では、その一般の評価核の点で、非局所自由な境界そのものを調べる。
以下、$C_0=X\cap\PP^2$ は滑らかな平面三次曲線、
$L\in\Pic^3(C_0)$、$L\ne\OO_{C_0}(1)$ とし、
$F=F_L$ を\eqref{eq:evaluation-kernel}の評価核とする。
文字 $C_0$ は曲線を表し、適合部分空間 $C$ と区別する。

\subsection{評価核の変形と三つの法方向}

$W_L=H^0(C_0,L)$ とし、曲線上の評価核、法束、および二つの線形束を
\[
 M_L=\ker(W_L\otimes\OO_{C_0}\to L),\qquad
 N=N_{C_0/X}=\OO_{C_0}(1)^{\oplus2},\qquad
 B=L^{-1}\otimes\det N
\]
と書く。さらに
\begin{equation}\label{eq:local-bn-bundle}
 \mathcal B_L=M_L\otimes B,\qquad
 \mathcal M_L=\det\mathcal B_L=\OO_{C_0}(4)\otimes L^{-3}
\end{equation}
とおく。$\mathcal B_L$ は階数2・次数3、$\mathcal M_L$ は次数3である。
ここで線形束の表示では、$\det W_L$ および $\det H^0(\mathcal B_L)$ という
基底からの1次元因子を局所自明化している。
自明化を選ばない場合、$\det M_L=\det W_L\otimes L^{-1}$ であり、
以下の評価列のkernelにも $\det H^0(\mathcal B_L)$ の因子を保持する。
これらの選択は射影埋め込みと次数に影響しない。

\begin{theorem}[評価核の滑らかな変形空間]\label{thm:evaluation-local-moduli}
$F$ のscalarを除いた変形空間 $U_F$ は滑らかな10次元の局所空間である。
平面三次曲線と次数3の線形束を動かす評価核の族は、
$F$ の近傍で滑らかな7次元の局所部分空間 $P\subset U_F$ を定める。
その法空間は自然に
\[
 N_{P/U_F,[F]}=H^0(C_0,\mathcal B_L),\qquad h^0(C_0,\mathcal B_L)=3
\]
と同定される。また $\Ext^2_X(F,F)=0$ である。
\end{theorem}

\begin{proof}
まず $M_L$ は安定である。実際、次数が $-1$ 以上の線部分束
$A\subset M_L\subset\OO_{C_0}^{\oplus3}$ を考える。
$\deg A>0$ なら $\Hom(A,\OO_{C_0})=0$ である。
$\deg A=0$ の場合、非零射が存在するには $A\simeq\OO_{C_0}$ が必要だが、
$H^0(M_L)=0$ に反する。
$\deg A=-1$ なら $h^0(A^*)=1$ なので、$A\to\OO_{C_0}^3$ は固定した
定数ベクトルを通る。その評価像が零であることは、この定数ベクトルに対応する
$L$ の切断が零であることを意味し、再び矛盾する。
従って任意の線部分束の次数は $-2$ 以下であり、$M_L$ は安定である。

$\mathcal B_L$ と $\mathcal B_L(-p)$ は正次数の安定束なので、
任意の $p\in C_0$ に対し $H^1(\mathcal B_L)=H^1(\mathcal B_L(-p))=0$ となる。
Riemann--Rochから $h^0(\mathcal B_L)=3$ であり、$\mathcal B_L$ は大域生成される。

評価列を $C_0$ に制限すると
\begin{equation}\label{eq:evaluation-restriction}
 0\longrightarrow L\otimes N^*\longrightarrow F|_{C_0}
   \longrightarrow M_L\longrightarrow0
\end{equation}
を得る。余次元2のregular immersionに対する局所Extの計算から
\[
 \mathcal Hom_X(F,\OO_X)=W_L^*\otimes\OO_X,
 \qquad \mathcal Ext^1_X(F,\OO_X)=j_*B
\]
であり、より高いsheaf Extは零である。また
\[
 0\longrightarrow M_L^*\otimes L\longrightarrow
 \mathcal Hom_X(F,j_*L)|_{C_0}\longrightarrow N\longrightarrow0,
 \qquad
 \mathcal Ext^1_X(F,j_*L)=j_*\det N.
\]
これらの束の $H^1$ は零なので
\[
 \dim\Hom_X(F,j_*L)=15,\quad
 \Ext^1_X(F,j_*L)=H^0(\OO_{C_0}(2)),\quad
 \Ext^1_X(F,W_L\otimes\OO_X)=W_L\otimes H^0(B).
\]
最後の二空間の次元は6と9である。それらの間の射は積写像であり、
\[
 0\longrightarrow\mathcal B_L\longrightarrow W_L\otimes B
   \longrightarrow\OO_{C_0}(2)\longrightarrow0
\]
と $H^1(\mathcal B_L)=0$ により全射である。
\eqref{eq:evaluation-kernel}に $\RHom_X(F,-)$ を適用し、$F$ の単純性を用いると
\[
 0\to\CC\to\CC^9\to\CC^{15}\to\Ext^1_X(F,F)
 \to\CC^9\to\CC^6\to\Ext^2_X(F,F)\to0
\]
を得る。従って $\dim\Ext^1_X(F,F)=10$、$\Ext^2_X(F,F)=0$ であり、
\begin{equation}\label{eq:normal-evaluation-sequence}
 0\longrightarrow T_{[F]}P\longrightarrow\Ext^1_X(F,F)
 \longrightarrow H^0(C_0,\mathcal B_L)\longrightarrow0
\end{equation}
の核は7次元である。

最後に、この核が実際の評価核の族の接空間であることを確認する。
平面の変形は $H^0(N)$ の6方向、線形束の変形は $H^1(\OO_{C_0})$ の1方向を持つ。
局所完全交叉の計算から $\Ext^1_X(j_*L,j_*L)$ の次元も7であり、
この平面・Picard族のKodaira--Spencer写像は同型である。
実際、評価列に $\RHom_X(-,j_*L)$ を適用し $H^1(L)=0$ を用いると
\[
 0\longrightarrow\operatorname{End}(W_L)/\CC
 \longrightarrow\Hom_X(F,j_*L)
 \longrightarrow\Ext^1_X(j_*L,j_*L)\longrightarrow0
\]
となる。
評価射の基底を選ぶと、この族の上に $\mathrm{PGL}(W_L)$ の8方向が加わる。
Quotの接空間 $\Hom_X(F,j_*L)$ は15次元なので、その全接方向が得られる。
核の変形への写像で基底変更の8方向を除くと、
\eqref{eq:normal-evaluation-sequence}の7次元の核になる。
従って評価核の族の接写像は単射である。
滑らかな族から滑らかな局所変形空間への射なので、
必要ならétale近傍に移り、$P$ を滑らかな局所閉部分空間として取れる。
形式局所環の水準では、これは接写像の単射と完備Nakayamaによる閉埋め込みである。
\end{proof}

\subsection{incidenceと三次曲線の別の埋め込み}

\begin{theorem}[評価核に沿う境界の三次錐]\label{thm:local-bn-cubic}
$U_F$ の十分小さい近傍で、非局所自由な層のincidenceを
$\mathcal I_F\subset X\times U_F$、そのscheme-theoretic imageを $D_F$ とする。
このとき以下が成立する。
\begin{enumerate}
 \item $\mathcal I_F$ は滑らかな9次元空間であり、
 $\mathcal I_F\times_{U_F}P=\mathscr C_P$ は評価核の曲線の相対族である。
 \item $D_F$ はintegralな因子であり、$P$ に沿う位数は3である。
 そのnormal coneの射影化は、各ファイバー上で
 \[
 C_0\hookrightarrow\PP H^0(C_0,\mathcal B_L)
 \]
 という完全線形系 $|\mathcal M_L|$ による三次埋め込みで与えられる。
 \item $\mathcal I_F\to D_F$ は $P$ に沿うblowupであり、
 その例外因子は $\mathscr C_P$ である。$D_F\setminus P$ は滑らかで、
 $\operatorname{Sing}(D_F)_{\mathrm{red}}=P$ である。
\end{enumerate}
ここでの等式はこの評価核の点の近傍についてであり、
$e=[E]$ を含む境界全体の特異部分の同定は主張していない。
\end{theorem}

\begin{proof}
まず $\mathcal B_L$ の評価列は
\begin{equation}\label{eq:local-bn-embedding}
 0\longrightarrow\mathcal M_L^{-1}
 \longrightarrow H^0(\mathcal B_L)\otimes\OO_{C_0}
 \longrightarrow\mathcal B_L\longrightarrow0
\end{equation}
となる。双対を取り $H^0(\mathcal B_L^*)=0$ を使うと、
3次元空間 $H^0(\mathcal B_L)^*$ と $H^0(\mathcal M_L)$ が同定される。
従って評価のkernel lineを送る写像は $|\mathcal M_L|$ による完全線形系であり、
次数3なので閉埋め込みである。

$p\in C_0$ で $F\simeq\mathcal I_{C_0}\oplus\OO_X^2$ と書けば、
普遍族の局所表示は $[\OO\to\OO^4]$ で、特殊点の列ベクトルは $(f,g,0,0)$ である。
ここで $(f,g)$ は $C_0$ のイデアルを生成する。
非自由incidenceは、この列ベクトルの四成分が全て零となるスキームである。
$X$ の法方向から $df,dg$ の2方向が得られ、残る2方向への変形の微分は
\[
 H^0(\mathcal B_L)\longrightarrow\mathcal B_L|_p
\]
という評価写像である。これは全射なので、incidenceの四方程式の微分は独立である。
従って $\mathcal I_F$ は $C_0$ に沿って滑らかである。
proper性を用いて近傍を縮小すれば、全体を滑らかな既約9次元空間としてよい。

評価核のFittingイデアルは曲線のイデアルそのものなので、スキームとして
\[
 \mathcal I_F\times_{U_F}P=\mathscr C_P
\]
が成立する。右辺は滑らかな8次元空間であるから、滑らかな $\mathcal I_F$ の
Cartier因子である。従ってblowupの普遍性によりproperな射
\[
 f:\mathcal I_F\longrightarrow\operatorname{Bl}_P U_F
\]
を得る。targetの例外因子を $E_P$ とすると
\[
 f^*E_P=\mathscr C_P
\]
がCartier因子として係数1で成立する。
例外因子上の射は\eqref{eq:local-bn-embedding}の相対版であり、閉埋め込みである。
その接写像の階数は8、法方向も係数1の等式により非零なので、$df$ の階数は9である。

このように $f$ は例外因子上でunramifiedで、その各像点のファイバーは被約な一点である。
proper性によりその近傍でfiniteにでき、finite加群に対するNakayamaにより
局所環の射は全射になる。従って $f$ はそこで閉埋め込みである。
さらにproper性で $[F]$ の近傍を縮小し、$f$ を境界のstrict transformへの同型にできる。
このstrict transformは滑らかで、$E_P$ と横断的に交わる。
交叉は $\PP_P(N_{P/U_F})$ 内の相対三次曲線なので、
$D_F$ の $P$ に沿う初期式は、重複のない相対三次式となる。
従って位数は3であり、normal coneの記述を得る。
strict transformは $\operatorname{Bl}_P D_F$ に一致する。
$P$ の外ではblowupが同型なので $D_F\setminus P$ は滑らかである。
一方、$P$ に沿う方程式の位数は3なので $P\subset\operatorname{Sing}(D_F)$ である。
これで最後の等式も従う。
\end{proof}

\subsection{境界の形式近傍は相対normal coneに一致する}

\begin{theorem}[三次項より高い項の除去]\label{thm:relative-cubic-formality}
$[F]$ の完備局所座標を $(s_1,\ldots,s_7,z_1,z_2,z_3)$、
$P=(z_1,z_2,z_3)$ と取る。$D_F$ の方程式を
\[
 d(s,z)=f_3(s,z)+O((z)^4)
\]
と書けば、$f_3$ は $z$ に関する斉次三次式で、$s=0$ の平面三次曲線は
$(C_0,\mathcal M_L)$ である。
$s$ を固定し $z$ の一次項を保つ形式座標変換により、$d$ は正確に $f_3$ へ移る。
従って $D_F$ の $P$ に沿う形式近傍は、相対normal coneの形式近傍に同型である。
\end{theorem}

\begin{proof}
$A=\CC[[s_1,\ldots,s_7]]$ とおく。
$f_3(0,z)$ は滑らかな平面三次曲線なので、その三つの偏微分は正則列をなし、
Jacobian環のHilbert seriesは $(1+t)^3$ である。
従って任意の $n\ge4$ に対し
\[
 A[z]_{n-2}^{\oplus3}\longrightarrow A[z]_n,
 \quad(h_1,h_2,h_3)\longmapsto
 \sum_i h_i\frac{\partial f_3}{\partial z_i}
\]
は、特殊ファイバーで全射であり、Nakayamaにより $A$ 上でも全射である。
最低次の剰余項 $d_n$ をこの像として書き、
$z_i\mapsto z_i-h_i$ で消す。$h_i$ の次数は $n-2\ge2$ なので、
それ以前の次数は変わらない。この操作を繰り返せば $(z)$-進位相で収束し、
所要の座標変換を得る。
\end{proof}

固定した横断sliceの解消では、例外曲線は $C_0$、その法束は
$\mathcal M_L^{-1}$ であり、自己交叉数は $-3$ である。
ここで得たのは相対三次錐であって、一定の楕円曲線の錐と
7次元滑らかな空間との積であるとは限らない。
平面を動かすと $C_0$ の $j$-不変量も動き得る。
また本節の高次項除去は、$[F_L]$ という評価核の点での三次横断方程式に対するものであり、
$e=[E]$ における六次方程式の高次項を消したものではない。

\section{六次境界からのTorelliと最小の局所次数}\label{sec:boundary-torelli}

第\ref{sec:sextic-boundary}節の六次超曲面は、もとの三次三次元多様体を
射影同型まで決定する。さらに、非局所自由境界の形式芽全体を用いなくても、
その6次までの局所環で復元に十分である。

$X=V(F)\subset\PP(V^*)$、$W=\Lambda^2V$ とし、
\[
 P_X(a)=F(q(a)),\qquad
 R_X=\operatorname{Sym}(W^*)/(P_X),\qquad
 \Gamma_X=\operatorname{Spec}R_X
\]
と書く。$q=(q_0,\ldots,q_4)$ は符号付きPfaffianである。
$H_X=\operatorname{Proj}R_X$ は正規な8次元六次超曲面で、
$(\operatorname{Sing}H_X)_{\mathrm{red}}=\Gr(2,V)$ である。

\subsection{偏極の内在性と射影Torelli}

\begin{lemma}[六次超曲面のPicard群]\label{lem:sextic-picard}
\[
 \operatorname{Pic}(H_X)=\mathbb Z[\OO_{H_X}(1)],\qquad
 \omega_{H_X}=\OO_{H_X}(-4).
\]
従って $\OO_{H_X}(1)$ は $\omega_{H_X}^{-1}$ の唯一の4乗根であり、
抽象多様体 $H_X$ から内在的に決まる。
\end{lemma}

\begin{proof}
解消を $r:J=\PP_X(\Lambda^2Q)\to H_X$、束射影を $\pi:J\to X$ と書く。
正規性とproper双有理性から $r_*\OO_J=\OO_{H_X}$ なので、
射影公式により $r^*:\operatorname{Pic}(H_X)\to\operatorname{Pic}(J)$ は単射である。
Lefschetzの定理\cite[Expos\'e XII, Corollaire 3.7]{SGA2}と射影束の公式から
\[
 \operatorname{Pic}(J)
 =\mathbb Z[\pi^*\OO_X(1)]\oplus\mathbb Z[\OO_J(1)],\qquad
 \OO_J(1)=r^*\OO_{H_X}(1).
\]
$r^*L=\pi^*\OO_X(b)\otimes\OO_J(a)$ とする。
階数2の一般の点のfiberは滑らかな平面三次曲線 $T$ であり、
$r^*L|_T$ と $\OO_J(1)|_T$ は自明である。
一方 $\deg\OO_X(1)|_T=3$ なので $3b=0$、従って $b=0$ となる。
引き戻しの単射性から $L=\OO_{H_X}(a)$ を得る。
超曲面のadjunctionは $\omega_{H_X}=\OO_{H_X}(6-10)$ を与える。
\end{proof}

\begin{theorem}[六次超曲面のTorelli]\label{thm:sextic-torelli}
滑らかな三次三次元多様体 $X,Y$ に対し、任意の抽象同型
$H_X\xrightarrow{\sim}H_Y$ は、一意の射影同型 $X\xrightarrow{\sim}Y$ から誘導される。
逆に任意の射影同型 $X\xrightarrow{\sim}Y$ は $H_X\xrightarrow{\sim}H_Y$ を誘導する。
特に、抽象代数多様体の自己同型群として
\[
                       \operatorname{Aut}(H_X)\simeq\operatorname{Aut}(X)
\]
である。
\end{theorem}

\begin{proof}
補題\ref{lem:sextic-picard}により同型は $\OO(1)$ を保つ。
制限 $H^0(\PP(W),\OO(1))\to H^0(H_X,\OO_{H_X}(1))$ は同型なので、
これは一意の射影線形同型 $\PP(W_X)\to\PP(W_Y)$ に延びる。
この同型は被約特異部分 $\Gr(2,V_X)$ と $\Gr(2,V_Y)$ を対応させる。
Chowの定理\cite[Section II.1]{ChowGrass}\cite[Theorem 3.6]{GhorpadeKaipa}により、
$5\ne2\cdot2$ なので、この射は $\PP(\Lambda^2g)$ の形を持つ。
ここで射影同型 $[g]:\PP(V_X)\to\PP(V_Y)$ は一意であり、
Grassmannianの双対性に由来する追加の同型は生じない。

階数4の部分上のkernel写像を
\[
 \kappa_X:\PP(\Lambda^2V_X)\setminus\Gr(2,V_X)\longrightarrow\PP(V_X^*)
\]
と書く。この射は全射であり、
\[
 H_X\setminus\Gr(2,V_X)=\kappa_X^{-1}(X),\qquad
 \kappa_Y(\Lambda^2g(a))=(g^{-1})^*\kappa_X(a).
\]
従って $\PP((g^{-1})^*)$ は $X$ を $Y$ へ移す。
逆向きは同じ公式から従う。
また $\operatorname{Pic}(X)=\mathbb Z[\OO_X(1)]$ なので、
$X$ の抽象自己同型も射影自己同型である。
\end{proof}

この証明は、次の座標に依存しない復元手順も与える。
$H=H_X$ をその内在的な $\OO_H(1)$ で $\PP^9$ に埋め込み、
\[
 Z=(\operatorname{Sing}H)_{\mathrm{red}},\qquad
 K=H^0(\PP^9,\mathcal I_Z(2))
\]
とおく。$K$ は5次元で、その二次式の全線形系による有理写像
\[
 \Phi_K:\PP^9\dashrightarrow\PP(K^*)
\]
は、Pfaffianによるkernel写像そのものである。従って
\begin{equation}\label{eq:boundary-reconstruct-cubic}
                 X\simeq\overline{\Phi_K(H\setminus Z)}
                 \subset\PP(K^*)
\end{equation}
と復元できる。
多項式の言葉では、代入写像
$\operatorname{Sym}^3K\to\operatorname{Sym}^6(W^*)$ は単射で、
六次方程式 $P_X$ の逆像が、スカラー倍を除いて一意な三次式 $F$ である。

\begin{theorem}[次数付き環と抽象アフィン錐]\label{cor:affine-sextic-torelli}
次の条件は同値である。
\begin{enumerate}
\item $X$ と $Y$ が射影同型である。
\item $H_X$ と $H_Y$ が抽象多様体として同型である。
\item $R_X$ と $R_Y$ が次数付き $\CC$-代数として同型である。
\item $\Gamma_X$ と $\Gamma_Y$ が、頂点の指定を忘れたアフィン多様体として同型である。
\end{enumerate}
また
\[
 1\longrightarrow\CC^*\longrightarrow\operatorname{Aut}_{\mathrm{gr}}(R_X)
 \longrightarrow\operatorname{Aut}(X)\longrightarrow1
\]
は完全列である。
\end{theorem}

\begin{proof}
前三条件の同値と完全列は定理\ref{thm:sextic-torelli}による。
アフィン錐については
\[
 (\operatorname{Sing}\Gamma_X)_{\mathrm{red}}
   =\operatorname{Cone}(\Gr(2,V)),\qquad
 \bigl(\operatorname{Sing}((\operatorname{Sing}\Gamma_X)_{\mathrm{red}})\bigr)_{\mathrm{red}}
   =\{0\}.
\]
実際、Pl\"ucker錐は頂点以外で滑らかで、頂点では次元7に対し接空間が10次元である。
従ってアフィン同型は頂点を保つ。
頂点の局所環の随伴次数環は、次数1で生成された $R_X$ 自身なので、
その同型は $R_X\simeq R_Y$ を誘導する。
\end{proof}

この完全列の分裂は主張しない。
また、アフィン錐や形式芽の全自己同型群では、接空間上恒等となる非線形自己同型も
考慮する必要があるので、それらを $\operatorname{Aut}(X)$ に同一視してはいない。

\subsection{形式境界と6次の局所環}

非局所自由境界の完備局所環を
\[
 A_X=\widehat{\OO}_{D_X,e_X},\qquad
 \mathfrak m_X\subset A_X
\]
とし、$r$ 次の局所環を $A_X^{(r)}=A_X/\mathfrak m_X^{r+1}$ と定める。

\begin{theorem}[境界の6次Torelli]\label{thm:boundary-sixth-jet}
抽象的な完備局所 $\CC$-代数 $A_X$ は、$X$ を射影同型まで決定する。
さらに
\[
 X\simeq Y\ \text{（射影同型）}
 \quad\Longleftrightarrow\quad
 A_X^{(6)}\simeq A_Y^{(6)}\ \text{（局所 $\CC$-代数同型）}.
\]
一方、全ての $r\le5$ について
\[
 A_X^{(r)}\simeq
 \CC[a_1,\ldots,a_{10}]/(a_1,\ldots,a_{10})^{r+1}
\]
であり、$X$ に依存しない。
従って最大イデアルによるこの切断の中では、6次が復元に必要な最小の次数である。
\end{theorem}

\begin{proof}
定理\ref{thm:sextic-boundary}から
$\operatorname{gr}_{\mathfrak m_X}A_X=R_X$ である。
より具体的に $T^*=\mathfrak m_X/\mathfrak m_X^2$ とおくと、
\[
 \ker\left(\operatorname{Sym}^6T^*
      \longrightarrow\mathfrak m_X^6/\mathfrak m_X^7\right)
 =\CC\cdot P_X
\]
は1次元である。この積写像は $A_X^{(6)}$ だけから内在的に決まるので、
六次超曲面 $H_X$、従って定理\ref{thm:sextic-torelli}により $X$ を復元できる。

逆に、正則な局所座標で境界の方程式を
$d_X=P_X+\text{（7次以上）}$ と書けば
\[
 A_X^{(6)}\simeq
 \CC[a_1,\ldots,a_{10}]/\bigl((a)^7,P_X\bigr).
\]
従って射影同型な $X,Y$ は同型な6次局所環を与える。
5次以下では関係式 $d_X$ 自体が消えるため、最後の主張も従う。
\end{proof}

\begin{remark}[このTorelliの入力]
ここでの入力は、gluing object $E$ における\emph{非局所自由境界の指定された形式芽}
である。その局所環を与えれば、周囲の滑らかな $U$、接座標、$X$ によって
添字付けられたincidenceの族は、復元の入力に残す必要がない。
一方、境界の構成には $\Ku(X)\subset D^b(X)$ という幾何的実現と
局所自由性、または点対象を用いたBrill--Noether条件を用いている。
従って、この結果だけから、境界を指定しない残余圏や滑らかな形式空間
$\Def_E\simeq\operatorname{Spf}\CC[[a_1,\ldots,a_{10}]]$ が
同じ復元データを内在的に選ぶとは結論しない。
\end{remark}

\subsection{境界そのものの形式斉次性を判定する条件}

6次Torelliは、境界がその接錐と形式同型であるという主張とは区別する必要がある。
後者には、universal対象のcohomology-jump方程式の高次項が関係する。
その判定には次の内在的な条件が使える。

\Needspace{8\baselineskip}
\begin{theorem}[形式Euler導分による判定]\label{prop:boundary-euler}
$A=\CC[[a_1,\ldots,a_{10}]]/(d)$ とし、$d$ の最低次項を
$P=P_X$ とする。次の条件は同値である。
\begin{enumerate}
\item $A\simeq\CC[[a_1,\ldots,a_{10}]]/(P)$ である。
\item $A$ に、最大イデアルを保ち、$\mathfrak m_A/\mathfrak m_A^2$ 上で恒等作用をする
連続 $\CC$-導分 $\delta$ が存在する。
\end{enumerate}
\end{theorem}

\begin{proof}
斉次な環にはEuler導分があるので $(1)\Rightarrow(2)$ は明らかである。
逆に $\delta(a_i)$ の代表元を選び、$\delta$ を形式冪級数環の導分へ持ち上げる。
その線形部分は $\mathcal E=\sum_i a_i\partial_{a_i}$ であり、
$\delta(d)=b d$、$b(0)=6$ を満たす。
次数 $n\ge2$ の座標補正に対する作用素が $n-1$ 倍であることを用い、
逐次 $\delta(y_i)=y_i$ となる形式座標を構成できる。
この座標で $\mathcal E(d')=b'd'$、$b'(0)=6$ である。
\[
                      \mathcal E(\log u)=6-b'
\]
を各次数 $n$ で $n$ により割って解けば、$u(0)=1$ のunit $u$ が得られ、
$\mathcal E(ud')=6ud'$ となる。従って $ud'$ は純粋な6次式である。
座標変更の線形部分と $u$ の定数項はいずれも恒等なので、その6次式は $P$ である。
\end{proof}

具体的には、$d=P+d_7+d_8+\cdots$ とし
$S=\CC[a_1,\ldots,a_{10}]$、$J_P=(\partial_{a_i}P)$ とおく。
最初の障害は
\[
                             [d_7]\in(S/J_P)_7
\]
である。二次の座標補正と一次のunit補正で変わる項は
$\sum_iS_2\partial_{a_i}P+S_1P=(J_P)_7$ だからである。
等号にはEulerの公式 $6P=\sum_i a_i\partial_{a_i}P$ を用いた。
7次から $N-1$ 次までを消した後には、同様に $(S/J_P)_N$ が次の障害を受け持つ。
高次の障害は、それ以前に行った消去の選択に依存し得る。
全次数で整合する消去を構成できることが、形式斉次化の条件である。
第\ref{sec:boundary-higher-jets}節では、全非零三次式上の正則で同変な座標を構成し、
実際の境界について $d_7=d_8=d_9=0$ とできることを示す。
第7版の段階で次に残った項は10次であった。
第\ref{sec:tate-boundary-tenth}節では、その実際の係数と消去不能性を決定する。

なお、この接錐には一般の有限決定性を適用できない。
任意の $M>6$ に対し
\[
                         P+\sum_{i=1}^{10}a_i^M
\]
は $P$ と $M-1$ 次まで一致するが、原点は孤立特異点である。
実際、次数と両立する単項式順序で、その10個の偏微分の先頭単項式は
$a_i^{M-1}$ なので、Jacobian環は有限次元になる。
一方、$P=0$ の特異部分は7次元のPl\"ucker錐である。
従って両者は形式同型ではない。
これは任意の超曲面芽の中での例であり、この摂動が実際の境界 $D_X$ から
生じるとは主張していない。

\clearpage
\section{境界の全特異部分と二段階の解消}
\label{sec:global-bn-singular}

本節では、$e=[E]$ を含む全形式近傍で、境界の被約な特異部分を同定する。
前節までの相対Picardの記述を使い、特異または非被約な平面三次曲線の方向も含める。
各曲線が滑らかであることは仮定せず、全平面を動かす普遍曲線の全空間が
滑らかであることを用いる。

$A=\widehat{\OO}_{\Def_E,e}\simeq\CC[[w_1,\ldots,w_{10}]]$、
$U=\operatorname{Spec}A$ とし、境界、適合部分空間、incidenceをそれぞれ
$D,C,\mathcal I$ と書く。
以下の主張はこの完備局所環上でのものである。
完備局所環上では正則性を用い、形式近傍の滑らかさとして記述する。
$G=\Gr(2,V)$、$W=\Lambda^2V\otimes T$、
$T=H^3(X,\OO_X(-2))$ とおく。

\begin{theorem}[境界の全特異部分]\label{thm:global-bn-singular}
$e$ における形式近傍で、閉被約部分スキームとして
\[
 \operatorname{Sing}(D)_{\mathrm{red}}=C
\]
が成立する。特に、$C$ に含まれない小次元の特異成分は存在しない。
また $D$ は正規なGorenstein因子である。
\end{theorem}

この命題は、$\operatorname{Sing}(D)$ の非被約なスキーム構造まで
$C$ と一致すると主張するものではない。
証明のため、まず射影された評価核の局所表示を全Picard族上で求める。

\subsection{射影対象の局所微分と平面曲線}

形式Picardの局所chartで零切断の方程式を $t=0$ とし、
曲線族を $j:\mathscr C\hookrightarrow X\times S$、線形束を $L$ と書く。
その評価核を
\[
 0\longrightarrow F\longrightarrow p_*L\otimes\OO_{X\times S}
 \longrightarrow j_*L\longrightarrow0
\]
で定める。$j$ は相対余次元2のregular immersionなので、
局所的には $F\simeq\mathcal I_{\mathscr C}\oplus\OO^2$ である。

\begin{lemma}[Picard族の局所行列]\label{lem:picard-local-matrix}
$F$ の局所free resolutionを
$[P^{-1}\xrightarrow{d_F}P^0]$、階数をそれぞれ1と4と取れる。
適当な局所自明化のもとで、射影対象 $\mathcal E=\pi F$ は
\begin{equation}\label{eq:picard-local-matrix}
 \mathcal E\simeq
 [P^{-1}\xrightarrow{\,t b^{-1}d_F\,}Q],
 \qquad b:Q\xrightarrow{\sim}P^0
\end{equation}
と表示される。従ってincidenceの四つの方程式のPicard族への制限は、
イデアルとして正確に $t\mathcal I_{\mathscr C}$ を生成する。
\end{lemma}

\begin{proof}
相対theta補題\ref{lem:relative-picard-theta}により、
単元を吸収すると
\[
 K=\RHom(Q,F)\simeq Rp_*L(-1)
   \simeq[\OO_S\xrightarrow{t}\OO_S]
\]
であり、右辺は次数0と1に置かれる。
$F\in\OO_X^\perp$ 上で $\pi=L_Q$ なので、mutation triangleは
\[
 K\otimes Q\xrightarrow{\mathrm{ev}}F\longrightarrow\mathcal E
\]
となる。零切断ではEuler列とKoszul列から
\[
 0\longrightarrow\OO_X(-1)\otimes\det U
 \longrightarrow Q\longrightarrow F_0\longrightarrow0
\]
を得る。evaluationの次数0成分はこの自然な全射 $Q\to F_0$ である。
局所resolutionをこの全射に合わせて選ぶと、
その持ち上げ $b:Q\to P^0$ は零切断上で可逆であり、
形式近傍でも可逆となる。ここでは基底の線形束を局所自明化している。

evaluationのconeの微分は
\[
 \begin{pmatrix}
  d_F&b\\
  0&-tI_4
 \end{pmatrix}
\]
である。可逆なblock $b$ を消去すると
\eqref{eq:picard-local-matrix}を得る。
$d_F$ の成分は $\mathcal I_{\mathscr C}$ を生成するため、イデアルの主張も従う。
この計算は相対perfect complexのconeに対するものであり、
各曲線が滑らかであることを必要としない。
\end{proof}

\subsection{二つのblowupと閉埋め込み}

\begin{theorem}[境界の埋め込み解消]\label{thm:bn-double-blowup}
次の二段階のblowupを考える：
\[
 Y_1=\operatorname{Bl}_eU,\qquad
 C_1=\operatorname{Bl}_eC\subset Y_1,\qquad
 Y_2=\operatorname{Bl}_{C_1}Y_1.
\]
$B=\operatorname{Bl}_X\mathcal I$ は $Y_2$ に閉埋め込み、
その像は $D$ のstrict transformに一致する。
$Y_2\to U$ は $D$ のembedded log resolutionである。
\end{theorem}

\begin{proof}
最初の例外因子を $E_1=\PP(W)$ と書く。
$\mathfrak m_e\OO_{\mathcal I}=\mathcal I_X$ なので、
blowupの普遍性からproperな射 $B\to Y_1$ があり、その閉ファイバーは
\[
 J=\PP_X(\Lambda^2Q).
\]
$B$ は滑らかな9次元空間であり、$E_1$ のpullbackは係数1で $J$ となる。

局所表示 $\mathcal E_U=[\OO\xrightarrow{d}\OO^4]$ を取る。
blowup chartで $w_j=\rho$、$w_i=\rho a_i$ とすると
$d_i(x,\rho a)=\rho r_i(x,a,\rho)$ と書ける。
$B$ の方程式は四つの $r_i=0$ である。
実際、$\rho=0$ の零点は滑らかな $J=\PP_X(\Lambda^2Q)$ で、
$r_i$ の四つの微分は独立である。
従って $V(r_1,\ldots,r_4)$ 上で $\rho$ は非零因子であり、
$\rho$-torsionを除く操作による余分な変化はない。
strict transformが $\operatorname{Bl}_X\mathcal I$ に一致する一般的性質は
\cite{StacksStrictTransform}も参照されたい。

相対Picardの記述により、$C_1$ の $G$ に沿う完備化は形式Picardそのものである。
従って $C_1$ は滑らかな7次元空間で、
$C_1\cap E_1=G$ は横断的な交叉である。
$\rho|_{C_1}$ はPicard parameter $t$ の単元倍であるから、
補題\ref{lem:picard-local-matrix}を四つの $r_i$ に適用すると、
スキームとして
\begin{equation}\label{eq:bn-blowup-picard-intersection}
 B\times_{Y_1}C_1=\mathscr C_{\widehat{\mathcal P}}
\end{equation}
を得る。これは非flatなbase changeとstrict transformの交換を仮定したものではなく、
明示した方程式を制限した結果である。

ここで普遍平面三次曲線は
\[
 \mathscr C_G=\{(p,U):p\in X,\ U\subset\ker(p:V\to\OO_X(1))\}
            =\Gr_X(2,Q)
\]
という $X$ 上のGrassmannian束であり、
その全空間は滑らかな7次元空間である。
形式Picardは $G$ 上の形式線形束なので、
\eqref{eq:bn-blowup-picard-intersection}の右辺は滑らかな8次元空間となる。
従って $Z=B\times_{Y_1}C_1$ は $B$ のCartier因子であり、
\[
 \mathcal I_{C_1}\OO_B=\mathcal I_Z.
\]
形式完備化で確認したCartier性は $A$ 上でも成立する。
その不成立箇所が非空なら、proper性により閉点ファイバーに会うからである。

よって $B\to Y_2$ へ持ち上がり、二番目の例外因子を $E_2$ とすると、
そのpullbackは係数1で $Z$ である。
$C_1$ は $E_1$ と横断的なので、閉ファイバーは正確に
\[
 (Y_2)_e=\operatorname{Bl}_G\PP(W).
\]
$G$ のidealは五つのPfaffianで生成されるから、
このblowupはPfaffianによる有理写像のgraph closureである。
rank4の稠密開集合でkernel pointを取ることにより、
\begin{equation}\label{eq:pfaffian-projective-graph}
 \operatorname{Bl}_G\PP(W)
 \simeq\PP_{\PP(V^*)}(\Lambda^2Q_{\PP(V^*)})
\end{equation}
を得る。右辺は滑らかな $\PP^5$-束であり、rank4のgraphの閉包に一致する。
$J=\PP_X(\Lambda^2Q)$ は右辺の $X\subset\PP(V^*)$ への逆像である。
従って $B\to Y_2$ の閉ファイバーは閉埋め込みである。

最後に、この閉埋め込みを全体へ上げる。
$A$ 上properな射 $f:B\to Y_2$ の非quasi-finite locusは閉集合である。
閉ファイバーが閉埋め込みなので、この閉集合は閉点ファイバーに会わない。
proper性により空であり、$f$ はproperかつquasi-finite、従ってfiniteである
\cite{StacksProperFinite}。
finite base changeとNakayamaにより
$\operatorname{coker}(\OO_{Y_2}\to f_*\OO_B)$ の台も閉点ファイバーに会わない。
再びproper性からこのcokernelは零で、$f$ は閉埋め込みとなる。
その像は、中心の外の境界の閉包、すなわち $D$ のstrict transformである。

最初の例外因子のstrict transformを $E_1'$ と書く。
$Y_2$ の次元は10であり、交叉はすべてスキームとして
\[
 \begin{aligned}
 B\cap E_1'&=J &&(\dim=8),\\
 B\cap E_2&=Z &&(\dim=8),\\
 B\cap E_1'\cap E_2&=\mathscr C_G &&(\dim=7),\\
 E_1'\cap E_2&=\PP_G(N_{G/\PP(W)}) &&(\dim=8)
 \end{aligned}
\]
と表される。各全空間は滑らかで、次元は期待次元に等しい。
従って $B+E_1'+E_2$ はsimple normal crossingsであり、
所要のembedded log resolutionを得る。
\end{proof}

\begin{proof}[命題\ref{thm:global-bn-singular}の証明]
二つのblowupの合成は $U\setminus C$ 上で同型である。
境界のstrict transformが滑らかな $B$ なので、
$D\setminus C$ は滑らかであり、
$\operatorname{Sing}(D)_{\mathrm{red}}\subset C$ を得る。
逆包含は、滑らかな平面切断の評価核を表す $C$ の稠密部分で既に証明した。
特異部分は閉集合で、$C$ は既約であるから、逆包含は $C$ 全体に延長される。

$D$ は整な超曲面なので $S_2$ を満たす。
その特異部分は $D$ 内で余次元2なので $R_1$ も満たし、正規である。
Gorenstein性は超曲面であることから従う。
\end{proof}

\subsection{対数標準特異点と二つの数値不変量}

discrepancyを $a$、log discrepancyを $A=a+1$ とする。
境界の全高次方程式を決定しなくても、上の解消からこれらを計算できる。
対数標準性とlog resolutionによるthresholdの計算には、
例えば\cite[Section 1]{MustataSing}の定義を用いる。

\begin{theorem}[対数標準性と最小log discrepancy]\label{thm:bn-log-invariants}
$D$ は対数標準特異点を持ち、そのnon-klt locusは $C$ である。
さらに
\[
 \operatorname{lct}_e(U,D)=1,\qquad
 \operatorname{mld}_e(D)=4
\]
となる。
\end{theorem}

\begin{proof}
$\pi:Y_2\to U$ とおく。
$e$ における境界の重複度は6であり、
$C_1$ の一般点におけるstrict transformの重複度は3である。
後者は滑らかな平面切断の評価核点での計算による。
従って
\begin{equation}\label{eq:bn-log-resolution-divisors}
 \pi^*D=B+6E_1'+3E_2,\qquad K_{Y_2/U}=9E_1'+2E_2.
\end{equation}
最初の中心は10次元空間の点、次の中心は余次元3である。
また $C_1$ は $E_1$ に含まれないため、
$E_1$ のpullbackに $E_2$ の追加成分は生じない。

adjunctionを $B$ に制限すると
\begin{equation}\label{eq:bn-adjunction-discrepancy}
 K_B=\pi_D^*K_D+3J-Z.
\end{equation}
例外因子のdiscrepancyは $3$ と $-1$ であり、$J+Z$ はsimple normal crossingsである。
従って $D$ は対数標準で、non-klt locusは $Z$ の像 $C$ となる。
また\eqref{eq:bn-log-resolution-divisors}から
\[
 \operatorname{lct}_e(U,D)
 =\min\left\{1,\frac{10}{6},\frac{3}{3}\right\}=1.
\]

$e$ を中心とする任意のdivisorial valuation $\nu$ の $B$ 上の中心は $J$ に入る。
$(B,J+Z)$ が対数標準であるため、
$A_B(\nu)\ge\nu(J)+\nu(Z)$ である。
よって\eqref{eq:bn-adjunction-discrepancy}から
\[
 A_D(\nu)=A_B(\nu)-\nu(Z)+3\nu(J)\ge4\nu(J)\ge4.
\]
$J$ 自身が $A_D(J)=4$ を達成するので、
$\operatorname{mld}_e(D)=4$ となる。
\end{proof}

\subsection{適合部分空間を一度に吹き上げる記述}

二段階の解消から、incidenceそのものにも埋め込みの記述が得られる。
以下では $C$ の標準Pl\"ucker座標を $A$ に持ち上げ、
$C$ を五つの斉次二次Pfaffianで定める形式座標を選ぶ。
これは $A\to\OO_C$ が余接空間の同型を誘導するため可能である。
traceの1次元因子は自明化して表示する。

\begin{theorem}[適合部分空間のblowup内のincidence]\label{thm:bn-single-blowup}
$H=\operatorname{Bl}_C U$ とすると、その閉ファイバーに沿う完備化は
\[
 \widehat H_{\PP(V^*)}
 \simeq
 \widehat{\operatorname{Tot}_{\PP(V^*)}
        (\Lambda^2Q_{\PP(V^*)}\otimes T)}_{\PP(V^*)}
\]
である。従って $H$ は滑らかで、
incidenceは $D$ のstrict transformとして $\mathcal I\hookrightarrow H$
と閉埋め込まれる。
射影を $p:H\to\PP(V^*)$ と書けば、そのCartier classは
\[
 \OO_H(\mathcal I)\simeq p^*\OO_{\PP(V^*)}(3)
\]
である。
さらに $\mathcal I\times_U C$ は被約なCartier因子 $S$ であり、
$S$ は $X=\mathcal I_e$ に沿って横断的なordinary double pointを持つ。
\end{theorem}

\begin{proof}
$S$ を $\mathcal I|_{C\setminus\{e\}}$ の閉包の素因子とする。
$\mathcal I$ は滑らかなので $S$ はCartierであり、
$\beta:B\to\mathcal I$ におけるstrict transformは $Z$ である。
$\beta^*S=Z+mJ$ と書く。
$J\to X$ の各 $\PP^5$-ファイバー上で $Z\cap J$ は
\[
 \Gr(2,Q_x)\subset\PP(\Lambda^2Q_x)
\]
というPl\"ucker二次超曲面であり、
$\OO_B(J)|_{\PP^5}=\OO_{\PP^5}(-1)$ である。
一方、$\beta^*\OO_{\mathcal I}(S)$ はこのファイバー上で自明だから、
$0=2-m$、すなわち $m=2$ となる。

$C$ の方程式は斉次二次式なので、
最初のblowupでは
\[
 \mathcal I_C\OO_{Y_1}
 =\OO_{Y_1}(-2E_1)\mathcal I_{C_1}.
\]
これを $B$ に引き戻すと
$\mathcal I_C\OO_B=\OO_B(-2J-Z)=\mathcal I_S\OO_B$ を得る。
$\mathcal I_C\OO_{\mathcal I}\subset\mathcal I_S$ だから、
invertible ideal $\mathcal I_S$ で割ったidealを $\mathcal K$ とおけば
$\mathcal K\OO_B=\OO_B$ である。
$\beta$ は全射なので $\mathcal K=\OO_{\mathcal I}$ となる。
実際、$\mathcal K$ がある点の極大idealに入るなら、
その逆像の点でもunit idealにはならない。
従って
\begin{equation}\label{eq:bn-preimage-compatibility}
 \mathcal I_C\OO_{\mathcal I}=\mathcal I_S
\end{equation}
がスキームとして成立し、$\mathcal I\to H$ が存在する。

$H$ はアフィンPfaffian写像のgraph closureである。
そのgraphは
\[
 \{(a,p):a\in\Lambda^2\ker p\otimes T\}
 =\operatorname{Tot}_{\PP(V^*)}(\Lambda^2Q\otimes T)
\]
であり、rank4の稠密開集合を含む滑らかなrank6のベクトル束となる。
これを原点でbase changeし、完備化すると命題の表示を得る。

また $\mathcal I_C\OO_{Y_2}=\OO_{Y_2}(-2E_1'-E_2)$ なので
$Y_2\to H$ も存在する。
閉ファイバー上では\eqref{eq:pfaffian-projective-graph}の $\PP(V^*)$ への射影である。
従って $J\to H_e$ は $J\to X\hookrightarrow\PP(V^*)$ に一致し、
$\mathcal I_e\to H_e$ は元の閉埋め込み $X\hookrightarrow\PP(V^*)$ である。
命題\ref{thm:bn-double-blowup}の証明で用いたproper性とNakayamaの議論により、
$\mathcal I\to H$ は全体で閉埋め込みとなる。

$H$ の例外因子を $R$ とする。
rank4束 $Q$ のskew form $a$ に対する $a\wedge a$ は
$\det Q\otimes T^2$ の切断なので
\[
 \OO_H(R)=p^*(\det Q\otimes T^2)\simeq p^*\OO(-1).
\]
最後の同型では定数の1次元因子を自明化した。
$C$ の一般点での境界の位数は3だから、$H\to U$ による総変換は
$\mathcal I+3R$ である。
$D$ は $U=\operatorname{Spec}A$ 上でprincipalなので
\[
 \OO_H(\mathcal I)=\OO_H(-3R)\simeq p^*\OO(3)
\]
を得る。ここで線形束のclassを決定したことは、
その切断の高次項を消去したことを意味しない。

最後に\eqref{eq:bn-preimage-compatibility}から $S=\mathcal I\cap R$ である。
零切断に沿う $R$ の二次項は、六つのskew coordinatesの非退化Pfaffianである。
$\mathcal I\cap\PP(V^*)=X$ は滑らかな三次超曲面なので、
$\mathcal I$ の $X$ に沿う法方向でもこの二次形式は非退化である。
形式Morse補題により、各 $x\in X$ の局所で $(S,X)$ は
\[
 \left(X\times
 \{z_1z_2+z_3z_4+z_5z_6=0\},\ X\times\{0\}\right)
\]
に形式同型となる。
これは $X$ 上局所的な主張であり、法束の大域的な自明性は仮定しない。
\end{proof}

$H$ の例外因子 $R$ 自体は零切断に沿って特異なので、
simple normal crossingsの計算には二段階のモデルを用いる。
実際、graphの束表示から $Y_2\simeq\operatorname{Bl}_{\PP(V^*)}H$ と分かり、
その中で $\mathcal I$ のstrict transformは $B=\operatorname{Bl}_X\mathcal I$ となる。

以上の結果は、特異平面切断を含む方向で別の特異成分が発生する可能性を除く。
一方、Jacobian idealが定める非被約な特異スキームの詳細や、
六次初期式 $F(q)$ より高い項を全て消去できるかは、別の問題として残る。


\subsection{境界からの偏極三次三次元多様体の復元}

\begin{theorem}[内在的なblowupによる復元]\label{thm:bn-intrinsic-reconstruction}
実際のBrill--Noether境界の、閉点 $e$ を指定した完備局所芽 $(D,e)$ に対して
$\widetilde D=\operatorname{Bl}_{\operatorname{Sing}(D)_{\mathrm{red}}}D$
とおく。偏極三次三次元多様体は
\[
 (X,\OO_X(1))
 \simeq\bigl(\widetilde D_e,\OO_{\widetilde D}(1)|_{\widetilde D_e}\bigr)
\]
として復元される。
従って二つのこの形式境界の芽が同型なら、
対応する三次三次元多様体は射影的に同型である。
\end{theorem}

\begin{proof}
strict transformの一般的性質\cite{StacksStrictTransform}から、
$H=\operatorname{Bl}_C U$ 内の $D$ のstrict transformは
$\operatorname{Bl}_C D$ に等しい。従って
\[
 \mathcal I\simeq\operatorname{Bl}_C D
 =\operatorname{Bl}_{\operatorname{Sing}(D)_{\mathrm{red}}}D.
\]
その閉ファイバーはスキームとして $X$ である。
またblowupのtautological line bundleは
$\mathcal I_C\OO_{\mathcal I}=\OO_{\mathcal I}(-S)$ であり、
$S=R|_{\mathcal I}$ と $\OO_H(R)\simeq p^*\OO(-1)$ から、
$X$ への制限は $\OO_X(1)$ となる。
全構成が形式境界のgermから内在的に定まり、所要の復元を得る。
\end{proof}

この復元は境界をその接錐へ形式同型にすることを必要としない。
入力は実際のBN境界の芽であり、境界を指定しない残余圏だけからの復元は主張しない。


\section{境界方程式の普遍共変量と高次項}
\label{sec:boundary-higher-jets}

前節では境界の特異部分と解消を決定した。本節では、頂点の方程式自体を調べる。
その係数を三次式 $F$ の多項式として扱える普遍座標を構成し、
7次から9次までを除外し、10次の候補を明示する。
その実際の係数は次節で計算する。

\subsection{特異三次式を含む普遍族上の形式座標と境界方程式}
\label{subsec:universal-bn-equation}

境界方程式の高次項へ群作用を用いるには、その係数が三次式 $F$ に関して
正則であることが必要である。滑らかな三次式の開集合だけで座標を取ると、
判別式の逆数が入り得る。本節では全ての非零三次式へ族を広げる。

\[
 S=\Sym^3V\setminus\{0\},\qquad
 T=(\det V)^*,\qquad W=\Lambda^2V\otimes T,
 \qquad L=(\det V)^{-3}
\]
とおく。trace線を保持すると、Pfaffian写像は
\[
 q:W\longrightarrow\Lambda^4V\otimes T^2
       \simeq V^*\otimes(\det V)^{-1},
 \qquad a\longmapsto\tfrac12a\wedge a
\]
であり、$P(F,a)=F(q(a))$ は $L$ に値を持つ。

\begin{theorem}[正則で同変な普遍境界方程式]
\label{thm:universal-bn-equation}
$S$ 上の対象 $E_F$ の変形空間に、cohomologyのedge同型を一次部分とする
$GL(V)$-同変な形式座標 $a\in W$ を同時に取れる。
この座標で、$L$ に値を持つ正則な同変形式級数
\begin{equation}\label{eq:universal-bn-series}
 d(F,a)=F(q(a))+\sum_{n\ge7}d_n(F,a)
\end{equation}
が存在し、各滑らかな $F$ での零点は実際の形式Brill--Noether境界である。
$d_n$ は $a$ に関して $n$ 次、$F$ に関して $n-5$ 次の多項式であり、
\begin{equation}\label{eq:universal-bn-covariants}
 d_n\in\mathcal C_n:=(\Hom_{GL(V)})\!\left(
 \Sym^{n-5}(\Sym^3V)\otimes\Sym^n(\Lambda^2V),
 (\det V)^{n-3}\right).
\end{equation}
この主張は、整合的に選べる普遍座標についての存在命題である。
任意のminimal modelで既に選んだ座標が、この正則性と同変性を持つとは主張しない。
\end{theorem}

\begin{proof}
\textbf{特異ファイバー上の対象モジュライ。}
$S$ は35次元affine空間から一点を除いたものであり、
\begin{equation}\label{eq:universal-bn-base-vanishing}
 \Gamma(S,\OO_S)=\CC[\Sym^3V],\qquad
 H^i(S,\OO_S)=0\quad(1\le i\le33),\qquad \Pic(S)=0.
\end{equation}
消滅は35個の座標関数の正則列に関する局所コホモロジーから従う。
定数ベクトル束についても同じ消滅が成り立つ。
さらに $GL(V)$ は標数零でlinearly reductiveなので、同変な
$H^1,H^2$ についても消滅する。

普遍超曲面 $\mathcal X\to S$ はflat、projective、Gorensteinである。
また全空間 $\mathcal X$ 自身は滑らかである。
実際、$F$ の係数方向の微分は、各 $p\in\PP(V^*)$ で非零の評価写像になる。
補題\ref{lem:universal-tangent-lift}の計算は、非被約な三次式も含めて
\[
 \mathcal E xt^0(E,E)=\OO_S,\qquad
 \mathcal E xt^1(E,E)=W\otimes\OO_S,
 \qquad \mathcal E xt^{r\ne0,1}(E,E)=0
\]
を与える。

Lieblichの対象モジュライの代数性とsimple locusのrigidification
\cite[Theorem 4.2.1, Corollary 4.3.3]{LieblichComplexes}を用い、
$E_F$ の切断 $\sigma:S\to M$ を含むalgebraic spaceを取る。
この定理の仮定はflat properな射であり、個々のファイバーの滑らかさは要らない。
$\Ext^2(E,E)=0$ と対象の障害理論
\cite[Theorem 3.1.1, Corollary 3.1.2]{LieblichComplexes}により、
$M\to S$ は切断の近傍で相対次元10の滑らかな射である。
ここでは切断が閉埋め込みになる\'etale局所germを用い、
全モジュライの分離性は仮定しない。以下の完備化とblowupもこのgermで行い、貼り合わせる。
その法束は $W\otimes\OO_S$ である。

\textbf{全次数の同変座標。}
滑らかな射と切断の形式座標は基底上局所的に存在する。
一次部分をedge同型に固定し、$n-1$ 次まで座標を貼り合わせたとする。
$n$ 次の二つの延長の差は
\[
 \Hom(W^*,\Sym^nW^*)\otimes\OO_S
       =W\otimes\Sym^nW^*\otimes\OO_S
\]
の切断である。この $H^1$ の消滅を同変圏で用いて、整合する形式同型
\[
                   \widehat M_{\sigma(S)}\simeq S\times\widehat W_0
\]
を次数ごとに構成する。必要ならuniversal対象のscalar gerbeも、
$S$ 上での対象 $E$ から始め、同変な $H^2$ の消滅により同時に自明化できる。
同型類の不変性からlinearizationの整合性を推論したのではなく、
最初から同変な貼り合わせの障害を消している。

\textbf{普遍incidenceとCartier像。}
$M$ の\'etale chartでuniversal複体を取る。
$\mathcal X\times_SM\to\mathcal X$ は滑らかなので、その全空間は正則であり、
このbounded coherentな複体はperfectである。
中心族に沿って可逆な部分を消すと、局所表示は
\[
                       [\,\OO\xrightarrow{\,b\,}\OO^4\,]
\]
となる。微分の四成分の零点はchartの変更と基底からの直線束によるtwistに不変なので、
scalar gerbeから降りてproperなincidence $\mathcal I\to M$ を与える。
その切断上のファイバーはスキームとして $\mathcal X$ である。
基底方向の微分は、全ての非零三次式上で全射な評価写像
\[
 W\otimes\OO_{\mathcal X}\longrightarrow Q(1)\otimes T
\]
であり、核は $N=\Lambda^2Q\otimes T$ である。
従って $\mathcal I$ は中心族に沿って滑らかで、$\mathcal X$ の法束は $N$ である。

proper性により近傍を縮小すると、$\mathcal I$ の中心族を通るgermはintegralである。
そのscheme-theoretic image $\mathcal D\subset M$ もintegralである。
滑らかな三次式の稠密な開集合では定理\ref{thm:sextic-boundary}から余次元1であり、
$M$ は正則なので $\mathcal D$ はCartier因子となる。
以下では、この通常の代数的な像を完備化する。

\textbf{特異三次式上でも先頭係数は消えない。}
$\mathcal I\times_MS=\mathcal X$ がスキームとして成立するので、properな射
\[
 f:\operatorname{Bl}_{\mathcal X}\mathcal I
           \longrightarrow\operatorname{Bl}_{\sigma(S)}M
\]
を得る。両辺の例外因子を $J,A$ とすると、$f^*A=J$ は係数1のCartier因子の等式である。
例外因子上の射は
\[
 J=\PP_{\mathcal X}(N)\longrightarrow S\times\PP(W)
\]
であり、その像は $\mathcal H=\{F(q(a))=0\}$ である。
階数4の部分ではkernel点が一意なので、この射はスキームとして同型になる。
階数2ではkernel平面と非零三次式の零点の交わりは常に非空である。
制限三次式が恒等的に零の場合も、ファイバーが平面全体になるだけである。
従って、この像の台の記述は全ての $F\in S$ で成立する。

$f$ の像は $\mathcal D$ のstrict transformであるから、その例外因子との交わりの台は
正確に $\mathcal H$ となる。
滑らかな三次式上では先頭項が $F(q)$ なので、6次未満のjetは $S$ 全体で零であり、
6次項は局所的に $c(F)F(q)$ である。
$F(q)$ は各非零 $F$ で非零なので、その張る直線は直線部分束であり、$c$ は正則である。
もし $c(F_0)=0$ なら、strict transformと例外因子の交わりは $F_0$ 上の
$\PP(W)$ 全体を含んでしまう。これは $\mathcal H_{F_0}$ が真の超曲面であることに反する。
従って $c$ は全ての点でunitである。

この先頭項は、同変な直線束の同型
\[
                    \mathcal I_{\mathcal D}|_S\simeq L^*\otimes\OO_S
\]
を与える。これを形式近傍へ延長する際の障害は
$H^1([S/GL(V)],\Sym^nW^*\otimes\OO_S)$ に入り、零である。
よって $L$ に値を持ち、先頭項が正確に $F(q)$ である同変な方程式 $d$ を取れる。

\textbf{係数の多項式性と重み。}
各 $d_n$ の係数は $S$ 上の正則関数なので、
\eqref{eq:universal-bn-base-vanishing}から三次式の係数の多項式である。
scalar行列 $z\id_V$ は $F,W,L$ にそれぞれ $z^3,z^{-3},z^{-15}$ で作用する。
従って $F$ の次数 $m$ は $3m-3n=-15$、すなわち $m=n-5$ を満たす。
traceのtwistを外すと\eqref{eq:universal-bn-covariants}を得る。

最後に、scheme-theoretic imageの任意の基底変換との可換性は使わない。
滑らかな $F$ へ制限した $d$ は実際のincidence上で零なので、その境界の方程式で割れる。
両者の先頭項は同じ非零の $F(q)$ であるから、商は形式局所環のunitである。
従って、滑らかな各ファイバーで実際の境界を定義する。
\end{proof}

\begin{lemma}[同じ普遍座標での適合部分空間]
\label{lem:universal-chart-plucker}
定理\ref{thm:universal-bn-equation}の座標では、各滑らかな $F$ の適合形式部分空間は
標準的なPl\"ucker錐そのものである。すなわち
\[
             \widehat C_c=\{a\in\widehat W_0:a\wedge a=0\}
\]
が閉形式部分スキームとして成立する。
\end{lemma}

\begin{proof}
$V=U\oplus Z$、$\dim U=2$ と固定し、
\[
 S_U=\{F:F|_Z\ne0\}
      \simeq\Aff^{25}\times(\Aff^{10}\setminus\{0\}),\qquad
 D_U=\det U\otimes T\simeq(\det Z)^*
\]
とおく。$\Gamma(S_U,\OO)=\CC[\Sym^3V]$ である。
普遍平面三次曲線のPicard群の形式完備化は、標数零の形式対数により
$D_U$ の零切断完備化と同一視できる。
形式対数は基底変換と自己同型に自然であり、各有限jetが $S_U$ 上で正則なので、
この同一視は $GL(U)\times GL(Z)$-同変である。
これは平面曲線が特異または非被約である場合も含む。

この形式Picard族からの射影族は
\[
                    D_U^{-1}\otimes i^*(j_*L)[-1]
\]
である。$D_U^{-1}$ の因子によって原点が $E_F$ に一致する。
局所的なuniversal直線束の選択の差は基底からの直線束によるtwistなので、
simple対象のrigidifiedモジュライへの射は貼り合わさる。
従って普遍座標への写像 $w(F,s)$ の各係数は $S_U$ 上で正則、従って $F$ の多項式である。

$GL(U)$ の中心 $r\id_U$ はPicardの接線 $D_U$ には重み0で作用し、
$W$ の三成分
\[
 (\Lambda^2U\otimes T)\ \oplus\ (U\otimes Z\otimes T)
                         \ \oplus\ (\Lambda^2Z\otimes T)
\]
にはそれぞれ重み $0,-1,-2$ で作用する。
一方、$F$ の各成分の重みは $0,1,2,3$ である。
多項式性と同変性から、$w(F,s)$ の負の重みの二成分は全次数で零であり、
像は常に $\Lambda^2U\otimes T$ の直線内にある。
Picard族のKodaira--Spencer写像はtautologicalな包含で非零なので、
この写像はパラメータの可逆な変更を除いて、その直線の形式完備化全体を与える。

滑らかな $X_F$ は平面を含まないので、全ての $U$ について $F|_Z\ne0$ である。
従ってこの議論はGrassmannian全体に適用でき、Picard族からの射は
tautological直線束の標準縮約と、各ファイバーで可逆な形式パラメータ変更との合成になる。
そのscheme-theoretic imageは標準Pl\"ucker錐である。
定理\ref{thm:picard-compatibility-equality}の全次数の像の等式から結論を得る。
\end{proof}


\subsection{高次境界項の共変量と第11次までの標準形}
\label{subsec:boundary-higher-jets}

以下では定理\ref{thm:universal-bn-equation}の正則な普遍同変座標を用いる。
$\Delta=\det V$、$W=\Lambda^2V\otimes\Delta^{-1}$ とし、
$P(F,a)=F(q(a))$ と書く。$O(a^r)$ は $a$ に関する次数が $r$ 以上という意味である。
本節の座標変換と単元は、特に断らない限り、全非零三次式上で正則かつ
$GL(V)$-同変なものに限る。

\begin{lemma}[境界共変量の有限次元計算]\label{lem:boundary-covariant-dimensions}
\eqref{eq:universal-bn-covariants}の空間は
\[
\begin{array}{c|rrrrrrr}
 n&6&7&8&9&10&11&12\\ \hline
 \dim\mathcal C_n&1&0&1&0&2&0&5
\end{array}
\]
という次元を持つ。従って普遍境界方程式の $d_7,d_9,d_{11}$ は零である。
これは表示した有限範囲の主張であり、全ての奇数次の消滅は主張しない。
\end{lemma}

\begin{proof}
Littlewood恒等式\cite[式(6)]{BeteaWheeler}は、5変数で
\[
 \Sym^n(\Lambda^2V)=
 \bigoplus_{\substack{r\ge s\ge0\\r+s=n}}
           \mathbb S_{(r,r,s,s,0)}V
\]
を与える。$m=n-5$ とおき、$(n-3)^5$ という長方形内で補分割を取ると、
\begin{equation}\label{eq:boundary-plethysm-candidates}
 \dim\mathcal C_n=
 \sum_{j=0}^{\lfloor(m-1)/2\rfloor}
 [\Sym^m(\Sym^3V):
   \mathbb S_{(m+2,m-1-j,m-1-j,j,j)}V]
\end{equation}
となる。例えば
\[
 \Sym^2\Sym^3V=\mathbb S_6V\oplus\mathbb S_{42}V,
\]
\[
 \Sym^3\Sym^3V=
 \mathbb S_9V\oplus\mathbb S_{72}V\oplus\mathbb S_{63}V
 \oplus\mathbb S_{522}V\oplus\mathbb S_{441}V
\]
である。後者は\cite[Example 1.6]{DeBoeckPagetWildon}にも明記されている。
これにより $n=7,8$ の主張を得る。

残る有限計算を再現可能にするため、計算方法も固定する。
$M_m(\beta)$ を
\[
 \prod_{\substack{\alpha\in\mathbb N^5\\|\alpha|=3}}
       (1-tx^\alpha)^{-1}
   =\sum_{m,\beta}M_m(\beta)t^mx^\beta
\]
の整数係数とし、$\rho=(4,3,2,1,0)$ とおく。Weylの交代式から
\begin{equation}\label{eq:boundary-character-alternant}
 [\Sym^m\Sym^3V:\mathbb S_\lambda V]
 =\sum_{\sigma\in S_5}\operatorname{sgn}(\sigma)
     M_m(\lambda+\rho-\sigma\rho).
\end{equation}
負の成分を持つ重みの係数は零とする。この有限和の結果は次の通りである。
\[
\begin{array}{c|l|l}
 n&\text{候補分割}&\text{各重複度}\\ \hline
 9 &(6,3,3),\ (6,2,2,1,1)&0,0\\
 10&(7,4,4),\ (7,3,3,1,1),\ (7,2,2,2,2)&1,0,1\\
 11&(8,5,5),\ (8,4,4,1,1),\ (8,3,3,2,2)&0,0,0\\
 12&(9,6,6),\ (9,5,5,1,1),\ (9,4,4,2,2),\ (9,3,3,3,3)&2,1,2,0
\end{array}
\]
$n=9$ の5行の分割は、$\Sym^4\Sym^3V$ が4個の一行分割のテンソル積の
直和因子であることから現れない。残る $(6,3,3)$ の重複度は3変数の交代式で
\[
                       18-12-16-6+8+8=0
\]
となる。$n=10$ の三つの候補について、\eqref{eq:boundary-character-alternant}の
正符号の和と負符号の和は、それぞれ
\[
                    85-84,\qquad780-780,\qquad2040-2039
\]
であり、$n=11$ では $187-187,\ 2279-2279,\ 7997-7997$ である。
添付の \texttt{nc\_cubic\_verify\_v7.py} は全係数を整数で計算し、
重みの総和とWeyl次元公式による独立な次元検査も行う。
\end{proof}

\begin{lemma}[唯一の八次候補]\label{lem:boundary-eighth-covariant}
$H=\operatorname{Hess}_F(q(a))$ とする。体積を選び、3元部分集合 $I$ に対し
$c_I=\epsilon_{I,I^c}a_{I^c}$ とおけば、
\begin{equation}\label{eq:boundary-theta-eight}
 \Theta_8(F,a)=\frac1{216}
       \sum_{|I|=|J|=3}c_Ic_J\det H_{I,J}
\end{equation}
は $\mathcal C_8$ の非零な元であり、その空間を生成する。
従って普遍境界方程式の八次項は $d_8=c_8\Theta_8$ と書ける。
ここで $c_8\in\CC$ は $F$ に依存しない定数である。
\end{lemma}

\begin{proof}
体積を選ばない表示では
$H\in\Sym^2V\otimes\Delta^{-1}$ である。
$\Lambda^3H$ の二つの $\Lambda^3V$ 因子を、二つの
$a\in\Lambda^2V\otimes\Delta^{-1}$ と外積で縮約すると、
\eqref{eq:boundary-theta-eight}は $\Delta^{-3}$ に値を持つ。
その $F$ の次数は3、$a$ の次数は8なので $\mathcal C_8$ に属する。
Fermat式 $F_0=\sum_{i=0}^4x_i^3$ では
\begin{equation}\label{eq:boundary-theta-fermat}
 \Theta_8(F_0,a)=
       \sum_{i<j}a_{ij}^2\prod_{k\ne i,j}q_k(a).
\end{equation}
下の横断切片でこれは非零なので、補題\ref{lem:boundary-covariant-dimensions}から結論を得る。
\end{proof}

\begin{theorem}[実際の八次項の幾何的消滅]\label{thm:boundary-eighth-vanishing}
実際のBrill--Noether境界について $c_8=0$ である。従って、正則な普遍同変座標で
\begin{equation}\label{eq:boundary-nine-jet}
                         d(F,a)=F(q(a))+O(a^{10})
\end{equation}
が成立する。
\end{theorem}

\begin{proof}
$F=F_0$ とし、$U=\langle x_3,x_4\rangle$、$Z=\langle x_0,x_1,x_2\rangle$ と分ける。
$\Gamma=\mu_3(x_3)\times\mu_3(x_4)$ の $W$ 上の固定部分は $a_{34}=t$ という直線である。
固定平面のFermat曲線
\[
 C_0=\{x_3=x_4=0,\quad x_0^3+x_1^3+x_2^3=0\}
\]
のPicard形式族は、この直線に入り、一次微分が非零なので形式直線全体を与える。
これは補題\ref{lem:universal-chart-plucker}の固定した $U$ に対する記述でもある。
その非零な一般点は、固定した $C_0$ 上の $L(t)$ の評価核に対応する。

命題\ref{thm:local-bn-cubic}により、この点での境界の法方向の三次曲線は
$C_0$ を $|\OO_{C_0}(4)\otimes L(t)^{-3}|$ で埋め込んだ曲線である。
従ってその抽象楕円曲線の $j$ は恒等的に零である。
適合部分空間は同じ座標で標準Pl\"ucker錐なので、$t\ne0$ の点でその法空間は
$\Lambda^2Z$ の三方向である。実際、混合する六成分はGrassmannianの接方向であり、
その混合成分を零にした切片では、適合部分空間は $\Lambda^2Z=0$ となる。
従って次の切片の三次項が、法方向の三次曲線をそのまま計算する。

混合成分を零とし、
\[
 a_{34}=t,\qquad (z_0,z_1,z_2)=(a_{12},-a_{02},a_{01})
\]
とおくと、直接計算により
\begin{equation}\label{eq:boundary-hesse-slice}
 P(F_0,a)=t^3(z_0^3+z_1^3+z_2^3),\qquad
 \Theta_8(F_0,a)=t^5z_0z_1z_2.
\end{equation}
$a$ に関する $n$ 次項の $z$ に関する三次部分は、$t^{n-3}$ の倍である。
法方向の三次式を $t^3$ で割ると
\[
 z_0^3+z_1^3+z_2^3+c_8t^2z_0z_1z_2+O(t^3)
\]
を得る。Fermat三次式の $GL_3$ 軌道の接空間は $z_0z_1z_2$ を除く
九つの三次単項式を張る。従って形式陰関数定理により、これは線形座標変換で
\[
 z_0^3+z_1^3+z_2^3+s(t)z_0z_1z_2,
 \qquad s(t)=c_8t^2+O(t^3)
\]
というHesse形になる。既に得たHesse式の公式から
\[
 j(s)=\frac{s^3(216-s^3)^3}{(s^3+27)^3}
     =512s^3+O(s^6)
\]
なので、$j=512c_8^3t^6+O(t^7)$ である。
これは $j\equiv0$ と両立するためには $c_8=0$ を要求する。
$d_7=d_9=0$ は補題\ref{lem:boundary-covariant-dimensions}で既に分かっている。
\end{proof}

\begin{theorem}[第11次までの標準形と残る普遍係数]\label{thm:boundary-eleven-jet}
正則で同変な単元を掛けることにより、実際の普遍境界方程式を
\begin{equation}\label{eq:boundary-eleven-jet}
 d(F,a)=F(q(a))+
 c_{10}\det\operatorname{Hess}_F(q(a))+O(a^{12})
\end{equation}
と取れる。$c_{10}\in\CC$ は普遍定数である。
本節では候補空間までを決定し、次節の命題\ref{thm:tate-tenth-coefficient}で
$c_{10}=1/96$ を証明する。
第10次の正則な普遍同変座標変換と単元による商は1次元であり、
$H_{10}=\det\operatorname{Hess}_F(q(a))$ で代表される。
\end{theorem}

\begin{proof}
単元の4次補正の空間を
\[
 \mathcal U_4=
 \Hom_{GL(V)}\left(
 \Sym^4\Sym^3V\otimes\Sym^4\Lambda^2V,\Delta^4\right)
\]
と定める。これは、$F$ と $a$ の両方に4次の、scalarに値を持つ普遍不変多項式の空間である。
交代式\eqref{eq:boundary-character-alternant}を二つのcharacterの積へ適用すると
$\dim\mathcal U_4=1$ を得る。その非零生成元を $U_4$ とする。
これは $U_4$ を非零定数倍の任意性まで定義し、以下に正規化は必要ない。
多項式環は整域なので $P U_4\ne0$ である。

一方、$H\in\Sym^2V\otimes\Delta^{-1}$ の行列式は
$(\det V)^2\otimes\Delta^{-5}=\Delta^{-3}$ に値を持つので、
$H_{10}\in\mathcal C_{10}$ である。これは $P$ の倍数ではない。
実際、Fermat式で $q=(1,1,1,1,\eta)$、$\eta^3=-4$ と取れば
\[
                 F_0(q)=0,\qquad
                 H_{10}(F_0,a)=6^5\prod_iq_i\ne0.
\]
この $q$ はある階数4の交代行列のPfaffianベクトルとして実現できる。
従って $P U_4,H_{10}$ は独立であり、$\dim\mathcal C_{10}=2$ から基底をなす。
$d_{10}=\alpha P U_4+c_{10}H_{10}$ と書き、単元 $1-\alpha U_4$ を掛ければよい。
$\mathcal C_{11}=0$ なので余りは $O(a^{12})$ である。

座標変換についても商の範囲を確認する。
$a$ に関する $r$ 次の普遍座標補正の空間は
\[
 \mathcal G_r=
 \Hom_{GL(V)}\left(
 \Sym^{r-1}\Sym^3V\otimes\Sym^r\Lambda^2V,
 \Lambda^2V\otimes\Delta^{r-1}\right).
\]
同じ整数計算により
\[
 \dim\mathcal G_2=\dim\mathcal G_3=\dim\mathcal G_4=0,
 \qquad \dim\mathcal G_5=2,
\]
また1次から3次までの単元補正も零である。
$\phi_5\in\mathcal G_5$ に対し $(dq)_a\phi_5$ は $F$ に4次、$a$ に6次であり、
その共変量空間の次元は
\[
 \dim\Hom_{GL(V)}\left(
 \Sym^4\Sym^3V\otimes\Sym^6\Lambda^2V,
 \mathbb S_{(5,5,5,5,4)}V\right)=1
\]
である。この空間は非零な $qU_4$ で張られるから
$(dP)_a\phi_5\in\CC P U_4$ である。
逆にEuler型補正 $\phi_5=U_4a$ は $(dP)_a\phi_5=6P U_4$ を与える。
従って座標変換と単元の像はちょうど $\CC P U_4$ である。
\end{proof}

\begin{remark}[標準形の適用範囲]\label{rem:boundary-jet-scope}
\eqref{eq:boundary-eleven-jet}は全次数の形式斉次性を証明しない。
また普遍同変な商と、固定した $F$ に対して任意の座標変換を許した
Jacobian商 $\CC[a]/J_{P_F}$ は異なる。
普遍な商の計算だけでは、$H_{10}(F,-)$ の固定した $F$ でのJacobian類は決まらない。
Fermatの場合には補題\ref{lem:fermat-tenth-jacobian}で通常のJacobian商での非零性を別に証明する。
実際の係数は命題\ref{thm:tate-tenth-coefficient}で $1/96$ と決定する。
また命題\ref{thm:all-smooth-hessian-jacobian}は、通常のJacobian類の非零性を
全ての滑らかな $F$ へ拡張する。
$H_{10}$ は $q$ に5次なので、階数2の部分に沿って少なくとも5次で消える。
従って法方向の三次曲線の $j$ を使った上の議論では、この係数は検出できない。
\end{remark}

\begin{remark}[厳密検算の範囲]
\texttt{nc\_cubic\_verify\_v7.py} はPython標準ライブラリのみを用い、
表の共変量、単元・座標補正の次元、Pfaffianの切片
\eqref{eq:boundary-hesse-slice}、および $j$ の先頭係数512を整数または有理数で検証する。
検算時の実行時間は約1秒であった。
普遍座標の存在、properなincidenceの像、Picard族の同定、
実際の三次曲線のisotrivialityは幾何的証明の入力であり、計算コードの検証対象ではない。
\end{remark}


\begin{remark}[9次の局所環と一段階解消での10次候補]
命題\ref{thm:boundary-eighth-vanishing}を局所環で述べると、
$A_X=\widehat{\OO}_{D,e}$ に対し
\[
 A_X/\mathfrak m_X^{10}
 \simeq\CC[a_1,\ldots,a_{10}]/\bigl((a)^{10},P_X\bigr)
\]
である。これは命題\ref{thm:boundary-sixth-jet}の6次Torelliを保ったまま、
接錐との一致を9次まで延ばす。

また補題\ref{lem:universal-chart-plucker}により、
命題\ref{thm:bn-single-blowup}の一段階解消を同じ普遍座標で使える。
$\operatorname{Tot}_{\PP(V^*)}(\Lambda^2Q\otimes T)$ の局所chartで、
基底点の代表を $p$、束の座標を $b$ とし、局所Pfaffian $r(b)$ を
$q(b)=r(b)p$ で正規化する。$r$ は $b$ に関して2次である。
このとき
\[
 P(F,b)=r(b)^3F(p),\qquad
 H_{10}(F,b)=r(b)^5\det\operatorname{Hess}_F(p).
\]
総変換の例外因子の係数が3なので、$r^3$ で割ったincidenceのstrict transformは
\[
 F(p)+c_{10}r(b)^2\det\operatorname{Hess}_F(p)+O(b^6)=0
\]
と書ける。余りの各斉次成分も $r^3$ で割れることは、
例外因子に沿う位数3と、$r$ の斉次性から従う。
従って境界の10次候補は、解消空間内では三次超曲面の方程式に対する
束方向4次の補正として現れる。
この式の係数は、次節で実際のincidenceをTate行列から構成して計算する。
\end{remark}

\subsection{Fermatにおける10次Hessian候補のJacobian類}

普遍な同変座標変更で消えないことと、固定した三次式上の任意の座標変更で
消えないことは区別する必要がある。次の計算は、Fermatに固定した通常のJacobian商についての主張である。

\begin{lemma}[10次候補の非零類]
\label{lem:fermat-tenth-jacobian}
$F_0=x_0^3+\cdots+x_4^3$ とし、符号付きPfaffianを $q_0,\ldots,q_4$ とする。
\[
 S_a=\CC[a_{ij}:0\le i<j\le4],\qquad
 P=\sum_{i=0}^4q_i^3,\qquad J_P=(\partial_{a_{ij}}P)
\]
とおくと、
\[
                  \Xi:=\prod_{i=0}^4q_i\notin(J_P)_{10}.
\]
従って $\det\operatorname{Hess}_{F_0}(q)=6^5\Xi$ も、
通常の次数10のJacobian商で非零である。
\end{lemma}

\begin{proof}
有限群 $\Gamma=\mu_3^5\rtimes S_5$ のtraceを含む作用を用いる。
対角元に対して
\[
 a_{ij}\longmapsto
       \frac{\zeta_i\zeta_j}{\prod_k\zeta_k}\,a_{ij}
\]
であり、$P$ と $\Xi$ は $\mu_3^5$ に不変である。
置換に対して両者のcharacterは $\operatorname{sgn}(\sigma)$ である。
次数10ではtraceからの符号の10乗は1なので、置換の計算には
\[
 \sigma(a_{ij})=a_{\sigma(i)\sigma(j)},\qquad a_{ji}=-a_{ij}
\]
という通常の交代行列の作用を使える。
符号付きReynolds和を
\[
                 R(h)=\sum_{\sigma\in S_5}
                          \operatorname{sgn}(\sigma)\,\sigma(h)
\]
と書く。

もし $\Xi\in(J_P)_{10}$ なら、次数5の係数を持つ表示を
$\Gamma$ の指定characterへ射影できる。
全ての偏微分は置換により $\partial_{a_{01}}P$ へ移るので、その像は
$R(M\partial_{a_{01}}P)$ で張られる。
ここで $M=\prod a_{ij}^{e_{ij}}$ の頂点次数を
$d_i=\sum_{j\ne i}e_{\min(i,j),\max(i,j)}$ とすると、対角群による射影は
\[
 d_0,d_1\equiv2\pmod3,\qquad d_2,d_3,d_4\equiv1\pmod3,
 \qquad \sum_i d_i=10
\]
を満たす単項式だけを残す。
次数5の全2002単項式のうち該当するものは20個であり、
$\{0,1\}$ を保つ部分群 $S_2\times S_3$ の下で次の四軌道に分かれる。

\[
                         m=a_{03}a_{04}^3a_{12}^2a_{13}a_{14}a_{23}^2
\]
の係数を抽出する。各軌道内のReynolds和は代表のものと符号だけ異なる。
Pfaffianの三項表示を代入した係数の計算は次のとおりである。
表の最後の欄には、120置換中の非零の寄与だけを、重複を含めて記した。

\begin{center}
\begin{tabular}{@{}clcl@{}}
\toprule
 & 代表 $M$ & 軌道の大きさ & $[m]R(M\partial_{a_{01}}P)$ の寄与\\
\midrule
$A$ & $a_{01}^2a_{02}a_{03}a_{04}$ & 2 & なし\\
$B$ & $a_{01}a_{02}a_{12}a_{23}a_{24}$ & 3 & なし\\
$C$ & $a_{02}^2a_{12}^2a_{34}$ & 3 & $6,6,-6,-6$\\
$D$ & $a_{02}^2a_{12}a_{13}a_{24}$ & 12 & $3,3,-3,-3$\\
\bottomrule
\end{tabular}
\end{center}

従って20列全てについて、この係数は零である。
一方 $[m]\Xi=1$ である。実際、$m$ を得る各Pfaffianの項の選択は
\[
 (a_{14}a_{23})(-a_{04}a_{23})(a_{04}a_{13})
                 (-a_{04}a_{12})(a_{03}a_{12})=m
\]
だけである。
$\Xi$ 自身は指定characterに属するので、$R(\Xi)=120\Xi$ である。
以上の係数抽出は、$\Xi$ が射影されたJacobian像に属する可能性を排除する。

なお $R_A:=R(A\partial_{a_{01}}P)$ などと略記すると、
同じ有限計算では $R_C=0$、$3R_D=-R_A$ となる。
\[
 r=a_{04}a_{14}^4a_{23}^3a_{24}a_{34},\qquad
 s=a_{03}a_{04}^3a_{12}^3a_{13}a_{23}a_{34}
\]
の二つの係数を取った行列は
\[
 \begin{pmatrix}{[r]R_A} &{[r]R_B}\\{[s]R_A} &{[s]R_B}\end{pmatrix}
                  =\begin{pmatrix}18&12\\0&24\end{pmatrix}
\]
なので、$R_A,R_B$ は独立である。
したがって指定characterにおけるJacobian像の次元は2で、
$\Xi$ を加えた次元は3になる。
全ての計算は整数係数の多項式と有理数上の線形代数で行い、
検算コードでも20列の係数と二つのrankを再現している。
\end{proof}

\begin{remark}[実際の境界の係数は別の問題]
この補題はFermatに固定した候補のJacobian類を計算している。
実際の境界方程式におけるこの候補の係数を決定したものではなく、
この補題だけでは境界の形式斉次性の否定も、一般の $F$ への非零性の拡張も結論しない。
両方を次節で証明する。
\end{remark}

\begin{remark}[係数が非零なら生じる形式的障害]
もし実際の普遍係数 $c_{10}$ が非零なら、Fermatの境界は六次錐と形式同型ではない。
ここでは、低次数の座標変更が10次の障害を隠さないことも確認できる。
形式同型の線形部分は六次式の線形自己同型として取り除ける。
接空間上恒等の座標変更と定数項1の単元の積はprounipotentな接触変換であり、
その対数はベクトル場と関数の乗法作用素の和 $Y+b$ である。
これらの作用素の交換子は再び同じ形で、次数を少なくとも1上げる。
従って $(u\Phi^*P-P)$ の最初の非零次数は $((Y+b)P)$ の最初の非零次数と一致する。
7次から9次まで変化がないなら、10次の変化は
\[
       \sum_{i<j}(Y_5)_{ij}\,\partial_{a_{ij}}P+b_4P
                         \in(J_P)_{10}
\]
に入る。Eulerの公式により $P$ 自身も $J_P$ に属する。
補題\ref{lem:fermat-tenth-jacobian}の非零類はこの変化では打ち消せない。
ここまでは $c_{10}\ne0$ を仮定した議論である。
次節で係数を計算し、同じ障害を全ての滑らかな $F$ で証明する。
\end{remark}

\clearpage
\section{Tate行列分解による境界の明示化と十次係数}
\label{sec:tate-boundary-tenth}

本節では、対象の変形を実際の有限階数行列分解から構成する。
出発点は剰余体の安定化を表すKoszul--Tate行列分解である
\cite[\S\S 2.3, 4.1]{DyckerhoffMF}。以下ではOrlovのwindowの選択に依存せず、
階数8束の評価写像からgluing対象そのものを回復する。
これにより、前節で未決定であった十次係数を計算できる。
以下の数値は、評価写像の接座標と通常の符号付きPfaffianを用いる
Tate規格化での値である。

\subsection{共変なMaurer--Cartan再帰}

$S=\CC[x_0,\ldots,x_4]$、$F\in S_3$ とし、
$\theta_i$ を外積作用、$\xi_i$ を対応する縮約作用とする。
従って $\{\theta_i,\xi_j\}=\delta_{ij}$ である。
\[
 f_i=\frac13F_i,\qquad B_{ij}=\frac16F_{ij},\qquad
 \mathsf D_0=\sum_i x_i\xi_i+\sum_i f_i\theta_i,
 \qquad \eta_i=\xi_i-\sum_jB_{ij}\theta_j
\]
とおく。Eulerの恒等式から
\[
 \mathsf D_0^2=F,\qquad \{\mathsf D_0,\theta_i\}=x_i,\qquad
 \{\mathsf D_0,\eta_i\}=0,\qquad \{\eta_i,\eta_j\}=-2B_{ij}.
 \tag{T1}\label{eq:tate-clifford-relations}
\]
$\eta_{[J]}$ は $\eta$ の積の反対称化を表す。
正規順序 $\theta_I\eta_{[J]}$ により、自己準同型dg代数の underlying
complex は $x_0,\ldots,x_4$ のKoszul複体と同定される。
この反対称化を使うことにより、同定は $GL(V)$ と可換である。

$x$ 次数 $d$、$\theta$ 次数 $|I|$ の成分に対して
\[
 h\bigl(p(x)\theta_I\eta_{[J]}\bigr)
 =\frac1{d+|I|}\sum_i
   \frac{\partial p}{\partial x_i}\,
   (\theta_i\wedge\theta_I)\eta_{[J]}
 \tag{T2}\label{eq:tate-homotopy}
\]
と定め、$d+|I|=0$ では零とする。
ここで微分するのは係数 $p(x)$ だけであり、$\eta$ は形式的な基底として扱う。
$\delta=[\mathsf D_0,-]_{\mathrm s}$ と書けば
\[
 \delta h+h\delta=1-\pi,
\]
ただし $\pi$ は $x=0$、$\theta$ 次数零への射影である。

体積元を固定し、$a=(a_{ij})\in\Lambda^2V\otimes(\det V)^*$ に対して
\[
 A_1(a)=\sum_{|I|=3}\epsilon(I,I^c)a_{I^c}\eta_{[I]},
 \qquad
 A_n(a)=-h\left(\sum_{i=1}^{n-1}A_i(a)A_{n-i}(a)\right)
 \quad(n\ge2)
 \tag{T3}\label{eq:tate-mc-recurrence}
\]
と定める。
これは $a$ の次数ごとの有限計算である。

\begin{lemma}[再帰の正当性]\label{lem:tate-mc-existence}
$\mathsf D_a=\mathsf D_0+\sum_{n\ge1}A_n(a)$ は形式的な次数付き行列分解であり、
$\mathsf D_a^2=F$ を満たす。
係数は $F$ の係数に多項式的に依存し、構成全体は $GL(V)$ と可換である。
\end{lemma}

\begin{proof}
重みを $\mathrm{wt}(x_i)=1$、$\mathrm{wt}(\theta_i)=-1/2$、
$\mathrm{wt}(\eta_i)=1/2$ とする。
$\mathsf D_0,A_1$ は重み $3/2$ の奇作用素であり、$h$ は重み $-3/2$ を持つ。
低次のMaurer--Cartan方程式を帰納的に仮定すると、
$R_n=\sum_{i=1}^{n-1}A_iA_{n-i}$ は $\delta$-closedで、重み3の偶作用素である。
$\pi R_n$ が非零なら、それは $\eta$ 次数6でなければならないが、
$\eta$ は5個しかない。従って $\pi R_n=0$、
$\delta(-hR_n)=-R_n$ となる。
これで全ての次数の方程式が解ける。
反対称化による正規順序、Koszul homotopy、および
$F_i/3,F_{ij}/6$ の構成は全て共変であり、除算は正整数によるものだけである。
\end{proof}

\subsection{rank 8束からgluing objectを作る}

行列分解の偶数側、奇数側をそれぞれ
\[
 \begin{aligned}
 P_0&=S(-9)\oplus S(-8)^{10}\oplus S(-7)^5,\\
 P_1&=S(-7)^5\oplus S(-6)^{10}\oplus S(-5)
 \end{aligned}
\]
と置く。各直和項は順に $\Lambda^0,\Lambda^2,\Lambda^4$、
および $\Lambda^1,\Lambda^3,\Lambda^5$ に対応する。
$\mathsf D_a$ の二つの部分を
\[
 \Phi_a:P_0\longrightarrow P_1,
 \qquad \Psi_a:P_1\longrightarrow P_0(3)
\]
と書く。$\Phi_a$ の $S(-7)^5\to S(-7)^5$ 成分を
$K_F(a)$ とする。これは $x$ に依存しない $5\times5$ 行列である。
$\Lambda^4$ の第 $j$ 列に $(-1)^j\theta_{\{0,\ldots,4\}\setminus j}$ を使うと
\[
                         K_F(a)=-a+O(a^2).
 \tag{T4}\label{eq:tate-K-linear}
\]

$X=V(F)$ として
\[
 i_*\mathcal M_a=\operatorname{coker}\bigl(\Psi_a(4)\bigr)
\]
を定める。その $\PP^4$ 上の表示は
\[
 0\longrightarrow
 \OO(-3)^5\oplus\OO(-2)^{10}\oplus\OO(-1)
 \xrightarrow{\ \Psi_a(4)\ }
 \OO(-2)\oplus\OO(-1)^{10}\oplus\OO^5
 \longrightarrow i_*\mathcal M_a\longrightarrow0.
 \tag{T5}\label{eq:tate-bundle-resolution}
\]
$\Psi_a$ の単射性は $\Phi_a\Psi_a=F$ と $F$ の非零因子性から従う。
任意のArtin基底へ還元しても $F$ は非零因子であり、同じ単射性が成立する。
従ってcokernelは基底上flatで、この表示はArtin基底変換と可換である。

\begin{lemma}[実際のuniversal対象]\label{lem:tate-universal-object}
$\mathcal M_a$ は $X$ 上のrank 8束の形式変形であり、
\[
 \mathcal E_a=\operatorname{Cone}\bigl(
 H^0(\mathcal M_a)\otimes\OO_X\longrightarrow\mathcal M_a\bigr)
 \tag{T6}\label{eq:tate-evaluation-object}
\]
は $E$ のuniversalな形式変形を与える。
接写像は $a\mapsto\pm a$ であり、接座標の比較に現れる定数は符号だけである。
\end{lemma}

\begin{proof}
中心でKoszul複体を比較すると、体積線を明記した完全列
\[
 0\longrightarrow\det V\otimes\OO_X(-1)
 \longrightarrow\Lambda^4V\otimes\OO_X
 \longrightarrow\mathcal M_0
 \longrightarrow Q\longrightarrow0
 \tag{T7}\label{eq:tate-central-evaluation}
\]
を得る。
具体的には、\eqref{eq:tate-bundle-resolution}の $\Lambda^4V\otimes\OO$
を除いた商を考える。$\Lambda^3V\otimes\OO(-2)$ からのKoszul写像を先に除くと、
残る標的は $\OO(-2)\oplus Q$ である。
$f=\sum_i f_i\theta_i$ と書けば、$\Lambda^1$ からの写像は
\[
 v\longmapsto\bigl(\iota_xv,-f\,\iota_xv\bigr)
\]
というgraphを通る。これは
$\iota_x(f\wedge v)=Fv-f\,\iota_xv$ と $F|_X=0$ による。
従ってcokernelは $Q$ であり、核は $\Lambda^5V\otimes\OO(-1)$
からのKoszul写像である。特に $\mathcal M_0$ は $Q$ と
$Q^*\otimes\det V$ の拡大なので局所自由である。
Artin基底上でもflatnessとNakayamaから局所自由性が従う。

\eqref{eq:tate-bundle-resolution}のコホモロジーから
$R\Gamma(\mathcal M_a)=\OO^{\,5}$ および
$R\Gamma(\mathcal M_a(-1))=0$ を得る。
従って $\mathcal E_a\in\Ku(X)$ である。
中心の三角は
\[
 \det V\otimes\OO_X(-1)[1]\longrightarrow\mathcal E_0
 \longrightarrow Q\longrightarrow\det V\otimes\OO_X(-1)[2].
\]
$R\Gamma(\mathcal E_0(-1))=0$ のため、最後の射は非零である。
これは $T=H^3(\OO_X(-2))\simeq(\det V)^*$ を含む
\eqref{eq:E-triangle}そのものである。

接写像もこの表示で追跡できる。
$\mathrm{vol}=\theta_0\cdots\theta_4$ とすると、
$A_1\mathrm{vol}$ の $\Lambda^2$ 成分は $-\sum a_{ij}\theta_i\theta_j$ である。
変形した関係式 $\Psi_a\mathrm{vol}=0$ を用い、商
$\Lambda^2V\otimes\OO(-1)\to Q$ に $\iota_x$ を適用すると、
コホモロジー間の一次微分は標準評価切断 $\pm\iota_x(a)$ になる。
この計算には $1/3$ などの追加の係数は現れない。
従って接写像は同型であり、$\Def_E$ の滑らかさから形式的同型となる。
\end{proof}


\begin{remark}[Serre traceとの規格化の比較]
\label{rem:tate-serre-trace-normalization}
Tate模型の二拡大は、通常のhypersurface residueで定めたtraceに対しても
符号を除いて係数1の生成元である。
実際、$U_i=(x_i\ne0)$ 上で $s_i=e_i/x_i$ とおき、
$f=\sum_i(F_i/3)e_i$ と書く。
$\mathcal M_0\to Q$ の局所分裂は
\[
 v\longmapsto(-s_i\wedge f\wedge v,v)
\]
で与えられる。
この差を二つのEuler列で順に持ち上げると、
$\Ext^2(Q,\OO_X(-1))\simeq H^3(\OO_X(-2))$ に対応する
\v Cech 3-cocycleは
\[
 \gamma_{ijkl}
   =\pm\det(s_i,s_j,s_k,s_l,f)
   =\pm\epsilon_{ijklm}\frac{f_m}{x_ix_jx_kx_l}
\]
となる。その $\PP^4$ 上への持ち上げのcoboundaryは
\[
 \pm\frac{\sum_m x_mf_m}{x_0x_1x_2x_3x_4}
   =\pm\frac{F}{x_0x_1x_2x_3x_4}.
\]
従って $0\to\OO_{\PP^4}(-5)\xrightarrow{F}\OO_{\PP^4}(-2)
\to\OO_X(-2)\to0$ のboundary mapは、この類を
$H^4(\OO_{\PP^4}(-5))$ の標準生成元
$\pm1/(x_0x_1x_2x_3x_4)$ に送る。
$f_i=F_i/3$ は $\sum_i x_if_i=F$ を満たすための規格化であり、
この比較に追加の係数3は現れない。
残る符号は接座標 $a$ の符号に対応するが、$P_F$ と $H_{10}$ の
$a$ 次数はいずれも偶数なので、係数 $1/96$ を変えない。
\end{remark}

\begin{proposition}[実際のincidenceの明示表示]\label{prop:tate-actual-incidence}
$\Phi_a$ を最後の $\Lambda^4$ 直和項に制限し、$\OO(7)$ でtwistした
$16\times5$ 行列を
\[
 \mathsf T_F(x,a):\OO^5\longrightarrow
 \OO^5\oplus\OO(1)^{10}\oplus\OO(2)
\]
と書く。その最初の5行は $K_F(a)$ である。
$X\times\widehat U$ 上の実際のBN incidenceは
\[
              \mathcal I=V\bigl(I_5(\mathsf T_F(x,a))\bigr)
 \tag{T7a}\label{eq:tate-actual-incidence}
\]
で与えられる。ここで $I_5$ は全ての $5\times5$ minorのイデアルである。
従って点変数 $x$ に関しては、次数0、1、2の行しか必要ない。
\end{proposition}

\begin{proof}
$X$ 上では $\ker\Phi_a=\operatorname{im}\Psi_a$ なので、$\Phi_a$ は
$\mathcal M_a$ を局所自由な像へ同型に送る。
この同型の下で $\mathsf T_F$ は評価写像
$\OO^5\to\mathcal M_a$ と、その像の局所的にsplitする包含の合成である。
従って両者のrank条件とdeterminantal idealは一致する。
局所的に可逆な $4\times4$ blockを除くと、評価写像は
$[\OO\to\OO^4]$ の四つの成分へ還元され、最大minorのイデアルは
前節までに用いたincidenceのイデアルとスキームとして一致する。
中心でのrankは4であり、近傍ではrankが3以下になる点はない。
$\mathcal E_a=[\OO^5\to\mathcal M_a]$ の導来ファイバーから
$\dim\operatorname{Hom}(\mathcal E_a,\OO_x)\ge4$ は評価rankが4以下という条件に一致する。
\end{proof}

\subsection{境界を定める有限行列式}

\begin{theorem}[境界の行列式表示]\label{thm:tate-boundary-determinant}
Tate座標において、BN境界 $D$ の定義式は、単元を除いて
\[
                         \det K_F(a)=0
 \tag{T8}\label{eq:tate-boundary-determinant}
\]
である。先頭項は $2F(q(a))$ である。
\end{theorem}

\begin{proof}
$K_F$ は $\mathsf T_F$ の最初の5行なので、その行列式は
$I_5(\mathsf T_F)$ に属する。従って $\det K_F$ は actual incidence
$\mathcal I$ 上でスキームとして零である。
$D$ は $\mathcal I\to\widehat U$ のスキーム論的像であり、
整なCartier因子であるから、その定義式は $\det K_F$ を割る。
$K_F$ の一次項は交代行列であり、六次の先頭項は
普遍共変式の一意性と対角三次式での計算から $2F(q)$ である。
実際 $F_0=\sum x_i^3$ と、下で使う $a_*$ を取ると
$K_{F_0}(ta_*)=-ta_*+2t^2I_5+O(t^3)$ であり、
$\operatorname{adj}(-a_*)=q(a_*)q(a_*)^{\mathsf t}$ から
六次係数は $2\sum_i q_i^2=10=2F_0(q(a_*))$ となる。
$D$ 自身の先頭項が $F(q)$ であることと比較すると、商は単元となる。
\end{proof}

行列分解から境界の外側も読める。
$K_F(a)$ が可逆なら、$\Phi_a$ の5個の定数pivotを消去でき、
対応する $\Psi_a$ の直和項は $F\cdot\mathrm{id}$ になる。
従って $\mathcal M_a$ は $\OO_X^5$ を直和因子として持つ。
残る $\Psi_a$ の $1\times10$ と $10\times1$ の定数blockは、
それぞれ一次項が $a$ の10座標なので、成分の生成するイデアルはNakayamaにより
極大イデアルである。最初の5組の消去は異なる内部次数を結ぶ正次数の行列で行われるため、
この二つの定数blockを変えない。従って原点以外で各一つのunitを消去できる。
残るのは $9\times9$ の線形行列であり、
$\mathcal E_a$ は
\[
 0\longrightarrow\OO_{\PP^4}(-2)^9
 \longrightarrow\OO_{\PP^4}(-1)^9
 \longrightarrow i_*\mathcal E_a\longrightarrow0
\]
を持つnormalizedなrank 3 Ulrich束である。

\subsection{小さい有理数行列による十次項の証明書}

\begin{theorem}[普遍十次係数]\label{thm:tate-tenth-coefficient}
上記の評価写像・Pfaffianの規格化では
\[
 \boxed{\displaystyle
 d(F,a)=F(q(a))+\frac1{96}\det\bigl(\operatorname{Hess}_F(q(a))\bigr)
       +O(a^{12}).}
 \tag{T9}\label{eq:tate-tenth-coefficient}
\]
ここでは定義式に掛かる単元の項を除いている。
特に十次Hessian類は零でない。
\end{theorem}

\begin{proof}
この再帰は全ての非零三次式上で正則かつ同変であり、接写像もedge同型に一致する。
従って第\ref{sec:boundary-higher-jets}節の普遍座標と同じクラスの座標である。
命題\ref{thm:boundary-eleven-jet}と
補題\ref{lem:tate-universal-object}により、普遍係数を一つの標本で決めればよい。
\[
 F_*=x_0^3+x_1^3+x_2^3+x_3^3-4x_4^3,\qquad
 a_*=(e_1-e_0)\wedge(e_2-e_0)+(e_3-e_0)\wedge(e_4-e_0)
\]
を選ぶ。このとき
\[
 q(a_*)=(1,1,1,1,1),\quad F_*(q(a_*))=0,\quad
 \det\operatorname{Hess}_{F_*}(q(a_*))=-31104.
\]
再帰\eqref{eq:tate-mc-recurrence}から、必要な定数blockは
\[
 K_{F_*}(t a_*)=tA+t^2B+t^4C+t^6D+O(t^7)
\]
となり、三次・五次blockは零である。具体的には
\[
 A=\begin{pmatrix}
 0&-1&1&-1&1\\1&0&-1&0&0\\-1&1&0&0&0\\
 1&0&0&0&-1\\-1&0&0&1&0
 \end{pmatrix},\qquad
 B=\operatorname{diag}(2,2,2,2,-8),
\]
\[
 C=\frac13\begin{pmatrix}
 0&-4&-4&16&16\\-4&0&8&12&32\\-4&8&0&12&32\\
 16&12&12&0&48\\16&32&32&48&0
 \end{pmatrix},
\]
\[
 D=\frac1{90}\begin{pmatrix}
 0&-3040&-3040&-3040&-3040\\
 -3280&-560&-4606&-2437&-1277\\
 -3280&-4606&-560&-4677&-5837\\
 -3280&-2437&-4677&-560&-1206\\
 -3280&-2477&-4637&-2406&2240
 \end{pmatrix}.
\]
これらは全て有理数での計算であり、各次数で
$[\mathsf D_0,A_n]_{\mathrm s}+\sum_{i+j=n}A_iA_j=0$ を検証できる。
十次の行列式には $B^5$、$A^2B^2C$、$A^4D$ の三種だけが寄与し、
その値は順に
\[
                    -128,\qquad\frac{640}{3},\qquad-\frac{2200}{3}
\]
である。従って
\[
                  [t^{10}]\det K_{F_*}(ta_*)=-648.
\]
六次先頭係数は $2F(q)$ である。
普遍式の10次項は $P U_4$ の倍数と $H_{10}$ の和であり、7--9次項は零である。
この標本では $F(q)=0$ なので、正規化単元に由来する $P U_4$ の項も零となる。
従って
\[
           2c_{10}(-31104)=-648,\qquad c_{10}=\frac1{96}.
\]
\end{proof}

補題\ref{lem:fermat-tenth-jacobian}と組み合わせると、
Fermat三次三次元多様体のBN境界は、その六次接錐と形式的に同型ではない。
また、一段blowup上では $q(a)=r(b)p$ なので、strict transformの方程式は
\[
 F(p)+\frac1{96}r(b)^2\det\operatorname{Hess}_F(p)+O(b^6)=0
\]
となる。この非零の四次法方向補正が、接錐では見えなかった境界の情報である。

\begin{remark}[計算証明書の範囲]
\texttt{nc\_cubic\_verify\_v8.py} は、上の標本のMaurer--Cartan方程式を
6次まで有理多項式として検証し、行列式の三つの寄与を再計算する。
Fermat式についても $q=(1,5,-7,-8,9)$ で独立に
$[t^{10}]\det K=408240$、$\det\operatorname{Hess}_{F_0}(q)=19595520$ を得る。
さらに32次元外積空間の実際の作用素を用い、別の積の実装で有限体上の恒等式も確認する。
数値だけから境界を同定したのではなく、実際の対象とscheme-theoretic imageの同定は
命題\ref{lem:tate-universal-object}--\ref{thm:tate-boundary-determinant}の証明による。
\end{remark}

\subsection{すべての滑らかな三次式に対するHessian類の非零性}
\label{subsec:all-smooth-hessian}

以下の議論では、普遍同変な座標変更への制限を外す。
$\mathfrak a\subset S=\CC[W]$ をPl\"ucker錐
$C=\widehat{\operatorname{Cone}\operatorname{Gr}(2,V)}$ の斉次素イデアルとし、
$P=P_F=F(q)$、$H=H_{10}(F,-)=\det\operatorname{Hess}_F(q)$ と書く。
ここでイデアルの計算は、同じ斉次方程式で定義される代数的な錐上で行う。

\begin{lemma}[三次の位数を保つベクトル場の接性]
\label{lem:cubic-order-forces-tangency}
$F$ を滑らかな三次式とし、$Y\in\operatorname{Der}_{\CC}(S)$ とする。
$Y(P)\in\mathfrak a^3$ ならば
\[
 Y(\mathfrak a)\subset\mathfrak a.
\]
\end{lemma}
\begin{proof}
$C$ の一般点は階数2の二形式であり、その点における法空間は3次元である。
$q$ の法方向の一次部分は、この法空間を対応する平面の3次元座標空間に同定する。
従って $P$ の法方向の三次初期式は、その平面への $F$ の制限である。
滑らかな $X=(F=0)$ の一般の平面切断はBertiniの定理により滑らかなので、
この三次初期式は一般点で滑らかな平面三次曲線を定義する。

より正確には、$R=S_{\mathfrak a}$ は剰余体 $K=\CC(C)$ を持つ
3次元正則局所環である。$Y$ を $R$ に延長すると、
その法方向の定数項は
\[
 v\colon\mathfrak aR/(\mathfrak aR)^2\longrightarrow K
\]
を定める。$f_3\in\operatorname{Sym}^3_K(\mathfrak aR/(\mathfrak aR)^2)$
を $P$ の初期式とすると、$Y(P)$ の二次初期式は $v\lrcorner df_3$ である。
$Y(P)\in\mathfrak a^3$ によりこれは零である。
滑らかな三次式 $f_3$ の三つの偏微分は一次独立なので $v=0$ となる。
実際、非零の定数方向で微分が零なら、その方向を座標に選ぶことで
$f_3$ は二変数の式になり、射影曲線はその方向の点で特異になる。
従って $Y(\mathfrak a)R\subset\mathfrak aR$ であり、
$\mathfrak a$ が素イデアルであることから $Y(\mathfrak a)\subset\mathfrak a$ が従う。
\end{proof}

\begin{theorem}[一般の三次式でも消えない10次Hessian類]
\label{thm:all-smooth-hessian-jacobian}
すべての滑らかな三次式 $F$ に対して
\[
 \boxed{\quad H_{10}(F,-)\notin(J_{P_F})_{10}.\quad}
\]
ここで $J_{P_F}$ は、固定した $F$ において10個の座標すべてで微分した
通常のJacobianイデアルである。
\end{theorem}
\begin{proof}
仮に $H\in(J_P)_{10}$ とすると、係数が5次の斉次ベクトル場 $Y_5$ を使って
$H=Y_5(P)$ と書ける。$P\in\mathfrak a^3$、$H\in\mathfrak a^5$ なので、
補題\ref{lem:cubic-order-forces-tangency}により $Y_5$ は $C$ に接する。

この接性により $Y_5$ はRees代数を保ち、従って吹上げ
\[
 \beta\colon\mathcal H=\operatorname{Bl}_C(W)
 \simeq\operatorname{Tot}_{\PP(V^*)}(\Lambda^2Q\otimes T)
 \longrightarrow W
\]
に正則ベクトル場 $\widetilde Y$ として持ち上がる。
このベクトル場の束方向の重みは4である。
また $\mathfrak a\OO_{\mathcal H}$ を保つので、例外因子 $R$ にも接する。

基底点の局所代表を $p$ とし、$q(b)=r(b)p$ と書く。
$r$ は例外因子の局所方程式で、束方向2次である。
すると
\[
 \beta^*P=r^3F(p),\qquad
 \beta^*H=r^5\det\operatorname{Hess}_F(p).
\]
$\widetilde Y(r)\in(r)$ を使って $Y_5(P)=H$ を $r^3$ で割ると
\[
 \widetilde Y(F)+3\frac{\widetilde Y(r)}rF
       =r^2\det\operatorname{Hess}_F.
\]
従って基底の $X=(F=0)$ 上では、$\widetilde Y$ の基底方向成分が
\[
 \xi\in H^0\bigl(X,T_{\PP^4}|_X\otimes\operatorname{Sym}^4N^*\bigr),
 \qquad
 dF(\xi)=r^2\det\operatorname{Hess}_F
 \tag{H4}\label{eq:hessian-blowup-vectorfield}
\]
を与える。束方向成分は $F$ に作用しない。

一方、$R_0=Q^*$ とおくと、定数直線因子を除いて
\[
 \operatorname{Sym}^4N^*
   =\mathbb S_{4400}R_0\oplus\mathbb S_{3311}R_0
      \oplus(\det R_0)^2.
\]
最後の直線因子は $\OO_X(2)$ であり、その包含は $r^2$ で与えられる。
定数因子を正確に戻すと、この因子は $\OO_X(2)\otimes T^{-2}$ である。
従って \eqref{eq:hessian-blowup-vectorfield}の最後の因子への射影は、normal sequence のconnecting mapの下で
\[
 [\det\operatorname{Hess}_F]\in
 H^1(T_X(2))\otimes T^{-2}
 \simeq\operatorname{Jac}(F)_5\otimes T^{-2}
\]
が零であることを要求する。
ここで $H^1(T_{\PP^4}|_X(2))=0$ はEuler sequenceから従い、
normal sequenceの大域切断の余核はEulerの等式 $3F=\sum p_iF_i$ により
通常のJacobian環の5次部分である。

しかし $F$ の五つの偏微分は二次の正則列をなし、
$\operatorname{Jac}(F)$ のHilbert級数は $(1+t)^5$ である。
その一次元のsocle $\operatorname{Jac}(F)_5$ はHessian行列式で生成される
\cite[p.~2]{BDSS}。
従って上の類は非零であり、矛盾する。
\end{proof}

\begin{remark}[主張の範囲]
この定理は、三次式を固定した後の任意の形式座標変更に関する非零性である。
普遍同変な共変量の商だけから導いた結論ではない。
また $C$ を保つという仮定をJacobian表示に付加していない。
$H\in\mathfrak a^5$ と法方向の三次式の滑らかさから、
可能なJacobian表示のベクトル場が自動的に $C$ に接することを先に証明している。
\end{remark}

\begin{corollary}[実際の境界の形式非斉次性]
\label{cor:all-smooth-nonhomogeneity}
すべての滑らかな三次三次元多様体 $X=V(F)$ に対して、実際のBN境界
$(D,e)$ は、その六次接錐 $(P_F=0,0)$ と形式同型ではない。
これは任意の形式座標変更と定義式の単元倍を許した主張である。
\end{corollary}
\begin{proof}
命題\ref{thm:tate-tenth-coefficient}により、実際の方程式は
$d=P_F+cH_{10}+O((a)^{12})$、$c=1/96\ne0$ である。
形式超曲面の同型は、周囲の形式座標変更と単元倍に持ち上がる。
その一次部分と単元の定数項は $P_F$ を保つので、それを右から合成して除去すると、
接空間上で恒等、単元の定数項1のcontact変換 $A$ が $A(P_F)=d$ を満たす。
その対数は $\mathfrak m$-進次数を上げる作用素 $Y+b$ であり、
$Y$ はベクトル場、$b$ は関数の乗法作用素である。
$\log(A)(P_F)$ の最初の非零項は $A(P_F)-P_F$ の最初の非零項に等しい。
よって10次部分を取ると
\[
 cH_{10}=Y_5(P_F)+b_4P_F\in(J_{P_F})_{10}.
\]
最後の包含にはEulerの等式 $6P_F=\sum a_i\partial_iP_F$ を使った。
定理\ref{thm:all-smooth-hessian-jacobian}に矛盾する。
\end{proof}


\begin{corollary}[最初の形式同型障害の次数]\label{cor:boundary-tenth-jet-obstruction}
$R_X=\widehat{\OO}_{D,e}$、$R_X^{\mathrm{cone}}=\CC[[a]]/(P_F)$ とし、
それぞれの極大イデアルを $\mathfrak m,\mathfrak n$ とする。
全ての滑らかな $X$ に対して
\[
 R_X/\mathfrak m^{10}\simeq R_X^{\mathrm{cone}}/\mathfrak n^{10},
 \qquad
 R_X/\mathfrak m^{11}\not\simeq R_X^{\mathrm{cone}}/\mathfrak n^{11}.
\]
すなわち、接錐と9次まで一致するが、10次jetで初めて同型が破れる。
\end{corollary}
\begin{proof}
最初の同型は $d=P_F+O((a)^{10})$ から従う。
仮に二番目にも同型があれば、極大イデアルの生成元の像を持ち上げることで、
$\CC[a]/(a)^{11}$ の座標変更が得られる。
両principal idealを比較すると、定数項が非零の単元倍を含めたcontact変換となる。
前の証明の対数をこの有限jet上で取るだけで
$cH_{10}\in(J_{P_F})_{10}$ が必要となり、同じ矛盾を得る。
\end{proof}

この次数10の障害は、六次jetから元の偏極三次三次元多様体を復元できるという
命題\ref{thm:boundary-sixth-jet}と両立する。六次は $X$ の復元に必要な最小次数であり、
十次は同じ $X$ の実際の境界と接錐を区別する最初の次数である。
また命題\ref{prop:boundary-euler}の、接空間上で恒等に作用するEuler導分は
実際の境界には存在しない。

\subsection{Incidenceの形式近傍は法束へ線形化できない}
\label{subsec:incidence-formal-neighborhood}

前節までの非斉次性には、滑らかなincidence自体についての帰結もある。
滑らかさと法束の決定から、形式近傍がその法束と大域的に一致するわけではない。

\begin{corollary}[形式法束模型の非存在]
\label{cor:incidence-not-formal-normal-bundle}
すべての滑らかな三次三次元多様体 $X$ に対し、実際のincidence
$\mathcal I$ の $X=\mathcal I_e$ に沿う形式完備化は、
法束 $N=\Lambda^2Q\otimes T$ の零切断に沿う形式完備化と同型ではない：
\[
 \widehat{\mathcal I}_X
 \not\simeq
 \widehat{\operatorname{Tot}_X(N)}_X.
\]
これは、$X$ への射や周囲の吹上げへの埋め込みを忘れた形式スキームとしても成立する。
\end{corollary}
\begin{proof}
$D_0=(P_F=0)$、$\mathcal I_0=\operatorname{Tot}_X(N)$ とおく。
実際の境界とその接錐には、それぞれproper birationalな射
\[
 \mathcal I\longrightarrow D,\qquad
 \mathcal I_0\longrightarrow D_0
\]
がある。両者の標的はnormalなので、構造層の順像はそれぞれ
$\OO_D$、$\OO_{D_0}$ である。
また両方の原点上のファイバーはスキームとして $X$ であり、
原点の極大イデアルの引き戻しは $X$ のイデアルに一致する。
形式関数定理\cite[Theorem 30.20.5, Tag 02OC]{StacksFormalFunctions}を適用すると
\[
 H^0(\widehat{\mathcal I}_X,\OO)
      \simeq\widehat{\OO}_{D,e},\qquad
 H^0(\widehat{\mathcal I_0}_X,\OO)
      \simeq\widehat{\OO}_{D_0,0}.
\]
従って二つの形式近傍が同型ならば、右辺の完備局所環も同型になる。
これは、実際の係数 $c_{10}=1/96\ne0$ と
定理\ref{thm:all-smooth-hessian-jacobian}から従う
系\ref{cor:all-smooth-nonhomogeneity}に反する。
\end{proof}

\begin{remark}[零切断への射との両立]
$\mathcal I\subset X\times U$ は、もちろん第一射影
$\pi_X\colon\mathcal I\to X$ を持つ。
上の系は、この射の存在を否定していない。
局所的には零切断を持つ滑らかな6次元の形式的な族であっても、
その座標を大域的に法束の線形座標へ揃えることはできない、という主張である。

実際、吹上げ
$\mathcal H=\operatorname{Tot}_{\PP^4}(\Lambda^2Q\otimes T)$ の
基底への射と $\pi_X$ には、既に一次で差が入り得る。
同型
$N\simeq\Lambda^2T_{\PP^4}|_X\otimes\OO_X(-3)$ と $dF$ による縮約から
\[
 0\longrightarrow\Omega_X^1(-1)\otimes T
 \longrightarrow N\longrightarrow T_X\longrightarrow0
\]
を得る。この最後の射が、そのような一次の接方向を与える。
従って $\pi_X$ の存在だけから、解消上の方程式
\[
 F(p)+\frac1{96}r(b)^2\det\operatorname{Hess}_F(p)+O(b^6)=0
\]
の4次項を零と結論することはできない。
\end{remark}


\clearpage
\section{Tate再帰の有限代数化と十二次項}
\label{sec:tate-algebraization-twelfth}
\label{sec:finite-algebraic-tate}

内部次数を固定すると、前節のTate再帰の係数空間は有限次元になり、
再帰全体を有限個の二次方程式の一つの局所解として表せる。
このことから、指定したTate座標における行列分解と境界方程式の
代数性および解析的収束を得る。

\subsection{有限次元の係数空間}

$\mathcal S=\Sym^3V$ を三次式の係数空間とする。
前節の反対称化した正規順序
$p(x)\theta_I\eta_{[J]}$ を用い、内部重み $w$ の係数空間を
$\mathsf E_w$ と書く。
$\eta_i=\xi_i-\sum_jF_{ij}\theta_j/6$ による基底変換は
$\xi$ 次数について対角成分1の三角変換なので、
$F$ に多項式的に依存し、すべての $F\in\mathcal S$ で可逆である。
従って以下の空間と $h$ の階数は、$F$ の特異性によって変化しない。

$m=\deg_xp$、$u=|I|$、$v=|J|$ とすると、重み $3/2$ の条件は
\[
 2m-u+v=3,\qquad 0\le u,v\le5,\qquad m\ge0.
 \tag{A1}\label{eq:tate-finite-weight}
\]
このとき $u+v$ は奇数で、$m\le4$ である。
特に $\mathsf E_{3/2}$ は有限次元である。
また
\[
 \mathsf B=\operatorname{im}
       (h\colon\mathsf E_3\longrightarrow\mathsf E_{3/2})
\]
とおく。

\begin{lemma}[補助係数の個数]
\label{lem:finite-tate-dimensions}
\[
 \dim\mathsf E_{3/2}=5890,\qquad
 \dim\mathsf B=4929,\qquad
 \dim\mathsf E_3=20330.
\]
$\mathsf E_{3/2}$ の $\ker h$ は、10次元の
$\Lambda^3\langle\eta\rangle$ と $\mathsf B$ の直和である。
\end{lemma}
\begin{proof}
条件\eqref{eq:tate-finite-weight}に対応する係数空間の次元は
$\binom5u\binom5v\binom{m+4}4$ である。
一方、$u\ge1$ の成分における $h$ の像の次元は
\[
 \binom5v
 \binom{m+5}{m+u}\binom{m+u-1}{u-1}.
\]
これは正の全次数を持つ多項式de Rham複体の完全性から従う。
例えば $\eta$ 成分を除いた像の次元 $r(m,u)$ は
\[
 r(m,0)=0,\qquad
 r(m,u)=\binom5{u-1}\binom{m+5}4-r(m+1,u-1)
\]
という再帰でも計算できる。
各 $u$ について和を取ると
\[
\begin{array}{c|rr}
u&\dim\mathsf E_{3/2}&\dim\mathsf B\\\hline
0&35&0\\
1&350&210\\
2&1260&935\\
3&2100&1765\\
4&1650&1524\\
5&495&495\\\hline
\text{合計}&5890&4929
\end{array}
\]
を得る。重み3の場合は $2m-u+v=6$ として同じ計算を行う。
最後の主張は $\delta h+h\delta=1-\pi$ と $h^2=0$ から従う。
\end{proof}

\subsection{4929個の二次方程式による\'etale模型}

$a\in W=\Lambda^2V\otimes(\det V)^*$、$b\in\mathsf B$ とし、
\[
 A=A_1(a)+b,\qquad
 \mathcal G_F(a,b)=b+h(A^2)\in\mathsf B
 \tag{A2}\label{eq:tate-finite-quadratic}
\]
とおく。$\mathsf B$ の基底を選べば、
$\mathcal G_F=0$ は $(a,b)$ に関する4929個の二次方程式である。
積の構造定数は $F$ の係数の多項式なので、
これは $\mathcal S\times W\times\mathsf B$ 上の代数的な方程式系になる。

\begin{theorem}[全次数のTate模型の代数化]
\label{thm:tate-algebraization}
方程式\eqref{eq:tate-finite-quadratic}の解スキームには、
零切断 $s(F)=(F,0,0)$ 全体を含む $GL(V)$-不変な開部分 $\mathcal Z$ があり、
射影
\[
 \mathcal Z\longrightarrow\mathcal S\times W
\]
は\'etaleである。
$\mathcal Z$ 上の $\mathsf D_0(F)+A_1(a)+b$ は代数的な次数付き行列分解である。
零切断に沿う形式完備化は、前節の共変なTate再帰による
形式行列分解と一致する。
\end{theorem}
\begin{proof}
零切断上では
\[
 \left.\frac{\partial\mathcal G_F}{\partial b}\right|_{a=b=0}
       =\operatorname{id}_{\mathsf B}.
\]
従って $J=\det(\partial\mathcal G_F/\partial b)$ を可逆にすれば、
Jacobian判定により解スキームの射影は\'etaleになる
\cite[Proposition 59.26.2(3)]{StacksEtaleCriterion}。

この段階で、homotopyで圧縮した方程式から元のMaurer--Cartan方程式を
代数的に回復する。曲率を
\[
 R=\delta A+A^2\in\mathsf E_3
\]
とおく。$b\in\operatorname{im}h$、$h^2=0$、$\pi A=A_1$ より
\[
 hR=b+h(A^2)=0.
\]
また $\pi R=0$ である。
実際、$x$ と $\theta$ を持たない重み3の項には
6個の $\eta$ が必要だが、$\eta$ は5個しかない。
Bianchi恒等式とcontractionの式から
\[
 \delta R=RA-AR,\qquad R=h(RA-AR)
\]
を得る。そこで有限次元空間 $\mathsf E_3$ 上の作用素
\[
 \mathcal T_A(U)=h(UA-AU),\qquad
 \Delta=\det(1-\mathcal T_A)
\]
を考える。$\Delta$ は零切断上で1である。
さらに $\Delta$ を可逆にすると、$R=\mathcal T_A(R)$ から $R=0$ が従う。
従って $J\Delta$ を可逆にした開部分を $\mathcal Z$ に取ればよい。
ここで20330次の行列式 $\Delta$ は開部分を指定する条件であり、
新しい変形パラメータを導入していない。

形式的陰関数の一意性により $b=O(a^2)$ の解は一意である。
その $a$ 次数ごとの係数を取ると、まさに
$A_n=-h\sum_{i=1}^{n-1}A_iA_{n-i}$ を得る。
従って有限次数だけの近似ではなく、元の再帰の全係数が一致する。
また正規順序、$h$、積はすべて共変であり、$J$ と $\Delta$ は
対応する有限次元表現上の共役で不変なので、構成は $GL(V)$ と両立する。
\end{proof}

\begin{remark}[\'etale近傍の意味]
固定した $F$ に対する $\mathcal Z_F$ の次元は10である。
4929は補助未知数と方程式の個数であり、モジュライの次元ではない。
必要なら零切断の上にある他の点を除去できる。
選んだ切断は\'etaleな零ファイバーの開かつ閉な部分なので、
この除去は零切断を含む開部分への制限である。
一方、$b$ が $a$ の多項式になることや、射影が
$W$ のZariski開部分への同型になることは主張していない。
\end{remark}

\subsection{収束半径の評価}

\begin{proposition}[Catalan数による収束評価]
\label{prop:tate-catalan-convergence}
有限次元空間 $\mathsf E_{3/2}$ と $W$ にノルムを選び、
\[
 \|h(UV)\|\le C_F\|U\|\|V\|,
 \qquad \|A_1(a)\|\le L\|a\|
\]
を満たす定数を取る。
$\operatorname{Cat}_k=\frac1{k+1}\binom{2k}k$ とすると
\[
 \|A_n(a)\|
 \le\operatorname{Cat}_{n-1}C_F^{n-1}(L\|a\|)^n.
 \tag{A3}\label{eq:tate-catalan-bound}
\]
従って $4C_FL\|a\|<1$ でTate再帰は絶対収束する。
この収束は、三次式 $F$ の任意のcompact集合上で一様な半径を持つ。
\end{proposition}
\begin{proof}
再帰式から
\[
 \|A_n\|\le C_F\sum_{i=1}^{n-1}\|A_i\|\|A_{n-i}\|
\]
となる。Catalan数の畳み込み恒等式を用いた帰納法で
\eqref{eq:tate-catalan-bound}を得る。
$C_F>0$ の場合、全体のノルムは
\[
 \sum_{n\ge1}\|A_n(a)\|
 \le\frac{1-\sqrt{1-4C_FL\|a\|}}{2C_F}
\]
で抑えられる。$C_F=0$ なら $A_1$ 以外の項は零である。
正規順序の構造定数は $F$ の多項式なので、$C_F$ はcompact集合上で
有界に選べる。$A_1$ はこの基底で $F$ に依存しないので、$L$ は共通に取れる。
\end{proof}

この命題が示すのは、指定したTate再帰自体の収束である。
また定理\ref{thm:tate-algebraization}の\'etale射は局所的にquasi-finiteなので、
その枝の座標 $b(a)$ は代数的冪級数である。
従って $K_F(a)$ と $\det K_F(a)$ も収束する代数的冪級数になる。

\subsection{代数的な対象と境界への帰結}

ここでは滑らかな三次式 $F$ を固定する。
$\mathcal Z_F$ は $W$ 上\'etaleなので滑らかであり、
$X\times\mathcal Z_F$ も滑らかである。
定理\ref{thm:tate-algebraization}の代数的行列分解を用いて
\[
 \mathcal M=\operatorname{coker}(\Psi(4)),\qquad
 \mathcal E=\operatorname{Cone}(\OO^5\to\mathcal M)
\]
を構成する。固定した長さ1の表示と行列分解の恒等式により、
$\mathcal M$ は $X\times\mathcal Z_F$ 上でmaximal Cohen--Macaulayであり、
全空間の正則性から局所自由になる。
零切断に沿って縮小すればrankは8であり、前節と同じ表示から
$R\Gamma(\mathcal M)=\OO^5$、$R\Gamma(\mathcal M(-1))=0$ を得る。
従って $\mathcal E$ は実際の代数的な対象族である。
単純対象のscalar rigidificationを取ったモジュライ
\cite[Theorem 4.2.1, Corollary 4.3.3]{LieblichComplexes}への射は、
零点でのKodaira--Spencer写像が同型なので、その点で\'etaleである。

\begin{corollary}[代数的・解析的な境界模型]
\label{cor:algebraic-analytic-tate-boundary}
固定した滑らかな $F$ に対し、$\mathcal Z_F$ の指定点のZariski近傍上で、
実際のincidenceは $I_5(\mathsf T_F)$、そのscheme imageは
$\det K_F=0$ で定義される。
従って第\ref{sec:tate-boundary-tenth}節の境界模型は、
収束する代数的冪級数による模型である。
\end{corollary}
\begin{proof}
Incidenceのminorによる表示は、代数的な評価写像に対しても同じである。
指定点の局所環を $R$ とする。
$R\to\widehat R$ はfaithfully flatである
\cite[Lemma 10.97.3, Tag 00MC]{StacksFaithfulCompletion}。
またincidenceの射はproperなので、そのscheme imageを定める
$\OO\to p_*\OO_{\mathcal I}$ の核は、このflat base changeと可換である
\cite[Lemma 30.5.2, Tag 02KH]{StacksFlatBaseChange}。
完備化後のイデアルは前節により $(\det K_F)$ なので、
faithful flatnessにより $R$ 上でも同じである。
最後にイデアル層のcoherenceから、この等式は指定点の近傍に広がる。
\end{proof}

この議論で使ったbase changeは局所環のflatな完備化である。
properなscheme imageが任意の三次式への特殊化と可換だという仮定は使っていない。
また、十次項を表示するための有限の単元による規格化は収束性を保つが、
任意の形式的な座標変更の収束性まで主張するものではない。

\begin{remark}[三次式を動かす場合]
有限方程式の構成自体は $F=0$ を含む係数空間全体で成立する。
普遍超曲面を扱う場合には非零三次式へ制限する。
その普遍超曲面はflatかつproperで、全空間は滑らかである。
実際、任意の射影点で非零になる三次単項式の係数について微分すれば、
普遍方程式の微分が非零になる。
$\mathcal Z\to\mathcal S$ も滑らかなので、
同じ表示は、この普遍族上でも零切断の近傍に代数的な $\mathcal M$ と
$\mathcal E$ を与える。
上の境界の系は、滑らかな $F$ を固定した場合について述べた。
\end{remark}

\subsection{Tate定数blockの二次項をPieri則で決定する}
\label{subsec:v9-tate-second-block}

\begin{proposition}\label{prop:v9-K2-hessian}
前節のTate行列分解の定数blockを
$K_F(a)=-a+K_2(F,a)+K_3(F,a)+\cdots$ と書く。
ここで $K_n$ は $a$ に関して次数 $n$ である。
すると、選んだ符号付き補集合基底において
\[
                     K_2(F,a)=\frac13\operatorname{Hess}_F(q(a)).
\]
特に
\[
                    [\det K_F]_6=2F(q(a)).
\]
\end{proposition}

\begin{proof}
$K_n$ は $V\otimes V\otimes(\det V)^{-1}$ に値を持つ。
中心の重みから、その $F$ 次数は $n-1$ である。
従って二次項の空間は
\[
 \operatorname{Hom}_{GL(V)}
 \left(\operatorname{Sym}^3V\otimes
       \operatorname{Sym}^2(\Lambda^2V),
       (V\otimes V)\otimes\det V\right)
\]
である。
$\operatorname{Sym}^2(\Lambda^2V)=S_{(2,2)}V\oplus S_{(1,1,1,1)}V$ を用いる。
標的の対称成分と交代成分の分割は、それぞれ
$(3,1,1,1,1)$ と $(2,2,1,1,1)$ である。
Pieri則によれば、三つの水平な箱を加えて前者を得る方法は
$(1,1,1,1)$ からの一通りだけであり、後者を得る方法はない。
従ってこのHom空間は一次元で、対称行列に値を持つ。
$\operatorname{Hess}_F(q)$ はその非零元なので、係数を一例で決めればよい。

$F=\sum x_i^3$、$a=e_0\wedge e_1+e_2\wedge e_3$ とする。
このとき $q=e_4$、$A_1=\eta_2\eta_3\eta_4+\eta_0\eta_1\eta_4$ である。
$A_1^2$ の $\eta$ 次数4の部分は $-2x_4\eta_0\eta_1\eta_2\eta_3$ であり、
$A_2=-h(A_1^2)$ の対応する項は
$2\theta_4\eta_0\eta_1\eta_2\eta_3$ である。
従って $(K_2)_{44}=2$、一方 $(\operatorname{Hess}_F(q))_{44}=6$ なので、
比例係数は $1/3$ である。
最後に $\operatorname{adj}(-a)=q\,q^{\mathsf T}$ およびEulerの恒等式から
\[
 [\det K_F]_6=q^{\mathsf T}K_2q
 =\frac13q^{\mathsf T}\operatorname{Hess}_F(q)q=2F(q)
\]
を得る。
\end{proof}

\subsection{十二次までの実計算と座標変更の作用}
\label{subsec:v9-twelfth-jet}

Tate再帰をblock次数8まで延ばすと、行列式を次数12まで厳密に計算できる。
各次数の多項式Maurer--Cartan方程式を有理数演算で検証し、
さらに実際の外積・縮約作用による $32$ 次元の作用素表示でも、
次数0から8までの行列分解方程式を独立に検証した。
以下は、定義式を $\frac12\det K_F$ に固定したときの値である。
十二次の数値を新しい不変量とは解釈しない。

$u_i=e_i-p_i e_0$ とし、$F_0=\sum_{i=0}^4x_i^3$、
$F_*=\sum_{i=0}^3x_i^3-4x_4^3$ と書く。次の三つの双ベクトルを用いる。
\[
 \begin{aligned}
 a^{(1)}&=u_1\wedge u_2+u_3\wedge u_4,\\
 a^{(2)}&=a^{(1)}+u_1\wedge u_3,\\
 a^{(3)}&=2u_1\wedge u_2+\tfrac12u_3\wedge u_4.
 \end{aligned}
\]
計算結果は次の通りである。表の係数は $\det K_F$ の係数であり、
正規化した方程式 $\frac12\det K_F$ の係数はその半分である。
\[
\begin{array}{c|c|c|r|r}
 F &p&a&[t^{10}]\det K_F(ta)&[t^{12}]\det K_F(ta)\\ \hline
 F_*&(1,1,1,1,1)&a^{(1)}&-648&-632\\
 F_*&(1,1,1,1,1)&a^{(2)}&-648&-4504/3\\
 F_*&(1,1,1,1,1)&a^{(3)}&-648&-3318\\
 F_0&(1,5,-7,-8,9)&a^{(1)}&408240&5295624\\
 F_0&(1,2,3,4,5)&a^{(1)}&-142560&-71472596/3
\end{array}
\]
各行で $q(a)=p$ である。最初の四行では $F(q)=0$、
最後の行では $F(q)=225$ である。
七次、八次、九次、十一次の行列式係数は全て零である。

最初の標本では $K_3=K_5=0$ だが、$K_7$ は零ではない。
$A=K_1$、$B=K_2$、$C=K_4$ と書くと、十二次に寄与する
五つのblockの組合せは
\[
\begin{array}{c|r}
 \text{blockの組合せ}&\text{寄与}\\ \hline
 A^4K_8&-285394/75\\
 A^3BK_7&12274/75\\
 A^2B^2K_6&78424/15\\
 A^2BC^2&-6656/3\\
 B^4C&0
\end{array}
\]
であり、その和は $-632$ となる。
ここで各記号は、行列式の対応する混合係数全体を表す。

\begin{proposition}[十二次の標本値は座標不変ではない]
\label{prop:v9-twelfth-gauge}
$P_F(a)=F(q(a))$ とおく。
七次の普遍座標変更
\[
 \phi_7(F,a)=\frac12\bigl(K_7(F,a)-K_7(F,a)^{\mathsf T}\bigr)
\]
は $W=\Lambda^2V\otimes(\det V)^{-1}$ に値を持ち、
$F$ 次数6、$a$ 次数7の正則な共変式である。
最初の標本 $(F_*,a_*)$ では
\[
                    (dP_{F_*})_{a_*}\phi_7=-\frac{6137}{50}\ne0.
\]
従って十二次係数の値 $-316=[t^{12}]\frac12\det K_{F_*}(ta_*)$ は、
十一次以下の既定の正規化を保つ普遍座標変更で変わる。
\end{proposition}

\begin{proof}
$K_7$ の値域と中心の重みから、$\phi_7$ が所定の普遍座標変更であることが分かる。
Pfaffianは二次なので、その微分は
\[
 (dq)_a(\phi_7)=\frac12\bigl(q(a+\phi_7)-q(a-\phi_7)\bigr)
\]
で厳密に計算できる。
最初の標本ではその値が
\[
 \left(0,-\frac{77221}{10800},
 \frac{224509}{10800},-\frac{6137}{225},\frac{6137}{900}\right)
\]
となる。これを $dF_*(q)$ と組み合わせると、主張の $-6137/50$ を得る。

実際、pullback $d'(a)=d(a+\lambda\phi_7(a))$ は次数12で
\[
                       d'_{12}=d_{12}+\lambda\,dP(\phi_7)
\]
を満たし、次数11以下を変えない。
$\lambda=-15800/6137$ とすれば、この標本での十二次係数だけは零にできる。
これは普遍な十二次項全体を消去したという主張ではない。
\end{proof}

普遍な正則 $GL(V)$-共変式の範囲では、十二次で調べるべき対象は
\[
 \mathcal C_{12}\Big/
 \left(dP(\mathcal G_7)+P\,\mathcal U_6\right)
\]
である。既計算の次元は
$\dim\mathcal C_{12}=5$、$\dim\mathcal G_7=7$、
$\dim\mathcal U_6=2$ だが、これらの次元だけから商の次元は決まらない。
四次の単元による十次項の正規化は、八次項が零であるため十二次項を変えない。
また低次の座標変更の非線形な交差項は、次数12には現れない。
従って上の商への射影を計算することが、十二次項を不変量として論じるための
次の必要な段階である。


\subsection{12次のcontact商について確定する上下界}
\label{subsec:degree-twelve-contact-bounds}

11次までの正規形を固定した後、12次に作用する普遍同変な変更は
\[
 dP_F(\phi_7)+u_6P_F,
 \qquad \phi_7\in\mathcal G_7,\quad u_6\in\mathcal U_6
\]
である。$\dim\mathcal G_7=7$、$\dim\mathcal U_6=2$、
$\dim\mathcal C_{12}=5$ だが、これらを差し引くことでは商の次元は決まらない。

ここで低次の変更が追加の項を生じないことも確認しておく。
既知の計算により $\mathcal G_2=\mathcal G_3=\mathcal G_4=0$ であり、
5次座標変更の $P_F$ への作用は4次単元の像に入る。
この作用を相殺して10次の正規形を保つと、$H_{10}$ への作用は14次以上になる。
6次座標変更の $P_F$ への作用は $\mathcal C_{11}=0$ により零である。
また二つの5次座標変更の交差項も14次から始まる。
従って12次に残る作用は上式の線形写像で尽くされる。

\begin{proposition}[Pfaffian微分の像]
\label{prop:degree-twelve-pfaffian-image}
$A=\Lambda^2V$ と書き、定数直線因子を戻した表示で
\[
 \mathcal D_8=\operatorname{Hom}_{\operatorname{GL}(V)}
 \left(\operatorname{Sym}^6(\operatorname{Sym}^3V)
       \otimes\operatorname{Sym}^8A,
       \Lambda^4V\otimes\det(V)^6\right)
\]
とおく。この空間は4次元であり、
\[
 \delta q\colon\mathcal G_7\longrightarrow\mathcal D_8,
 \qquad\phi\longmapsto a\wedge\phi(a)
\]
は全射である。従ってその核は3次元である。
\end{proposition}
\begin{proof}
双対を取り、共通の $\det(V)^7$ を掛けると、係数写像の双対は
\[
 \mathsf d\colon V\otimes\operatorname{Sym}^8A
       \longrightarrow\Lambda^3V\otimes\operatorname{Sym}^7A,
 \qquad
 v\otimes f\longmapsto
 \sum_{i<j}(v\wedge e_i\wedge e_j)\otimes
                \frac{\partial f}{\partial a_{ij}}
\]
となる。非零定数倍は以下の階数に影響しない。
Littlewoodの分解、Pieriの公式、および前節までと同じ整数alternant計算により、
必要な成分は次の表で尽くされる。ここで $\lambda$ は
$\Sym^6(\Sym^3V)$ の分割、$\mu=(7-\lambda_5,\ldots,7-\lambda_1)$ であり、
後ろの二列は対応する重複度である。
\[
\begin{array}{c|c|c|c}
 \lambda&\mu&\mathcal D_8\text{側}&\mathcal G_7\text{側}\\ \hline
 (6,6,6,0,0)&(7,7,1,1,1)&1&2\\
 (7,5,4,1,1)&(6,6,3,2,0)&1&1\\
 (7,4,4,2,1)&(6,5,3,3,0)&1&1\\
 (6,4,4,2,2)&(5,5,3,3,1)&1&2\\
 (6,5,4,2,1)&(6,5,3,2,1)&0&1
\end{array}
\]
表中の $\lambda$ はplethysmにそれぞれ重複度1で現れる。

第一行と第四行は $\mathcal U_6$ の二つの成分である。
$\phi=a u_6$ に対し $\delta q(\phi)=2q u_6$ なので、
この二成分で写像は非零である。

残る二成分については最高重みベクトルを直接使う。
\[
 Q=a_{01}a_{23}-a_{02}a_{13}+a_{03}a_{12},
 \qquad w=e_0a_{12}-e_1a_{02}+e_2a_{01}
\]
とおくと、$e_0a_{01}^2Q^3$ は重み $(6,5,3,3,0)$ の最高重みベクトルであり、
その $\mathsf d$ による像の $e_0\wedge e_2\wedge e_3$ 係数は
$3a_{01}^3Q^2\ne0$ である。
また $w a_{01}^3Q^2$ は重み $(6,6,3,2,0)$ の最高重みベクトルで、
\[
 \mathsf d(w a_{01}^3Q^2)
       =8(e_0\wedge e_1\wedge e_2)a_{01}^3Q^2\ne0.
\]
従って必要な四つの既約成分のすべてで双対写像が単射となり、主張が従う。
\end{proof}

\begin{corollary}[12次商の厳密な上界と下界]
\label{cor:degree-twelve-contact-bounds}
\[
 \mathcal Q_{12}=
 \mathcal C_{12}/\bigl(dP_F(\mathcal G_7)+P_F\mathcal U_6\bigr)
\]
とおく。これは $F$ も変数に含む普遍共変式の商である。
このとき
\[
 \boxed{
 3\le\operatorname{rank}(dP_F\colon\mathcal G_7\to\mathcal C_{12})\le4,
 \qquad 1\le\dim\mathcal Q_{12}\le2.
 }
\]
\end{corollary}
\begin{proof}
$dP_F$ は $\delta q$ と $dF$ の合成なので、命題より階数は4以下である。
Eulerの等式により
\[
 dP_F\left(\frac16a u_6\right)=P_Fu_6.
\]
従って単元の像はすでに座標変更の像に含まれ、その次元は2である。

さらにTate再帰の定数ブロック $K_7$ から得られる
$\phi_7=\operatorname{Alt}(K_7)$ は $\mathcal G_7$ の元である。
$\operatorname{Hom}(\Lambda^4V,V)$ の交代化は、
$\Lambda^4V\simeq V^*\otimes\det V$ を使った自然な
$V\otimes V\otimes\det(V)^{-1}\to
\Lambda^2V\otimes\det(V)^{-1}$ である。
従ってこれは座標依存の行列操作ではない。

命題\ref{prop:v9-twelfth-gauge}の標本 $(F_*,a_*)$ では
\[
 P_{F_*}(a_*)=0,\qquad
 dP_{F_*}(\phi_7)(a_*)=-\frac{6137}{50}\ne0
\]
である。従ってこの像は $P_F\mathcal U_6$ に含まれず、階数は3以上である。
この最後の値は保存された $K_7$ の有理数行列から交代化とPfaffianの微分を
別に計算し直すことで検証できる。
\end{proof}

\begin{remark}[純粋な$q$の項と残る計算]
12次共変量の $C$ に沿う位数は、Littlewood分解の
$(r,r,s,s,0)$ における $s$ で読める。
実際、最高重みベクトル $a_{01}^{r-s}Q^s$ は
$I_C^s$ に属し、一般の階数2点でちょうど $s$ 次で消える。
12次の成分は、$s=3,4,5$ にそれぞれ次元 $2,1,2$ で現れる。
$s=6$ の純粋な $q$ に6次の成分は、
$\mathbb S_{(9,3,3,3,3)}V$ のplethysm重複度が零なので存在しない。

一段階吹上げで $r^3$ を割った後は、束方向6次の
$\mathbb S_{6600}Q^*$、$\mathbb S_{5511}Q^*$、$\mathbb S_{4422}Q^*$
の三成分が残り、最後の $\OO_X(3)$ 成分は現れない。
従って $\operatorname{Jac}(F)_6=0$ だけでは12次の全成分は消去できない。
上の上下界は、4次元の $\mathcal D_8$ から5次元の $\mathcal C_{12}$ への
$dF$ の階数を完全に決めたものではない。
また、この段階では実際の12次項の $\mathcal Q_{12}$ における類も未決定である。
これら二点は次節で解決する。
\end{remark}



\clearpage
\section{12次・13次の普遍接触商と境界の標準形}
\label{sec:contact-twelve-thirteen}

本節では、次数12の普遍接触商を完全に決定し、実際の境界の類が零であることを示す。
さらに次数13の普遍接触商も零であり、境界は
$P_F+H_{10}/96+O(a^{14})$ に正規化される。
ここで $P_F=F(q(a))$、$H_{10}=\det\operatorname{Hess}_F(q(a))$ である。
候補の空間、座標変更の像、実際の係数の三つを順に決定する。
共変量の次元計算は、前節までと同じLittlewood分解とSchur指標の
有限計算\cite{BeteaWheeler,DeBoeckPagetWildon}に基づく。
有限個の評価値が普遍的な結論を与える理由も、各段階で明示する。

\subsection{12次contact写像の階数の確定}
\label{subsec:v10-contact-exact-rank}

前節の上界4は等号である。これを証明するために、Tate再帰の
一つの部分項を独立の座標変換として使う。
$\kappa$ を $x=0$ における $\Lambda^4V\to V$ ブロックの抽出とし、
$\operatorname{Alt}$ は
$\operatorname{Hom}(\Lambda^4V,V)\simeq
V\otimes V\otimes\det(V)^{-1}$ の自然な交代化とする。
\[
 \Phi_7=\operatorname{Alt}(K_7),\qquad
 \Psi_7=\operatorname{Alt}\,\kappa
       \bigl(h(A_1\,h(A_3^2))\bigr).
\]
後者は、$A_6$ の部分項 $-h(A_3^2)$ を
$A_7$ の項 $-h(A_1A_6)$ に代入して得られる。

\begin{lemma}[部分項の同変性と次数]
$\Psi_7$ は $\mathcal G_7$ の元である。
\end{lemma}
\begin{proof}
$A_1$、Chevalley反対称化による正規表示、積、$h$、$\kappa$、
交代化はすべて普遍的な $\operatorname{GL}(V)$ 同変操作である。
各操作は $F$ の係数に多項式的であり、$a$ の次数は葉の数に等しい。
従って表示された部分項は $a$ に7次である。

$F$ の次数は次の補助重みで確認できる：
\[
 \mathrm{wt}(F)=1,\quad\mathrm{wt}(a)=-1,\quad
 \mathrm{wt}(x)=0,\quad
 \mathrm{wt}(\theta)=-\tfrac12,\quad
 \mathrm{wt}(\eta)=\tfrac12.
\]
Clifford関係は斉次で、$h$ の重みは $-1/2$、$A_1$ の重みは
$1/2$ である。従って二項操作 $-h(UV)$ で作られる各部分項の重みは
$1/2$ になる。$x=0$ の $\Lambda^4V\to V$ ブロックでは
$\eta$ の個数から $\theta$ の個数を引いたものが3なので、
$a$ に $n$ 次の部分項の係数は $F$ に $n-1$ 次である。
ここでは $n=7$ であり、所要の $F$ に6次、$a$ に7次という次数を得る。
部分項自体がMaurer--Cartan方程式を満たす必要はない。
\end{proof}

\begin{proposition}[12次contact商は1次元]
\label{prop:v10-contact-exact-rank}
普遍正則な同変写像について
\[
 \operatorname{rank}(dF\colon\mathcal D_8\to\mathcal C_{12})=4,
 \qquad
 \operatorname{rank}(dP_F\colon\mathcal G_7\to\mathcal C_{12})=4.
\]
従って
\[
 \dim\mathcal Q_{12}=1,\qquad
 \ker(dP_F)=\ker(\delta q),\qquad\dim\ker(dP_F)=3.
\]
\end{proposition}
\begin{proof}
$F_*=x_0^3+x_1^3+x_2^3+x_3^3-4x_4^3$ とし、
$u_i=e_i-e_0$ に対して
\[
 a^{(1)}=u_1\wedge u_2+u_3\wedge u_4,
 \qquad a^{(2)}=a^{(1)}+u_1\wedge u_3
\]
とおく。両方で $q(a^{(j)})=(1,1,1,1,1)$ なので $P_{F_*}=0$ である。
有理数上のTate演算は
\[
\begin{array}{c|cc}
 &dP_{F_*}(\Phi_7)&dP_{F_*}(\Psi_7)\\\hline
 a^{(1)}&-6137/50&0\\
 a^{(2)}&-166211/600&-48
\end{array}
\]
を与える。第二列の値は、より小さい計算
\[
 dq(\Psi_7)(a^{(1)})=(0,40/3,-40/3,0,0),
 \qquad
 dq(\Psi_7)(a^{(2)})=(-40/3,0,32/3,-40/3,0)
\]
と $dF_*(1,1,1,1,1)=(3,3,3,3,-12)$ からも確認できる。
評価行列の行列式は
\[
 \frac{147288}{25}\ne0
\]
である。従って $dP_F(\Phi_7)$ と $dP_F(\Psi_7)$ は
$P_F\mathcal U_6$ を法として一次独立である。
$P_F\mathcal U_6$ 自体は2次元で、Eulerの等式によりすでに
$dP_F(\mathcal G_7)$ に含まれるため、全体の階数は4以上である。
命題\ref{prop:degree-twelve-pfaffian-image}による上界4と合わせて等号を得る。
同命題の $\delta q$ の全射性により、$dF$ は単射である。
\end{proof}

ここで示したのは $F$ を変数とする普遍正則な同変写像の階数であり、
固定した個々の三次式について任意の非同変座標変換を許した
Jacobian商の次元を主張するものではない。


\subsection{12次商の生成元と実際の12次項の消滅}
\label{subsec:v10-actual-twelfth-class}

次元が1であることと、実際の項が零であることは別の主張である。
ここでは商の非零生成元と、その係数を読み取る線形汎関数を与えたうえで、
実際の12次項を判定する。

まず $d^{\mathrm T}=\tfrac12\det K$ をTate座標での未正規化の方程式とし、
$d^{\mathrm T}_{10}$、$d^{\mathrm T}_{12}$ をその斉次成分とする。
定理\ref{thm:tate-tenth-coefficient}と $\dim\mathcal C_{10}=2$ により
\begin{equation}\label{eq:v10-normalized-U4}
 U_4=\frac{H_{10}/96-d^{\mathrm T}_{10}}{3P_F}
       \in\mathcal U_4
\end{equation}
は多項式として定義される。分子の $P_F$ による可除性は既に証明した
普遍的な $\mathcal C_{10}$ の分解から従い、点ごとの除算を意味しない。

\begin{lemma}[四次単元の対角三次式への制限]
\label{lem:v10-U4-diagonal}
$F=\sum_i c_i x_i^3$ に対して、\eqref{eq:v10-normalized-U4}の正規化では
\[
 U_4(F,a)=\sum_{i=0}^4\prod_{j\ne i}(c_j a_{ij}).
\]
ここで $a_{ji}=-a_{ij}$ とする。
\end{lemma}
\begin{proof}
対角トーラスの重みを用いる。$c^\alpha a^\beta$ は $c$ と $a$ の
両方に4次であり、頂点 $i$ の $a$ の次数を $\deg_i(\beta)$ とすると
\[
 3\alpha_i+\deg_i(\beta)=4.
\]
従って $\alpha_i$ は0または1で、ちょうど一つだけ0である。
その頂点の次数は4、残る四頂点の次数は1だから、対応する単項式は
その頂点を中心とする星型に限られる。置換対称性により
五つの行積 $\prod_{j\ne i}(c_j a_{ij})$ の係数は等しい。

係数の正規化には既存のFermatの錐外の計算を使う。
$q=(1,2,3,4,5)$、$a=u_1\wedge u_2+u_3\wedge u_4$、
$u_i=e_i-q_i e_0$ において
\[
 P_F=225,\qquad H_{10}/96=9720,\qquad
 d^{\mathrm T}_{10}=-71280.
\]
従って\eqref{eq:v10-normalized-U4}の値は120である。
上の星型の和も同じ点で120なので、係数は1である。
\end{proof}

\begin{theorem}[12次商の生成元と係数汎関数]
\label{thm:v10-twelfth-generator}
$\Theta_8$ を\eqref{eq:boundary-theta-eight}で定義した八次共変量とする。
普遍共変量空間は
\[
 \mathcal C_{12}
 =P_F\mathcal U_6
   \oplus\CC\,dP_F(\Phi_7)
   \oplus\CC\,dP_F(\Psi_7)
   \oplus\CC\,\Theta_8U_4
\]
と直和分解される。特に $[\Theta_8U_4]$ は $\mathcal Q_{12}$ の基底である。

$F_*=x_0^3+x_1^3+x_2^3+x_3^3-4x_4^3$、$u_i=e_i-e_0$ とし、
\[
 a^{(1)}=u_1\wedge u_2+u_3\wedge u_4,\qquad
 a^{(3)}=2u_1\wedge u_2+\tfrac12u_3\wedge u_4
\]
とおくと、商における $[\Theta_8U_4]$ の係数は
\begin{equation}\label{eq:v10-rho-functional}
 \rho(f)=\frac{21f(F_*,a^{(1)})-4f(F_*,a^{(3)})}{400}
\end{equation}
によって読み取れる。
\end{theorem}
\begin{proof}
三つ目の点 $a^{(2)}=a^{(1)}+u_1\wedge u_3$ も含めて値を並べると、
全ての点で $P_{F_*}=0$ であり、厳密計算は
\[
\begin{array}{c|rrr}
 &a^{(1)}&a^{(2)}&a^{(3)}\\\hline
 dP(\Phi_7)&-6137/50&-166211/600&-128877/200\\
 dP(\Psi_7)&0&-48&0\\
 U_4&-4&0&-4\\
 \Theta_8&-14&-26&-97/2\\
 \Theta_8U_4&56&0&194
\end{array}
\]
を与える。$U_4$ の行は補題\ref{lem:v10-U4-diagonal}、$\Theta_8$ の行は
\[
 \Theta_8(F,a)=\sum_{i<j}a_{ij}^2\prod_{k\ne i,j}c_kq_k(a)
\]
から直接得られる。$\rho$ は $P_F\mathcal U_6$ を消し、
表示された二つの座標変更の像も消す。一方、
\[
 \rho(\Theta_8U_4)=\frac{21\cdot56-4\cdot194}{400}=1.
\]
命題\ref{prop:v10-contact-exact-rank}から接触写像の像は4次元であり、
$P_F\mathcal U_6$ と二つの表示された像で張られる。
従って $\rho$ はその像を消す非零の汎関数であり、
$\dim\mathcal C_{12}=5$ と合わせて直和分解と係数公式を得る。
\end{proof}

\begin{theorem}[実際の12次接触類は零]
\label{thm:v10-actual-twelfth-zero}
実際のTate族について
\[
 [d^{\mathrm T}_{12}]=0\quad\text{in }\mathcal Q_{12}.
\]
従って正則な普遍同変座標変換と単元変更により、境界方程式を
\begin{equation}\label{eq:v10-twelve-normal-form}
 d(F,a)=P_F(a)+\frac1{96}H_{10}(F,a)+O(a^{13})
\end{equation}
とできる。
\end{theorem}
\begin{proof}
前節の厳密な行列因子分解の計算から
\[
 \bigl(d^{\mathrm T}_{12}(F_*,a^{(1)}),
       d^{\mathrm T}_{12}(F_*,a^{(2)}),
       d^{\mathrm T}_{12}(F_*,a^{(3)})\bigr)
 =\left(-316,-\frac{2252}{3},-1659\right).
\]
従って
\[
 \rho(d^{\mathrm T}_{12})
 =\frac{21(-316)-4(-1659)}{400}=0.
\]
命題\ref{thm:v10-twelfth-generator}により $\rho$ は1次元商の係数を完全に判定するから、接触類は零である。
これは個々の評価値だけを根拠とする多項式恒等式の推測ではなく、
共変量空間と接触写像の像を先に決定したうえでの係数計算である。

変更も有限な普遍式で指定できる。
\[
 \alpha=\frac{15800}{6137},\qquad
 \beta=\frac{689855}{883728},\qquad
 V_6=\frac{d^{\mathrm T}_{12}
       -\alpha\,dP_F(\Phi_7)-\beta\,dP_F(\Psi_7)}{P_F}.
\]
最初の二つの評価から $\alpha,\beta$ が決まり、上の直和分解により
$V_6\in\mathcal U_6$ は多項式である。
まず $1+3U_4$ を掛けて10次の $P_F$ の倍数を消す。
$d^{\mathrm T}_8=0$ なので、この変更は12次項を変えない。
次に
\[
 a\longmapsto a-\alpha\Phi_7(F,a)-\beta\Psi_7(F,a)
\]
で引き戻し、$1-V_6$ を掛けると12次項が消える。
この座標変更が10次の $H_{10}$ に作用するのは16次以降である。
全ての変更は有限な多項式で、座標変更の線形項は恒等、単元の定数項は1である。
従って前節で構成した代数的エタール近傍・収束する複素解析的近傍にも適用できる。
\end{proof}

12次で新しい普遍候補が存在することと、実際の族がその方向に動くことは異なる。
ここでは $[\Theta_8U_4]\ne0$ である一方、実際の係数は零である。
この消去は10次のHessian障害を取り除かず、形式的非斉次性についての
前節までの結論を変えない。また、\eqref{eq:v10-twelve-normal-form}は
13次以上の項の消滅を主張しない。


\subsection{13次の普遍項はすべて除去できる}
\label{subsec:v10-degree-thirteen}

まず実際の13次係数を計算することなく、13次項の除去可能性を確定する。
続いて実際の係数も計算し、消去に使う普遍定数を明示する。
\begin{proposition}[13次contact商の消滅]
\label{prop:v10-degree-thirteen}
\[
 \dim\mathcal C_{13}=1,\qquad
 \dim\mathcal G_8=7,\qquad
 \dim\mathcal U_7=0,
\]
かつ $dP_F\colon\mathcal G_8\to\mathcal C_{13}$ は全射である。
従って13次の普遍正則なcontact商は零である。
\end{proposition}
\begin{proof}
前節までと同じ整数alternant計算により、
$\operatorname{Sym}^8(\operatorname{Sym}^3V)$ における候補
\[
 (10,7,7,0,0),\quad(10,6,6,1,1),\quad
 (10,5,5,2,2),\quad(10,4,4,3,3)
\]
の重複度は順に $0,0,1,0$ である。
また $\mathcal G_8$ と $\mathcal U_7$ の重複度計算はそれぞれ
\[
 [(8,8,7,7,7)]\bigl(
 \operatorname{Sym}^7(\operatorname{Sym}^3V)
 \otimes\operatorname{Sym}^8(\Lambda^2V)\bigr)=7,
\]
\[
 [(7,7,7,7,7)]\bigl(
 \operatorname{Sym}^7(\operatorname{Sym}^3V)
 \otimes\operatorname{Sym}^7(\Lambda^2V)\bigr)=0
\]
となる。

Tate族の既に計算された定数ブロックから
\[
 \Phi_8=\operatorname{Alt}(K_8)\in\mathcal G_8
\]
を取る。前の補助重みと同変性の議論により、これは普遍正則な
$F$ に7次、$a$ に8次の座標変換である。
先の $F_*$ と $a^{(1)}$ において、保存済みの有理数行列から直接
\[
 dq(\Phi_8)=
 \left(\frac{20}{3},\frac{101839}{8640},
 -\frac{85889}{8640},-\frac{25801}{1440},
 \frac{25801}{1440}\right)
\]
を得る。従って
\[
 \boxed{dP_{F_*}(\Phi_8)(a^{(1)})=-\frac{2189}{9}\ne0.}
\]
標的が1次元なので $dP_F$ は全射となり、核は6次元である。
\end{proof}

12次項を消した後の方程式を
\[
 d=P_F+\frac1{96}H_{10}+d_{13}+O(a^{14})
\]
と書くと、命題より $d_{13}=\lambda\,dP_F(\Phi_8)$ と書ける。
座標変更 $a\mapsto a-\lambda\Phi_8(a)$ は13次項を消し、
12次以下を変えない。従って
\[
 \boxed{d\sim P_F+\frac1{96}H_{10}+O(a^{14}).}
\]
ここで示したのは13次項の除去可能性であり、
元のTate座標における13次係数自体が零であるとは主張していない。


\begin{proposition}[実際の13次係数と有限な正規化]
\label{prop:v10-explicit-thirteenth}
未正規化のTate方程式 $d^{\mathrm T}=\tfrac12\det K$ について、普遍恒等式
\begin{equation}\label{eq:v10-actual-thirteenth}
 d^{\mathrm T}_{13}=-\frac{4600}{2189}\,dP_F(\Phi_8)
\end{equation}
が成立する。前に定義した $U_4,V_6,\alpha,\beta$ を用いて
\[
 \chi(F,a)=a-\alpha\Phi_7(F,a)-\beta\Psi_7(F,a)
                         +\frac{4600}{2189}\Phi_8(F,a),
 \qquad u(F,a)=1+3U_4(F,a)-V_6(F,a)
\]
とおけば
\begin{equation}\label{eq:v10-final-normal-form}
 \boxed{
 u(F,a)\,d^{\mathrm T}(F,\chi(F,a))
   =P_F(a)+\frac1{96}H_{10}(F,a)+O(a^{14}).
 }
\end{equation}
ここで $\chi$ と $u$ は $F,a$ に多項式的であり、
$d_a\chi|_{a=0}=\id$、$u(F,0)=1$ である。
\end{proposition}
\begin{proof}
標準点 $(F_*,a^{(1)})$ でTate再帰を $A_9$ まで計算する。
Clifford正規形の有理多項式演算により、9次までのMaurer--Cartan等式を
確認し、$K_1,\ldots,K_8$ は前節の保存値と全成分で一致する。
新しい $A_9$ の項数は5766であり、その定数ブロックを追加した行列式は
\[
 [t^{13}]\det K(F_*,t a^{(1)})=\frac{9200}{9},\qquad
 d^{\mathrm T}_{13}(F_*,a^{(1)})=\frac{4600}{9}
\]
となる。次数13の行列式に寄与するブロックは $K_1,\ldots,K_9$ で尽くされる。
これらの有理数行列と係数は添付の証明書にすべて保存し、
別実装の32次元外積・縮約作用素でも次数0--9の恒等式を照合する。

$\mathcal C_{13}$ は1次元であり、$dP_F(\Phi_8)$ はその非零生成元である。
従って一つの値の比
\[
 \frac{4600/9}{-2189/9}=-\frac{4600}{2189}
\]
により\eqref{eq:v10-actual-thirteenth}が普遍的に決まる。

12次までの変更は13次項を変えない。
実際、4次単元が13次に与え得る項は $3U_4d^{\mathrm T}_9=0$、
6次単元からは $V_6d^{\mathrm T}_7=0$ である。
7次座標変更が13次に与え得る項も $d(d^{\mathrm T}_7)(\phi_7)=0$ であり、
$P_F$ への二重作用は18次からである。
従って $+(4600/2189)\Phi_8$ で引き戻せば13次項が消える。
二つの単元の積の10次以上の部分や、単元を座標変更で引き戻す差は、
方程式の14次未満には寄与しないので、表示した $u$ を用いてよい。
これで\eqref{eq:v10-final-normal-form}を得る。

線形項と定数項の条件により、固定した $F$ に対して $\chi$ は原点でエタールであり、
複素解析的にも原点の近傍で可逆になる。$u$ も同じ近傍で単元である。
従ってこれは形式的な等式に加え、前節で得た代数的・収束する解析的族に
適用できる有限な正規化である。大域的な多項式自己同型であるとは仮定しない。
\end{proof}

これまでの結果は次のように整理できる。各行では下の次数の正規形を固定し、
普遍正則な同変変更を用いている。
\[
\begin{array}{c|c|c|c|c}
 a\text{次数}&\dim\mathcal C_n&\text{接触像の次元}&\text{商の次元}
                                      &\text{実際の類}\\\hline
 8&1&0&1&0\\
 10&2&1&1&[H_{10}]/96\ne0\\
 12&5&4&1&0\\
 13&1&1&0&0
\end{array}
\]
10次の非零性は、全ての滑らかな $F$ で通常のJacobian商にも残るという
前節までの主張を含む。12次・13次の消去はこの障害を変えない。


\begin{remark}[普遍的消去と固定した三次式での消去]
\label{rem:v10-universal-fixed-scope}
上の接触商は、$F$ を変数として多項式的に定義された
$\operatorname{GL}(V)$ 共変量の商である。
固定した $F$ に対する任意の座標変更を許した商とは区別する必要がある。
普遍商の非零元が、すべての特殊化でも非零であるとは限らない。
一方、普遍的な座標変更 $\chi(F,a)$ と単元 $u(F,a)$ による
恒等式\eqref{eq:v10-final-normal-form}は、各 $F$ に特殊化してそのまま使える。
この方向には、固定した $F$ の全自己同型や全接触変換の分類は不要である。

固定した $F$ には、普遍的な共変量空間からは見えない低次の安定化群もある。
その具体例として、体積を固定して $\beta\in\Lambda^2V^*$ を任意に取り、
$g=\nabla F(q(a))$、$H=\operatorname{Hess}F(q(a))$ とおくと、
\[
 \phi_{\beta,3}(a)=\frac12\bigl(H\beta a-a\beta H\bigr),
 \qquad \delta q(\phi_{\beta,3})=\beta g,
 \qquad dP_F(\phi_{\beta,3})=0.
\]
最初の等式は行列の縮約を表し、$\phi_{\beta,3}$ は $a$ に3次である。
Pfaffianの共変性から得られる恒等式
\[
 \delta q(Ma+aM^{\mathsf T})
    =\operatorname{tr}(M)q-M^{\mathsf T}q
\]
に $M=H\beta/2$ を代入すれば、上の主張が従う。

さらに $h(p)=\det\operatorname{Hess}F(p)$ とおけば、$h$ は $p$ に5次なので、
\[
 \psi_{\beta,7}(a)
   =\frac14\bigl((\operatorname{Hess}h)(q)\beta a
                       -a\beta(\operatorname{Hess}h)(q)\bigr)
\]
は $a$ に7次であり、
\[
 \delta q(\psi_{\beta,7})=\beta\nabla h(q),\qquad
 d(h\circ q)(\phi_{\beta,3})=-dP_F(\psi_{\beta,7}).
\]
従って、このHamilton型の低次安定化群が10次のHessian項に与える12次の変化は、
通常のJacobianイデアルにすでに含まれる。
この計算は、固定した $F$ の安定化群を無視してよいと主張するものではなく、
上記の明示的な族について作用を計算したものである。
\end{remark}



\clearpage
\section{14次の接触商と実際の項の消去}
\label{sec:contact-fourteen}

第10版の正規形の次の次数を決定する。本節の結論は
\[
 \mathcal Q^{\mathrm{eff}}_{14}=\mathbb C\cdot[\Theta_8J_6],\qquad
 [R_{14}]=0,\qquad
 d\sim P_F+\frac1{96}H_{10}+O(a^{15})
\]
である。ここで $J_6$ は以下で定義する六次単元共変量であり、最後の式は
各滑らかな $F$ の実際の境界方程式に適用される。

共変量と座標変更は、これまでと同じく $F$ の係数に多項式的に正則で、
指定した双次数を持つ普遍 $\mathrm{GL}(V)$ 同変なものに限る。
接触変換では $F$ を固定し、$a$ の線形項と方程式の先頭係数を恒等に固定する。
以下で用いる $F$ の方向微分は共変量を判別するための線形汎関数である。

\subsection{14次の共変量とPfaffian微分}
\label{subsec:v11-fourteenth-characters}

$A=\Lambda^2V$ と書き、定数直線因子を戻した表示で
\[
 \mathcal D_{10}=\operatorname{Hom}_{\operatorname{GL}(V)}
 \left(\operatorname{Sym}^8(\operatorname{Sym}^3V)
       \otimes\operatorname{Sym}^{10}A,
       \Lambda^4V\otimes\det(V)^8\right)
\]
とおく。14次に関係する空間は、前の次数より大きくなる。

\begin{proposition}[14次の次元とPfaffian微分の全射性]
\label{prop:v11-fourteenth-characters}
\[
 \dim\mathcal C_{14}=15,\qquad
 \dim\mathcal G_9=28,\qquad
 \dim\mathcal U_8=7,\qquad
 \dim\mathcal D_{10}=17.
\]
さらに
\[
 \delta q:\mathcal G_9\longrightarrow\mathcal D_{10},
 \qquad \phi\longmapsto a\wedge\phi
\]
は全射であり、その核は11次元である。
\end{proposition}
\begin{proof}
Littlewood分解と、前節までと同じ整数alternantによる重複度計算を用いる。
$\mathcal C_{14}$ について必要な
$\operatorname{Sym}^9(\operatorname{Sym}^3V)$ のSchur重複度は、
$\lambda_j=(11,8-j,8-j,j,j)$、$0\le j\le4$ として
\[
\begin{array}{c|ccccc}
 j&0&1&2&3&4\\\hline
 \text{重複度}&2&2&7&2&2
\end{array}
\]
である。また $\mathcal U_8$ では
$\lambda=(8,8-j,8-j,j,j)$、$0\le j\le4$ の重複度が
$1,1,3,1,1$ となる。$\mathcal G_9$ と $\mathcal D_{10}$ はそれぞれ
\[
 [(9,9,8,8,8)]\bigl(
   \operatorname{Sym}^8(\operatorname{Sym}^3V)
   \otimes\operatorname{Sym}^{9}A\bigr)=28,
\]
\[
 [(9,9,9,9,8)]\bigl(
   \operatorname{Sym}^8(\operatorname{Sym}^3V)
   \otimes\operatorname{Sym}^{10}A\bigr)=17
\]
から得られる。

全射性を確認する。双対を取り、共通の $\det(V)^9$ を掛けると、
必要な写像は
\[
 \mathsf d:V\otimes\operatorname{Sym}^{10}A
       \longrightarrow\Lambda^3V\otimes\operatorname{Sym}^{9}A,
 \qquad
 v\otimes f\longmapsto
 \sum_{i<j}(v\wedge e_i\wedge e_j)\otimes
                         \frac{\partial f}{\partial a_{ij}}
\]
となる。左辺はPieri則により重複度なしで分解される。
必要な成分は、$r+s=10$ として
\[
 \mathrm A_{r,s}=(r+1,r,s,s,0),\qquad
 \mathrm B_{r,s}=(r,r,s+1,s,0),\qquad
 \mathrm C_{r,s}=(r,r,s,s,1)
\]
の三種である。

$\mathrm C_{r,s}$ は $s=1,\ldots,5$ に現れ、対応する
$\operatorname{Sym}^8(\operatorname{Sym}^3V)$ 側の重複度は
$1,1,3,1,1$ である。これらは7次元の単元方向に対応する。
$\delta q(a u_8)=2q u_8$ が多項式として単射であるため、この七方向を
すべて含む。

残る成分では
\[
 Q=a_{01}a_{23}-a_{02}a_{13}+a_{03}a_{12},\qquad
 w=e_0a_{12}-e_1a_{02}+e_2a_{01}
\]
を使う。$\mathrm A_{r,s}$ の最高重みベクトル
$e_0a_{01}^{r-s}Q^s$ の像の $e_0\wedge e_2\wedge e_3$ 係数は
\[
 s\,a_{01}^{r-s+1}Q^{s-1}
\]
であり、$\mathrm B_{r,s}$ の最高重みベクトルについては
\[
 \mathsf d\bigl(w a_{01}^{r-s-1}Q^s\bigr)
   =(r+2)(e_0\wedge e_1\wedge e_2)a_{01}^{r-s-1}Q^s
\]
である。必要な成分と重複度は
\[
\begin{array}{c|ccc|ccc}
 \text{型}&\mathrm A_{7,3}&\mathrm A_{6,4}&\mathrm A_{5,5}
           &\mathrm B_{8,2}&\mathrm B_{7,3}&\mathrm B_{6,4}\\\hline
 \text{$F$側の重複度}&2&1&1&1&4&1\\
 \text{上の非零係数}&3&4&5&10&9&8
\end{array}
\]
となる。従って残る10方向でも双対写像は単射であり、
合わせて17方向すべてで単射になる。これで元の写像の全射性と核の次元を得る。
\end{proof}

なお $dP_F=dF\circ\delta q$ であるが、上の全射性は $dF$ の階数を
決定していない。Eulerの等式により、7次元の $P_F\mathcal U_8$ は
すでに $dP_F(\mathcal G_9)$ に含まれる。

\subsection{三つの独立なKoszul共変量による14次階数の上界}
\label{subsec:v11-koszul-three-syzygies}

前節の $\delta q$ の全射性だけでは、
$dF:\mathcal D_{10}\to\mathcal C_{14}$ の階数は決まらない。
ここでは三つの多項式syzygyを具体的に作り、独立性を整数計算で示す。
この方法は、$q=0$ に沿ったKoszul複体の完全性を仮定しない。

\begin{proposition}[独立な三つのKoszul像]
\label{prop:v11-three-koszul-images}
\[
 \mathcal B_{14}=
 \operatorname{Hom}_{\operatorname{GL}(V)}
 \left(\operatorname{Sym}^7(\operatorname{Sym}^3V)
       \otimes\operatorname{Sym}^6(\Lambda^2V),
       \Lambda^2V^*\otimes\det(V)^7\right)
\]
とおく。写像
\[
 \mathcal B_{14}\longrightarrow\mathcal D_{10},
 \qquad \beta\longmapsto\beta\,\nabla F(q)
\]
の像は少なくとも三次元である。従って
\[
 \operatorname{rank}(dF:\mathcal D_{10}\to\mathcal C_{14})\le14.
\]
別途得られる階数14の下界と合わせれば、$dF$ および
$dP:\mathcal G_9\to\mathcal C_{14}$ の階数は正確に14となる。
\end{proposition}
\begin{proof}
体積を固定した成分表示で
\[
 c_{ijk}=\epsilon_{ijklm}a_{lm}\quad(l<m,\ \{l,m\}=\{0,\ldots,4\}
                                  \setminus\{i,j,k\})
\]
と書く。これは自然な同型
$\Lambda^2V\simeq\Lambda^3V^*\otimes\det(V)$ による $a$ の表示である。
七つの対称テンソル $F$ と六つの交代テンソル $c$ を縮約し、
三つの $F$ 側の添字を自由に残す。その出力を交代化して双対化すると、
\[
 \Lambda^3V\otimes\det(V)^6
       \simeq\Lambda^2V^*\otimes\det(V)^7
\]
に値を持つ、双次数 $(7,6)$ の共変量 $\beta$ を得る。
したがって $\beta\in\mathcal B_{14}$ である。

具体的な三つの縮約を次の表で指定する。左頂点 $0,\ldots,6$ は
それぞれ一つの $F$ を表す。各列の三つ組は、その $c$ 頂点に接続する
三つの $F$ 頂点を、$c$ の添字順で並べたものを表す。
最後の列は縮約しない三つの自由添字の順序である。
各 $F$ 頂点は、自由添字を含めてちょうど三回現れる。
\[
\begin{array}{c|cccccc|c}
 &c_0&c_1&c_2&c_3&c_4&c_5&\text{自由添字}\\\hline
 \Gamma_1&024&046&023&456&136&135&125\\
 \Gamma_2&036&016&234&145&256&124&035\\
 \Gamma_3&016&235&245&136&045&134&026
\end{array}
\]
例えば $024$ は順序付き三つ組 $(0,2,4)$ を意味する。
縮約後の三添字テンソルを $M_\Gamma^{klm}$ とすると、ここでの規格化は
\[
 (\beta_\Gamma)_{ij}
     =\sum_{k,l,m=0}^4\epsilon_{ijklm}M_\Gamma^{klm}
\]
である。交代化の $1/6$ は入れない。

$g=\nabla F(q)$ とおくと、$g$ は
$V\otimes\det(V)^2$ に値を持つ双次数 $(1,4)$ の共変量である。
従って $h_\Gamma=\beta_\Gamma g$ は
$V^*\otimes\det(V)^9\simeq\Lambda^4V\otimes\det(V)^8$
に値を持つ双次数 $(8,10)$ の共変量、すなわち
$\mathcal D_{10}$ の元である。$\beta_\Gamma$ の交代性により
\[
 dF(h_\Gamma)=g^{\mathsf t}\beta_\Gamma g=0.
\]

三つの像の独立性は、次の一つの整数評価で確認できる。
\[
 a_*=
 \begin{pmatrix}
 0&1&-1&1&-1\\
 -1&0&1&0&0\\
 1&-1&0&0&0\\
 -1&0&0&0&1\\
 1&0&0&-1&0
 \end{pmatrix},
 \qquad
 F_{\mathrm K}=3x_0x_4^2+6x_1x_2x_3.
\]
ここでは $q(a_*)=(1,1,1,1,1)$、
$g_{\mathrm K}=(3,6,6,6,6)$ である。
対称テンソルの規格化は $F_{044}=F_{123}=1$ とその置換で指定される。
三つの $h_{\Gamma}$ の第 $0,1,2$ 成分を列に並べると
\[
 \begin{pmatrix}
 192&0&-48\\
 336&66&-96\\
 -816&-6&48
 \end{pmatrix},
 \qquad
 \det=-1990656\ne0
\]
を得る。この計算は標準ライブラリだけを使う
\texttt{v11\_filtered\_koszul\_certificate.py} で、全中間値を整数として
再現できる。
評価に用いた $F_{\mathrm K}$ が特異であることは問題にならない。
これは普遍的な多項式共変量の線形独立性を証明する評価であり、
非零の小行列式は滑らかな三次式からなる稠密開集合でも恒等的には消えない。

従って $\dim\ker dF\ge3$ であり、
$\dim\mathcal D_{10}=17$ から $\operatorname{rank}dF\le14$ を得る。
さらに $\dim\mathcal B_{14}=3$ なので、この具体的な写像は単射である。
階数14の下界が得られた後には
$\ker dF=\operatorname{im}\mathcal B_{14}$ も従うが、
この等式を上界の証明に先取りして使ってはいない。
\end{proof}

\subsection{Fermat 型評価で失われる成分}\label{subsec:v11-fermat-scope}

\(F_{\mathrm{Fer}}=x_0^3+\cdots+x_4^3\) とし、その有限安定化部分群
\[
 H=\mu_3^5\rtimes\mathfrak S_5\subset\mathrm{GL}(V)
\]
を考える。以下は対角型三次式での評価の限界を与えるものであり、
接触写像の階数そのものの計算ではない。
\(
 \operatorname{Sym}^{d}(\operatorname{Sym}^{3}V)
 =\bigoplus_\lambda M_\lambda\otimes S_\lambda V
\)
と分解すると、\(F_{\mathrm{Fer}}^d\) の各成分は
\(M_\lambda\otimes(S_\lambda V)^H\) に属する。
従って、多重度空間から Fermat 軌道上の共変式への制限写像の階数は
\[
 \min\{\dim M_\lambda,\dim(S_\lambda V)^H\}
\]
以下である。非零の対角係数を持つ三次式はすべて Fermat 三次式の
\(\mathrm{GL}(V)\)-軌道に属し、零係数の場合も多項式恒等式として従う。

\begin{center}
\begin{tabular}{c|rrrrr|c}
 &\(j=0\)&\(j=1\)&\(j=2\)&\(j=3\)&\(j=4\)&制限階数の上界\\\hline
\(C_{14}\) の多重度&2&2&7&2&2&\\
\(S_{(11,8-j,8-j,j,j)}V\) の \(H\)-固定次元&21&8&7&0&0&11\\\hline
\(U_8\) の多重度&1&1&3&1&1&\\
\(S_{(8,8-j,8-j,j,j)}V\) の \(H\)-固定次元&3&1&2&0&0&4
\end{tabular}
\end{center}

特に \(C_{14}\) の分割 \((11,5,5,3,3)\) と \((11,4,4,4,4)\) に対応する
二次元ずつ、計四次元の成分は、すべての対角三次式上で恒等的に消える。
後者については指標計算を行わずにも分かる。
\(S_{(11,4,4,4,4)}V=\det^4V\otimes\operatorname{Sym}^7V\) の重みは
各成分が四以上であり、\(\mu_3^5\)-固定なら各成分が三の倍数なので六以上となる。
その総和が二十七であることに反する。
同じ理由で \(U_8\) の分割 \((8,4,4,4,4)\) の成分も消える。
表全体からは \(U_8\) の対角制限の核が少なくとも三次元であることが分かる。
また \(D_{10}\) の十一個の関連分割について同じ計算を行うと、
対角制限の階数は十以下、核は少なくとも七次元である。
いずれも上界・下界であり、制限写像の階数がこれらの値に一致するとは主張しない。

計算は有限群の指標平均だけで再現できる。
\(\pi\vdash5\) を巡回型、\(\ell(\pi)\) を巡回数、
\(z_\pi=\prod_r r^{m_r}m_r!\) とする。
各巡回の位相積を \(\zeta^{k_c}\), \(k_c\in\{0,1,2\}\) とすると
\[
 \sum_{n\geq0}h_n(g)t^n
 =\prod_{c\in\pi}(1-\zeta^{k_c}t^{|c|})^{-1},\qquad
 \chi_\lambda(g)=\det(h_{\lambda_i-i+j}(g)).
\]
従って
\[
 \dim(S_\lambda V)^H
 =\sum_{\pi\vdash5}\frac{1}{z_\pi3^{\ell(\pi)}}
   \sum_{(k_c)}\chi_\lambda(g).
\]
添付コードは \(\mathbb Z[\zeta]/(\zeta^2+\zeta+1)\) の整数対で計算し、
恒等元での指標が Weyl 次元公式に一致することも検査する。

対角軌道の法空間は
\[
 N=\operatorname{Jac}(F_{\mathrm{Fer}})_3
 =\left\langle x_ix_jx_k:i<j<k\right\rangle
\]
である。巡回型表示におけるその指標は
\(\chi_N(g)=[z^3]\prod_{c\in\pi}(1+\zeta^{k_c}z^{|c|})\) で与えられる。
同じ平均に \(\overline{\chi_N}\) を掛けると、上記二つの消える
\(C_{14}\) 成分について
\[
 \dim\operatorname{Hom}_H(N,S_{(11,5,5,3,3)}V)=3,\qquad
 \dim\operatorname{Hom}_H(N,S_{(11,4,4,4,4)}V)=2
\]
を得る。このため非対角方向の一階微分による検出は表現論的には可能である。
これは実際の微分写像の非零性や単射性を証明するものではない。
一方 \(U_8\) の分割 \((8,5,5,3,3)\) では固定次元とこの \(\operatorname{Hom}\) 次元が
ともに零であるため、対角軌道に沿って少なくとも二次で消える。
以上により、対角型三次式だけの有限評価の階数を普遍接触写像の階数と
同一視することはできない。今回の普遍階数の証明では、非対角の法方向の
微分を加えることで、対角型評価で見えない方向を検出する。

\subsection{14次の接触写像の階数}
\label{subsec:v11-contact-rank-fourteenth}

$\kappa$ と $\operatorname{Alt}$ は第10版と同じブロック抽出と交代化とする。
次の七つの普遍座標ベクトル場を用いる：
\[
 \Gamma_0=\operatorname{Alt}(K_9),\qquad
 \Gamma_i=\operatorname{Alt}\,\kappa\!\left(
       h\bigl(A_1h(A_iA_{8-i})\bigr)\right),\quad 1\le i\le6.
\]
各 $\Gamma_i$ は $\mathcal G_9$ に属する。
実際、Chevalley反対称化による正規表示、積、$h$、$\kappa$、交代化は
いずれも普遍的な $GL(V)$ 同変操作であり、$F$ に多項式的である。
$a$ の次数は9で、第\ref{sec:contact-twelve-thirteen}節の補助重み
\[
 \mathrm{wt}(F)=1,\quad\mathrm{wt}(a)=-1,\quad
 \mathrm{wt}(x)=0,\quad\mathrm{wt}(\theta)=-\tfrac12,
 \quad\mathrm{wt}(\eta)=\tfrac12
\]
を用いると、定数 $\Lambda^4V\to V$ ブロックの係数は $F$ に8次となる。
部分項自身にMaurer--Cartan方程式を要求しているわけではない。

\begin{proposition}[14次の接触階数]
\label{prop:v11-contact-rank-fourteenth}
上で指定した普遍正則な同変写像について
\[
 \operatorname{rank}(dP\colon\mathcal G_9\longrightarrow\mathcal C_{14})=14.
\]
従って $\dim\mathcal C_{14}=15$ と合わせると、通常の線形接触商は1次元である。
\end{proposition}
\begin{proof}
上界14は独立に構成した命題\ref{prop:v11-three-koszul-images}の三つのKoszulベクトルと
$\dim\mathcal D_{10}=17$ による。ここでは下界14を証明する。

$F_*=x_0^3+x_1^3+x_2^3+x_3^3-4x_4^3$、$u_i=e_i-e_0$ とし、
\begin{align*}
 a^{(1)}&=u_1\wedge u_2+u_3\wedge u_4,\\
 a^{(2)}&=a^{(1)}+u_1\wedge u_3,\\
 a^{(3)}&=2u_1\wedge u_2+\tfrac12u_3\wedge u_4,\\
 a^{(4)}&=u_1\wedge u_2+u_1\wedge u_3
              +u_1\wedge u_4+u_2\wedge u_3.
\end{align*}
いずれも $q(a^{(j)})=(1,1,1,1,1)$、$P(F_*,a^{(j)})=0$ を満たす。
さらに
\[
 G=6x_0x_1x_2-6x_4^3
\]
とおけば $G(q(a^{(j)}))=0$ なので、$F_*+tG$ に沿う方向微分も
$P\mathcal U_8$ を消す線形汎関数になる。

七つの線形汎関数を
\[
 L_1,\ldots,L_4=\operatorname{ev}_{(F_*,a^{(1)})},\ldots,
                   \operatorname{ev}_{(F_*,a^{(4)})},
 \qquad
 L_{4+j}=\left.\frac{d}{dt}\right|_{t=0}
             \operatorname{ev}_{(F_*+tG,a^{(b_j)})},
 \quad (b_1,b_2,b_3)=(1,2,4)
\]
と定める。これは普遍多項式を区別するための評価であり、
接触変換において $F$ を変更するという意味ではない。

正確な有理数計算で得られる行列
$M=(L_j(dP(\Gamma_i)))_{1\le j\le7,\,0\le i\le6}$ は
\[
 \det M=
 \frac{60571203387672046774905265545089621}{26572050000000000}\ne0
\]
を満たす。独立な検算として、分母が可逆である素数 $1000003$ を法とした
行列式は $916511$ である。

従って七つの像 $dP(\Gamma_i)$ は $P\mathcal U_8$ を法として独立である。
一方、$\dim\mathcal U_8=7$、多項式環は整域であり、Eulerの等式
\[
 dP(a u_8/6)=P u_8
\]
から $P\mathcal U_8$ は既に接触写像の像に含まれる7次元の部分空間である。
従って階数は少なくとも $7+7=14$ であり、上界と一致する。
\end{proof}

\paragraph{方向微分の計算規約と再現性.}
Clifford計量を $B=\operatorname{Hess}_F/6$、方向微分を
$C=\operatorname{Hess}_G/6$ と書く。
$\theta$ を正規順序に置き、$\eta$ はChevalley反対称化した外積記号として扱う。
この記号の積 $\mu_B$ に対し
\[
 \left.\frac{d}{dt}\right|_{t=0}\mu_{B+tC}(U,V)
 =-\sum_{i,j} C_{ij}\,
       \mu_B(D^R_{\eta_i}U,D^L_{\eta_j}V)
\]
である。これはClifford積を与える、$\theta,\eta$ について定係数の縮約作用素を微分した恒等式である。
右微分と左微分の符号を区別して用い、$h$ は $F$ に依存しないので、
Tate再帰全体を $\mathbb Q[t]/(t^2)$ 上でそのまま微分できる。
方向微分行の値には
\[
 dG(q)\,(dq)_a(\Gamma_i)
 +dF(q)\,(dq)_a\bigl(D_F\Gamma_i[G]\bigr)
\]
の両項を含めている。

ブロック抽出のための $x$ 次数の切断も正確である。
Clifford積の縮約は $x$ 次数を下げず、$h$ だけがそれを1下げる。
従って $\kappa h(UV)$ には積の $x$ 次数1以下、
$\kappa h(A_1h(UV))$ の内側の積には2以下を残せば十分である。
方向微分の縮約 $C_{ij}$ も $x$ に1次なので同じ議論が適用される。

\texttt{v11\_contact\_rank\_certificate.py} は保存した七行について
Pfaffian、その微分、三次式の勾配、上の行列式を正確に再計算する。
\texttt{--recompute} を付けると、外部の計算キャッシュを読まずに
四点の $A_1,\ldots,A_8$ と三つの方向微分を再構成する。
各次数のMaurer--Cartan方程式は有理数で検証され、別実装の32次元外積作用素でも
0次から8次まで検証される。
さらに全ての使用ブロックを $x=0$ の独立な外積・縮約作用素と照合する。


特に $\dim\ker(dP|_{\mathcal G_9})=14$ であり、
$\delta q$ の11次元の核に加えて三つのKoszul像の持ち上げが核を与える。
$\ker(dF|_{\mathcal D_{10}})=\operatorname{im}\mathcal B_{14}$ もこの段階で従う。
これは当該共変量成分の結論であり、Pfaffian代入後のKoszul複体全体の
完全性を主張するものではない。

\subsection{四次単元の普遍的なテンソル公式}
\label{subsec:v11-tetrahedral-U4}

$W=\Lambda^2V\otimes(\det V)^{-1}\simeq\Lambda^3V^*$ という
自然な同型により、$a$ を三形式 $\tau$ と見なす。
体積を固定した座標では
\[
 \tau_{ijk}=\epsilon_{ijklm}a_{lm},\qquad l<m,
\]
であり、添字が重複するときは零とする。
$F=f^{ijk}x_ix_jx_k$ では順序付き三つ組全体を和に取る規約とする。
例えば普通の多項式係数が $6x_0x_1x_2$ なら、その6個の対称成分は全て1である。

\begin{proposition}[四面体縮約による $U_4$]
\label{prop:v11-tetrahedral-U4}
第10版の正規化を持つ四次単元は
\[
 U_4=\frac1{24}
 f^{a_1a_2a_3}f^{b_1b_2b_3}f^{c_1c_2c_3}f^{d_1d_2d_3}
 \tau_{a_1b_1c_1}\tau_{a_2b_2d_1}
 \tau_{a_3c_2d_2}\tau_{b_3c_3d_3}
\]
で与えられる。右辺では全ての反復添字を和に取る。
\end{proposition}
\begin{proof}
$F\in\operatorname{Sym}^3V$ の全ての上添字が
$\tau\in\Lambda^3V^*$ の下添字と縮約されているので、分子 $S$ は
$GL(V)$ 不変な多項式である。その二重次数は $(F^4,a^4)$ であり、
$S\in\mathcal U_4$ となる。既証の $\dim\mathcal U_4=1$ から
$S=cU_4$ となる定数 $c$ が存在する。

対角三次式 $F=\sum_i c_i x_i^3$ では四つの頂点の添字は全て相異なる
場合だけが寄与し、各四つ組の24通りの順序は同じ符号を与える。
従って
\[
 S=24\sum_{i=0}^4\prod_{j\ne i}(c_j a_{ij}).
\]
特に $F_*=\operatorname{diag}(1,1,1,1,-4)$ と
$a=u_1\wedge u_2+u_3\wedge u_4$ で $S=-96$、$U_4=-4$ なので
$c=24$ であり、主張を得る。
\end{proof}

この公式は $P=0$ でも除算を必要とせず、一般の非対角三次式について
$U_4$ とその任意の方向微分を直接与える。
\texttt{v11\_tetrahedral\_u4.py} は外部キャッシュを用いずに正規化を検証し、
四次多項式の正確な補間による方向微分と、非対角三次式を生む
行列式1の剪断による不変性も検算する。

\subsection{14次では低次の安定化群も作用する}
\label{subsec:v11-filtered-contact}

以下では $P=P_F$、$H=H_{10}$、$c=1/96$ と書き、
第10版で得た正規形を
\[
 f=P+cH+R_{14}+O(a^{15}),\qquad R_{14}\in\mathcal C_{14}
\]
とする。座標の線形項と方程式の先頭係数を固定し、$F$ 自体は変更しない。
14次の項を比較する際には、13次までを保つ普遍正則な接触変換の全体を用いる。

\begin{proposition}[14次の有効な接触作用]
\label{prop:v11-effective-contact-fourteen}
二つの14次項 $R_{14},R'_{14}$ が、上記の変換により14次まで同値であるための
必要十分条件は
\[
 R'_{14}-R_{14}\in
 dP(\mathcal G_9)+P\mathcal U_8+\CC\,HU_4
\]
である。従って有効な商は
\begin{equation}\label{eq:v11-effective-fourteen-quotient}
 \mathcal Q^{\mathrm{eff}}_{14}
 =\frac{\mathcal C_{14}}
 {dP(\mathcal G_9)+P\mathcal U_8+\CC\,HU_4}
 =\frac{\mathcal C_{14}}
 {dP(\mathcal G_9)+\CC\,HU_4}.
\end{equation}
最後の等号はEulerの等式
$dP(au_8/6)=Pu_8$ による。
\end{proposition}
\begin{proof}
普遍共変量の低次の消滅から、必要なjetは
\[
 \chi=a+\phi_5+\phi_6+\phi_7+\phi_8+\phi_9+O(a^{10}),
 \qquad
 u=1+b_4+b_6+b_8+O(a^9)
\]
と書ける。ここで $\phi_r\in\mathcal G_r$、$b_s\in\mathcal U_s$ である。
$\mathcal G_2=\mathcal G_3=\mathcal G_4=0$、
$\mathcal U_1=\mathcal U_2=\mathcal U_3=\mathcal U_5=\mathcal U_7=0$
を用いた。$X=\phi_5$ を導分としても記す。

$u(f\circ\chi)$ が13次まで $P+cH$ と一致する条件は、順に
\[
 XP+b_4P=0,\qquad dP(\phi_6)=0,\qquad
 dP(\phi_7)+b_6P=0,\qquad dP(\phi_8)=0
\]
である。その14次の変化は
\begin{align*}
 R'_{14}-R_{14}
 ={}&dP(\phi_9)+b_8P
       +\frac12d^2P(\phi_5,\phi_5)+b_4XP\\
   &\quad+c\bigl(XH+b_4H\bigr).
\end{align*}
他の混合項は15次以上に現れる。

一見すると第1行の二次項が非線形な作用を与えるが、これは高次の変更へ
吸収できる。実際、$XP=-b_4P$ と
\[
 X^2P=d^2P(\phi_5,\phi_5)+dP(D\phi_5\cdot\phi_5)
\]
から
\[
 \frac12d^2P(\phi_5,\phi_5)+b_4XP
 =-\frac12dP(D\phi_5\cdot\phi_5)
  -\frac12(Xb_4+b_4^2)P
\]
を得る。$D\phi_5\cdot\phi_5\in\mathcal G_9$、
$Xb_4+b_4^2\in\mathcal U_8$ なので、これらはそれぞれ
$\phi_9$、$b_8$ の変更に含められる。

残る項を計算する。第7版の5次座標変更の計算から
\[
 \delta q(\phi_5)=\lambda qU_4
\]
となる $\lambda\in\CC$ が存在する。$F$ は $q$ に3次、
$\det\operatorname{Hess}_F(q)$ は $q$ に5次なので
\[
 XP=3\lambda PU_4,\qquad XH=5\lambda HU_4.
\]
多項式環が整域であることと $XP+b_4P=0$ から
$b_4=-3\lambda U_4$ が従い、従って
\[
 c(XH+b_4H)=2c\lambda HU_4.
\]
これで必要性を得る。

逆に、高次の $dP(\phi_9)$ と $b_8P$ はそのまま座標変更と単元で実現する。
$HU_4$ の方向については、任意の $t\in\CC$ に対し
\[
 \chi_t=a+t\,aU_4+\frac52t^2aU_4^2,
 \qquad u_t=1-6tU_4+6t^2U_4^2
\]
を取る。この有限な多項式変換は13次までを保ち、14次を
\[
 R_{14}\longmapsto R_{14}+4ct\,HU_4
\]
と変える。$c\ne0$ なので任意の係数を実現でき、十分性も従う。
\end{proof}

\begin{remark}[径方向の変更を全次数で書く]
上の有限なjetは、原点で1となる枝を選んだ代数的冪級数
\[
 \chi_t=a(1-4tU_4)^{-1/4},\qquad
 u_t=(1-4tU_4)^{3/2}
\]
の必要部分である。$P$ は $a$ に6次、$H$ は $a$ に10次なので、
\[
 u_tP(\chi_t)=P,
 \qquad u_tH(\chi_t)=\frac{H}{1-4tU_4}
\]
が全次数で成立する。従って $HU_4$ の方向は、先頭の6次式を厳密に保つ
径方向の接触変換から生じている。
\end{remark}

\begin{lemma}[Hessian項の単元係数は径方向に正規化できる]
\label{lem:v11-radial-hessian-unit}
$A(F,a)$ を $A(F,0)=1$ を満たす普遍不変な形式冪級数とする。
このとき
\[
 P_F+cH_{10}A(F,a)\ \sim\ P_F+cH_{10}
\]
となる普遍同変な径方向の座標変更と単元が存在する。
$A$ が代数的冪級数なら変更も代数的であり、$A$ が収束する場合には
変更も原点近傍で収束する。
\end{lemma}
\begin{proof}
$\lambda(F,0)=1$ であるscalarな形式冪級数を
\[
 \lambda^4 A(F,\lambda a)=1
\]
によって定める。左辺から1を引いた式の、$a=0,\lambda=1$ における
$\lambda$ 微分は4なので、係数を次数ごとに解いて一意に構成できる。
各段階で割るのは非零定数4だけであり、係数の $F$ に関する正則性と
同変性を保つ。$A$ が収束するときは解析的陰関数定理が適用できる。
$A$ が代数的で $Q(F,a,A)=0$ を満たすなら、$\lambda$ は
$Q(F,\lambda a,\lambda^{-4})=0$ を満たす。
分母を払って得る多項式は非零である。実際、$a^\alpha Z^k$ は
$a^\alpha\lambda^{|\alpha|-4k}$ に移り、この単項式の対応は単射である。
従って $\lambda$ も代数的冪級数である。

座標変更 $\chi=\lambda a$ と単元 $u=\lambda^{-6}$ を取れば、
斉次性と定義式により
\[
 u\bigl(P_F(\chi)+cH_{10}(F,\chi)A(F,\chi)\bigr)
 =P_F+cH_{10}\lambda^4A(F,\lambda a)=P_F+cH_{10}.
\]
これで主張を得る。一般の境界方程式の残余項が $H_{10}$ の倍数に
まとまることまでは、この補題は主張していない。
\end{proof}

\subsection{第10版の正規化が実際の14次項に与える補正}
\label{subsec:v11-actual-fourteen-correction}

\begin{proposition}[14次項の正規化補正]
\label{prop:v11-actual-fourteen-correction}
未正規化のTate方程式を $d^{\mathrm T}=\tfrac12\det K$ とする。
命題\ref{prop:v10-explicit-thirteenth}の具体的な $\chi,u$ により得られる
14次項は
\begin{equation}\label{eq:v11-fourteen-corrected}
 R_{14}
 =d^{\mathrm T}_{14}+3U_4d^{\mathrm T}_{10}
 =d^{\mathrm T}_{14}+\frac1{32}HU_4-9PU_4^2.
\end{equation}
従って有効な商\eqref{eq:v11-effective-fourteen-quotient}では
\[
 [R_{14}]=[d^{\mathrm T}_{14}].
\]
\end{proposition}
\begin{proof}
第10版の座標変更の非恒等部分は7次と8次であり、$P$ に対してそれぞれ
12次、13次から作用する。14次への一次寄与には元の8次項または7次項が必要だが、
どちらも零である。二次寄与は18次以上であり、10次項への作用も16次以上である。

単元 $1+3U_4-V_6$ については、14次に寄与するのは
$3U_4d^{\mathrm T}_{10}$ だけである。$-V_6$ の14次への寄与は
$-V_6d^{\mathrm T}_8=0$ となる。従って最初の等号を得る。
$d^{\mathrm T}_{10}=H/96-3PU_4$ を代入すれば第2の等号となる。
最後の主張は、$PU_4^2\in P\mathcal U_8$ と
$HU_4$ が有効な接触像に含まれることから従う。
\end{proof}

さらに、上で得た正規化方程式に径方向の変更 $t=-3/4$ と単元
$1+9U_4^2$ を続けて適用すれば、$4ct=-1/32$ であることから
\[
 P+cH+R_{14}+O(a^{15})
 \ \sim\ P+cH+d^{\mathrm T}_{14}+O(a^{15})
\]
とできる。従って、未正規化の14次係数をそのまま14次の代表元とする
正規化も、具体的な普遍接触変換によって選べる。

この正規化補正の計算だけからは、$HU_4$ の $dP(\mathcal G_9)$ への所属は従わない。
次の汎関数の計算で、その所属も決定する。

\subsection{十四次商の生成元と実際の係数の消滅}\label{subsec:v11-fourteenth-complement}

まず \(\mathcal U_6\) の一つの元 \(J_6\) をテンソル縮約として定義する。
\(F(x)=\sum_{i,j,k}f^{ijk}x_ix_jx_k\) とし、
\[
 \tau(a)_{ijk}=\epsilon_{ijklm}a_{lm}\quad(l<m,
       \ \{l,m\}=\{0,\ldots,4\}\setminus\{i,j,k\})
\]
と置く。これは \(W\simeq\Lambda^3V^*\) による自然な三形式である。
六個の \(F\)-頂点を \(0,\ldots,5\) とし、六個の \(\tau\)-頂点の接続先を
\begin{equation}\label{eq:v11-unit-six-graph}
 (024),\ (034),\ (035),\ (124),\ (125),\ (135)
\end{equation}
とする。各 \(F\)-頂点には三本の辺が入る。
各辺に添字を一つ割り当て、\(F\)-頂点で対称テンソル \(f\) を、
各三つ組で昇順に \(\tau\) を置き、十八個の添字をすべて縮約した値を
\(J_6(F,a)\) とする。追加の正規化係数は入れない。
この定義から自然性と次数 \((\deg_F,\deg_a)=(6,6)\) が従い、\(J_6\in\mathcal U_6\) である。
対角三次式では六個の頂点の色 \(i_v\in\{0,\ldots,4\}\) を用いて
\[
 J_6(F,a)=\sum_{i_0,\ldots,i_5}
      \left(\prod_{v=0}^5c_{i_v}\right)
      \prod_{(uvw)\text{ in }\eqref{eq:v11-unit-six-graph}}
                   \tau_{i_ui_vi_w}
\]
となる。

既存の正規化 \(\Theta_8\) と合わせて
\(W_{14}=\Theta_8J_6\in\mathcal C_{14}\) と置く。
この \(\Theta_8\) も、必要なら
\[
 T^{ijk}=f^{irs}f^{jtu}f^{kvw}\tau_{rtv}\tau_{suw},\qquad
 \Theta_8(F,a)=\frac16 T^{ijk}q_iq_jq_k
\]
という自然な縮約で記述できる（反復添字はすべて縮約する）。
対角式でこれは
\(\sum_{i<j}a_{ij}^2\prod_{k\ne i,j}c_kq_k\) となり、従来の正規化に一致する。

以下の五点評価で \(W_{14}\) の商における非零性と実際の係数の消滅を同時に証明する。
\(F_*=x_0^3+x_1^3+x_2^3+x_3^3-4x_4^3\), \(u_i=e_i-e_0\) とし、
\[
\begin{split}
 a(b,p,r,s,t)={}&b\,u_1\wedge u_2+p\,u_1\wedge u_3
       +r\,u_1\wedge u_4+s\,u_2\wedge u_3+t\,u_2\wedge u_4\\
 &+\frac{1+pt-rs}{b}\,u_3\wedge u_4
\end{split}
\]
に、順に
\[
 (1,0,0,0,0),\ (1,1,0,0,0),\ (2,0,0,0,0),\
 (1,1,1,1,0),\ (1,1,2,0,1)
\]
を代入した点を \(a_0,\ldots,a_4\) とする。
すべて \(q(a_i)=(1,1,1,1,1)\), \(P_{F_*}(a_i)=0\) を満たす。
\begin{center}
\begin{tabular}{c|r|r|r|r}
\(i\)&\(J_6\)&\(\Theta_8\)&\(H_{10}U_4\)&\(d^{\mathrm T}_{14}\)\\\hline
0&888&\(-14\)&124416&\(-931186/225\)\\
1&10392&\(-26\)&0&\(-17749271/5400\)\\
2&\(442371/8\)&\(-97/2\)&124416&\(-208492553/10800\)\\
3&135096&\(-62\)&\(-373248\)&\(12734093/900\)\\
4&82776&\(-42\)&\(-373248\)&\(166164259/16200\)
\end{tabular}
\end{center}
\(f_i=f(F_*,a_i)\) と略記し、評価の有理線形結合を
\begin{equation}\label{eq:v11-rho-fourteen}
 \begin{split}
 \rho_{14}(f)=\frac{1}{3473309792250}\bigl(&-1353058 f_0+789525 f_1
                  -173792 f_2\\
                  &-292755 f_3-216195 f_4\bigr)
 \end{split}
\end{equation}
と定義する。
正確な有理数計算により、前節の七個の画像基底
\[
 dP(\operatorname{Alt}K_9),\qquad
 dP\bigl(\operatorname{Alt}\kappa h(A_1h(A_iA_{8-i}))\bigr)
 \quad(1\leq i\leq6)
\]
のすべてを \(\rho_{14}\) が消し、また
\begin{equation}\label{eq:v11-rho-fourteen-values}
 \rho_{14}(W_{14})=1,\qquad
 \rho_{14}(H_{10}U_4)=0,\qquad
 \rho_{14}(d^{\mathrm T}_{14})=0
\end{equation}
を満たす。添付コードはこの七個の rooted 共変式について、
反対称行列から Pfaffian 微分と \(dP\) を再構成し、すべてを消すことを検査する。
最初の四点は階数証明用の独立な証明書と照合し、第五点にも保存キャッシュを
用いない Tate 再帰と独立な外積作用素による再計算を用意している。
上に並べた最初の四個の像と \(W_{14}\) の五点評価行列の行列式は
\[
 -\frac{2679359403669770814551549251}{15552000}\ne0
\]
である。

\(P\mathcal U_8\) は全五点で消え、命題~\ref{prop:v11-contact-rank-fourteenth} の
普遍階数十四の証明により、
\(P\mathcal U_8\) の七次元と上記七個の非単元画像は \(dP(\mathcal G_9)\) 全体を張る。
従って \(\rho_{14}\) は普遍商上の線形形式であり、
\[
 \mathcal C_{14}=dP(\mathcal G_9)\oplus\mathbb C W_{14},\qquad
 \mathcal C_{14}/dP(\mathcal G_9)\simeq\mathbb C,
 \quad [f]\longmapsto\rho_{14}(f)
\]
を与える。対角評価の階数を普遍階数と同一視しているのではなく、
非対角一階微分と三個の独立な Koszul 関係による普遍階数十四を
先に用いている点が本質的である。

さらに \eqref{eq:v11-rho-fourteen-values} から
\(H_{10}U_4\in dP(\mathcal G_9)\) が従う。
したがって低次安定化群の追加作用を含む有効十四次商も一次元であり、
実際の Tate 行列式の十四次類はその中で零である。
低次正規化後の係数
\[
 R_{14}=d^{\mathrm T}_{14}+\frac1{32}H_{10}U_4-9PU_4^2
\]
についても \(\rho_{14}(R_{14})=0\) となる。

\begin{theorem}[実際の境界の14次までの普遍正規形]
\label{thm:v11-fourteenth-normal-form}
$d^{\mathrm T}=\frac12\det K$ をTate座標での普遍境界方程式とする。
普遍的に多項式な座標変更と単元を順次適用することにより、
\begin{equation}\label{eq:v11-fourteenth-normal-form}
 d^{\mathrm T}\ \sim\
 P_F+\frac1{96}\det\operatorname{Hess}_F(q(a))+O(a^{15})
\end{equation}
とできる。従って任意の滑らかな三次式 $F$ について、実際の境界はこの形を持つ。
また
\[
 \mathcal Q^{\mathrm{eff}}_{14}
   =\mathcal C_{14}/dP(\mathcal G_9)
   =\mathbb C\cdot[\Theta_8J_6]
\]
であり、$\rho_{14}$ はこの商の係数を与える。
\end{theorem}
\begin{proof}
命題\ref{prop:v11-contact-rank-fourteenth}と
\eqref{eq:v11-rho-fourteen-values}から
\[
 \ker\rho_{14}=dP(\mathcal G_9),\qquad
 H_{10}U_4,\ d^{\mathrm T}_{14}\in\ker\rho_{14}
\]
である。従って命題\ref{prop:v11-effective-contact-fourteen}の追加方向は
線形接触像に既に含まれ、有効な商も同じ一次元商である。
さらに\eqref{eq:v11-fourteen-corrected}から $R_{14}\in dP(\mathcal G_9)$ を得る。

より具体的には、階数の証明で選んだ七つの像 $dP(\Gamma_i)$ が張る空間と
$P\mathcal U_8$ が接触像を直和に分解するので、
\[
 R_{14}=dP(\Xi_9)+P B_8,\qquad
 \Xi_9\in\langle\Gamma_0,\ldots,\Gamma_6\rangle,\quad B_8\in\mathcal U_8
\]
と書ける。七つの係数は、同じ七つの線形汎関数 $L_j$ を用いて
$M^{-1}(L_j(R_{14}))_j$ と一意に定まる。
この分解は共変量空間内の分解なので、$B_8$ に $P$ による極は生じない。
第10版で正規化した方程式に
\[
 a\longmapsto a-\Xi_9(F,a),\qquad
 u=1-B_8(F,a)
\]
を適用すれば14次が消える。$H_{10}$ へのこの座標変更の寄与は18次以上であり、
既に正規化した13次以下を変えない。これで\eqref{eq:v11-fourteenth-normal-form}を得る。
\end{proof}

この有限次数の正規化は原点で微分が恒等な多項式座標変更と定数項1の単元からなり、
各固定 $F$ についてétaleかつ解析的局所可逆である。
普遍恒等式なので、固定した滑らかな $F$ に特殊化して同じ消去が成立する。
一方、上記の一次元という次元は指定した普遍共変量の商の次元であり、
固定 $F$ の全Jacobian商や、判別式を逆にした係数環上の商の次元と
同一視してはいない。

14次の消去によって、10次のHessian障害が消えることはない。
第\ref{sec:tate-boundary-tenth}節の全ての滑らかな $F$ に対する非所属は維持される。
また補題\ref{lem:v11-radial-hessian-unit}はHessian項の単元係数を全次数で
正規化するが、実際の全残余項がその形にまとまることは未証明である。

\clearpage
\section{Clifford反転と境界の全次数の偶正規形}
\label{sec:all-order-parity}

前節までの有限次数の消去を、実際のTate境界が持つ対称性に結び付ける。
$P_F=F(q(a))$、$H_{10}=\det\operatorname{Hess}_F(q(a))$ と書く。
本節の主張は、一つの普遍形式接触変換により全奇数次を同時に取り除き、
残る正規形の問題を16次以上の偶数次へ絞れることである。
10次のHessian障害は維持される。

\subsection{Clifford反転から生じる全次数の対称性}
\label{subsec:v12-parity}

以下の議論では $F$ を固定したまま変形する。
係数環は三次式の全係数の多項式環とし、$a$ に関する完備化だけを行う。
特に判別式その他の $F$ の多項式を逆にしない。
まず、Tate再帰が与える座標を、行列分解の同型類から一意に回復できることを示す。

\begin{lemma}[Tate切断への形式的な帰着と一意性]
\label{lem:v12-marked-tate-slice}
$\delta=[\mathsf D_0,-]_{\mathrm s}$ とし、
$h,\pi$ を\eqref{eq:tate-homotopy}のcontractionとする。
$\mathfrak m$ を変形パラメータのイデアルとする。
内部重み $3/2$ の奇作用素 $B\in\mathfrak m\mathsf E_{3/2}$ が
$(\mathsf D_0+B)^2=F$ を満たすなら、
内部重み0の偶可逆作用素 $g\equiv1\pmod{\mathfrak m}$ と
一意な形式パラメータ $b\in W\widehat\otimes\mathfrak m$ が存在して
\[
 g^{-1}(\mathsf D_0+B)g=\mathsf D_0+A(b)
\]
となる。ここで $A(b)$ はTate再帰\eqref{eq:tate-mc-recurrence}による切断である。
帰着を次数ごとの $h,\pi$ で行えば、各係数の $F$ に関する多項式性と
$GL(V)$ 同変性を保つ。
\end{lemma}
\begin{proof}
必要なcohomologyはcontractionから直接読める。
$\pi$ は $x$ と $\theta$ を消し、$\eta$ の外積記号だけを残す。
従って内部重み $3/2$ の奇部分では $\pi$ の像は
$\Lambda^3\langle\eta\rangle=A_1(W)$ であり、
内部重み0の偶部分では $\pi$ の像はscalarだけである。
また $\pi\delta=0$、$\delta h+h\delta=1-\pi$ が成立する。

存在を $\mathfrak m$ の次数について構成する。
既に次数 $n-1$ まで現在の作用素 $B$ と $A(b)$ が一致したとする。
その次数 $n$ の差を $R_n$ と書く。
二つのMaurer--Cartan等式を引くと、低次の項は一致しているので
$\delta R_n=0$ である。
$b$ の次数 $n$ の補正を
$A_1(\Delta b_n)=\pi R_n$ によって選ぶ。
この変更後に残る差は
\[
 R_n-\pi R_n=\delta hR_n
\]
である。$g_n=1-hR_n$ による共役
$D\mapsto g_n^{-1}Dg_n$ の次数 $n$ の変化は
$-\delta hR_n$ なので、この差を消す。
これを繰り返すと、$\mathfrak m$ 進完備性から $b,g$ を得る。

一意性では、$B=A(b),\ C=A(c)$ とし、
\[
 D_Bg=gD_C,\qquad D_B=\mathsf D_0+B,\quad
 D_C=\mathsf D_0+C,\quad g\equiv1\pmod{\mathfrak m}
\]
を仮定する。$b,c$ が次数 $N-1$ まで一致しているとき、
まず $g$ の次数 $N$ 未満の項を、同じ共役を保ちながら消せることを示す。
$g=1+g_m+O(\mathfrak m^{m+1})$、$m<N$ とすれば
上の等式の次数 $m$ は $\delta g_m=0$ を与える。
従って
\[
 g_m=c_m+\delta h(g_m),\qquad c_m=\pi g_m
\]
であり、$c_m$ は内部重み0のscalarである。
ここで、変形後の微分 $\delta_B=[D_B,-]_{\mathrm s}$ を用いて
\[
 s=1+c_m+\delta_B\bigl(h(g_m)\bigr)
\]
とおく。$D_B^2=F$ は中心元なので
\[
 \delta_B^2(T)=[F,T]=0.
\]
従って $s$ は $D_B$ と\emph{正確に}可換する。
また $s=1+g_m+O(\mathfrak m^{m+1})$ だから可逆であり、
$g$ を $s^{-1}g$ に置き換えても $D_Bg=gD_C$ は保たれ、
次数 $m$ の項だけが消える。

これを $m<N$ に繰り返した後は $g=1+g_N+O(\mathfrak m^{N+1})$ である。
次数 $N$ の等式は
\[
 B_N-C_N=-\delta g_N
\]
となる。両辺に $\pi$ を適用すると、Tate切断の
$\pi A(b)=A_1(b)$ と $\pi\delta=0$ から
$A_1(b_N-c_N)=0$、従って $b_N=c_N$ を得る。
$N$ に関する帰納法で $b=c$ が従う。
ここで一意なのは切断上のパラメータであり、共役作用素 $g$ 自体ではない。

すべての操作は積、$h,\pi$、および定数項1の作用素の形式逆元からなる。
各次数で行う操作は有限個であり、$h$ が割るのは正整数だけである。
従って $F$ に関する多項式性と同変性も保たれる。
\end{proof}

\begin{theorem}[普遍境界の形式的な偶性]
\label{thm:v12-all-order-parity}
Tate境界 $d^{\mathrm T}(F,a)=\tfrac12\det K(F,a)$ に対し、
$F$ の係数に多項式的な係数を持つ普遍同変な形式座標変更と形式単元により、
\[
 d^{\mathrm T}(F,a)\ \sim\ d_{\mathrm{ev}}(F,b),\qquad
 d_{\mathrm{ev}}(F,-b)=d_{\mathrm{ev}}(F,b)
\]
とできる。座標変更の線形部は恒等写像、単元の定数項は1である。
従って奇数次の項を、全次数で同時に消せる。
\end{theorem}
\begin{proof}
Clifford代数上の\emph{通常の}反準同型 $T\mapsto T^*$ を、
$x_i,\theta_i,\eta_i$ と $F$ の係数を固定し、積の順序を逆転させることにより定める。
これはsuper反準同型ではない。
Clifford関係の対称性から定義は整合的で、
\[
 (T^*)^*=T,\qquad \mathsf D_0^*=\mathsf D_0,\qquad
 A_1(a)^*=-A_1(a)
\]
である。最後の符号は三個の反対称化した $\eta$ を逆順にする符号
$(-1)^{3\cdot2/2}=-1$ である。
$(\mathsf D_0+A)^2=F$ なら $(\mathsf D_0+A^*)^2=F$ も成立する。
また偶作用素 $g$ に対して
\[
 (g^{-1}Dg)^*=g^*D^*(g^*)^{-1}
\]
なので、反転は $g\equiv1$ による共役同値類にも作用する。

補題\ref{lem:v12-marked-tate-slice}で $A(a)^*$ を切断へ戻すと、
形式写像 $\sigma$ と $g\equiv1$ が得られて
\begin{equation}\label{eq:v12-reversion-regauging}
 \mathsf D_0+A(\sigma(a))
 =g(a)^{-1}\bigl(\mathsf D_0+A(a)^*\bigr)g(a)
\end{equation}
となる。切断上のパラメータの一意性により
\[
 \sigma^2=\operatorname{id},\qquad \sigma(a)=-a+O(a^2).
\]
存在の構成から $\sigma$ と $g$ の係数は $F$ に多項式的で、構成は同変である。
特に中心的homothetyの重みを比較すれば、
$\sigma$ の $a$ に関する $r$ 次項は $F$ に $r-1$ 次であり、
既に用いている $\mathcal G_r$ に属する。

次に境界の方程式を追跡する。外積代数に双線形形式
\[
 \langle u,v\rangle
 =\bigl[\operatorname{rev}(u)\wedge v\bigr]_{\Lambda^5V}
\]
を入れる。ここで $\operatorname{rev}$ は外積の順序を逆転する。
5次元ではこの形式は対称で、外積作用 $\theta_i$ と縮約作用 $\xi_i$ は
ともに自己随伴である。従って $\eta_i$ も自己随伴であり、
上の $*$ はこの形式についての随伴と一致する。
実際、次数 $p$ の外積の反転は $(-1)^{p(p-1)/2}$ 倍であり、
外積と縮約の自己随伴性は外積のLeibniz則から直接従う。

$\Lambda^4V$ の符号付き補基底
$v_j=(-1)^j e_0\wedge\cdots\wedge\widehat e_j\wedge\cdots\wedge e_4$
と $\Lambda^1V$ の基底 $e_i$ の間のpairingは $\delta_{ij}$ である。
従って $x=0$ での $\Lambda^4V\to\Lambda^1V$ ブロックは、
$*$ により正確に $K^t$ に移り、その行列式は変わらない。

内部重み0の偶共役 $g$ の $x=0$ での作用も重要である。
この場合 $\eta_i=\xi_i$ であり、残る単項式では
$\theta$ の個数と $\eta$ の個数が等しい。
従って $g|_{x=0}$ は各 $\Lambda^kV$ を保存する。
その可逆なブロックを $g_k$ とすると、
\eqref{eq:v12-reversion-regauging}は
\[
 K(\sigma(a))=g_1^{-1}K(a)^t g_4
\]
を与える。ゆえに
\begin{equation}\label{eq:v12-boundary-involution-unit}
 d^{\mathrm T}(F,\sigma(a))
 =v(F,a)d^{\mathrm T}(F,a),\qquad
 v=\frac{\det g_4}{\det g_1},\quad v(F,0)=1.
\end{equation}
この単元は共役作用素から直接得ており、境界方程式による除算を必要としない。

標数0なので
\[
 b=\frac{a-\sigma(a)}2
\]
は線形部が恒等写像である形式座標変更となる。
$b(\sigma(a))=-b(a)$ だから、この座標では involution は $b\mapsto-b$ である。
逆写像を $a=a(b)$、$\widetilde d(F,b)=d^{\mathrm T}(F,a(b))$ と書けば
\eqref{eq:v12-boundary-involution-unit}から
\[
 \widetilde d(F,-b)=\widetilde v(F,b)\widetilde d(F,b),\qquad
 \widetilde v(F,0)=1.
\]
そこで
\begin{equation}\label{eq:v12-even-average}
 d_{\mathrm{ev}}(F,b)
 =\frac{\widetilde d(F,b)+\widetilde d(F,-b)}2
 =\frac{1+\widetilde v(F,b)}2\,\widetilde d(F,b)
\end{equation}
とおく。右の係数は定数項1の単元で、左は厳密に偶関数である。
形式逆写像と単元の構成も各次数で有限回の多項式操作しか必要としないので、
普遍的な正則性と同変性を保つ。
\end{proof}

\begin{corollary}[既存の有限正規化と両立する偶正規形]
\label{cor:v12-even-hessian-normal-form}
実際のTate境界には普遍同変な形式正規形
\begin{equation}\label{eq:v12-even-hessian-normal-form}
 \boxed{\displaystyle
 d^{\mathrm T}\ \sim\
 P_F+\frac1{96}H_{10}
       +\sum_{\substack{n\ge16\\ n\ \mathrm{even}}}R_n,
 \qquad R_n\in\mathcal C_n
 }
\end{equation}
が存在する。特に12次と14次の消去、および全奇数次の消去を同時に実現できる。
\end{corollary}
\begin{proof}
$\mathcal G_2=\mathcal G_3=\mathcal G_4=0$ である。
$\sigma^2=\operatorname{id}$ の最低次の非線形項を比較すると、
5次の項は奇数次なので零である。
6次以降の非線形項どうしの合成は11次から現れるため、同じ比較により
7次と9次の項も零である。
従って、上で構成した座標の逆写像は
\[
 a(b)=b+\phi_6(b)+\phi_8(b)+\phi_{10}(b)+O(b^{11})
\]
となる。

未正規化のTate方程式では
$d^{\mathrm T}_7=d^{\mathrm T}_8=d^{\mathrm T}_9=d^{\mathrm T}_{11}=0$
である。
上の座標変更による $P$ の二次のTaylor項は16次から現れる。
一次の項についても、14次以下の偶数次に寄与するには
元の7次または9次の項が必要であり、いずれも零である。
従って $a(b)$ で引き戻しても、6、8、10、12、14次の係数は変わらない。
算術平均\eqref{eq:v12-even-average}も、これらの偶数次の係数をそのまま保つ。
特にその10次係数は従来と同じ
$H_{10}/96-3PU_4$ である。

そこで命題\ref{prop:v10-explicit-thirteenth}の変更のうち、
7次の座標変更と4、6次の単元だけを用いる。
これは10次を $H_{10}/96$ にし、12次を消す。
8次の座標変更は13次の消去に使っていたものなので、
既に偶関数となった方程式には不要である。
これを省いても14次の係数は変わらない。
実際、14次へのその一次寄与は元の7次係数の微分であり、零である。
従って14次の係数は従来と同じ
\[
 R_{14}=d^{\mathrm T}_{14}+3U_4d^{\mathrm T}_{10}
\]
となる。定理\ref{thm:v11-fourteenth-normal-form}の
9次座標変更と8次単元で、この項も消せる。

7次と9次の座標変更はいずれも奇写像であり、
4、6、8次の単元はいずれも偶関数である。
従ってこれらの有限な変更は、方程式の全次数における偶性を保つ。
これで表示した正規形を得る。
\end{proof}

\begin{remark}[偶性の主張の範囲]
この定理は、全ての奇数次共変量空間が零であるという主張ではない。
例えば $\mathcal C_{13}$ と $\mathcal C_{15}$ は非零であり、
元のTate座標における奇数次係数も非零になり得る。
ここで得たのは、実際の境界の反転対称性に由来する、単一の形式接触変換である。
この第12版の証明は形式的な構成であったが、
第\ref{sec:algebraic-parity}節で同じ偶化の代数性と収束性を証明する。
係数が $F$ に多項式的であることは、$a$ に関する級数が多項式であることを意味しない。
偶数次の接触類の消滅は別の問題であり、16次については
第\ref{sec:contact-sixteen}節で決定する。
\end{remark}

\clearpage
\section{Pfaffian微分の全射性と15次の接触商}
\label{sec:contact-fifteen}

本節では、全次数で使えるPfaffian微分の構造を証明した後、
15次の共変量空間、低次安定化群、接触写像の階数を決定する。
前節の偶性は実際の境界に対する構造的な結果である。
ここで証明する15次接触写像の全射性は、許された普遍共変な15次候補全体についての結果であり、
別の有限計算による証明を持つ。

\subsection{Pfaffian微分の全次数での全射性}
\label{subsec:v12-uniform-pfaffian}

15次の計算に先立ち、前の各次数で個別に検証していた
$\delta q$ の全射性を、5次元の一つの表現論的恒等式にまとめる。
$A=\Lambda^2V$ とし、$r\ge1$ に対して
\[
 M_r=\operatorname{Sym}^{r-1}(\operatorname{Sym}^3V)
\]
および
\begin{align*}
 \mathcal G_r&=\operatorname{Hom}_{\mathrm{GL}(V)}
 \left(M_r\otimes\operatorname{Sym}^r A,
        \Lambda^2V\otimes\det(V)^{r-1}\right),\\
 \mathcal D_{r+1}&=\operatorname{Hom}_{\mathrm{GL}(V)}
 \left(M_r\otimes\operatorname{Sym}^{r+1} A,
        \Lambda^4V\otimes\det(V)^{r-1}\right)
\end{align*}
と書く。

\begin{proposition}[普遍共変量上のPfaffian微分は全次数で全射]
\label{prop:v12-uniform-pfaffian-surjectivity}
$\dim V=5$ のとき、全ての $r\ge1$ について
\[
 \delta q:\mathcal G_r\longrightarrow\mathcal D_{r+1}
\]
は全射である。
\end{proposition}
\begin{proof}
まず $F$ の係数を取り除いた多項式表現の写像
\[
 \mathsf d:V\otimes\operatorname{Sym}^{r+1}A
       \longrightarrow\Lambda^3V\otimes\operatorname{Sym}^r A,
 \qquad
 \mathsf d(vf)=\sum_{i<j}(v\wedge e_i\wedge e_j)
                   \frac{\partial f}{\partial a_{ij}}
\]
を考える。これは対称代数の標準的な余微分と外積の合成なので同変である。
Littlewood分解とPieri則により、始域は重複度1の三種類の成分からなる。
$u+v=r+1$、$u\ge v\ge0$ として、その最高重みは
\[
\begin{array}{c|c|c}
 \text{型}&\text{最高重み}&\text{条件}\\\hline
 \mathrm A_{u,v}&(u+1,u,v,v,0)&v\ge0\\
 \mathrm B_{u,v}&(u,u,v+1,v,0)&u>v\ge0\\
 \mathrm C_{u,v}&(u,u,v,v,1)&v\ge1
\end{array}
\]
であり、異なる行・添字から同じ重みは生じない。

\[
 Q=a_{01}a_{23}-a_{02}a_{13}+a_{03}a_{12},\qquad
 w=e_0a_{12}-e_1a_{02}+e_2a_{01},
\]
\[
 \Theta=\sum_{i=0}^4(-1)^i e_i\operatorname{Pf}_{\widehat i}(a)
       \in V\otimes\operatorname{Sym}^2A
\]
とおく。$\Theta$ はこの空間の $\det(V)$ 成分である。
三種類の最高重みベクトルはそれぞれ
\[
 e_0a_{01}^{u-v}Q^v,\qquad
 wa_{01}^{u-v-1}Q^v,\qquad
 \Theta a_{01}^{u-v}Q^{v-1}
\]
と選べる。
$e_{ijk}=e_i\wedge e_j\wedge e_k$ と書けば、直接微分により
\begin{align*}
 [e_{023}]\,\mathsf d(e_0a_{01}^{u-v}Q^v)
  &=v\,a_{01}^{u-v+1}Q^{v-1}\quad(v\ge1),\\
 \mathsf d(wa_{01}^{u-v-1}Q^v)
  &=(u+2)e_{012}a_{01}^{u-v-1}Q^v,\\
 [e_{234}]\,\mathsf d(\Theta a_{01}^{u-v}Q^{v-1})
  &=(v+2)a_{01}^{u-v+1}Q^{v-1}.
\end{align*}
第2式では $\mathsf dw=3e_{012}$、
$w\wedge e_{01}=a_{01}e_{012}$、
$w\wedge\sum_{i<j}e_{ij}\partial_{ij}Q=Qe_{012}$ を用いた。
第3式では $[e_{234}]\mathsf d\Theta=3a_{01}$ であり、
$Q^{v-1}$ の微分からさらに $(v-1)a_{01}$ が加わる。
従ってこれらの係数はいずれも非零である。
一方、$\mathrm A_{r+1,0}$ の最高重みベクトル
$e_0a_{01}^{r+1}$ の像は零である。

始域が重複度1であることと既約性から、これで核全体が決まり、
\begin{equation}\label{eq:v12-uniform-pfaffian-kernel}
 \ker\mathsf d=S_{(r+2,r+1,0,0,0)}V
\end{equation}
となる。最小の $r=1$ においても、残る型
$\mathrm A_{1,1},\mathrm B_{2,0},\mathrm C_{1,1}$ の上の係数は
 それぞれ1、4、3であり、同じ結論が成立する。

最後にこの核を普遍共変量へ戻す。
$\Lambda^2V\simeq(\Lambda^3V)^*\otimes\det V$、
$\Lambda^4V\simeq V^*\otimes\det V$ により
\begin{align*}
 \mathcal G_r&\simeq\operatorname{Hom}_{\mathrm{GL}(V)}
 \left(M_r\otimes\Lambda^3V\otimes\operatorname{Sym}^rA,
       \det(V)^r\right),\\
 \mathcal D_{r+1}&\simeq\operatorname{Hom}_{\mathrm{GL}(V)}
 \left(M_r\otimes V\otimes\operatorname{Sym}^{r+1}A,
       \det(V)^r\right).
\end{align*}
この表示で $\delta q$ は、対称化規約による非零定数倍を除いて
$\operatorname{id}_{M_r}\otimes\mathsf d$ への前合成である。
核\eqref{eq:v12-uniform-pfaffian-kernel}と $M_r$ から
$\det(V)^r$ へのpairingが存在するには、$M_r$ に最高重み
\[
 (r,r,r,-1,-2)
\]
の表現が現れなければならない。しかし $M_r$ は多項式表現なので、そのような成分はない。
従って全ての $\mathcal D_{r+1}$ の元は $M_r\otimes\ker\mathsf d$ を消す。
標数0の $\mathrm{GL}(V)$ の完全可約性により、その像上の同変写像は
全終域へ同変に延長できる。これが求める全射性である。
なお $r=1$ の共変量は $\mathcal G_1=\CC\,a$、
$\mathcal D_2=\CC\,q$ で、写像は $\delta q(a)=2q$ そのものである。
\end{proof}

この命題は $dF$ の核や、Pfaffian代入後のKoszul複体の完全性を主張しない。
今後の各次数で必要な階数の問題は
$dP=dF\circ\delta q$ のうち、$dF$ の部分に集約される。

\subsection{15次の共変量空間とPfaffian微分}
\label{subsec:v12-fifteenth-characters}

$A=\Lambda^2V$ とし、定数直線因子を戻した表示で
\[
 \mathcal D_{11}=\operatorname{Hom}_{\mathrm{GL}(V)}
 \left(\operatorname{Sym}^{9}(\operatorname{Sym}^3V)
       \otimes\operatorname{Sym}^{11}A,
       \Lambda^4V\otimes\det(V)^9\right)
\]
とおく。

\begin{proposition}[15次の次元とPfaffian微分]
\label{prop:v12-fifteenth-characters}
\[
 \dim\mathcal C_{15}=7,\qquad
 \dim\mathcal G_{10}=35,\qquad
 \dim\mathcal U_9=1,\qquad
 \dim\mathcal D_{11}=22.
\]
また $\delta q:\mathcal G_{10}\to\mathcal D_{11}$ は全射であり、
その核は13次元である。
\end{proposition}
\begin{proof}
Littlewood分解と整数alternantによるSchur重複度計算を用いる。
$\mathcal C_{15}$ で必要となる
$\lambda_j=(12,9-j,9-j,j,j)$、$0\le j\le4$ の
$\operatorname{Sym}^{10}(\operatorname{Sym}^3V)$ における重複度と、
$\mathcal U_9$ で必要となる
$\mu_j=(9,9-j,9-j,j,j)$ の
$\operatorname{Sym}^{9}(\operatorname{Sym}^3V)$ における重複度は
\[
\begin{array}{c|rrrrr}
 j&0&1&2&3&4\\\hline
 {}[S_{\lambda_j}V]&0&0&3&3&1\\
 {}[S_{\mu_j}V]&0&0&0&1&0
\end{array}
\]
である。$\mathcal G_{10}$ と $\mathcal D_{11}$ の次元は、それぞれ
\[
 [(10,10,9,9,9)]
 \bigl(\operatorname{Sym}^{9}(\operatorname{Sym}^3V)
                  \otimes\operatorname{Sym}^{10}A\bigr)=35,
\]
\[
 [(10,10,10,10,9)]
 \bigl(\operatorname{Sym}^{9}(\operatorname{Sym}^3V)
                  \otimes\operatorname{Sym}^{11}A\bigr)=22
\]
から得られる。

全射性は命題\ref{prop:v12-uniform-pfaffian-surjectivity}を $r=10$ に適用すれば従う。
この次数に限った検証も記録する。双対写像
\[
 \mathsf d:V\otimes\operatorname{Sym}^{11}A
       \longrightarrow\Lambda^3V\otimes\operatorname{Sym}^{10}A
\]
のPieri成分で、非単元方向に関係する最高重みと非零係数は次の通りである。
前の次数と同様に
$\mathrm A_{u,v}=(u+1,u,v,v,0)$、
$\mathrm B_{u,v}=(u,u,v+1,v,0)$ と書く。
\[
\begin{array}{c|rrrr|rrrr}
 \text{型}&\mathrm A_{9,2}&\mathrm A_{8,3}&\mathrm A_{7,4}&\mathrm A_{6,5}
 &\mathrm B_{9,2}&\mathrm B_{8,3}&\mathrm B_{7,4}&\mathrm B_{6,5}\\\hline
 \text{$F$側の重複度}&1&2&6&2&1&6&1&2\\
 \text{非零係数}&2&3&4&5&11&10&9&8
\end{array}
\]
これらは21方向を与え、残る一方向は
$\delta q(au_9)=2qu_9$ で与えられる単元方向である。
添付コードは十二個のPieri・plethysm成分と上の八つの最高重みの像を
厳密な多項式として検証する。従って階数は22、核の次元は $35-22=13$ である。
\end{proof}

低次安定化群に関係する空間についても、同じ計算で
\[
 \dim\mathcal G_6=\dim\mathcal D_7=\dim\mathcal B_{11}=1
\]
を得る。ここで
\[
 \mathcal B_{11}
 =\operatorname{Hom}_{\mathrm{GL}(V)}
 \left(\operatorname{Sym}^{4}(\operatorname{Sym}^3V)
       \otimes\operatorname{Sym}^{3}A,
       \Lambda^2V^*\otimes\det(V)^4\right)
\]
である。

\paragraph{この次数でも対角評価には核がある.}
第11版と同じ有限Fermat安定化群
$H=\mu_3^5\rtimes\mathfrak S_5$ の指標平均から、
$\mathcal C_{15}$ に関係する三つの非零成分について
\[
\begin{array}{c|ccc}
 \lambda&(12,7,7,2,2)&(12,6,6,3,3)&(12,5,5,4,4)\\\hline
 \text{普遍重複度}&3&3&1\\
 \dim(S_\lambda V)^H&4&3&0
\end{array}
\]
を得る。従って最後の一成分は全ての対角三次式で恒等的に消え、
対角制限の階数は高々6である。
この最後の成分については
$N=\operatorname{Jac}(F_{\mathrm{Fer}})_3$ として
\[
 \dim\operatorname{Hom}_H(N,S_{(12,5,5,4,4)}V)=1
\]
であり、法方向の一階微分による検出は表現論的には可能である。
一方 $\mathcal U_9$ の一成分 $(9,6,6,3,3)$ の固定次元は1である。
これらは制限写像の階数に対する上界であり、実際の評価写像の階数は
以下の具体的な接触像の計算で別に決定する。

なお、15次に関係するKoszul関係の二つの源の次元は
\[
\begin{split}
 &\dim\operatorname{Hom}_{\mathrm{GL}(V)}
 \left(\operatorname{Sym}^{8}(\operatorname{Sym}^3V)
       \otimes\operatorname{Sym}^{7}A,
       \Lambda^2V^*\otimes\det(V)^8\right)=19,\\
 &\dim\operatorname{Hom}_{\mathrm{GL}(V)}
 \left(\operatorname{Sym}^{7}(\operatorname{Sym}^3V)
       \otimes\operatorname{Sym}^{3}A,
       \Lambda^3V^*\otimes\det(V)^6\right)=3.
\end{split}
\]
この二つの数から $dF$ の階数を引き算だけで決めることはできない。
Pfaffian代入後のKoszul複体の完全性は仮定していない。

\subsection{6次の安定化ベクトル場とHessian項への作用}
\label{subsec:v12-sixth-stabilizer}

$P=F(q(a))$、$h(p)=\det\operatorname{Hess}_F(p)$、$H=h(q(a))$ と書く。
四次単元 $U_4$ の $a$ に関する微分を交代行列
\[
 B_{ij}=\frac{\partial U_4}{\partial a_{ij}}\quad(i<j),\qquad
 B_{ji}=-B_{ij}
\]
で表す。
これは $(F,a)$ に関して次数 $(4,3)$、値域
$\Lambda^2V^*\otimes\det(V)^4$ の普遍共変量である。
\[
 T=(\operatorname{Hess}_F)(q),\qquad
 S=(\operatorname{Hess}h)(q)
\]
とおき、二つのベクトル場を
\begin{equation}\label{eq:v12-sixth-tenth-fields}
 \Phi_6=\frac12(TBa-aBT),\qquad
 \Psi_{10}=\frac14(SBa-aBS)
\end{equation}
で定める。

\begin{lemma}[6次安定化群の作用も15次のJacobian像に入る]
\label{lem:v12-sixth-hessian-action}
$\Phi_6$ は1次元の $\mathcal G_6$ の非零生成元であり、
\[
 \delta q(\Phi_6)=B\nabla F(q),\qquad dP(\Phi_6)=0.
\]
また $\Psi_{10}\in\mathcal G_{10}$ であり、
\[
 \delta q(\Psi_{10})=B\nabla h(q),\qquad
 dH(\Phi_6)=-dP(\Psi_{10}).
\]
特に $dH(\mathcal G_6)\subset dP(\mathcal G_{10})$ である。
\end{lemma}
\begin{proof}
Pfaffianの共変性を微分すると、任意の行列 $M$ について
\[
 (dq)_a(Ma+aM^t)=\operatorname{tr}(M)q-M^tq
\]
を得る。この式は $a$ の各点での接ベクトルの式なので、
$M$ が $a$ に依存する場合にも同じように使える。
$M=TB/2$ とすると対称行列と交代行列の積なのでtraceは零である。
$Tq=2\nabla F(q)$ から $\delta q(\Phi_6)=B\nabla F(q)$ を得る。
交代性により $\nabla F(q)^tB\nabla F(q)=0$、従って $dP(\Phi_6)=0$ である。

$h$ は $p$ に5次なので $Sq=4\nabla h(q)$ である。
$M=SB/4$ に同じ式を適用すると
$\delta q(\Psi_{10})=B\nabla h(q)$ となる。
従って
\[
 dH(\Phi_6)=\nabla h(q)^tB\nabla F(q)
 =-\nabla F(q)^tB\nabla h(q)=-dP(\Psi_{10}).
\]
次数は $\Phi_6$ が $(F,a)$ に $(5,6)$、
$\Psi_{10}$ が $(9,10)$ であり、値域のdeterminant twistも
それぞれ $\det(V)^5$、$\det(V)^9$ となる。
従って表示した共変量空間に属する。

非零性を具体的に確認する。
標準点 $F_*=x_0^3+x_1^3+x_2^3+x_3^3-4x_4^3$ と
$a^{(1)}=(e_1-e_0)\wedge(e_2-e_0)+(e_3-e_0)\wedge(e_4-e_0)$ では
\[
 q=(1,1,1,1,1),\qquad U_4=-4,
\]
\[
 (B_{01},B_{02},B_{03},B_{04})=(-4,4,-4,4),\qquad
 B_{ij}=0\quad(1\le i<j\le4).
\]
よって
\[
 \delta q(\Phi_6)=(-60,12,-12,12,-12)\ne0.
\]
命題\ref{prop:v12-fifteenth-characters}の $\dim\mathcal G_6=1$ から
$\Phi_6$ が生成元であることも従う。
\end{proof}

同じ標準点では
\[
 dH(\Phi_6)=1866240,\qquad dP(\Psi_{10})=-1866240.
\]
従ってHessian項への作用そのものが零なのではなく、15次のJacobian像に吸収される。
またTateブロックの交代部分について、標準点で
\[
 \delta q(\operatorname{Alt}K_6)
   =\left(\frac{20}{3},-\frac43,\frac43,-\frac43,\frac43\right)
\]
であることと1次元性を用いると、普遍恒等式
\[
 \operatorname{Alt}K_6=-\frac19\Phi_6,
 \qquad dH(\operatorname{Alt}K_6)=\frac19dP(\Psi_{10})
\]
を得る。

\paragraph{小さい独立計算証明書.}
\texttt{v12\_filtered\_contact\_certificate.py} は有理数演算だけを用い、
既存の五つの保存点で\eqref{eq:v12-sixth-tenth-fields}の二つのPfaffian微分、
$dH(\Phi_6)=-dP(\Psi_{10})$、
$\operatorname{Alt}K_6=-\Phi_6/9$ を全成分で確認する。
同じコードは前節の $\mathrm C_{u,v}$ 系列の係数 $v+2$ を
$1\le r\le12$ の全該当成分について整数多項式として検証する。
全次数・普遍的な主張そのものは、上記の表現論と行列恒等式の証明による。

\subsection{15次の有効な接触作用と正規化補正}
\label{subsec:v12-filtered-contact}

\begin{proposition}[15次の低次安定化群は新しい商方向を与えない]
\label{prop:v12-effective-contact-fifteen}
$c\in\CC$ を固定し、
\[
 f=P+cH+R_{15}+O(a^{16})
\]
を考える。$F$、座標の線形項、先頭係数を固定し、14次までの正規形を保つ
全ての普遍正則な接触変換を用いても、有効な15次商は
\[
 \mathcal Q^{\mathrm{eff}}_{15}
   =\mathcal C_{15}/dP(\mathcal G_{10})
\]
である。
\end{proposition}
\begin{proof}
必要なjetを
\[
 \chi=a+\phi_5+\phi_6+\phi_7+\phi_8+\phi_9+\phi_{10}+O(a^{11}),
 \qquad u=1+b_4+b_6+b_8+b_9+O(a^{10})
\]
と書き、$X=\phi_5$、$Y=\phi_6$ とおく。
低次の共変量の消滅は14次の議論と同じである。
13次までを保つ条件は
\[
 XP+b_4P=0,\qquad YP=0,\qquad
 dP(\phi_7)+b_6P=0,\qquad dP(\phi_8)=0.
\]
さらに14次を保つ式は、前節で求めた
$\phi_9,b_8,X,b_4$ の式である。
15次の変化はTaylor展開から正確に
\[
 R'_{15}-R_{15}
 =dP(\phi_{10})+b_9P+d^2P(X,Y)+b_4YP+c\,dH(Y)
\]
となる。他の二次項・単元の混合項はこの次数に達しない。

$YP=0$ を $X$ で微分すると
\[
 d^2P(X,Y)=-dP(DY\cdot X),\qquad DY\cdot X\in\mathcal G_{10}
\]
を得る。また $b_4YP=0$ であり、
補題\ref{lem:v12-sixth-hessian-action}により
$dH(Y)\in dP(\mathcal G_{10})$ である。
単元項もEulerの等式 $b_9P=dP(ab_9/6)$ によって同じ像に含まれる。
これで必要性を得る。
逆に任意の $dP(\phi_{10})$ は10次の座標変更により実現し、
14次以下を変えない。従って十分性も成立する。
\end{proof}

\begin{proposition}[既存の具体的な正規化は15次の係数を変えない]
\label{prop:v12-actual-fifteen-correction}
命題\ref{prop:v10-explicit-thirteenth}と
定理\ref{thm:v11-fourteenth-normal-form}で選んだ具体的な有限正規化の後では
\[
 R_{15}=d^{\mathrm T}_{15}.
\]
\end{proposition}
\begin{proof}
これらの正規化の座標変更は7、8、9次からなり、単元は4、6、8次からなる。
座標変更が15次に与え得る一次寄与は、それぞれ元の9、8、7次係数の微分であり、
全て零である。
二次の座標変更が $P$ に与える寄与は18次からである。
単元からの寄与も、それぞれ元の11、9、7次係数の倍数であり、全て零である。
また最小の単元と最小の座標変更の混合項は16次から現れる。
従って15次の係数は変わらない。
ここでは主正規化で使った具体的な変更を対象とし、
任意の5次・6次の安定化変更を加えた正規化について同じ等式を主張してはいない。
\end{proof}

\subsection{15次の接触写像は全射である}
\label{subsec:v12-contact-rank-fifteenth}

15次では、前の次数とは異なり、普遍的な線形接触商そのものが消える。
第10・11版と同じブロック抽出 $\kappa$ と交代化を用い、
六つの普遍座標ベクトル場を
\[
 \Gamma_0=\operatorname{Alt}(K_{10}),\qquad
 \Gamma_i=\operatorname{Alt}\,\kappa\!\left(
      h\bigl(A_1h(A_iA_{9-i})\bigr)\right),\quad1\le i\le5
\]
と定める。いずれも $\mathcal G_{10}$ に属する。
実際、Chevalley反対称化、Clifford積、$h$、$\kappa$、交代化は
普遍的な $\mathrm{GL}(V)$ 同変操作であり、$F$ に多項式的である。
$a$ の次数は10で、前版と同じ補助重み
\[
 \mathrm{wt}(F)=1,\quad\mathrm{wt}(a)=-1,\quad
 \mathrm{wt}(x)=0,\quad\mathrm{wt}(\theta)=-\tfrac12,
 \quad\mathrm{wt}(\eta)=\tfrac12
\]
から、定数 $\Lambda^4V\to V$ ブロックの係数は $F$ に9次となる。
これらの部分項自身にMaurer--Cartan方程式を課しているわけではない。

\begin{proposition}[15次の普遍接触商の消滅]
\label{prop:v12-contact-rank-fifteenth}
普遍正則な同変写像
\[
 dP:\mathcal G_{10}\longrightarrow\mathcal C_{15}
\]
は全射である。従って
\[
 \operatorname{rank}dP=7,\qquad \dim\ker dP=28,
 \qquad\mathcal Q^{\mathrm{lin}}_{15}
       =\mathcal Q^{\mathrm{eff}}_{15}=0.
\]
また $dF:\mathcal D_{11}\to\mathcal C_{15}$ は階数7、核は15次元である。
\end{proposition}
\begin{proof}
$F_*=x_0^3+x_1^3+x_2^3+x_3^3-4x_4^3$、$u_i=e_i-e_0$ とし、
\begin{align*}
 a^{(1)}&=u_1\wedge u_2+u_3\wedge u_4,\\
 a^{(2)}&=a^{(1)}+u_1\wedge u_3,\\
 a^{(3)}&=2u_1\wedge u_2+\tfrac12u_3\wedge u_4,\\
 a^{(4)}&=u_1\wedge u_2+u_1\wedge u_3
                    +u_1\wedge u_4+u_2\wedge u_3.
\end{align*}
全ての点で $q=(1,1,1,1,1)$、$P=F_*(q)=0$ である。
さらに $G=6x_0x_1x_2-6x_4^3$ とおけば $G(q)=0$ も成り立つ。
六つの線形汎関数を
\[
 L_j=\operatorname{ev}_{(F_*,a^{(j)})}\quad(1\le j\le3),
 \qquad
 L_{3+j}=\left.\frac{d}{dt}\right|_{t=0}
       \operatorname{ev}_{(F_*+tG,a^{(b_j)})},
 \quad (b_1,b_2,b_3)=(1,2,4)
\]
とする。これらは全て $P\mathcal U_9$ を消す。
方向微分は普遍多項式を区別するための評価であり、
接触変換における $F$ の変更を許すという意味ではない。

行列 $M=(L_j(dP(\Gamma_i)))_{1\le j\le6,\,0\le i\le5}$ は
正確な有理数計算により
\[
 \det M=
 \frac{4583737618198478767481154324593650656534581}
      {55540601303040000000000}\ne0
\]
を満たす。分母が可逆である素数 $1000003$ を法として、
行列から直接計算した行列式も $147271\ne0$ である。
従って六つの像 $dP(\Gamma_i)$ は $P\mathcal U_9$ を法として独立である。

一方、$\dim\mathcal U_9=1$、$P$ は非零多項式であり、
多項式環が整域であることから $P\mathcal U_9$ は1次元である。
Eulerの等式
\[
 dP(au_9/6)=Pu_9
\]
により、この方向も接触像に含まれる。従って
$\operatorname{rank}dP\ge6+1=7$ となる。
$\dim\mathcal C_{15}=7$ なので全射であり、
$\dim\ker dP=35-7=28$ である。
有効な接触同値もこの線形の座標変更を含むので、両方の商が零になる。

最後に $dP=dF\circ\delta q$ で、$\delta q$ は全射である。
従って $dF$ も階数7をもち、その核は $22-7=15$ 次元である。
\end{proof}

\paragraph{計算証明の範囲と再現性.}
方向微分の規約は第11版と同じである。
$B=\operatorname{Hess}_F/6$、$C=\operatorname{Hess}_G/6$ と書けば、
Chevalley反対称化した $\eta$ 記号のClifford積の微分は
\[
 D_B\mu_B(U,V)[C]
 =-\sum_{i,j}C_{ij}\,
       \mu_B(D^R_{\eta_i}U,D^L_{\eta_j}V)
\]
で与えられる。方向微分行には
\[
 dG(q)\,(dq)_a(\Gamma_i)
 +dF(q)\,(dq)_a\bigl(D_F\Gamma_i[G]\bigr)
\]
の両項を含めている。
Clifford積は $x$ 次数を下げず、$h$ だけが1下げるので、
定数ブロックに必要な積を $x$ 次数1以下、二重の $h$ の内側を2以下で
切断しても正確である。

\texttt{v12\_contact\_rank\_certificate.py} は、六行分の全ての
交代行列、Pfaffian、その微分、勾配との縮約、行列式を正確に検証する。
\texttt{--recompute} は私的な計算キャッシュを読まず、
四点の $A_1,\ldots,A_8$ と三つの一階変分を再構成する。
Maurer--Cartan方程式を有理数で検証し、別実装の32次元外積作用素と照合する。
使用する定数ブロックも $x=0$ の外積・縮約作用素で検算される。
\texttt{--compare-fresh} は別に保存した新規計算の六行と完全一致することを確認する。

\begin{corollary}[Pfaffian代入後のKoszul完全性は破れる]
\label{cor:v12-koszul-homology}
15次に関係する普遍同変なKoszul複体の成分で、
二番目のKoszulホモロジーは少なくとも1次元である。
\end{corollary}
\begin{proof}
$g=\nabla F(q)$ とし、前節の19次元の空間を $\mathcal B_{15}$、
その前にある3次元の空間を $\mathcal K_{3,15}$ と書く。
縮約による複体は
\[
 \mathcal K_{3,15}\xrightarrow{\iota_g}
 \mathcal B_{15}\xrightarrow{\beta\mapsto\beta g}
 \mathcal D_{11}\xrightarrow{h\mapsto g\cdot h}\mathcal C_{15}
\]
である。交代性から隣接する写像の合成は零となる。
最後の写像の核は15次元なので、
\[
 \dim\ker(\mathcal B_{15}\to\mathcal D_{11})\ge19-15=4.
\]
前の写像の像は高々3次元であり、従って
\[
 \dim\frac{\ker(\mathcal B_{15}\to\mathcal D_{11})}
 {\operatorname{im}(\mathcal K_{3,15}\to\mathcal B_{15})}
 \ge4-3=1.
\]
ここではホモロジーの正確な次元を主張していない。
この下界だけでも、二つの源の次元を交互に引き算して接触階数を決める方法は
この次数では成立しないことが分かる。
\end{proof}

\subsection{十一番目のTateブロックと実際の十五次項}

以下でも $a$ に関する次数を用い、$x$ に関する次数とは区別する。
Tate再帰を
\[
 A_n=-hR_n,\qquad R_n=\sum_{i=1}^{n-1}A_iA_{n-i},
 \qquad K_n=\kappa(A_n)
\]
と書く。$\kappa$ は $x=0$ における符号付き補集合基底の
$\Lambda^4V\to V$ ブロックであり、
$K(a)=\sum_{n\ge1}K_n(a)$、$d^{\mathrm T}=\tfrac12\det K$ である。

\begin{lemma}[十五次係数のための正確な切断]
\label{lem:v12-exact-eleventh-truncation}
$d^{\mathrm T}_{15}$ を計算するには $K_1,\ldots,K_{11}$ があれば十分である。
さらに、完全な $A_1,\ldots,A_8$ から出発して、
\[
\begin{array}{c|c|c}
 \text{求める項}&\text{残す曲率の $x$ 次数}&
                    \text{得られる $A_n$ の $x$ 次数}\\ \hline
 A_9&3\text{以下}&2\text{以下}\\
 A_{10}&2\text{以下}&1\text{以下}\\
 A_{11}&1\text{以下}&0
\end{array}
\]
の順に再帰を計算すれば、必要な十一個の定数ブロックをすべて正確に得る。
\end{lemma}

\begin{proof}
行列式は五個の成分の積からなる。各成分の $a$ 次数が1以上なので、
総次数15に寄与する各因子の次数は11以下である。

Chevalley正規表示でのClifford積は $x$ 次数を下げない。
$\eta_i,\eta_j$ の縮約に現れる
$B_{ij}=\operatorname{Hess}_{ij}F/6$ は $x$ に1次であり、
$\theta_i,\eta_j$ の縮約は $x$ 次数0である。
一方、$h$ は係数を一回微分するので $x$ 次数を1下げる。
従って $A_9$ の $x$ 次数2以下は $R_9$ の3次以下だけから決まり、
ここでは $A_1,\ldots,A_8$ だけを使う。
$A_{10}$ の1次以下を求める際は $R_{10}$ の2次以下が必要であり、
既に得た $A_9$ の2次以下で十分である。
最後に $A_{11}$ の定数項には $R_{11}$ の1次以下だけが必要なので、
$A_{10}$ の1次以下と $A_9$ の1次以下があればよい。
切り捨てた高次の項が後から必要な定数項へ戻ることはない。
\end{proof}

\begin{proposition}[標準点での未正規化十五次係数]
\label{prop:v12-actual-fifteenth-value}
\[
 F_*=x_0^3+x_1^3+x_2^3+x_3^3-4x_4^3,\qquad
 a_*=
 \begin{pmatrix}
 0&1&-1&1&-1\\
 -1&0&1&0&0\\
 1&-1&0&0&0\\
 -1&0&0&0&1\\
 1&0&0&-1&0
 \end{pmatrix}
\]
とする。$q(a_*)=(1,1,1,1,1)$、$P_{F_*}(a_*)=0$ であり、
\[
 \boxed{\displaystyle
 d^{\mathrm T}_{15}(F_*,a_*)=
 \frac{4805924857}{972000}.}
\]
同じ点で
\[
 dP\bigl(\operatorname{Alt}K_{10}\bigr)(F_*,a_*)
 =\frac{461328704861}{181440000}
\]
である。
\end{proposition}

\begin{proof}
\texttt{v12\_fifteenth\_tate.py} は有理数演算でTate再帰を行い、
補題\ref{lem:v12-exact-eleventh-truncation}の切断から
$K_1,\ldots,K_{11}$ と $\det K$ の15次までを求める。
標準点では切断しない $A_1,\ldots,A_{11}$ も別途作り、
各次数でMaurer--Cartan方程式を有理数上で検証した。
この完全な再帰と切断した再帰の十一個のブロックはすべて一致する。
得られる行列式の15次係数は
\[
 [t^{15}]\det K_{F_*}(ta_*)=\frac{4805924857}{486000}.
\]
また $\operatorname{Alt}K_{10}=(K_{10}-K_{10}^{\mathsf t})/2$ を
$dP=dF(q)\circ dq$ に代入すると、二番目の値を得る。

すべてのブロックと有理数値は
\texttt{v12\_fifteenth\_full.json} に収録した。
\texttt{--audit --output v12\_fifteenth\_full.json} は保存した十一個の
ブロックから行列式と上の接触像を有理数で再計算する。
\texttt{--full --indices 0 --output v12\_fifteenth\_full.json} は
外部の計算キャッシュを読まずに完全な再帰を再構成する。
さらに独立に実装した32次元の外積・縮約作用素で、
$x=(1,2,3,4,5)$、法 $1000003$ における各次数0から11の
Maurer--Cartan方程式を照合し、
$x=0$ で各定数ブロックも照合している。
この最後の有限体での照合は実装の独立な整合性検査であり、
表示した数値そのものは有理数上の計算から得ている。
\end{proof}

\Needspace{22\baselineskip}
\begin{theorem}[実際の境界の十五次までの普遍正規形]
\label{thm:v12-fifteenth-normal-form}
普遍的に多項式な座標変更と定数項1の単元により、
\[
 \boxed{\displaystyle
 d^{\mathrm T}\ \sim\
 P_F+\frac1{96}\det\operatorname{Hess}_F(q(a))+O(a^{16}).}
\]
この消去は各固定三次式 $F$ に特殊化して成立し、
接触変換の間に $F$ を動かす必要はない。
\end{theorem}

\begin{proof}
命題\ref{prop:v12-actual-fifteen-correction}により、
第11版の具体的な正規化後の15次項は $R_{15}=d^{\mathrm T}_{15}$ である。

命題\ref{prop:v12-contact-rank-fifteenth}より
$dP:\mathcal G_{10}\to\mathcal C_{15}$ は全射である。
各 $K_n$ の $F$ 次数は $n-1$ なので、行列式の15次係数は
$F$ 次数10の共変量であり、$R_{15}\in\mathcal C_{15}$ である。
従って $R_{15}=dP(\Xi_{10})$ となる普遍共変な
$\Xi_{10}\in\mathcal G_{10}$ を選べる。
既に得た正規形に
\[
 a\longmapsto a-\Xi_{10}(F,a)
\]
を適用すれば15次項が消える。
この変更の $H_{10}$ への最初の寄与は19次であり、
$P$ への二次寄与は24次なので、15次以下の他の項を変えない。
係数は有限次元の普遍共変量空間で解いた定数の線形方程式から得られ、
$P$ や判別式による除算を含まない。
\end{proof}

従って上の非零な数
$4805924857/972000$ は指定したTate座標の係数であり、
十五次の接触不変量ではない。
この次数では $\mathcal Q^{\mathrm{lin}}_{15}
=\mathcal Q^{\mathrm{eff}}_{15}=0$ であるため、
実際の境界の十五次項に限らず、許された普遍共変な十五次項はすべて消去できる。


\clearpage
\section{全次数の偶化の代数性と収束性}
\label{sec:algebraic-parity}

前版の反転対称性は、無限の形式接触変換を与えた。
本節ではその共役を有限方程式で構成し、代数的かつ収束する変更として実現する。
Tate族そのものの有限代数化と、偶化する座標変更の代数化を分けて証明する。


\subsection{偶化に必要な共役を有限個の方程式で構成する}
\label{subsec:v13-algebraic-parity}

第12版の偶化は全次数の形式接触変換として構成した。
ここではその共役を有限次元の方程式へ戻すことにより、
同じ変換が代数的冪級数であり、複素解析的にも収束することを示す。
係数環を三次式の全係数の多項式環とし、$F$ の判別式は逆にしない。

\begin{lemma}[内部重み0の代数と共役方向の補空間]
\label{lem:v13-gauge-weight-spaces}
Chevalley正規表示による内部重み空間について
\[
 \dim\mathsf E_{-3}=0,\qquad
 \dim\mathsf E_{-3/2}=50,\qquad
 \dim\mathsf E_0=1002,\qquad
 \dim\mathsf E_{3/2}=5890
\]
である。$\mathsf E_0$ は単位元を持つ有限自由な結合代数である。
\[
 \mathsf L=\operatorname{im}
       (h:\mathsf E_{3/2}\longrightarrow\mathsf E_0),\qquad
 \mathsf B=\operatorname{im}
       (h:\mathsf E_3\longrightarrow\mathsf E_{3/2})
\]
とおくと
\[
 \dim\mathsf L=951,\qquad \dim\mathsf B=4929,
 \qquad h\delta|_{\mathsf L}=\operatorname{id}_{\mathsf L},
\]
および自然な直和分解
\begin{equation}\label{eq:v13-full-gauge-decomposition}
 \mathsf E_{3/2}
 =A_1(W)\oplus\mathsf B\oplus\delta\mathsf L,
 \qquad 5890=10+4929+951
\end{equation}
が成立する。正規記号による固定した係数空間の上で、
$h,\pi,\delta$ は $F$ に依存しない。
具体的な作用素への同定と積の構造定数は $F$ に多項式的である。
\end{lemma}
\begin{proof}
正規単項式 $x^\alpha\theta_I\eta_{[J]}$ で
$m=|\alpha|$、$u=|I|$、$v=|J|$ とすれば、内部重みの条件は
\[
 2m-u+v=2w.
\]
対応する成分の次元は
$\binom{m+4}{4}\binom5u\binom5v$ である。
$w=0$ では $m=0,1,2$ だけが可能で、その寄与は
\[
 252+600+150=1002.
\]
$w=-3/2$ では $m=0,1$ の寄与が $45+5=50$ であり、
$w=-3$ では成分が存在しない。
$w=3/2$ の和は既存の5890である。
積は内部重みを保つので $\mathsf E_0$ は積について閉じている。
Clifford関係の構造定数は $F$ に多項式的であり、
正規基底は全ての $F$ にわたって自由な基底を与える。
従って $\mathsf E_0$ は上記の有限自由代数である。

負の内部重みでは $\pi=0$ であり、
$\mathsf E_{-3}=0$ とcontractionから
$\delta:\mathsf E_{-3/2}\to\mathsf E_0$ は単射となる。
また内部重み0で $\pi$ が残すのはscalarだけなので
\[
 \ker(\delta|_{\mathsf E_0})
 =\CC\cdot1\oplus\delta\mathsf E_{-3/2}.
\]
従って
\[
 \operatorname{rank}(\delta:\mathsf E_0\to\mathsf E_{3/2})
 =1002-1-50=951.
\]
$\ell\in\mathsf L$ に対しては $h\ell=\pi\ell=0$ なので
$h\delta\ell=\ell$ である。
逆に $\delta h$ は $\operatorname{im}\delta$ 上で恒等写像となるため、
$\delta:\mathsf L\to\operatorname{im}\delta$ は同型である。
よって $\dim\mathsf L=951$ を得る。

$\mathsf E_{3/2}$ にcontractionを適用すると
$T=\pi T+h\delta T+\delta hT$ となる。
$\pi T\in A_1(W)$、$h\delta T\in\mathsf B$、
$\delta hT\in\delta\mathsf L$ であり、
$h^2=\pi h=0$ と $h\delta|_{\mathsf L}=1$ からこの和は直和である。
従って\eqref{eq:v13-full-gauge-decomposition}と
$\dim\mathsf B=5890-10-951=4929$ が従う。
\end{proof}

\paragraph{重み空間の小さい厳密証明書.}
\texttt{v13\_weight\_space\_certificate.py} は通常のPython標準ライブラリだけで、
これらの単項式基底を列挙する。
内部重み $-3/2,0,3/2$ の全6942個の基底記号について
$\delta h+h\delta=1-\pi$、$\delta^2=h^2=0$ を有理数で確認し、
$\delta:\mathsf E_{-3/2}\to\mathsf E_0$ と
$\delta:\mathsf E_0\to\mathsf E_{3/2}$ の階数50、951を
疎な有理数消去法でも直接検証する。
後の代数性・収束性は、この有限計算だけから推論するのではなく、
次の方程式とJacobianの同型性から証明する。

\begin{proposition}[反転したTate族の代数的な切断への帰着]
\label{prop:v13-algebraic-marked-regauging}
Tate族の反転 $\mathsf D_A^*=\mathsf D_0+A(a)^*$ に対して、
第12版の形式写像 $\sigma$ と同じ形式展開を持つ代数的冪級数 $\sigma(a)$、
および代数的冪級数の偶可逆作用素 $g=1+\ell$ が存在し、
\begin{equation}\label{eq:v13-algebraic-intertwining}
 \mathsf D_A^*g=g\bigl(\mathsf D_0+A(\sigma(a))\bigr),
 \qquad \ell\in\mathsf L\widehat\otimes(a),
 \qquad g\equiv1\pmod{(a)}
\end{equation}
を満たす。両者の係数は $F$ に多項式的で、構成は $GL(V)$ と可換である。
各固定 $F$ について、これらの級数は $a=0$ の複素解析的近傍で収束する。
\end{proposition}
\begin{proof}
入力として定理\ref{thm:tate-algebraization}の
4929個の二次方程式によるTate族の\'etale近傍 $\mathcal Z$ を用いる。
その上で $A(a)$ は補助変数を含む代数的な作用素であり、
零切断に沿う形式展開が元のTate再帰と一致する。
反転 $*$ もこの有限次元空間上の代数的な操作である。

追加の未知数を
\[
 b\in W,\qquad z\in\mathsf B,\qquad \ell\in\mathsf L
\]
とし、$C=A_1(b)+z$ と書く。未知数の総数は
$10+4929+951=5890$ である。
内部重み $3/2$ の5890成分の方程式
\begin{equation}\label{eq:v13-finite-regauging-system}
 (\mathsf D_0+A(a)^*)(1+\ell)
       -(1+\ell)(\mathsf D_0+C)=0
\end{equation}
を課す。これは
\[
 \delta\ell+A(a)^*-C+A(a)^*\ell-\ell C=0
\]
という有限個の多項式方程式である。
入力の零切断 $a=0$ では $(b,z,\ell)=(0,0,0)$ が解となり、
未知数に関する微分は
\begin{equation}\label{eq:v13-regauging-jacobian}
 (\dot b,\dot z,\dot\ell)
   \longmapsto\delta\dot\ell-A_1(\dot b)-\dot z.
\end{equation}
補題\ref{lem:v13-gauge-weight-spaces}の直和分解により、これは同型である。
Chevalley正規基底では、この同型は $F$ にも依存しない。
従ってJacobian判定\cite[Proposition 59.26.2(3)]{StacksEtaleCriterion}により、
解空間の零切断を含む開部分は
$\mathcal Z$ 上\'etaleとなる。

$g=1+\ell$ の可逆性も代数的な開条件で確保できる。
有限自由代数 $\mathsf E_0$ 上の左乗算 $L_g$ を取り、
$\det L_g$ を可逆にする。この行列式は零切断上で1である。
$L_g$ が可逆なら $gw=1$ を満たす $w$ があり、
$g(wg-1)=0$ と $L_g$ の単射性から $wg=1$ も従う。
従って $g$ は代数 $\mathsf E_0$ の可逆元となる。

式\eqref{eq:v13-finite-regauging-system}から
$\mathsf D_0+C=g^{-1}\mathsf D_A^*g$ である。
よって $(\mathsf D_0+C)^2=F$ が成立する。
一方、$z\in\mathsf B$ なので $hC=0$、$\pi C=A_1(b)$ であり、
Maurer--Cartan方程式に $h$ を適用すると
\[
 C-A_1(b)+h(C^2)=0
\]
を得る。これは既存のTate切断の有限方程式そのものである。
従って零切断に沿う形式展開では $C=A(b)$ となる。

補題\ref{lem:v12-marked-tate-slice}のパラメータの一意性により、
ここで得られた $b(a)$ は第12版の $\sigma(a)$ と一致する。
特に $\sigma^2=\operatorname{id}$ と
$d\sigma|_{a=0}=-\operatorname{id}$ を継承する。
ここで共役作用素の選択に一意性を課しているのは
$g=1+\ell$、$\ell\in\mathsf L$ という切断の中だけであり、
全てのmarked gauge作用素が一意であるという主張ではない。

上の解空間は有限型の\'etale近傍なので、$b,z,\ell$ は
零切断に沿う代数的冪級数である。
その形式係数は、入力のTate係数と
定数の可逆写像\eqref{eq:v13-regauging-jacobian}を使って次数ごとに求まる。
従って $F$ による除算は生じず、各係数は $F$ の多項式である。
また全ての空間・写像・方程式は同変であり、形式解の一意性から
得られた写像も同変である。

固定した $F$ に対しては、入力のTate枝が収束し、
\eqref{eq:v13-regauging-jacobian}が可逆なので、
複素解析的陰関数定理が同じ方程式に適用できる。
得られる解析解のTaylor展開は一意な形式解と一致する。
従って $\sigma,g$ は実際に原点の近傍で収束する。
\end{proof}

\begin{theorem}[代数的かつ収束する全次数の偶正規形]
\label{thm:v13-algebraic-even-normal-form}
Tate境界には、普遍同変な代数的冪級数による座標変更と単元変更を用いて
\begin{equation}\label{eq:v13-algebraic-even-normal-form}
 \boxed{\displaystyle
 d^{\mathrm T}\ \sim\
 P_F+\frac1{96}H_{10}
       +\sum_{\substack{n\ge16\\n\ \mathrm{even}}}R_n,
 \qquad R_n\in\mathcal C_n
 }
\end{equation}
という全次数の正規形が存在する。
座標変更の線形部は恒等写像、単元の定数項は1であり、
それらの形式係数は $F$ に多項式的である。
各固定 $F$ では、これらの変更は複素解析的にも原点近傍で可逆であり、
右辺の級数も収束する。
\end{theorem}
\begin{proof}
命題\ref{prop:v13-algebraic-marked-regauging}で得た $g$ は、
第12版の規約と同じ
\[
 \mathsf D_0+A(\sigma(a))=g^{-1}\mathsf D_A^*g
\]
を満たす。
内部重み0により $x=0$ では $g$ が各 $\Lambda^kV$ を保存することと、
反転が定数ブロックを $K^t$ にすることから、
\[
 d^{\mathrm T}(F,\sigma(a))
 =v(F,a)d^{\mathrm T}(F,a),\qquad
 v=\frac{\det g_4}{\det g_1},\qquad v(F,0)=1
\]
となる。この $v$ も代数的で収束する単元である。

新しい座標を
\[
 y=\frac{a-\sigma(a)}2
\]
とする。これは代数的冪級数で、線形部が恒等写像なので、
その逆 $a=a(y)$ も代数的冪級数である。
実際、これを定義する写像は零切断上でJacobianが可逆なため、
さらに\'etaleな近傍を取れば逆の枝を代数的に表せる。
同じ理由から、解析的逆関数定理により逆写像も収束する。
$\sigma^2=1$ から $y(\sigma(a))=-y(a)$ となる。

$\widetilde d(F,y)=d^{\mathrm T}(F,a(y))$ と書き、
第12版と同じ平均
\[
 d_{\mathrm{ev}}(F,y)
 =\frac{\widetilde d(F,y)+\widetilde d(F,-y)}2
 =\frac{1+v(F,a(y))}{2}\,\widetilde d(F,y)
\]
を取る。平均に現れる係数は定数項1の代数的・収束的単元であり、
$d_{\mathrm{ev}}$ は全次数で偶関数となる。
代数的冪級数の合成・逆写像・有限和は、必要なら\'etale近傍を共通に細分して
代数的に表せるので、ここまでの構成は全て同じ種類の代数的germの中にある。

系\ref{cor:v12-even-hessian-normal-form}の次数計算により、
この偶化は従来の偶数次係数を14次まで変えない。
従って従来の7次・9次の多項式座標変更と4次・6次・8次の多項式単元で、
10次を $H_{10}/96$ に正規化し、12次と14次を消せる。
これらの座標変更は奇写像、単元は偶関数なので、全次数の偶性も保たれる。
有限な多項式変更は原点で可逆であり、上の代数的・収束的変更との合成も
代数的で収束する。これで\eqref{eq:v13-algebraic-even-normal-form}を得る。
\end{proof}

\begin{remark}[改善された範囲と残る問題]
本節は、第12版で留保した無限偶化の代数性・収束性を解決する。
代数性は零切断近傍の\'etaleなgerm、または代数的冪級数としての意味であり、
$W$ のZariski開部分上の単一の多項式自己同型を主張するものではない。
構成に使った5890個は、既存のTate模型に追加した有限表示の補助未知数であり、
変形パラメータの個数や最小の表示の大きさではない。

本節の議論だけでは、16次以降の偶数次の残余項の消滅は従わない。
次節では16次も消せることを証明するが、18次以上は残る問題である。
境界を全次数で $P_F+H_{10}/96$ だけへ正規化できるかどうかは、依然として未決定である。
\end{remark}


\clearpage
\section{16次の接触商と実際の境界係数}
\label{sec:contact-sixteen}

本節では、15次まで $P_F+H_{10}/96$ に正規化した境界の次の偶数次を調べる。
普遍共変量の次元、接触写像の階数、低次安定化群による追加作用を順に決定する。
実際のTate係数の評価には、12番目の定数ブロックまでを使う。
以下の商は、全ての三次式に多項式的かつ同変に依存する共変量の商である。
固定した $F$ の任意の座標変更によるJacobian商と同一視しない。


\subsection{16次の普遍共変量空間}
\label{subsec:v13-sixteenth-characters}

\begin{proposition}[16次の次元]
\label{prop:v13-sixteenth-characters}
$\dim V=5$、$A=\Lambda^2V$ とする。このとき
\[
 \dim\mathcal C_{16}=40,\qquad
 \dim\mathcal G_{11}=103,\qquad
 \dim\mathcal U_{10}=18,\qquad
 \dim\mathcal D_{12}=59.
\]
また $\delta q:\mathcal G_{11}\to\mathcal D_{12}$ は全射であり、
その核は44次元である。
\end{proposition}
\begin{proof}
Littlewood分解と整数alternantによる重複度計算を用いる。
$\mathcal C_{16}$ に関係する
$\lambda_j=(13,10-j,10-j,j,j)$ の
$\operatorname{Sym}^{11}(\operatorname{Sym}^3V)$ における重複度と、
$\mathcal U_{10}$ に関係する
$\mu_j=(10,10-j,10-j,j,j)$ の
$\operatorname{Sym}^{10}(\operatorname{Sym}^3V)$ における重複度は
\[
\begin{array}{c|rrrrrr}
 j&0&1&2&3&4&5\\\hline
 {}[S_{\lambda_j}V]&2&2&14&10&12&0\\
 {}[S_{\mu_j}V]&1&0&6&5&6&0
\end{array}
\]
である。$\mathcal G_{11}$ と $\mathcal D_{12}$ の次元はそれぞれ
\[
 [(11,11,10,10,10)]
 \bigl(\operatorname{Sym}^{10}(\operatorname{Sym}^3V)
                    \otimes\operatorname{Sym}^{11}A\bigr)=103,
\]
\[
 [(11,11,11,11,10)]
 \bigl(\operatorname{Sym}^{10}(\operatorname{Sym}^3V)
                    \otimes\operatorname{Sym}^{12}A\bigr)=59
\]
から得られる。添付コードは、対称冪の総次元とWeyl次元公式を照合し、
非零のPieri・plethysm成分16行を記録する。
$\delta q$ の全射性は命題\ref{prop:v12-uniform-pfaffian-surjectivity}の
$r=11$ の場合であり、核の次元は $103-59=44$ である。
\end{proof}

\paragraph{Koszul複体の源の次元.}
$g=\nabla F(q)$ による縮約に関係する空間は
\begin{align*}
 \mathcal B_{16}&=\operatorname{Hom}_{\mathrm{GL}(V)}
 \left(\operatorname{Sym}^{9}(\operatorname{Sym}^3V)\otimes\operatorname{Sym}^{8}A,
        \Lambda^2V^*\otimes\det(V)^9\right),\\
 \mathcal K_{3,16}&=\operatorname{Hom}_{\mathrm{GL}(V)}
 \left(\operatorname{Sym}^{8}(\operatorname{Sym}^3V)\otimes\operatorname{Sym}^{4}A,
        \Lambda^3V^*\otimes\det(V)^7\right),\\
 \mathcal K_{4,16}&=\operatorname{Hom}_{\mathrm{GL}(V)}
 \left(\operatorname{Sym}^{7}(\operatorname{Sym}^3V),
        \Lambda^4V^*\otimes\det(V)^5\right)
\end{align*}
であり、同じ整数指標計算から
\[
 \dim\mathcal B_{16}=22,\qquad
 \dim\mathcal K_{3,16}=2,\qquad
 \dim\mathcal K_{4,16}=0
\]
を得る。これらの数の交代和から接触階数を決めることはしない。
実際のKoszul像は、次節で明示的なテンソル縮約によって構成する。
低次安定化群に必要な値も
$\dim\mathcal U_6=2$、$\dim\mathcal G_7=7$、$\dim\mathcal D_8=4$ と再確認できる。

\paragraph{対角三次式への制限では多くの方向が消える.}
Fermat安定化群 $H=\mu_3^5\rtimes\mathfrak S_5$ と
$N=\operatorname{Jac}(F_{\mathrm{Fer}})_3$ を用いる。
$\mathcal C_{16}$ の非零成分について、指標平均による値は
\[
\begin{array}{c|r|r|r}
 \lambda&\text{普遍重複度}&\dim(S_\lambda V)^H
         &\dim\operatorname{Hom}_H(N,S_\lambda V)\\\hline
 (13,10,10,0,0)&2&50&462\\
 (13,9,9,1,1)&2&23&234\\
 (13,8,8,2,2)&14&19&142\\
 (13,7,7,3,3)&10&3&48\\
 (13,6,6,4,4)&12&2&16
\end{array}
\]
となる。各成分の制限階数は普遍重複度と固定次元の小さい方以下なので、
$\mathcal C_{16}$ の対角制限の階数は高々 $2+2+14+3+2=23$ である。
同じ計算による他の空間の上界も合わせると
\[
\begin{array}{c|rrrr}
 &\mathcal C_{16}&\mathcal U_{10}&\mathcal D_{12}&\mathcal G_{11}\\\hline
 \text{普遍次元}&40&18&59&103\\
 \text{対角制限の階数上界}&23&8&30&49\\
 \text{制限核の次元下界}&17&10&29&54
\end{array}
\]
である。従って対角三次式の評価のみでは、今回の普遍接触階数は決まらない。
以下では一般の三次テンソルの整数点を用いる。
これらの表は制限階数の上界であり、法方向の一階微分が実際に各成分を検出することまでは主張しない。

\subsection{一般の三次式を用いた16次の接触階数}
\label{subsec:v13-sixteenth-tensor-rank}

以下で $J_{16}=dP(\mathcal G_{11})$ とおく。
普遍共変量を明示するため、体積形式の基底を選び、
\[
 F(x)=\sum_{i,j,k=0}^4 F_{ijk}x_ix_jx_k,\qquad F_{ijk}=F_{(ijk)},
 \qquad \epsilon_{01234}=1
\]
という対称テンソルの規約を用いる。和は順序付き三つ組全体にわたる。
また、$i,j,k$ が相異なり、$l<m$ がその補集合であるとき
\[
 \tau_{ijk}=\epsilon_{ijklm}a_{lm}
\]
とし、添字が重なる場合は零とする。
この段落では $a\in A=\Lambda^2V$ の表示を用いており、
座標に依存しない表示では
$\tau\in\Lambda^3V^*\otimes\det V$ である。
実際の変形変数 $W=A\otimes\det(V)^{-1}$ に戻すときは、
各 $a$ 因子についてこの定数直線の捻りを戻す。

\paragraph{グラフが定める普遍共変量.}
用いるグラフは、左右に同数の頂点を持つ三価二部グラフである。
各右頂点に並べた三つの左隣接頂点が $\tau$ の添字の順序を指定し、
左頂点における三本の脚は右頂点番号の昇順に並べる。
辺ごとの添字を0から4まで縮約する。
全ての隣接リストと全ての入力は
\texttt{v13\_tensor\_rank\_certificate.json} に記録している。

まず左側に九個の $F$、右側に八個の $\tau$ を置き、
九番目の右頂点の三本の脚だけを自由にする。
得られる三階テンソルを $R^{ijk}$ として
\[
 \beta_{lm}=\sum_{i,j,k}\epsilon_{lmijk}R^{ijk}
\]
と交代化する。この出力は
$\Lambda^3V\otimes\det^8\simeq\Lambda^2V^*\otimes\det^9$ に属し、
$\beta$ は $F$ に9次、$a$ に8次なので $\mathcal B_{16}$ の元である。
従って $g=\nabla F(q)$ に対し
\[
 h=\beta g\in\mathcal D_{12},\qquad g\cdot h=0
\]
が恒等的に成立する。これはKoszul関係の明示的な構成である。

次に左側に十個の $F$、右側に十一個の $\tau$ を置き、
十一番目の左頂点の三本の脚を自由にする。
得られる $R_{ijk}$ から
\begin{equation}\label{eq:v13-graph-coordinate-field}
 \Gamma^{lm}=\frac16\sum_{i,j,k}\epsilon_{lmijk}R_{ijk}
\end{equation}
と定める。この出力は
$\Lambda^3V^*\otimes\det^{11}\simeq\Lambda^2V\otimes\det^{10}$ に属し、
$F^{10}a^{11}$ の次数を持つので $\Gamma\in\mathcal G_{11}$ である。
$b_{lm}=\partial P/\partial a_{lm}$、$l<m$ とし、これを交代行列へ拡張する。
$T_{ijk}=\epsilon_{ijklm}b_{lm}$ と置けば、閉じたグラフの縮約は
\[
 dP(\Gamma)=\frac16\sum_{i,j,k}R_{ijk}T_{ijk}
\]
である。添付コードはこの閉じた縮約だけでなく、
自由脚から\eqref{eq:v13-graph-coordinate-field}の交代行列を取り出して
Pfaffian微分を直接適用する方法とも照合する。

最後に、左に $s$ 個の $F$、右に $s$ 個の $\tau$ を置いた閉じたグラフは、
$F^sa^s$ の $\det(V)^s$ 値共変量である。
$W$ の変数に戻すとスカラー値となり、$\mathcal U_s$ に属する。
以下では $s=6$ の二つのグラフと、既出の四面体グラフを用いる。
四面体の右隣接リストは
\[
 (012),\quad(013),\quad(023),\quad(123)
\]
であり、その縮約を24で割ったものが既出の $U_4$ である。

\begin{proposition}[16次の普遍接触階数]
\label{prop:v13-contact-rank-sixteenth}
\[
 \operatorname{rank}(dP:\mathcal G_{11}\to\mathcal C_{16})=39,
 \qquad \dim\ker dP=64.
\]
さらに、公開データの二十一個のグラフに対応する場を
$\Gamma_1,\ldots,\Gamma_{21}$ とすると
\begin{equation}\label{eq:v13-contact-image-decomposition}
 J_{16}=P\mathcal U_{10}
       \oplus\operatorname{Span}_{\CC}\{dP(\Gamma_1),\ldots,dP(\Gamma_{21})\}.
\end{equation}
$dF:\mathcal D_{12}\to\mathcal C_{16}$ の核は20次元であり、
構成した二十個のKoszulベクトルがその基底を与える。
\end{proposition}
\begin{proof}
上界と下界を独立に証明する。
上の第一の型のグラフを二十個選び、六個の整数入力で得られる
ベクトル $\beta g$ の成分を並べる。
その $20\times20$ 小行列の、素数1009を法とする行列式は
\[
 \det M_{\mathrm K}\equiv930\pmod{1009}
\]
である。従って二十個の普遍Koszulベクトルは独立である。
これらは全て $\ker dF$ に含まれるので、$\dim\mathcal D_{12}=59$ から
\[
 \operatorname{rank}dP\le\operatorname{rank}dF\le59-20=39
\]
を得る。Koszul複体の完全性は用いていない。

下界には、第二の型の二十一個のグラフを用いる。
評価点は $q(a)=(1,1,1,1,1)$ を満たす整数交代行列と整数対称テンソルである。
$F_{444}$ を他の順序付き三つ組の係数和の負数に取ることで、
\[
 P=F(q)=\sum_{i,j,k}F_{ijk}=0
\]
を有理数体上で厳密に満たす。
従ってこれらの評価汎関数は $P\mathcal U_{10}$ を消す。
二十一個の $dP(\Gamma_j)$ の評価行列には
\[
 \det M_{\mathrm G}\equiv175\pmod{1009}
\]
となる $21\times21$ 小行列がある。
全ての入力は既に整数であり、場の規格化に現れる分母6も1009で可逆なので、
これは有理数体上の独立性を証明する非零小行列である。

一方、$\dim\mathcal U_{10}=18$、多項式環は整域であり、Eulerの等式
\[
 dP(au_{10}/6)=Pu_{10}
\]
から $P\mathcal U_{10}$ は接触像に含まれる18次元の部分空間である。
従って $\operatorname{rank}dP\ge21+18=39$ となり、上界と一致する。
これにより直和分解\eqref{eq:v13-contact-image-decomposition}と
$\dim\ker dP=103-39=64$ を得る。
全射 $\delta q$ による因子分解 $dP=dF\circ\delta q$ から、
$\operatorname{rank}dF=39$、$\dim\ker dF=20$ も従う。
既に構成した二十個のベクトルは、この核全体を張る。
\end{proof}

\subsection{16次の補元とHessian単元の所属}
\label{subsec:v13-sixteenth-complement}

既出の八次共変量 $\Theta_8$ は、今回のテンソル規約では
\begin{equation}\label{eq:v13-theta-general-tensor}
 B_{ij}=\sum_kF_{ijk}q_k=\frac16\operatorname{Hess}_F(q)_{ij},
 \qquad
 \Theta_8=\sum_{\substack{|I|=3\,,\ |J|=3}}
          \tau_I\tau_J\det B_{I,J}
\end{equation}
と計算できる。三つ組 $I,J$ は昇順に並べる。
$F$ が対角のときは既出の
$\sum_{i<j}a_{ij}^2\prod_{k\ne i,j}c_kq_k$ に戻る。
また $H=\det\operatorname{Hess}_F(q)=6^5\det B$ とする。

\begin{proposition}[16次の通常接触商の生成元]
\label{prop:v13-sixteenth-complement}
\[
 \mathcal C_{16}=J_{16}\oplus\CC\,\Theta_8U_4^2,
 \qquad
 \mathcal Q^{\mathrm{lin}}_{16}=\CC\,[\Theta_8U_4^2].
\]
\end{proposition}
\begin{proof}
前節と同じ種類の整数入力を二十二個用いる。
二十一個の接触像と $W=\Theta_8U_4^2$ を並べた評価行列は
\[
 \det M_{\mathrm{full}}\equiv457\pmod{1009}
\]
を満たす。従って $W$ は
$P\mathcal U_{10}+\operatorname{Span}\{dP(\Gamma_j)\}$ に属さない。
既に $\dim J_{16}=39$、$\dim\mathcal C_{16}=40$ と分かっているので、
これが求める補元である。

特に、この二十二個の評価汎関数からなる写像を
$L:\mathcal C_{16}\to\CC^{22}$ と書けば、その階数は22であり、
\begin{equation}\label{eq:v13-generic-evaluation-kernel}
 \ker L=P\mathcal U_{10}
\end{equation}
である。この等号は次の所属判定にも用いる。
\end{proof}

\begin{proposition}[Hessianの六次単元倍も通常接触像に入る]
\label{prop:v13-hessian-unit-membership}
\[
 H\mathcal U_6\subset J_{16}.
\]
従って命題\ref{prop:v13-effective-contact-sixteen}の有効な接触商も
\[
 \mathcal Q^{\mathrm{eff}}_{16}
 =\mathcal Q^{\mathrm{lin}}_{16}
 =\CC\,[\Theta_8U_4^2]
\]
である。
\end{proposition}
\begin{proof}
六頂点の二つの閉じたグラフが定める単元を $S_1,S_2$ とする。
その評価行列の $2\times2$ 小行列は法1009で869となるので、
$\dim\mathcal U_6=2$ と合わせて $S_1,S_2$ はこの空間の基底である。

この所属の証明には、有限体上で行列式が零となることを使わない。
前の二十二点における二十一個の接触像、$W$、$HS_1,HS_2$ を、
多倍長整数のテンソル縮約で全て再計算する。
得られた正確な有理数行列で
\[
 L(HS_i)=\sum_{j=1}^{21}c_{ij}L(dP(\Gamma_j))+b_iL(W)
\]
を解くと、
\[
 b_1=b_2=0
\]
が有理数体上で厳密に成立する。
二組の係数 $c_{ij}$ は公開JSONに全て記録されており、検証コードは
二十二行それぞれの等式を正確に検算する。
よって\eqref{eq:v13-generic-evaluation-kernel}から
\[
 HS_i-\sum_{j=1}^{21}c_{ij}dP(\Gamma_j)\in P\mathcal U_{10}\subset J_{16}
\]
となり、主張が従う。
有効商の分母は $J_{16}+H\mathcal U_6$ なので、この分母は $J_{16}$ と一致する。
\end{proof}

\paragraph{証明データの独立監査.}
\texttt{v13\_tensor\_rank\_certificate.json} は、二十個のKoszulグラフ、
二十一個の座標グラフ、二つの六次単元グラフ、全ての入力と小行列、
Hessian単元の有理係数を含む。
入力における三次式の滑らかさは証明の仮定ではない。
非零評価は普遍多項式の独立性を証明する。
多項式族として滑らかな三次式の稠密開集合へ制限しても、この独立性は保たれる。
各三次式を固定した後の独立性を、この計算だけから主張しているわけではない。

探索時の有限体縮約は右頂点で辺を番号付けし、整数の\texttt{einsum}を用いた。
公開された正確な再計算は、左頂点で辺番号を作り直し、
\texttt{NumPy}の\texttt{object}配列による明示的な軸縮約を用いる。
成分はPythonの多倍長整数または有理数であり、浮動小数点演算を用いない。
さらに全ての二十一個の座標グラフについて、第一の一般入力で
自由三脚から交代行列を復元し、Pfaffian微分による $dP$ と閉じた縮約を一致させる。
これにより自由脚の順序、交代化、$1/6$ の規格化を別経路で確認する。

\texttt{v13\_tensor\_rank\_certificate.py} は、正確な等式と非零小行列を検証する。

\texttt{--recompute} では全テンソル縮約を再計算する。
非公開キャッシュや探索用ファイルは読まない。

\subsection{16次では6次単元に伴うHessian方向も作用する}
\label{subsec:v13-filtered-contact}

以下では $P=P_F$、$H=H_{10}$、$c=1/96$ と書く。
第12版で得た15次までの正規形を用い、
\[
 f=P+cH+R_{16}+O(a^{17}),\qquad R_{16}\in\mathcal C_{16}
\]
を考える。$F$、座標の線形項、方程式の先頭係数を固定する。
線形接触像を $J_{16}=dP(\mathcal G_{11})$ とおく。

\begin{lemma}[7次の接触安定化群のHessian作用]
\label{lem:v13-seventh-stabilizer}
$Z\in\mathcal G_7$、$b_6\in\mathcal U_6$ が
\[
 ZP+b_6P=0
\]
を満たすなら
\[
 \delta q(Z)=-\frac13b_6q,\qquad
 ZH+b_6H=-\frac23b_6H.
\]
特に $ZP=0$ なら $ZH=0$ である。
\end{lemma}
\begin{proof}
命題\ref{prop:v10-contact-exact-rank}により
$dF:\mathcal D_8\to\mathcal C_{12}$ は単射である。
また $dF(q)=3P$ なので
\[
 dF\left(\delta q(Z)+\frac13b_6q\right)=ZP+b_6P=0.
\]
括弧の中は $\mathcal D_8$ に属するので零である。
$H$ は $q$ に5次であるから
$ZH=-5b_6H/3$ となり、主張が従う。
ここで用いた単射性は普遍正則な共変量空間上のものであり、
固定した $F$ の任意のベクトル場についての主張ではない。
\end{proof}

\begin{proposition}[16次の有効な接触商]
\label{prop:v13-effective-contact-sixteen}
15次までの $P+cH$ を保つ全ての普遍正則な接触変換を用いた有効な商は
\begin{equation}\label{eq:v13-effective-sixteen-quotient}
 \boxed{\displaystyle
 \mathcal Q^{\mathrm{eff}}_{16}
 =\frac{\mathcal C_{16}}{J_{16}+H\mathcal U_6}.
 }
\end{equation}
従って、この段階で線形接触商 $\mathcal C_{16}/J_{16}$ だけを用いるには、
$H\mathcal U_6\subset J_{16}$ を別に示す必要がある。
\end{proposition}
\begin{proof}
必要なjetを
\begin{align*}
 \chi&=a+\phi_5+\phi_6+\phi_7+\phi_8+\phi_9
                +\phi_{10}+\phi_{11}+O(a^{12}),\\
 u&=1+b_4+b_6+b_8+b_9+b_{10}+O(a^{11})
\end{align*}
と書き、$X=\phi_5$、$Y=\phi_6$、$Z=\phi_7$ とおく。
$\mathcal G_2=\mathcal G_3=\mathcal G_4=0$ および
$\mathcal U_1=\mathcal U_2=\mathcal U_3=\mathcal U_5=\mathcal U_7=0$
を用いた。
15次までを保つ条件のうち、10、11、12次の式は
\[
 XP+b_4P=0,\qquad YP=0,\qquad ZP+b_6P=0
\]
である。13、14、15次の残る条件は、順に
\begin{align*}
 &dP(\phi_8)=0,\\
 &dP(\phi_9)+b_8P+\frac12d^2P(X,X)+b_4XP
                         +c(XH+b_4H)=0,\\
 &dP(\phi_{10})+b_9P+d^2P(X,Y)+b_4YP+cYH=0
\end{align*}
である。$\phi_8,\phi_9,\phi_{10},b_8,b_9$ は
16次の展開には直接寄与しない。

16次の変化を全て書くと
\begin{align*}
 R'_{16}-R_{16}
 ={}&dP(\phi_{11})+b_{10}P
       +d^2P(X,Z)+\frac12d^2P(Y,Y)\\
    &+b_4ZP+b_6XP+c\bigl(ZH+b_6H\bigr).
\end{align*}
実際、$P$ の二次Taylor項では $5+7$ と $6+6$ だけが16次に達する。
$\phi_5$ の三重作用は18次からであり、この次数には現れない。
$H$ の一次Taylor項は $\phi_7$ だけが16次を与え、
単元との交差項は表示した二つだけである。

非線形な $P$ の寄与は高次の接触像に吸収できる。
$ZP=-b_6P$ を $X$ で微分し、$YP=0$ を $Y$ で微分すると
\[
 d^2P(X,Z)=-dP(DZ\cdot X)-(Xb_6)P-b_6XP,
\]
\[
 \frac12d^2P(Y,Y)=-\frac12dP(DY\cdot Y)
\]
を得る。さらに $b_4ZP=-b_4b_6P$ なので
\begin{align*}
 &d^2P(X,Z)+\frac12d^2P(Y,Y)+b_4ZP+b_6XP\\
 &\hspace{1em}=-dP\left(DZ\cdot X+\frac12DY\cdot Y\right)
               -(Xb_6+b_4b_6)P.
\end{align*}
ここで $DZ\cdot X,DY\cdot Y\in\mathcal G_{11}$、
$Xb_6+b_4b_6\in\mathcal U_{10}$ である。
また $P\mathcal U_{10}\subset J_{16}$ は
$Pu_{10}=dP(au_{10}/6)$ から従う。
補題\ref{lem:v13-seventh-stabilizer}を用いると、残る変化は
\[
 R'_{16}-R_{16}\equiv-\frac{2c}{3}b_6H\pmod{J_{16}}.
\]
これで必要性を得る。

逆に $J_{16}$ の方向は11次の座標変更だけで実現できる。
$w_6\in\mathcal U_6$ と $t\in\CC$ に対して
\[
 \chi_t=a+taw_6,\qquad u_t=1-6tw_6
\]
を取ると、$P$ の12次変化は相殺され、次の変化は18次からである。
従ってこの変換は15次までを保ち、16次を
\[
 R_{16}\longmapsto R_{16}+4ctHw_6
\]
と変える。$c\ne0$ なので $H\mathcal U_6$ の全方向を実現できる。
これで十分性も従う。
\end{proof}

\subsection{具体的な有限正規化による16次補正}
\label{subsec:v13-sixteenth-correction}

第10版で実際に固定した規約を用いる。
\[
 X_7=-\alpha\Phi_7-\beta\Psi_7,\qquad
 X_8=\frac{4600}{2189}\Phi_8,
\]
\[
 \alpha=\frac{15800}{6137},\qquad
 \beta=\frac{689855}{883728},\qquad
 \chi=a+X_7+X_8,\qquad u=1+3U_4-V_6.
\]
この規約は、単元を逐次掛けた積を16次まで保持する規約とは異なるので、
命題\ref{prop:v10-explicit-thirteenth}で指定した式をそのまま用いる。

\begin{proposition}[16次の正規化補正式]
\label{prop:v13-actual-sixteen-correction}
上の第10版の変更、第11版の9次座標変更と8次単元、
第12版の15次消去に必要な10次座標変更と9次単元を順に適用する。
得られる16次係数は
\begin{equation}\label{eq:v13-raw-sixteen-correction}
 \boxed{\begin{aligned}
 R_{16}={}&d^{\mathrm T}_{16}+d(d^{\mathrm T}_{10})(X_7)
        +3U_4\bigl(d^{\mathrm T}_{12}+dP(X_7)\bigr)
        -V_6d^{\mathrm T}_{10}.
 \end{aligned}}
\end{equation}
さらに
\begin{equation}\label{eq:v13-sixteen-linear-class}
 R_{16}\equiv d^{\mathrm T}_{16}
                  +c\bigl(dH(X_7)-V_6H\bigr)\pmod{J_{16}},
\end{equation}
従って有効な商では
\begin{equation}\label{eq:v13-sixteen-effective-representative}
 \boxed{\displaystyle
 [R_{16}]
 =\left[d^{\mathrm T}_{16}
       -\frac{\alpha}{96}dH(\Phi_7)
       -\frac{\beta}{96}dH(\Psi_7)\right]
 \quad\text{in }\mathcal Q^{\mathrm{eff}}_{16}.
 }
\end{equation}
\end{proposition}
\begin{proof}
第10版の座標変更の16次への寄与は
$d(d^{\mathrm T}_{10})(X_7)$ だけである。
$X_8$ の寄与には元の9次係数が必要だが零であり、
$P$ への二次Taylor項は18次からである。
単元 $3U_4$ は、元の12次係数だけでなく
$X_7$ による $P$ の12次補正にも掛かる。
単元 $-V_6$ は元の10次係数に掛かる。
これで\eqref{eq:v13-raw-sixteen-correction}を得る。

第11版の9次座標変更と8次単元が16次に作用するには、
いずれも8次係数が必要だが零である。
第12版の10次座標変更と9次単元についても7次係数が必要で、やはり零である。
従ってこれらの後続の変更は16次係数を変えない。

既知の恒等式
\[
 d^{\mathrm T}_{10}=cH-3PU_4,\qquad
 d^{\mathrm T}_{12}+dP(X_7)=PV_6
\]
を代入して整理すると、正確な等式
\[
 R_{16}=d^{\mathrm T}_{16}
       +c\bigl(dH(X_7)-V_6H\bigr)
       +3U_4d^{\mathrm T}_{12}
       +3P\bigl(U_4V_6-X_7U_4\bigr)
\]
となる。
$U_4d^{\mathrm T}_{12}=-dP(U_4X_7)+PU_4V_6$ は $J_{16}$ に入り、
最後の $P$ の倍数も $P\mathcal U_{10}\subset J_{16}$ に入る。
これで\eqref{eq:v13-sixteen-linear-class}が従う。
有効な商では $V_6H\in H\mathcal U_6$ も消えるので、
\eqref{eq:v13-sixteen-effective-representative}を得る。
\end{proof}

\begin{corollary}[偶化の明示計算を必要としない16次判定]
\label{cor:v13-even-normalization-comparison}
第12版の全次数偶正規形における16次係数の有効な接触類は、
\eqref{eq:v13-sixteen-effective-representative}によって計算できる。
従ってこの類を決定するために、形式involution $\sigma$ やその線形化を
16次まで明示計算する必要はない。
\end{corollary}
\begin{proof}
有限正規化による方程式と全次数の偶正規形は、どちらも普遍的な形式接触変換で
同じTate方程式から得られ、15次まで $P+cH$ と一致する。
一方の変換の形式逆を他方と合成すると、15次までを保つ普遍接触変換となる。
命題\ref{prop:v13-effective-contact-sixteen}により、二つの16次係数の差は
$J_{16}+H\mathcal U_6$ に含まれる。
従って有効な接触類は一致する。
\end{proof}

16次の係数そのものを求める際には、$dH(X_7)$ の補正を省けない。
$X_7$ は12次項を消すために選んだベクトル場であり、
$X_7P=0$ という安定化条件を満たさず、標準点でもこの補正は非零となる。
ただし、後の系\ref{cor:v13-seventh-Hessian-in-image}で
$dH(\mathcal G_7)\subset J_{16}$ を証明するので、有効な類ではこの補正も消える。
この消滅は追加の所属判定の結果である。
$V_6$ の点ごとの値は有効な接触類には不要であり、
判定に必要な追加量は $d^{\mathrm T}_{16}$ と、
既存の $\Phi_7,\Psi_7$ に対する $dH(\Phi_7),dH(\Psi_7)$ だけである。

\subsection{第十二Tateブロックと実際の十六次係数}
\label{subsec:v13-actual-sixteenth}

以下の計算では、切断しない $A_{12}$ を作る必要はない。
$x$ 次数の切断を一段延ばすことにより、
十六次の行列式を正確に計算できる。

\begin{lemma}[十六次係数のための正確な切断]
\label{lem:v13-exact-twelfth-truncation}
$d^{\mathrm T}_{16}=\tfrac12[\deg_a16]\det K$ に寄与する定数ブロックは
$K_1,\ldots,K_{12}$ で尽くされる。
完全な $A_1,\ldots,A_8$ から出発して、曲率
$R_n=\sum_{i=1}^{n-1}A_iA_{n-i}$ をそれぞれ
\[
 R_9:\deg_x\le4,\qquad
 R_{10}:\deg_x\le3,\qquad
 R_{11}:\deg_x\le2,\qquad
 R_{12}:\deg_x\le1
\]
で切断し、$A_n=-hR_n$ を順に求めれば、
\[
 A_9:\deg_x\le3,\qquad
 A_{10}:\deg_x\le2,\qquad
 A_{11}:\deg_x\le1,\qquad
 A_{12}:\deg_x=0
\]
を正確に得る。これらで必要な十二個のブロックがすべて決まる。
\end{lemma}

\begin{proof}
五行五列の行列式の各項は、$a$ 次数1以上の成分五個の積である。
従って総次数16に寄与する因子の次数は12以下である。
Clifford積は $x$ 次数を下げず、$h$ は1下げるので、
最後の定数項へ到達できる項だけを逆順に追えば、表示した切断限度を得る。
例えば $A_{12}$ の定数項は $R_{12}$ の1次以下から決まり、
これには $A_{11}$ の1次以下があればよい。
その部分は $R_{11}$ の2次以下から決まり、$A_{10}$ の2次以下で十分である。
同じ議論を続ければ $A_9$ の3次以下と完全な $A_1,\ldots,A_8$ に帰着する。

また $\delta$ は正規表示で $x$ 次数を1上げるので、
曲率の各 $x$ 斉次成分は別々に $\delta$-closed である。
従って表示した切断された曲率と補正項についても、
閉性と対応するMaurer--Cartan等式を有理数上で検証できる。
\end{proof}

\begin{proposition}[標準点での十六次項と低次正規化の補正]
\label{prop:v13-actual-sixteenth-values}
\[
 F_*=x_0^3+x_1^3+x_2^3+x_3^3-4x_4^3,\qquad
 a_*=(e_1-e_0)\wedge(e_2-e_0)
       +(e_3-e_0)\wedge(e_4-e_0)
\]
とする。このとき $q(a_*)=(1,1,1,1,1)$、$P_{F_*}(a_*)=0$ であり、
\[
 \boxed{\displaystyle
 d^{\mathrm T}_{16}(F_*,a_*)
 =\frac{3557764176857}{102060000}.}
\]
第10版の座標ベクトル場
$\Phi_7=\operatorname{Alt}K_7$、
$\Psi_7=\operatorname{Alt}\kappa h(A_1h(A_3^2))$ については、
\[
 dH_{10}(\Phi_7)(F_*,a_*)=\frac{5302368}{25},
 \qquad
 dH_{10}(\Psi_7)(F_*,a_*)=0.
\]
従って
\[
 X_7=-\frac{15800}{6137}\Phi_7
          -\frac{689855}{883728}\Psi_7
\]
とおくと、$dH_{10}(X_7)(F_*,a_*)=-546048$ である。
有効な十六次商で使う代表元
\[
 E_{16}=d^{\mathrm T}_{16}+\frac1{96}dH_{10}(X_7)
\]
の値は
\[
 \boxed{\displaystyle
 E_{16}(F_*,a_*)=\frac{2977246896857}{102060000}.}
\]
\end{proposition}

\begin{proof}
\texttt{v13\_tate\_sixteenth.py} は標準点の $A_1,\ldots,A_8$ を
外部の計算キャッシュを読まずに再構成する。
この部分のMaurer--Cartan等式を有理数で検証し、
独立な32次元の外積・縮約作用素でも照合する。
補題\ref{lem:v13-exact-twelfth-truncation}に従って残り四段を計算し、
各段の切断された曲率の閉性とMaurer--Cartan等式を有理数で検証する。
十二個の定数ブロックは、$x=0$ における独立な外積・縮約作用素でも照合される。
既存の $K_1,\ldots,K_{11}$ および15次までの行列式係数は、
第12版の新規計算による公開データと全成分で一致する。

行列式の新しい係数は
\[
 [t^{16}]\det K_{F_*}(ta_*)=\frac{3557764176857}{51030000}
\]
である。$\Phi_7,\Psi_7$ の行列、$dq$ による像も同じ再帰から求める。
対角三次式 $F=\sum_i c_ix_i^3$ については
\[
 H_{10}=6^5\prod_{i=0}^4c_iq_i,\qquad
 dH_{10}(Y)
 =6^5\sum_{i=0}^4 c_i(dq(Y))_i
                      \prod_{j\ne i}c_jq_j.
\]
この式から表示した二つの微分値を得る。
$\alpha=15800/6137$、$\beta=689855/883728$ を用いると、
\[
 \frac1{96}dH_{10}(X_7)
 =-\frac{\alpha}{96}\frac{5302368}{25}
 =-5688.
\]
これを未正規化の十六次係数に加えれば最後の値を得る。

すべての有理数ブロック、行列式、ベクトル場および補正量は
\texttt{v13\_tate\_sixteenth.json} に収録した。
\texttt{--audit} は保存した十二個のブロックから行列式を再計算し、
Pfaffian微分、Hessianの微分値、補正後の代表元も有理数で再確認する。
既定の実行は新規再帰であり、完全な $A_{12}$ や非公開の入力データを必要としない。
\end{proof}

命題\ref{prop:v13-actual-sixteen-correction}と
\eqref{eq:v13-sixteen-effective-representative}により、
$E_{16}$ は既存の有限正規化後の十六次項の有効類を表す。
上の非零な評価値だけから、その有効類が非零であるとは結論できない。
有効な接触像への所属は、共変量空間と接触写像の階数を使って別に判定する。

\subsection{十六次商の係数を読み取る十二点の汎関数}
\label{subsec:v13-sixteenth-rho}

命題\ref{prop:v13-sixteenth-complement}で得た補元を
$T_{16}=\Theta_8U_4^2$ と書く。
普遍商が一次元であることを使うと、実際の十六次類は
固定した対角三次式の上の有限個の評価だけで決定できる。

$F_*=x_0^3+x_1^3+x_2^3+x_3^3-4x_4^3$、$u_i=e_i-e_0$ とし、
\begin{align*}
 a(b,x,y,z,w)={}&b\,u_1\wedge u_2+x\,u_1\wedge u_3
       +y\,u_1\wedge u_4+z\,u_2\wedge u_3+w\,u_2\wedge u_4\\
 &+\frac{1+xw-yz}{b}\,u_3\wedge u_4
\end{align*}
とおく。$b\ne0$ なら $q(a)=(1,1,1,1,1)$ であり、
$P_{F_*}(a)=0$ となる。
以下の表で十二点 $a_j$ と整数 $N_j$ を指定する。

\begin{center}
\small
\begin{tabular}{r|rrrrr|r}
 $j$&$b$&$x$&$y$&$z$&$w$&$N_j$\\ \hline
 0&1&0&0&0&0&1510767119831929052613\\
 1&2&0&0&0&0&-662247272760455119872\\
 2&3&0&0&0&0&50489862555173579235\\
 3&1&1&0&0&0&78606347649894763400\\
 4&1&1&1&1&0&-114640079068578143080\\
 5&1&1&2&0&1&59225726829655272120\\
 6&1&1&-1&-1&0&363557452143472899000\\
 7&1&-1&1&-1&-1&-474107275671038100\\
 8&1&0&-1&0&-1&-592980195704384042100\\
 9&1&-1&1&-1&1&-54925626926618003400\\
 10&1&0&-1&-1&-1&220089974146168130160\\
 11&1&1&0&0&1&-366444447183161167600
\end{tabular}
\end{center}

\begin{proposition}[十六次類の係数汎関数]
\label{prop:v13-diagonal-rho-sixteenth}
\[
 D=2884048771027944336273600,\qquad
 \rho_{16}(f)=\frac1D\sum_{j=0}^{11}N_j f(F_*,a_j)
\]
とおく。すると
\[
 \rho_{16}(J_{16})=0,\qquad
 \rho_{16}(\Theta_8U_4^2)=1.
\]
従って任意の $f\in\mathcal C_{16}$ について
\[
 [f]=\rho_{16}(f)[\Theta_8U_4^2]
 \quad\text{in }\mathcal Q^{\mathrm{lin}}_{16}
                   =\mathcal Q^{\mathrm{eff}}_{16}.
\]
\end{proposition}

\begin{proof}
全ての点で $P=0$ なので、$\rho_{16}$ は
$P\mathcal U_{10}$ を消す。
命題\ref{prop:v13-contact-rank-sixteenth}の二十一個のグラフによる
ベクトル場を $\Gamma_1,\ldots,\Gamma_{21}$ と書く。
十二点での値を整数および有理数のテンソル縮約で求めると、
\[
 \sum_{j=0}^{11}N_j\,dP(\Gamma_i)(F_*,a_j)=0
 \qquad(1\le i\le21)
\]
をすべて厳密に満たす。
一方、対角三次式上の公式
\[
 \Theta_8=\sum_{i<j}a_{ij}^2\prod_{k\ne i,j}c_kq_k,\qquad
 U_4=\sum_i\prod_{k\ne i}c_k a_{ik}
\]
を使うと
\[
 \sum_{j=0}^{11}N_j(\Theta_8U_4^2)(F_*,a_j)=D
\]
となる。

計算に使った十二行、最初の十一個の接触グラフの列、
および $\Theta_8U_4^2$ の列からなる正方行列の行列式は
\[
 1356890223458353730451986976679070269440000000\ne0
\]
である。
ただし残り十個の接触グラフについても、上の消滅を有理数上で個別に確認している。
命題\ref{prop:v13-contact-rank-sixteenth}の直和分解
\[
 J_{16}=P\mathcal U_{10}\oplus
           \langle dP(\Gamma_1),\ldots,dP(\Gamma_{21})\rangle
\]
により、$\rho_{16}$ は接触像全体を消す。
命題\ref{prop:v13-sixteenth-complement}と
命題\ref{prop:v13-hessian-unit-membership}から、最後の係数公式を得る。
\end{proof}

\paragraph{有理数上での再現方法.}
\texttt{v13\_diagonal\_rho\_certificate.json} は、
十二点、二十一個のグラフ、全評価行列、整数重みを含む。
\texttt{v13\_diagonal\_rho\_certificate.py} の既定動作は、
上の二十二個の線形等式と非零行列式を有理数で確認する。
\texttt{--recompute} は公開されたグラフから全評価値を再構成する。

分数成分を持つ点では共通分母 $d$ を払い、$a'=da$ を使う。
このとき $\tau(a')=d\tau(a)$、
$dP$ の交代勾配テンソルは $d^5$ 倍なので、
十一次ベクトル場の接触像の縮約値は $d^{16}$ 倍される。
従って、任意精度整数による縮約の結果を $d^{16}$ で割れば元の有理数値に戻る。
この実装では有限体上の消滅だけを根拠に所属を推測していない。

\subsection{実際の十六次類は零である}
\label{subsec:v13-actual-sixteenth-zero}

\begin{theorem}[実際の十六次接触類の消滅]
\label{thm:v13-actual-sixteenth-class}
実際のTate方程式について
\[
 \rho_{16}(d^{\mathrm T}_{16})=0,\qquad
 \rho_{16}\bigl(dH_{10}(\Phi_7)\bigr)=0,\qquad
 \rho_{16}\bigl(dH_{10}(\Psi_7)\bigr)=0
\]
が成立する。従って、この三つの共変量はそれぞれ $J_{16}$ に含まれ、
低次正規化後の実際の十六次項は
\[
 \boxed{[R_{16}]=0\quad\text{in }\mathcal Q^{\mathrm{eff}}_{16}}
\]
を満たす。
\end{theorem}

\begin{proof}
命題\ref{prop:v13-diagonal-rho-sixteenth}の十二点について、
補題\ref{lem:v13-exact-twelfth-truncation}の再帰を行う。
標準点 $a_0$ の新規計算済みデータを再利用し、残る十一点では
$A_1,\ldots,A_8$ をそれぞれ新たに作る。
全ての点で有理数上のMaurer--Cartan等式を検証し、
独立な外積作用素と定数ブロックの照合を行った。
この後に $K_9,\ldots,K_{12}$ を正確な $x$ 切断で求め、
行列式の十六次係数を計算する。
同時に既存の二つのベクトル場の $dq$ による像から
$dH_{10}(\Phi_7)$、$dH_{10}(\Psi_7)$ を求める。

各値を表示した整数重みで足すと、有理数上で厳密に
\begin{align*}
 \sum_{j=0}^{11}N_jd^{\mathrm T}_{16}(F_*,a_j)&=0,\\
 \sum_{j=0}^{11}N_jdH_{10}(\Phi_7)(F_*,a_j)&=0,\\
 \sum_{j=0}^{11}N_jdH_{10}(\Psi_7)(F_*,a_j)&=0
\end{align*}
となる。従って
\[
 \rho_{16}\left(
 d^{\mathrm T}_{16}-\frac{\alpha}{96}dH_{10}(\Phi_7)
                   -\frac{\beta}{96}dH_{10}(\Psi_7)\right)=0
\]
でもある。
$\ker\rho_{16}=J_{16}$ は既に普遍的に決定されているので、
最初の三つの共変量は個別に $J_{16}$ に入る。
\eqref{eq:v13-sixteen-effective-representative}より
実際の正規化後の十六次類も零である。

\texttt{v13\_actual\_sixteenth\_certificate.json} は十二点の
全ての定数ブロック、行列式係数、低次正規化のベクトル場、
Hessian微分値および重み付き和を一つにまとめた公開データである。
\texttt{v13\_actual\_sixteenth\_certificate.py} は十二個の行列式を再計算し、
三つの表示した和、補正項の和、補正後の代表元の和が全て零となることを確認する。
従って、この結論は有限体上での零や近似計算による判定ではない。
\end{proof}

\begin{corollary}[七次座標変更のHessian作用も接触像に含まれる]
\label{cor:v13-seventh-Hessian-in-image}
\[
 dH_{10}(\mathcal G_7)\subset J_{16}
\]
である。特に、十六次類に関しては
$dH_{10}(X_7)$ の補正も接触像に吸収できる。
\end{corollary}
\begin{proof}
第10版で示した $dF:\mathcal D_8\to\mathcal C_{12}$ の単射性と
十二次接触像の四次元の基底から、
\[
 \mathcal D_8
 =q\mathcal U_6\oplus
       \mathbb C\,\delta q(\Phi_7)\oplus
       \mathbb C\,\delta q(\Psi_7)
\]
となる。実際、$q\mathcal U_6$ は $dF$ によって
$3P\mathcal U_6$ に写り、残る二つの像と合わせて四次元の接触像を張る。
$H_{10}=h(q)$、$h=\det\operatorname{Hess}F$ は $q$ に五次なので、
$q\mathcal U_6$ のHessian作用は $5H_{10}\mathcal U_6$ に入る。
これは命題\ref{prop:v13-hessian-unit-membership}により $J_{16}$ に含まれる。
残る二つの方向も定理\ref{thm:v13-actual-sixteenth-class}によって
$J_{16}$ に含まれる。
また $\delta q$ の核は $H_{10}=h(q)$ に零作用するので、主張が従う。
\end{proof}

\Needspace{24\baselineskip}
\begin{theorem}[残余項が18次以降の代数的・収束的な偶正規形]
\label{thm:v13-eighteenth-even-normal-form}
実際のTate境界には、普遍同変な代数的冪級数による接触変換の下で
\begin{equation}\label{eq:v13-eighteenth-even-normal-form}
 \boxed{\displaystyle
 d^{\mathrm T}\ \sim\
 P_F+\frac1{96}\det\operatorname{Hess}_F(q(a))
       +\sum_{\substack{n\ge18\\n\ \mathrm{even}}}R_n,
 \qquad R_n\in\mathcal C_n
 }
\end{equation}
という正規形が存在する。
座標変更の線形部は恒等写像、単元の定数項は1であり、
各形式係数は $F$ の係数に多項式的である。
各固定 $F$ では、変更と右辺の級数は原点近傍で収束する。
\end{theorem}
\begin{proof}
定理\ref{thm:v13-algebraic-even-normal-form}の偶正規形から出発する。
系\ref{cor:v13-even-normalization-comparison}と
定理\ref{thm:v13-actual-sixteenth-class}により、その十六次係数は
$J_{16}+H_{10}\mathcal U_6=J_{16}$ に含まれる。
従ってそれを $dP(\Xi_{11})$ と書き、
\[
 a\longmapsto a-\Xi_{11}(F,a),\qquad \Xi_{11}\in\mathcal G_{11}
\]
と変更すれば十六次項が消える。
$\Xi_{11}$ は奇数次なので、この座標変更は奇写像であり、
方程式の全次数における偶性を保つ。
その $H_{10}$ への最初の寄与は20次であり、
既に正規化した15次以下を変えない。
有限な多項式座標変更は原点で可逆であり、
既存の代数的・収束的な変更と合成しても同じ性質を持つ。
これで表示した正規形を得る。
\end{proof}

十六次の普遍商は零ではなく、
$[\Theta_8U_4^2]$ という非零の候補を持つ。
今回確定したのは、実際の境界がこの候補方向に成分を持たないことである。
十八次以降の偶数次項が全て消えるという主張は、まだ得ていない。


\clearpage
% BEGIN characters/v14_character_notes.tex
% 第14版への追加用LaTeX断片。既存本文の定理環境とAMSパッケージを用いる。
\section{18次の普遍共変量と対角制限の限界}
\label{sec:v14-eighteenth-characters}

この節では、18次の接触問題の有限次元の入力を決定する。
以下の次元は整数指標計算による確定値である。
接触写像の階数、Koszul複体の完全性、実際の18次境界係数は、
ここでの計算のみからは決まらない。

\subsection{次数とdeterminant twistの導出}

$\dim V=5$、$A=\Lambda^2V$、$\Delta=\det V$、
$W=A\otimes\Delta^{-1}$ とする。
$F\in\operatorname{Sym}^3V$、$a\in W$ に対し、
$q(a)=\tfrac12a\wedge a\in V^*\otimes\Delta^{-1}$、
$P(F,a)=F(q(a))\in\Delta^{-3}$ である。
中心元 $t\,\mathrm{id}_V$ は $F$、$a$、$P$ にそれぞれ
$t^3$、$t^{-3}$、$t^{-15}$ で作用する。
したがって境界の $a$ 次数 $n$ の項は $F$ 次数 $n-5$ を持ち、
$a$ 次数 $r$ の座標補正は $F$ 次数 $r-1$、
$a$ 次数 $s$ の単元補正は $F$ 次数 $s$ を持つ。
trace線を外した表示は
\begin{align*}
\mathcal C_n&=\operatorname{Hom}_{\mathrm{GL}(V)}
 \left(\operatorname{Sym}^{n-5}(\operatorname{Sym}^3V)
 \otimes\operatorname{Sym}^n A,\Delta^{n-3}\right),\\
\mathcal G_r&=\operatorname{Hom}_{\mathrm{GL}(V)}
 \left(\operatorname{Sym}^{r-1}(\operatorname{Sym}^3V)
 \otimes\operatorname{Sym}^r A,A\otimes\Delta^{r-1}\right),\\
\mathcal D_{r+1}&=\operatorname{Hom}_{\mathrm{GL}(V)}
 \left(\operatorname{Sym}^{r-1}(\operatorname{Sym}^3V)
 \otimes\operatorname{Sym}^{r+1} A,\Lambda^4V\otimes\Delta^{r-1}\right),\\
\mathcal U_s&=\operatorname{Hom}_{\mathrm{GL}(V)}
 \left(\operatorname{Sym}^s(\operatorname{Sym}^3V)
 \otimes\operatorname{Sym}^s A,\Delta^s\right).
\end{align*}
18次では $n=18$、$r=13$、$s=12$ である。

さらに $g=\nabla F(q)$ は、trace線を外した変数では
$F$ 次数1、$a$ 次数4の $V\otimes\Delta^2$ 値共変量となる。
従ってKoszulの第 $k$ 項は
\begin{equation}\label{eq:v14-koszul-general-term}
\mathcal K_{k,n}=
\operatorname{Hom}_{\mathrm{GL}(V)}
\left(\operatorname{Sym}^{n-5-k}(\operatorname{Sym}^3V)
 \otimes\operatorname{Sym}^{n-4k}A,
 \Lambda^kV^*\otimes\Delta^{n-3-2k}\right).
\end{equation}
$\mathcal K_{0,n}=\mathcal C_n$、
$\mathcal K_{1,n}=\mathcal D_{n-4}$、
$\mathcal K_{2,n}=\mathcal B_n$ と一致する。
これは各段で $g$ と縮約すると、$F$ 次数が1、$a$ 次数が4、
determinant指数が2だけ増えることによる。
この次数の確認を経ずに、前の版の添字だけを置き換えることはしない。

\begin{proposition}[18次の次元]\label{prop:v14-eighteenth-dimensions}
次の次元が成り立つ。
\[
\begin{array}{c|rrrrrrr}
 &\mathcal C_{18}&\mathcal G_{13}&\mathcal U_{12}&\mathcal D_{14}
 &\mathcal B_{18}&\mathcal K_{3,18}&\mathcal K_{4,18}\\\hline
\dim&119&372&58&212&112&19&1
\end{array}
\]
\end{proposition}

\begin{proof}
Littlewood分解
\[
\operatorname{Sym}^d A=
\bigoplus_{j=0}^{\lfloor d/2\rfloor}
 S_{(d-j,d-j,j,j,0)}V
\]
を用いる。長方形内の補分割を取ると、
$\mathcal C_{18}$ は
$\operatorname{Sym}^{13}(\operatorname{Sym}^3V)$ の
$\lambda_j=(15,12-j,12-j,j,j)$ 成分、
$\mathcal U_{12}$ は
$\operatorname{Sym}^{12}(\operatorname{Sym}^3V)$ の
$\mu_j=(12,12-j,12-j,j,j)$ 成分の重複度の和となる。
それぞれの重複度は
\[
\begin{array}{c|rrrrrrr}
j&0&1&2&3&4&5&6\\\hline
[S_{\lambda_j}]&3&3&20&27&47&11&8\\
{}[S_{\mu_j}]&2&1&8&12&24&6&5
\end{array}
\]
であり、和は119と58である。

残る五空間も、次の表のhighest weightの重複度として計算する。
\[
\begin{array}{c|rr|c|r}
\text{空間}&F\text{次数}&a\text{次数}&\text{値域のhighest weight}&\dim\\\hline
\mathcal G_{13}&12&13&(13,13,12,12,12)&372\\
\mathcal D_{14}&12&14&(13,13,13,13,12)&212\\
\mathcal B_{18}&11&10&(11,11,11,10,10)&112\\
\mathcal K_{3,18}&10&6&(9,9,8,8,8)&19\\
\mathcal K_{4,18}&9&2&(7,6,6,6,6)&1
\end{array}
\]
各行では入力の中心重み $3\deg_F+2\deg_a$ と、
値域のhighest weightの総和も一致する。

計算証明書の算法を明示する。
35個の三次単項式の重みを $\alpha$ として
\[
 \prod_{|\alpha|=3}(1-tx^\alpha)^{-1}
 =\sum_{m,\beta}M_m(\beta)t^mx^\beta
\]
を整数の動的計画法で計算する。
$\rho=(4,3,2,1,0)$ とすると、重複度は
\[
 [\operatorname{Sym}^m(\operatorname{Sym}^3V):S_\lambda V]
 =\sum_{\sigma\in\mathfrak S_5}
  \operatorname{sgn}(\sigma)\,
  M_m(\lambda+\rho-\sigma\rho)
\]
で与えられる。左辺の記法は $S_\lambda V$ の出現重複度を表す。
一般の二つのcharacterの積に対しても同じ交代式を用いる。

独立な照合として、$h_m=\operatorname{ch}\operatorname{Sym}^m
(\operatorname{Sym}^3V)$ を
\[
 m h_m=\sum_{k=1}^m p_kh_{m-k},\qquad
 p_k=\sum_{|\alpha|=3}x^{k\alpha},\qquad h_0=1
\]
というNewton--Adams再帰でも計算した。
$1\le m\le13$ の全重みで両算法の整数係数が一致し、
$m=13$ では123410個の重みを照合した。
さらに $9\le m\le13$ の全Schur分解について
\[
 \sum_\lambda m_\lambda\dim S_\lambda V=\binom{m+34}{m}
\]
をWeyl次元公式で確認した。
$\mathcal G_{13}$ と $\mathcal D_{14}$ は22行のPieri分解とも一致する。
三つのKoszul源は、別実装の縦strip Pieri算法でも同じ値になる。
計算は全て整数または有理数によるものであり、標本評価や浮動小数点近似を用いない。
\end{proof}

\paragraph{Pfaffian微分についての帰結.}
第12版の全次数Pfaffian全射定理を $r=13$ に適用すると
\[
 \delta q:\mathcal G_{13}\twoheadrightarrow\mathcal D_{14},
 \qquad\operatorname{rank}\delta q=212,\qquad
 \dim\ker\delta q=160
\]
となる。この全射性は指標の一致だけから導いているのではなく、
既存の全次数定理に依存する。
18次接触写像の残る問題は $dF:\mathcal D_{14}\to\mathcal C_{18}$ の階数である。

\paragraph{KoszulのEuler標数は階数を決めない.}
18次の該当成分は
\[
 0\longrightarrow\mathcal K_{4,18}
 \longrightarrow\mathcal K_{3,18}
 \longrightarrow\mathcal B_{18}
 \longrightarrow\mathcal D_{14}
 \longrightarrow\mathcal C_{18}\longrightarrow0
\]
という複体であり、各空間の次元は順に $1,19,112,212,119$ である。
$k=5$ では $a$ 次数が $-2$ となるので項は零である。
従ってEuler標数は
\[
 119-212+112-19+1=1.
\]
この等式から $\dim\mathcal C_{18}/dF(\mathcal D_{14})=1$ とは結論できない。
そのためには、他のホモロジーを別に決定する必要がある。

\subsection{Fermat値と法方向一階微分で検出できない成分}

$F_{\mathrm{Fer}}=x_0^3+\cdots+x_4^3$、
$H=\mu_3^5\rtimes\mathfrak S_5$、
$N=\operatorname{Jac}(F_{\mathrm{Fer}})_3$ とおく。
$N$ は平方因子を持たない三次単項式10個で張られる。
$S_\lambda V$ の普遍出現重複度を $m_\lambda$ とし、
\[
 h_0(\lambda)=\dim(S_\lambda V)^H,\qquad
 h_1(\lambda)=\dim\operatorname{Hom}_H(N,S_\lambda V)
\]
と書く。値の制限階数は $\min(m_\lambda,h_0)$ 以下であり、
値と全法方向一階微分を合わせても $\min(m_\lambda,h_0+h_1)$ 以下である。
Pieri側の重複度がある場合は、この上界にその重複度を掛けて足す。
これは評価写像の実際の階数ではなく、表現論的な上界である。

\begin{proposition}[18次の対角制限と一階jetの上界]
\label{prop:v14-eighteenth-visibility}
次の上界と下界が成り立つ。
\[
\begin{array}{c|rrrr}
&\mathcal C_{18}&\mathcal U_{12}&\mathcal D_{14}&\mathcal G_{13}\\\hline
\text{普遍次元}&119&58&212&372\\
\text{対角値の階数上界}&61&22&93&154\\
\text{対角値の核の次元下界}&58&36&119&218\\
\text{値と法方向一階jetの階数上界}&119&54&208&365\\
\text{同じ一階jetの核の次元下界}&0&4&4&7
\end{array}
\]
特に単元空間・Pfaffian微分の終域・座標補正空間については、
Fermat値と全法方向一階微分だけを用いても普遍空間全体を検出できない。
\end{proposition}

\begin{proof}
有限群の指標平均を $\mathbb Q(\zeta_3)$ 上で正確に計算する。
$\mathcal C_{18}$ に必要な表は
\[
\begin{array}{c|r|r|r}
\lambda&m_\lambda&h_0&h_1\\\hline
(15,12,12,0,0)&3&122&993\\
(15,11,11,1,1)&3&60&629\\
(15,10,10,2,2)&20&52&452\\
(15,9,9,3,3)&27&20&180\\
(15,8,8,4,4)&47&11&79\\
(15,7,7,5,5)&11&1&12\\
(15,6,6,6,6)&8&3&5
\end{array}
\]
となり、値の上界は $3+3+20+20+11+1+3=61$ である。
単元空間については
\[
\begin{array}{c|r|r|r}
\lambda&m_\lambda&h_0&h_1\\\hline
(12,12,12,0,0)&2&18&79\\
(12,11,11,1,1)&1&8&81\\
(12,10,10,2,2)&8&12&97\\
(12,9,9,3,3)&12&5&41\\
(12,8,8,4,4)&24&4&23\\
(12,7,7,5,5)&6&0&3\\
(12,6,6,6,6)&5&2&2
\end{array}
\]
である。最後の二成分には一階jetでもそれぞれ少なくとも3方向と1方向の核がある。
この二成分は $\mathcal D_{14}$ 側にはそれぞれ外側重複度1、1で、
$\mathcal G_{13}$ 側には2、1で現れる。
したがって同じ核の下界はそれぞれ4、7である。
他の成分を含む22行のPieri表と指標平均はJSON証明書に全て保存した。
\end{proof}

\paragraph{18次研究への含意.}
接触像の厳密な階数を決める際には、一般の三次テンソルを含む評価を使うか、
必要な場合は二階以上の法方向jetを加える必要がある。
$\mathcal C_{18}$ の一階jetの上界119は普遍次元と等しいが、
このことだけでは一階jetによる単射性を意味しない。
また $P\mathcal U_{12}\subset dP(\mathcal G_{13})$ はEulerの等式により
58次元の単元方向を与えるが、それだけでは残る接触像や18次商は決まらない。

\subsection{再現用ファイルと確認範囲}

\texttt{v14\_character\_certificate.py} は、原第13版の次の三ファイルを変更せず読み込む。
\begin{itemize}
 \item \texttt{nc\_cubic\_verify\_v7.py}
 \item \texttt{v9\_verify\_contact\_bounds.py}
 \item \texttt{v11\_fermat\_scope\_certificate.py}
\end{itemize}
原コードのEuler積、交代式、有限安定化群の指標平均に加え、
Newton--Adams再帰、全Schur次元、22行のPieri整合性を検査する。
\texttt{v14\_koszul\_pieri\_check.py} は旧Pieri実装を使わず、
行の部分集合を列挙してKoszul源の三つの次元を再検算する。
\texttt{v14\_jet\_visibility\_certificate.py} は上の一階jet上界を再計算する。
新しい計算証明書とログは旧版の証明書から分離して保存する。

本節の結論は、新たに確認した有限次元表現論の命題と、
既存のPfaffian全射定理からの帰結に限る。
実際の18次Tate項の接触像への所属は、後続の節で独立に検証する。
20次以上の全消去は未決定である。

% END characters/v14_character_notes.tex


% BEGIN rank/rank18_section_ja.tex
\section{十八次の通常接触商}
\label{sec:v14-ordinary-contact}

\subsection{18次の通常接触階数と一つの補元}
\label{subsec:v14-eighteenth-contact-rank}

以下では $P=P_F=F(q(a))$、$g=\nabla F(q(a))$ とし、
$J_{18}=dP(\mathcal G_{13})$ とおく。
前節の整数指標計算から
\[
 \dim\mathcal C_{18}=119,\quad
 \dim\mathcal G_{13}=372,\quad
 \dim\mathcal D_{14}=212,\quad
 \dim\mathcal U_{12}=58
\]
を用いる。また命題\ref{prop:v12-uniform-pfaffian-surjectivity}により
$\delta q:\mathcal G_{13}\twoheadrightarrow\mathcal D_{14}$ は全射である。
本小節は、Koszul複体の完全性を仮定せず、像と核の有限証明書を独立に構成する。

\begin{proposition}[18次の通常接触階数]
\label{prop:v14-contact-rank-eighteenth}
普遍共変量の空間上で
\[
 \operatorname{rank}(dP:\mathcal G_{13}\longrightarrow\mathcal C_{18})=118,
 \qquad \dim\ker dP=254.
\]
さらに
\[
 \operatorname{rank}(dF:\mathcal D_{14}\longrightarrow\mathcal C_{18})=118,
 \qquad \dim\ker dF=94.
\]
公開データの六十個の座標場を $\Gamma_1,\ldots,\Gamma_{60}$ とすると
\[
 J_{18}=P\mathcal U_{12}\oplus
    \operatorname{Span}_{\CC}\{dP(\Gamma_1),\ldots,dP(\Gamma_{60})\}.
\]
また、九十四個の明示的なKoszulベクトルが $\ker dF$ の基底となる。
\end{proposition}

\begin{proof}
\textbf{普遍的なkernel方向の構成。}
第\ref{subsec:v13-sixteenth-tensor-rank}節と同じテンソル規約を用いる。
$F_{ijk}=F_{(ijk)}$、$\tau_{ijk}=\epsilon_{ijklm}a_{lm}$ とし、
$l<m$ は $i,j,k$ の補集合とする。
三価二部グラフの左に十一個の $F$、右に十個の $\tau$ を置き、
第十一右頂点の三本の脚を自由にする。
その三階テンソルを $R^{ijk}$ として
\[
 \beta_{lm}=\sum_{i,j,k}\epsilon_{lmijk}R^{ijk},\qquad h=\beta g
\]
と定める。$\beta$ は $F$ に十一、$a$ に十の次数を持ち、
$\mathcal B_{18}$ の元となる。
交代性から $g\cdot h=g^{\mathrm t}\beta g=0$ が恒等的に成立するので、
$h\in\ker(dF:\mathcal D_{14}\to\mathcal C_{18})$ である。

この方法で構成した九十四個の $h$ を三十個の整数入力で評価し、
五成分を並べた行列には、素数1009を法として
\[
                 \det M_{\mathrm K}\equiv7\pmod{1009}
\]
となる $94\times94$ 小行列がある。
全グラフ・全入力・選択行は
\texttt{koszul18\_search.json} に、
有理数上で再縮約した小行列の全成分は
\texttt{koszul18\_exact\_minor.json} に記録した。
従って九十四個の普遍ベクトルは特性零で独立であり、
\[
                         \operatorname{rank}dF\le212-94=118.
\]

\textbf{接触像の下界。}
今度は左に十二個の $F$、右に十三個の $\tau$ を置き、
第十三左頂点の三本の脚を自由にする。
得られる $R_{ijk}$ から
\[
              \Gamma^{lm}=\frac16\sum_{i,j,k}\epsilon_{lmijk}R_{ijk}
\]
と定める。これは $F^{12}a^{13}$ の次数を持つ
$\mathcal G_{13}$ の元である。
$T_{ijk}=\epsilon_{ijklm}(\partial P/\partial a_{lm})$ と置けば、
閉じたグラフの縮約
\[
                         \frac16\sum_{i,j,k}R_{ijk}T_{ijk}
\]
が $dP(\Gamma)$ を与える。

六十二個の整数入力を
\[
 q(a)=(1,1,1,1,1),\qquad
 P=\sum_{i,j,k}F_{ijk}=0
\]
となるように取る。この等式は有限体だけでなく有理数体上で成立する。
これらの評価は $P\mathcal U_{12}$ を消す。
選んだ六十個の場に対する評価行列には
\[
                    \det M_{\mathrm G}\equiv889\pmod{1009}
\]
となる $60\times60$ 小行列がある。
全ての縮約は整数演算であり、場の正規化の唯一の分母6は1009で可逆である。
従って、これは特性零での、$P\mathcal U_{12}$ を法とする独立性の証明書となる。
データと有理数による再縮約は
\texttt{contact18\_search.json} と
\texttt{contact18\_exact\_minor.json} に記録した。

一方、Eulerの関係
\[
                         dP(au/6)=Pu
\]
から $P\mathcal U_{12}\subset J_{18}$ である。
多項式環は整域なので、この部分空間は58次元である。
よって $\operatorname{rank}dP\ge60+58=118$。
全射 $\delta q$ による因子分解 $dP=dF\circ\delta q$ と合わせると
上下界が一致する。核の次元と、構成したKoszul像が全kernelを尽くすことも従う。
\end{proof}

\begin{remark}[有限体の役割と独立な再検算]
探索に擬似乱数を用いたが、採用したグラフと整数入力は全て固定して公開している。
非零の小行列式を有限体で確認する操作は、有理数体上の独立性を証明するものであり、
確率的な階数の推測ではない。
発見時の縮約は左頂点で辺を番号付けし、明示的な軸の縮約を用いた。
独立監査では右頂点で番号付けし、整数 \texttt{einsum} によって全入力を再計算した。
各縮約の和の上界を確認して整数のoverflowを排除し、演算ごとに1009で還元している。
両方法で小行列式 $7,889$ を再現し、さらに小行列の全成分を
多倍長整数・有理数で再計算した。
\end{remark}

\begin{proposition}[18次の通常商]
\label{prop:v14-ordinary-eighteenth-quotient}
第\ref{subsec:v13-sixteenth-tensor-rank}節の四面体unit $U_4$ と、
第28節の六頂点unit $J_6$ を用いると
\[
                   \mathcal C_{18}/J_{18}
                         =\CC\cdot[\Theta_8U_4J_6].
\]
\end{proposition}
\begin{proof}
$\Theta_8U_4J_6$ は $F^{13}a^{18}$ の次数を持つので
$\mathcal C_{18}$ に属する。
上の六十列にこの共変量の評価列を加えると、
$61\times61$ 小行列式が1009を法として948となる。
従って、この共変量は $J_{18}$ に属さない。
前命題と $\dim\mathcal C_{18}=119$ から商は一次元である。
$\Theta_8,U_4,J_6$ の六十二点での値は全て有理数で再計算し、
\texttt{complement18\_certificate.json} に記録した。
\end{proof}

\subsection{有限評価が普遍所属を判定する理由}
\label{subsec:v14-evaluation-and-divisibility}

後で接触像への所属を評価値から判定するため、
$P\mathcal U_{12}$ で割った商の大きさを確認する。
三次式の係数を $F_\alpha$ として
\[
 R=\CC[a_{ij},F_\alpha],\qquad
 P=\sum_{|\alpha|=3}F_\alpha q(a)^\alpha
\]
と書く。テンソル係数の規約で現れる非零の整数係数は、以下の議論を変えない。
Pl\"uckerイデアルは高さ3なので、五つの $q_i$ の共通因子は1である。
$q_i^3$ を含む上の係数のcontentも1であり、
$P$ は $\CC[a]$ 上primitiveである。
$\CC(a)[F_\alpha]$ 上では非零一次式として既約なので、Gaussの補題により
$P$ は $R$ でも既約となる。

従って $\mathcal C_{18}$ の元が $P=0$ の全スキーム上で消えれば、
その元は $P$ で割り切れる。
商の $F,a$ に関する次数はともに12であり、
$P$ と元の式の変換則を割るとunitの変換則を得る。
よってその商は $\mathcal U_{12}$ に属する。
このことから、制限の核はちょうど $P\mathcal U_{12}$ であり、
\[
                         \dim\bigl(\mathcal C_{18}/P\mathcal U_{12}\bigr)=61.
\]
前小節の六十個の接触像と一つの補元は、この商全体の基底を与える。
そのため、この六十像を消し、$\Theta_8U_4J_6$ を1へ送る
\emph{有理数上の}有限評価汎関数は、通常接触商の係数を正確に読み取る。
ある追加列が1009を法として像に入るという観察だけでは、
特性零での所属を結論しない。

% END rank/rank18_section_ja.tex


% BEGIN contact/contact_log_reduction_ja.tex
% 第14版追加候補。第13版原稿のプリアンブル・記号・既存ラベルを用いる。
% 作成日：2026-10-29。原稿内の既存命題を入力として明記する。

\section{接触変換の対数と18次の有効作用}
\label{sec:v14-log-contact}

本節では、低次の正規形を保つ接触変換を、接触作用素の対数によって
有限個の線形条件へ帰着する。これにより18次では、通常のJacobian像に
加えて調べるべき方向を、Hessianの八次単元倍と一つの追加方向に絞れる。
この段階では、追加方向が通常のJacobian像に含まれるとは仮定しない。

以下、$P=P_F=F(q(a))$、$H=H_{10}=\det\operatorname{Hess}_F(q(a))$、
$c=1/96$ と書く。$F$ は接触変換の間に固定し、全ての係数は
$F$ に多項式的かつ $\mathrm{GL}(V)$ 同変とする。
座標変更の線形項は恒等写像、単元の定数項は1に固定する。
$\mathcal G_r,\mathcal U_s,\mathcal C_n$ の規約は既存の節と同じである。

\subsection{全次数で使えるfilteredな対数補題}

\begin{lemma}[接触作用素の対数]
\label{lem:v14-contact-logarithm}
$R$ を $\mathbb Q$-代数、$A=R[[a_1,\ldots,a_m]]$、
$\mathfrak m=(a_1,\ldots,a_m)$ とする。
$\chi=a+O(\mathfrak m^{s+1})$、$u=1+O(\mathfrak m^s)$、$s\ge1$ に対し
\[
 T(f)=u\,f\circ\chi
\]
と置く。このとき $\log T$ は一意に
\[
 L=Y+b
\]
と書ける。ここで $Y$ は連続な $R$-導分、$Y(a_i)\in\mathfrak m^{s+1}$、
$b\in\mathfrak m^s$ は乗算作用素である。
逆にこの形の $L$ の指数は、上と同じ型の接触作用素を与える。

$f\in A$、$N\ge0$ に対し
\[
 T(f)-f\in\mathfrak m^N\quad\Longleftrightarrow\quad
 L(f)\in\mathfrak m^N.
\]
これらが成立するとき
\[
 T(f)-f\equiv L(f)\pmod{\mathfrak m^{N+s}}.
\]
従って、指定された次数未満のjetを保つ条件と次の次数における作用は、
$L(f)$ の斉次成分を用いる線形条件で完全に記述できる。
\end{lemma}

\begin{proof}
$T-1$ は $\mathfrak m$-進次数を少なくとも $s$ 上げるので、
$\log T=\sum_{j\ge1}(-1)^{j+1}(T-1)^j/j$ は収束する。
対数が一階作用素であることを確認する。
各有限jet上で
\[
 T^t=\sum_{j\ge0}\binom tj(T-1)^j
\]
は $t$ の多項式である。非負整数 $t$ では $T^t$ は接触作用素の反復なので
\[
 T^t(fg)\,T^t(1)=T^t(f)\,T^t(g)
\]
が成立する。有限jetごとの多項式恒等式として全ての $t$ で成立し、
$t=0$ で微分すると
\[
 L(fg)+fgL(1)=L(f)g+fL(g)
\]
を得る。従って $b=L(1)$、$Y=L-b$ と置けば $Y$ は導分である。
次数の条件も $T-1$ の条件から従い、$b=L(1)$ により分解は一意である。

逆に $L=Y+b$ について、$u_t=\exp(tL)(1)$ と置く。
微分方程式を比較すると
\[
 \exp(tL)(f)=u_t\,\exp(tY)(f)
\]
である。$\exp(tY)$ は形式自己同型だから、$t=1$ で接触作用素を得る。
これらの構成では有限jetごとに有限回の和・積と有理数による除算だけを用いる。

最後に、$T(f)-f\in\mathfrak m^N$ なら対数の第2項以降は
$\mathfrak m^{N+s}$ に入る。逆向きは指数級数に同じ議論を適用する。
これで同値と合同式を得る。
\end{proof}

\begin{remark}[単元とその対数の規約]
$L$ の乗算成分 $b$ は、一般には $\log u$ そのものではない。
座標変更と単元の合成が対数に含まれるためである。
以下の $b_k$ は常に $\log T$ の乗算成分を表し、
具体的な有限正規化の単元の斉次成分とは区別する。
この区別を保てば、mergedな単元表示も逐次合成も同じ補題で扱える。
\end{remark}

\begin{corollary}[次のjetの有効作用は線形条件の核からの像である]
\label{cor:v14-filtered-contact-linearization}
$L=\sum_{k\ge s}Q_k$ と内部の $a$ 次数増分で分解し、
$Q_k=Y_{k+1}+b_k$ とする。$f=\sum_j f_j$ に対し、
$N$ 次未満を保つ接触変換が $N$ 次に与える差の集合は、
\[
 \sum_{k+j=n}Q_kf_j=0\quad(n<N)
\]
という有限個の線形条件を満たす $Q_k$ による
\[
 \sum_{k+j=N}Q_kf_j
\]
の像である。特にこれは線形部分空間である。
普遍共変量の場合、対数・指数の各係数も $F$ に多項式的で同変である。
必要な有限jetは、多項式の座標変更と単元で実現できる。
\end{corollary}

\begin{proof}
補題を $N$ 次まで適用する。$Q_k$ は指定次数を $k$ 上げるので、
表示された条件はちょうど $L(f)$ の $N$ 次未満の消滅である。
各有限次数での対数・指数は有限和であり、同変性と係数の多項式性を保つ。
最後に座標変更と単元の不要な高次項を切り捨てても、対象の有限jetは変わらない。
\end{proof}

\subsection{18次での全接触作用と偶性を保つ作用の一致}

既存の低次計算により $\mathcal G_2=\mathcal G_3=\mathcal G_4=0$、
$\mathcal U_1=\mathcal U_2=\mathcal U_3=0$ なので、最小の次数増分は4である。
$Q_k\in\mathcal G_{k+1}\oplus\mathcal U_k$ と書く。
$f=P+cH+R_{18}+O(a^{19})$ の17次までを保つ条件は、対数を用いると
\[
 Q_kP+cQ_{k-4}H=0\quad(4\le k\le11),
 \qquad Q_j=0\ (j<4)
\]
であり、18次の変化は
\[
 Q_{12}P+cQ_8H
\]
である。
偶数の次数増分だけを並べると
\begin{equation}\label{eq:v14-even-stabilizer-chain}
 \begin{split}
 Q_4P&=0,\qquad Q_6P=0,\\
 Q_8P+cQ_4H&=0,\qquad Q_{10}P+cQ_6H=0.
 \end{split}
\end{equation}
奇数の増分についても独立に
\[
 Q_5P=Q_7P=0,\qquad
 Q_9P+cQ_5H=Q_{11}P+cQ_7H=0
\]
となる。

\begin{proposition}[18次では偶性保存への制限で有効像が変わらない]
\label{prop:v14-parity-restriction}
17次までの $P+cH$ を保つ全ての普遍接触変換の18次像は、
奇写像の座標変更と偶関数の単元だけを許した18次像と一致する。
\end{proposition}

\begin{proof}
$P,H$ はともに偶関数であり、$Q_k$ は次数を $k$ 増やすため、
上の低次条件は偶増分と奇増分で分離する。18次の変化には
$Q_{12},Q_8$ だけが現れる。従って任意の解の奇増分成分を捨てても、
低次条件と18次作用は変わらない。
偶増分だけを持つ $L$ は $a\mapsto-a$ と可換する。
その指数から得られる座標変更は奇写像、単元は偶関数である。
逆の包含は明らかである。
\end{proof}

この証明は、任意の次数で同変接触分類と非同変接触分類が一致するという主張ではない。
ここでは $F$ に普遍同変な変換の範囲を保ち、$P+cH$ の18次作用だけを扱っている。
同じ対数の偶奇分離は、他の偶数次でも下の正規形が偶関数であれば利用できる。

\subsection{九次の接触安定化場のHessian作用}

$J_{18}=dP(\mathcal G_{13})\subset\mathcal C_{18}$ と置く。
Eulerの恒等式により $P\mathcal U_{12}\subset J_{18}$ である。
以下で用いる重要な既存入力は、第28節で証明した
\[
 \ker(dF:\mathcal D_{10}\to\mathcal C_{14})
       =\{\beta\nabla F(q):\beta\in\mathcal B_{14}\}
\]
と、第30節のPfaffian微分の全次数の全射性である。
このKoszul像による核の記述は、当該14次成分だけの結論である。

\begin{lemma}[九次安定化場の作用]
\label{lem:v14-ninth-stabilizer-hessian}
$S_8=Y_9+b_8$、$Y_9\in\mathcal G_9$、$b_8\in\mathcal U_8$ が
$S_8P=0$ を満たすなら
\[
 S_8H\equiv-\frac23b_8H\pmod{J_{18}}.
\]
特に $dP(Y_9)=0$ なら $dH(Y_9)\in J_{18}$ である。
\end{lemma}

\begin{proof}
$g=\nabla F(q)$、$h(p)=\det\operatorname{Hess}_F(p)$ と置く。
$dF(q)=3P$ より
\[
 dF\left(\delta q(Y_9)+\frac13b_8q\right)=0.
\]
括弧内は $\mathcal D_{10}$ に属するため、上記の既存核の記述により
\[
 \delta q(Y_9)=-\frac13b_8q+\beta g
\]
となる普遍交代共変量 $\beta\in\mathcal B_{14}$ が存在する。
$h$ は $q$ に5次なので
\[
 S_8H=-\frac23b_8H+(\nabla h(q))^t\beta g
     =-\frac23b_8H-g^t\beta\nabla h(q).
\]
$\beta\nabla h(q)$ は $F$ に12次、$a$ に14次の $\mathcal D_{14}$ の元である。
$\delta q:\mathcal G_{13}\to\mathcal D_{14}$ は全射なので、この元を持ち上げれば
最後の項は $J_{18}$ に入る。
\end{proof}

\subsection{18次の有効像の正確な記述}

第28節の所属 $HU_4\in dP(\mathcal G_9)$ により、
\[
 Z_9\in\mathcal G_9,\qquad dP(Z_9)=HU_4
\]
を一つ選ぶ。ここで $U_4$ は第27節の規格化を持つ四次単元である。

\begin{theorem}[18次有効商の有限な縮約]
\label{thm:v14-effective-eighteenth}
$c\ne0$ とする。17次までの正規形 $P+cH$ を保つ全ての普遍接触変換の
18次像は正確に
\begin{equation}\label{eq:v14-effective-image}
 I_{18}=J_{18}+H\mathcal U_8+\CC\,dH(Z_9)
\end{equation}
であり、従って
\[
 \mathcal Q^{\mathrm{eff}}_{18}
 =\mathcal C_{18}/\bigl(J_{18}+H\mathcal U_8+\CC dH(Z_9)\bigr).
\]
$dH(Z_9)$ の $\mathcal C_{18}/J_{18}$ における類は、
$dP(Z_9)=HU_4$ を満たす $Z_9$ の選択に依存しない。
この式は $H\mathcal U_8$ や $dH(Z_9)$ の $J_{18}$ への所属を主張しない。
\end{theorem}

\begin{proof}
第24節の低次計算により、$Y_5\in\mathcal G_5$ のPfaffian像は
$\delta q(Y_5)=\lambda qU_4$ の形である。
$Q_4=Y_5+b_4$、$Q_4P=0$ なら $b_4=-3\lambda U_4$ であり、
\[
 Q_4H=2\lambda HU_4.
\]
$t=\lambda/2$ と置けば $Q_4H=4tHU_4$ となる。
この値は $Q_4=tU_4(\mathcal E-6)$ で実現される。
ここで $\mathcal E=\sum a_i\partial_{a_i}$ はEuler導分である。
$\ker\delta q$ の五次場は $P,H$ の両方に零作用する。

\eqref{eq:v14-even-stabilizer-chain}により
$Q_8P=-4ctHU_4$ だから
\[
 Q_8=-4ctZ_9+S_8,\qquad S_8P=0
\]
と書ける。従って18次作用は
\[
 Q_{12}P+cQ_8H
 \in J_{18}+H\mathcal U_8+\CC dH(Z_9)
\]
となる。ここで $Q_{12}P\in J_{18}$ はEulerの恒等式、
$S_8H\in J_{18}+H\mathcal U_8$ は前補題による。
これで一方の包含を得る。

逆に $J_{18}$ は13次の座標場だけで実現できる。
任意の $w_8\in\mathcal U_8$ に対し
$Q_8=w_8(\mathcal E-6)$、他の成分を零とすれば、
全低次条件を満たし、18次作用は $4cHw_8$ である。
最後に
\[
 Q_4=tU_4(\mathcal E-6),\qquad Q_8=-4ctZ_9,
 \qquad Q_6=Q_{10}=0
\]
とすれば全低次条件を満たし、18次作用は $-4c^2t\,dH(Z_9)$ である。
$c\ne0$ なので、表示された全方向が実現される。
対数補題により、これらは実際の有限jetの接触変換を与える。

二つの選択 $Z_9,Z'_9$ の差は $dP$ の核に入り、
補題\ref{lem:v14-ninth-stabilizer-hessian}の $b_8=0$ の場合から
$dH(Z_9-Z'_9)\in J_{18}$ となる。
\end{proof}

\begin{remark}[グラフによる計算に適した持ち上げ]
実装では、純粋な持ち上げの代わりに
\[
 HU_4=dP(Z'_9)+PB_8
\]
という分解を求めればよい。このとき
$Z_9=Z'_9+aB_8/6$ は純粋な持ち上げであり、
\[
 dH(Z_9)-dH(Z'_9)=\frac53HB_8\in H\mathcal U_8.
\]
従って\eqref{eq:v14-effective-image}の $dH(Z_9)$ を
$dH(Z'_9)$ に置き換えても像は変わらない。
既存の $J_{14}/P\mathcal U_8$ は7次元なので、
$P=0$ 上で独立な七つの九次グラフの接触像を用いれば、
七元の線形方程式から $Z'_9$ を構成できる。
\end{remark}

\begin{proposition}[八次単元のうち $P$ の倍数となる部分]
\label{prop:v14-pdivisible-eighth-unit}
$\mathcal U_8$ のうち $P$ で割り切れる部分は1次元であり、
この部分の $H$ 倍は $J_{18}$ に含まれる。
\end{proposition}

\begin{proof}
商の共変量は
\[
 \operatorname{Hom}_{GL(V)}\left(
 \operatorname{Sym}^7(\operatorname{Sym}^3V)
 \otimes\operatorname{Sym}^2(\Lambda^2V),\det(V)^5\right)
\]
に属する。Littlewood分解による二つの補分割は
$(5,5,5,3,3)$ と $(5,4,4,4,4)$ である。
既存の整数指標計算をこの二つに適用すると、重複度はそれぞれ1と0となる。
従ってこの空間は1次元である。
多項式環は整域なので $P$ 倍は単射であり、商の共変性は
$P$ と元の式の共変性から従う。
$u_8=Pv$ と書けば、$Hv\in\mathcal U_{12}$ だから
$Hu_8=P(Hv)\in P\mathcal U_{12}\subset J_{18}$ である。
\end{proof}

\section{十八次の実際の係数に入る低次正規化の補正}
\label{sec:v14-eighteenth-corrections}

ここでは第27節のmerged規約をそのまま固定する。
\[
 X_7=-\alpha\Phi_7-\beta\Psi_7,
 \quad\alpha=\frac{15800}{6137},\quad
 \beta=\frac{689855}{883728},\qquad
 X_8=\frac{4600}{2189}\operatorname{Alt}K_8,
\]
\[
 \chi_{10}=a+X_7+X_8,\qquad u_{10}=1+3U_4-V_6.
\]
$u_{10}$ を $(1+3U_4)(1-V_6)$ に交換してはならない。
$d^{\mathrm T}_n$ は未正規化の $\tfrac12\det K$ の斉次成分とし、
\[
 d^{\mathrm T}_{10}=cH-3PU_4,
 \qquad d^{\mathrm T}_{12}+dP(X_7)=PV_6
\]
を用いる。
第28節の十四次消去を
\[
 R_{14}=d^{\mathrm T}_{14}+3U_4d^{\mathrm T}_{10}
       =dP(\Xi_9)+PB_8,
 \qquad\chi_{11}=a-\Xi_9,\quad u_{11}=1-B_8
\]
という規約で行う。

\begin{proposition}[十八次の正確な補正式]
\label{prop:v14-raw-eighteen-correction}
上記二段階の変更の後の18次係数は正確に
\begin{equation}\label{eq:v14-raw-eighteen-correction}
 \begin{split}
 R_{18}={}&d^{\mathrm T}_{18}
   +d(d^{\mathrm T}_{12})(X_7)
   +\frac12d^2P(X_7,X_7)
   +3U_4d^{\mathrm T}_{14}-PV_6^2\\
   &\hspace{2em}-c\bigl(dH(\Xi_9)+B_8H\bigr).
 \end{split}
\end{equation}
その後に行う十五次消去の10次座標変更・9次単元、および
十六次消去の11次座標変更・10次単元は、18次係数を変えない。
ここで $d^2P(X_7,X_7)$ はHessianの双線形評価であり、
変数に依存するベクトル場の反復 $X_7^2P$ とは区別する。
\end{proposition}

\begin{proof}
最初の引き戻しの18次には、
$d^{\mathrm T}_{18}$、$d(d^{\mathrm T}_{12})(X_7)$、
$\tfrac12d^2P(X_7,X_7)$ のみが寄与する。
$X_8$ の一次寄与には $d^{\mathrm T}_{11}$ が必要だが零であり、
$X_7,X_8$ の混合二次項は19次からである。
$3U_4$ が掛かる14次係数は $d^{\mathrm T}_{14}$、
$-V_6$ が掛かる12次係数は
$d^{\mathrm T}_{12}+dP(X_7)=PV_6$ である。
従って最初の変更の18次は
\[
 d^{\mathrm T}_{18}+d(d^{\mathrm T}_{12})(X_7)
 +\tfrac12d^2P(X_7,X_7)+3U_4d^{\mathrm T}_{14}-PV_6^2.
\]
この変更後の10次係数は $cH$ である。
次の $-\Xi_9$ による18次への寄与は $-c\,dH(\Xi_9)$、
$-B_8$ による寄与は $-cB_8H$ であり、二次寄与は22次以降である。
これで式を得る。

その後の10次座標変更・9次単元が18次へ作用するには9次係数、
11次座標変更・10次単元が作用するには8次係数が必要である。
いずれも零であり、二次Taylor項はさらに高次なので、18次は変わらない。
\end{proof}

\subsection{有効商における二項への簡約}

第28節により $d^{\mathrm T}_{14}\in J_{14}$ である。
従って $Z_{14}\in\mathcal G_9$ を
\[
 dP(Z_{14})=d^{\mathrm T}_{14}
\]
となるように選べる。

\begin{theorem}[実際の十八次類の計算に十分な代表元]
\label{thm:v14-eighteenth-effective-representative}
原稿の既存正規化の後の実際の18次類は
\begin{equation}\label{eq:v14-eighteen-effective-representative}
 \boxed{\displaystyle
 [R_{18}]=\left[
 d^{\mathrm T}_{18}-\frac12d^2P(X_7,X_7)-c\,dH(Z_{14})
 \right]\quad\text{in }\mathcal C_{18}/I_{18}.
 }
\end{equation}
$Z_{14}$ の選択はこの類に影響しない。
$P\mathcal U_8$ を残す持ち上げ
$d^{\mathrm T}_{14}=dP(Z'_{14})+PC_8$ でも、
式の $Z_{14}$ を $Z'_{14}$ に置き換えてよい。
この定理は、表示された補正項や実際の18次類が零であることを主張しない。
\end{theorem}

\begin{proof}
恒等式 $d^{\mathrm T}_{12}+X_7P=PV_6$ を $X_7$ で微分する。
ベクトル場の係数の微分を含めると
\[
 X_7d^{\mathrm T}_{12}+d^2P(X_7,X_7)
       +dP(DX_7\cdot X_7)
 =V_6X_7P+P X_7V_6.
\]
$DX_7\cdot X_7$ と $V_6X_7$ は $\mathcal G_{13}$ に属し、
$X_7V_6\in\mathcal U_{12}$ である。従って
\[
 d(d^{\mathrm T}_{12})(X_7)
    \equiv-d^2P(X_7,X_7)\pmod{J_{18}}.
\]
また $d^{\mathrm T}_{14}\in J_{14}$ より
$U_4d^{\mathrm T}_{14}\in J_{18}$、
$PV_6^2\in P\mathcal U_{12}\subset J_{18}$ である。
補正式は従って
\[
 R_{18}\equiv d^{\mathrm T}_{18}-\tfrac12d^2P(X_7,X_7)
                       -c\,dH(\Xi_9)\pmod{J_{18}+H\mathcal U_8}
\]
に簡約される。

$\widetilde\Xi_9=\Xi_9+aB_8/6$ と置くと
$dP(\widetilde\Xi_9)=R_{14}$ である。
一方、
\[
 R_{14}=d^{\mathrm T}_{14}+3cHU_4-9PU_4^2
\]
なので
\[
 dP\left(Z_{14}+3cZ_9-\tfrac32aU_4^2\right)=R_{14}.
\]
二つの持ち上げの差は $dP$ の核に入るから、
補題\ref{lem:v14-ninth-stabilizer-hessian}によりHessianへの作用の差は
$J_{18}$ に入る。Euler補正のHessian作用は $H\mathcal U_8$ に入り、
$3c\,dH(Z_9)$ は $I_{18}$ に含まれる。
従って $dH(\Xi_9)$ を $dH(Z_{14})$ に置き換えられる。

純粋な二つの持ち上げ $Z_{14}$ の差には同じ補題を適用する。
$P\mathcal U_8$ を残す場合は $aC_8/6$ を加えれば純粋な持ち上げとなり、
そのHessian作用は $5HC_8/3$ であるため、有効商で消える。
\end{proof}

\begin{corollary}[十七次の明示的な消去を省いてよい]
\label{cor:v14-eighteen-without-seventeenth}
既存の有限正規化を16次まで実行した方程式と、全次数の偶正規形の
18次有効類は一致する。従って式\eqref{eq:v14-eighteen-effective-representative}を
用いるために、17次項や偶化する対合を明示計算する必要はない。
\end{corollary}

\begin{proof}
有限正規化後を
$f=P+cH+R_{17}+R_{18}+O(a^{19})$ と書き、
全次数偶正規形と比較する接触作用素の対数を取る。
次数増分が4以上なので $L(R_{17})$ は21次以上であり、
18次までの条件は $L(P+cH)$ だけによる。
17次での非零右辺は奇増分の条件だけを変える。
偶増分の四条件\eqref{eq:v14-even-stabilizer-chain}はそのまま成立し、
18次の差はその条件を満たす $Q_{12}P+cQ_8H$ である。
従って差は $I_{18}$ に入る。
\end{proof}

\subsection{さらに必要な評価を減らす二つの判定}

次の命題は既存の核の記述から従う。
これにより、後の有限証明書で確認すべき所属を十六列に限定できる。

\begin{proposition}[全九次場のHessian作用を調べる有限基底]
\label{prop:v14-test-all-ninth-hessian}
$\Gamma_1,\ldots,\Gamma_7\in\mathcal G_9$ の接触像が
$J_{14}/P\mathcal U_8$ の基底を与えるとする。
$H\mathcal U_8\subset J_{18}$ および
$dH(\Gamma_i)\in J_{18}$ が全て成り立てば
\[
 dH(\mathcal G_9)\subset J_{18}.
\]
\end{proposition}

\begin{proof}
任意の $Y_9$ に対して、その $dP$ 像から七つの像の線形結合と
$P\mathcal U_8$ の部分を除く。後者はEuler場で持ち上がり、
そのHessian作用は $H\mathcal U_8$ に含まれる。
残る場は $dP$ の核に属し、補題\ref{lem:v14-ninth-stabilizer-hessian}により
Hessianへの作用は $J_{18}$ に入る。
\end{proof}

\begin{proposition}[七次場の二次補正を調べる三つの評価]
\label{prop:v14-test-seventh-quadratic}
$\Gamma_1,\Gamma_2\in\mathcal G_7$ の像が
$J_{12}/P\mathcal U_6$ の基底を与えるとする。
\[
 B_{ij}=\operatorname{Hess}_F(q)
             \bigl(\delta q(\Gamma_i),\delta q(\Gamma_j)\bigr)
 \quad(1\le i\le j\le2)
\]
の三つが $J_{18}$ に含まれれば、全ての $Y,Z\in\mathcal G_7$ について
$d^2P(Y,Z)\in J_{18}$ である。
\end{proposition}

\begin{proof}
合成関数の二次微分は
\[
 d^2P(Y,Z)=\operatorname{Hess}_F(q)(\delta qY,\delta qZ)
                      +dF\bigl(d^2q(Y,Z)\bigr)
\]
である。$d^2q(Y,Z)\in\mathcal D_{14}$ であり、
全次数のPfaffian微分の全射性から最後の項は $J_{18}$ に入る。
第27節の $dF:\mathcal D_8\to\mathcal C_{12}$ の単射性により
\[
 \mathcal D_8=q\mathcal U_6\oplus
       \CC\delta q(\Gamma_1)\oplus\CC\delta q(\Gamma_2).
\]
$\operatorname{Hess}_F(q)q=2\nabla F(q)$ なので、
$q\mathcal U_6$ を一方の因子に含む双線形項は $J_{18}$ に入る。
従って残る三つの組だけを判定すればよい。
\end{proof}

上の二つの有限判定により、必要な十六列の所属が確認されれば、
実際の18次有効類は未正規化の $[d^{\mathrm T}_{18}]$ だけから読める。
後の有理数上の係数汎関数の証明書が、この有限判定を完了する。

% END contact/contact_log_reduction_ja.tex


% BEGIN detector/v14_detector_notes_ja.tex
\section{五点汎関数と低次補正の普遍消滅}
\label{sec:v14-five-point}

% 原稿追加用断片。通常の18次接触階数の証明の後、実際のTate係数の節の前に置く。
\subsection{五つの対角値による18次商の検出}
\label{subsec:v14-five-point-detector}

この節では、命題\ref{prop:v14-ordinary-eighteenth-quotient}で決定した普遍商
\[
 \mathcal Q_{18}=\mathcal C_{18}/J_{18},\qquad
 J_{18}=dP(\mathcal G_{13}),\qquad
 \dim\mathcal Q_{18}=1
\]
の係数を、五つの明示的な有理数点から読む。
全ての共変量を対角三次式だけで復元する必要はない。
既に普遍商の次元が1であると分かっているので、
$J_{18}$ を消し、補元で非零となる一つの汎関数を構成すれば十分である。

基底 $e_0,\ldots,e_4$ と体積形式を固定し、
\[
 F_*=x_0^3+x_1^3+x_2^3+x_3^3-4x_4^3,
\]
\[
 a_b=b(e_1-e_0)\wedge(e_2-e_0)
       +b^{-1}(e_3-e_0)\wedge(e_4-e_0)
 \qquad (b\ne0)
\]
とおく。直接計算で
\[
 q(a_b)=(1,1,1,1,1),\qquad P(F_*,a_b)=0
\]
である。普遍共変量 $f$ のこの点での値を $\operatorname{ev}_b(f)$ と書く。

\begin{proposition}[18次商の五点汎関数]
\label{prop:v14-rho-eighteenth}
\begin{align*}
 D={}&3295420924723200000,\\
 (n_1,n_2,n_3,n_4,n_5)={}&
 (7455703449849690,\ -159838529212047360,\\
 &\hspace{1em}232937478407243349,\ -92218652889186304,\
 11664000244140625)
\end{align*}
として
\[
 \rho_{18}=D^{-1}\sum_{b=1}^{5}n_b\operatorname{ev}_b
\]
と定める。このとき
\[
 \rho_{18}(J_{18})=0,\qquad
 \rho_{18}(\Theta_8U_4J_6)=1.
\]
従って $\ker(\rho_{18}:\mathcal C_{18}\to\mathbb C)=J_{18}$ であり、
任意の $f\in\mathcal C_{18}$ の普遍接触類は
\[
 [f]=\rho_{18}(f)[\Theta_8U_4J_6]
\]
で与えられる。
\end{proposition}

\begin{proof}
通常の接触階数の証明で得た60個のグラフ共変量を
$\Gamma_1,\ldots,\Gamma_{60}$ と書くと
\[
 J_{18}=P\mathcal U_{12}\oplus
        \operatorname{Span}\{dP(\Gamma_1),\ldots,dP(\Gamma_{60})\}.
\]
五点は全て $P=0$ 上にあるため、$\rho_{18}$ は $P\mathcal U_{12}$ を消す。
全60列を五点で整数・有理数により再縮約すると、その評価行列の階数は4であり、
補元 $\Theta_8U_4J_6$ を加えた行列の階数は5となる。
元のグラフ順序で第1、2、3、20列と補元列を取った $5\times5$ 行列式は
\[
 -190124533925774622720\ne0
\]
である。表示した重みはこの行列の転置に対する有理数連立方程式を解いて得る。
選んだ4列だけでなく、残る56列についても重み付き和が厳密に零であることを確認した。

例として $b=1$ での値は
\[
 \Theta_8=-14,\quad U_4=-4,\quad J_6=888,\quad
 \Theta_8U_4J_6=49728,
\]
選んだ4接触列の値は $(496,-480,576,3520)$ である。
五点全ての60列、追加列、グラフの定義、非零行列式、汎関数を
\texttt{v14\_rho18\_exact\_certificate.json} に保存した。
これで $\rho_{18}$ は $J_{18}$ を消して補元を1へ送る。
普遍商の次元が1であることから、核と接触像の等式が従う。
\end{proof}

\paragraph{再縮約と保存証明書監査の区別.}
上の証明書は全 $5\times60$ 成分を新たに縮約して得た。
$a_b$ の分母を $b$ 倍で除き、多倍長整数で縮約した後、
18次の斉次性により $b^{18}$ で割っている。
再計算用の \texttt{v14\_rho18\_exact\_certificate.py} はこの操作を実行する。
別の \texttt{v14\_rho18\_audit.py} は、保存された行列に対する
有理数線形代数、五点の補元値、60列全ての消滅を高速に再検算する。
後者の成功を新規の全テンソル再縮約と同一視しない。

\subsection{低次正規化の寄与は18次接触像に含まれる}
\label{subsec:v14-lower-normalizers-vanish}

五点汎関数を、低次安定化群と実際の正規化補正にも適用する。
以下の主張は、特殊な五点での値だけを普遍恒等式と見なすものではない。
命題\ref{prop:v14-rho-eighteenth}の $\ker\rho_{18}=J_{18}$ と、各空間の既知の有限次元分解を用いる。

\begin{proposition}[単元・Hessian作用・二次補正の消滅]
\label{prop:v14-all-lower-corrections-in-image}
$H=H_{10}$ とすると
\[
 H\mathcal U_8\subset J_{18},\qquad
 dH(\mathcal G_9)\subset J_{18},\qquad
 d^2P(\mathcal G_7,\mathcal G_7)\subset J_{18}
\]
が成り立つ。
\end{proposition}

\begin{proof}
\textbf{単元方向。}
$\dim\mathcal U_8=7$ である。
$P$ の倍数となる方向を作る共変量空間は
\[
 \mathcal E=
 \operatorname{Hom}_{\mathrm{GL}(V)}
 \left(\operatorname{Sym}^7(\operatorname{Sym}^3V)
 \otimes\operatorname{Sym}^2(\Lambda^2V),\Delta^5\right).
\]
二つのLittlewood成分に対応する
$(5,5,5,3,3)$ と $(5,4,4,4,4)$ の重複度は1、0なので
$\dim\mathcal E=1$ である。
また六つの明示的な単元グラフ $u_1,\ldots,u_6$ は
$P=0$ 上の六点での有理数行列式が非零である。
従って
\[
 \mathcal U_8=P\mathcal E\oplus
             \operatorname{Span}\{u_1,\ldots,u_6\}.
\]
$H(P\mathcal E)\subset P\mathcal U_{12}\subset J_{18}$ であり、
五点の厳密縮約は全ての $j$ について $\rho_{18}(Hu_j)=0$ を与える。
これで第一の包含を得る。
この分解の独立な指標計算と $6\times6$ の有理数minorは
\texttt{v14\_unit8\_span\_certificate.py} に記録した。

\textbf{9次座標場のHessian作用。}
ここでは $h(p)=\det\operatorname{Hess}_F(p)$ と置き、
$H(a)=h(q(a))$ と書く。
$dH$ は $a$ 方向の微分、$d_qh(q)$ は $q$ の値域での微分を表す。
既存の14次階数証明から、$\mathcal D_{10}$ は
$q\mathcal U_8$、七つの非単元グラフのPfaffian微分、
および $dF$ の核である三つのKoszul方向によって生成される。
Pfaffian微分の核は $dH$ にも作用しない。
径方向には $d_qh(q)(q u)=5Hu$ が成り立ち、既に第一の包含で処理した。
Koszul方向を $\beta g$、$g=\nabla F(q)$ と書けば、$\beta$ の交代性から
\[
 d_qh(q)(\beta g)=-dF\bigl(\beta\nabla h(q)\bigr)
\]
である。右辺の共変ベクトルは $\mathcal D_{14}$ に属するので $J_{18}$ に入る。
残る七方向については、五点での新規縮約が
\[
 \rho_{18}(dH(\Gamma_j^{(9)}))=0\qquad(1\le j\le7)
\]
を与える。これで第二の包含が従う。

\textbf{7次座標場の二次作用。}
$\mathcal D_8$ は、二次元の $q\mathcal U_6$ と
二つの非単元グラフのPfaffian微分によって生成される。
$F$ が三次であることから
\[
 \operatorname{Hess}_F(q)(q,v)=2\,\nabla F(q)\cdot v.
\]
従って一方が径方向である項は $J_{18}$ に入る。
残る二つの方向の対称積は三方向であり、その全てについて
\[
 \rho_{18}\bigl(\operatorname{Hess}_F(q)
       (\delta q\Gamma_i^{(7)},\delta q\Gamma_j^{(7)})\bigr)=0
 \qquad(1\le i\le j\le2)
\]
を厳密に確認した。
最後に合成関数の二階微分は
\[
 d^2P(X,Y)=\operatorname{Hess}_F(q)(\delta qX,\delta qY)
          +dF\bigl(d^2q(X,Y)\bigr)
\]
であり、第二項は $d^2q(X,Y)\in\mathcal D_{14}$ により $J_{18}$ に属する。
これで第三の包含が得られる。

計算証明書は、以上の六つの単元方向、七つのHessian方向、
三つの二次方向の計16列について、汎関数の値が有理数として全て零であることを保存する。
\end{proof}

\paragraph{実際の18次係数への適用.}
低次の正規化を展開したときに現れる有効代表は
\[
 R_{18}^{\mathrm{eff}}
 =d_{18}^{\mathrm T}
  -\frac12d^2P(X_7,X_7)-\frac1{96}dH(Z_{14})
 \pmod{J_{18}},
\]
ただし $X_7$ は12次を消す座標補正、$Z_{14}$ は14次項の9次liftである。
命題\ref{prop:v14-all-lower-corrections-in-image}により
\[
 \rho_{18}(R_{18}^{\mathrm{eff}})=\rho_{18}(d_{18}^{\mathrm T})
\]
となる。
従って実際の18次類を決めるために必要な新規Tate計算は、
公開した $b=1,2,3,4,5$ の五点だけである。
$Z_{14}$ を新たに選び直したり、その全係数を再構成したりする必要はない。
18次有効商に対する低次安定化群の追加方向も全て $J_{18}$ に含まれるため、
有効商と通常商は一致する。

% END detector/v14_detector_notes_ja.tex


% BEGIN tate/v14_tate_notes_ja.tex
\section{十八次Tate係数の厳密な再構成}
\label{sec:v14-tate-eighteenth}

前節までの接触商の計算と、実際の境界の係数の計算を区別する。
ここでは第13版の切断を証明なしに延長せず、18次に必要な範囲を改めて決める。
$d^{\mathrm T}=\tfrac12\det K$、$K_n=\kappa(A_n)$ とし、
Tate再帰を $A_n=-h\sum_{i=1}^{n-1}A_iA_{n-i}$ とする。

\begin{lemma}[十八次に必要な正確な切断]
\label{lem:v14-exact-fourteenth-truncation}
$d^{\mathrm T}_{18}$ は $K_1,\ldots,K_{14}$ だけから決まる。
完全な $A_1,\ldots,A_{10}$ から出発し、曲率を
\[
 R_{11}:\deg_x\le4,\quad
 R_{12}:\deg_x\le3,\quad
 R_{13}:\deg_x\le2,\quad
 R_{14}:\deg_x\le1
\]
で順に切断すれば、必要な補正の範囲
\[
 A_{11}:\deg_x\le3,\quad
 A_{12}:\deg_x\le2,\quad
 A_{13}:\deg_x\le1,\quad
 A_{14}:\deg_x=0
\]
を正確に得る。
\end{lemma}

\begin{proof}
行列式は五つの成分の積であり、各成分の $a$ 次数は1以上である。
従って18次に寄与する一因子の次数は $18-4=14$ 以下である。
Clifford関係による縮約は係数の $x$ 次数を下げず、
Euler homotopy $h$ はそれを1下げる。
$A_{14}$ の定数項に必要な $R_{14}$ は1次以下、
そこに現れる $A_{13}$ は1次以下でよい。
これを逆向きにたどると、表示された各切断次数が得られる。
初めの十項を完全に作るので、切断した項が後で必要な部分に戻ることはない。

内部重み $3/2$ の正規単項式では
$2\deg_x-|I|+|J|=3$ なので、$A_n$ の $x$ 次数は高々4である。
また $\delta$ は $x$ 次数を1上げるため、曲率の閉性を各 $x$ 次数で判定できる。
従って切断した曲率と補正に対しても
$\delta R_n=0$、$\delta A_n+R_n=0$ を厳密に照合できる。
\end{proof}

\subsection{共通分母による整数演算と独立な照合}

各 $A_n$ を整数係数の正規記号多項式 $N_n$ と正整数 $D_n$ を用いて
$A_n=N_n/D_n$ と保存する。積の分母 $D_iD_{n-i}$ の最小公倍数を取れば、
曲率の加算を多倍長整数だけで実行できる。
$h$ の係数に現れる $\deg_x+|I|$ の最小公倍数をさらに掛け、
最後に全係数と分母の最大公約数で割る。
これは有理数演算の表現を変えるだけであり、丸めを伴わない。

\texttt{v14\_tate\_eighteenth.py} は各入力から $A_1$ を作り直し、
上の範囲を再構成する。全ての再帰段階で内部重み、閉性、対応する
Maurer--Cartan等式を確認する。
完全な $A_1,\ldots,A_{10}$ については、別実装の32次元外積・収縮作用素で
次数0から10の等式も検査する。さらに十四個の定数ブロックを
独立な外積作用素で照合する。
行列式の係数は、120個の置換の和と、使用済み列集合による動的計画法の
二方式で有理数として一致することを確認する。
これは後半の切断された $A_n$ を完全な作用素として評価したという意味ではない。

\begin{proposition}[標準点での十八次係数]
\label{prop:v14-standard-eighteenth-value}
\[
 F_*=x_0^3+x_1^3+x_2^3+x_3^3-4x_4^3,
 \qquad
 a_*=(e_1-e_0)\wedge(e_2-e_0)
       +(e_3-e_0)\wedge(e_4-e_0)
\]
とすると、$q(a_*)=(1,1,1,1,1)$、$P(F_*,a_*)=0$ であり、
\[
 d^{\mathrm T}_{18}(F_*,a_*)
     =\frac{2908218666274031183}{1875352500000}.
\]
第27節の $X_7=-\alpha\Phi_7-\beta\Psi_7$ に対する二次補正は
\[
 \frac12d^2P(X_7,X_7)(F_*,a_*)=\frac{23856412}{1377}.
\]
\end{proposition}

\begin{proof}
公開した入力を用い、補題\ref{lem:v14-exact-fourteenth-truncation}の再帰を
二方式で初めから実行した。
通常のPython \texttt{Fraction} による方法と、上記の共通分母整数演算は、
十四個のブロック全て、0次から18次までの行列式係数全て、
および低次正規化のデータ全てで一致した。
0次から16次は第13版の保存証明書とも一致する。

二次補正は、$q(a+tX)=q(a)+t\,\delta q(X)+t^2q(X)$ を用いて
\[
 [t^2]P(a+tX)
  =\sum_i c_i\left(3q_i^2q_i(X)+3q_i(\delta q_i(X))^2\right)
\]
と計算した。$t=0,1,\ldots,6$ の七つの正確な値からLagrange補間で求めた
$t^2$ 係数とも一致する。
\end{proof}

以上の標準点での非零値は未正規化の係数の値である。
普遍接触商における実際の類は、接触像を消す汎関数を適用して決定する。

\subsection{五つの新規Tate値と実際の十八次類の消滅}

\begin{proposition}[実際の十八次類は零である]
\label{prop:v14-actual-eighteenth-class}
命題\ref{prop:v14-rho-eighteenth}の五点汎関数に対して
\[
 \rho_{18}(d^{\mathrm T}_{18})=0,\qquad
 \rho_{18}\left(\tfrac12d^2P(X_7,X_7)\right)=0.
\]
通常商と有効商はともに $\CC\cdot[\Theta_8U_4J_6]$ であり、
既存の低次正規化後の実際の十八次係数はその零元を表す。
\end{proposition}

\begin{proof}
公開した $a_b$ の五つの入力から、Tate再帰を各々初めから再構成した。
得られた未正規化の係数は次の通りである。
\[
\begin{array}{c|c}
b&d^{\mathrm T}_{18}(F_*,a_b)\\\hline
1&\dfrac{2908218666274031183}{1875352500000}\\[0.7em]
2&\dfrac{396921899835705304591}{11854080000000}\\[0.7em]
3&\dfrac{259707810011438926969409}{384505867265625}\\[0.7em]
4&\dfrac{1590036351104669653664022401}{245806202880000000}\\[0.7em]
5&\dfrac{2791927762015243689203604361}{73255957031250000}
\end{array}
\]
命題\ref{prop:v14-rho-eighteenth}の整数重みで足し合わせると、
\[
 \sum_{b=1}^5 n_b\,d^{\mathrm T}_{18}(F_*,a_b)=0
\]
が有理数として正確に成立する。
二次補正の五つの値は順に
\[
 \left(
 \frac{23856412}{1377},\quad
 \frac{501277499}{2754},\quad
 \frac{33576631924}{37179},\quad
 \frac{125350459859}{44064},\quad
 \frac{5971728868828}{860625}
 \right)
\]
であり、同じ重み付き和は零である。
後者は命題\ref{prop:v14-all-lower-corrections-in-image}の普遍的な所属とも整合する。

五つの各計算について、十四個の定数ブロック、18次までの全行列式係数、
閉性とMaurer--Cartan等式の検査結果、独立なブロック抽出の検査結果を公開した。

\texttt{v14\_actual\_eighteenth\_certificate.py} は、
汎関数の五点とTate計算の入力を厳密一致で照合し、各行列式を再計算してから
上の和を有理数で検証する。

別の \texttt{v14\_detector\_independent\_audit.py} は、
独自の有理数消去法により評価行列と重み付き和をもう一度確認する。

命題\ref{prop:v14-all-lower-corrections-in-image}により、
有効像の追加方向と低次正規化の二項の補正は全て $J_{18}$ に入る。
従って定理\ref{thm:v14-eighteenth-effective-representative}は
$[R_{18}]=[d^{\mathrm T}_{18}]$ を与える。
一方、$\ker\rho_{18}=J_{18}$ であり、今計算した値が零なので
この類は零である。
\end{proof}

% END tate/v14_tate_notes_ja.tex


% BEGIN contact/twentieth_normal_form_ja.tex
% 第14版の確定稿。2026-10-29、rank・16補正方向・fresh Tate 5点を有理数上で認証。
% contact_log_reduction_ja.tex の命題、第13版の結果、今回の有限証明書を使用する。

\section{十八次判定から二十次開始の正規形への帰結}
\label{sec:v14-twentieth}

本節では、十八次の通常階数、五点係数汎関数、低次正規化の補正、
実際のTate係数の有限証明書を組み合わせる。
これらの計算により、通常商に残る一つの普遍方向は有効商にも残るが、
実際の境界が定める類は零であることが確定する。
$P=P_F$、$H=H_{10}$、$c=1/96$、$J_{18}=dP(\mathcal G_{13})$ とする。

\begin{theorem}[十八次の消去と二十次開始の代数的収束偶正規形]
\label{thm:v14-twentieth-normal-form}
十八次の通常商と有効商は一致し、
\[
 \mathcal Q^{\mathrm{lin}}_{18}
 =\mathcal Q^{\mathrm{eff}}_{18}
 =\CC\cdot[W_{18}],\qquad
 W_{18}=\Theta_8U_4J_6
\]
と一次元である。ここで $J_6$ は六頂点単元グラフによる既存の六次共変量である。
実際の境界の十八次類は零である。
任意の滑らかな三次式 $F$ の実際の境界は、普遍同変な代数的冪級数による
接触変換の下で
\begin{equation}\label{eq:v14-twentieth-normal-form}
 \boxed{\displaystyle
 d^{\mathrm T}\sim P_F+\frac1{96}\det\operatorname{Hess}_F(q(a))
       +\sum_{\substack{n\ge20\\n\ \mathrm{even}}}R_n,
 \qquad R_n\in\mathcal C_n
 }
\end{equation}
と正規化できる。座標変更の線形項は恒等写像、単元の定数項は1である。
各形式係数は $F$ の係数に多項式的であり、
各固定 $F$ では変更・逆変更・方程式の級数は原点近傍で収束する。
この結論は20次項以降の消去を主張しない。
\end{theorem}

\begin{proof}
まず、今回の有限証明書から使う事実を整理する。
整数指標と普遍接触階数の証明書により
\[
 \dim\mathcal C_{18}=119,\qquad
 \dim J_{18}=118,\qquad \dim P\mathcal U_{12}=58
\]
であり、選ばれた六十の接触グラフ像は
$J_{18}/P\mathcal U_{12}$ の基底である。
五つの有理点で定義された係数汎関数 $\rho_{18}$ は、
この六十列を全て消し、$\rho_{18}(W_{18})=1$ を満たす。
五点は全て $P=0$ 上にあるため、$P\mathcal U_{12}$ も自動的に消える。
従って
\[
 \ker\rho_{18}=J_{18}.
\]
これは119次元の共変量を五点から復元するという主張ではなく、
既に余次元1と確定した接触像を消す、非零な係数汎関数の構成である。

命題\ref{prop:v14-all-lower-corrections-in-image}の
同じ汎関数による有理数上の計算は、
六つの八次単元の $H$ 倍、七つの九次場のHessian作用、
二つの七次場に対する三つのHessian双線形評価の全十六列を消す。
最初の六つは $\mathcal U_8$ の $P$ 倍部分を法とする基底であり、
残る一次元の $H$ 倍は命題\ref{prop:v14-pdivisible-eighth-unit}により
$J_{18}$ に属する。従って $H\mathcal U_8\subset J_{18}$ である。
七つの場 $\Gamma_i\in\mathcal G_9$ の接触像は
$J_{14}/P\mathcal U_8$ の基底を、二つの場
$\Lambda_i\in\mathcal G_7$ の接触像は
$J_{12}/P\mathcal U_6$ の基底を与える。
以上の所属と命題\ref{prop:v14-test-all-ninth-hessian}、
\ref{prop:v14-test-seventh-quadratic}から
\[
 dH(\mathcal G_9)\subset J_{18},\qquad
 d^2P(\mathcal G_7,\mathcal G_7)\subset J_{18}
\]
を得る。特に $dP(Z_9)=HU_4$ を満たす場のHessian作用も $J_{18}$ に入る。
定理\ref{thm:v14-effective-eighteenth}により有効像は通常像 $J_{18}$ と一致する。
従って通常商と有効商はいずれも $[W_{18}]$ を基底とする一次元空間である。

定理\ref{thm:v14-eighteenth-effective-representative}の補正項は、
今得た二つの所属によりともに $J_{18}$ に入る。
従って、既存の有限正規化後の18次係数 $R_{18}^{\mathrm{fin}}$ は
\[
 [R_{18}^{\mathrm{fin}}]=[d^{\mathrm T}_{18}]
      \quad\text{in }\mathcal C_{18}/J_{18}
\]
を満たす。
命題\ref{prop:v14-actual-eighteenth-class}で
実際のTate作用素から新たに計算した五点の18次係数は、
同じ有理数係数の汎関数に対して厳密に
\[
 \rho_{18}(d^{\mathrm T}_{18})=0
\]
を満たす。五点の入力は $(F,a)$ の有理数値で照合されており、
既存の補正係数や零という結論から18次係数を逆算していない。
よって $d^{\mathrm T}_{18}\in J_{18}$ であり、実際の18次類は零である。

第13版の定理\ref{thm:v13-eighteenth-even-normal-form}の
全次数偶正規形を
\[
 f_{\mathrm{ev}}=P+cH+R_{18}^{\mathrm{ev}}+R_{20}^{\mathrm{ev}}+\cdots
\]
と書く。
有限正規化後には17次の項が残っていてもよい。
系\ref{cor:v14-eighteen-without-seventeenth}の対数による偶奇分離から
\[
 R_{18}^{\mathrm{ev}}-R_{18}^{\mathrm{fin}}\in I_{18}=J_{18}
\]
である。従って $R_{18}^{\mathrm{ev}}\in J_{18}$ であり、
17次の明示的な消去や、対合 $\sigma$ の18次までの計算を必要としない。

$J_{18}=dP(\mathcal G_{13})$ の定義から、普遍多項式な
$\Xi_{13}\in\mathcal G_{13}$ を
$dP(\Xi_{13})=R_{18}^{\mathrm{ev}}$ となるように選べる。
これは有限次元の普遍共変量空間の線形写像に対する持ち上げであり、
$F$ の判別式や $P$ の点ごとの値による除算を含まない。
\[
 a\longmapsto a-\Xi_{13}(F,a)
\]
で引き戻すと18次が消える。
$\Xi_{13}$ は奇数次なので座標変更は奇写像であり、
全次数の偶性を保つ。
$H$ への最初の寄与は22次、$P$ への二次Taylor項は30次からであるため、
既定の低次正規形は変わらない。次に残り得る次数は20である。

既存の全次数偶正規形への変換は、代数的冪級数であり、
各固定 $F$ において解析的に収束していた。
今回加えた変更は有限な多項式で、原点でのJacobianは恒等写像である。
従ってその局所逆は代数的冪級数で、解析的にも原点近傍で収束する。
既存の変更との有限合成はこれらの性質を保つ。
普遍係数の多項式性と同変性も保たれ、所要の正規形を得る。
\end{proof}

\begin{remark}[有限計算の範囲と残る問題]
普遍階数の上下界は、一般三次テンソルによる独立性と
恒等的なKoszul核の構成から証明している。
五点の評価は、その後で一次元の商を検出するために用いている。
所属と実際の類の消滅には有理数上の零を使っており、
有限体上の零という観察だけには依存しない。
標準点における未正規化の $d^{\mathrm T}_{18}$ は非零だが、
接触像を除いた普遍類は零である。

定理は、十八次の実際の障害が消えることと、二十次以上にのみ項が残り得る
代数的収束偶正規形の存在を述べる。
二十次以上の普遍接触商、実際の次の非零係数、
$P+H/96$ への全次数の接触同値は、ここでは未決定である。
\end{remark}

% END contact/twentieth_normal_form_ja.tex


% BEGIN rank/v14_unit_structure_ja.tex
\section{次の次数に向けた二つの単元の構造}
\label{sec:v14-unit-structure}

十八次までの商に現れた $U_4,J_6$ について、
全次数の議論へ進むための具体的な制限計算を記録する。
ここでは普遍多項式の性質を調べるため、補助入力として特異な三次式も用いる。

\begin{proposition}[階数二の点における単元の制限と代数的独立性]
\label{prop:v14-rank-two-unit-restriction}
体積形式を固定して
\[
 a=e_0\wedge e_1,\qquad
 F_t=x_2^3+x_3^3+x_4^3+6t x_2x_3x_4
\]
とすると、本稿の規格化で
\[
 U_4(F_t,a)=t^4-t,\qquad
 J_6(F_t,a)=6-120t^3-48t^6
\]
であり、
\[
 J_6(F_t,a)^2-2304U_4(F_t,a)^3=36(1+8t^3)^3
\]
となる。また $U_4,J_6$ は普遍多項式として代数的に独立である。
\end{proposition}

\begin{proof}
$U_4$ と $J_6$ のテンソルグラフをこの入力で縮約する。
両者の $F$ 次数はそれぞれ4、6なので、$t$ についての次数は高々4、6である。
$t=0,1,\ldots,6$ の七つの整数入力で得た値は
\[
\begin{array}{c|rrrrrrr}
t&0&1&2&3&4&5&6\\\hline
U_4&0&0&14&78&252&620&1290\\
J_6&6&-162&-4026&-38226&-204282&-764994&-2265402
\end{array}
\]
であり、有理数の補間は表示した多項式を与える。
二つ目の恒等式はそれらの多項式を直接展開して確認できる。

さらに $F=sF_t$ とすると、次数から
\[
 (U_4,J_6)=\bigl(s^4(t^4-t),\ s^6(6-120t^3-48t^6)\bigr).
\]
この写像の $(s,t)=(1,0)$ におけるJacobian行列式は36である。
従って二つの成分の間に非零の多項式関係はなく、普遍多項式としても独立である。
計算は \texttt{v14\_rank2\_unit\_restriction.py} で再現できる。
\end{proof}

\paragraph{二十次に向けた作業仮説.}
既に決定した通常商の生成元
\[
 [\Theta_8],\quad[\Theta_8U_4],\quad[\Theta_8J_6],\quad
 [\Theta_8U_4^2],\quad[\Theta_8U_4J_6]
\]
は、それぞれ8、12、14、16、18次に現れる。
この規則から、20次では $\Theta_8U_4^3$ と $\Theta_8J_6^2$ を
まず調べるのが自然である。
上の単元の代数的独立性だけでは、この二つが接触商で独立であることも、
商を生成することも従わない。
全次数の接触商の構造とHilbert級数は、ここでは未決定である。
特に、今回消去した実際の十八次類が零であることから、
二十次以降の実際の類の消滅を推論することはできない。

% END rank/v14_unit_structure_ja.tex

\clearpage
% BEGIN characters/v15_character_notes_ja.tex
% 第15版用。第13・14版の本文・証明書は変更せず参照する。
\section{20次の普遍共変量とKoszul終端}
\label{sec:v15-twentieth-characters}

20次の研究の入力として、普遍境界、座標補正、単元補正、Pfaffian微分の
終域とKoszul各項の次元を確定する。
本節の計算は整数指標計算であり、Koszul完全性や接触写像の階数を仮定しない。
特に18次とは異なり、20次ではKoszulの第5項が消えない。

\subsection{次数とdeterminant twist}

$\dim V=5$、$A=\Lambda^2V$、$\Delta=\det V$、
$W=A\otimes\Delta^{-1}$ とする。
$F\in\operatorname{Sym}^3V$、$a\in W$、
$q(a)=\tfrac12a\wedge a\in V^*\otimes\Delta^{-1}$、
$P=F(q(a))\in\Delta^{-3}$ という規約を保つ。
中心元 $t\,\mathrm{id}_V$ に対する重みは
$F$ が3、$a$ が $-3$、$P$ が $-15$ である。
従って $a$ 次数 $n$ の境界項は $F$ 次数 $n-5$、
$a$ 次数 $r$ の座標補正は $F$ 次数 $r-1$、
$a$ 次数 $s$ の単元補正は $F$ 次数 $s$ を持つ。
trace線を外すと
\begin{align*}
\mathcal C_n&=\operatorname{Hom}_{\mathrm{GL}(V)}
 \left(\operatorname{Sym}^{n-5}(\operatorname{Sym}^3V)
 \otimes\operatorname{Sym}^n A,\Delta^{n-3}\right),\\
\mathcal G_r&=\operatorname{Hom}_{\mathrm{GL}(V)}
 \left(\operatorname{Sym}^{r-1}(\operatorname{Sym}^3V)
 \otimes\operatorname{Sym}^r A,A\otimes\Delta^{r-1}\right),\\
\mathcal D_{r+1}&=\operatorname{Hom}_{\mathrm{GL}(V)}
 \left(\operatorname{Sym}^{r-1}(\operatorname{Sym}^3V)
 \otimes\operatorname{Sym}^{r+1}A,\Lambda^4V\otimes\Delta^{r-1}\right),\\
\mathcal U_s&=\operatorname{Hom}_{\mathrm{GL}(V)}
 \left(\operatorname{Sym}^{s}(\operatorname{Sym}^3V)
 \otimes\operatorname{Sym}^{s}A,\Delta^{s}\right).
\end{align*}
20次では $n=20$、$r=15$、$s=14$ となる。

$g=\nabla F(q)$ は $F$ 次数1、$a$ 次数4を持つ。
trace線を外したKoszulの第 $k$ 項は
\begin{equation}\label{eq:v15-koszul-term}
\mathcal K_{k,20}=
\operatorname{Hom}_{\mathrm{GL}(V)}
\left(\operatorname{Sym}^{15-k}(\operatorname{Sym}^3V)
 \otimes\operatorname{Sym}^{20-4k} A,
 \Lambda^k V^*\otimes\Delta^{17-2k}\right).
\end{equation}
$\mathcal K_{0,20}=\mathcal C_{20}$、
$\mathcal K_{1,20}=\mathcal D_{16}$、
$\mathcal K_{2,20}=\mathcal B_{20}$ である。
$k=5$ では $a$ 次数は0、$F$ 次数は10、
値域は $\Lambda^5V^*\otimes\Delta^7=\Delta^6$ となる。
従って、この項を次数の理由で零とすることはできない。

\begin{proposition}[20次の厳密次元]\label{prop:v15-twentieth-dimensions}
次の次元が成り立つ。
\[
\begin{array}{c|rrrrrrrr}
 &\mathcal C_{20}&\mathcal G_{15}&\mathcal U_{14}&\mathcal D_{16}
 &\mathcal B_{20}&\mathcal K_{3,20}&\mathcal K_{4,20}&\mathcal K_{5,20}\\\hline
\dim&330&1244&164&688&468&117&10&1
\end{array}
\]
\end{proposition}

\begin{proof}
Littlewood分解
\[
\operatorname{Sym}^d A=
 \bigoplus_{j=0}^{\lfloor d/2\rfloor}S_{(d-j,d-j,j,j,0)}V
\]
とdeterminant長方形の補分割を用いる。
$\mathcal C_{20}$ は $\operatorname{Sym}^{15}(\operatorname{Sym}^3V)$ の
$\lambda_j=(17,14-j,14-j,j,j)$ 成分の重複度の和、
$\mathcal U_{14}$ は $\operatorname{Sym}^{14}(\operatorname{Sym}^3V)$ の
$\mu_j=(14,14-j,14-j,j,j)$ 成分の重複度の和となる。
ここで $0\le j\le7$ であり、重複度は
\[
\begin{array}{c|rrrrrrrr}
j&0&1&2&3&4&5&6&7\\\hline
[S_{\lambda_j}]&3&4&28&49&118&72&55&1\\
{}[S_{\mu_j}]&1&1&11&21&59&38&32&1
\end{array}
\]
である。和はそれぞれ330、164となる。
残る空間は、次のhighest weightの重複度として求める。
\[
\begin{array}{c|rr|c|r}
\text{空間}&F\text{次数}&a\text{次数}&\text{値域のhighest weight}&\dim\\\hline
\mathcal G_{15}&14&15&(15,15,14,14,14)&1244\\
\mathcal D_{16}&14&16&(15,15,15,15,14)&688\\
\mathcal B_{20}&13&12&(13,13,13,12,12)&468\\
\mathcal K_{3,20}&12&8&(11,11,10,10,10)&117\\
\mathcal K_{4,20}&11&4&(9,8,8,8,8)&10\\
\mathcal K_{5,20}&10&0&(6,6,6,6,6)&1
\end{array}
\]
各行で $3\deg_F+2\deg_a$ はhighest weightの総和と一致する。

計算算法は有限な整数演算である。
三次単項式の35個の重み $\alpha$ について
\[
 \prod_{|\alpha|=3}(1-tx^\alpha)^{-1}
  =\sum_{m,\beta}M_m(\beta)t^mx^\beta
\]
をEuler積の動的計画法で計算した。
$\rho=(4,3,2,1,0)$ とするとSchur重複度は
\[
 \sum_{\sigma\in\mathfrak S_5}\operatorname{sgn}(\sigma)
 M_m(\lambda+\rho-\sigma\rho)
\]
である。二つのcharacterの積にも同じ交代式を適用する。
第5項の終端については、正符号項と負符号項の和がそれぞれ
$35085580$、$35085579$ であり、その差が1となる。

独立な検算としてNewton--Adams再帰
\[
 mh_m=\sum_{k=1}^m p_kh_{m-k},\qquad
 p_k=\sum_{|\alpha|=3}x^{k\alpha},\qquad h_0=1
\]
でも $h_m=\operatorname{ch}\operatorname{Sym}^m(\operatorname{Sym}^3V)$ を求めた。
$1\le m\le15$ の全重みの整数係数が一致し、
$m=15$ では211876個の重みを照合した。
全体の次元も
\[
 \dim\operatorname{Sym}^{15}(\operatorname{Sym}^3V)
 =\binom{49}{15}=1575580702584
\]
と一致する。
$10\le m\le15$ の全Schur分解を列挙し、
Weyl次元公式による和と $\binom{m+34}{m}$ の一致も確認した。
この範囲の非零Schur成分数は順に
$550,793,1107,1469,1931,2476$ である。

$\mathcal G_{15}$ と $\mathcal D_{16}$ は27行の独立なPieri表で再計算した。
四つのKoszul源は、旧Pieri実装を読み込まず、
5行のうち $k$ 行を選んで1を加える縦strip列挙でも再計算した。
各段でPieri分解全体の次元と
$\binom5k\dim\operatorname{Sym}^{20-4k}A$ も一致する。
これらの検算は標本評価や浮動小数点近似を用いない。
\end{proof}

\paragraph{Pfaffian微分とEuler標数.}
既存の全次数Pfaffian全射定理を $r=15$ に適用すると
\[
 \delta q:\mathcal G_{15}\twoheadrightarrow\mathcal D_{16},\qquad
 \operatorname{rank}\delta q=688,\qquad
 \dim\ker\delta q=556
\]
を得る。ここでの全射性は指標の次元比較ではなく、既存の定理に依存する。
20次のKoszul成分は
\[
0\longrightarrow\mathcal K_{5,20}\longrightarrow\mathcal K_{4,20}
 \longrightarrow\mathcal K_{3,20}\longrightarrow\mathcal B_{20}
 \longrightarrow\mathcal D_{16}\longrightarrow\mathcal C_{20}
 \longrightarrow0
\]
であり、そのEuler標数は
\[
 330-688+468-117+10-1=2.
\]
この等式だけでは $\dim\mathcal C_{20}/dF(\mathcal D_{16})=2$ とは結論できない。
他のホモロジーまたは写像の階数を別に決める必要がある。

\subsection{対角値と一階jetの普遍的な不足}

$F_{\mathrm{Fer}}=\sum_i x_i^3$、
$H=\mu_3^5\rtimes\mathfrak S_5$、
$N=\operatorname{Jac}(F_{\mathrm{Fer}})_3$ とする。
$N$ は平方因子を持たない三次単項式10個で張られる。
$m_\lambda$ を普遍出現重複度、
$h_0(\lambda)=\dim(S_\lambda V)^H$、
$h_1(\lambda)=\dim\operatorname{Hom}_H(N,S_\lambda V)$ と書く。
対角値の制限階数は $\min(m_\lambda,h_0)$ 以下、
値と全法方向一階jetを合わせた階数も $\min(m_\lambda,h_0+h_1)$ 以下である。

\begin{proposition}[20次の制限上界]\label{prop:v15-diagonal-visibility}
次の上界と下界が成り立つ。
\[
\begin{array}{c|rrrr}
&\mathcal C_{20}&\mathcal U_{14}&\mathcal D_{16}&\mathcal G_{15}\\\hline
\text{普遍次元}&330&164&688&1244\\
\text{対角値の階数上界}&137&43&230&382\\
\text{対角値の核の次元下界}&193&121&458&862\\
\text{値と全法方向一階jetの階数上界}&307&130&624&1102\\
\text{同じ一階jetの核の次元下界}&23&34&64&142
\end{array}
\]
\end{proposition}

\begin{proof}
有限群の指標平均を $\mathbb Q(\zeta_3)$ の整数対表示で計算する。
$\mathcal C_{20}$ の表は
\[
\begin{array}{c|r|r|r}
\lambda&m_\lambda&h_0&h_1\\\hline
(17,14,14,0,0)&3&218&2019\\
(17,13,13,1,1)&4&151&1529\\
(17,12,12,2,2)&28&133&1105\\
(17,11,11,3,3)&49&52&564\\
(17,10,10,4,4)&118&38&323\\
(17,9,9,5,5)&72&10&79\\
(17,8,8,6,6)&55&5&28\\
(17,7,7,7,7)&1&0&0
\end{array}
\]
となる。最後の二成分には、一階jetまででもそれぞれ少なくとも22方向、1方向の核が残る。
単元空間の表は
\[
\begin{array}{c|r|r|r}
\lambda&m_\lambda&h_0&h_1\\\hline
(14,14,14,0,0)&1&21&165\\
(14,13,13,1,1)&1&22&219\\
(14,12,12,2,2)&11&33&232\\
(14,11,11,3,3)&21&12&133\\
(14,10,10,4,4)&59&13&98\\
(14,9,9,5,5)&38&3&22\\
(14,8,8,6,6)&32&2&10\\
(14,7,7,7,7)&1&0&0
\end{array}
\]
であり、最後の三成分から一階jetの核の次元下界 $13+20+1=34$ を得る。
$\mathcal D_{16}$ と $\mathcal G_{15}$ では、同じ上界にPieri側の外側重複度を掛けて和を取る。
全27行の表と指標平均はJSON証明書に保存した。
\end{proof}

これらは制限写像の階数の上界であり、実際の階数ではない。
20次では境界空間そのものも、対角値と全法方向一階jetだけでは単射に検出できない。
一方、普遍接触商の次元と補元が別途確定した場合には、
その商だけを少数の対角点で検出できる可能性は残る。
普遍空間の検出と、既に決定された商の検出を区別する必要がある。

\subsection{計算証明書}

\nolinkurl{v15_character_certificate.py} は原第13版の
\nolinkurl{nc_cubic_verify_v7.py}、
\nolinkurl{v9_verify_contact_bounds.py}、
\nolinkurl{v11_fermat_scope_certificate.py} を不変入力として使用する。
入力のSHA256、全交代式、全Schur分解、Pieri表、有限安定化群の平均、
全重みNewton照合の記録を出力する。
\nolinkurl{v15_koszul_pieri_check.py} は四つのKoszul源を別実装で再検算する。
いずれもassertionを無効にする実行設定を拒否する。
本節では、通常接触商・有効接触商の次元と実際の20次係数をまだ主張しない。

% END characters/v15_character_notes_ja.tex


% BEGIN rank/rank20_section_ja.tex
\section{二十次の接触部分空間と階数の未決定部分}
\label{sec:v15-contact-subspace}

\subsection{20次の明示的接触部分空間と通常階数の上下界}
\label{subsec:v15-twentieth-contact-subspace}

本小節で決定するのは、通常接触像に含まれる明示的な328次元部分空間である。
通常接触像そのものの次元には一方向の未決定部分を残すが、
後述の実際の20次項の消去には、この部分空間への所属だけで十分である。
以下では
\[
 P=P_F=F(q(a)),\qquad g=\nabla F(q(a)),\qquad
 J_{20}=dP(\mathcal G_{15})
\]
とおく。前節の独立な整数指標計算から
\[
 \dim\mathcal C_{20}=330,\qquad
 \dim\mathcal G_{15}=1244,\qquad
 \dim\mathcal U_{14}=164,\qquad
 \dim\mathcal D_{16}=688
\]
を用いる。命題\ref{prop:v12-uniform-pfaffian-surjectivity}により
$\delta q:\mathcal G_{15}\twoheadrightarrow\mathcal D_{16}$ は全射である。
従って $dP=dF\circ\delta q$ の二つの写像 $dP,dF$ は同じ階数を持つ。
以下のテンソルの値域は、前小節と同じくtrace線を外した表示で記す。

\begin{proposition}[確定した20次接触部分空間]
\label{prop:v15-twentieth-contact-subspace}
公開された164個の普遍座標場
$\Gamma_1,\ldots,\Gamma_{164}\in\mathcal G_{15}$ を用いて
\[
 S_{20}=P\mathcal U_{14}\oplus
     \operatorname{Span}_{\CC}\{dP(\Gamma_1),\ldots,dP(\Gamma_{164})\}
\]
とおくと、
\[
 \dim S_{20}=328,\qquad S_{20}\subset J_{20}.
\]
さらに、第13・14版と同じ正規化の $U_4,J_6,\Theta_8$ に対して
\[
 \mathcal C_{20}
   =S_{20}\oplus
       \CC\,\Theta_8U_4^3\oplus\CC\,\Theta_8J_6^2.
 \tag{20S}\label{eq:v15-twentieth-certified-splitting}
\]
特に $\mathcal C_{20}/S_{20}$ は二次元であり、右辺の二つの類を基底とする。
この段階では、これらの類が $\mathcal C_{20}/J_{20}$ でも独立であるとは主張しない。
\end{proposition}

\begin{proof}
既存のテンソル規約を用い、
$F_{ijk}=F_{(ijk)}$、
$\tau_{ijk}=\epsilon_{ijklm}a_{lm}$ とする。
ここで $l<m$ は $i,j,k$ の補集合である。
左・右に各15頂点を持つ三価二部グラフについて、
最後の左頂点を除き、左の14頂点に $F$、
右の15頂点に $\tau$ を置く。
三本の自由脚を持つ縮約を $R_{ijk}$ として
\[
 \Gamma^{lm}=\frac16\sum_{i,j,k}\epsilon_{lmijk}R_{ijk}
\]
と定める。$\Gamma$ は $F^{14}a^{15}$ の次数を持ち、
determinant twist を含めて $\mathcal G_{15}$ に属する。
具体的には値域は $\bigwedge^2V\otimes(\det V)^{14}$ である。
また
\[
 T_{ijk}=\epsilon_{ijklm}\frac{\partial P}{\partial a_{lm}}
\]
を最後の左頂点に置いた閉じたグラフの縮約を6で割ると、
$dP(\Gamma)$ が得られる。

整数入力を180個取り、各入力で
\[
              q(a)=(1,1,1,1,1),\qquad P=0
\]
が有理数上で成立するようにする。
各対称テンソル係数を整数として選び、
$F_{444}$ によって $P=\sum_{i,j,k}F_{ijk}=0$ を課している。
この制限で $P\mathcal U_{14}$ は消える。
選んだ164グラフの評価行列は、1009を法として
\[
                  \det M_{\rm G}\equiv342\pmod{1009}
\]
となる $164\times164$ 小行列を持つ。
したがって $dP(\Gamma_i)$ は $P\mathcal U_{14}$ を法として特性零で独立である。
入力、グラフ、選択行、全評価値の剰余は
\texttt{contact20\_search.json} に記録した。

Eulerの関係
\[
                         dP(au/6)=Pu
\]
から $P\mathcal U_{14}\subset J_{20}$ である。
普遍係数と $a$ の多項式環は整域なので、乗法 $u\mapsto Pu$ は単射であり、
この部分空間の次元は164である。
従って $S_{20}$ は上記の直和を持ち、次元は328となる。

一方、
\[
             W_1=\Theta_8U_4^3,\qquad W_2=\Theta_8J_6^2
\]
はいずれも $F^{15}a^{20}$ の次数と $(\det V)^{17}$ の変換則を持つので
$\mathcal C_{20}$ に属する。
同じ180入力で $W_1,W_2$ を有理数により評価し、上の164列に追加すると、
\[
                  \det M_{\rm aug}\equiv667\pmod{1009}
\]
となる $166\times166$ 小行列を得る。
値と正規化は \texttt{complement20\_certificate.json} に記録した。
したがって、この166方向は $P\mathcal U_{14}$ を法として独立である。
さらに $164+166=330=\dim\mathcal C_{20}$ なので、
式\eqref{eq:v15-twentieth-certified-splitting}が従う。
\end{proof}

\begin{proposition}[通常20次階数の確定した範囲]
\label{prop:v15-twentieth-contact-rank-bounds}
通常接触写像について
\[
 328\le
 \operatorname{rank}(dP:\mathcal G_{15}\to\mathcal C_{20})
 =
 \operatorname{rank}(dF:\mathcal D_{16}\to\mathcal C_{20})
 \le329.
\]
従って
\[
 \begin{split}
 \dim\ker dP&\in\{915,916\},\\
 \dim\ker dF&\in\{359,360\},\\
 \dim(\mathcal C_{20}/J_{20})&\in\{1,2\}.
 \end{split}
\]
$\mathcal C_{20}/J_{20}$ は $[W_1],[W_2]$ によって生成される。
これらの間に通常接触像を法とする一つの関係があるかどうかは未決定である。
\end{proposition}

\begin{proof}
下界328は命題\ref{prop:v15-twentieth-contact-subspace}から従う。
上界には359個の明示的な核ベクトルを用いる。
左・右に各13頂点を持つ三価二部グラフで最後の右頂点を除き、
左に13個の $F$、右に12個の $\tau$ を置く。
三本の自由脚を持つ縮約を $R^{ijk}$ として
\[
       \beta_{lm}=\sum_{i,j,k}\epsilon_{lmijk}R^{ijk},
       \qquad h=\beta g
\]
と定める。$\beta$ は $F^{13}a^{12}$ の次数と
$\bigwedge^2V^*\otimes(\det V)^{13}$ の変換則を持つ。
従って $h$ は $F^{14}a^{16}$ の次数と
$V^*\otimes(\det V)^{15}$ の変換則を持つ $\mathcal D_{16}$ の元である。
交代性から
\[
                      dF(h)=g^{\rm t}\beta g=0
\]
が多項式恒等式として成立する。

この方法で得た359個の $h$ を90個の整数入力で評価し、
五成分を並べる。この評価では $P=0$ を課さない。
評価行列には
\[
                   \det M_{\rm K}\equiv269\pmod{1009}
\]
となる $359\times359$ 小行列がある。
全データは \texttt{koszul20\_search.json} に記録した。
従って核には特性零で独立な359方向があり、
\[
                        \operatorname{rank}dF\le688-359=329
\]
となる。$\delta q$ の全射性と各空間の次元を用いれば、残りの主張も従う。
\end{proof}

\begin{remark}[非零minorの意味と再現可能性]
\label{rem:v15-twentieth-rank-certificate-scope}
採用したグラフと整数入力を固定した後は、
上記の小行列式は全て厳密な有限計算である。
テンソル公式は整係数であり、接触場の正規化に現れる分母6は1009で可逆である。
従って非零の法1009小行列式は、有理数上の対応する小行列式の非零性を証明する。
これは階数の確率的推測ではない。

独立監査では、発見時とは異なる右頂点による辺の番号付けと
整数 \texttt{einsum} を用いて、全ての採用グラフの全入力を再縮約した。
各縮約では加算項数を数え、64 bit整数のoverflowが起こらないことを確認してから、
演算ごとに1009で還元した。
この監査は $342,269$ を再現する。
二つの補元の180点での値は有理数上で独立に再計算した。
\nolinkurl{v15_rank_certificate.py} は、これらの非零minor、
独立監査の入力hash、指標次元、および上下界をまとめて検証する。

本計算は、Koszul複体の完全性も、Euler和からの階数決定も用いない。
特に359は、構成したKoszul像の\emph{独立方向数}であり、
Koszul像全体の次元が359だと証明したものではない。
さらに多くのグラフを探索して階数が増えなかったという事実も、
数学的な上界としては用いていない。
\end{remark}

\subsection{通常商を完全決定しなくても20次消去を証明できる理由}
\label{subsec:v15-certified-subspace-suffices}

後述する三点の有理評価汎関数を $\rho_1,\rho_2$ とする。
それらが全ての $dP(\Gamma_i)$ と $P\mathcal U_{14}$ を消し、
\[
                  \rho_i(W_j)=\delta_{ij}
\]
を満たすことを有理数上で確認すれば、
式\eqref{eq:v15-twentieth-certified-splitting}から直ちに
\[
 S_{20}=\ker\rho_1\cap\ker\rho_2,\qquad
 R\in S_{20}\ \Longleftrightarrow\
       \rho_1(R)=\rho_2(R)=0\quad(R\in\mathcal C_{20})
 \tag{20rho}\label{eq:v15-twentieth-subspace-detection}
\]
が従う。
ここでは $\rho_i$ が $J_{20}$ 全体を消すことを仮定しない。
二つの汎関数は、確定した商 $\mathcal C_{20}/S_{20}$ の座標を読み取る。

特に、rawなTate係数 $d_{20}^{\mathrm T}$ に対して
$\rho_1(d_{20}^{\mathrm T})=\rho_2(d_{20}^{\mathrm T})=0$ が厳密に分かれば、
\[
                    d_{20}^{\mathrm T}\in S_{20}\subset J_{20}
\]
となる。低次の座標補正を全て含む正規化された20次項を $R_{20}$ とすると、
本版では別途、接触作用の計算から
\[
                         R_{20}-d_{20}^{\mathrm T}\in J_{20}
\]
を証明する。この二つの所属を合わせて $R_{20}\in J_{20}$ が従い、
20次項を通常接触変換で消去できる。
補正全体が $S_{20}$ に入ることや、補正全体に両汎関数が零となることは、
この結論に必要ないし、ここでは主張しない。
本版で後に示す22次から始まる正規形はこの道筋を用い、
通常20次商が一次元か二次元かという未解決問題とは論理的に独立している。

有限評価の一般的な背景も記しておく。
$R=\CC[a_{ij},F_\alpha]$ において
$P=\sum F_\alpha q(a)^\alpha$ は三次係数 $F_\alpha$ に関する一次式である。
Pfaffianの共通因子は1であり、係数に $q_i^3$ が含まれるため、
この一次式は $\CC[a]$ 上primitiveとなる。
$\CC(a)[F_\alpha]$ 上での既約性とGaussの補題から、$P$ は $R$ でも既約である。
従って $\mathcal C_{20}$ の元が $P=0$ 上で恒等的に消えれば $P$ で割り切れ、
商の次数と変換則からその商は $\mathcal U_{14}$ に属する。
よって
\[
                       \dim(\mathcal C_{20}/P\mathcal U_{14})=166.
\]
この166次元空間の基底を上記164接触像と $W_1,W_2$ が与えている。
ただし追加の式が法1009でその像に入るという観察だけから、
特性零での所属を結論してはならない。

\begin{remark}[残る一方向の問題]
通常商を完全決定するには、例えば次のいずれかが十分である。
特性零で独立な360番目の $\ker dF$ の元を構成すれば、
$\dim J_{20}=328$ と $\mathcal C_{20}/J_{20}\simeq\CC^2$ が従う。
逆に $P\mathcal U_{14}$ を法として165番目の独立接触像を構成すれば、
$\dim J_{20}=329$ と $\dim(\mathcal C_{20}/J_{20})=1$ が従う。
あるいは、上記二汎関数が $\mathcal G_{15}$ 全体の像を消すことを
普遍的に証明しても、前者の結論を得る。
これらは今後の課題であり、本版の実際の20次消去の前提ではない。
\end{remark}

% END rank/rank20_section_ja.tex


% BEGIN contact/contact_twentieth_ja.tex
% 第15版の研究用追記。2026-10-29。
% 構造命題は第13・14版の確定結果から証明する。
% 二十次の新しい有限所属判定は、このファイルでは仮定として区別する。

\section{二十次の接触作用と低次正規化の補正}
\label{sec:v15-contact-twentieth}

$P=P_F=F(q(a))$、$H=\det\operatorname{Hess}_F(q(a))$、$c=1/96$、
$J_{20}=dP(\mathcal G_{15})$ とする。
全ての場と単元は $F$ に多項式的かつ普遍同変とし、
座標変更の線形項と単元の定数項をそれぞれ恒等写像と1に固定する。
本節では、十八次における追加方向の消滅を二十次へ外挿しない。

\subsection{低次条件を保つ対数作用の二本の鎖}

第14版の接触対数補題\ref{lem:v14-contact-logarithm}を用いる。
低次の共変量の消滅により、$L=\log T$ の最小の次数増分は4であり、
\[
 L=\sum_{k\ge4}Q_k,\qquad
 Q_k=Y_{k+1}+b_k\in\mathcal G_{k+1}\oplus\mathcal U_k
\]
と書ける。$P+cH$ の19次までを保つ条件は
\[
 Q_kP+cQ_{k-4}H=0\quad(4\le k\le13),\qquad Q_j=0\ (j<4),
\]
二十次の作用は $Q_{14}P+cQ_{10}H$ である。
偶数増分の条件は
\begin{equation}\label{eq:v15-even-contact-chain}
 \begin{split}
 Q_4P&=0,\qquad Q_6P=0,\\
 Q_8P+cQ_4H&=0,\qquad Q_{10}P+cQ_6H=0,\\
 Q_{12}P+cQ_8H&=0.
 \end{split}
\end{equation}
このうち $Q_4,Q_8,Q_{12}$ の鎖は二十次の作用に現れず、全て零に選べる。
従って有効像を決めるのは $Q_6,Q_{10}$ の鎖である。
奇数増分の条件は独立であり、奇数増分を零にしてもこの像は変わらない。
よって全ての普遍接触変換と、奇座標・偶単元による接触変換の
二十次有効像は一致する。

\begin{lemma}[十一次の接触安定化場のHessian作用]
\label{lem:v15-eleventh-stabilizer-hessian}
$S_{10}=Y_{11}+b_{10}$ が $S_{10}P=0$ を満たせば
\[
 S_{10}H\equiv-\frac23b_{10}H\pmod{J_{20}}.
\]
特に $dP(Y_{11})=0$ なら $dH(Y_{11})\in J_{20}$ である。
\end{lemma}

\begin{proof}
$g=\nabla F(q)$、$h(p)=\det\operatorname{Hess}_F(p)$ と置く。
第13版の命題\ref{prop:v13-contact-rank-sixteenth}により
$\ker(dF:\mathcal D_{12}\to\mathcal C_{16})$ は二十のKoszul像で張られる。
従って $S_{10}P=0$ から
\[
 \delta q(Y_{11})=-\frac13b_{10}q+\beta g,
 \qquad \beta\in\mathcal B_{16}
\]
を得る。$h$ が $q$ に五次であることと $\beta$ の交代性から
\[
 S_{10}H=-\frac23b_{10}H-g^t\beta\nabla h(q).
\]
$\beta\nabla h(q)$ は $F^{14}a^{16}$ の次数を持つ
$\mathcal D_{16}$ の元である。
全次数のPfaffian微分の全射性
$\delta q:\mathcal G_{15}\twoheadrightarrow\mathcal D_{16}$ を用いれば、
最後の項は $J_{20}$ に入る。
\end{proof}

第13版の $H\mathcal U_6\subset J_{16}$ により、線形な持ち上げ
\[
 Z_{11}:\mathcal U_6\longrightarrow\mathcal G_{11},\qquad
 dP(Z_{11}(w))=Hw
\]
を一つ選べる。ここで $\dim\mathcal U_6=2$ である。

\begin{theorem}[二十次の有効像の有限縮約]
\label{thm:v15-effective-twentieth}
$c\ne0$ のとき、19次までの $P+cH$ を保つ普遍接触変換の二十次像は
\begin{equation}\label{eq:v15-effective-image}
 I_{20}=J_{20}+H\mathcal U_{10}
             +dH\bigl(Z_{11}(\mathcal U_6)\bigr)
\end{equation}
である。最後の項は高々二次元であり、各 $w$ に対する
$dH(Z_{11}(w))$ の $\mathcal C_{20}/J_{20}$ における類は純粋な持ち上げに依存しない。
\end{theorem}

\begin{proof}
$Q_6=Y_7+b_6$、$Q_6P=0$ を考える。
第10版の $dF:\mathcal D_8\to\mathcal C_{12}$ の単射性から
\[
 \delta q(Y_7)=-\frac13b_6q,\qquad Q_6H=-\frac23b_6H.
\]
$w=-b_6/6$ と置けば $Q_6H=4Hw$ であり、
この作用は $Q_6=w(\mathcal E-6)$ によって実現する。
ここで $\mathcal E=\sum a_i\partial_{a_i}$ はEuler導分である。
条件 $Q_{10}P+cQ_6H=0$ の一般解は
\[
 Q_{10}=-4cZ_{11}(w)+S_{10},\qquad S_{10}P=0
\]
と書ける。前補題と $P\mathcal U_{14}\subset J_{20}$ から
二十次作用は表示した右辺に含まれる。

逆に、$J_{20}$ は $Q_{14}$ で、$H\mathcal U_{10}$ は
$Q_{10}=u_{10}(\mathcal E-6)$ で実現する。
$w\in\mathcal U_6$ に対し
\[
 Q_6=w(\mathcal E-6),\qquad Q_{10}=-4cZ_{11}(w),
 \qquad Q_4=Q_8=Q_{12}=0
\]
とすれば全低次条件を満たし、二十次作用は $-4c^2dH(Z_{11}(w))$ である。
$c\ne0$ より全方向が実現する。
持ち上げの差については前補題の $b_{10}=0$ の場合を適用する。
\end{proof}

\begin{remark}[単元を残した持ち上げ]
$Hw=dP(Z'_{11}(w))+PB_{10}(w)$ であれば、
$Z_{11}(w)=Z'_{11}(w)+aB_{10}(w)/6$ と純粋な持ち上げに直せる。
そのHessian作用の差は $(5/3)HB_{10}(w)$ なので、
\eqref{eq:v15-effective-image}では $Z_{11}$ を $Z'_{11}$ に置き換えてよい。
現時点では右辺の追加部分が $J_{20}$ に含まれるとは主張していない。
\end{remark}

\begin{proposition}[十次単元の $P$ 倍部分]
\label{prop:v15-pdivisible-tenth-unit}
$\mathcal U_{10}$ の $P$ で割り切れる部分は二次元であり、
その $H$ 倍は $J_{20}$ に入る。
したがって $H\mathcal U_{10}$ の判定には、$P$ 倍部分を法とする
十六次元の商の基底を用いれば十分である。
\end{proposition}

\begin{proof}
商となる共変量の空間は
\[
 \operatorname{Hom}_{\mathrm{GL}(V)}\left(
 \operatorname{Sym}^{9}(\operatorname{Sym}^3V)
 \otimes\operatorname{Sym}^{4}(\Lambda^2V),\det(V)^7\right)
\]
である。Littlewood分解の補分割は
$(7,7,7,3,3)$、$(7,6,6,4,4)$、$(7,5,5,5,5)$ となる。
第14版の九次plethysmの全分解の証明書から、重複度は順に1、1、0である。
よって次元は二であり、整域における $P$ 倍の単射性と商の変換則により
これが $\mathcal U_{10}$ の $P$ 倍部分に一致する。
$u_{10}=Pv$ なら $Hv\in\mathcal U_{14}$ なので
$Hu_{10}=P(Hv)\in J_{20}$ である。
$\dim\mathcal U_{10}=18$ から商の次元は十六となる。
\end{proof}

\paragraph{二つの二次安定化方向の具体的な入力.}
第13版の \texttt{v13\_tensor\_rank\_certificate.json} の
\texttt{HU6\_coefficients} は、二つの $\mathcal U_6$ の基底について、
二十一の接触グラフ像と一つの補元による係数を保存している。
補元の係数はいずれも零である。先頭二十一係数で $Z'_{11}$ を作れば、
上の単元を残した持ち上げを得る。
\nolinkurl{v15_lower_contact_inputs.py} は全二十二行の有理数恒等式を再検算し、
この二つの持ち上げと十次単元の指標入力を集約する。
これは二十次の新しい所属判定とは区別する。

\subsection{二十次までの正確な補正台帳}

未正規化Tate係数を $f_n=d_n^{\mathrm T}$ と書く。
第10版のmerged規約を固定して
\[
 X_7=-\frac{15800}{6137}\Phi_7-\frac{689855}{883728}\Psi_7,
 \qquad X_8=\frac{4600}{2189}\operatorname{Alt}K_8,
\]
\[
 T_{10}(f)=(1+3U_4-V_6)f(a+X_7+X_8)
\]
とする。既存の普遍恒等式は
\[
 f_{10}=cH-3PU_4,\qquad
 f_{12}+dP(X_7)=PV_6,\qquad f_{13}+dP(X_8)=0
\]
である。単元を $(1+3U_4)(1-V_6)$ に置き換えることはしない。

第14次消去は $a-\Xi_9$ と $1-B_8$、第15次消去は10次場と9次単元を用いる。
これらの後の16次係数は
\begin{equation}\label{eq:v15-sixteenth-input}
 A_{16}=f_{16}+d(f_{10})(X_7)
             +3U_4\bigl(f_{12}+dP(X_7)\bigr)-V_6f_{10}
\end{equation}
であり、第13版により $A_{16}\in J_{16}$ である。
純粋な持ち上げ $\Xi_{11}\in\mathcal G_{11}$、
$dP(\Xi_{11})=A_{16}$ を選んで $a-\Xi_{11}$ で16次を消す。

\begin{proposition}[二十次の有限正規化補正式]
\label{prop:v15-exact-twentieth-correction}
上の有限正規化の後の二十次係数は正確に
\begin{equation}\label{eq:v15-exact-twentieth-correction}
 \begin{split}
 R_{20}^{\mathrm{fin}}={}&f_{20}+d(f_{14})(X_7)+d(f_{13})(X_8)
       +\frac12d^2P(X_8,X_8)\\
   &+3U_4\bigl(f_{16}+d(f_{10})(X_7)\bigr)-V_6f_{14}
       -c\,dH(\Xi_{11}).
 \end{split}
\end{equation}
16次を $A_{16}=dP(\Xi'_{11})+PB_{10}$ の規約で消す場合は、
最後の項を $-c(dH(\Xi'_{11})+B_{10}H)$ に置き換える。
その後の18次消去は二十次係数を変えない。
\end{proposition}

\begin{proof}
初回の座標変更で二十次に到達する項は
$f_{20}$、$d(f_{14})(X_7)$、$d(f_{13})(X_8)$、
$\tfrac12d^2P(X_8,X_8)$ である。
二つの $X_7$ は $P$ には18次、$X_7,X_8$ の混合は19次を与える。
それらが20次に達する他の元の係数は7、8次等であり零である。
三つ以上の座標補正による項は24次以上となる。
四次単元が掛かる16次係数は $f_{16}+d(f_{10})(X_7)$、
六次単元が掛かる14次係数は $f_{14}$ なので、最初の六項を得る。

14次消去が二十次へ作用するには12次係数が必要だが、初回変更により零である。
15次消去には11次係数が必要だが同様に零である。
16次消去は10次項 $cH$ へ作用するので、表示した最後の項を与える。
その二次Taylor項は26次以上である。
18次消去の13次座標場と12次単元が二十次へ作用するには8次係数が必要で、零である。
従って18次を消しても表示式は変わらない。
\end{proof}

ここには、13次項を消す偶数次座標場 $X_8$ の二階作用が現れる。
元の13次項が奇数次であることや、後で全次数偶化できることから、
この二十次の補正を省くことはできない。

\subsection{有効商における三つの補正への簡約}

$f_{14}\in J_{14}$ より、$Z_{14}\in\mathcal G_9$ を
$dP(Z_{14})=f_{14}$ と選ぶ。

\begin{theorem}[実際の二十次類の代表元]
\label{thm:v15-twentieth-effective-representative}
第14版の正規形における実際の二十次類は
\begin{equation}\label{eq:v15-twentieth-effective-representative}
 \boxed{\displaystyle
 [R_{20}]=\left[f_{20}+d^2P(X_7,Z_{14})
       -\frac12d^2P(X_8,X_8)-c\,dH(\Xi_{11})\right]
       \quad\text{in }\mathcal C_{20}/I_{20}.
 }
\end{equation}
純粋な持ち上げ $Z_{14}$、$\Xi_{11}$ の選択はこの類に影響しない。
$P$ 倍の単元項を残した持ち上げを用いても有効類は同じである。
\end{theorem}

\begin{proof}
恒等式 $f_{13}=-dP(X_8)$ を $X_8$ で微分すると
\[
 d(f_{13})(X_8)=-d^2P(X_8,X_8)-dP(DX_8\cdot X_8).
\]
最後の場は $\mathcal G_{15}$ に属するので、奇数次項を消した二階補正は
$-\tfrac12d^2P(X_8,X_8)$ に簡約される。
同様に $f_{14}=dP(Z_{14})$ を $X_7$ で微分すれば
\[
 d(f_{14})(X_7)\equiv d^2P(X_7,Z_{14})\pmod{J_{20}}.
\]

既存の所属 $f_{16}\in J_{16}$、$f_{14}\in J_{14}$ から
$U_4f_{16},V_6f_{14}\in J_{20}$ である。
さらに $f_{10}=cH-3PU_4$ と、第13版の
$dH(\mathcal G_7)\subset J_{16}$ を用いると
$U_4d(f_{10})(X_7)\in J_{20}$ である。
これで表示した三補正の式を得る。

二つの純粋な $Z_{14}$ の差 $K_9$ には $dP(K_9)=0$ が成り立ち、
これを $X_7$ で微分すれば $d^2P(X_7,K_9)\in J_{20}$ である。
$\Xi_{11}$ の差については補題\ref{lem:v15-eleventh-stabilizer-hessian}を適用する。
単元項を残した持ち上げと純粋なものとの差はEuler場であり、
$P$ の二階微分への作用は $J_{20}$、$H$ への作用は
$H\mathcal U_{10}\subset I_{20}$ に入る。
\end{proof}

\begin{corollary}[十七次・十九次の明示消去は不要]
\label{cor:v15-twentieth-without-odd-normalizers}
上の有限正規化と18次消去の後に17次・19次の項が残っていても、
式\eqref{eq:v15-twentieth-effective-representative}は全次数偶正規形の
二十次有効類を与える。
\end{corollary}

\begin{proof}
有限正規化後を $f=P+cH+R_{17}+R_{19}+R_{20}+O(a^{21})$ とし、
第14版の偶正規形と比較する接触作用素を $T$ とする。
$Tf-f$ は17次以上なので、対数補題から
$Tf-f\equiv Lf\pmod{(a)^{21}}$ である。
最小次数増分が4であるため、$L(R_{17}+R_{19})$ も21次以上である。
従って二十次までの比較は $L(P+cH)$ だけによる。
偶数増分の条件は\eqref{eq:v15-even-contact-chain}のままであり、
17次・19次消去の非零な右辺は奇数増分の条件にだけ入る。
二十次の差は $I_{20}$ に含まれる。

一般の偶化に含まれ得る六次や八次の座標場をTaylor展開すると、
それぞれの二階作用が二十次に見える場合がある。
ここでは低次を保つ全条件を満たした作用素全体の対数を用いており、
そのような項の相殺も含んだ結論である。
\end{proof}

\subsection{二階補正を小さい有限族へ落とす一般補題}

\begin{lemma}[Jacobian像を法とした二階作用]
\label{lem:v15-hessian-pairing-quotient}
$Y\in\mathcal G_r$、$Z\in\mathcal G_s$、$N=r+s+4$ とする。
$d^2P(Y,Z)$ の $\mathcal C_N/J_N$ における類は、
\[
 [dP(Y)]\in J_{r+5}/P\mathcal U_{r-1},\qquad
 [dP(Z)]\in J_{s+5}/P\mathcal U_{s-1}
\]
だけに依存する。したがって、これらの商からの対称な双線形写像を定める。
\end{lemma}

\begin{proof}
$dP(K)=0$ を $Z$ で微分すると
$d^2P(K,Z)=-dP(DK\cdot Z)\in J_N$ である。
また $dP(Y)=Pu$ のとき、$Y-au/6$ は $dP$ の核に入る。
Eulerの関係を微分して
\[
 d^2P(a u/6,Z)=\frac56u\,dP(Z)\in J_N
\]
を得る。よって表示した各商を通る。
\end{proof}

\begin{proposition}[二十次の補正を全て判定する有限族]
\label{prop:v15-finite-lower-tests}
以下の四群の全てが $J_{20}$ に含まれると仮定する。
\begin{enumerate}
 \item $H\mathcal U_{10}$。$P$ 倍の単元方向は恒等的に $J_{20}$ に入るので省ける。
 \item 第13版の $J_{16}/P\mathcal U_{10}$ の基底を持ち上げる
 二十一の場 $\Gamma_j^{(11)}$ に対する $dH(\Gamma_j^{(11)})$。
 \item 第14版の $J_{12}/P\mathcal U_6$ の基底を持ち上げる二つの場
 $\Gamma_i^{(7)}$ と $J_{14}/P\mathcal U_8$ の基底を持ち上げる七つの場
 $\Gamma_j^{(9)}$ に対する、十四の共変量
 \[
 \operatorname{Hess}_F(q)
    (\delta q\Gamma_i^{(7)},\delta q\Gamma_j^{(9)}).
 \]
 \item $J_{13}=\mathcal C_{13}$ の非零生成元を持ち上げる一つの場
 $\Gamma^{(8)}$ に対する
 \[
 \operatorname{Hess}_F(q)(\delta q\Gamma^{(8)},\delta q\Gamma^{(8)}).
 \]
\end{enumerate}
このとき $I_{20}=J_{20}$ であり、実際の二十次類は
$[R_{20}]=[d_{20}^{\mathrm T}]$ となる。
本命題は二十次の有限所属判定が完了するまでは条件付きである。
\end{proposition}

\begin{proof}
(1),(2)と補題\ref{lem:v15-eleventh-stabilizer-hessian}から
$dH(\mathcal G_{11})\subset J_{20}$ である。
実際、任意の場から二十一の基底像とEuler単元場を除けば $dP$ の核に入る。
従って二つの安定化方向も $dH(\Xi_{11})$ も $J_{20}$ に入り、$I_{20}=J_{20}$ となる。

合成関数の微分により
\[
 d^2P(Y,Z)=\operatorname{Hess}_F(q)(\delta qY,\delta qZ)
                      +dF(d^2q(Y,Z)).
\]
ここでは $d^2q(Y,Z)\in\mathcal D_{16}$ なので最後の項は $J_{20}$ に入る。
前補題と $\dim(J_{12}/P\mathcal U_6)=2$、
$\dim(J_{14}/P\mathcal U_8)=7$ により混合補正は(3)の十四方向で尽くされる。
また $\dim J_{13}=1$、$\mathcal U_7=0$ なので、
二つの八次場による対称二階補正は(4)の一方向だけで尽くされる。
これで式\eqref{eq:v15-twentieth-effective-representative}の三補正が全て消える。
\end{proof}

\begin{remark}[判定が零でない場合]
上の有限族のいずれかが通常商で非零なら、その所属を仮定してはならない。
その場合でも有効像は\eqref{eq:v15-effective-image}で決まり、
実際の類は\eqref{eq:v15-twentieth-effective-representative}で計算できる。
特に八次場の平方補正は元のTateの $K_8$ から直接評価できる。
混合補正は $d(f_{14})(X_7)$ として方向微分してもよく、
その違いは $J_{20}$ に入る。
\end{remark}

% END contact/contact_twentieth_ja.tex


% BEGIN detector/v15_detector_notes_ja.tex
% 第15版用。S20の328次元独立性はrank証明書の節を参照する。
\section{20次の既知接触像を判定する三点汎関数}
\label{sec:v15-three-point-detectors}

20次の普遍境界空間は330次元である。
本節では、既に構成した328次元の接触像の部分空間への所属を、
三つの値だけで判定する二つの汎関数を与える。
この所属判定は、通常接触写像の全階数やKoszul完全性を必要としない。
通常接触像全体との一致がまだ証明されていない段階では、
汎関数の非零値だけから非零接触類を結論しない。

\subsection{既知像と二つの補元}

164個の15次座標場を $\Gamma_{15,\alpha}$ とし、
\[
 \mathcal S_{20}:=
 P\mathcal U_{14}
 +\left\langle dP(\Gamma_{15,\alpha}):1\le\alpha\le164\right\rangle
 \subset\mathcal J_{20}:=dP(\mathcal G_{15})
\]
とおく。ここで $P\mathcal U_{14}\subset\mathcal J_{20}$ は
$dP(a u/6)=Pu$ というEulerの等式による。
一般の $P=0$ の整数入力における164列の非零minorにより、
これらの列は $P\mathcal U_{14}$ を法として独立である。
$\dim\mathcal U_{14}=164$ と $P$ 倍の単射性から
\begin{equation}\label{eq:v15-known-image-dimension}
 \dim\mathcal S_{20}=328
\end{equation}
を得る。
この下界の証明と、通常像 $\mathcal J_{20}$ の上界の証明は区別する。

第11版と同じ正規化の6頂点単元グラフを $J_6$ とし、
\[
 W_1=\Theta_8U_4^3,\qquad W_2=\Theta_8J_6^2
\]
とおく。両者は $F$ 次数15、$a$ 次数20を持つ自然な共変量であり、
$\mathcal C_{20}$ に属する。

\subsection{三つの入力}

\begin{align*}
 F_D&=x_0^3+x_1^3+x_2^3+x_3^3-2x_4^3,\\
 F_H&=x_0^3+x_1^3+x_2^3+x_3^3+4x_4^3-6x_2x_3x_4,\\
 a_b&=b\,e_0\wedge e_1
       +b^{-1}(e_2-e_4)\wedge(e_3-e_4)
\end{align*}
とおく。Pfaffianの規約は前節までと同じである。
任意の $b\ne0$ に対し
\[
 q(a_b)=(0,0,1,1,1),\qquad F_D(q(a_b))=F_H(q(a_b))=0
\]
となる。三つの評価点を
\[
 z_1=(F_D,a_1),\qquad z_2=(F_D,a_2),\qquad z_3=(F_H,a_1)
\]
とする。$F_D$ は非零対角係数を持つ滑らかな三次式である。
$F_H$ の最後の三変数の部分はHesse型であり、
対角係数の積と混合tensor係数の判別条件は $4+8(-1)^3=-4\ne0$ である。
最初の二変数は分離しているので $F_H$ も滑らかである。

この族では、グラフの整数縮約から
\[
\begin{array}{c|ccc}
 &\Theta_8&U_4&J_6\\\hline
(F_D,a_b)&-2b^2&0&24b^6\\
(F_H,a_b)&-4b^2&5b^4&528b^6
\end{array}
\]
を得る。$F_H$ のtensorの規約は
$F_{234}=F_{243}=\cdots=-1$ であり、
多項式の混合項は $-6x_2x_3x_4$ となる。
この係数6を省略しない。

\begin{proposition}[三点の厳密な双対汎関数]
\label{prop:v15-exact-dual-detectors}
$\operatorname{ev}_i(f)=f(z_i)$ とおき、
\begin{align}
 \rho_1&=\frac{-22168\operatorname{ev}_1+82\operatorname{ev}_2
                    -1365\operatorname{ev}_3}{682500},
                    \label{eq:v15-rho1}\\
 \rho_2&=\frac{4\operatorname{ev}_1-\operatorname{ev}_2}{18869760}
                    \label{eq:v15-rho2}
\end{align}
と定める。このとき
\[
 \rho_i(\mathcal S_{20})=0,\qquad
 \rho_i(W_j)=\delta_{ij},\qquad
 \ker(\rho_1,\rho_2)=\mathcal S_{20}
\]
が成り立つ。
\end{proposition}

\begin{proof}
$P(z_i)=0$ なので、全ての $P\mathcal U_{14}$ の方向は三点とも零である。
164個のグラフについて、三つの値を有理数上で再計算した。
各列の値は厳密に $(c,4c,-16c)$ の形を持つ。
従って両汎関数は全ての164列を消す。
この等式は有限体での試験に留めず、非有界整数のtensor縮約と有理数演算で検証した。

選択した一つの接触列と二つの補元の評価行列は
\[
 \begin{pmatrix}
 -96&0&-1152\\
 -384&0&-18874368\\
 1536&-500&-1115136
 \end{pmatrix}
\]
であり、その行列式は
\[
 905748480000\ne0
\]
である。逆行列の対応する二行が
式\eqref{eq:v15-rho1}--\eqref{eq:v15-rho2}を与え、
$\rho_i(W_j)=\delta_{ij}$ を満たす。
両汎関数は $\mathbb Q$ 上で定義され、複素数へ係数拡大して用いる。
従って $(\rho_1,\rho_2):\mathcal C_{20}\to\mathbb C^2$ は階数2で、
その核の次元は $330-2=328$ である。
式\eqref{eq:v15-known-image-dimension}と包含
$\mathcal S_{20}\subset\ker(\rho_1,\rho_2)$ から核の等式を得る。
ここでは $\mathcal S_{20}=\mathcal J_{20}$ は仮定していない。
\end{proof}

従って任意の $f\in\mathcal C_{20}$ に対し
\[
 f-\rho_1(f)W_1-\rho_2(f)W_2\in\mathcal S_{20}\subset\mathcal J_{20}
\]
である。特に二つの値がともに零であることは、接触像への所属の十分条件となる。
通常接触像が既知像より大きい可能性を残しても、この十分条件は有効である。

\subsection{低次正規化から来る補正の消滅}

$h_F(p)=\det\operatorname{Hess}_F(p)$、$H(F,a)=h_F(q(a))$ と書く。
$dH$ は $a$ 変数方向の微分、$d_qh_F$ は $q$ 変数方向の微分とする。
三つの入力において、$\operatorname{Hess}_F(q)$ の最初の二行と二列は零であり、
残りの $3\times3$ ブロックは可逆である。
従って全Hessianの階数は3で、
\[
 H(z_i)=0,\qquad d_qh_F(q(a))(z_i)=0
\]
となる。後者は $5\times5$ 行列の全ての $4\times4$ 余因子が零であることによる。
従って任意の $u\in\mathcal U_{10}$ と $Z\in\mathcal G_{11}$ に対し
\[
 (Hu)(z_i)=0,\qquad
 dH(Z)(z_i)=d_qh_F(q(a))(dq(Z))(z_i)=0
\]
である。これらは $\mathcal C_{20}$ の元なので、
命題\ref{prop:v15-exact-dual-detectors}から
\begin{equation}\label{eq:v15-hessian-secondary-in-known-image}
 H\mathcal U_{10}+dH(\mathcal G_{11})\subset\mathcal S_{20}
\end{equation}
を得る。ここでは $\mathcal U_{10}$ や $\mathcal G_{11}$ の評価基底を列挙する必要はない。

次に $\mathcal J_{12}/P\mathcal U_6$ の2個の代表場と、
$\mathcal J_{14}/P\mathcal U_8$ の7個の代表場を、それぞれ
$X_{7,i}$、$Y_{9,j}$ とする。
これらは第14版で用いた同じグラフ場である。
14個の普遍境界共変量
\[
 B_{ij}:=\operatorname{Hess}_F(q)
          \bigl(dq(X_{7,i}),dq(Y_{9,j})\bigr)
 \quad (1\le i\le2,\ 1\le j\le7)
\]
について三点で厳密評価し、全ての $i,j$ に対して
\[
 \rho_1(B_{ij})=\rho_2(B_{ij})=0
\]
を確認した。従って $B_{ij}\in\mathcal S_{20}$ である。
同様に、既知の1次元像 $\mathcal J_{13}$ を生成する8次場 $X_8$ の代表を一つ選ぶと、
\[
 \operatorname{Hess}_F(q)(dq(X_8),dq(X_8))\in\mathcal S_{20}
\]
となる。この代表場の $dP$ が非零であることも、
別の一般整数入力で有理数上の縮約により確認した。

これらの有限検査から、全ての必要な場へ広げる理由を記す。
$dP(X)=0$ なら、その恒等式を場 $Y$ 方向に微分することで
\[
 P''(X,Y)=-dP(DX[Y])\in\mathcal J_{20}
\]
を得る。Euler単元場が一方に入る場合も、
$\operatorname{Hess}_F(q)(q,v)=2\,dF(v)$ により接触像に入る。
残る方向は上の $2\times7$ 個と8次像の1個で尽くされる。
さらに
\[
 P''(X,Y)=\operatorname{Hess}_F(q)(dq(X),dq(Y))
            +dF\bigl(d^2q(X,Y)\bigr)
\]
の最後の項は $d^2q(X,Y)\in\mathcal D_{16}$ と
全次数Pfaffian全射定理から $\mathcal J_{20}$ に入る。
従って
\begin{equation}\label{eq:v15-quadratic-corrections-contact}
 P''(\mathcal G_7,\mathcal G_9)
  +P''(\mathcal G_8,\mathcal G_8)\subset\mathcal J_{20}
\end{equation}
を得る。
この議論の途中で生じる任意の接触項について、
まだ両 $\rho_i$ が零になると仮定しているのではない。
それらは直接 $\mathcal J_{20}$ に属するため、接触同値の計算で捨ててよい。

\subsection{実際の20次項への適用と再現範囲}

filtered contact作用の計算で得られた代表は、既存の規約では
\[
 R_{20}^{\mathrm{eff}}
 =d_{20}^{T}+P''(X_7,Z_{14})
       -\tfrac12P''(X_8,X_8)-\tfrac1{96}dH(\Xi_{11})
       \pmod{\mathcal J_{20}}
\]
という形に書ける。
式\eqref{eq:v15-hessian-secondary-in-known-image}と
\eqref{eq:v15-quadratic-corrections-contact}により
\[
 [R_{20}^{\mathrm{eff}}]=[d_{20}^{T}]
 \quad\text{in }\mathcal C_{20}/\mathcal J_{20}
\]
である。三点で独立に再計算したfresh Tate値は、
\[
 \bigl(d_{20}^{T}(z_1),d_{20}^{T}(z_2),d_{20}^{T}(z_3)\bigr)
 =(A,4A,-16A),\qquad
 A=\frac{1112989093269228481}{187535250000000}
\]
である。式\eqref{eq:v15-rho1}--\eqref{eq:v15-rho2}に代入すると
$\rho_1(d_{20}^{T})=\rho_2(d_{20}^{T})=0$ となる。
従って $d_{20}^{T}\in\mathcal S_{20}\subset\mathcal J_{20}$ であり、
実際の20次項は接触像に入り消去できる。
この結論には、通常商の次元が2であることの証明を先に必要としない。
fresh Tateブロックの再構成と次数の切断の検査は、Tate計算の節とその証明書で行う。

\texttt{v15\_rho20\_three\_cases.json} は三つの入力を保持する。
\nolinkurl{v15_rho20_dual_certificate.py} は
164接触列、二つの補元、14個の混合Hessianと1個の平方を有理数上で再計算する。
分母を持つ $a$ は一度 $Da$ にして整数縮約し、20次の斉次性により $D^{20}$ で割り戻す。
$U_4,J_6,\Theta_8$ についても各々の次数で同じ処理を行う。
探索時には、先に完成した二つの対角点の厳密値を再利用してHesse点を計算した。
最終証明書では全三点をfreshに再計算した。
既定の再実行も全三点をfreshに再計算する。
値を入力とする監査と、tensor縮約そのものの再計算は区別して記録する。

% END detector/v15_detector_notes_ja.tex


% BEGIN tate/v15_tate_notes_ja.tex
\section{二十次Tate係数の新規再構成とHesse型Clifford計算}
\label{sec:v15-tate-twentieth}

本節では、前節の三点における未正規化係数 $d_{20}^{\mathrm T}$ を、
TateのMaurer--Cartan再帰から新たに計算する。
既存の十八次ブロックを二十次に流用するのではなく、
必要な内部変数の次数を改めて導き、全ての入力点で最初から再構成した。
二つの対角点に加え、混合項 $-6x_2x_3x_4$ を持つHesse点も扱う。

\subsection{二十次に必要な切断の範囲}

補正を $A=\sum_{r\ge1}A_r$、定数ブロックを
$K_r:\Lambda^4V\to V$ と書く。規約は
\[
 R_n=\sum_{i=1}^{n-1}A_iA_{n-i},\qquad A_n=-hR_n,
 \qquad d^{\mathrm T}=\tfrac12\det\!\left(\sum_{r\ge1}K_r\right)
\]
である。各 $K_r$ は $a$ に $r$ 次で、$K_1=-a$、
$K_2=\operatorname{Hess}_F(q)/3$ となる。

\begin{lemma}[二十次の厳密な内部次数切断]
\label{lem:v15-tate-twentieth-truncation}
二十次の行列式係数には $K_1,\ldots,K_{16}$ だけが必要である。
これらは、$A_1,\ldots,A_{12}$ を全て求めた後、
$R_{13},R_{14},R_{15},R_{16}$ をそれぞれ内部 $x$ 次数
$4,3,2,1$ 以下で計算すれば厳密に得られる。
そのとき $A_{13},A_{14},A_{15},A_{16}$ はそれぞれ
$x$ 次数 $3,2,1,0$ 以下まで正確である。
\end{lemma}

\begin{proof}
行列式の各単項式は、次数が少なくとも1の五つのブロックの積である。
その総次数を20以下にするには、各因子の次数は高くとも16でよい。
Clifford積の収縮は $x$ 次数を下げず、Eulerホモトピー $h$ はちょうど1下げる。
また内部重みの関係
\[
 2\deg_x-\#\theta+\#\eta=3
\]
により、$A_r$ の $x$ 次数は常に4以下である。
従って $A_1,\ldots,A_{12}$ の全計算で前半を確保できる。
後半で既に求めた $A_i$ は $x$ 次数 $16-i$ 以下まで正確である。
$R_n$ の積で $i\le n-1$ なら $16-i\ge17-n$ なので、
$x$ 次数 $17-n$ 以下の曲率に必要な係数は全て保存されている。
これを帰納的に用いれば、表示した切断で各定数ブロックを失わない。
\end{proof}

第14版の $A_{12}$ は $x$ 次数2以下までの計算であった。
今回の全 $A_{12}$ は新たに計算しており、前版の切断を全項と見なしてはいない。

\subsection{三変数の混合項を持つ自然なClifford積}

$F=x_0^3+x_1^3+f(x_2,x_3,x_4)$ とし、
\[
 B_{ij}(x)=\tfrac16\frac{\partial^2F}{\partial x_i\partial x_j}
\]
とおく。自然なClifford生成元は
\[
 \eta_i=\iota_i-\sum_jB_{ij}(x)\theta_j,
 \qquad \{\eta_i,\eta_j\}=-2B_{ij},\qquad
 \{\theta_i,\eta_j\}=\delta_{ij}
\]
を満たす。ここで $\theta_i$ は外積作用、$\iota_i$ は収縮作用である。
記号としての $\eta$ は外積で反対称化し、$\theta$ を通常順に置く。
これが既存の普遍同変なEulerホモトピーを適用する記号基底である。

\begin{lemma}[Hesse型の外積記号とPBW基底]
\label{lem:v15-hesse-symbol-product}
混合項が $x_2,x_3,x_4$ のみに含まれる場合、
$0,1$ の各一変数Clifford代数と、最後の三変数のClifford代数との
次数付きテンソル積によって全代数を計算できる。
最後の三変数で外積記号から順序付きPBW基底への変換を $S$ とすると
\[
 S=1+C,\qquad C^2=0,\qquad S^{-1}=1-C.
\]
従って自然な記号の積は
\[
 \mu_{\mathrm{ext}}=S^{-1}\circ\mu_{\mathrm{PBW}}\circ(S\otimes S)
\]
で計算できる。
\end{lemma}

\begin{proof}
異なるブロック間のClifford係数は零なので、ブロック間には次数付き符号だけが現れる。
三変数ブロックは六つの奇生成元を持ち、PBW基底は $2^6=64$ 個である。
$i<j<k$ について外積記号は
\[
 \eta_i\wedge\eta_j=\eta_i\eta_j+B_{ij},
\]
\[
 \eta_i\wedge\eta_j\wedge\eta_k
  =\eta_i\eta_j\eta_k+B_{ij}\eta_k-B_{ik}\eta_j+B_{jk}\eta_i
\]
となる。$C$ はこの一対の収縮であり、最後の三変数には
$\eta$ が高々三つしかないので二対の収縮はできない。
従って $C^2=0$ であり、表示した逆変換が厳密に成立する。
$C$ は $x$ 次数を上げるため、前補題の切断の議論も保たれる。
$x=0$ では $S=1$ なので、定数ブロックの取り出しは従来の規約と一致する。
\end{proof}

実装では、三変数ブロックの語をClifford関係で正規順序へ還元してから、
上の $S^{-1}$ を適用する。
PBW順序のままEulerホモトピーを適用すると自然な外積記号と異なるため、
この変換を省略しない。
基本作用素は
\[
 D=\sum_i x_i\eta_i+\frac23\sum_i\partial_iF\,\theta_i
\]
であり、$F_H$ の混合項からは各対応する二次単項式に係数 $-4$ が現れる。
これを含む完全な $B_{ij}$ を用いて、独立な32次元外積表現も実装した。
独立表現ではPBW還元も $S$ の式も使わず、実際の $\eta_i$ 作用を反対称化する。

\subsection{三点の二十次係数}

\begin{proposition}[新規に再構成した三点Tate値]
\label{prop:v15-fresh-tate-twentieth}
第\ref{sec:v15-three-point-detectors}節の三点 $z_1,z_2,z_3$ に対し、
\[
 A_*=\frac{1112989093269228481}{187535250000000}
\]
とおくと
\[
 \boxed{\bigl(d_{20}^{\mathrm T}(z_1),d_{20}^{\mathrm T}(z_2),
                       d_{20}^{\mathrm T}(z_3)\bigr)
          =A_*(1,4,-16).}
\]
従って $\rho_1(d_{20}^{\mathrm T})=\rho_2(d_{20}^{\mathrm T})=0$ であり、
$d_{20}^{\mathrm T}\in\mathcal S_{20}\subset\mathcal J_{20}$ である。
\end{proposition}

\begin{proof}
各点で $A_1$ から再帰を開始し、前補題の範囲を有理数上で再構成した。
各段の曲率の閉性とMaurer--Cartan方程式を保存した次数範囲で直接照合した。
二十次係数は
\[
\begin{array}{c|c}
 z_1&1112989093269228481/187535250000000\\
 z_2&1112989093269228481/46883812500000\\
 z_3&-1112989093269228481/11720953125000
\end{array}
\]
である。五行の行列式を120置換による展開と列部分集合による動的計画法で
それぞれ計算し、0次から20次まで全係数の一致を確認した。
最後の二つの分母は最初の分母のそれぞれ $1/4,1/16$ である。
両汎関数の整数分子は、この比に対して
\[
 -22168+82\cdot4-1365(-16)=0,\qquad4-4=0
\]
となる。命題\ref{prop:v15-exact-dual-detectors}を用いて所属を得る。
\end{proof}

\paragraph{独立照合の範囲.}
三点以外にも、前版で用いた $F=x_0^3+x_1^3+x_2^3+x_3^3-4x_4^3$ の
六点を二十次まで新たに再構成した。
その標準点では、各係数を個別の有理数として計算する方法と、
共通分母を分離して整数で累積する方法の両方を実行した。
全16ブロック、0次から20次までの全行列式係数、および低次場の値が完全に一致した。
同じ点の18次以下は第14版の値にも一致した。
従って新規計算は異なる九点、標準点の独立再構成を含めて十回である。

Maurer--Cartan再帰と行列式の照合は有理数上で行った。
独立な32次元外積作用素による0次から12次の照合、および定数ブロックの
独立作用素との照合は素数 $1000003$ 上で行っている。
この二つの検証範囲を混同しない。
Hesse型Clifford生成元の反交換関係は有理数上で全て確認した。

\paragraph{再現用ソース.}
\texttt{source\_v15/tate/} 内に、対角入力の
\nolinkurl{v15_tate_twentieth.py}、混合項の積を実装する
\nolinkurl{v15_hesse_clifford.py} とその再帰、
三点の最終所属を確認する
\nolinkurl{v15_actual_twentieth_certificate.py} を置いた。
公開入力、全定数ブロック、有理数の行列式係数、各再帰段の記録を添付する。
最終所属の監査は保存済みブロックから行列式を再計算する。
それと、最初からTate再帰を再構成する実行とは区別した。

% END tate/v15_tate_notes_ja.tex


% BEGIN contact/twenty_second_normal_form_ja.tex
% 第15版の確定稿。三点のfresh Tate係数と既知像への所属に基づく。
% 通常接触像の全階数が確定したことは仮定せず、主張もしない。

\section{二十次の消去と二十二次開始の正規形}
\label{sec:v15-twenty-second}

本節では、通常接触像の一つの既知部分空間への所属から、
実際の二十次項を消去する。
以下、$P=P_F$、$H=H_{10}$、$c=1/96$、
$J_{20}=dP(\mathcal G_{15})$ とし、
$\mathcal S_{20}\subset J_{20}$ は
命題\ref{prop:v15-exact-dual-detectors}の328次元の既知像とする。

\begin{theorem}[二十二次開始の普遍代数的収束偶正規形]
\label{thm:v15-twenty-second-normal-form}
二十次の有効接触像は通常接触像と一致し、$I_{20}=J_{20}$ である。
実際の境界の二十次類は零であり、任意の滑らかな三次式 $F$ に対して
\begin{equation}\label{eq:v15-twenty-second-normal-form}
 \boxed{\displaystyle
 d^{\mathrm T}\sim
 P_F+\frac1{96}\det\operatorname{Hess}_F(q(a))
       +\sum_{\substack{n\ge22\\ n\ \mathrm{even}}}R_n,
 \qquad R_n\in\mathcal C_n
 }
\end{equation}
という正規形が存在する。
接触変換は普遍同変な代数的冪級数で、座標変更の線形項は恒等写像、
単元の定数項は1であり、各形式係数は $F$ の係数に多項式的である。
各固定 $F$ では、変換・逆変換・正規形の級数は原点近傍で収束する。
この結論は $\mathcal S_{20}=J_{20}$ や通常商の次元の確定を必要とせず、
二十二次以降の消去を主張しない。
\end{theorem}

\begin{proof}
命題\ref{prop:v15-exact-dual-detectors}により
\[
 \ker(\rho_1,\rho_2)=\mathcal S_{20}\subset J_{20}.
\]
三点で $H$ と $d_qh$ が点ごと零であることから、
式\eqref{eq:v15-hessian-secondary-in-known-image}は
\[
 H\mathcal U_{10}+dH(\mathcal G_{11})
       \subset\mathcal S_{20}\subset J_{20}
\]
を与える。従って定理\ref{thm:v15-effective-twentieth}の
二つの追加安定化方向も通常像に入り、$I_{20}=J_{20}$ である。
さらに式\eqref{eq:v15-quadratic-corrections-contact}により
\[
 d^2P(\mathcal G_7,\mathcal G_9)
       +d^2P(\mathcal G_8,\mathcal G_8)\subset J_{20}.
\]
定理\ref{thm:v15-twentieth-effective-representative}の
混合二階補正、八次場の平方補正、十一次場のHessian作用は、
以上により全て $J_{20}$ に属する。
よって、有限正規化後の二十次係数について
\begin{equation}\label{eq:v15-raw-and-normalized-twentieth-class}
 [R_{20}^{\mathrm{fin}}]=[d_{20}^{\mathrm T}]
           \quad\text{in }\mathcal C_{20}/J_{20}
\end{equation}
である。

次に、三点 $z_1,z_2,z_3$ の実際のTate再帰をそれぞれ再構成する。
対角三次式の二点とHesse型三次式の一点から得られた未正規化値は
\begin{equation}\label{eq:v15-three-fresh-twentieth-values}
 \bigl(d_{20}^{\mathrm T}(z_1),d_{20}^{\mathrm T}(z_2),
                d_{20}^{\mathrm T}(z_3)\bigr)
       =(\lambda_{20},4\lambda_{20},-16\lambda_{20}),
 \qquad
 \lambda_{20}=\frac{1112989093269228481}{187535250000000}.
\end{equation}
これらは $A_1,\ldots,A_{12}$ と必要な $A_{13},\ldots,A_{16}$ の切断を
新たに作り、十六個の定数ブロックから求めた値である。
Hesse型入力では自然なClifford積と独立な外積作用素を用いており、
三値の比を仮定して係数を決めてはいない。
命題\ref{prop:v15-exact-dual-detectors}の整数重みを用いると
\[
 -22168\lambda_{20}+82(4\lambda_{20})
       -1365(-16\lambda_{20})=0,\qquad
 4\lambda_{20}-4\lambda_{20}=0.
\]
従って $\rho_1(d_{20}^{\mathrm T})=\rho_2(d_{20}^{\mathrm T})=0$ であり、
\[
 d_{20}^{\mathrm T}\in\mathcal S_{20}\subset J_{20}.
\]
式\eqref{eq:v15-raw-and-normalized-twentieth-class}から
実際の二十次類は零である。

第14版の定理\ref{thm:v14-twentieth-normal-form}の全次数偶正規形を
\[
 f_{\mathrm{ev}}=P+cH+R_{20}^{\mathrm{ev}}
                       +R_{22}^{\mathrm{ev}}+\cdots
\]
と書く。
系\ref{cor:v15-twentieth-without-odd-normalizers}により
$R_{20}^{\mathrm{ev}}-R_{20}^{\mathrm{fin}}\in I_{20}=J_{20}$ である。
従って $R_{20}^{\mathrm{ev}}\in J_{20}$ であり、
普遍多項式な場 $\Xi_{15}\in\mathcal G_{15}$ を
\[
 dP(\Xi_{15})=R_{20}^{\mathrm{ev}}
\]
となるように選べる。
これは有限次元の普遍共変量空間での線形な持ち上げであり、
$F$ の判別式や評価点の値による除算を伴わない。

奇座標変更 $a\mapsto a-\Xi_{15}(F,a)$ で引き戻せば二十次が消える。
$\Xi_{15}$ は奇数次なので、全次数の偶性は保たれる。
$H$ への最初の寄与は24次、$P$ の二階Taylor項は34次からであり、
既定の低次正規形は変わらない。
従って次に残り得る次数は22である。

今回加えた座標変更は有限な多項式で、原点でのJacobianが恒等写像である。
その形式逆は代数的冪級数であり、固定した $F$ では解析的にも収束する。
第14版で構成した代数的・収束的な接触変換との有限合成は、
これらの性質および普遍同変性と係数の $F$ 多項式性を保つ。
これで所要の正規形を得る。
\end{proof}

\begin{remark}[通常商の同定との区別]
二つの共変量 $W_1=\Theta_8U_4^3$ と $W_2=\Theta_8J_6^2$ は、
$\mathcal C_{20}/\mathcal S_{20}$ の基底を与える。
通常接触像が $\mathcal S_{20}$ より大きい可能性を残す段階では、
これらを $\mathcal C_{20}/J_{20}$ の基底とは呼ばない。
同様に、$\rho_i$ の非零値から非零接触類を結論しない。

本定理で必要なのは、二つの汎関数値 $\rho_1,\rho_2$ がともに零になると
既知像 $\mathcal S_{20}$ に入るという十分条件である。
これにより実際の二十次項の消去は確定する。
通常接触商の完全な同定と、二十二次以降の普遍商・実際の類の決定は
別の問題として残る。
\end{remark}

% END contact/twenty_second_normal_form_ja.tex

\clearpage
% BEGIN theory/v16_plucker_hessian_structure_ja.tex
% 第16版：第15版で未採用だった Pluecker/Hessian メモの証明。
% source_v13, source_v14, source_v15 は変更しない。
% 外部入力：Fisher, The invariants of a genus one curve, Theorem 4.4(i), n=3。
% 追加 bibliography key: Fisher2008（別ファイルに収録）。

\section{Pl\"ucker錐とHessian作用の全次数簡約}
\label{sec:v16-plucker-hessian-structure}

本節では、第15版では未採用だった構造案を証明する。
通常接触像の階数や、Pfaffian代入後のKoszul複体の完全性は仮定しない。
全ての係数は $F$ に多項式的であり、接触作用の間に $F$ を動かさない。
以下の全次数の包含は、実際の高次Tate係数の消滅を意味するものではない。

\subsection{線束の規約とPl\"uckerイデアル}

$\dim V=5$、$\Delta=\det V$ とし、
\[
 A=\Lambda^2V,\qquad W=A\otimes\Delta^{-1},\qquad
 E=V^*\otimes\Delta^{-1},\qquad L=\Delta^{-3}
\]
と置く。変数は $F\in\operatorname{Sym}^3V$、$a\in W$ であり、
\[
 q(a)=\frac12a\wedge a\in E,\qquad P=F(q(a))\in L.
\]
自然な同一視
$\operatorname{Sym}^3V=\operatorname{Sym}^3E^*\otimes L$
を用いれば、$F$ は $E$ 上の $L$ 値三次式である。
既存の $\mathcal G_r,\mathcal D_{r+1},\mathcal U_s,\mathcal C_n$ は、
それぞれ $W,E,\mathbb C,L$ 値の普遍同変多項式であり、その二重次数は
\[
\begin{array}{c|cc|c}
 &\deg_F&\deg_a&\text{値域}\\\hline
 \mathcal G_r&r-1&r&W\\
 \mathcal D_{r+1}&r-1&r+1&E\\
 \mathcal U_s&s&s&\mathbb C\\
 \mathcal C_n&n-5&n&L
\end{array}
\]
である。$J_n=dP(\mathcal G_{n-5})$ と書く。
Euler場によって $P\mathcal U_{n-6}\subset J_n$ である。

体積形式を選んで $a_{ij}$ を座標とし、
\[
 \Sigma=\{a\in W:\operatorname{rank}a\le2\},\qquad
 \mathfrak I=(q_0,q_1,q_2,q_3,q_4)
\]
と置く。以下での $\mathfrak I$ は、必要に応じて
$\mathbb C[F,a]$ へ拡大したイデアルも表す。

\begin{lemma}[Pl\"ucker錐の被約なイデアル]
\label{lem:v16-plucker-prime-ideal}
$\mathfrak I$ は素イデアルであり、$\Sigma$ のイデアルに等しい。
従って、普遍多項式が全ての $(F,a)$、$a\in\Sigma$ で消えることと、
それが $\mathbb C[F,a]$ の $\mathfrak I$ に入ることは同値である。
\end{lemma}

\begin{proof}
次の五式を $g_1,\ldots,g_5$ とする。
\[
\begin{split}
 g_1&=a_{12}a_{34}-a_{13}a_{24}+a_{14}a_{23},\\
 g_2&=a_{02}a_{34}-a_{03}a_{24}+a_{04}a_{23},\\
 g_3&=a_{03}a_{14}-a_{04}a_{13}-a_{01}a_{34},\\
 g_4&=a_{02}a_{14}-a_{04}a_{12}-a_{01}a_{24},\\
 g_5&=a_{02}a_{13}-a_{03}a_{12}-a_{01}a_{23}.
\end{split}
\]
これらは順に $q_0,-q_1,-q_2,q_3,-q_4$ である。
変数の順序を
\[
 a_{02}>a_{03}>a_{04}>a_{12}>a_{13}>a_{14}>
 a_{23}>a_{24}>a_{34}>a_{01}
\]
とする辞書式順序では、この五式はGr\"obner基底となる。
実際、十組の $S$ 多項式を表示した順に除法すると全て余りが零になる。
その十組の除法の商全体を
\nolinkurl{v16_structure_certificate.json} に記録し、
\nolinkurl{v16_structure_certificate.py} は各等式を整数・有理数で再構成する。
初期単項式はいずれも $a_{01}$ を含まないので、
$\mathfrak I:a_{01}^{\infty}=\mathfrak I$ である。

$a_{01}$ を可逆にすると、最後の三式から
\[
 a_{ij}=\frac{a_{0i}a_{1j}-a_{0j}a_{1i}}{a_{01}}
 \quad(2\le i<j\le4)
\]
を得る。これを最初の二式へ代入すると恒等的に零となるため、局所化した商は
\[
 \mathbb C[a_{01}^{\pm1},a_{02},a_{03},a_{04},a_{12},a_{13},a_{14}]
\]
である。飽和性により元の商はこの整域へ単射するので、$\mathfrak I$ は素である。
この開集合の点は二つのベクトルの外積であり、その閉包は階数二以下の錐である。
また五つの4行4列Pfaffianが同時に零となる交代行列の階数は二以下である。
従って零点集合は正確に $\Sigma$ である。
$F$ の係数を加える操作は多項式拡大なので、最後の主張も従う。
\end{proof}

\begin{lemma}[普遍Pfaffian方向は錐上で消える]
\label{lem:v16-plucker-vanishing-D}
全ての $r\ge1$ と $w\in\mathcal D_{r+1}$ について
\[
 w(F,a)|_{a\in\Sigma}=0.
\]
さらに、$F,a$ にそれぞれ $r-1$ 次の、$\operatorname{End}(E)$ 値の
普遍同変多項式 $M(F,a)$ を選んで
\begin{equation}\label{eq:v16-equivariant-Mq}
 w=Mq
\end{equation}
と書ける。そのトレースは $\operatorname{tr}M\in\mathcal U_{r-1}$ である。
\end{lemma}

\begin{proof}
$V=L_2\oplus Q_3$、$L_2=\langle e_0,e_1\rangle$、
$Q_3=\langle e_2,e_3,e_4\rangle$ と分け、
$a_0=e_0\wedge e_1\otimes\Delta^{-1}$ を取る。
ここで右辺の $\Delta^{-1}$ の元は選んだ体積による基底を意味する。
一パラメータ部分群
\[
 \lambda(t)=\operatorname{diag}(t,t,1,1,1)
\]
は、外積の重み2と $\Delta^{-1}$ の重み $-2$ が相殺するため $a_0$ を固定する。
入力 $F\in\operatorname{Sym}^3V$ の $\lambda$ 重みは0、1、2、3である。
従って $F$ に多項式的な式には負の重みは現れない。
一方、出力 $E=V^*\otimes\Delta^{-1}$ の重みは、
$L_2^*$ 成分で $-3$、$Q_3^*$ 成分で $-2$ である。
同変性から $w(F,a_0)=0$ を得る。
非零の階数二の点は一つの $\mathrm{GL}(V)$ 軌道をなすので、$w$ は $\Sigma$ 全体で消える。
これは determinant twist を含めた証明であり、固定した滑らかな $F$ のみに関する議論ではない。

前補題から各成分は $\mathfrak I$ に入り、二重次数をそろえて $w=Mq$ と書ける。
この表示を与える有限次元の係数空間の全射は $\mathrm{GL}(V)$ 同変である。
標数零の完全可約性により、同変な $w$ の持ち上げ $M$ を同変に選べる。
つまり $M$ の値域は $E\otimes E^*=\operatorname{End}(E)$ であり、
変換則は共役作用である。従ってトレースはスカラー値で、その二重次数から
$\mathcal U_{r-1}$ に属する。$M$ の一意性は主張しない。

表現論的にも、錐の $a$ 次数 $r+1$ の成分は $S_{(r+1,r+1)}V$ であるのに対し、
trace線を外した出力の最高重みは $(r,r,r,r,r-1)$ である。
前者の第一行は後者より大きいので、多項式 $F$ 表現とのテンソル積から出力は現れない。
これは上の一パラメータ部分群による消失と整合する。
\end{proof}

\subsection{adjugate恒等式と同変な持ち上げ}

$p\in E$ を独立変数として
\[
 g(p)=dF_p,\qquad T(p)=\operatorname{Hess}_F(p),\qquad
 h(p)=\det T(p),\qquad B(p)=\operatorname{adj}T(p)
\]
と置き、$H=h(q(a))$ と書く。
ここで各テンソルの値域を省略してはならない。
\[
\begin{split}
 g&\in E^*\otimes L,\qquad
 T\in\operatorname{Sym}^2E^*\otimes L,\\
 h&\in(\det E^*)^2\otimes L^5=L,\\
 B&\in\operatorname{Sym}^2E\otimes(\det E^*)^2\otimes L^4
      =\operatorname{Sym}^2E.
\end{split}
\]
最後の二つの同一視には $\det E^*=\Delta^6$ を用いた。
特に $B$ は自然な写像 $E^*\to E$ であり、ここで選んだ体積に依存する
余分な determinant factor は残らない。

三次式のEuler恒等式 $Tp=2g$ と $BT=h\,\operatorname{id}$ から
\begin{equation}\label{eq:v16-adjugate-euler}
 B g=\frac h2p
\end{equation}
を得る。これを $p_k$ で微分すると、座標表示では
\begin{equation}\label{eq:v16-differentiated-adjugate}
 p_j\partial_k h
   =h\delta_{jk}+2\sum_i(\partial_k B_{ji})g_i
\end{equation}
となる。これは $\operatorname{End}(E)\otimes L$ 値のテンソル恒等式である。

\begin{proposition}[Hessianの一次作用の全次数簡約]
\label{prop:v16-hessian-derivation-reduction}
全ての $r\ge1$ について
\begin{equation}\label{eq:v16-all-hessian-action-mod-units}
 dH(\mathcal G_r)\subset J_{r+9}+H\mathcal U_{r-1}.
\end{equation}
より正確には、$w\in\mathcal D_{r+1}$ に対し
\[
 dh_{q(a)}(w)=H\operatorname{tr}M+dF_{q(a)}(v),\qquad
 v\in\mathcal D_{r+5},
\]
と書ける。最後の項は $J_{r+9}$ に入る。
\end{proposition}

\begin{proof}
前補題の $w=Mq$ を用いる。
式\eqref{eq:v16-differentiated-adjugate}に $M_{kj}$ を掛けて $k,j$ を縮約すれば
\[
 dh_q(Mq)=H\operatorname{tr}M+\sum_i g_i(q)v_i,
 \qquad
 v_i=2\sum_{k,j}M_{kj}(F,a)(\partial_kB_{ji})(F,q(a)).
\]
ここで微分 $\partial_k$ は $B$ にだけ作用し、$M$ を微分しない。
$dB\in E^*\otimes\operatorname{Sym}^2E$ と $M\in E\otimes E^*$ の縮約なので、
$v$ は $E$ 値の普遍同変多項式である。
$B$ は $F$ と $p$ にそれぞれ4次だから、$v$ の二重次数は
\[
 (\deg_Fv,\deg_av)=(r+3,r+5).
\]
従って $v\in\mathcal D_{r+5}$ である。
命題\ref{prop:v12-uniform-pfaffian-surjectivity}の全次数全射
\[
 dq:\mathcal G_{r+4}\twoheadrightarrow\mathcal D_{r+5}
\]
を用いて $v=dq(Z)$ と持ち上げれば $dF_q(v)=dP(Z)\in J_{r+9}$ となる。
$w=dq(Y)$、$Y\in\mathcal G_r$ とすれば所望の包含を得る。
\end{proof}

\begin{proposition}[錐上で消える単元のHessian倍]
\label{prop:v16-plucker-unit-annihilation}
$u\in\mathcal U_s$ が $\Sigma$ 上で消えるなら
\begin{equation}\label{eq:v16-ideal-unit-annihilation}
 Hu\in J_{s+10}.
\end{equation}
従って $H\mathcal U_s$ の通常接触商における像は、
$\mathcal U_s$ のPl\"ucker錐への制限だけに依存する。
\end{proposition}

\begin{proof}
$s=0,1$ の場合はイデアルの最小次数が二であるため $u=0$ である。
$s\ge2$ では、補題\ref{lem:v16-plucker-prime-ideal}と完全可約性から
\[
 u=\langle b,q\rangle,\qquad b(F,a)\in E^*,\qquad
 (\deg_Fb,\deg_ab)=(s,s-2)
\]
という普遍同変な表示を選べる。
式\eqref{eq:v16-adjugate-euler}と $B$ の対称性により
\[
 Hu=2\langle b,Bg\rangle=\langle g,2Bb\rangle.
\]
$2Bb$ は $E$ 値で、その二重次数は $(s+4,s+6)$ なので
$\mathcal D_{s+6}$ に属する。
全射 $dq:\mathcal G_{s+5}\to\mathcal D_{s+6}$ で持ち上げれば結論を得る。
\end{proof}

\begin{remark}[持ち上げの代数性と有理性]
ここでの二つの持ち上げは、有限次数の普遍係数空間の線形写像に対するものである。
$F$ の判別式、$P$ の値、$H$ の値、行列 $T$ の行列式による除算はない。
adjugate は最初から多項式である。
入力が有理係数なら、同変性も含めた線形条件は有理数体上の条件であり、
複素数体上での可解性から有理数体上の持ち上げを選べる。
従ってこれらは固定した $F$ ごとの存在主張ではなく、$F$ に多項式的な普遍恒等式を与える。
\end{remark}

\subsection{単元制限環と三元三次不変式}

$\mathcal U=\bigoplus_{s\ge0}\mathcal U_s$ とし、その斉次イデアルを
\[
 \mathcal N_s=\{u\in\mathcal U_s:u|_\Sigma=0\},\qquad
 \mathcal N=\bigoplus_{s\ge0}\mathcal N_s
\]
と定める。

\begin{theorem}[単元のPl\"ucker制限環]
\label{thm:v16-unit-restriction-ring}
階数二の点 $a_0=e_0\wedge e_1\otimes\Delta^{-1}$ での制限により、次数付き環の同型
\begin{equation}\label{eq:v16-unit-restriction-ring}
 \mathcal U/\mathcal N
 \simeq\mathbb C[\operatorname{Sym}^3Q_3]^{\mathrm{SL}(Q_3)}
 =\mathbb C[S_4,T_6]
 =\mathbb C[\overline U_4,\overline J_6]
\end{equation}
を得る。従って $\mathcal U_s/\mathcal N_s$ の基底は
\[
 \overline U_4^{\,a}\overline J_6^{\,b},\qquad
 a,b\ge0,\quad 4a+6b=s
\]
である。
\end{theorem}

\begin{proof}
補題\ref{lem:v16-plucker-vanishing-D}と同じ一パラメータ部分群は $a_0$ を固定する。
スカラー値の $u(F,a_0)$ の重みは零であり、$F$ の各成分の重みは非負なので、
この多項式は $F$ の $\operatorname{Sym}^3Q_3$ 成分だけに依存する。
$f\in\operatorname{Sym}^3Q_3$ に対し $\overline u(f)=u(f,a_0)$ と置く。
$h\in\mathrm{GL}(Q_3)$ を $V$ 上で恒等写像と直和して作用させると
\[
 h a_0=(\det h)^{-1}a_0,
 \qquad \overline u(hf)=(\det h)^s\overline u(f).
\]
特に $\overline u$ は $\mathrm{SL}(Q_3)$ 不変式であり、その通常の $f$ 次数は $s$ である。
逆に、三次式の係数に $s$ 次の $\mathrm{SL}_3$ 不変式の $\mathrm{GL}_3$ 重みは
$\det^s$ なので、次数とdeterminant twistは一致する。
非零の階数二の点が一軌道であることから、制限写像の核は正確に $\mathcal N$ である。

三元三次式の全係数空間の $\mathrm{SL}_3$ 不変式環は、次数4、6の二つの生成元の
多項式環である。ここではFisherの一次資料\cite[\S3.3,\S4,
Lemma~4.3, Theorem~4.4(i)]{Fisher2008}を、$n=3$、標数零で適用する。
この定理は判別式を可逆化した環ではなく、特異三次式も含む全係数空間の環についてのものである。

命題\ref{prop:v14-rank-two-unit-restriction}によれば、
$f_t=x_2^3+x_3^3+x_4^3+6t x_2x_3x_4$ 上で
\[
 \overline U_4(f_t)=t^4-t,\qquad
 \overline J_6(f_t)=6-120t^3-48t^6.
\]
両者はそれぞれ次数4、6の非零不変式である。
その次数の不変式空間はそれぞれ一次元なので、$S_4,T_6$ の非零定数倍を与える。
従って制限写像の像は不変式環の生成元を両方含み、全射である。
これで環の同型が従う。
既存の制限計算では、$(s,t)\mapsto(s^4\overline U_4(f_t),s^6\overline J_6(f_t))$
の $(1,0)$ におけるJacobian行列式が36であり、代数的独立性も直接検証されている。
\end{proof}

\begin{theorem}[全ての正次数単元倍の消去]
\label{thm:v16-positive-unit-hessian-annihilation}
全ての整数 $s>0$ について
\begin{equation}\label{eq:v16-all-positive-unit-hessian}
 H\mathcal U_s\subset J_{s+10}.
\end{equation}
\end{theorem}

\begin{proof}
低次の入力は、既存の次の二つの所属だけである。
\[
 HU_4\in J_{14},\qquad HJ_6\in J_{16}.
\]
第一の所属は、式\eqref{eq:v11-rho-fourteen-values}と
命題\ref{prop:v11-contact-rank-fourteenth}から
小節\ref{subsec:v11-fourteenth-complement}で証明されている。
第二の所属は、$J_6\in\mathcal U_6$ と
命題\ref{prop:v13-hessian-unit-membership}から従う。

任意の $u\in\mathcal U_s$ に対し、前定理によって
\[
 u-\sum_{4a+6b=s}c_{ab}U_4^aJ_6^b\in\mathcal N_s
\]
と書ける。左辺の $H$ 倍は命題\ref{prop:v16-plucker-unit-annihilation}により
$J_{s+10}$ に入る。
$s>0$ なので各単項式では $a>0$ または $b>0$ である。
前者なら $HU_4$、後者なら $HJ_6$ の既知の持ち上げに残る単元を掛ければよい。
実際、$v\in\mathcal U_t$、$Y\in\mathcal G_r$ について
\[
 vY\in\mathcal G_{r+t},\qquad v\,dP(Y)=dP(vY)
\]
なので、通常接触像は単元共変量環の加群である。
従って全ての項が $J_{s+10}$ に入り、結論を得る。
\end{proof}

\begin{corollary}[全次数のHessian一次作用と22次への帰結]
\label{cor:v16-all-hessian-first-actions}
全ての $r\ge2$ について
\begin{equation}\label{eq:v16-all-hessian-first-actions}
 dH(\mathcal G_r)\subset J_{r+9}.
\end{equation}
特に
\[
 H\mathcal U_{12}\subset J_{22},\qquad
 dH(\mathcal G_{13})\subset J_{22},\qquad
 dH(\mathcal G_8)\subset J_{17}.
\]
$r=1$ はこの消去の主張に含めない。
実際、$\mathcal G_1=\mathbb C a$ であり $dH(a)=10H$ である。
\end{corollary}

\begin{proof}
命題\ref{prop:v16-hessian-derivation-reduction}と
定理\ref{thm:v16-positive-unit-hessian-annihilation}を、$s=r-1>0$ として適用する。
22次については制限環の12次部分が $\overline U_4^3,\overline J_6^2$ で張られるため、
全次数の記述を使わずこの二方向と錐イデアル部分に分けても同じ結論を得る。
また $r=8$ の場合は、既存の $\mathcal U_7=0$ と
式\eqref{eq:v16-all-hessian-action-mod-units}から直接にも従う。
\end{proof}

\subsection{下位jetが二項正規形である場合の有効像}

\begin{theorem}[二項正規形の後の有効像は全次数で通常像に等しい]
\label{thm:v16-all-effective-images}
$N>10$、$c\in\mathbb C$ とする。
$N-1$ 次までのjetが $P+cH$ である方程式について、
その下位jetを保ち、線形部分が恒等写像、単元の定数項が1である普遍接触変換の
$N$ 次有効像は
\begin{equation}\label{eq:v16-all-effective-images}
 I_N=J_N
\end{equation}
である。偶数 $N$ に対しては、奇座標変更と偶単元だけに制限しても同じ像となる。
\end{theorem}

\begin{proof}
接触対数補題\ref{lem:v14-contact-logarithm}により、下位jetを保つ変換の
最初の差は、対数作用素 $\sum_{k>0}Q_k$ の一次作用で計算できる。
ここで $Q_k=Y_{k+1}+b_k\in\mathcal G_{k+1}\oplus\mathcal U_k$ である。
$N$ 次の差は
\[
 Q_{N-6}P+cQ_{N-10}H
\]
となる。第一項はEuler場も含めて $J_N$ に入る。
第二項は $N-10>0$ なので、前定理と前系により
\[
 dH(\mathcal G_{N-9})+H\mathcal U_{N-10}\subset J_N
\]
に属する。従って $I_N\subset J_N$ である。
逆に $Y\in\mathcal G_{N-5}$ による座標変更 $a\mapsto a+Y$ は
$N-1$ 次までを保ち、$N$ 次に $dP(Y)$ を与えるため、$J_N\subset I_N$ である。
$N$ が偶数なら $N-5$ は奇数なので、この逆包含は奇座標変更だけで実現する。
\end{proof}

\begin{remark}[適用範囲]
本定理の仮定は、対象の下位jetが既に $P+cH$ であることである。
消去されていない別の低次項がある場合には、その項への接触作用を追加する必要がある。
また $I_N=J_N$ は $\mathcal C_N/J_N$ の次元も、実際の $R_N$ の類も決定しない。
とりわけ、この構造結果だけから実際の22次係数の零性を結論しない。
\end{remark}

\subsection{有限検算と証明への依存}

\nolinkurl{v16_structure_certificate.py} はPython標準ライブラリだけを使う。
五つのPl\"ucker関係の十組の $S$ 多項式の正確な除法、
二重次数と中心重み、既存の階数二単元制限のJacobian36を検算する。
さらに三つの明示的な整数係数 $F$ について、$p$ を記号変数のまま保持し、
式\eqref{eq:v16-adjugate-euler}の五成分と
式\eqref{eq:v16-differentiated-adjugate}の25成分を全て正確に検算する。
最後の三例は添字・係数の監査であり、有限個の代入から普遍恒等式を推論するものではない。
普遍恒等式の証明は $Tp=2g$ とadjugateの恒等式を微分する上の議論である。
三元三次不変式環については一次資料の定理を用い、
既存の接触階数やTate計算を今回独立に再認証したとは主張しない。

% END theory/v16_plucker_hessian_structure_ja.tex


% BEGIN theory/v16_normal_slice_injection_ja.tex
% 第16版：rank-two normal slice による通常接触商への全次数単射。
% この節の単射は Fisher の不変式環生成定理には依存しない。
% v14 の rank-two U4,J6 独立性、v15 の既知328次元部分空間を継承する。

\section{階数二の横断三次項と通常接触商の独立な方向}
\label{sec:v16-normal-slice-injection}

前節のPl\"ucker錐上の消失を使うと、通常接触像を消す新しい検出方法が得られる。
それは階数二の点の横断三次項を、三元三次式のJacobian空間を法として取る方法である。
この方法は、有限体の零やグラフ探索の停滞による上界を用いない。
特に第15版で残っていた20次通常階数の二択を解消する。

\subsection{横断切片と三次Jacobian空間}

体積形式を一つ固定し、$p=(p_2,p_3,p_4)$、$\varepsilon$ を独立変数として
\begin{equation}\label{eq:v16-normal-slice}
 \begin{split}
 a_\varepsilon(p)&=e_0\wedge e_1+
 \varepsilon\bigl(p_4e_2\wedge e_3-p_3e_2\wedge e_4+p_2e_3\wedge e_4\bigr),\\
 F_f&=x_0^3+x_1^3+f(x_2,x_3,x_4)
 \end{split}
\end{equation}
と置く。$a$ の $\Delta^{-1}$ 因子はこの体積形式で表示している。
Pfaffianは正確に
\begin{equation}\label{eq:v16-normal-pfaffian}
 q(a_\varepsilon(p))=(0,0,\varepsilon p_2,\varepsilon p_3,\varepsilon p_4)
\end{equation}
となる。$f$ は三元三次式の全係数空間を動く。
\[
 \operatorname{Jac}(f)_3
 =\operatorname{Span}_{\mathbb C}
 \{p_j\partial_i f: i,j\in\{2,3,4\}\}
 \subset\mathbb C[p_2,p_3,p_4]_3
\]
と書く。

\begin{lemma}[通常接触像の横断三次項]
\label{lem:v16-contact-normal-cubic-jet}
$N\ge6$、$Y\in\mathcal G_{N-5}$ とする。
この切片における $dP(Y)$ の $\varepsilon^3$ 係数は
$\operatorname{Jac}(f)_3$ に入る。
\end{lemma}

\begin{proof}
$w=dq(Y)\in\mathcal D_{N-4}$ と置く。
補題\ref{lem:v16-plucker-vanishing-D}により、
$a_0=e_0\wedge e_1$ で $w(F_f,a_0)=0$ である。
従って、$p$ に線形なベクトル $w_1(f,p)$ を用いて
\[
 w(F_f,a_\varepsilon(p))=\varepsilon w_1(f,p)+O(\varepsilon^2)
\]
と展開できる。
一方、$F_f$ の最初の二つの偏微分は $3q_0^2,3q_1^2$ なので、
切片上では正確に零である。残る三成分は
\[
 \partial_iF_f(q(a_\varepsilon))=\varepsilon^2\partial_i f(p)
 \qquad(i=2,3,4)
\]
となる。よって
\[
 [\varepsilon^3]\,dP(Y)(F_f,a_\varepsilon(p))
   =\sum_{i=2}^4\partial_i f(p)\,(w_1)_i(f,p)
   \in\operatorname{Jac}(f)_3.
\]
場 $Y$ の接方向の成分を零と仮定する必要はない。
最初の二成分の寄与が消えるのは、分離した $F_f$ と切片の明示式による。
\end{proof}

\begin{lemma}[$\Theta_8$ の切片上の正確な値]
\label{lem:v16-theta-normal-hessian}
$B_f(p)=\tfrac16\operatorname{Hess}_f(p)$ と置くと
\begin{equation}\label{eq:v16-exact-theta-slice}
 \Theta_8(F_f,a_\varepsilon(p))
   =\varepsilon^3\det B_f(p)
   =\frac{\varepsilon^3}{216}\det\operatorname{Hess}_f(p).
\end{equation}
これは最初の係数だけでなく、切片上での正確な多項式恒等式である。
また $u\in\mathcal U_s$ について
\[
 [\varepsilon^0]\,u(F_f,a_\varepsilon(p))=\overline u(f)
\]
であり、右辺は前節の階数二制限である。
\end{lemma}

\begin{proof}
既存の式\eqref{eq:boundary-theta-eight}を用いる。
切片上の $\operatorname{Hess}_{F_f}(q)$ は、最初の二行・二列が零であり、
残る三行三列のブロックは $\varepsilon\operatorname{Hess}_f(p)$ である。
従って3行3列minorの和に残るのは
$I=J=\{2,3,4\}$ の一項だけである。
その係数は $c_{234}^2=a_{01}^2=1$ なので、表示した恒等式となる。
単元の定数項は $u(F_f,a_0)$ であり、
前節と同じ一パラメータ部分群の重みにより、これは $F_f$ の三元部分 $f$ だけに依存する。
\end{proof}

三元三次式 $f$ について、十個の三次式
\begin{equation}\label{eq:v16-ten-ternary-cubics}
 \{p_j\partial_i f: i,j=2,3,4\},\qquad \det B_f
\end{equation}
の係数行列を $M(f)$ とする。
Fermat式 $f=p_2^3+p_3^3+p_4^3$ では、行を
九個の単項式 $p_i^2p_j$（$i,j$ の順）と $p_2p_3p_4$ の順に並べると
\begin{equation}\label{eq:v16-fermat-normal-minor}
 M(f)=\operatorname{diag}(3,3,3,3,3,3,3,3,3,1),
 \qquad\det M(f)=3^9=19683\ne0.
\end{equation}
従って三元三次式の非空なZariski開集合で
\[
 \mathbb C[p_2,p_3,p_4]_3
   =\operatorname{Jac}(f)_3\oplus\mathbb C\det B_f
\]
となる。この非零性には、Jacobian環の完全交叉性やsocleについての外部定理は不要である。

独立検算 \nolinkurl{v16_normal_slice_certificate.py} は、
上の十行十列行列とPfaffianの切片式を正確な整数・有理数で再構成する。
さらに $f_t=p_2^3+p_3^3+p_4^3+6t p_2p_3p_4$ について
\[
 \det M(f_t)=3^9(1+8t^3)^3
\]
および
\[
 \det B_{f_t}
 =(1+8t^3)p_2p_3p_4-\frac{t^2}{3}\sum_{i=2}^4p_i\partial_i f_t
\]
を検証する。前者は次数上界12と13個の整数値、後者は係数ごとの多項式恒等式による。
これらの追加のHesse計算を用いなくても、Fermatの一つの非零minorと開性で以下の証明は閉じる。

\subsection{通常接触商への全次数単射}

\begin{theorem}[$\Theta_8$ と二つの単元が与える自由部分加群]
\label{thm:v16-hessian-socle-injection}
$\deg U_4=4$、$\deg J_6=6$、$\deg\Theta_8=8$ とすると、
\begin{equation}\label{eq:v16-theta-module-injection}
 \mathbb C[U_4,J_6]\,[\Theta_8]
 \hookrightarrow
 \bigoplus_N\mathcal C_N/J_N
\end{equation}
は次数付き $\mathbb C[U_4,J_6]$ 加群の単射である。
すなわち、各 $s\ge0$ について
\[
 \{[\Theta_8U_4^aJ_6^b]:4a+6b=s\}
 \quad\text{は }\mathcal C_{s+8}/J_{s+8}\text{ で独立である。}
\]
\end{theorem}

\begin{proof}
重み $s$ の多項式 $A(U_4,J_6)$ が
$\Theta_8 A(U_4,J_6)\in J_{s+8}$ を満たすと仮定する。
前の二補題により、全ての三元三次式 $f$ に対して
\[
 A(\overline U_4(f),\overline J_6(f))\det B_f
      \in\operatorname{Jac}(f)_3
\]
となる。式\eqref{eq:v16-fermat-normal-minor}が定める非空な開集合では
$[\det B_f]\ne0$ なので
\[
 A(\overline U_4(f),\overline J_6(f))=0
\]
である。これは三元三次式の係数の多項式だから、全係数空間で恒等的に零となる。
命題\ref{prop:v14-rank-two-unit-restriction}の代数的独立性から $A=0$ が従う。
また通常接触像は単元共変量による積で保たれるため、この単射は加群の単射である。

この証明は、$\overline U_4,\overline J_6$ の既存の独立性だけを使う。
三元三次不変式環の全生成定理\cite{Fisher2008}は本定理の依存には入らない。
\end{proof}

\begin{corollary}[20次通常接触階数の確定]
\label{cor:v16-exact-twentieth-contact-rank}
第15版の記号を保つと
\begin{equation}\label{eq:v16-exact-twentieth-rank}
 J_{20}=S_{20},\qquad\dim J_{20}=328,
 \qquad\dim\mathcal C_{20}/J_{20}=2.
\end{equation}
通常商の基底は
\[
 [\Theta_8U_4^3],\qquad[\Theta_8J_6^2]
\]
であり、
\[
 \dim\ker(dP:\mathcal G_{15}\to\mathcal C_{20})=916,
 \qquad
 \dim\ker(dF:\mathcal D_{16}\to\mathcal C_{20})=360.
\]
\end{corollary}

\begin{proof}
命題\ref{prop:v15-twentieth-contact-subspace}の既存の整数証明書により
$S_{20}\subset J_{20}$、$\dim S_{20}=328$、$\dim\mathcal C_{20}=330$ である。
前定理を $s=12$ に適用すると、二つの候補は $\mathcal C_{20}/J_{20}$ でも独立なので
$\dim J_{20}\le328$ となる。従って $J_{20}=S_{20}$ である。
既存の $\dim\mathcal G_{15}=1244$ から最初の核の次元を得る。
$dq:\mathcal G_{15}\twoheadrightarrow\mathcal D_{16}$ が全射であり
$dP=dF\circ dq$ なので、$dF$ の階数も328である。
$\dim\mathcal D_{16}=688$ から最後の次元を得る。
\end{proof}

\begin{remark}[第15版の検出子とKoszul核の読み替え]
第15版の二つの三点汎関数は $S_{20}$ を正確に消していた。
今回の $S_{20}=J_{20}$ により、それらは通常接触商上の双対基底でもあることが新たに確定した。
これは第15版の時点で仮定していた性質ではない。
第15版の原データ、三点値、実際の20次消去の証明を変更する必要はない。

また既存の359個の独立なKoszul核ベクトルは、確定した360次元核の真部分空間である。
残る一方向が別のKoszul像で生成されるのか、Koszul像の外にあるのかは、
この論証では決まらない。従ってKoszul完全性や、全Koszul像の階数は引き続き主張しない。
\end{remark}

\begin{corollary}[通常接触商の全次数Hilbert級数の下界]
\label{cor:v16-contact-hilbert-lower-bound}
$\mathbb C[U_4,J_6]$ 加群として
\[
 \mathbb C[U_4,J_6]\,[\Theta_8]\oplus\mathbb C\cdot[H]
 \hookrightarrow\bigoplus_N\mathcal C_N/J_N
\]
であり、第二項では正次数イデアル $(U_4,J_6)$ が零作用する。
従ってHilbert級数には係数ごとの下界
\begin{equation}\label{eq:v16-contact-hilbert-lower-bound}
 \sum_N\dim(\mathcal C_N/J_N)t^N
 \ \ge_{\mathrm{coeff}}\
 \frac{t^8}{(1-t^4)(1-t^6)}+t^{10}
\end{equation}
が成り立つ。
\end{corollary}

\begin{proof}
$[H]\ne0$ は定理\ref{thm:boundary-eleven-jet}の十次通常商の計算による。
定理\ref{thm:v16-positive-unit-hessian-annihilation}から、正次数単元は $[H]$ を消す。
$\Theta_8$ 部分加群の次数は $8+4a+6b$ であり、10次は含まないので直和となる。
残りは前定理による。
\end{proof}

\begin{remark}[下界と全商の同定を区別する]
式\eqref{eq:v16-contact-hilbert-lower-bound}を全商のHilbert級数との等式とは主張しない。
横断三次項の検出で見えない追加の接触類が、高次に存在する可能性は残る。
例えば22次にこの方法が与えるのは $[\Theta_8U_4^2J_6]$ の非零性と一次元の下界であり、
22次通常商の全次元も実際のTate係数の類も、これだけでは確定しない。
\end{remark}

% END theory/v16_normal_slice_injection_ja.tex


% BEGIN characters/v16_character_notes_ja.tex
\section{二十二次の普遍共変量とKoszul各項}
\label{sec:v16-characters}

二十二次では、候補となる単元の積だけから通常接触商の次元を決めない。
本節では普遍空間とKoszul源の次元を先に確定し、写像の階数とは区別する。
計算は整数の重み指標、Weyl交代和、Pieri則による。

\subsection{次数とdeterminant twist}

$\dim V=5$、$\Delta=\det V$、$A=\Lambda^2V$、
$W=A\otimes\Delta^{-1}$ とする。
$F\in\operatorname{Sym}^3V$、$a\in W$、
$q(a)=\tfrac12a\wedge a\in V^*\otimes\Delta^{-1}$、
$P=F(q(a))\in\Delta^{-3}$ という既存の規約を保つ。
中心元 $t\operatorname{id}_V$ の重みから、$a$ 次数 $n$ の境界項は
$F$ 次数 $n-5$、$a$ 次数 $r$ の座標補正は $F$ 次数 $r-1$、
$a$ 次数 $s$ の単元補正は $F$ 次数 $s$ を持つ。
従って、trace線を外した普遍空間は
\begin{align*}
 C_n&=\operatorname{Hom}_{\mathrm{GL}(V)}
 \bigl(\operatorname{Sym}^{n-5}(\operatorname{Sym}^3V)
       \otimes\operatorname{Sym}^{n}A,\Delta^{n-3}\bigr),\\
 G_r&=\operatorname{Hom}_{\mathrm{GL}(V)}
 \bigl(\operatorname{Sym}^{r-1}(\operatorname{Sym}^3V)
       \otimes\operatorname{Sym}^{r}A,A\otimes\Delta^{r-1}\bigr),\\
 D_{r+1}&=\operatorname{Hom}_{\mathrm{GL}(V)}
 \bigl(\operatorname{Sym}^{r-1}(\operatorname{Sym}^3V)
       \otimes\operatorname{Sym}^{r+1}A,\Lambda^4V\otimes\Delta^{r-1}\bigr),\\
 U_s&=\operatorname{Hom}_{\mathrm{GL}(V)}
 \bigl(\operatorname{Sym}^{s}(\operatorname{Sym}^3V)
       \otimes\operatorname{Sym}^{s}A,\Delta^s\bigr).
\end{align*}
二十二次のKoszul第 $k$ 項は
\[
 K_{k,22}=\operatorname{Hom}_{\mathrm{GL}(V)}
 \left(\operatorname{Sym}^{17-k}(\operatorname{Sym}^3V)
 \otimes\operatorname{Sym}^{22-4k}A,
 \Lambda^kV^*\otimes\Delta^{19-2k}\right),\qquad 1\leq k\leq5.
\]
$K_{1,22}=D_{18}$、$K_{2,22}=B_{22}$ と書く。
第5項にも $a$ 次数2が残るので、前の次数の終端をそのまま切り捨ててはならない。

\begin{proposition}[二十二次の厳密次元]
\label{prop:v16-character-dimensions}
各普遍空間の次数・標的分割・次元は次の通りである。
\[
\begin{array}{c|r|r|c|r}
\text{空間}&\deg_F&\deg_a&\text{標的分割}&\dim\\\hline
C_{22}&17&22&(19,19,19,19,19)&916\\
G_{17}&16&17&(17,17,16,16,16)&3929\\
U_{16}&16&16&(16,16,16,16,16)&486\\
D_{18}&16&18&(17,17,17,17,16)&2140\\
B_{22}&15&14&(15,15,15,14,14)&1706\\
K_{3,22}&14&10&(13,13,12,12,12)&534\\
K_{4,22}&13&6&(11,10,10,10,10)&57\\
K_{5,22}&12&2&(8,8,8,8,8)&2
\end{array}
\]
これらの値は接触写像の階数やKoszul完全性を仮定しない。
\end{proposition}

\begin{proof}
$\operatorname{Sym}^m(\operatorname{Sym}^3V)$ の指標を、35個の三次単項式重み
$\alpha$ に対するEuler積
\[
 \prod_{|\alpha|=3}(1-tx^\alpha)^{-1}
\]
の $t^m$ 係数として整数演算で計算する。
$\rho=(4,3,2,1,0)$ とすると、Schur多重度は
\[
 [S_\lambda V:\operatorname{Sym}^m(\operatorname{Sym}^3V)]
 =\sum_{\sigma\in S_5}\operatorname{sgn}(\sigma)
 [x^{\lambda+\rho-\sigma\rho}]
 \operatorname{ch}\operatorname{Sym}^m(\operatorname{Sym}^3V)
\]
である。外積二乗側の指標も10個の重みから独立に作り、標的分割に対する
テンソル積の交代和を直接評価した。

また
\[
 \operatorname{Sym}^{d}(\Lambda^2V)
 =\bigoplus_{j=0}^{\lfloor d/2\rfloor}S_{(d-j,d-j,j,j,0)}V
\]
の各多重度1とWeyl次元和を照合した。
$C_{22}$ で必要な $F$ 側の分割は
$(19,16-j,16-j,j,j)$、$0\leq j\leq8$ であり、各多重度は
\[
 (3,4,36,78,226,234,254,61,20).
\]
和は916となる。
$U_{16}$ の分割は $(16,16-j,16-j,j,j)$、$0\leq j\leq8$ で、各多重度は
\[
 (2,1,14,33,112,126,148,37,13).
\]
和は486となる。

$G_{17}$ と $D_{18}$ はそれぞれ外積第3冪・第1冪を用いる垂直帯Pieri計算、
$K_{k,22}$ は外積第 $k$ 冪を用いる垂直帯Pieri計算でも計算し、
テンソル積指標からの値と照合した。
Koszul源で非零となるPieri成分数は $k=2,3,4,5$ の順に25、20、8、1である。
第5項の非零成分は $F$ 側分割 $(8,8,8,6,6)$ で、多重度2である。

全ての $F$ 次数 $1\leq m\leq17$ と $a$ 次数 $1\leq d\leq22$ について、
Euler積とは異なるNewton--Adams再帰
\[
 m h_m=\sum_{k=1}^m p_kh_{m-k}
\]
により全重み係数を再計算し、辞書全体が一致することを確認した。
必要な $F$ 次数12から17までの全Schur分解についても、
多重度とWeyl次元の積の和が $\binom{m+34}{m}$ に一致する。
これらは全て有限の整数恒等式であり、浮動小数点階数を用いない。
\end{proof}

\subsection{Euler交代和と制限計算の範囲}

二十二次のKoszul各項の交代和は
\[
 916-2140+1706-534+57-2=3.
\]
従って $H_k$ を当該有限次数成分のKoszulホモロジーとすると
\[
 \dim H_0-\dim H_1+\dim H_2-\dim H_3+\dim H_4-\dim H_5=3.
\]
この等式だけから $\dim(C_{22}/J_{22})=3$、あるいは1と結論しない。
全次数Pfaffian微分全射性を用いる場合には
\[
 \dim\ker(\delta q:G_{17}\longrightarrow D_{18})=3929-2140=1789
\]
を得るが、$dF:D_{18}\to C_{22}$ の核の次元は別に求める必要がある。

\begin{proposition}[Pl\"ucker錐への制限の標的次元]
\label{prop:v16-plucker-dimensions}
分解可能な $a$ の錐への制限では、$\operatorname{Sym}^dA$ をCartan成分
$S_{(d,d)}V$ に置き換える。
この置換後のHom空間の次元は次の通りである。
\[
\begin{array}{c|rrrrrrrr}
 &C_{22}&G_{17}&U_{16}&D_{18}&B_{22}&K_{3,22}&K_{4,22}&K_{5,22}\\\hline
 \dim&0&2&2&0&0&12&17&2
\end{array}
\]
また $U_{14}$ の制限標的は1次元である。
\end{proposition}

\begin{proof}
前のLittlewood分解の $j=0$ のみを残して同じPieri計算を行う。
$U_{16}$ では $F$ 側分割 $(16,16,16,0,0)$ の多重度2、
$U_{14}$ では $(14,14,14,0,0)$ の多重度1を得る。
$D_{18}$ では第一行18が比較する矩形の幅17を超えるため全成分が消える。
各成分と交代和は保存証明書の \texttt{Plucker} 欄に収録した。
\end{proof}

ここで計算したのは制限先の普遍Hom空間である。
特定の候補がその空間を尽くすことや、通常接触商全体を尽くすことを
この表から自動的に結論してはならない。
特に $U_{14}$ の錐上成分が1次元であっても、それだけで
$C_{22}/J_{22}$ が $[\Theta_8U_4^2J_6]$ の1次元空間になるとはいえない。

Fermat固定群による従来の成分別制限計算からも、次の上界を得た。
\[
\begin{array}{c|r|r|r}
 &\dim&\text{対角値の階数上界}&\text{対角値と一階法jetの階数上界}\\\hline
C_{22}&916&284&803\\
G_{17}&3929&972&3244\\
U_{16}&486&106&371\\
D_{18}&2140&576&1834
\end{array}
\]
これは固定群指標から得る上界であり、実際の評価写像の階数の等号を主張しない。
対角値や一階jetだけを調べて普遍空間の全てを見たとみなすことはできない。

\subsection{再現用証明書}

\texttt{characters/v16\_character\_certificate.py} が全重みのEuler積・Newton再帰、
全Schur次元和、テンソル積の交代和とPieri照合を新規に実行する。
\texttt{v16\_koszul\_pieri\_check.py} は独自の垂直帯列挙を使い、
保存されたSchur分解から全Koszul源を再計算する。
\texttt{v16\_character\_audit.py} は保存証明書の次数、twist、交代和、
Weyl次元、Pieri、制限上界を短時間で監査する。
保存証明書の監査と全重みの新規再計算は区別する。

% END characters/v16_character_notes_ja.tex


% BEGIN rank/rank22_section_ja.tex
\section{二十二次の通常接触商とKoszul非完全性}
\label{sec:v16-contact-rank}

本節では、指標による源の次元とは独立に通常接触像の階数を決める。
下界には普遍場の明示的な評価小行列を、上界には
定理\ref{thm:v16-hessian-socle-injection}の横断三次jetによる非所属を使う。
KoszulのEuler交代和から写像の階数を推測する必要はない。

\begin{proposition}[二十二次の通常接触商]
\label{prop:v16-contact-rank-twenty-second}
通常接触像を $J_{22}=dP(G_{17})=dF(D_{18})$ とすると
\[
 \dim J_{22}=915,\qquad
 C_{22}/J_{22}=\mathbf C\cdot[\Theta_8U_4^2J_6].
\]
従って
\[
 \dim\ker(dP:G_{17}\to C_{22})=3014,\qquad
 \dim\ker(dF:D_{18}\to C_{22})=1225.
\]
さらに429個の明示的な普遍場 $\Gamma_{17,\alpha}$ により
\[
 J_{22}=P U_{16}\oplus
 \operatorname{span}\{dP(\Gamma_{17,\alpha}):1\leq\alpha\leq429\}
\]
と書ける。
\end{proposition}

\begin{proof}
17対17の三価二部グラフを考え、左側の最後の頂点を開口とする。
残り16個の左頂点へ対称三階テンソル $F$、17個の右頂点へ
$\tau_{ijk}=\epsilon_{ijklm}a_{lm}$ を置いて縮約し、自由三脚 $R_{ijk}$ を得る。
\[
 \Gamma^{lm}=\frac16\sum_{i,j,k}\epsilon_{lmijk}R_{ijk}
\]
は $F$ 次数16、$a$ 次数17の普遍場で、$G_{17}$ に属する。
開口へ $T_{ijk}=\epsilon_{ijklm}(\partial P/\partial a_{lm})$ を戻した閉縮約を
6で割ると $dP(\Gamma)$ になる。

採用した429個のグラフのうち286個は連結グラフの探索から、143個は
既存の低次場と単元の積から得た。
後者に用いた候補は、非単元 $J_{18}$ の60方向に $U_4$ を掛けた60列、
$J_{16}$ の21方向と $U_6$ の2方向から作る42列、
$J_{14}$ の7方向と $U_8$ の6方向から作る42列、
$J_{12}$ の2方向と $U_{10}$ の18方向から作る36列の計180列である。
そのうち143列を採用した。
単元 $u$ と場 $\Gamma$ に対して $dP(u\Gamma)=u\,dP(\Gamma)$ なので、
全て二十二次の通常接触像に属する。
積は単元グラフと場のグラフの非連結和でも表せるため、
最終証明書には各列を単独で再構成できる17対17のグラフを記録した。
四頂点単元グラフは未除算の値 $24U_4$ を使うが、非零定数倍は張る空間を変えない。

$q=(1,1,1,1,1)$、$P=0$ を有理数上で満たす450個の整数入力で評価する。
$F_{444}$ を調整して $P=0$ を課しており、有限体上だけの条件ではない。
429列の評価行列には、法1009で行列式527となる429行の小行列がある。
各グラフは整係数縮約で、分母6は1009で可逆なので、これは特性0での独立性を証明する。
選択した小行列の全 $429^2=184041$ 成分を、別の辺番号付けと整数縮約で再計算し、
二種類のGaussian消去でも非零行列式を照合した。

$P$ は非零多項式であり、普遍係数環は整域なので
$\dim(PU_{16})=\dim U_{16}=486$ である。
Eulerの等式 $dP(au/6)=Pu$ により $PU_{16}\subset J_{22}$。
前の429方向は $P=0$ 制限で独立なので $PU_{16}$ と交わらず、
\[
 \dim J_{22}\geq486+429=915.
\]

一方、定理\ref{thm:v16-hessian-socle-injection}により
$\Theta_8U_4^2J_6\notin J_{22}$。
$\dim C_{22}=916$ なので $\dim J_{22}\leq915$ となり、両界が一致する。
全次数Pfaffian微分全射性により $dP(G_{17})=dF(D_{18})$ であるから、
核の次元は $3929-915=3014$ および $2140-915=1225$ となる。
\end{proof}

\begin{proposition}[補元の独立な有限証明書]
\label{prop:v16-twenty-second-complement}
前の429列へ $W_{22}=\Theta_8U_4^2J_6$ の評価列を加えた行列には、
法1009で行列式591となる430行の小行列がある。
従って、明示的な直和分解
\[
 C_{22}=P U_{16}\oplus
 \operatorname{span}\{dP(\Gamma_{17,\alpha}):1\leq\alpha\leq429\}
 \oplus\mathbf C W_{22}
\]
が得られる。
\end{proposition}

\begin{proof}
$U_4$ は四頂点グラフを24で割った既存の正規化、$J_6$ は六頂点グラフの
既存の正規化を使う。450入力全てで $\Theta_8,U_4,J_6$ と $W_{22}$ を
有理数上で新規に評価し、選択小行列の行列式591を得る。
$486+430=916$ という独立な指標次元との一致から直和が従う。
この有限証明書単独は既知915次元部分空間への補元を示すものだが、
前命題の非所属による上界と合わせることで通常商そのものの基底となる。
\end{proof}

\begin{corollary}[二十二次の正次数Koszulホモロジー]
\label{cor:v16-positive-koszul-homology}
同一の二重次数とdeterminant twistを持つ二十二次Koszul成分のホモロジーを
$H_k$、その次元を $h_k$ とすると
\[
 h_2+h_4=2+h_1+h_3+h_5\geq2.
\]
従ってこの成分は正のホモロジー次数で完全ではなく、$H_2$ または $H_4$ は非零である。
\end{corollary}

\begin{proof}
命題\ref{prop:v16-character-dimensions}の全Koszul源次元からEuler交代和は3。
前命題により $h_0=\dim(C_{22}/J_{22})=1$ なので
\[
 3=1-h_1+h_2-h_3+h_4-h_5
\]
を整理すればよい。
\end{proof}

さらに定理\ref{thm:v16-koszul-high-homology-vanishing}を用いると、
$H_i=0$ ($i>2$) により
\[
 h_2=h_1+2\geq2
\]
へ強化される。これは次節で示す普遍勾配の三成分の正則列性に依存する帰結であり、
Euler交代和のみの帰結と区別する。
対応する第5・第4・第3Koszul微分の階数2、55、479は
系\ref{cor:v16-twentysecond-koszul-defect}にまとめる。
第1・第2ホモロジー群の個別の次元は、なおこの計算だけでは決まらない。
$\ker dF$ が1225次元であっても、それが第2Koszul微分の像全体に一致するとは
主張していない。補助的なKoszulグラフ探索で構成した323独立候補は、
通常階数915の証明にも上の非完全性の証明にも使っていない。
同様に、二十次でEuler交代和と通常商の次元がともに2であることは、
正次数ホモロジーの個別消滅を意味しない。

再現用の主入力は \texttt{rank/contact22\_combined.json}、
\texttt{rank/complement22\_certificate.json} と独立縮約監査である。
\texttt{rank/v16\_rank\_certificate.py} は保存された小行列・入力hash・補元値を監査する。
\texttt{--recompute} を付けると選択小行列の全成分を別の縮約方法で作り直す。
幾何学的な横断jetの非所属定理と全次数Pfaffian微分全射性は、
有限計算とは区別された原稿の証明入力である。

% END rank/rank22_section_ja.tex


% BEGIN theory/v16_koszul_height_ja.tex
% 第16版：普遍勾配イデアルのheightによるKoszul高次homologyの消失。
% 全Koszul完全性は主張しない。H1,H2が残り得ることが本質である。

\section{普遍勾配の余次元とKoszul複体の残る二つのhomology}
\label{sec:v16-universal-koszul-height}

22次では、全ての源次元の交代和と通常商の次元が一致しない。
この差の位置を特定するため、$F$ と $a$ を共に変数とする普遍係数環を用いる。
体積形式と基底を選んで
\[
 R=\mathbb C[F,a],\qquad\dim R=35+10=45,
 \qquad g_i=(\partial_iF)(q(a))\quad(0\le i\le4)
\]
と置く。固定した $F$ のJacobian環に置き換えることはしない。

\begin{proposition}[固定した三つの普遍勾配は正則列]
\label{prop:v16-universal-three-gradient-regular}
イデアル $(g_0,g_1,g_2)$ のheightは3であり、
$g_0,g_1,g_2$ は $R$ 上の正則列である。
また五つ全部の勾配イデアル $(g_0,\ldots,g_4)$ のheightも3である。
\end{proposition}

\begin{proof}
$q\ne0$ を固定すると、線形写像
\[
 \operatorname{Sym}^3V\longrightarrow V,\qquad F\longmapsto dF_q
\]
は全射である。例えば $q=e_0^*$ に基底変換すれば、
対称テンソル係数 $f_{000},f_{001},f_{002},f_{003},f_{004}$ をそれぞれ変えることで
$g_0,\ldots,g_4$ を独立に変えられる。
本稿の $F(p)=\sum f_{ijk}p_ip_jp_k$ の規約では、この部分行列は $3I_5$ である。
座標に依存しない表示として、$q(v)=1$ となる $v\in V$ を選べば、任意の $z\in V$ に対し
\[
 F_z=v^2z-\frac23q(z)v^3
 \qquad\Longrightarrow\qquad d(F_z)_q=z
\]
とも書ける。
従って $q(a)\ne0$ の開集合上では、固定した三成分への射影
$F\mapsto(g_0,g_1,g_2)$ も常に階数3である。

$a$ の開集合 $W\setminus\Sigma$ は10次元なので、
その上の $g_0=g_1=g_2=0$ の零点は階数32のベクトル束の全空間であり、次元は42である。
一方、$q=0$ 上では全ての $g_i$ が恒等的に零である。
補題\ref{lem:v16-plucker-prime-ideal}のチャートから $\dim\Sigma=7$ であるため、
この部分の次元は $35+7=42$ である。
両部分を合わせた零点集合の次元は42であり、$R$ 内のheightは3となる。
五成分を同時に零とする場合、開集合上の次元は $10+(35-5)=40$、
錐上の次元は42なので、やはりheightは3である。

多項式環は正則なCohen--Macaulay環である。
ここで大域的な正則列へ戻すために斉次性を使う。
全ての $F,a$ 変数に次数1を与えると各 $g_i$ は5次斉次であり、
斉次極大イデアル $\mathfrak m=(F,a)$ に属する。
斉次性と上のheight計算により、$R_{\mathfrak m}$ でも
$(g_0,g_1,g_2)$ のheightは3である。
従って局所的な正則列判定から、この固定した順序の三元は
$R_{\mathfrak m}$ 上の正則列となる
\cite[Corollary~3.8]{SwanMKN}。
各乗算写像の核は次数付き加群であり、非零なら非零な斉次元を含む。
そのような元は、非零定数項を持つ多項式に掛けても最低次数の項が消えないため、
$\mathfrak m$ での局所化後も非零である。
局所化後の単射性から各大域的な乗算写像も単射となる。
また商は斉次極大で非零なので、$g_0,g_1,g_2$ は $R$ 上の正則列である。
同じ局所的な可換環論の入力は\cite[Tags 00NQ, 02JN]{StacksCMRegular}でも確認できる。
ここで正則列は三つの固定した成分についての結論であり、
五つ全部の成分を正則列と仮定しているのではない。
\end{proof}

\begin{theorem}[普遍Koszul複体の高いhomologyは消える]
\label{thm:v16-koszul-high-homology-vanishing}
$K_\bullet=K(g_0,g_1,g_2,g_3,g_4;R)$ を普遍勾配のKoszul複体とする。
全ての $i>2$ について
\begin{equation}\label{eq:v16-koszul-high-homology-zero}
 H_i(K_\bullet)=0.
\end{equation}
従って、本稿で使う任意の有限二重次数とdeterminant twistの同変成分でも
$H_3,H_4,H_5$ は零である。
\end{theorem}

\begin{proof}
前命題によって $K(g_0,g_1,g_2;R)$ は
$R/(g_0,g_1,g_2)$ の自由分解である。
\[
 K_\bullet=K(g_0,g_1,g_2;R)\otimes_R K(g_3,g_4;R)
\]
と分ける。後者は自由加群からなる長さ二の複体なので、テンソル後の複体は
\[
 R/(g_0,g_1,g_2)\otimes_R K(g_3,g_4;R)
\]
と擬同型である。これはhomological degree $0,1,2$ にしか項を持たない。
これで全 $R$ 加群としての高次homologyの消失が従う。
この証明のために選んだ三成分の分割が $\mathrm{GL}(V)$ 不変である必要はない。
得られた結論は、もとの同変なKoszul複体のhomologyそのものについての消失である。

twistと二重次数も明示する。勾配は
$V\otimes\Delta^{-2}$ 値で、$F,a$ の次数は $(1,4)$ である。
従って、$k$ 番目の外積の値域を
\[
 \Lambda^k(V^*\otimes\Delta^2)\otimes\Delta^{-3}
   =\Lambda^kV^*\otimes\Delta^{2k-3}
\]
とすれば、勾配による縮約が正しい同変微分となる。
全二重次数を $(n-5,n)$ に固定すると、係数多項式の次数は
$(n-5-k,n-4k)$ であり、$a$ のtrace線を外した同変成分は正確に
\[
 \mathcal K_{k,n}=\operatorname{Hom}_{\mathrm{GL}(V)}
 \left(\operatorname{Sym}^{n-5-k}(\operatorname{Sym}^3V)
 \otimes\operatorname{Sym}^{n-4k}A,
 \Lambda^kV^*\otimes\Delta^{n-3-2k}\right)
\]
となる。負の次数を持つ項は零とする。
各二重次数の空間は有限次元であり、標数零の完全可約性によって
同変成分を取る操作は完全である。
従って全複体の式\eqref{eq:v16-koszul-high-homology-zero}は、
この有限次元複体にもそのまま降りる。
\end{proof}

\begin{corollary}[22次のKoszul defectは第二homologyに現れる]
\label{cor:v16-twentysecond-koszul-defect}
22次の同変Koszul複体の微分を
$\delta_k:\mathcal K_{k,22}\to\mathcal K_{k-1,22}$ とし、
$h_i=\dim H_i(\mathcal K_{\bullet,22})$ と書く。
命題\ref{prop:v16-character-dimensions}と
命題\ref{prop:v16-contact-rank-twenty-second}の結果から
\begin{equation}\label{eq:v16-twentysecond-koszul-defect}
 h_2=h_1+2\ge2,
 \qquad
 \operatorname{rank}\delta_5=2,
 \quad\operatorname{rank}\delta_4=55,
 \quad\operatorname{rank}\delta_3=479
\end{equation}
となる。特に22次のKoszul複体は正次数で完全ではない。
\end{corollary}

\begin{proof}
確定した次元は順に
\[
 \dim(\mathcal K_{0,22},\ldots,\mathcal K_{5,22})
   =(916,2140,1706,534,57,2)
\]
であり、その交代和は3である。
$dq$ の全射性と通常階数915により $h_0=916-915=1$ である。
前定理によって $h_3=h_4=h_5=0$ なので、Euler等式は
$1-h_1+h_2=3$ となる。
また $H_5=0$ から $\delta_5$ は単射であり、その階数は2である。
$H_4=0$ より $\ker\delta_4=\operatorname{im}\delta_5$ なので
$\operatorname{rank}\delta_4=57-2=55$、
$H_3=0$ より同様に $\operatorname{rank}\delta_3=534-55=479$ を得る。
\end{proof}

\begin{remark}[まだ決まらない階数]
この結論は $\delta_2:\mathcal B_{22}\to\mathcal D_{18}$ の全階数を決めない。
$b=\operatorname{rank}\delta_2$ と置けば
\[
 h_1=1225-b,\qquad h_2=1227-b
\]
である。構成した少数のKoszulベクトルや探索の停止を、$b$ の上界や正確な値と取り違えない。
また20次では同じ高次homology消失と今回確定した通常商次元2により
$h_2=h_1$ となるが、その共通値の零性はなお決まらない。
既存の359独立Koszulベクトルと360次元核からは、20次のこの共通値が0または1であることまで従う。
\end{remark}


% END theory/v16_koszul_height_ja.tex


% BEGIN contact/contact_twenty_second_ja.tex
% 第16版：22次の filtered contact 作用と低次正規化の補正。
% 普遍的な補正式・有限縮約と、まだ別証明を要する像所属を区別する。

\section{二十二次の接触作用と正規化補正の有限縮約}
\label{sec:v16-contact-twenty-second}

$P=P_F=F(q(a))$、$H=\det\operatorname{Hess}_F(q(a))$、$c=1/96$ とし、
$J_n=dP(\mathcal G_{n-5})$ と書く。
全ての場、単元、持ち上げは $F$ に多項式的な普遍共変量の空間で選ぶ。
固定した $F$ ごとのJacobian商を用いて持ち上げるのではない。
本節は、前版で未決定であった通常20次階数を仮定しない。

\subsection{二十二次の対数作用}

接触作用素 $T(f)=u\,f\circ\phi$ の対数を
\[
 L=\log T=\sum_{k\ge4}Q_k,\qquad
 Q_k=Y_{k+1}+b_k\in\mathcal G_{k+1}\oplus\mathcal U_k
\]
とする。第14版の補題\ref{lem:v14-contact-logarithm}により、
これは導分と乗算作用素の和であり、最小の次数増分は4である。

\begin{proposition}[二十二次の保存条件と有効作用]
\label{prop:v16-effective-twenty-second-chain}
$P+cH$ の21次までを保存する条件は
\begin{equation}\label{eq:v16-twenty-second-preservation}
 Q_kP+cQ_{k-4}H=0\quad(4\le k\le15),\qquad Q_j=0\ (j<4)
\end{equation}
であり、22次の作用は $Q_{16}P+cQ_{12}H$ である。
特に、線形条件の核
\[
 \mathcal L_{22}=\{(Q_4,Q_8,Q_{12}):
 Q_4P=0,\ Q_8P+cQ_4H=0,\ Q_{12}P+cQ_8H=0\}
\]
を用いれば
\begin{equation}\label{eq:v16-effective-twenty-second-kernel}
 I_{22}=J_{22}
    +c\{Q_{12}H:(Q_4,Q_8,Q_{12})\in\mathcal L_{22}\}.
\end{equation}
全ての普遍接触変換と、奇座標・偶単元に制限した変換は同じ22次像を与える。
さらに系\ref{cor:v16-all-hessian-first-actions}と
定理\ref{thm:v16-positive-unit-hessian-annihilation}を用いると
\begin{equation}\label{eq:v16-effective-twenty-second-ordinary}
                         I_{22}=J_{22}.
\end{equation}
\end{proposition}

\begin{proof}
$T(P+cH)-(P+cH)$ が22次以上であることと、
$L(P+cH)$ が22次以上であることは対数補題により同値である。
この場合、22次では両者が一致する。
従って次数ごとに比較すれば保存条件と作用式を得る。

偶数増分の条件は
\[
\begin{array}{lll}
Q_4P=0,&Q_8P+cQ_4H=0,&Q_{12}P+cQ_8H=0,\\
Q_6P=0,&Q_{10}P+cQ_6H=0,&Q_{14}P+cQ_{10}H=0
\end{array}
\]
という二本の鎖に分かれる。
22次作用に現れるのは第一の鎖だけで、第二の鎖は全て零に選べる。
奇数増分の条件も独立なので、それらも零に選べる。
また $P\mathcal U_{16}\subset J_{22}$ より $Q_{16}P$ の像は $J_{22}$ である。
これで\eqref{eq:v16-effective-twenty-second-kernel}を得る。

一般の $Q_{12}$ に対して
\[
 Q_{12}H=dH(Y_{13})+b_{12}H
       \in dH(\mathcal G_{13})+H\mathcal U_{12}\subset J_{22}
\]
である。逆包含 $J_{22}\subset I_{22}$ は $Q_{16}$ の自由な作用から従う。
ここでは通常20次商の次元も、22次の通常像の階数も用いない。
\end{proof}

\subsection{merged規約からの正確な二十二次台帳}

未正規化Tate係数を $f_n=d_n^{\mathrm T}$ とし、前版と同じ
\[
 X_7=-\frac{15800}{6137}\Phi_7-\frac{689855}{883728}\Psi_7,
 \qquad X_8=\frac{4600}{2189}\operatorname{Alt}K_8
\]
を用いる。初回変更は厳密に
\[
 g=(1+3U_4-V_6)f(a+X_7+X_8)
\]
とし、単元を積 $(1+3U_4)(1-V_6)$ に変更しない。
既存の恒等式は
\[
 f_{10}=cH-3PU_4,\quad
 f_{12}+dP(X_7)=PV_6,\quad f_{13}+dP(X_8)=0
\]
である。従って $g_6=P$、$g_{10}=cH$、$g_{12}=g_{13}=0$ であり、
\[
 A_{14}:=g_{14}=f_{14}+3U_4f_{10},\qquad g_{15}=f_{15},
\]
\[
 A_{16}:=g_{16}=f_{16}+df_{10}(X_7)
       +3U_4(f_{12}+dP(X_7))-V_6f_{10}.
\]
これらの14、15、16次項は、それぞれ $J_{14},J_{15},J_{16}$ に属する。

以下では、Euler場で単元項を吸収した純粋な持ち上げを選び、
\[
 dP(\Xi_9)=A_{14},\qquad dP(\Xi_{10})=f_{15},\qquad
 dP(\Xi_{11})=A_{16}
\]
により、順に $a-\Xi_9$、$a-\Xi_{10}$、$a-\Xi_{11}$ で消去する。
ここで
\[
 g_{18}=f_{18}+df_{12}(X_7)+\tfrac12P''(X_7,X_7)
           +3U_4f_{14}-V_6(f_{12}+dP(X_7))
\]
であり、18次消去の入力は
\begin{equation}\label{eq:v16-eighteenth-normalizer-input}
 A_{18}=g_{18}-c\,dH(\Xi_9).
\end{equation}
第14版の所属から $A_{18}\in J_{18}$ なので、
$dP(\Xi_{13})=A_{18}$ として $a-\Xi_{13}$ で消す。
20次も前版の定理により $J_{20}$ に入り、15次場で消せる。

\begin{proposition}[二十二次の正確な有限補正式]
\label{prop:v16-exact-twenty-second-correction}
初回変更の22次係数は
\begin{equation}\label{eq:v16-initial-twenty-second}
\begin{split}
 g_{22}={}&f_{22}+df_{16}(X_7)+df_{15}(X_8)
                +\tfrac12 f_{10}''(X_7,X_7)\\
 &+3U_4\bigl(f_{18}+df_{12}(X_7)+\tfrac12P''(X_7,X_7)\bigr)
             -V_6\bigl(f_{16}+df_{10}(X_7)\bigr).
\end{split}
\end{equation}
上の純粋な有限正規化の後の22次係数は正確に
\begin{equation}\label{eq:v16-exact-twenty-second}
 R_{22}^{\mathrm{fin}}
 =g_{22}-dA_{14}(\Xi_9)+\tfrac12P''(\Xi_9,\Xi_9)
                         -c\,dH(\Xi_{13}).
\end{equation}
15、16、20次消去は22次係数に寄与しない。
\end{proposition}

\begin{proof}
初回の座標場の次数増分は6と7である。
22次に到達する一次項は $df_{16}(X_7)$ と $df_{15}(X_8)$、
二次項は $\tfrac12 f_{10}''(X_7,X_7)$ である。
混合二次項に必要な $f_9$ と、$X_8$ の平方に必要な $f_8$ は零である。
三階以上のTaylor項は、最小次数の $P$ に作用しても24次以上となる。
四次単元を掛けるには18次係数、六次単元を掛けるには16次係数を用いるので、
\eqref{eq:v16-initial-twenty-second}の残る二項を得る。
特に $3U_4P''(X_7,X_7)/2$ も22次に含まれている。

14次消去の増分は8である。
22次には14次項の一次微分と、$P$ の二階微分が寄与する。
15次、16次消去の増分は9、10であるが、対応する13次、12次係数は既に零である。
18次消去の増分12は10次項 $cH$ に作用するので、
$-c\,dH(\Xi_{13})$ を加える。
20次消去の増分14に対応する8次係数も零である。
各段の二階以上の残る寄与は、表示した14次消去の平方を除いて22次を超える。
\end{proof}

\nolinkurl{v16_degree_ledger.py} は、元の係数の次数、各座標場の次数増分、
merged単元の三項を独立に列挙し、初回の九項、単元を残す14次消去の五項、
後段で18次消去のHessian作用だけが残ることを有理数のTaylor係数とともに検算する。
この次数検算は、接触像への所属や全次数偶化の証明とは別の確認である。

\begin{remark}[単元を残す規約での同じ台帳]
14次を $A_{14}=dP(\Xi'_9)+PB_8$ とし、
$(1-B_8)g(a-\Xi'_9)$ で消す場合、その22次係数は正確に
\[
 g_{22}-dA_{14}(\Xi'_9)+\tfrac12P''(\Xi'_9,\Xi'_9)-PB_8^2
\]
となる。最後の項を落としてはならない。
この規約での18次入力は $g_{18}-c(dH(\Xi'_9)+B_8H)$ である。
その18次入力を $dP(\Xi'_{13})+PB_{12}$ として消すなら、
22次へ加える項は $-c(dH(\Xi'_{13})+B_{12}H)$ となる。
15、16、20次の単元からの22次寄与は、それぞれ13、12、8次係数が零なので消える。
これらの規約差を商で捨てる際は、$P\mathcal U_{16}\subset J_{22}$、
二階作用のEuler方向、および $H\mathcal U_{12}\subset J_{22}$ を用いる。
\end{remark}

\subsection{通常商に残る三つの二階作用}

$f_{16}\in J_{16}$、$f_{15}\in J_{15}$ により
\[
 dP(Z_{16})=f_{16},\quad Z_{16}\in\mathcal G_{11},\qquad
 dP(Z_{15})=f_{15},\quad Z_{15}\in\mathcal G_{10}
\]
と選ぶ。添字 $Z_{16},Z_{15}$ は持ち上げる係数の次数を表す。

\begin{theorem}[二十二次の実際の類の代表]
\label{thm:v16-twenty-second-correction-representative}
通常像 $J_{22}$ を法として
\begin{equation}\label{eq:v16-twenty-second-correction-representative}
\begin{split}
 [R_{22}^{\mathrm{fin}}]=\bigl[ f_{22}
  &+P''(X_7,Z_{16})+P''(X_8,Z_{15})
    +\tfrac c2H''(X_7,X_7)\\
  &-\tfrac12P''(\Xi_9,\Xi_9)-c\,dH(\Xi_{13})\bigr].
\end{split}
\end{equation}
最後の項は $J_{22}$ に属する。
さらに $dH(\mathcal G_7)\subset J_{16}$ の純粋な線形持ち上げ
\[
 W:\mathcal G_7\longrightarrow\mathcal G_{11},\qquad
 dP(W(X))=dH(X)
\]
を選ぶと
\begin{equation}\label{eq:v16-three-quadratic-families}
\boxed{\displaystyle
 [R_{22}^{\mathrm{fin}}]=
 \left[f_{22}+P''\!\left(X_7,Z_{16}+\tfrac c2W(X_7)\right)
       +P''(X_8,Z_{15})-\tfrac12P''(\Xi_9,\Xi_9)\right].}
\end{equation}
従って必要な残余補正は
$P''(\mathcal G_7,\mathcal G_{11})$、
$P''(\mathcal G_8,\mathcal G_{10})$、
$P''(\mathcal G_9,\mathcal G_9)$ の三族だけである。
\end{theorem}

\begin{proof}
恒等式 $dP(Z_{16})=f_{16}$ を $X_7$ で微分し、
$dP(Z_{15})=f_{15}$ を $X_8$ で微分すると
\[
 df_{16}(X_7)\equiv P''(X_7,Z_{16}),\qquad
 df_{15}(X_8)\equiv P''(X_8,Z_{15})\pmod{J_{22}}
\]
を得る。同様に $A_{14}=dP(\Xi_9)$ から
\[
 -dA_{14}(\Xi_9)+\tfrac12P''(\Xi_9,\Xi_9)
       \equiv-\tfrac12P''(\Xi_9,\Xi_9)\pmod{J_{22}}
\]
である。

第14版の命題\ref{prop:v14-all-lower-corrections-in-image}により
$P''(\mathcal G_7,\mathcal G_7)\subset J_{18}$ である。
また $f_{18}\in J_{18}$、$f_{16}\in J_{16}$ であり、
$f_{12}=PV_6-dP(X_7)$ を微分すると
\[
 df_{12}(X_7)+\tfrac12P''(X_7,X_7)\in J_{18}.
\]
さらに $df_{10}(X_7)\in J_{16}$ は
$f_{10}=cH-3PU_4$ と $dH(\mathcal G_7)\subset J_{16}$ から従う。
従って\eqref{eq:v16-initial-twenty-second}の単元を含む最後の二群は全て $J_{22}$ に入る。

$f_{10}''=cH''-3(PU_4)''$ において、
\[
 (PU_4)''(X_7,X_7)
 =U_4P''(X_7,X_7)+2dP(X_7)dU_4(X_7)
                        +P\,U_4''(X_7,X_7)
\]
の三項も $J_{22}$ に入る。
これで\eqref{eq:v16-twenty-second-correction-representative}を得る。

最後に $dP(W(X_7))=dH(X_7)$ を $X_7$ で微分すると
\[
 H''(X_7,X_7)+dH(DX_7\cdot X_7)
 =P''(W(X_7),X_7)+dP(D(W(X_7))\cdot X_7).
\]
$DX_7\cdot X_7\in\mathcal G_{13}$ と
$dH(\mathcal G_{13})\subset J_{22}$ により、
$H''(X_7,X_7)\equiv P''(W(X_7),X_7)$ である。
同じ包含により $dH(\Xi_{13})$ も消える。
全ての純粋な持ち上げの差は $dP$ の核に入り、
第15版の補題\ref{lem:v15-hessian-pairing-quotient}により表示した類に影響しない。
\end{proof}

\subsection{十七次の奇項と全次数偶化の比較}

22次では、前版の「17次消去を明示しなくてよい」という議論を、
次数の確認なしにそのまま使うことはできない。
そこで、17次の奇項が通常像に入ることを先に示す。

\begin{lemma}[残る十七次項の純粋消去]
\label{lem:v16-seventeenth-odd-removal}
上の有限正規化の後に残る17次項 $R_{17}$ は $J_{17}$ に入る。
12次の普遍座標場によってこれを消しても、22次係数は変わらない。
\end{lemma}

\begin{proof}
有限正規化の境界と、第15版の全次数偶正規形を比較する。
両者の差は17次以上なので、対数補題により17次の差は
\[
 -R_{17}=Q_{11}P+cQ_7H
\]
である。$Q_{11}P\in J_{17}$ であり、
系\ref{cor:v16-all-hessian-first-actions}の $r=8$ の場合と
$\mathcal U_7=0$ から $Q_7H=dH(Y_8)\in J_{17}$ となる。
従って $R_{17}\in J_{17}$ である。

$dP(\Xi_{12})=R_{17}$ を選び $a-\Xi_{12}$ で消す。
22次への一次寄与には11次係数が必要だが零である。
$P$ の二階寄与は28次からなので、22次を変えない。
18次、20次の既に消した係数も変わらない。
\end{proof}

\begin{corollary}[二十二次の偶正規形との一致]
\label{cor:v16-evenization-twenty-second-class}
第15版の全次数偶正規形における22次係数 $R_{22}^{\mathrm{ev}}$ と、
上の有限正規化の係数は
\[
 R_{22}^{\mathrm{ev}}-R_{22}^{\mathrm{fin}}\in I_{22}=J_{22}
\]
を満たす。従って式\eqref{eq:v16-three-quadratic-families}は、
実際の全次数偶正規形の22次類を与える。
\end{corollary}

\begin{proof}
前補題により17次を消した後は
\[
 f^{\mathrm{fin}}=P+cH+R_{19}+R_{21}+R_{22}^{\mathrm{fin}}+O(a^{23})
\]
である。偶正規形と比較する $T$ に対し $Tf^{\mathrm{fin}}-f^{\mathrm{fin}}$ は19次以上なので、
対数補題により $Tf^{\mathrm{fin}}-f^{\mathrm{fin}}\equiv Lf^{\mathrm{fin}}\pmod{(a)^{23}}$。
最小増分4から $L(R_{19}+R_{21})$ も23次以上である。
19次、21次の非零な比較式は奇数増分の条件にだけ入り、
偶数増分の鎖は\eqref{eq:v16-twenty-second-preservation}のままである。
従って22次の差は\eqref{eq:v16-effective-twenty-second-kernel}の像に入る。
\end{proof}

\begin{remark}[十七次を残した場合の二階対数項]
17次を残して直接比較する場合、22次には
\[
 (Tf-f)_{22}=Q_{16}P+cQ_{12}H+\tfrac12Q_5R_{17}
\]
という追加項が現れる。
実際、$Lf$ の17次項は $-R_{17}$、18次項は零なので、
$\tfrac12L^2f$ の22次項は $-\tfrac12Q_5R_{17}$ である一方、
$Lf$ 自身には $Q_5R_{17}$ が入る。
ただし $Q_5P=0$ なので $Q_5$ は通常像を保存し、
$R_{17}\in J_{17}$ から $Q_5R_{17}\in J_{22}$ である。
具体的に $Q=Y+b$、$QP=0$、$R=dP(Z)$ なら
\[
 QR=dP(DZ\cdot Y-DY\cdot Z)-P\,db(Z)
\]
となる。
従って直接比較しても同じ類を得るが、追加項の次数を無視してよいわけではない。
任意の低次安定化を含む座標変更の三階以上のTaylor項についても、
ここでは作用素全体の対数と低次保存条件がその相殺を含めて扱っている。
\end{remark}

\subsection{七十六列による十分な所属検査}

既存の低次像の次元は
\[
\begin{array}{c|ccccc}
 &J_{12}/P\mathcal U_6&J_{13}/P\mathcal U_7&J_{14}/P\mathcal U_8
 &J_{15}/P\mathcal U_9&J_{16}/P\mathcal U_{10}\\\hline
\dim&2&1&7&6&21
\end{array}
\]
である。各商の基底を持ち上げた場を、それぞれ
$G_{7,i}$、$G_8$、$G_{9,j}$、$G_{10,k}$、$G_{11,l}$ とする。

\begin{proposition}[二十二次補正の有限族への帰着]
\label{prop:v16-seventy-six-lower-pairings}
以下の76個の共変量が全て $J_{22}$ に属すると仮定する。
\[
\begin{array}{ll}
 \operatorname{Hess}_F(q)(dq(G_{7,i}),dq(G_{11,l})),&1\le i\le2,\ 1\le l\le21,\\
 \operatorname{Hess}_F(q)(dq(G_8),dq(G_{10,k})),&1\le k\le6,\\
 \operatorname{Hess}_F(q)(dq(G_{9,i}),dq(G_{9,j})),&1\le i\le j\le7.
\end{array}
\]
このとき
\begin{equation}\label{eq:v16-raw-twenty-second-suffices}
 [R_{22}^{\mathrm{ev}}]=[d_{22}^{\mathrm T}]
                         \quad\text{in }\mathcal C_{22}/J_{22}.
\end{equation}
ここで仮定は22次の普遍所属であり、少数点での零値だけに置き換えてはならない。
\end{proposition}

\begin{proof}
第15版の二階作用の補題により、$P''(Y,Z)$ の通常類は各 $dP$ 像を
$P$ 倍の単元像で割った商だけに依存する。
また三族のいずれも $d^2q(Y,Z)\in\mathcal D_{18}$ なので、
\[
 P''(Y,Z)=\operatorname{Hess}_F(q)(dqY,dqZ)+dF(d^2q(Y,Z))
\]
の最後の項は、全次数Pfaffian全射によって $J_{22}$ に入る。
従って76個の代表の所属から三族全体の所属が従う。
式\eqref{eq:v16-three-quadratic-families}と
系\ref{cor:v16-evenization-twenty-second-class}を適用すればよい。
\end{proof}

\paragraph{計算入力と現在の検査範囲.}
7、9、11次場は既存のグラフ基底、8次場は第15版で非零像を確認したグラフを用いる。
10次については、旧版のTate部分項の代わりに六つのグラフ場を新たに作った。
一般の $P=0$ 入力での $6\times6$ 行列を有理数上で再構成し、
非零行列式と、その1009を法とする値934を確認した。
既存の $\dim(J_{15}/P\mathcal U_9)=6$ と合わせ、これらも基底を与える。
入力と行列は \nolinkurl{v16_g10_graph_basis.json} に保存した。

\nolinkurl{v16_lower_pairings_exact.py} は第15版の $F_D,F_H,a_b$ を用いる。
四点
\[
 (F_D,a_1),\quad(F_D,a_2),\quad(F_H,a_1),\quad(F_H,a_2)
\]
で76列を有理数上で再計算する。
全76列は、この四点の全てで点ごと厳密に零となった。
従って候補汎関数
\[
 \widetilde\rho_{22}
 =\frac{16\operatorname{ev}_{(F_H,a_1)}-
                \operatorname{ev}_{(F_H,a_2)}}{3459456000}
\]
は全76列を消し、$\Theta_8U_4^2J_6$ に対する値は1である。
この評価結果だけでは、上の76列の $J_{22}$ 所属や、
この汎関数が通常接触像全体を消すことは従わない。
本稿では命題\ref{prop:v16-exact-rho22}で通常像と核の一致を証明し、
系\ref{cor:v16-all-lower-corrections-in-image}で76列の普遍所属を確定する。

% END contact/contact_twenty_second_ja.tex


% BEGIN detector/v16_detector_notes_ja.tex
% 第16版：全429接触方向のQ再縮約による22次通常商の厳密検出子。

\section{二十二次通常商を読む二点汎関数}
\label{sec:v16-exact-twenty-second-detector}

命題\ref{prop:v16-contact-rank-twenty-second}により
\[
 \dim\mathcal C_{22}=916,\qquad
 \dim J_{22}=915,\qquad
 J_{22}=P\mathcal U_{16}
    \oplus\operatorname{Span}\{dP(\Gamma_1),\ldots,dP(\Gamma_{429})\}
\]
である。ここで $\dim\mathcal U_{16}=486$、
429方向は286個の連結グラフと143個の下位場・単元の積からなる。
通常商の生成元を
\[
                      W_{22}=\Theta_8U_4^2J_6
\]
とする。本節では、これら全ての方向の二点での値を有理数上で再構成し、
通常商の係数を読む汎関数を確定する。

\subsection{全429方向の有理数による再縮約}

第15版と同じ
\[
 F_H=x_0^3+x_1^3+x_2^3+x_3^3+4x_4^3-6x_2x_3x_4,
\]
\[
 a_b=b\,e_0\wedge e_1+b^{-1}(e_2-e_4)\wedge(e_3-e_4)
\]
を用い、$z_1=(F_H,a_1)$、$z_2=(F_H,a_2)$ とする。
両点で $q=(0,0,1,1,1)$、$P=0$ であり、
\[
 \Theta_8(F_H,a_b)=-4b^2,\qquad
 U_4(F_H,a_b)=5b^4,\qquad J_6(F_H,a_b)=528b^6.
\]
従って
\begin{equation}\label{eq:v16-complement-two-values}
 W_{22}(z_1)=-52800,\qquad
 W_{22}(z_2)=-3460300800.
\end{equation}

\begin{proposition}[二十二次の厳密な係数汎関数]
\label{prop:v16-exact-rho22}
$\operatorname{ev}_i(R)=R(z_i)$ とし、
\begin{equation}\label{eq:v16-exact-rho22}
 \boxed{\displaystyle
 \rho_{22}=
 \frac{16\operatorname{ev}_1-\operatorname{ev}_2}{3459456000}}
\end{equation}
と置く。このとき
\[
 \rho_{22}(J_{22})=0,\qquad
 \rho_{22}(W_{22})=1,\qquad
                       \ker\rho_{22}=J_{22}.
\]
特に任意の $R\in\mathcal C_{22}$ に対して
\[
 [R]=\rho_{22}(R)[W_{22}]
                        \quad\text{in }\mathcal C_{22}/J_{22}
\]
である。汎関数は $\mathbb Q$ 上で定義し、$\mathbb C$ へ係数拡大して用いる。
\end{proposition}

\begin{proof}
429個の公開グラフを全て有理数上で縮約し、各列について
\[
                      dP(\Gamma_j)(z_2)=16dP(\Gamma_j)(z_1)
\]
を確認した。従って両評価の表示した結合は全429方向を消す。
$P(z_i)=0$ なので $P\mathcal U_{16}$ の486方向も点ごと零である。
式\eqref{eq:v16-complement-two-values}から
\[
 16W_{22}(z_1)-W_{22}(z_2)=3459456000
\]
となり、正規化は1である。
具体的に公開配列の33列目の接触方向と補元を並べた評価行列は
\[
 \begin{pmatrix}
 1920&-52800\\
 30720&-3460300800
 \end{pmatrix},\qquad
 \det=-6642155520000\ne0
\]
である。
この検出子は非零なので $\dim\ker\rho_{22}=916-1=915$。
通常像の次元と包含 $J_{22}\subset\ker\rho_{22}$ から核の一致を得る。
\end{proof}

\paragraph{積方向と正規化の検査.}
143個の積方向では、下位場 $\Gamma_r$ と単元グラフ $u_s$ は $r+s=17$ を満たし、
\[
 dP(u_s\Gamma_r)=u_s\,dP(\Gamma_r)
\]
という普遍恒等式を用いる。
単元のグラフを先に、下位場のグラフを後に置いた非連結17頂点グラフは、
最後の左頂点を接触方向の自由頂点に保つ。
\nolinkurl{v16_rho22_exact_certificate.py} は各積について、
この全17頂点グラフの値と、二つの因子を別々に縮約した値の積を、
両点で正確に照合した。
特に四頂点の単元はグラフそのものの値 $24U_4$ を用いており、
143列の途中で24による除算を挿入していない。
一方、補元 $W_{22}$ には既存の正規化された $U_4$ を用いる。

有理数の $a$ は一度 $Da$ にして整数縮約し、
全接触列は22次の斉次性により $D^{22}$ で割り戻す。
積の各因子もそれぞれの次数で割り戻し、同じ値になることを確認する。
これは全450点の有限体minorによる独立性の監査とは異なる、有理数上の所属判定である。

\subsection{低次補正の全消滅}

\begin{corollary}[二十二次の全低次正規化補正は通常像に入る]
\label{cor:v16-all-lower-corrections-in-image}
\[
 \begin{split}
 P''(\mathcal G_7,\mathcal G_{11})&\subset J_{22},\\
 P''(\mathcal G_8,\mathcal G_{10})&\subset J_{22},\\
 P''(\mathcal G_9,\mathcal G_9)&\subset J_{22}.
 \end{split}
\]
従って実際の22次偶正規形の係数は
\begin{equation}\label{eq:v16-all-corrections-vanish}
 R_{22}^{\mathrm{ev}}-d_{22}^{\mathrm T}\in J_{22}
\end{equation}
を満たす。
\end{corollary}

\begin{proof}
命題\ref{prop:v16-seventy-six-lower-pairings}の76列は、
$z_1,z_2$ のそれぞれで有理数上の再縮約により全て零であった。
従って前命題により全76列は $J_{22}$ に入る。
同命題のkernel方向、Euler方向、および $d^2q$ の持ち上げの議論によって、
三族全体が $J_{22}$ に入る。

Hessian一次作用 $dH(\mathcal G_{13})$ と $H\mathcal U_{12}$ の全体は、
前節までの構造定理により既に $J_{22}$ に入っている。
$H''(\mathcal G_7,\mathcal G_7)$ も
定理\ref{thm:v16-twenty-second-correction-representative}で第一の族へ還元した。
従って有限正規化の全補正が消え、
系\ref{cor:v16-evenization-twenty-second-class}により全次数偶化後の係数にも同じ結論が成り立つ。
\end{proof}

\subsection{実際の二十二次類}

\begin{proposition}[二十二次の実際の境界類の消滅]
\label{prop:v16-actual-twenty-second-class}
未正規化Tate係数と実際の正規化後の係数は、ともに $J_{22}$ に属する。
\[
                d_{22}^{\mathrm T}\in J_{22},\qquad
                [R_{22}^{\mathrm{ev}}]=0.
\]
\end{proposition}

\begin{proof}
命題\ref{prop:v16-fresh-tate22-values}による新規Tate再帰の値を
\[
 A_{22}=\frac{1329012984671053041316939}{630118440000000000}
\]
と書けば、
\[
             d_{22}^{\mathrm T}(z_1)=A_{22},\qquad
             d_{22}^{\mathrm T}(z_2)=16A_{22}.
\]
従って $\rho_{22}(d_{22}^{\mathrm T})=0$ であり、
前命題の核の同定から $d_{22}^{\mathrm T}\in J_{22}$ となる。
式\eqref{eq:v16-all-corrections-vanish}を用いれば正規化後の類も零である。
対角入力 $(F_D,a_1),(F_D,a_2)$ の新規値はいずれも零であり、
四点の全データは同じTate命題とその証明書に収めた。
\end{proof}

\paragraph{証明書の範囲.}
\nolinkurl{v16_rho22_exact_certificate.json} には、429グラフの両点での有理数値、
143積方向の各因子と積、補元値、選択した小行列、76補正の汎関数値を保存した。
Tate入力については全三次テンソルと交代行列 $a$ を一致させ、
保存された18個の定数ブロックから0次〜22次の行列式係数を再計算した。
この保存ブロックの監査を、最初からのMaurer--Cartan再帰生成とは区別する。
通常像の独立性、三元切片による補元の非所属、低次補正の数学的な縮約は、
それぞれの命題を別の入力として用いる。

% END detector/v16_detector_notes_ja.tex


% BEGIN tate/v16_tate_notes_ja.tex
\section{二十二次Tate係数の新規再構成}
\label{sec:v16-tate-twentysecond}

二十二次では、前版の十六個の定数ブロックを延長するだけでは足りない。
十八個の定数ブロックに必要な内部次数を確保するため、四つの入力点で
$A_1$ からMaurer--Cartan再帰を再実行した。
さらに一つの対角点では有理数の実装を変えて再実行し、
全十八ブロックと全行列式係数を比較した。

\subsection{十八次ブロックを得る切断}

\begin{lemma}[二十二次の厳密な内部次数切断]
\label{lem:v16-tate22-truncation}
二十二次以下の行列式係数には $K_1,\ldots,K_{18}$ だけが必要である。
これらは $A_1,\ldots,A_{14}$ を全て計算し、続いて
$R_{15},R_{16},R_{17},R_{18}$ をそれぞれ内部 $x$ 次数
$4,3,2,1$ 以下で求めれば厳密に得られる。
対応する $A_{15},A_{16},A_{17},A_{18}$ は、それぞれ
$x$ 次数 $3,2,1,0$ 以下まで正確である。
\end{lemma}

\begin{proof}
行列式の各項は正次数の五つのブロックの積なので、総次数22以下では
各ブロックの次数は18以下である。
Clifford積は $x$ 次数を下げず、Eulerホモトピーは1下げる。
内部重み $2\deg_x-\#\theta+\#\eta=3$ により、各 $A_r$ の
$x$ 次数は4以下である。
$n=15,\ldots,18$ の曲率に必要な $x$ 次数は $19-n$ 以下であり、
その積の各因子の添字 $i\le n-1$ に対して
$18-i\ge19-n$ である。従って、それ以前の切断が必要な項を失わないことが
帰納的に従う。ホモトピーを適用すると $18-n$ 次以下の補正を得る。
\end{proof}

第15版で保存した $A_{13},A_{14}$ は内部次数を切断していたため、
本節では両者の全項を改めて計算している。
Hesse型の計算では補題\ref{lem:v15-hesse-symbol-product}の
$S=1+C$、$S^{-1}=1-C$ を用い、常に自然な外積記号に戻してから
Eulerホモトピーを適用する。
このClifford積の実装は第15版からバイト単位で変更していない。

\subsection{四点の厳密な二十二次値}

\[
\begin{split}
 F_D&=x_0^3+x_1^3+x_2^3+x_3^3-2x_4^3,\\
 F_H&=x_0^3+x_1^3+x_2^3+x_3^3+4x_4^3-6x_2x_3x_4,\\
 a_b&=b\,e_0\wedge e_1+
       b^{-1}(e_2-e_4)\wedge(e_3-e_4)\qquad(b=1,2)
\end{split}
\]
とする。いずれの点でも $q(a_b)=(0,0,1,1,1)$、$P=0$ である。
この四つの入力点では $H_{10}=0$ であり、$q$ 変数に関する
Hessian行列式の勾配も零となる。

\Needspace{7\baselineskip}
\begin{proposition}[新規再構成による四点の二十二次係数]
\label{prop:v16-fresh-tate22-values}
\[
 A_{22}=\frac{1329012984671053041316939}{630118440000000000}
\]
とおくと、未正規化のTate行列式係数は
\[
\boxed{\bigl(d_{22}^{\mathrm T}(F_D,a_1),d_{22}^{\mathrm T}(F_D,a_2),
             d_{22}^{\mathrm T}(F_H,a_1),d_{22}^{\mathrm T}(F_H,a_2)\bigr)
       =(0,0,A_{22},16A_{22}).}
\]
従って命題\ref{prop:v16-exact-rho22}の検出子について
$\rho_{22}(d_{22}^{\mathrm T})=0$ である。
\end{proposition}

\begin{proof}
各点の再帰を最初から有理数上で実行し、保存した内部次数の全範囲で
曲率の閉性とMaurer--Cartan方程式を照合した。
定数ブロックの規約は $K_1=-a$、$K_2=\operatorname{Hess}_F(q)/3$、
$d^{\mathrm T}=\frac12\det(\sum K_r)$ である。
行列式を120置換による展開と列部分集合による動的計画法で別々に計算し、
0次から22次までの全23係数を有理数上で比較した。
得られた非零の二十二次値は
\[
 d_{22}^{\mathrm T}(F_H,a_1)
   =\frac{1329012984671053041316939}{630118440000000000},\qquad
 d_{22}^{\mathrm T}(F_H,a_2)
   =\frac{1329012984671053041316939}{39382402500000000}.
\]
後者は前者の16倍であるため、検出子の分子
$16d_{22}^{\mathrm T}(F_H,a_1)-d_{22}^{\mathrm T}(F_H,a_2)$ は零である。
\end{proof}

\subsection{独立検証と証明の範囲}

各再帰では、PBW積とその基底変換を用いない32次元外積作用素を
$x=(1,2,3,4,5)$、素数 $1000003$ で別に構成し、
完全な $A_1,\ldots,A_{14}$ から得る0次から14次までの作用素関係を検証した。
全十八の定数ブロックにも独立な外積作用素による照合を行った。
これは有理数上の再帰と切断証明に加える実装検査である。

$(F_D,a_1)$ では、共通分母による整数実装と各係数を
\texttt{Fraction} とする実装を独立に走らせた。
両者の全十八ブロック、全23行列式係数、二十二次値が一致した。
さらに、前版に同じ入力がある $(F_D,a_1),(F_D,a_2),(F_H,a_1)$ では、
前版の十六ブロックと0次から20次までの全係数に一致した。
これらは保存データに対する標準ライブラリだけの独立監査でも再確認した。

配布する \texttt{v16\_tate\_saved\_audit.py} は保存ブロックから
行列式を再計算する監査であり、新規Maurer--Cartan再帰そのものを再実行しない。
新規再帰用の二つのドライバ、四点五回の出力、ログを別に同梱する。
局所的なチェックポイントは配布物に含めず、ドライバも既存の
チェックポイントを入力として読まない。

本節の四点での値だけから普遍接触像への所属は従わない。
命題\ref{prop:v16-exact-rho22}の完全な商の検出、
系\ref{cor:v16-all-lower-corrections-in-image}の低次正規化補正の消滅と
系\ref{cor:v16-evenization-twenty-second-class}の偶正規形との比較を合わせて初めて、
実際の二十二次接触類が零であることを結論する。

% END tate/v16_tate_notes_ja.tex


% BEGIN contact/twenty_fourth_normal_form_ja.tex
% 第16版の確定結論。22次像、全低次補正、fresh Tate値に依存する。

\section{二十二次の消去と二十四次開始の正規形}
\label{sec:v16-twenty-fourth-normal-form}

\begin{theorem}[二十四次開始の普遍代数的収束偶正規形]
\label{thm:v16-twenty-fourth-normal-form}
実際の境界の22次類は零であり、任意の滑らかな三次式 $F$ に対して
\begin{equation}\label{eq:v16-twenty-fourth-normal-form}
 \boxed{\displaystyle
 d^{\mathrm T}\sim
 P_F+\frac1{96}\det\operatorname{Hess}_F(q(a))
       +\sum_{\substack{n\ge24\\n\ \mathrm{even}}}R_n,
 \qquad R_n\in\mathcal C_n}
\end{equation}
という全次数の偶正規形が存在する。
接触変換は普遍同変な代数的冪級数であり、座標変更の線形項は恒等写像、
単元の定数項は1である。各形式係数は $F$ の係数に多項式的で、
各固定 $F$ では変換・逆変換・正規形の級数は原点近傍で収束する。
24次以降の実際の類の消滅は、この定理の主張に含めない。
\end{theorem}

\begin{proof}
第15版の定理\ref{thm:v15-twenty-second-normal-form}により
\[
 f_{\mathrm{ev}}=P_F+\frac1{96}H+R_{22}^{\mathrm{ev}}
                             +R_{24}^{\mathrm{ev}}+\cdots
\]
という普遍同変な代数的収束偶正規形がある。
系\ref{cor:v16-all-lower-corrections-in-image}は
\[
 R_{22}^{\mathrm{ev}}-d_{22}^{\mathrm T}\in J_{22}
\]
を示す。この比較には、17次の奇項の純粋消去と、
全偶化との対数作用による比較も含まれている。
命題\ref{prop:v16-exact-rho22}と
命題\ref{prop:v16-fresh-tate22-values}により
\[
 \rho_{22}(d_{22}^{\mathrm T})
   =\frac{16A_{22}-16A_{22}}{3459456000}=0,
 \qquad\ker\rho_{22}=J_{22}.
\]
従って $R_{22}^{\mathrm{ev}}\in J_{22}$ である。
ここで $I_{22}=J_{22}$ は構造定理でも確定しており、
異なる下位正規化による有効類と通常類の区別も解消している。

有限次元の普遍共変量空間で
\[
 \Xi_{17}\in\mathcal G_{17},\qquad
                       dP(\Xi_{17})=R_{22}^{\mathrm{ev}}
\]
と選ぶ。これは $F$ の判別式等で割る操作ではなく、
$F$ に16次、$a$ に17次の普遍多項式な場の持ち上げである。
奇座標変更 $a\mapsto a-\Xi_{17}(F,a)$ を行うと22次項が消える。
Hessian項への最初の寄与は26次、$P$ の二階Taylor項は38次からであるため、
既存の低次正規形を保ち、次に残り得る次数は24となる。
座標変更が奇関数なので全次数の偶性も保たれる。

今回の追加変更は有限多項式で、原点でのJacobianは恒等写像である。
その形式逆は代数的冪級数であり、固定した $F$ に対しては解析的にも局所収束する。
第15版の変換との有限合成は、代数性、局所収束性、普遍同変性、
各形式係数の $F$ 多項式性を保つ。
これで式\eqref{eq:v16-twenty-fourth-normal-form}を得る。
\end{proof}

\begin{remark}[確定した商と残る課題]
通常22次商は一次元で、非零な $[\Theta_8U_4^2J_6]$ によって生成される。
その商で実際のTate境界の係数は零であった。
また、下位jetが $P+cH$ である場合の $I_N=J_N$ は全次数の構造定理であるが、
これ自体は実際の高次係数の零性を与えない。
通常商の全次数の構造と、24次以降の実際の類は引き続き調べる必要がある。
特に全境界が $P_F+H/96$ と接触同値であることは、まだ結論していない。
\end{remark}

% END contact/twenty_fourth_normal_form_ja.tex

\clearpage
% BEGIN theory/v17_theta_annihilator_ja.tex
% 第17版：Theta8の全単元環におけるannihilator。
% 新しい外部入力は不要。第16版のPluecker素性・dq全射・横断単射を継承。

\section{行列束の恒等式と八次共変量の消滅イデアル}
\label{sec:v17-theta-annihilator}

第16版の記号
\[
 \Delta=\det V,\quad E=V^*\otimes\Delta^{-1},\quad
 L=\Delta^{-3},\quad
 \mathcal U=\bigoplus_{s\ge0}\mathcal U_s,\quad
 \mathcal N_s=\{u\in\mathcal U_s:u|_{\Sigma}=0\}
\]
を用いる。通常像を $J_n=dP(\mathcal G_{n-5})$ とする。
本節では、$H$ に対して既知だったPl\"uckerイデアルによる消去を
$\Theta_8$ に対しても示し、その消滅イデアルを正確に決める。
通常商の全ての元がこの二つの共変量から生成されるとは、まだ主張しない。

\subsection{対称行列と交代行列の束}

$T=\operatorname{Hess}_F(q(a))$、$g=dF_{q(a)}$ と置く。
交代テンソル $a$ も、$T$ と同じ型の双線形形式とみなせる。実際
\[
 \Lambda^2E^*\otimes L=\Lambda^2V\otimes\Delta^{-1}=W,
 \qquad
 \operatorname{Sym}^2E^*\otimes L=\operatorname{Sym}^2V\otimes\Delta^{-1}.
\]
従って独立なスカラー変数 $z$ に対して $T+za:E\to E^*\otimes L$
およびその行列式・余因子行列を自然に定義できる。
$\det E^*=\Delta^6$ なので
\[
 \det(T+za)\in L[z],\qquad
 \operatorname{adj}(T+za)\in(E\otimes E)[z].
\]
以後の行列表示は体積形式を選んだものであり、
このdeterminant twistを省略した別のテンソルを導入しているわけではない。

\begin{lemma}[行列束の三つの係数]
\label{lem:v17-matrix-pencil-coefficients}
普遍多項式恒等式
\begin{equation}\label{eq:v17-determinant-pencil}
 \det(T+za)=H+216z^2\Theta_8+6z^4P
\end{equation}
が成立する。さらに
\[
 B_2=[z^2]\operatorname{adj}(T+za)
 \in\operatorname{Sym}^2E
\]
は二重次数 $(\deg_F,\deg_a)=(2,6)$ の普遍共変量である。
\end{lemma}

\begin{proof}
まず $T$ を任意の対称五行五列行列、$a$ を任意の交代行列としてよい。
転置によって $z$ は $-z$ に変わるので、行列式は $z$ の偶多項式である。
定数項は $\det T$ であり、五次の係数は交代奇数行列の行列式なので零である。
対角な $T=\operatorname{diag}(t_0,\ldots,t_4)$ では、
二次の係数は
\[
 \sum_{i<j}a_{ij}^2\prod_{k\ne i,j}t_k
\]
であり、式\eqref{eq:boundary-theta-eight}の右辺の216倍と一致する。
四次の係数は $\sum_i t_iq_i^2=q^{\mathsf T}Tq$ である。
両辺は同じ合同変換則を持つ多項式であり、非退化対称行列は合同変換で対角化できる。
その集合は稠密なので、任意の $T$ に対して係数の等式が成立する。
三次式のEuler恒等式から $Tq=2g$、
$q^{\mathsf T}Tq=6F(q)=6P$ となり、表示した式を得る。

余因子行列についても
$\operatorname{adj}(T+za)^{\mathsf T}=\operatorname{adj}(T-za)$
なので偶数係数 $B_2$ は対称である。
四つの因子のうち二つが $T$、二つが $a$ であり、
$T$ の二重次数は $(1,2)$ なので、$B_2$ の二重次数は $(2,6)$ である。
\end{proof}

\begin{proposition}[$\Theta_8$ のPl\"ucker消去恒等式]
\label{prop:v17-theta-plucker-identity}
分母に $F$ や $a$ の多項式を用いず、
\begin{equation}\label{eq:v17-theta-plucker-identity}
 216\Theta_8\,q=2B_2g
\end{equation}
が普遍的に成立する。特に、普遍係数環 $R=\mathbb C[F,a]$ で
\[
 \Theta_8(q_0,\ldots,q_4)\subset(g_0,\ldots,g_4)
\]
である。
\end{proposition}

\begin{proof}
Pfaffianの恒等式 $aq=0$ と $Tq=2g$ を用いると
\[
 \det(T+za)\,q
 =\operatorname{adj}(T+za)(T+za)q
 =2\operatorname{adj}(T+za)g.
\]
$z^2$ の係数を比較すればよい。
\end{proof}

\subsection{全単元環に対する正確な消滅イデアル}

\begin{theorem}[$\Theta_8$ に作用する単元環]
\label{thm:v17-theta-unit-annihilator}
全ての $s\ge0$ について
\begin{equation}\label{eq:v17-theta-unit-annihilator}
 \Theta_8u\in J_{s+8}
 \quad\Longleftrightarrow\quad
 u\in\mathcal N_s
 \qquad(u\in\mathcal U_s).
\end{equation}
従って、通常接触商における $[\Theta_8]$ の $\mathcal U$-消滅イデアルは
正確に $\mathcal N=\bigoplus_s\mathcal N_s$ であり、
\[
 \mathcal U[\Theta_8]\simeq(\mathcal U/\mathcal N)[-8]
 \simeq\mathbb C[\overline U_4,\overline J_6][-8].
\]
ここで $[-8]$ は生成元を8次に置くという意味である。
\end{theorem}

\begin{proof}
$u\in\mathcal N_s$ とする。
補題\ref{lem:v16-plucker-prime-ideal}と有限次数での完全可約性によって、
$E^*$ 値の同変多項式 $b$ を用い
\[
 u=\langle b,q\rangle,\qquad
 (\deg_Fb,\deg_ab)=(s,s-2)
\]
と書ける。$s<2$ なら、この次数のイデアル部分は零である。
式\eqref{eq:v17-theta-plucker-identity}と $B_2$ の対称性から
\[
 \Theta_8u=\left\langle g,\frac1{108}B_2b\right\rangle.
\]
$v=B_2b/108$ は $E$ 値で、二重次数は $(s+2,s+4)$、
従って $v\in\mathcal D_{s+4}$ である。
命題\ref{prop:v12-uniform-pfaffian-surjectivity}の全射
$dq:\mathcal G_{s+3}\twoheadrightarrow\mathcal D_{s+4}$
を用いれば、$\Theta_8u\in J_{s+8}$ を得る。

逆に $\Theta_8u\in J_{s+8}$ とする。
第16版の横断切片に制限すると、補題\ref{lem:v16-contact-normal-cubic-jet}と
補題\ref{lem:v16-theta-normal-hessian}により
\[
 \overline u(f)\det(\operatorname{Hess}_f/6)
 \in\operatorname{Jac}(f)_3
\]
である。式\eqref{eq:v16-fermat-normal-minor}が与える非空開集合では
Hessian類は非零なので、$\overline u(f)=0$ となる。
多項式性によって全ての三元三次式で零であり、
階数二制限の記述から $u\in\mathcal N_s$ を得る。
環の最後の同型は定理\ref{thm:v16-unit-restriction-ring}による。
\end{proof}

\begin{remark}[この結論が決めていないこと]
$\mathcal U$ 全体による $\Theta_8$ の生成部分加群は正確に決まった。
ただし、通常商全体に別の生成元が存在しないことはこの定理から従わない。
第16版で得た $\mathbb C\cdot[H]$ は依然として別の非零方向であり、
さらに高次の未知の方向を排除するには別の議論が必要である。
\end{remark}

\subsection{独立な全係数検算}

\nolinkurl{theory/v17_structure_certificate.py} はPython標準ライブラリだけを用い、
対称 $T$ の15変数と交代 $a$ の10変数を全て独立な記号変数として計算する。
行列式の六つの係数、$aq=0$ の五成分、
$B_2Tq=216\Theta_8q$ の五成分を整数多項式の全係数で検証する。
さらに $[z^4]\operatorname{adj}(T+za)=qq^{\mathsf T}$ も25成分で照合する。
別に35個の三次係数と5個の独立変数でEuler恒等式を検証する。
これは有限標本による推測ではなく、表示した普遍恒等式そのものの係数検算である。

% END theory/v17_theta_annihilator_ja.tex


% BEGIN theory/v17_normal_projection_ja.tex
% 第17版：normal三次jetの全像・直和分解・普遍二階作用。
% 外部入力：Fisher, The Hessian of a genus one curve, Theorem 2.7(i), n=3。

\section{横断三次項の射影と通常商に残る未知の部分}
\label{sec:v17-normal-projection}

$\mathcal A=\mathcal U/\mathcal N
=\mathbb C[\overline U_4,\overline J_6]$ と置く。
第16版の単射を、本節では通常商からの射影と直和分解へ強める。
一方、横断三次項に現れない部分を明示して残すため、
この議論だけで全通常商の次元を決定したと取り違えない。

\subsection{全境界共変量の横断三次項}

$V=L_2\oplus Q_3$、$a_0=e_0\wedge e_1$ とし、
第16版の切片
\[
 a_\varepsilon=a_0+
 \varepsilon(p_4e_2\wedge e_3-p_3e_2\wedge e_4+p_2e_3\wedge e_4),
 \qquad F_f=x_0^3+x_1^3+f
\]
を使う。体積線の自明化は第16版と同じである。

\begin{lemma}[三次未満の横断項は現れない]
\label{lem:v17-covariant-normal-order}
$C\in\mathcal C_n$、$n\ge6$ とする。
$C(F,a_\varepsilon)$ の $\varepsilon^0,\varepsilon^1,\varepsilon^2$ 係数は、
全ての $F$ に対して零である。
三次係数は $F$ の $f\in\operatorname{Sym}^3Q_3$ 成分だけに依存し、
\[
 \gamma_n(C)(f,p)
 =[\varepsilon^3]C(F_f,a_\varepsilon)
\]
は $f$ の次数 $n-5$、$p$ の次数3の多項式である。
$f\mapsto\gamma_n(C)(f,-)$ は $\mathrm{SL}(Q_3)$ 共変写像
$\operatorname{Sym}^3Q_3\to\operatorname{Sym}^3Q_3$ となる。
\end{lemma}

\begin{proof}
$\lambda(t)=\operatorname{diag}(t,t,1,1,1)$ は $a_0$ を固定する。
横断変数 $b\in\Lambda^2Q_3\otimes\Delta^{-1}$ の重みは $-2$、
出力 $L=\Delta^{-3}$ の重みは $-6$ である。
従って $\varepsilon^k$ の係数に現れる $F$ の重みは $2k-6$ でなければならない。
$F$ の重みは非負なので $k<3$ は零であり、
$k=3$ では重み零の三元成分 $f$ だけが現れる。
この段階では $F_f$ に分離してから論じる必要はない。

$h\in\mathrm{SL}(Q_3)$ を $L_2$ 上の恒等写像で延長すると $a_0$ は固定され、
横断Pfaffian座標 $p$ は $h^{-\mathsf T}p$ に変わる。
従って
$\gamma_n(C)(hf,h^{-\mathsf T}p)=\gamma_n(C)(f,p)$
であり、出力三次式は $f$ と同じ表現で変換する。
\end{proof}

三元Hessianを
\[
 h_f=\frac1{216}\det\operatorname{Hess}_f
     =\det(\operatorname{Hess}_f/6)
\]
と正規化する。これは五変数の共変量 $H=H_{10}$ とは異なる。

\begin{theorem}[横断三次項が与える全次数の分裂射影]
\label{thm:v17-split-normal-projection}
各 $C\in\mathcal C_n$ は、唯一の
$A_C\in\mathcal A_{n-6}$、$B_C\in\mathcal A_{n-8}$ によって
\begin{equation}\label{eq:v17-normal-covariant-decomposition}
 \gamma_n(C)=A_C(\overline U_4,\overline J_6)\,f
            +B_C(\overline U_4,\overline J_6)\,h_f
\end{equation}
と書ける。負の次数の成分は零とする。
写像 $C\mapsto B_C$ は通常商に降り、次数付き $\mathcal U$ 加群の全射
\begin{equation}\label{eq:v17-normal-projection}
 \pi:\bigoplus_{n\ge6}\mathcal C_n/J_n
       \longrightarrow\mathcal A[-8]
\end{equation}
となる。右辺への $\mathcal U$ の作用は $\mathcal U\to\mathcal A$ を通じる。
$\overline u\mapsto[\Theta_8u]$ はこの全射の切断であり、
\begin{equation}\label{eq:v17-normal-splitting}
 \bigoplus_{n\ge6}\mathcal C_n/J_n
   =\mathcal U[\Theta_8]\oplus\mathcal K^{\mathrm{nor}},
 \qquad \mathcal K^{\mathrm{nor}}=\ker\pi
\end{equation}
という直和分解を得る。
\end{theorem}

\begin{proof}
Fisherの一次資料\cite[Definition~2.5, Lemma~2.6, Theorem~2.7(i)]{FisherHessian2012}
を $n=3$、標数零で用いる。
この定理は三元三次式の全係数空間上の多項式共変写像について、
恒等共変量 $f$ とHessian共変量が
$\mathbb C[c_4,c_6]$ 上の自由基底をなすというものである。
判別式を可逆にした環に限定した主張ではない。
Hessianを非零定数倍して $h_f$ とし、
定理\ref{thm:v16-unit-restriction-ring}の同じ不変式環
$\mathcal A=\mathbb C[\overline U_4,\overline J_6]$ を用いれば、
前補題から式\eqref{eq:v17-normal-covariant-decomposition}を得る。
$f$ と $h_f$ の係数次数は1と3なので、表示した次数も従う。

$C\in J_n$ なら、補題\ref{lem:v16-contact-normal-cubic-jet}により
$\gamma_n(C)\in\operatorname{Jac}(f)_3$ である。
$f=\frac13\sum p_i\partial_i f$ もこの空間に属する。
第16版の非空開集合では $[h_f]\ne0$ なので $B_C=0$ となり、
多項式性によって恒等的に零である。従って $\pi$ は通常商に降りる。

$u\in\mathcal U_s$ とすると、
前補題で $C$ の三次未満の横断項が消えることから
\[
 \gamma_{n+s}(uC)=\overline u\,\gamma_n(C)
\]
である。これが $\mathcal U$-線形性を示す。
補題\ref{lem:v16-theta-normal-hessian}により $\gamma_8(\Theta_8)=h_f$ なので、
$\pi([\Theta_8u])=\overline u$ となる。
定理\ref{thm:v17-theta-unit-annihilator}により切断は持上げ $u$ の選択に依存しない。
よって全射性と直和分解が従う。
\end{proof}

\begin{proposition}[未知の部分は横断四次以上で表せる]
\label{prop:v17-normal-invisible-representatives}
$\mathcal C_n^{\ge4}$ を、全 $F$ に対して上の切片で横断三次係数も零になる
$C\in\mathcal C_n$ の部分空間とする。この条件は非零階数二軌道の各点で
横断次数が少なくとも4であることと同値である。
すると
\begin{equation}\label{eq:v17-normal-invisible-representatives}
 \mathcal K_n^{\mathrm{nor}}
 =(\mathcal C_n^{\ge4}+J_n)/J_n.
\end{equation}
特に $[H]\in\mathcal K_{10}^{\mathrm{nor}}$ である。
\end{proposition}

\begin{proof}
$B_C=0$ なら $\gamma_n(C)=A_C f$ である。
$A_C$ を同じ多項式 $A_C(U_4,J_6)\in\mathcal U_{n-6}$ で持ち上げ、
$C-PA_C(U_4,J_6)$ と置き換える。
$PA_C(U_4,J_6)\in J_n$ であり、その三次係数は $A_Cf$ なので、
この代表は $\mathcal C_n^{\ge4}$ に入る。逆は定義から明らかである。
三次係数が $f$ だけに依存することと非零階数二軌道の推移性により、
全 $F$・軌道上の全点で同じ消失を得る。
$H$ は $q$ の五次式なので、錐に沿う次数は少なくとも5である。
\end{proof}

\begin{remark}[全商の生成問題の正確な残り]
第16版の単射は全次数の直和因子へ強まったが、
$\mathcal K^{\mathrm{nor}}=\mathbb C\cdot[H]$ はまだ証明していない。
これが全通常商の生成予想の残る主張である。
例えば24次では
\[
 \mathcal C_{24}/J_{24}
 =\mathbb C\cdot[\Theta_8U_4^4]\oplus
   \mathbb C\cdot[\Theta_8U_4J_6^2]\oplus\mathcal K_{24}^{\mathrm{nor}}.
\]
従って上界2は $\mathcal K_{24}^{\mathrm{nor}}=0$ と同値であり、
三元共変量の自由生成定理だけではこの上界は得られない。
本版では別の明示階数計算でこの消失を確定する
（命題\ref{prop:v17-ordinary-contact-rank24}）。
本節の $\mathcal K^{\mathrm{nor}}$ はKoszul複体の項
$\mathcal K_{k,n}$ やそのhomologyとは別の記号である。
\end{remark}

\subsection{普遍な場の横断一次作用はスカラーである}

\begin{lemma}[三元三次式からのEnd共変量]
\label{lem:v17-ternary-endomorphism-scalar}
三元三次式の係数に多項式的な
$\mathrm{SL}(Q_3)$ 同変写像
$M:\operatorname{Sym}^3Q_3\to\operatorname{End}(Q_3^*)$
は、全ての $f$ に対してスカラー行列である。
同変写像が $\mathrm{GL}(Q_3)$ のdeterminant指標でねじれていても同じ結論となる。
\end{lemma}

\begin{proof}
まず $M$ を $f$ の次数 $d$ の斉次成分 $M_d$ に分ける。
各成分も $\mathrm{SL}_3$ 同変であり、$h=\lambda h_0$、
$h_0\in\mathrm{SL}_3$ と書くと
\[
 M_d(hf)=(\det h)^d h^{-\mathsf T}M_d(f)h^{\mathsf T}.
\]
従って各斉次成分は自動的にdeterminant指標でねじれた
$\mathrm{GL}_3$ 共変写像へ延びる。以下は各 $M_d$ に適用してから和を取ればよい。
$f_t=x^3+y^3+z^3+6txyz$ とし、$\omega$ を原始三乗根とする。
$D=\operatorname{diag}(1,\omega,\omega^2)$ と三つの座標の巡回置換 $C$
は $\mathrm{SL}_3$ に属し、全ての $f_t$ を固定する。
$Q_3^*$ 上で $D$ は異なる三固有値を持つため、これと可換な行列は対角である。
さらに $C$ と可換なら三対角成分は等しい。
従って $M(f_t)$ はスカラーである。

写像
\[
 \mathrm{GL}_3\times\mathbb A^1\longrightarrow\operatorname{Sym}^3Q_3,
 \qquad(h,t)\longmapsto hf_t
\]
は優勢である。実際 $(1,0)$ での九つの行列方向の微分は
$3p_i^2p_j$（$i,j=1,2,3$）、最後の方向の微分は $6p_1p_2p_3$ である。
この十個は三次単項式の基底であり、対応する行列式は
$6\cdot3^9=118098\ne0$ である。
同変性とdeterminantによるスカラー倍は行列のスカラー性を保つ。
従って稠密な集合で $M$ はスカラーであり、多項式性により全 $f$ でそうなる。
\end{proof}

\begin{proposition}[Pfaffian方向の正確な横断一次項]
\label{prop:v17-scalar-normal-pfaffian-action}
$w\in\mathcal D_{r+1}$ とする。全 $F$ に対して
$w=Mq$ と同変に表示でき、
分離切片上の $Q_3^*$ 成分は
\begin{equation}\label{eq:v17-scalar-normal-pfaffian-action}
 w_Q(F_f,a_\varepsilon)
  =\varepsilon\alpha_w(f)p+O(\varepsilon^2)
\end{equation}
となる。$\alpha_w$ は $f$ に多項式的である。
この一次作用は $M$ の非一意性によらず、$w$ から定まる。
\end{proposition}

\begin{proof}
$w=Mq$ は補題\ref{lem:v16-plucker-vanishing-D}による。
$\lambda(t)$ に対する $Q_3^*\to Q_3^*$ ブロックの重みは零なので、
$M_{QQ}(F,a_0)$ は $F$ の三元成分 $f$ だけに依存する。
$\mathrm{SL}(Q_3)$ ではこれは共役作用に関して同変である。
$h\in\mathrm{GL}(Q_3)$ まで広げると、
$a_0\mapsto(\det h)^{-1}a_0$ と $a$ 次数 $r-1$ による
因子 $(\det h)^{r-1}$ が付くが、それは前補題で許したdeterminant指標である。
従って $M_{QQ}(f)=\alpha_w(f)\operatorname{id}$ となる。

切片では $q=(0,0,\varepsilon p)$ が正確に成立するので表示した一次項を得る。
二つの $M$ が同じ $w$ を表すなら、その $QQ$ ブロックの差は
全 $p$ を零に送るため零である。
これは $dq_{a_0}$ が三次元の横断空間から $Q_3^*$ への同型を与えることでも説明できる。
\end{proof}

\begin{theorem}[全二階接触補正は横断射影で消える]
\label{thm:v17-quadratic-actions-normal-zero}
$r,s\ge1$、$Y\in\mathcal G_r$、$Z\in\mathcal G_s$、$n=r+s+4$ とすると
\begin{equation}\label{eq:v17-quadratic-actions-normal-zero}
 \pi_n([P''(Y,Z)])=0.
\end{equation}
従って $P''(Y,Z)$ の通常商における類は
$\mathcal K_n^{\mathrm{nor}}$ に入る。
特に $\mathcal K_n^{\mathrm{nor}}=0$ と別に分かっている次数では、
全てのこの二階補正が $J_n$ に入る。
\end{theorem}

\begin{proof}
連鎖律により
\[
 P''(Y,Z)=\operatorname{Hess}_F(q)(dqY,dqZ)
              +dF_q\bigl(d^2q(Y,Z)\bigr).
\]
$d^2q(Y,Z)$ は $E$ 値で、二重次数は
$(r+s-2,r+s)$、従って $\mathcal D_{r+s}$ に属する。
全次数全射 $dq:\mathcal G_{r+s-1}\twoheadrightarrow\mathcal D_{r+s}$
から、第二項は全体として $J_{r+s+4}$ に入る。
単なる横断評価上の零性に置き換えているのではない。

第一項を分離切片に制限する。
$\operatorname{Hess}_{F_f}(q)$ の最初の二行・二列は零であり、
残りのブロックは $\varepsilon\operatorname{Hess}_f(p)$ である。
前命題からその横断三次係数は
\[
 \alpha_{dqY}(f)\alpha_{dqZ}(f)\,
       p^{\mathsf T}\operatorname{Hess}_f(p)p
 =6\alpha_{dqY}(f)\alpha_{dqZ}(f)f(p)
\]
となる。従ってHessian共変量の係数は零である。
\end{proof}

\begin{corollary}[二十四次の全二階補正の通常像所属]
\label{cor:v17-twentyfourth-corrections-conditional}
二十四次の四つの二階補正族について
\[
 P''(\mathcal G_7,\mathcal G_{13}),\quad
 P''(\mathcal G_8,\mathcal G_{12}),\quad
 P''(\mathcal G_9,\mathcal G_{11}),\quad
 P''(\mathcal G_{10},\mathcal G_{10})
 \ \subset J_{24}.
\]
これは四族の代表を個別の評価点で全て消すことを仮定しない。
\end{corollary}

\begin{proof}
命題\ref{prop:v17-ordinary-contact-rank24}により、
通常二十四次商は $[\Theta_8U_4^4],[\Theta_8U_4J_6^2]$ の二つで張られる。
従って $\mathcal K_{24}^{\mathrm{nor}}=0$ であり、
各組は $r+s+4=24$ を満たすので前定理を適用する。
この次数の核の消失は、後続節の明示像の独立性と
既存のnormal-slice単射を合わせた結果であり、
本節の定性的な分裂だけから推論したものではない。
\end{proof}

\begin{remark}[証明と有限監査の範囲]
Heisenberg二生成元の共通commutantの方程式を
$\mathbb Q(\omega)$ から有理数の方程式へ展開すると階数8であり、
核は単位行列の一次元空間となる。
\nolinkurl{theory/v17_structure_certificate.py} はこの有理行列と
優勢性の十行十列行列・行列式118098を独立に記録する。
三元共変量加群の全生成は有限標本から推測せず、明示したFisherの定理を入力とする。
全通常商の上界、$\mathcal K^{\mathrm{nor}}$ の消失、
実際のTate係数の消滅をこの監査だけから結論してはいない。
\end{remark}

% END theory/v17_normal_projection_ja.tex


% BEGIN theory/v17_discriminant_torsion_ja.tex
% 第17版：rank-four support、Schur補チャート、normal核の判別式torsion。
% 固定したsmooth Fの全Jacobian商や、原点での正則な単元変更とは区別する。

\section{横断射影の核は三元判別式上のtorsionである}
\label{sec:v17-discriminant-torsion}

本節では、前節で残した $\mathcal K^{\mathrm{nor}}$ の支持をさらに限定する。
普遍係数環を $R=\mathbb C[F,a]$、二つのイデアルを
\[
 I=(q_0,\ldots,q_4),\qquad G=(g_0,\ldots,g_4),\quad
 g_i=(\partial_iF)(q(a))
\]
と書く。通常商は有限同変成分を取った $R/G$ と一致する。
実際、$C\in\mathcal C_n$ が $G$ に入るなら、二重次数をそろえ、
完全可約性により同変な係数 $w\in\mathcal D_{n-4}$ を選べる。
全次数 $dq$ 全射から $C=dF(w)\in J_n$ を得る。
逆包含は定義から従う。

\subsection{通常共変量の類はPl\"ucker錐に支持される}

\begin{lemma}[階数四開集合での通常像所属]
\label{lem:v17-rank-four-support}
全ての $C\in\mathcal C_n$ と各 $i$ について
\[
 C\in G R_{q_i}.
\]
従って、$R/G$ における $C$ の類は $I$-torsionである。
すなわち、$C$ に依存する整数 $m$ が存在して
$I^mC\subset G$ となる。
\end{lemma}

\begin{proof}
階数四の点
$a_*=e_0\wedge e_1+e_2\wedge e_3$ を取る。
一パラメータ部分群
\[
 \mu(t)=\operatorname{diag}(t,t,t,t,t^{-2})
\]
では $\det\mu=t^2$ であり、$W=\Lambda^2V\otimes\Delta^{-1}$ の
$a_*$ は固定される。$q(a_*)=e_4^*$（体積線を自明化した表示）である。
$F$ のテンソル単項式の重みは、添字4の個数を $k$ として $3-3k$ となる。
$g(F,a_*)=0$ は $f_{i44}=0$（$i<4$）および $f_{444}=0$ を意味し、
負の重み $-3,-6$ の全成分を消す。
残る $F$ の重みは3と0であるのに対し、出力 $L$ の重みは $-6$ である。
従って多項式同変性により $C(F,a_*)=0$ となる。
階数四の点は一つの $\mathrm{GL}(V)$ 軌道なので、
$C$ は $q\ne0,\ g=0$ の全点で零である。

$q_i\ne0$ の開集合では、$F\mapsto g(F,a)$ は五成分全射である。
これらは $F$ の係数に線形な方程式なので、零点スキームは
$a$ の開集合上のベクトル束の全空間であり、被約である。
例えば対称テンソル係数 $f_{iii},f_{jii}$（$j\ne i$）を選べば、
対応する五行五列係数行列は、非零定数倍の $q_i^{10}$ を行列式に持つ。
従って零点上の消失から $C\in GR_{q_i}$ が従う。
各 $i$ についてある $m_i$ が $q_i^{m_i}C\in G$ を満たす。
$m=\sum_i(m_i-1)+1$ とすれば、$I^m$ の各単項式に
ある $q_i^{m_i}$ が因数として現れ、最後の主張を得る。
\end{proof}

\begin{remark}[支持を意味する範囲]
全ての $R/G$ が $I$-torsionだとは述べていない。
例えば $q\ne0$ にも $g=0$ の普遍incidenceは存在する。
上の結論は、指定した値域 $L$ の普遍同変成分を取ったときの消失である。
また $I$-torsion性だけから $IC\subset G$ や $C\in J_n$ は従わない。
\end{remark}

\subsection{Schur補チャートではPfaffianが正確に線形になる}

$x=a_{01}$ を可逆にし、$i,j\in\{2,3,4\}$ について
\[
 l_i=a_{0i},\qquad m_i=a_{1i},\qquad
 b_{ij}=a_{ij}-\frac{l_im_j-l_jm_i}{x}
\]
と置く。これらは $R_x$ の座標変換であり、
$I R_x=(b_{23},b_{24},b_{34})$ となる。
$p=(b_{34},-b_{24},b_{23})$ とすれば
\begin{equation}\label{eq:v17-schur-linear-pfaffian}
 q_0=\sum_{i=2}^4m_ip_i,\qquad
 q_1=-\sum_{i=2}^4l_ip_i,\qquad
 q_i=xp_i\quad(i=2,3,4).
\end{equation}
これは初期項ではなく正確な恒等式である。

行列 $h$ を
\[
 he_0=e_0-\sum_{i=2}^4\frac{m_i}{x}e_i,\qquad
 he_1=e_1+\sum_{i=2}^4\frac{l_i}{x}e_i,\qquad
 he_i=e_i\quad(i=2,3,4)
\]
で定義すると $\det h=1$ であり、
$a=h(xe_0\wedge e_1+\sum b_{ij}e_i\wedge e_j)$ である。
$F'=h^{-1}F$ と置き、その三元部分を $f=F'|_{Q_3}$ とすると
\begin{equation}\label{eq:v17-schur-exact-ternary}
 P(F,a)=x^3f(p),\qquad
 x^3\partial_i f(p)=m_ig_0-l_ig_1+xg_i.
\end{equation}
従って、三つの三元偏微分は $G R_x$ に入る。

\begin{lemma}[滑らかな三元三次式では横断四次以上を消せる]
\label{lem:v17-smooth-ternary-fourth-power}
上のチャートの係数環で、三元式 $f$ の判別式を可逆にすると
\begin{equation}\label{eq:v17-normal-fourth-in-gradient}
 (b_{23},b_{24},b_{34})^4\subset G
\end{equation}
が成立する。
\end{lemma}

\begin{proof}
まず固定した滑らかな三元三次式 $f$ を考える。
Euler恒等式から、その三つの偏微分の射影共通零点は曲線の特異点に等しく、
したがって存在しない。アフィン共通零点は原点だけなので、
$\mathbb C[p_2,p_3,p_4]$ で三つの二次式はheight3の正則列となる。
ここでも斉次極大での局所判定と斉次性を用いて大域へ戻せる。
三つの二次式のKoszul自由分解からHilbert級数は
\[
 \frac{(1-t^2)^3}{(1-t)^3}=(1+t)^3
\]
である。従って四次成分は零で、全ての四次単項式が三偏微分のイデアルに入る。

次に $F'$ と $x^{\pm1},l_i,m_i$ を変数とする係数環を $B$ とする。
$R_x=B[p_2,p_3,p_4]$ であり、上の三元式の係数も $B$ に属する。
判別式を可逆にした $B$ 上で、三偏微分と二次単項式の積による写像
\[
 B^{18}\longrightarrow B^{15}
\]
を考える。前段の議論から、この写像は各閉点の剰余体上で全射である。
有限生成の余核が非零ならその支持は閉点を持ち、
Nakayamaの補題に反する。従って係数環上でも全射である。
全ての高い単項式は四次単項式の倍数なので、
式\eqref{eq:v17-schur-exact-ternary}から主張を得る。
無限回の $I$-進近似や完備化による持上げは用いていない。
\end{proof}

\subsection{未知の通常商成分の判別式torsion}

\begin{lemma}[正規化済み単元からの三元判別式]
\label{lem:v17-unit-discriminant}
\[
 \mathscr D=J_6^2-2304U_4^3\in\mathcal U_{12}
\]
と置く。$\overline{\mathscr D}\in\mathcal A_{12}$ は三元三次式の通常の判別式の
非零定数倍である。特に $\overline{\mathscr D}(f)\ne0$ と
$f$ が滑らかであることは同値である。
\end{lemma}

\begin{proof}
Fisherの不変式定理\cite[Theorem~4.4]{Fisher2008}では、
三元判別式は次数12であり、その非消失は滑らかさと同値である。
既知の制限環 $\mathcal A$ の12次成分は
$\overline U_4^3,\overline J_6^2$ の二つで張られる。
第14版のHesse制限では
\[
 \overline U_4(f_t)=t^4-t,\qquad
 \overline J_6(f_t)=6-120t^3-48t^6.
\]
$t=0$ は滑らかであり、$t=-1/2$ の
$x^3+y^3+z^3-3xyz$ は三直線の積である。
後者での値は $U_4=9/16$、$J_6=81/4$ で、$J_6^2/U_4^3=2304$ となる。
従って判別式の二つの係数の比が決まり、
滑らかな $t=0$ での非消失により定数倍は非零である。
係数の追加検算として
\[
 \overline{\mathscr D}(f_t)=36(1+8t^3)^3
\]
も成り立つ。
\end{proof}

\begin{theorem}[横断射影は判別式可逆化後に同型となる]
\label{thm:v17-normal-kernel-discriminant-torsion}
$\mathcal K^{\mathrm{nor}}=\ker\pi$ の任意の元 $\xi$ に対して、
ある整数 $m=m(\xi)\ge0$ が存在して
\begin{equation}\label{eq:v17-normal-kernel-discriminant-torsion}
 \mathscr D^m\xi=0
\end{equation}
となる。従って
\begin{equation}\label{eq:v17-discriminant-localized-quotient}
 \left(\bigoplus_{n\ge6}\mathcal C_n/J_n\right)[\mathscr D^{-1}]
 \simeq
 \mathcal A[\overline{\mathscr D}^{-1}]\,[\Theta_8]
\end{equation}
である。この局所化は次数を保ち、$\mathscr D^{-1}$ の次数を $-12$ とする。
\end{theorem}

\begin{proof}
元は有限個の斉次成分の和なので、各成分について示せば十分である。
$\xi$ を横断四次以上の $C\in\mathcal C_n^{\ge4}$ で表す。
補題\ref{lem:v17-rank-four-support}から、その $R/G$ における類は $I$-torsionである。
従って、$I$ を含まない素イデアルで局所化すれば零になる。

$I$ を含み $\mathscr D$ を含まない素イデアル $\mathfrak p$ を考える。
$\mathscr D$ の $a$ 次数は12なので、$\mathfrak p$ は全ての $a_{ij}$ を含むことはできない。
ある $a_{ij}$ を可逆にしたSchur補チャートで論じられる。
添字を替えて $x=a_{01}$ が可逆であるとしてよい。
錐上の点での同変性と $a$ 次数により
\[
 \mathscr D(F,a)\equiv x^{12}\overline{\mathscr D}(f)\pmod I.
\]
従って $\overline{\mathscr D}(f)$ は $R_{\mathfrak p}$ で可逆である。

横断四次以上という条件は、このチャートで
$C\in(b_{23},b_{24},b_{34})^4$ を意味する。
実際、$\det h=1$ と $a$ に関する斉次性により
\[
 C(F,a)=C(F',xe_0\wedge e_1+b)
       =x^n C(F',e_0\wedge e_1+b/x)
\]
であり、定義から四次未満の $b$ 係数は全て零である。
補題\ref{lem:v17-smooth-ternary-fourth-power}を適用すると、
$C\in GR_{\mathfrak p}$ となる。
これで $\mathscr D$ を可逆にした環の全素点で $C$ の類が零となり、
$C\in GR_{\mathscr D}$ を得る。
局所化の定義からある $m$ に対して $\mathscr D^mC\in G$ である。
同変な有限次数成分を取り、全次数 $dq$ 全射を使えば
$\mathscr D^mC\in J_{n+12m}$ となる。
直和分解\eqref{eq:v17-normal-splitting}を局所化して最後の同型を得る。
\end{proof}

\begin{remark}[局所化で見えなくなる部分は依然として未決定]
この定理は $\mathcal K^{\mathrm{nor}}=0$ を意味しない。
既知の $[H]$ も $\mathscr D[H]=0$ を満たす非零元である。
ほかの判別式torsionが高次に現れないかは、引き続き全商の生成問題として残る。
また $\mathscr D$ は原点 $a=0$ で零なので、
この局所化を原点で許される定数項1の単元変更と解釈してはならない。
$\mathscr D^mR_n$ が通常像に入ることだけでは、実際の $R_n$ を同じ次数で消去できない。
本定理は全ての次数で使える構造上の制限であり、24次商の上界2の代わりではない。
\end{remark}

% END theory/v17_discriminant_torsion_ja.tex


% BEGIN theory/v17_vertex_torsion_ja.tex
% 第17版：Cech--Koszulによるnormal核のvertex支持と有限次元性。
% 有限補正の全次元・終端はまだ計算しない。

\section{通常商の未知部分は頂点に支持される有限次元加群である}
\label{sec:v17-vertex-finiteness}

前節では未知の部分 $\mathcal K^{\mathrm{nor}}$ が判別式torsionであることを示した。
さらに普遍勾配のKoszul複体と横断変数の\v{C}ech複体を同時に使うと、
この部分は $a=0$ の頂点に支持され、有限次元であることまで分かる。
したがって十分高い次数では通常商は $\Theta_8$ の自由部分加群と一致する。
ただし、その「十分高い」次数の明示的な上限や、有限補正の全体は本節では決定しない。

\subsection{横断三変数での局所コホモロジーの次数上界}

\begin{lemma}[\v{C}ech--Koszulの零次端]
\label{lem:v17-cech-koszul-normal-degree-bound}
$B$ をNoether環、$S=B[p_1,p_2,p_3]$、
$\mathfrak m=(p_1,p_2,p_3)$ とする。
$g_1,\ldots,g_5$ は $p$ に関して二次斉次とし、
$K_\bullet=K(g_1,\ldots,g_5;S)$ と置く。
正次数Koszulホモロジー $H_i(K_\bullet)$ が全て
$\mathfrak m$-torsionであると仮定する。
このとき
\begin{equation}\label{eq:v17-local-torsion-degree-bound}
 \bigl(H^0_{\mathfrak m}(S/(g_1,\ldots,g_5))\bigr)_d=0
 \qquad(d>3)
\end{equation}
である。下付き $d$ は $p$ の次数であり、係数環 $B$ の次数ではない。
\end{lemma}

\begin{proof}
拡大\v{C}ech複体を $\check C^\bullet_{\mathfrak m}$ と書き、
$K_i$ をコホモロジー次数 $-i$ に置いて
\[
 \operatorname{Tot}
  \bigl(\check C^\bullet_{\mathfrak m}\otimes_S K_\bullet\bigr)
\]
を考える。両方向に有限な双複体であり、二つのスペクトル系列は収束する
\cite[Tag 012X, Lemmas 12.25.1 and 12.25.3]{StacksCechKoszul}.
局所化は平坦なので、先にKoszulホモロジーを取った系列の第二項は
\[
 E_2^{u,v}=H^u_{\mathfrak m}(H_{-v}(K_\bullet)),
 \qquad u\ge0,\quad v\le0,
\]
総次数は $u+v$ である。
$i>0$ の $H_i(K_\bullet)$ は $\mathfrak m$-torsionなので、
その\v{C}ech複体は次数0以外で零となる。
従って総次数0の非零になり得る項は
\[
 E_2^{0,0}=H^0_{\mathfrak m}(S/(g))
\]
だけである。$r\ge2$ の微分の出先は $(r,1-r)$ であり、
正次数Koszulホモロジーの正次数\v{C}echコホモロジーだから零である。
入って来る項は $u=-r<0$ なので存在しない。
よって総複体の零次コホモロジーはこの $E_2^{0,0}$ と同型である。

逆に、先に\v{C}echコホモロジーを取る。
$p_1,p_2,p_3$ は多項式環の正則変数なので、自由 $S$ 加群の
\v{C}echコホモロジーは次数3以外で零である
\cite[Tag 0G6L]{StacksCechKoszul}。
ここではさらに
\[
 H^3_{\mathfrak m}(S)
  =\bigoplus_{\alpha_1,\alpha_2,\alpha_3\ge1}
       B\,p_1^{-\alpha_1}p_2^{-\alpha_2}p_3^{-\alpha_3}
\]
と明示でき、その最高 $p$ 次数は $-3$ である。
従って同じ総複体の零次コホモロジーは
\[
 H_3\bigl(H^3_{\mathfrak m}(K_\bullet)\bigr)
\]
である。$K_3=S(-6)^{10}$ なので、その第三項
$H^3_{\mathfrak m}(K_3)$ の最高 $p$ 次数は $6-3=3$ である。
ホモロジーはこの項の部分商だから、次数 $d>3$ では零となる。
両順序の計算を比較して結論を得る。
\end{proof}

\subsection{普遍勾配への適用}

\begin{theorem}[横断射影の核は頂点に支持される]
\label{thm:v17-normal-kernel-vertex-support}
$R=\mathbb C[F,a]$、$G=(g_0,\ldots,g_4)$、
$\mathfrak a=(a_{ij}:i<j)$ とする。
任意の $\xi\in\mathcal K^{\mathrm{nor}}$ は $R/G$ において
$\mathfrak a$-torsionである。
すなわち、$a=0$ の頂点を除く全ての点でその類は零である。
\end{theorem}

\begin{proof}
斉次な $\xi$ を、命題\ref{prop:v17-normal-invisible-representatives}による
$C\in\mathcal C_n^{\ge4}$ で表す。一般の元は有限個の斉次成分の和として扱えばよい。
補題\ref{lem:v17-rank-four-support}により、その類は $I$-torsionである。
各 $a_{ij}$ を可逆にしたSchur補チャートを考え、
添字を替えて $x=a_{01}$ の場合で示す。

前節の変数変更により $R_x=B[p_2,p_3,p_4]$ と書ける。
五つの勾配の各成分は $p$ の二次斉次式である。
$p\ne0$ では $q\ne0$ であり、$F$ から五つの勾配への写像は全射である。
従って五つの勾配はこの開集合上で、$F$ の五つの独立な線形方程式として
正則列をなす。よって正次数Koszulホモロジーは
$\mathfrak m=(p_2,p_3,p_4)$ に支持され、有限生成性から
$\mathfrak m$-torsionである。
ここでは五つの勾配が $p=0$ でも正則列であるとは仮定していない。

$C$ の局所化は $\mathfrak m$-torsionであり、
\[
 C(F,a)=x^nC(F',e_0\wedge e_1+b/x)
\]
により $p$ の次数が4以上の項だけからなる。
勾配イデアルはこの次数に関して斉次なので、
各斉次成分も $\mathfrak m$-torsionである。
前補題によりその全てが零であり、$C\in GR_x$ となる。
全ての $a_{ij}$ チャートについて同じことが成り立つから、
頂点以外の全開集合で $\xi$ は零である。
局所化の定義から、ある $\ell$ で $\mathfrak a^\ell\xi=0$ となる。
\end{proof}

\begin{theorem}[未知部分の有限次元性と十分高い次数の通常商]
\label{thm:v17-finite-normal-kernel}
$\mathcal K^{\mathrm{nor}}$ は有限次元の次数付き複素ベクトル空間である。
従って、ある整数 $n_0$ が存在し、全ての $n\ge n_0$ について
\begin{equation}\label{eq:v17-eventual-contact-quotient}
 \mathcal C_n/J_n
  =\bigoplus_{\substack{a,b\ge0\\4a+6b=n-8}}
      \mathbb C[\Theta_8U_4^aJ_6^b]
\end{equation}
となる。特にHilbert級数は
\begin{equation}\label{eq:v17-contact-hilbert-finite-correction}
 \sum_{n\ge6}\dim(\mathcal C_n/J_n)t^n
   =\frac{t^8}{(1-t^4)(1-t^6)}+E_{\mathrm{vert}}(t),
 \qquad E_{\mathrm{vert}}(t)\in\mathbb Z_{\ge0}[t]
\end{equation}
と書ける。$E_{\mathrm{vert}}(t)$ は有限な多項式で、
既知の十次類によりその $t^{10}$ 係数は1である。
$E_{\mathrm{vert}}(t)=t^{10}$ という全生成予想は、なお未決定である。
\end{theorem}

\begin{proof}
$T_0=H^0_{\mathfrak a}(R/G)$ は、Noether加群 $R/G$ の部分加群なので有限生成である。
従って有限個の二重斉次生成元 $v_1,\ldots,v_e$ と整数 $\ell$ を選び、
\[
 \mathfrak a^\ell T_0=0
\]
とできる。$v_j$ の $a$ 次数を $b_j$ とすると、
$T_0$ の $a$ 次数は $\max_j b_j+\ell-1$ 以下に限られる。
実際、任意の元を生成元の $F,a$ 多項式倍で表し、
$a$ 次数が $\ell$ 以上の係数単項式を除けばよい。

前定理から $\mathcal K_n^{\mathrm{nor}}$ は $T_0$ の指定された同変成分に入り、
その二重次数は $(n-5,n)$ である。
従って有限個の $n$ についてしか非零にならない。
各 $n$ では $F$ 次数も固定されているので、その空間は有限次元である。
ここで $T_0$ 全体が複素ベクトル空間として有限次元であるとは述べていない。
必要なのは $a$ 次数の有界性と、この対角的な二重次数条件の組合せである。

分裂\eqref{eq:v17-normal-splitting}から十分高い次数の等式とHilbert級数を得る。
十次では $\mathcal A_2=0$ であり、
定理\ref{thm:boundary-eleven-jet}の通常商が既知の $[H]$ 一方向なので、
$E_{\mathrm{vert}}$ の該当係数は1である。
\end{proof}

\begin{corollary}[未知部分は正次数単元についての有限torsion]
\label{cor:v17-positive-unit-torsion}
$\mathcal U_+=\bigoplus_{s>0}\mathcal U_s$ とすると、
通常商の $\mathcal U_+$-torsion部分は正確に $\mathcal K^{\mathrm{nor}}$ である。
ここでtorsionとは、ある $k$ に対して $(\mathcal U_+)^k\xi=0$ となるという意味である。
単一の正次数単元による消失とは区別する。
その全体を消す $\mathcal U_+$ の冪が存在するが、この冪の最小値は未計算である。
\end{corollary}

\begin{proof}
有限個の次数しか持たない $\mathcal K^{\mathrm{nor}}$ は、
十分大きい正次数単元の積で消える。
一方、$\Theta_8\mathcal A$ では $\overline U_4$ のどの冪も非零な元を消さない。
$\mathcal A$ が整域であることと分裂から結論を得る。
\end{proof}

\begin{remark}[有効な有限計算の範囲はまだ別に必要である]
本節は、十分高い次数での通常商を定性的に決める。
しかし、閾値 $n_0$、$E_{\mathrm{vert}}$ の最高次数と係数、
および $\mathcal K^{\mathrm{nor}}=\mathbb C\cdot[H]$ かどうかは決めていない。
従って、個別の24次上界を本節だけから得たとは述べない。
また実際のTate係数が自由部分へ零に写ることも、別の計算または幾何学的論証を要する。
\end{remark}

% END theory/v17_vertex_torsion_ja.tex


% BEGIN characters/v17_character_notes_ja.tex
\section{二十四次の普遍共変量とKoszul各項}
\label{sec:v17-characters}

二十四次では、normal sliceから得られる二つの商の元を出発点とし、
普遍空間の次元と接触写像の階数を別々に求める。
本節の計算は整数の重み指標、Weyl交代和、Pieri則だけに依存する。

\subsection{次数とdeterminant twist}

$\Delta=\det V$、$A=\Lambda^2V$、$W=A\otimes\Delta^{-1}$ とし、
$F\in\operatorname{Sym}^3V$、$a\in W$、
$q(a)=\tfrac12a\wedge a\in V^*\otimes\Delta^{-1}$、
$P=F(q(a))\in\Delta^{-3}$ という規約を保つ。
中心元の重みから、$a$ 次数 $n$ の境界項は $F$ 次数 $n-5$、
$a$ 次数 $r$ の座標補正は $F$ 次数 $r-1$、
$a$ 次数 $s$ の単元補正は $F$ 次数 $s$ を持つ。
従って二十四次のKoszul第 $k$ 項は
\[
 K_{k,24}=\operatorname{Hom}_{\mathrm{GL}(V)}
 \left(\operatorname{Sym}^{19-k}(\operatorname{Sym}^3V)
 \otimes\operatorname{Sym}^{24-4k}A,
 \Lambda^kV^*\otimes\Delta^{21-2k}\right),\qquad 0\leq k\leq5.
\]
$K_{0,24}=C_{24}$、$K_{1,24}=D_{20}$、$K_{2,24}=B_{24}$ と書く。
第5項にも $a$ 次数4が残るので、終端の多重度は前の次数から引き継がない。

\begin{proposition}[二十四次の厳密次元]
\label{prop:v17-character-dimensions}
各普遍空間の次数・標的分割・次元は次の通りである。
\[
\begin{array}{c|r|r|c|r}
\text{空間}&\deg_F&\deg_a&\text{標的分割}&\dim\\\hline
C_{24}&19&24&(21,21,21,21,21)&2458\\
G_{19}&18&19&(19,19,18,18,18)&11676\\
U_{18}&18&18&(18,18,18,18,18)&1365\\
D_{20}&18&20&(19,19,19,19,18)&6262\\
B_{24}&17&16&(17,17,17,16,16)&5615\\
K_{3,24}&16&12&(15,15,14,14,14)&2079\\
K_{4,24}&15&8&(13,12,12,12,12)&282\\
K_{5,24}&14&4&(10,10,10,10,10)&11
\end{array}
\]
補助的な十七次の空間についても
\[
 \dim C_{17}=31,\qquad \dim G_{12}=155,\qquad
 \dim U_{11}=9,\qquad \dim D_{13}=92
\]
を得る。これらの値は接触写像の階数を仮定しない。
\end{proposition}

\begin{proof}
$\operatorname{Sym}^m(\operatorname{Sym}^3V)$ の指標を35個の三次単項式重み
$\alpha$ に対するEuler積
\[
 \prod_{|\alpha|=3}(1-tx^\alpha)^{-1}
\]
の $t^m$ 係数として計算する。
$\rho=(4,3,2,1,0)$ とするとSchur多重度は
\[
 \sum_{\sigma\in S_5}\operatorname{sgn}(\sigma)
 [x^{\lambda+\rho-\sigma\rho}]
 \operatorname{ch}\operatorname{Sym}^m(\operatorname{Sym}^3V)
\]
である。外積二乗側も10個の重みから指標を作り、
テンソル積の交代和を直接評価した。
さらにLittlewood分解
\[
 \operatorname{Sym}^{d}(\Lambda^2V)
 =\bigoplus_{j=0}^{\lfloor d/2\rfloor}S_{(d-j,d-j,j,j,0)}V
\]
の多重度1とWeyl次元和を照合した。

$C_{24}$ で必要な $F$ 側の分割は
$(21,18-j,18-j,j,j)$、$0\leq j\leq9$ であり、多重度は
\[
 (4,5,42,108,369,521,765,425,213,6).
\]
和は2458である。
$U_{18}$ の分割は $(18,18-j,18-j,j,j)$、$0\leq j\leq9$ で、多重度は
\[
 (2,1,16,44,179,275,447,261,136,4).
\]
和は1365である。

$G_{19}$ と $D_{20}$ はそれぞれ外積第3冪・第1冪を用いる垂直帯Pieri計算、
$K_{k,24}$ は外積第 $k$ 冪を用いるPieri計算でも計算し、
直接のテンソル積指標と照合した。
Koszul源の非零Pieri成分数は $k=2,3,4,5$ の順に29、24、13、3である。
第5項の $F$ 側分割と多重度は
\[
 (10,10,10,6,6):7,\qquad
 (10,9,9,7,7):2,\qquad
 (10,8,8,8,8):2
\]
となる。十七次の補助四空間も同じ二通りの計算で一致する。

全ての $F$ 次数 $1\leq m\leq19$ と $a$ 次数 $1\leq d\leq24$ について、
独立のNewton--Adams再帰
\[
 m h_m=\sum_{k=1}^{m}p_kh_{m-k}
\]
で全重み係数を再計算し、Euler積から得た辞書全体と一致することを確認した。
必要な $F$ 次数 $11,12,14,15,16,17,18,19$ の全Schur分解についても、
多重度とWeyl次元の積の和が $\binom{m+34}{m}$ に一致する。
この検証は浮動小数点階数を使わない有限の整数恒等式である。
\end{proof}

\subsection{Euler交代和と高次Koszul微分}

交代和は
\[
 2458-6262+5615-2079+282-11=3.
\]
前版の高次ホモロジー消失定理
\ref{thm:v16-koszul-high-homology-vanishing}を適用すると
\[
 \dim H_0-\dim H_1+\dim H_2=3,
 \qquad H_i=0\quad(i>2).
\]
従って上位三つの微分については
\[
 \operatorname{rank}\delta_5=11,\qquad
 \operatorname{rank}\delta_4=282-11=271,\qquad
 \operatorname{rank}\delta_3=2079-271=1808
\]
が確定する。一方、$H_1,H_2$ の消失はこの計算からは従わない。
通常商の次元は別の接触像の計算で求める必要がある。
Pfaffian微分の全次数全射性を用いれば
\[
 \dim\ker(\delta q:G_{19}\longrightarrow D_{20})=11676-6262=5414
\]
を得るが、$dF:D_{20}\to C_{24}$ の階数はこの差からは決まらない。

\begin{proposition}[Pl\"ucker錐への制限の標的次元]
\label{prop:v17-plucker-dimensions}
分解可能な $a$ の錐への制限で、$\operatorname{Sym}^dA$ をCartan成分
$S_{(d,d)}V$ に置き換えたHom空間の次元は
\[
\begin{array}{c|rrrrrrrr}
 &C_{24}&G_{19}&U_{18}&D_{20}&B_{24}&K_{3,24}&K_{4,24}&K_{5,24}\\\hline
 \dim&0&2&2&0&0&18&36&7
\end{array}
\]
である。また $U_{16}$ の制限標的は2次元である。
\end{proposition}

\begin{proof}
Littlewood分解の $j=0$ のみを残して同じPieri計算を行う。
$U_{18}$ では $F$ 側分割 $(18,18,18,0,0)$ の多重度2、
$U_{16}$ では $(16,16,16,0,0)$ の多重度2を得る。
$D_{20}$ は第一行20が比較する矩形の幅19を超えるため全成分が消える。
全成分を保存証明書の \texttt{Plucker} 欄に記録した。
\end{proof}

この表は制限先のHom空間の次元であり、特定の候補が通常商全体を尽くすことを意味しない。
二つの候補 $\Theta_8U_4^4$、$\Theta_8U_4J_6^2$ の通常商での独立性は
normal sliceの単射定理に依存し、それらが全商を尽くすかどうかは接触階数によって判定する。

Fermat固定群の成分別計算から得る評価階数の上界は
\[
\begin{array}{c|r|r|r}
 &\dim&\text{対角値の階数上界}&\text{対角値と一階法jetの階数上界}\\\hline
C_{24}&2458&621&1962\\
G_{19}&11676&2253&8319\\
U_{18}&1365&224&863\\
D_{20}&6262&1314&4594
\end{array}
\]
である。実際の評価階数との等号は主張しない。

\subsection{再現用証明書}

\texttt{characters/v17\_character\_certificate.py} が全重みEuler積・Newton再帰、
全Schur次元和、テンソル積の交代和とPieri照合を新規に実行する。
\texttt{v17\_koszul\_pieri\_check.py} は別実装の垂直帯列挙を使い、
保存Schur分解から全Koszul源を再計算する。
\texttt{v17\_character\_audit.py} は保存証明書のtwist、交代和、Weyl次元、
Pieri、補助十七次、制限上界を監査する。
保存監査と全重みの新規再計算は別の操作として記録する。

% END characters/v17_character_notes_ja.tex


% BEGIN rank/rank24_section_ja.tex
\section{二十四次の通常接触商の厳密決定}
\label{sec:v17-contact-rank}

二十四次の計算では、19頂点の新しい場だけを列挙する代わりに、
下位接触像の単元倍と既知のHessian作用の全次数定理を併用する。
全ての追加列が通常像に属することを先に保証し、
二階Taylor補正の未証明な所属を階数の下界へ混ぜない。

\begin{proposition}[二十四次の通常像と二次元商]
\label{prop:v17-ordinary-contact-rank24}
二十四次では
\[
 \dim J_{24}=2456,\qquad
 \dim(C_{24}/J_{24})=2,
\]
かつ
\[
 C_{24}/J_{24}
 =\mathbb C\cdot[\Theta_8U_4^4]
   \oplus\mathbb C\cdot[\Theta_8U_4J_6^2].
\]
従って
\[
 \dim\ker(dP:G_{19}\to C_{24})=9220,\qquad
 \dim\ker(dF:D_{20}\to C_{24})=3806.
\]
\end{proposition}

\begin{proof}
命題\ref{prop:v17-character-dimensions}より
$\dim C_{24}=2458$、$\dim U_{18}=1365$ である。
普遍係数環が整域で $P\neq0$ なので、$u\mapsto Pu$ は単射であり、
$\dim(PU_{18})=1365$ となる。
以下では $P=0$ 上の評価を用いて、この部分空間を除いた像の下界を求める。

まず、既存の低次の場 $\Gamma$ と単元グラフ $u$ に対する
$u\,dP(\Gamma)=dP(u\Gamma)$ を作る。
次数の組は $(\deg_a\Gamma,\deg_a u)=(15,4),(13,6),(11,8),$
$(9,10),(7,12),(10,9),(8,11)$ である。
629個の候補から502個の独立な列を選んだ。
この積は、単元グラフと場のグラフの互いに素な和を19頂点のグラフとして書いたものでもある。
四頂点の生の単元グラフは $24U_4$ であり、
これを $U_4$ と取り違えない。像の基底の構成では生のグラフを使う。

次に、定理\ref{thm:v16-positive-unit-hessian-annihilation}により
$HU_{14}\subset J_{24}$、系\ref{cor:v16-all-hessian-first-actions}により
$dH(G_{15})\subset J_{24}$ である。
これらから、それぞれ97列、241列を追加する。
最後に、19頂点の連結グラフで定義される $G_{19}$ の場から251列を追加した。
開いた交代頂点への二重辺も探索対象に含めた。
採用された列数の内訳は
\[
 \begin{array}{c|r}
 \text{通常像への所属が保証された列}&\text{採用数}\\\hline
 u\,dP(\Gamma)\text{ の下位積}&502\\
 H\,u,\ u\in U_{14}&97\\
 dH(\Gamma),\ \Gamma\in G_{15}&241\\
 dP(\Gamma),\ \Gamma\in G_{19}\text{ の連結グラフ}&251\\\hline
 \text{計}&1091
 \end{array}
\]
となる。各列は普遍的な多項式であり、Hessianの列も拡大した有効像ではなく通常像に属する。

評価点は1120個の有理点で、各点で $q(a)=(1,1,1,1,1)$、$P=0$ が有理数上で成立する。
選んだ1091行と1091列の行列式は
\[
 \det M_{24}\equiv579\pmod{1009}
\]
である。全分母は1009と互いに素であり、従ってこの有理行列式は非零である。
全ての $PU_{18}$ の元は評価点で消えるので、1091列は $PU_{18}$ を法として独立であり、
\[
 \dim J_{24}\ge1365+1091=2456
\]
を得る。

他方、normal sliceの単射定理\ref{thm:v16-hessian-socle-injection}より
$[\Theta_8U_4^4]$ と $[\Theta_8U_4J_6^2]$ は通常商で独立である。
従って $\dim J_{24}\le2458-2=2456$ であり、両側の評価が一致する。
$PU_{18}$ はEuler場で $dP(G_{19})$ に吸収でき、
Pfaffian微分 $\delta q:G_{19}\to D_{20}$ は全射なので、
二つの核の次元はそれぞれ $11676-2456=9220$、$6262-2456=3806$ となる。
\end{proof}

\subsection{独立監査と補元の有限証明書}

探索側は左の頂点で辺を番号付けし、サンプル軸を持つ整数行列積で縮約する。
独立監査側は右の頂点で番号付けし直し、元の貪欲な縮約経路と
\texttt{einsum} による別の整数縮約を用いる。
19頂点のグラフは57個の辺添字を持つため、各二項縮約で添字を局所的に振り直す。
各加算の前には
\[
 (1009-1)^2\,5^s<2^{63}
\]
を確認する。ここで $s$ はその二項縮約の総和添字の個数である。
有限体への剰余を各段階で取り、浮動小数点計算は使わない。

非連結グラフの成分は実際の二部グラフから分解し、
同じ成分の値は一度だけ独立縮約する。
905個の相異なる因子から、採用された小行列の
\[
 1091^2=1190281
\]
成分を全て復元して探索側と照合する。
Hessianの入力微分は探索側では余因子を使い、独立側では
$h(q+te_i)$ の次数5の補間から求める。
また有限体の行列式は、配列演算による消去とPython整数による別の消去で照合する。
最終独立監査では1190281成分が全て一致し、二つの消去はいずれも579を返した。
成分分解はグラフの定義から従う因数分解であり、行列の一部の検査を省略する操作ではない。

二つの補元の値は全1120点で有理数上に新規計算した。
$U_4$ は四頂点グラフを24で割ったもの、$J_6$ は既存の固定した六頂点グラフを使う。
補元を追加した1093次の小行列式は
\[
 \det(M_{24},W_{24,1},W_{24,2})\equiv56\pmod{1009}
\]
である。この有限行列は明示した部分空間からの独立性を再確認する。
全通常像に対する非所属は、有限体で列が増えたことだけからではなく、
前の証明で引用したnormal sliceの定理により保証される。

探索が停止したこと、有限体の値が零であったこと、あるいはEuler交代和は、
いずれも通常像の上界には使っていない。

\begin{corollary}[二十四次のnormal核とKoszulホモロジー]
\label{cor:v17-rank24-normal-koszul}
normal射影の直和分解\eqref{eq:v17-normal-splitting}に現れる核は
\[
 \mathcal K_{24}^{\mathrm{nor}}=0
\]
である。従って二十四次ではnormal三次jetの二つのスカラー係数が通常商の完全な検出子となる。
また $b=\operatorname{rank}\delta_2$ と書けば、
\[
 \dim H_1=3806-b,\qquad
 \dim H_2=3807-b,\qquad
 \dim H_2=\dim H_1+1\ge1.
\]
\end{corollary}

\begin{proof}
通常商の次元2と、直和因子
$\mathbb C\cdot[\Theta_8U_4^4]\oplus\mathbb C\cdot[\Theta_8U_4J_6^2]$
の次元2が一致するため、残るnormal核は零である。
Koszul第1微分の核の次元は3806であり、
第3微分の階数は前節の高次ホモロジー消失から1808である。
従って $\dim H_1=3806-b$、
$\dim H_2=(5615-b)-1808=3807-b$ となる。
\end{proof}

これは二十四次でもKoszul複体が正次数で完全でないことを意味する。
第2微分の階数 $b$ と $H_1,H_2$ の個別の次元は本計算では確定していない。

\subsection{再現の範囲}

\texttt{rank/contact24\_combined.json} は、全入力、各列の型とグラフ、
有限体値、採用行を収録する。
\texttt{v17\_independent\_modular\_audit.py} が採用小行列の全成分を新規に独立縮約し、
\texttt{contact24\_combined\_independent\_audit.json} に結果と実行脚本のhashを保存する。
\texttt{v17\_rank\_certificate.py} は保存された独立因子から全成分を再構成し、
二つの小行列式と入力・脚本のhashを監査する。
保存監査は大きなテンソル縮約を新規に実行する操作とは区別する。

% END rank/rank24_section_ja.tex


% BEGIN contact/contact_twenty_fourth_ja.tex
% 第17版：24次の厳密な有限台帳と四つの二階族。
% 本節単独では通常24次商の次元も、四族の像所属も仮定しない。

\section{二十四次の接触作用と低次補正の四族への縮約}
\label{sec:v17-contact-twenty-fourth}

$P=P_F=F(q(a))$、$H=\det\operatorname{Hess}_F(q(a))$、$c=1/96$、
$J_n=dP(\mathcal G_{n-5})$ とする。
第16版の定理\ref{thm:v16-twenty-fourth-normal-form}により、
実際の境界には $P+cH+\sum_{n\ge24,\ n\ {\rm even}}R_n$ という偶正規形がある。
本節では未正規化Tate係数 $f_n=d_n^{\mathrm T}$ からその24次類へ至る
補正を正確に追跡する。通常24次商が二次元であることも、
以下の二階補正が零であることも、この段階では用いない。

\subsection{二十四次の対数作用と十七次の補助基底}

\begin{proposition}[二十四次の保存条件と有効像]
\label{prop:v17-effective-twenty-fourth}
最小増分4の接触対数 $L=\sum_{k\ge4}Q_k$、
$Q_k=Y_{k+1}+b_k$ に対し、$P+cH$ の23次までを保存する条件は
\begin{equation}\label{eq:v17-twenty-fourth-preservation}
 Q_kP+cQ_{k-4}H=0\quad(4\le k\le17),\qquad Q_j=0\quad(j<4)
\end{equation}
である。24次作用は $Q_{18}P+cQ_{14}H$ であり、
\[
                            I_{24}=J_{24}.
\]
全普遍接触変換と、奇座標・偶単元に制限した変換は同じ24次像を与える。
\end{proposition}

\begin{proof}
補題\ref{lem:v14-contact-logarithm}により、
差が24次以上である条件は $L(P+cH)$ が24次以上である条件と同値で、
24次では指数作用と対数作用が一致する。
当該作用に関わる偶数増分の鎖は
\[
 Q_6P=0,\qquad Q_{10}P+cQ_6H=0,\qquad
 Q_{14}P+cQ_{10}H=0.
\]
$Q_4,Q_8,Q_{12},Q_{16}$ の鎖と奇数増分は独立なので零に選べる。
一方、構造定理\ref{thm:v16-positive-unit-hessian-annihilation}と
系\ref{cor:v16-all-hessian-first-actions}により
\[
 Q_{14}H\in dH(\mathcal G_{15})+H\mathcal U_{14}\subset J_{24}.
\]
$Q_{18}P$ の像はEuler場を含めて $J_{24}$ であり、
自由な19次場が逆包含を与える。
\end{proof}

\begin{proposition}[十七次通常像と二十二個の代表場]
\label{prop:v17-seventeenth-contact-basis}
\[
 \dim\mathcal C_{17}=31,\quad
 \dim\mathcal G_{12}=155,\quad
 \dim\mathcal U_{11}=9,\quad
 \dim\mathcal D_{13}=92,\qquad
                           J_{17}=\mathcal C_{17}.
\]
特に $J_{17}/P\mathcal U_{11}$ は22次元である。
公開する22個のグラフ場 $\Gamma_{12,k}$ の像はその基底となる。
$dP$ と $dF$ の核の次元は、それぞれ124と61である。
\end{proposition}

\begin{proof}
四つの源の次元は
\nolinkurl{v17_character_certificate.json} の
\texttt{auxiliary\_degree17} に記録した全重み指標と独立Pieri計算による。
$P=0$ の一般整数入力で22グラフの像を有理数上で再縮約すると、
公開した $22\times22$ 行列の行列式は非零で、1009を法とする値は529である。
この行列と有理数Gaussian消去を
\nolinkurl{v17_g12_graph_basis.py} とそのJSONで再現する。
従って22方向は $P\mathcal U_{11}$ を法として独立である。
$P\mathcal U_{11}$ は9次元で $J_{17}$ に含まれるので、
$\dim J_{17}\ge22+9=31=\dim\mathcal C_{17}$ となる。
全次数Pfaffian全射により $dF$ も同じ像を持つため、核の次元も従う。
\end{proof}

\subsection{merged変更の十三項と純粋な有限消去}

前版と同じ $X_7,X_8$ および $U_4,V_6$ を用い、
\[
               g=(1+3U_4-V_6)f(a+X_7+X_8)
\]
とする。単元を $(1+3U_4)(1-V_6)$ に置き換えない。
また $X_7,X_8$ は一つの座標変更の和であり、二つの変更の逐次合成ではない。
既存の恒等式は
\[
 f_6=P,\quad f_7=f_8=f_9=f_{11}=0,\qquad
 f_{10}=cH-3PU_4,\quad f_{12}=PV_6-dP(X_7),\quad f_{13}=-dP(X_8).
\]
従って $g_{12}=g_{13}=0$ である。
$A_{14}=g_{14}$、$A_{16}=g_{16}$ と置き、
\[
 dP(\Xi_9)=A_{14},\qquad dP(\Xi_{10})=f_{15},\qquad
 dP(\Xi_{11})=A_{16}
\]
という純粋な持ち上げで14、15、16次を順に消す。
続く18、20、22次の入力も既存の結果により通常像に入るので、
それぞれ $\Xi_{13},\Xi_{15},\Xi_{17}$ で消す。

\begin{proposition}[二十四次の正確な有限補正式]
\label{prop:v17-exact-twenty-fourth-correction}
初回変更の24次係数は
\begin{equation}\label{eq:v17-initial-twenty-fourth}
\begin{split}
 g_{24}={}&f_{24}+df_{18}(X_7)+df_{17}(X_8)
       +\tfrac12 f_{12}''(X_7,X_7)+\tfrac12 f_{10}''(X_8,X_8)\\
 &+\tfrac16P'''(X_7,X_7,X_7)\\
 &+3U_4\bigl(f_{20}+df_{14}(X_7)+df_{13}(X_8)
                         +\tfrac12P''(X_8,X_8)\bigr)\\
 &-V_6\bigl(f_{18}+df_{12}(X_7)+\tfrac12P''(X_7,X_7)\bigr).
\end{split}
\end{equation}
上の純粋な有限消去後は正確に
\begin{equation}\label{eq:v17-finite-twenty-fourth}
\begin{split}
 R_{24}^{\mathrm{fin}}={}&g_{24}-dA_{16}(\Xi_9)-df_{15}(\Xi_{10})
                  +\tfrac12P''(\Xi_{10},\Xi_{10})
                  -c\,dH(\Xi_{15})
\end{split}
\end{equation}
となる。純粋な17次・19次消去を挿入しても24次係数は変わらない。
\end{proposition}

\begin{proof}
初回場の次数増分は6と7である。24次へ到達する一次項は
$df_{18}(X_7)$、$df_{17}(X_8)$、二次項は
$f_{12}''(X_7,X_7)/2$、$f_{10}''(X_8,X_8)/2$ である。
混合項には零の $f_{11}$ が必要となる。
三階項は $6+3\cdot6=24$ の $P'''(X_7,X_7,X_7)/6$ が初めて現れ、
四階以上は30次以上である。
四次単元を掛ける前の20次係数と、六次単元を掛ける前の18次係数を展開すると、
\eqref{eq:v17-initial-twenty-fourth}の全十三項を得る。

14次消去の増分8は16次項へ作用する。
15次消去の増分9は15次項へ作用し、さらに $6+2\cdot9=24$ の
$P$ 二階項を生む。この段の15次係数は元の $f_{15}$ のままである。
16、18、22次消去の対応先は、既に零である14、12、8次係数である。
20次消去の増分14は $cH$ に作用する。
17次と19次の純粋消去の対応先は零の13次と11次であり、
各々の二階Taylor項も24次を超える。
これで\eqref{eq:v17-finite-twenty-fourth}を得る。
\end{proof}

\nolinkurl{v17_degree_ledger.py} はこの十三項と、
14次から23次までの各純粋消去・単元を残す消去の次数を独立に列挙する。
これは次数の検査であり、通常像への所属の証明とは区別する。

\begin{remark}[単元を別に残す場合]
14次消去を $(1-B_8)g(a-\Xi'_9)$ とするなら、
24次に $-B_8A_{16}$ が追加される。
15次の入力を $f_{15}=dP(\Xi'_{10})+PB_9$ として消すと、
自己作用から $-PB_9^2$ が追加される。
20次消去の単元は $-cB_{14}H$ を加える。
それぞれ $A_{16}\in J_{16}$、$P\mathcal U_{18}\subset J_{24}$、
$H\mathcal U_{14}\subset J_{24}$ によって通常商では消える。
別規約の場を用いた二階項は、Euler方向を明示して比較する。
\end{remark}

\subsection{三階Taylor項を含む四つの二階族}

$f_{18}\in J_{18}$ は第14版で確定しており、
命題\ref{prop:v17-seventeenth-contact-basis}により
$f_{17}\in\mathcal C_{17}=J_{17}$ である。
従って
\[
 dP(Z_{18})=f_{18},\quad Z_{18}\in\mathcal G_{13},\qquad
 dP(Z_{17})=f_{17},\quad Z_{17}\in\mathcal G_{12}
\]
と選べる。系\ref{cor:v16-all-hessian-first-actions}から
\[
 W_{12}:\mathcal G_8\longrightarrow\mathcal G_{12},\qquad
                       dP(W_{12}(X))=dH(X)
\]
という純粋な線形持ち上げも選べる。
第14版の $P''(\mathcal G_7,\mathcal G_7)\subset J_{18}$ を用いて
\[
 V_{13}\in\mathcal G_{13},\qquad dP(V_{13})=P''(X_7,X_7),
 \qquad A=DX_7\cdot X_7\in\mathcal G_{13}
\]
とする。これらは全て普遍共変量の有限次元空間で選び、固定した $F$ ごとの除算はしない。

\begin{theorem}[二十四次の実際の類の四族代表]
\label{thm:v17-four-family-twenty-fourth-representative}
通常像を法として
\begin{equation}\label{eq:v17-four-family-twenty-fourth}
\begin{split}
 [R_{24}^{\mathrm{fin}}]=\bigl[f_{24}
 &+P''\bigl(X_7,Z_{18}-\tfrac13(V_{13}+A)\bigr)\\
 &+P''\bigl(X_8,Z_{17}+\tfrac c2W_{12}(X_8)\bigr)\\
 &-P''(\Xi_9,\Xi_{11})-\tfrac12P''(\Xi_{10},\Xi_{10})\bigr]
 \quad\text{in }\mathcal C_{24}/J_{24}.
\end{split}
\end{equation}
従って必要な補正は
\begin{equation}\label{eq:v17-four-pairing-families}
 P''(\mathcal G_7,\mathcal G_{13}),\quad
 P''(\mathcal G_8,\mathcal G_{12}),\quad
 P''(\mathcal G_9,\mathcal G_{11}),\quad
 P''(\mathcal G_{10},\mathcal G_{10})
\end{equation}
の四族に含まれる。
\end{theorem}

\begin{proof}
第14版と式\eqref{eq:v15-quadratic-corrections-contact}の既存の所属により
\[
 P''(\mathcal G_7,\mathcal G_7)\subset J_{18},\qquad
 P''(\mathcal G_7,\mathcal G_9)+P''(\mathcal G_8,\mathcal G_8)\subset J_{20}.
\]
$f_{20}\in J_{20}$、$f_{18}\in J_{18}$、
$f_{14}\in J_{14}$、$f_{13}=-dP(X_8)$、
$f_{12}=PV_6-dP(X_7)$ を用いて微分すると、
\eqref{eq:v17-initial-twenty-fourth}の $3U_4$ 群と $-V_6$ 群は全て $J_{24}$ に入る。
ここでは $uJ_n\subset J_{n+s}$（$u\in\mathcal U_s$）を用いている。
同じ理由で $(PU_4)''(X_8,X_8)$ と $(PV_6)''(X_7,X_7)$ も $J_{24}$ に入る。

持ち上げの恒等式を微分して
\[
 df_{18}(X_7)\equiv P''(X_7,Z_{18}),\qquad
 df_{17}(X_8)\equiv P''(X_8,Z_{17})\pmod{J_{24}}
\]
を得る。$dH(\mathcal G_{15})\subset J_{24}$ を用いると
\[
 H''(X_8,X_8)+dH(DX_8\cdot X_8)
   =P''(X_8,W_{12}(X_8))+dP(D(W_{12}(X_8))\cdot X_8)
\]
から $H''(X_8,X_8)\equiv P''(X_8,W_{12}(X_8))$ となる。

三階項の符号と係数を明記する。
点ごとの二階微分に対して
\[
 (dP(X_7))''(X_7,X_7)
  =P'''(X_7,X_7,X_7)+2P''(A,X_7)+dP(D^2X_7(X_7,X_7)).
\]
従って
\[
 \tfrac12f_{12}''(X_7,X_7)+\tfrac16P'''(X_7,X_7,X_7)
 \equiv-\tfrac13P'''(X_7,X_7,X_7)-P''(A,X_7).
\]
一方、$dP(V_{13})=P''(X_7,X_7)$ を $X_7$ で微分すると
\[
 P'''(X_7,X_7,X_7)+2P''(A,X_7)
 =P''(X_7,V_{13})+dP(DV_{13}\cdot X_7).
\]
二式を合わせて
\[
 \tfrac12f_{12}''(X_7,X_7)+\tfrac16P'''(X_7,X_7,X_7)
 \equiv-\tfrac13P''(X_7,V_{13}+A)\pmod{J_{24}}
\]
を得る。捨てた $dP$ 項の場は全て $\mathcal G_{19}$ に属する。

最後に $A_{16}=dP(\Xi_{11})$、$f_{15}=dP(\Xi_{10})$ により
\[
 -dA_{16}(\Xi_9)\equiv-P''(\Xi_9,\Xi_{11}),\qquad
 -df_{15}(\Xi_{10})+\tfrac12P''(\Xi_{10},\Xi_{10})
       \equiv-\tfrac12P''(\Xi_{10},\Xi_{10}).
\]
$-c\,dH(\Xi_{15})\in J_{24}$ と合わせると主張を得る。
\end{proof}

\subsection{十九次奇項と全偶化との比較}

\begin{lemma}[十九次の純粋消去]
\label{lem:v17-nineteenth-odd-removal}
上の有限正規化に17次の純粋消去を加えた後の19次係数 $R_{19}$ は $J_{19}$ に入る。
純粋な14次場でこれを消しても24次係数は変わらない。
\end{lemma}

\begin{proof}
17次消去は補題\ref{lem:v16-seventeenth-odd-removal}で可能である。
第16版の全偶正規形との接触変換を比較すると、差は19次以上であり、
最小増分4の対数補題から
\[
                        -R_{19}=Q_{13}P+cQ_9H.
\]
$dH(\mathcal G_{10})+H\mathcal U_9\subset J_{19}$ により右辺は $J_{19}$ に入る。
$dP(\Xi_{14})=R_{19}$ として $a-\Xi_{14}$ で消すと、
24次への一次作用には零の11次係数が必要である。
$P$ の二階項は32次からなので、24次は変わらない。
\end{proof}

\begin{corollary}[全偶正規形の二十四次類]
\label{cor:v17-evenization-twenty-fourth-class}
第16版の全偶正規形の24次係数を $R_{24}^{\mathrm{ev}}$ とすると
\[
                   R_{24}^{\mathrm{ev}}-R_{24}^{\mathrm{fin}}\in J_{24}.
\]
従って式\eqref{eq:v17-four-family-twenty-fourth}は実際の24次類を与える。
\end{corollary}

\begin{proof}
17次と19次を純粋に消した有限形は
$P+cH+R_{21}+R_{23}+R_{24}^{\mathrm{fin}}+O(a^{25})$ である。
全偶形との差は21次以上なので、指数作用と対数作用は24次まで一致する。
$L(R_{21}+R_{23})$ は25次以上であり、24次の差は $Q_{18}P+cQ_{14}H$ となる。
この二項はともに $J_{24}$ に属する。
\end{proof}

\begin{remark}[十九次を残す直接比較]
19次を残したまま比べると、24次には
$\tfrac12Q_5R_{19}$ が加わり得る。
$Q_5P=0$ から $Q_5$ は通常像を保存するので、この項も $J_{24}$ に入る。
ただし前版の「19次消去を明示しなくてよい」という結論だけを引用して、
二階対数項そのものを無視することはできない。
\end{remark}

\subsection{三百十組による有限の所属条件}

\begin{proposition}[二十四次補正の有限検査]
\label{prop:v17-three-hundred-ten-pairings}
$J_{12}/P\mathcal U_6$、$J_{13}/P\mathcal U_7$、
$J_{14}/P\mathcal U_8$、$J_{15}/P\mathcal U_9$、
$J_{16}/P\mathcal U_{10}$、$J_{17}/P\mathcal U_{11}$、
$J_{18}/P\mathcal U_{12}$ の基底を持ち上げる場を選ぶ。
それらの個数は順に
\[
                         2,\ 1,\ 7,\ 6,\ 21,\ 22,\ 60
\]
である。次の310個の共変量が全て $J_{24}$ に入るなら、
$R_{24}^{\mathrm{ev}}-d_{24}^{\mathrm T}\in J_{24}$ となる。
\[
\begin{array}{ll}
 \operatorname{Hess}_F(q)(dq(\Gamma_{7,i}),dq(\Gamma_{13,j})),
       &1\le i\le2,\quad1\le j\le60,\\
 \operatorname{Hess}_F(q)(dq(\Gamma_8),dq(\Gamma_{12,k})),
       &1\le k\le22,\\
 \operatorname{Hess}_F(q)(dq(\Gamma_{9,i}),dq(\Gamma_{11,j})),
       &1\le i\le7,\quad1\le j\le21,\\
 \operatorname{Hess}_F(q)(dq(\Gamma_{10,i}),dq(\Gamma_{10,j})),
       &1\le i\le j\le6.
\end{array}
\]
\end{proposition}

\begin{proof}
$dP(K)=0$ なら
$P''(K,Y)=-dP(DK\cdot Y)$ は通常像に入る。
Euler単元場 $au/6$ については
$P''(au/6,Y)=(5/6)u\,dP(Y)$ も通常像に入る。
従って各二階族の通常類は、各 $dP$ 像を $P$ 倍単元で割った商にのみ依存する。
さらに
\[
 P''(Y,Z)=\operatorname{Hess}_F(q)(dq(Y),dq(Z))
                 +dF(d^2q(Y,Z))
\]
の第二項は、$d^2q(Y,Z)\in\mathcal D_{20}$ と全次数Pfaffian全射により $J_{24}$ に入る。
必要な組数は $2\cdot60+1\cdot22+7\cdot21+6\cdot7/2=310$ である。
各代表の所属から四族全体の所属を得て、前定理と偶化比較を適用する。
\end{proof}

この命題の仮定は普遍な $J_{24}$ 所属であり、有限点での零値と置き換えてはならない。
実際、今回のHesse入力では二階補正の一部が点ごとには非零となる。
後の検出子による消去、または通常商全体を尽くす構造定理が別に必要である。
本版では通常24次階数の確定とnormal射影を合わせ、
系\ref{cor:v17-actual-twenty-fourth-class-zero}で全四族の $J_{24}$ 所属を結論する。

% END contact/contact_twenty_fourth_ja.tex


% BEGIN tate/v17_tate_notes_ja.tex
\section{二十四次Tate係数の新規再帰}
\label{sec:v17-tate24}
この節では、通常接触商への射影とは独立に、四つの明示的入力で
raw Tate係数を有理数体上で計算する。
規約は従来通り $d^{\mathrm T}=\tfrac12\det K$ とする。

\begin{proposition}[二十四次に必要な再帰の範囲]\label{prop:v17-tate-truncation}
二十四次以下の行列式を決めるには $K_1,\ldots,K_{20}$ で十分である。
これらの定数ブロックは、$A_1,\ldots,A_{16}$ を全項で求め、
曲率 $R_{17},R_{18},R_{19},R_{20}$ をそれぞれ内部$x$次数
$4,3,2,1$ 以下へ切断してからEulerホモトピーを適用すれば厳密に求まる。
対応する $A_{17},A_{18},A_{19},A_{20}$ の保持範囲は $3,2,1,0$ である。
\end{proposition}
\begin{proof}
五つの行列成分はいずれも$a$次数一以上から始まるため、
二十四次以下の積に現れる一因子の次数は高々 $24-4=20$ である。
内部重み $2\deg_x-\#\theta+\#\eta=3$ から全 $A_i$ の$x$次数は高々四である。
Clifford積は$x$次数を下げず、Eulerホモトピーはそれを一つ下げる。
第$n$段階の定数ブロックへ最終的に影響し得る範囲を逆算すると、
$n\ge17$ では曲率の$x$次数 $21-n$ 以下、補正の$x$次数 $20-n$ 以下を
残せばよい。既に切断した第$i$段階 ($i<n$) の上限 $20-i$ は
次の曲率に必要な上限 $21-n$ 以上であり、捨てた項は戻らない。
第16版で切断した $A_{15},A_{16}$ はここで全項として流用せず、
本節の再帰では最初から全項を新たに生成する。
\end{proof}

\subsection{二つのHesse型入力}
$m\in\{-1,-2\}$ とし、
\[
 F_m=x_0^3+x_1^3+x_2^3+x_3^3+(-2-6m)x_4^3+6m x_2x_3x_4,
 \qquad
 a_b=b\,e_{01}+b^{-1}(e_2-e_4)\wedge(e_3-e_4)
\]
を $b=1,2$ で評価する。対称三階テンソルの規約では $F_{234}=m$ である。
$q(a_b)=(0,0,1,1,1)$ かつ $F_m(q)=0$ であり、
三元部分の滑らかさの判別条件に現れる数
$(-2-6m)+8m^3=2(m-1)(2m+1)^2$ はそれぞれ $-4,-54$ と非零である。
従ってどちらの五変数三次式も滑らかである。

\begin{proposition}[四入力のraw二十四次係数]\label{prop:v17-tate24-values}
次の有理数を置く。
\begin{align*}
 A_{24,1}&=-\frac{2504247107738347979241135831211}
                      {1458881718210000000000000},\\
 A_{24,2}&=-\frac{6596059248103221614470055785281}
                      {131766880000000000000}.
\end{align*}
四入力での係数は
\[
 \begin{array}{c|cc}
       &b=1&b=2\\\hline
  m=-1&A_{24,1}&64A_{24,1}\\
  m=-2&A_{24,2}&64A_{24,2}
 \end{array}
\]
である。この表はraw係数 $f_{24}=d^{\mathrm T}_{24}$ の値であり、
正規化後の類との比較は接触作用の節で別に行う。
\end{proposition}
\begin{proof}
\texttt{v17\_hesse\_tate\_twentyfourth.py} は各入力から $A_1$ を構成し、
既存の再帰データを入力として読み込まず、命題
\ref{prop:v17-tate-truncation} の全段階を有理数体上で計算する。
共通分母をもつ整数多項式として積とEulerホモトピーを実装し、
各段階の閉性とMaurer--Cartan方程式を保持範囲で厳密に検査する。
混合Hesse項については自然な外積$\eta$記号とPBW記号の変換を先に適用し、
PBW順序のままホモトピーを取らない。
独立な32次元外積作用素で、素数 $1000003$ を法として
全 $0,\ldots,16$ 次の方程式と全20定数ブロックを照合する。
行列式は有理数体上で置換展開と部分集合動的計画の二方法により照合した。
四入力の全テンソル、交代行列、20ブロック、25行列式係数、
各段階の記録は同梱のJSONとログに保存している。
\end{proof}

独立な保存値監査は、上の再帰とは別の標準ライブラリだけを用いて
20ブロックから全25係数を再構成する。
$m=-2,b=1$ では共通分母法とは別に各係数を\texttt{Fraction}で扱う再帰も行い、
全ブロックと全行列式係数を比較する。
$m=-1,b=1,2$ は第16版と同一入力なので、旧18ブロックと22次以下の全係数が
一致することを確認する。保存値を再読する監査と、最初からの再帰の再実行とは区別する。

\subsection{法方向の三次jetを直接生成する再帰}
\label{subsec:v17-tate-normal-jets}
上の四点値とは別に、$\varepsilon$ を独立変数として
\[
 a(\varepsilon)=e_{01}+\varepsilon(e_2-e_4)\wedge(e_3-e_4),
 \qquad q(a(\varepsilon))=\varepsilon(0,0,1,1,1)
\]
と置く。係数環を $\mathbb Q[\varepsilon]/(\varepsilon^4)$ に置き換えて
同じMaurer--Cartan再帰を最初から行う。
Clifford積、微分、Eulerホモトピーのいずれも$\varepsilon$次数を下げないため、
この切断は行列式の$\varepsilon^3$係数を厳密に保持する。
$x$次数に関する切断は命題\ref{prop:v17-tate-truncation}と独立に併用できる。

\begin{proposition}[二十四次の二つの法三次jet]
\label{prop:v17-tate24-normal-jets}
$m=-1,-2$ の両方について
\[
 [\varepsilon^3]\,d_{24}^{\mathrm T}(F_m,a(\varepsilon))=0.
\]
これは四つの有限点の値からの補間ではなく、
$\varepsilon$を形式変数とした再帰の厳密な計算結果である。
\end{proposition}
\begin{proof}
\nolinkurl{tate/v17_tate_normal_jet.py} は各 $A_n$ を
$\varepsilon^0,\ldots,\varepsilon^3$ の四つの係数多項式として保持する。
共通分母整数演算によって全 $A_1,\ldots,A_{16}$ と所定の$x$切断をもつ
$A_{17},\ldots,A_{20}$ を新たに生成し、各$\varepsilon$係数で閉性と
Maurer--Cartan方程式を有理数体上で確認する。
独立な32次元外積作用素では $a$次数 $0,\ldots,16$ と
$\varepsilon$次数 $0,\ldots,3$ の全組を素数 $1000003$ 上で確認する。
さらに20ブロックの四係数それぞれを独立な外積作用素で照合する。

$a$次数と$\varepsilon$次数の二重多項式として行列式を展開し、
置換展開と部分集合動的計画の二方法で全 $25\times4$ 係数を比較する。
両入力で二十四次・$\varepsilon^3$ の係数は零である。
入力、全80係数行列、100行列式係数、再帰段階の記録を
\nolinkurl{v17_tate_normal_m1.json} と
\nolinkurl{v17_tate_normal_m2.json} に保存した。
\end{proof}

この二点は三元三次式自身の値が零であり、そのHessianの値は非零である。
したがって横断射影定理と二つの不変式単項式の評価行列を組み合わせれば、
raw二十四次係数の $\Theta_8$ 成分を決められる。
命題\ref{prop:v17-ordinary-contact-rank24}によって
$\mathcal K_{24}^{\mathrm{nor}}=0$ も確定している。
この階数定理と低次正規化の補正を合わせる段階は、次の二節で説明する。

% END tate/v17_tate_notes_ja.tex


% BEGIN detector/v17_detector_notes_ja.tex
% 第17版：二点normal三次jetでTheta二方向を読む。
% 通常24次像の全階数とは独立に、normal射影の検出を確定する。

\section{横断三次項による二十四次の二つの係数検出}
\label{sec:v17-normal-third-jet-detectors}

通常商全体を評価グラフの大行列から復元する前に、
定理\ref{thm:v17-split-normal-projection}の射影
\[
 \pi_{24}:\mathcal C_{24}/J_{24}\longrightarrow
 \mathcal A_{16}
 =\mathbb C\,\overline U_4^4\oplus
   \mathbb C\,\overline U_4\overline J_6^2
\]
を二つの横断三次項で読む。
この方法による全接触像の消去は構造定理から従い、
有限個の接触グラフを検査して全てを尽くしたと推測するものではない。

\subsection{二つのHesse入力と正規化}

$m=-1,-2$ に対して
\[
\begin{split}
 F_m&=x_0^3+x_1^3+x_2^3+x_3^3+(-2-6m)x_4^3+6m x_2x_3x_4,\\
 a_\varepsilon&=e_0\wedge e_1+
             \varepsilon(e_2-e_4)\wedge(e_3-e_4)
\end{split}
\]
と置く。$q(a_\varepsilon)=(0,0,\varepsilon,\varepsilon,\varepsilon)$ であり、
三元部分 $f_m$ は $f_m(1,1,1)=0$ を満たす。
$h_f=\det(\operatorname{Hess}_f/6)$ とする。
錐上の単元制限と三元Hessianの値は
\[
\begin{array}{c|rrr}
 m&\overline U_4(f_m)&\overline J_6(f_m)&h_{f_m}(1,1,1)\\ \hline
 -1&5&528&-4\\
 -2&36&7128&-54
\end{array}
\]
である。\nolinkurl{v17_normal_detector_certificate.py} は
$a_0=e_0\wedge e_1$ において二つの単元グラフを有理数上で再縮約し、
この表を検証する。$U_4$ はraw四頂点グラフの $1/24$ 倍であり、
その正規化も固定している。

線形汎関数
\[
 j_m(C)=[\varepsilon^3]C(F_m,a_\varepsilon)
 \qquad(C\in\mathcal C_{24})
\]
を定める。$W_1=\Theta_8U_4^4$、$W_2=\Theta_8U_4J_6^2$ とすると
\begin{equation}\label{eq:v17-normal-jet-matrix}
 M=\begin{pmatrix}
 j_{-1}(W_1)&j_{-1}(W_2)\\
 j_{-2}(W_1)&j_{-2}(W_2)
 \end{pmatrix}
 =
 \begin{pmatrix}
 -2500&-5575680\\
 -90699264&-98771498496
 \end{pmatrix}.
\end{equation}
\[
 \det M=-258781326059520\ne0.
\]

\begin{proposition}[二点の横断三次係数検出子]
\label{prop:v17-normal-jet-detectors}
\begin{align}
 \rho_1^{\mathrm{nor}}
   &=\frac{177147\,j_{-1}-10\,j_{-2}}{464125140},
       \label{eq:v17-normal-rho-one}\\
 \rho_2^{\mathrm{nor}}
   &=\frac{-22674816\,j_{-1}+625\,j_{-2}}{64695331514880}
       \label{eq:v17-normal-rho-two}
\end{align}
とすると、$\rho_i^{\mathrm{nor}}(J_{24})=0$、
$\rho_i^{\mathrm{nor}}(W_j)=\delta_{ij}$ である。
任意の $C\in\mathcal C_{24}$ に対して
\[
 \pi_{24}([C])
  =\rho_1^{\mathrm{nor}}(C)\,\overline U_4^4
       +\rho_2^{\mathrm{nor}}(C)\,\overline U_4\overline J_6^2.
\]
汎関数は $\mathbb Q$ 上で定義され、$\mathbb C$ へ係数拡大して用いる。
この命題は通常商全体の次元を仮定しない。
\end{proposition}

\begin{proof}
構造定理によって
\[
 \gamma_{24}(C)=A_C f+B_C h_f,\qquad
 B_C=b_1\overline U_4^4+b_2\overline U_4\overline J_6^2
\]
と書ける。二入力では $f_m(1,1,1)=0$ なので
\[
 \begin{pmatrix}j_{-1}(C)\\j_{-2}(C)\end{pmatrix}
       =M\begin{pmatrix}b_1\\b_2\end{pmatrix}.
\]
式\eqref{eq:v17-normal-rho-one}--\eqref{eq:v17-normal-rho-two}は
この $2\times2$ 行列の逆行列を有理数で計算したものである。
$C\in J_{24}$ では構造定理により $B_C=0$ なので、
二つの $j_m$ 自体が $J_{24}$ を消す。
これで全ての主張が従う。
\end{proof}

\begin{remark}[通常商との区別]
常に
\[
 \frac{\ker(\rho_1^{\mathrm{nor}},\rho_2^{\mathrm{nor}})}{J_{24}}
                      =\mathcal K_{24}^{\mathrm{nor}}
\]
である。命題\ref{prop:v17-ordinary-contact-rank24}は
$\dim J_{24}=2456$ と、通常商が $[W_1],[W_2]$ で尽くされることを証明した。
従って24次では
$\ker(\rho_1^{\mathrm{nor}},\rho_2^{\mathrm{nor}})=J_{24}$ となる。
これは横断検出子の構成に、別に証明した通常商の次元を合わせた結論である。
\end{remark}

\subsection{実際の係数と低次補正の横断射影}

\begin{proposition}[二十四次の実際のnormal射影は零である]
\label{prop:v17-actual-normal-projection-zero}
未正規化Tate係数と、第16版の偶正規形の24次係数は
\[
 \rho_1^{\mathrm{nor}}(d_{24}^{\mathrm T})
 =\rho_2^{\mathrm{nor}}(d_{24}^{\mathrm T})=0,\qquad
 \pi_{24}([R_{24}^{\mathrm{ev}}])
 =\pi_{24}([d_{24}^{\mathrm T}])=0
\]
を満たす。従って両者の通常類は $\mathcal K_{24}^{\mathrm{nor}}$ に入る。
この主張は通常24次階数の決定を必要としない。
\end{proposition}

\begin{proof}
命題\ref{prop:v17-tate24-normal-jets}のfreshな横断再帰により
\[
 [\varepsilon^3]d_{24}^{\mathrm T}(F_{-1},a_\varepsilon)
 =
 [\varepsilon^3]d_{24}^{\mathrm T}(F_{-2},a_\varepsilon)=0
\]
である。前命題の逆行列を適用すれば二つの検出値が零になる。

定理\ref{thm:v17-four-family-twenty-fourth-representative}と
系\ref{cor:v17-evenization-twenty-fourth-class}により、
$R_{24}^{\mathrm{ev}}-d_{24}^{\mathrm T}$ は $J_{24}$ と四つの二階族の和に入る。
定理\ref{thm:v17-quadratic-actions-normal-zero}は
全てのこの二階族の $\pi_{24}$ 射影を消す。
従って両者のnormal射影は一致し、零である。
\end{proof}

\begin{corollary}[実際の二十四次通常類の消滅]
\label{cor:v17-actual-twenty-fourth-class-zero}
\[
 d_{24}^{\mathrm T}\in J_{24},\qquad
 R_{24}^{\mathrm{ev}}\in J_{24},\qquad
 R_{24}^{\mathrm{ev}}-d_{24}^{\mathrm T}\in J_{24}.
\]
また24次を持つ全二階族 $P''(\mathcal G_r,\mathcal G_s)$ は $J_{24}$ に入る。
\end{corollary}

\begin{proof}
命題\ref{prop:v17-ordinary-contact-rank24}から
$\mathcal K_{24}^{\mathrm{nor}}=0$ であり、$\pi_{24}$ は同型である。
前命題のnormal射影零は、従って通常像への所属を意味する。
全二階族についても定理\ref{thm:v17-quadratic-actions-normal-zero}により
同じ同型の核に入り、零となる。
\end{proof}

\paragraph{保存ブロックの独立監査.}
\nolinkurl{v17_normal_saved_audit.py} は研究計算のPythonモジュールをimportせず、
標準ライブラリの有理数だけで保存した20次数・各4係数のブロックを読む。
全125個の三次テンソル成分と $a$ の定数・$\varepsilon$ 係数を検出側入力と照合し、
$K_1=-a$、$K_2=\operatorname{Hess}_F(q)/3$ も確認する。
120置換による疎多項式展開と、列部分集合による密係数の動的計画法を別に実装し、
各入力でTate次数 $0,\ldots,24$、$\varepsilon$ 次数 $0,\ldots,3$ の全100係数を
再計算して保存値と比較する。二入力の全200係数が一致し、
$d_{24}^{\mathrm T}$ のnormal三次値と二検出値は零であった。
入力JSONと生成ソースのSHA256も照合する。
これは保存ブロックからの監査であり、新規Maurer--Cartan再帰とは別である。

\subsection{有限評価による別の補助検査}

横断三次係数を使わない比較用データとして
\[
 a_b=b\,e_0\wedge e_1+b^{-1}(e_2-e_4)\wedge(e_3-e_4),
 \qquad m=-1,-2,\quad b=1,2
\]
の四点でも評価した。
$q(a_b)=(0,0,1,1,1)$、$P=0$ であり、五変数Hessianとその $q$ 勾配は点ごとに零である。
二候補 $W_1,W_2$ は $b^{18}$ に比例する。
選択した接触グラフには $b^6$ 型の値が現れ、
$64\operatorname{ev}_{(F_m,a_1)}-\operatorname{ev}_{(F_m,a_2)}$
はその検査済みの列を消す。
この有限グラフ検査だけを、全通常像を消す証明には用いない。

命題\ref{prop:v17-three-hundred-ten-pairings}の全310組を四点で有理数再縮約した。
非零なのは $G_7$ の第1代表と $G_{13}$ の第17、第46代表の二組だけである。
この二組をそれぞれ $B_{17},B_{46}$ と書くと、$b=1$ の値は
\[
\begin{array}{c|rr}
 m&B_{17}&B_{46}\\ \hline
 -1&-67584&-135168\\
 -2&-2244806784&-4489613568
\end{array}
\]
となった。$b=2$ ではそれぞれ64倍であり、残る308組は四点全てで零である。
従ってこの二階補正は点ごとに全て零なのではなく、表示した二点差分で消える。
グラフの順序と値は \nolinkurl{v17_lower_pairings_exact.json} に記録する。
全四族のnormal射影零は前述の構造定理によって証明し、
この有限評価はそれとは別の検査として保持する。

% END detector/v17_detector_notes_ja.tex


% BEGIN contact/twenty_sixth_normal_form_ja.tex
% 第17版：通常24次階数2456の独立監査を含む確定最終定理。

\section{二十四次の消去から二十六次開始の正規形へ}
\label{sec:v17-twenty-sixth-normal-form}

\begin{theorem}[二十六次開始の普遍代数的収束偶正規形]
\label{thm:v17-twenty-sixth-normal-form}
実際の境界の24次類は零であり、任意の滑らかな三次式 $F$ に対して
\begin{equation}\label{eq:v17-twenty-sixth-normal-form}
 \boxed{\displaystyle
 d^{\mathrm T}\sim
 P_F+\frac1{96}\det\operatorname{Hess}_F(q(a))
      +\sum_{\substack{n\ge26\\n\ \mathrm{even}}}R_n,
 \qquad R_n\in\mathcal C_n}
\end{equation}
という全次数の偶正規形が存在する。
接触変換は普遍同変な代数的冪級数であり、各形式係数は $F$ に多項式的である。
座標変更の線形項は恒等写像、単元の定数項は1であり、
各固定 $F$ では変換・逆変換・正規形は原点近傍で収束する。
26次以降の実際の類の消滅は主張しない。
\end{theorem}

\begin{proof}
第16版の定理\ref{thm:v16-twenty-fourth-normal-form}から
\[
 f_{\mathrm{ev}}=P_F+\frac1{96}H+R_{24}^{\mathrm{ev}}
                                      +R_{26}^{\mathrm{ev}}+\cdots
\]
という普遍代数的収束偶正規形を得る。
命題\ref{prop:v17-actual-normal-projection-zero}により
$\pi_{24}([R_{24}^{\mathrm{ev}}])=0$ である。
ここには十三項の初回台帳、三階Taylor項の二階族への還元、
十九次奇項の消去と全偶化との比較が含まれる。

命題\ref{prop:v17-ordinary-contact-rank24}により通常24次商は二次元であり、
$[\Theta_8U_4^4]$、$[\Theta_8U_4J_6^2]$ が基底である。
従って $\mathcal K_{24}^{\mathrm{nor}}=0$ であるため、
通常商上の $\pi_{24}$ は同型である。
系\ref{cor:v17-actual-twenty-fourth-class-zero}のとおり
$R_{24}^{\mathrm{ev}}\in J_{24}$ となる。
これはraw係数だけを正規化後の係数と同一視する議論ではなく、
四つの二階補正族のnormal射影零を経由した結論である。
また下位jetは $P+cH$ なので、構造定理の $I_{24}=J_{24}$ が適用できる。

普遍共変量の有限次元空間で
\[
 \Xi_{19}\in\mathcal G_{19},\qquad dP(\Xi_{19})=R_{24}^{\mathrm{ev}}
\]
と選び、奇多項式変更 $a\mapsto a-\Xi_{19}(F,a)$ を行う。
この場は $F$ 次数18、$a$ 次数19であり、判別式等による除算を必要としない。
$P$ への一次作用が24次を消し、
$H$ への最初の作用は28次、$P$ の二階項は42次からである。
従って既存の低次正規形を保ち、次に残り得る次数は26となる。
奇座標変更なので全次数の偶性も保たれる。

追加変更は有限多項式で原点でのJacobianが恒等写像であり、
形式逆は代数的、固定した $F$ では解析的な局所逆も存在する。
第16版の変換との有限合成によって、
代数性、局所収束性、普遍同変性と各係数の $F$ 多項式性が保たれる。
\end{proof}

\begin{remark}[次に残る次数と全Hessian形]
次に検討すべき実際の偶数次数は26である。
本定理は26次以降の類の消滅や、境界全体が $P+H/96$ に接触同値であることを
まだ主張しない。
また定理\ref{thm:v17-finite-normal-kernel}は未知部分の有限次元性を与えるが、
その全次数、安定域の開始次数、および
$\mathcal K^{\mathrm{nor}}=\mathbb C\cdot[H]$ かどうかは未決定である。
\end{remark}

% END contact/twenty_sixth_normal_form_ja.tex

\clearpage
% BEGIN contact/v18_normal_contact_functoriality_ja.tex
% 第18版：normal三次射影の全形式接触変換則。
% 第17版の二階P作用だけを扱う定理を、全高階作用へ強める。

\section{横断三次射影の全次数接触変換則}
\label{sec:v18-normal-contact-functoriality}

第17版では全二階 $P$ 補正のnormal射影が零であることを示した。
本節では任意の有限・形式接触変換に対する変換則を求める。
これにより、最初の非零normal係数は接触不変であることが分かり、
高次のraw Tate係数と正規化後の係数をnormal射影上で直接比較できる。
通常商の残る頂点部分の消失を、この変換則だけから結論するものではない。

$V=L_2\oplus Q_3$、$a_0=e_0\wedge e_1$ とし、
\[
 B(p)=p_4e_2\wedge e_3-p_3e_2\wedge e_4+p_2e_3\wedge e_4,
 \qquad a_\varepsilon=a_0+\varepsilon B(p),
 \qquad F_f=x_0^3+x_1^3+f
\]
と置く。第17版の記号
$\mathcal A=\mathbb C[\overline U_4,\overline J_6]$、
$\gamma_n(C)=A_Cf+B_Ch_f$、$h_f=\det(\operatorname{Hess}_f)/216$
を引き継ぐ。

\subsection{場の接線成分と法一次成分}

\begin{lemma}[rank二切片上の二つのスカラー]
\label{lem:v18-two-normal-field-scalars}
$Y_r\in\mathcal G_r$ に対して、一意な
$\eta_{Y_r},\beta_{Y_r}\in\mathcal A_{r-1}$ が存在し、
\begin{equation}\label{eq:v18-two-normal-field-scalars}
 Y_r(F_f,a_\varepsilon)
   =\eta_{Y_r}(f)a_0
     +\varepsilon\beta_{Y_r}(f)B(p)+O(\varepsilon^2)
\end{equation}
となる。右辺の誤差は $W$ の全成分についての誤差である。
第17版の $\alpha_{dqY_r}$ とは
\[
 \alpha_{dqY_r}=\eta_{Y_r}+\beta_{Y_r}
\]
という関係にある。特に $r$ が偶数なら両スカラーは零である。
\end{lemma}

\begin{proof}
$\lambda(t)=\operatorname{diag}(t,t,1,1,1)$ は $a_0$ を固定する。
$W=\Lambda^2V\otimes\Delta^{-1}$ の
$\Lambda^2L_2$、$L_2\otimes Q_3$、$\Lambda^2Q_3$ 成分の重みは順に
$0,-1,-2$、横断変数 $B$ の重みは $-2$ である。
従って $\varepsilon^k$ 係数に必要な $F$ の重みは、これらに $2k$ を加えた数である。
$k=0$ で非零になり得るのは $\Lambda^2L_2$ 成分だけで、
それは三元成分 $f$ にのみ依存する。
よって $Y_r(F_f,a_0)=\eta_{Y_r}(f)a_0$ と書ける。

$k=1$ の $\Lambda^2Q_3$ 成分は、$f$ のみに依存する $B$ の線形写像である。
$\mathrm{SL}(Q_3)$ に関してこれは $Q_3^*$ のEnd共変量だから、
補題\ref{lem:v17-ternary-endomorphism-scalar}によってスカラーである。
これを $\beta_{Y_r}$ と書く。
残る二成分の一次係数に必要な $F$ の重みは2と1である。
一方 $F_f$ の非零係数の重みは0と3だけなので、その両係数は零である。
両スカラーは $f$ の次数 $r-1$ の $\mathrm{SL}_3$ 不変式であり、
$\mathcal A_{r-1}$ に属する。
\[
 dq_{a_\varepsilon}(Y_r)
   =a_\varepsilon\wedge Y_r
   =\varepsilon(\eta_{Y_r}+\beta_{Y_r})a_0\wedge B+O(\varepsilon^2)
\]
が $\alpha$ との関係を与える。
$\mathcal A$ は次数4と6で生成されるので、奇数次成分は零である。
\end{proof}

\begin{proposition}[一つの斉次共変量の変換公式]
\label{prop:v18-homogeneous-normal-pullback}
$C_n\in\mathcal C_n$、$Y_i\in\mathcal G_{r_i}$ とし、$t_i$ を独立な形式変数とする。
\[
 \Psi_t(a)=a+\sum_i t_iY_i(F,a),\qquad
 s(t)=1+\sum_i t_i\eta_{Y_i},\qquad
 b(t)=1+\sum_i t_i\beta_{Y_i}
\]
と置けば
\begin{equation}\label{eq:v18-homogeneous-normal-pullback}
 [\varepsilon^3]C_n(F_f,\Psi_t(a_\varepsilon))
    =s(t)^{n-3}b(t)^3\gamma_n(C_n)(f,p)
\end{equation}
が成り立つ。
\end{proposition}

\begin{proof}
前補題から $\Psi_t(a_\varepsilon)=s(t)a_0+
\varepsilon b(t)B+O(\varepsilon^2)$ である。
$s(t)$ は形式単元なので、$a_{01}$ を可逆にしたSchur補チャートを用いる。
交差成分を除くdeterminant一の基底変更は $\varepsilon=0$ で恒等であり、
変換後の $F$ の三元成分は $f+O(\varepsilon)$ である。
Schur補の横断成分の一次項は正確に $\varepsilon b(t)B$ であり、
基点の $a_{01}$ 成分は $s(t)+O(\varepsilon)$ である。

補題\ref{lem:v17-covariant-normal-order}により $C_n$ の横断次数は三以上である。
従って三次項の計算では、基点の $O(\varepsilon)$ 変動、
$F$ の $O(\varepsilon)$ 変動、横断成分の $O(\varepsilon^2)$ 変動は寄与しない。
$a$ についての斉次性と横断三次性から
\[
 s(t)^n\left(\frac{b(t)}{s(t)}\right)^3\gamma_n(C_n)
\]
を得る。これが主張である。指数 $n-3$ は非負なので、
結果そのものは $s(t)$ の逆数を含まない多項式恒等式になっている。
\end{proof}

\begin{corollary}[全高階Taylor作用のnormal射影]
\label{cor:v18-all-higher-actions-normal-zero}
任意の $k\ge1$ と $Y_i\in\mathcal G_{r_i}$ に対し、
$n=6+\sum_i(r_i-1)$ と置くと
\[
 \pi_n\bigl([P^{(k)}(Y_1,\ldots,Y_k)]\bigr)=0.
\]
同様に、$H$ の全高階Taylor作用のnormal三次項は零である。
より一般に、$\pi([C_m])=0$ なら、$C_m$ の任意の点ごとの高階Taylor作用も
対応する次数でnormal射影が零となる。
\end{corollary}

\begin{proof}
前命題で $t_1\cdots t_k$ の係数を比較する。
右辺は $\gamma_m(C_m)$ の不変式倍なので、Hessian成分が零ならその性質を保つ。
$\gamma_6(P)=f$、$\gamma_{10}(H)=0$ を代入すれば二つの特例を得る。
ここで微分は点ごとの多重微分であり、場を逐次合成した微分とは区別する。
逐次合成した場合も連鎖律で現れる場の微分は普遍共変場であるため、同じ結論が従う。
\end{proof}

\subsection{全形式境界に対する変換則}

形式境界 $\mathfrak f=\sum_{n\ge6}C_n$ に対し、補助変数 $z$ を使って
\[
 \mathcal B_{\mathfrak f}(z)=\sum_{n\ge6}z^nB_{C_n}\in\mathcal A[[z]]
\]
と定める。これはnormal三次項のHessian係数を全次数でまとめたものである。

\begin{theorem}[normal係数の接触変換則]
\label{thm:v18-normal-contact-law}
普遍同変な接触変換を
\[
 \widetilde{\mathfrak f}(F,a)
   =u(F,a)\mathfrak f(F,\Phi(F,a)),\qquad
 u=1+\sum_{s>0}u_s,\quad
 \Phi=a+\sum_{r>1}Y_r
\]
と書く。各 $u_s\in\mathcal U_s$、$Y_r\in\mathcal G_r$ とし、
\[
 \overline u(z)=1+\sum_{s>0}z^s\overline u_s,
 \qquad
 \sigma(z)=1+\sum_{r>1}z^{r-1}\eta_{Y_r},
 \qquad
 \tau(z)=1+\sum_{r>1}z^{r-1}\beta_{Y_r}
\]
と置く。このとき
\begin{equation}\label{eq:v18-normal-contact-law}
 \boxed{\displaystyle
 \mathcal B_{\widetilde{\mathfrak f}}(z)
 =\overline u(z)\tau(z)^3\sigma(z)^{-3}
      \mathcal B_{\mathfrak f}\bigl(z\sigma(z)\bigr).}
\end{equation}
\end{theorem}

\begin{proof}
$\Phi(F,z a_\varepsilon)/z$
$=a_\varepsilon+\sum_{r>1}z^{r-1}Y_r(F,a_\varepsilon)$ と展開する。
前命題を各 $C_n$ に適用すると、そのnormal三次項は
$z^n\sigma(z)^{n-3}\tau(z)^3\gamma_n(C_n)$ である。
各 $C_n$ の横断次数は三以上なので、単元を掛けるときは
$u(F,z a_0)=\overline u(z)$ のみが寄与する。
Hessian側の係数を取り、$n$ について加えると公式を得る。
各 $z$ 次数で有限個の項しか寄与しないので、全ての操作は完備環上で定義される。
\end{proof}

\begin{corollary}[最初の非零normal係数は接触不変である]
\label{cor:v18-first-normal-contact-invariant}
上の接触変換について、$\mathcal B_{\mathfrak f}=0$ と
$\mathcal B_{\widetilde{\mathfrak f}}=0$ は同値である。
さらに一方の $N$ 次未満の係数が全て零なら他方もそうであり、
\[
 [z^N]\mathcal B_{\widetilde{\mathfrak f}}
 =[z^N]\mathcal B_{\mathfrak f}
\]
となる。
\end{corollary}

\begin{proof}
$\overline u,\sigma,\tau$ はいずれも定数項1の形式単元であり、
$z\sigma(z)=z+O(z^2)$ である。
公式の各係数は、同じ次数の元の係数と、それより低い次数の係数の和である。
次数に関する帰納法により結論を得る。
\end{proof}

\begin{corollary}[二十六次の実際のnormal係数]
\label{cor:v18-actual26-normal-equals-raw}
第17版の偶正規形を
\[
 \mathfrak f_{\mathrm{ev}}=P+\frac1{96}H+
          R_{26}^{\mathrm{ev}}+R_{28}^{\mathrm{ev}}+\cdots
\]
と書く。このとき未正規化Tate係数 $f_n=d_n^{\mathrm T}$ について
\[
 \pi_n([f_n])=0\quad(n<26),\qquad
 \pi_{26}([R_{26}^{\mathrm{ev}}])=\pi_{26}([f_{26}])
\]
が成り立つ。これは通常26次商の次元を仮定しない。
\end{corollary}

\begin{proof}
定理\ref{thm:v17-twenty-sixth-normal-form}の普遍形式接触変換へ前系を適用する。
正規形の26次未満には $P,H$ しかなく、両者のHessian側normal係数は零である。
よってraw側の同じ次数の係数も零であり、最初に残り得る26次係数は一致する。
\end{proof}

\begin{remark}[通常像への所属と区別する]
本節は $\mathcal K^{\mathrm{nor}}$ の消失も、通常26次商の二次元性も使わない。
従って $\pi_{26}([f_{26}])=0$ と分かれば実際の26次類は
$\mathcal K_{26}^{\mathrm{nor}}$ に入るが、その類自体の消滅には別の証明が必要である。
全次数でrawのnormal係数が零と証明できれば、この性質は任意の形式正規化後も保たれる。
\end{remark}

% END contact/v18_normal_contact_functoriality_ja.tex


% BEGIN theory/v18_actual_normal_vanishing_ja.tex
% 第18版：実際の境界の全次数normal消滅。有限標本からの推測ではない。
\section{実際の境界の横断Hessian係数は全次数で消える}
\label{sec:v18-actual-normal-vanishing}

第17版までの有限次数の消滅計算を、ここでは幾何学的な全次数の命題へ強める。
入力は、固定した平面三次曲線上のPicard族が階数二の直線を与えることと、
その非零点での境界の横断三次曲線が元の抽象曲線に一致することである。
Hesse pencilの $j$ 不変量が定数でないことを組み合わせると、
通常商の自由な直和因子に現れる実際の境界の係数が全て零になる。

\subsection{固定した平面のPicard族と横断三次式}

$V=L_2\oplus Q_3$、$a_0=e_0\wedge e_1$ とし、前節と同じ $B(p)$ を用いる。
$f\in\operatorname{Sym}^3Q_3$ を滑らかな平面三次曲線を定める式とし、
\[
 F_f=x_0^3+x_1^3+f,\qquad C_f=V(f)\subset\mathbb P(Q_3^*)
\]
と置く。この $F_f$ も滑らかである。
実際、五つの偏微分が同時に零なら $x_0=x_1=0$ となり、
残る三つの式は $f$ の滑らかさに反する。

先頭項が $P$ の普遍境界方程式を $d=\sum_{n\ge6}d_n$ とする。
定理\ref{thm:v17-split-normal-projection}により
\[
 \gamma_n(d_n)=A_n(f)f+B_n(f)h_f,
 \qquad h_f=\det(\operatorname{Hess}_f)/216
\]
と一意に書ける。以下では二つの形式級数
\begin{equation}\label{eq:v18-normal-pencil-series}
 \mathsf A_f(t)=\sum_{n\ge6}t^{n-6}A_n(f),\qquad
 \mathsf B_f(t)=\sum_{n\ge6}t^{n-6}B_n(f)
\end{equation}
を使う。$A_6=1$、$B_6=0$ なので
$\mathsf A_f(0)=1$、$\mathsf B_f(0)=0$ である。
$a$ に関する斉次性から、直線 $t a_0$ に沿う横断三次項は
\begin{equation}\label{eq:v18-geometric-normal-cubic}
 [\varepsilon^3]d(F_f,t a_0+\varepsilon B(p))
 =t^3\bigl(\mathsf A_f(t)f(p)+\mathsf B_f(t)h_f(p)\bigr).
\end{equation}

\begin{lemma}[横断三次曲線の抽象同型類]
\label{lem:v18-fixed-plane-normal-j}
式\eqref{eq:v18-geometric-normal-cubic}の括弧内が定める
$\mathbb C((t))$ 上の平面三次曲線の $j$ 不変量は $j(C_f)$ に等しい。
\end{lemma}

\begin{proof}
補題\ref{lem:universal-chart-plucker}の証明を、固定分解
$V=L_2\oplus Q_3$ に適用する。
固定平面の曲線 $C_f$ 上のPicard形式族は、可逆な形式パラメータ変更を除いて
直線 $t a_0$ 全体を与える。その非零形式一般点は $C_f$ 上の直線束の評価核に対応する。
適合部分空間は同じ普遍座標で標準Pl\"ucker錐であり、
$t\ne0$ では $\Lambda^2Q_3$ の三方向がその法空間を与える。
従って式\eqref{eq:v18-geometric-normal-cubic}は、その点のnormal coneの三次方程式である。

定理\ref{thm:local-bn-cubic}により、射影化したnormal coneは
$C_f$ を次数三の直線束によって埋め込んだものである。
直線束はPicardパラメータとともに変化してよいが、抽象曲線は同じ $C_f$ である。
従って $j$ は一定である。
形式一般点への適用は、定理\ref{thm:boundary-eighth-vanishing}で用いたものと同じであり、
ここではFermat三次式を一般の滑らかな $f$ に置き換えただけである。
方程式に掛かる単元と $t^3$ は、射影化した三次曲線を変えない。
\end{proof}

\subsection{Hesse pencilからの全次数消滅}

\begin{lemma}[Hesse pencilでの形式的な一意性]
\label{lem:v18-hesse-formal-rigidity}
$\lambda$ を超越元とし、$k=\mathbb C(\lambda)$ とする。
\[
 D=p_2^3+p_3^3+p_4^3,\qquad M=p_2p_3p_4,\qquad
 f_\lambda=D+6\lambda M
\]
と置く。$A(t),B(t)\in k[[t]]$、$A(0)=1$、$B(0)=0$ とし、
$A(t)f_\lambda+B(t)h_{f_\lambda}$ の $j$ 不変量が
$f_\lambda$ の $j$ 不変量に等しいなら、$B(t)=0$ である。
\end{lemma}

\begin{proof}
三行三列の行列式を展開すると
\[
 h_{f_\lambda}=-\lambda^2D+(1+2\lambda^3)M
\]
である。従って $A f_\lambda+B h_{f_\lambda}$ は、単元 $A-\lambda^2B$ を除いて
$D+6\mu(t)M$ と書ける。ただし
\begin{equation}\label{eq:v18-hesse-parameter-change}
 \mu(t)=\lambda+
 \frac{(1+8\lambda^3)B(t)}{6(A(t)-\lambda^2B(t))}.
\end{equation}
第\ref{sec:boundary-higher-jets}節で用いたHesse式を同じ規約に直すと
\begin{equation}\label{eq:v18-hesse-j-function}
 J(\lambda)=
 \frac{110592\lambda^3(1-\lambda^3)^3}{(1+8\lambda^3)^3}.
\end{equation}
これは定数でない有理関数なので、標数零の体 $k$ で $J'(\lambda)\ne0$ である。
$\mu(0)=\lambda$ と $J(\mu(t))=J(\lambda)$ から $\mu(t)=\lambda$ を得る。
例えば $\mu-\lambda$ の最初の非零項が $c t^r$ なら、
$J(\mu)-J(\lambda)$ の同じ次数の係数は $J'(\lambda)c\ne0$ となり矛盾する。
$1+8\lambda^3\ne0$ と式\eqref{eq:v18-hesse-parameter-change}から $B(t)=0$ が従う。
\end{proof}

\begin{theorem}[実際の境界の全次数normal消滅]
\label{thm:v18-all-order-normal-vanishing}
実際のBrill--Noether境界の、先頭項を $P$ に規格化した普遍方程式
$d=\sum_{n\ge6}d_n$ に対して
\begin{equation}\label{eq:v18-all-order-normal-vanishing}
 \boxed{\displaystyle
  \pi_n([d_n])=0\quad\text{for every }n\ge6,\qquad
  \mathcal B_d(z)=0.}
\end{equation}
特に未正規化Tate方程式 $d^{\mathrm T}=\frac12\det K$ にもこの結論が成立し、
任意の普遍形式接触変換の後も成立する。
\end{theorem}

\begin{proof}
まず幾何学的な普遍座標で考える。
前二補題を一般のHesse式に適用すると
$\mathsf B_{f_\lambda}(t)=0$ となる。
ここで超越的な $\lambda$ への移行は、各滑らかな複素パラメータで得た
$j$ の等式を各 $t$ 係数ごとの有理関数恒等式として読むことで正当化できる。
従って各 $n$ について $B_n(f_\lambda)=0$ が $\mathbb C(\lambda)$ で成立する。
各 $B_n$ は三元三次式の全係数空間上の斉次多項式不変量である。
補題\ref{lem:v17-ternary-endomorphism-scalar}の証明で確認した
\[
 \mathrm{GL}_3\times\mathbb A^1\longrightarrow\operatorname{Sym}^3Q_3,
 \qquad(h,\lambda)\longmapsto h f_\lambda
\]
の優勢性と、$\mathrm{GL}_3$ に対する相対不変性から、$B_n$ は恒等的に零となる。
この最後の多項式恒等式は、特異な三元三次式も含めて成立する。
滑らかな開集合上で $j$ を用いた段階と、この多項式的延長の段階を区別している。

定理\ref{thm:tate-boundary-determinant}および
定理\ref{thm:tate-tenth-coefficient}の証明が示すように、Tate座標も
全非零三次式上で正則かつ同変で、接写像を同じedge同型に固定した普遍座標である。
従ってその方程式と上の方程式は、定数項1の単元と一次項が恒等の
普遍形式座標変換で移り合う。
定理\ref{thm:v18-normal-contact-law}により $\mathcal B=0$ はこの変換で保存される。
これでTate方程式と、さらに任意の普遍形式接触変換後について結論を得る。
\end{proof}

\begin{corollary}[全ての実際の剰余類が入る部分]
\label{cor:v18-all-actual-classes-vertex}
実際の境界を有限次数まで普遍形式接触正規化したとき、
残る各斉次係数の通常商における類は $\mathcal K_n^{\mathrm{nor}}$ に入る。
特に、第17版の偶正規形について
\[
 [R_{26}^{\mathrm{ev}}]\in\mathcal K_{26}^{\mathrm{nor}}.
\]
これは $\mathcal K_{26}^{\mathrm{nor}}=0$ も $R_{26}^{\mathrm{ev}}\in J_{26}$ も主張しない。
\end{corollary}

\begin{proof}
全次数消滅と、定理\ref{thm:v17-split-normal-projection}で定義した核の等式による。
\end{proof}

\begin{remark}[全次数の証明と有限計算の役割]
本定理は、有限個のTate係数の零性から外挿したものではない。
幾何学的なnormal coneの同定、Hesse pencilの非定数な $j$、
および多項式共変量の分裂射影を用いた証明である。
後続節の26次の二入力再計算は、この全次数命題と実装規約との独立な整合性検査になる。
また、全てのnormal係数の消滅は、実際の境界が $P+H/96$ と形式接触同値であることを
まだ意味しない。その最後の問題は、有限次元の $\mathcal K^{\mathrm{nor}}$ に残る。
\end{remark}

% END theory/v18_actual_normal_vanishing_ja.tex


% BEGIN theory/v18_effective_vertex_bound_ja.tex
% 第18版：有限頂点部分の明示次数上界。
% 一次資料：Dubé, SIAM J. Comput. 19 (1990), Theorem 8.2, p.773。

\section{頂点部分の明示的な次数上界}
\label{sec:v18-effective-vertex-bound}

第17版で得た分裂を
\[
 \mathcal Q=\bigoplus_{n\ge6}\mathcal C_n/J_n
 =\mathcal A[\Theta_8]\oplus\mathcal K^{\mathrm{nor}},
 \qquad \mathcal A=\CC[\overline U_4,\overline J_6]
\]
と書く。また
\[
 R=\CC[F_{ijk},a_{ij}],\qquad
 G=(g_0,\ldots,g_4),\qquad g_i=(\partial_iF)(q(a)),
 \qquad\mathfrak m=(a_{ij})
\]
とし、$T_0=H^0_{\mathfrak m}(R/G)$ と置く。
$F$ の独立係数は35個、$a$ の座標は10個である。
各 $g_i$ の二重次数は $(\deg_F,\deg_a)=(1,4)$ なので、
45個の変数すべてを次数1とした通常次数では5となる。

第17版では $\mathcal K^{\mathrm{nor}}\subset T_0$ と
$T_0$ の $a$ 次数の有界性を証明したが、その上界は存在的であった。
本節は、巨大ながら具体的な整数上界を与える。
係数体を $\CC(F)$ に変えると $F$ 方向のtorsionを失う可能性があるため、
以下では終始45変数の多項式環を使う。

\subsection{初項の指数から相対torsionの次数を抑える}

\begin{lemma}[初項の箱に入る相対torsion]
\label{lem:v18-initial-box-torsion}
$S=k[y_1,\ldots,y_s,x_1,\ldots,x_r]$ とし、
$J\subset S$ は $y$ 次数と $x$ 次数に関して二重斉次なイデアルとする。
単項式順序を固定し、$L=\operatorname{in}(J)$ と置く。
$L$ の最小単項式生成元の $x_i$ 指数がすべて $D_i$ 以下であるとする。
各 $D_i\ge1$ なら
\begin{equation}\label{eq:v18-initial-box-bound}
 \bigl(H^0_{(x_1,\ldots,x_r)}(S/J)\bigr)_d=0
 \qquad\left(d>\sum_{i=1}^r(D_i-1)\right),
\end{equation}
ここで添字 $d$ は $x$ 次数であり、$y$ 次数は任意である。
ある $D_i=0$ なら、この相対torsion加群全体が零である。
\end{lemma}

\begin{proof}
非零の $x$ 斉次torsion類を取り、その剰余の標準形を $f$ とする。
二重斉次イデアルの被約Gr\"obner基底は二重斉次なので、
標準形も同じ $x$ 次数 $d$ を持つ。
$m=\operatorname{lm}(f)$ と置くと $m\notin L$ である。
各 $i$ に対し、torsion性からある $e_i$ について $x_i^{e_i}f\in J$、
従って $x_i^{e_i}m\in L$ となる。

もし $m$ の $x_i$ 指数が $D_i$ 以上なら、
$x_i^{e_i}m$ を割る $L$ の最小生成元はすでに $m$ を割る。
実際、他の変数の指数は変わらず、その生成元の $x_i$ 指数も $D_i$ 以下だからである。
これは $m\notin L$ に反する。従ってすべての $i$ について
$\deg_{x_i}m\le D_i-1$ であり、式\eqref{eq:v18-initial-box-bound}を得る。
ある $D_i=0$ の場合には同じ議論がどの非零標準単項式にも矛盾を与える。
\end{proof}

\begin{remark}[実際の初項計算で使える、より精密な形]
\label{rem:v18-initial-torsion-certificate}
証明は、torsion類の先頭標準単項式 $m$ が
\[
 m\notin L,\qquad
 \text{各 }i\text{ について、ある }e_i\text{ が存在して }x_i^{e_i}m\in L
\]
を満たすことも示す。
従って完全な初項イデアルが得られれば、単なる最大指数の和より強い上界を
$H^0_{(x)}(S/L)$ の標準単項式から得られる。
ただし、計算途中のGr\"obner基底や一部の初項だけでは、
上記の $D_i$ が全初項を抑えることにはならない。
有限体で得た候補についても、特性零でのGr\"obner基底証明を別に要する。
\end{remark}

\subsection{45変数の斉次Gr\"obner次数上界の適用}

\begin{theorem}[有限頂点部分の明示上界]
\label{thm:v18-effective-vertex-cutoff}
次の整数を定める。
\begin{equation}\label{eq:v18-effective-cutoff-definition}
 D_*=2\cdot18^{\,2^{43}},\qquad
 N_*=10(D_*-1)=20\cdot18^{\,2^{43}}-10.
\end{equation}
すると
\[
 (T_0)_d=0\quad(d>N_*),\qquad
 \mathcal K_n^{\mathrm{nor}}=0\quad(n>N_*).
\]
従って第17版の有限多項式 $E_{\mathrm{vert}}$ は
\begin{equation}\label{eq:v18-Evert-explicit-degree}
 \deg E_{\mathrm{vert}}\le N_*
\end{equation}
を満たす。
\end{theorem}

\begin{proof}
Dub\'eの斉次版\cite[Theorem~8.2, p.~773]{Dube1990}
によれば、$v$ 変数の多項式環において次数 $d$ 以下の斉次多項式で生成された
イデアルは、任意の単項式順序について総次数
\[
 2\left(\frac{d^2}{2}+d\right)^{2^{v-2}}
\]
以下のGr\"obner基底を持つ。
本件は $v=45,d=5$ なので、その上界は
\[
 2(35/2)^{2^{43}}<2\cdot18^{2^{43}}=D_*.
\]
特に各初項の、個々の $a_{ij}$ の指数も $D_*$ 以下である。
補題\ref{lem:v18-initial-box-torsion}を10個の $a$ 変数に適用すると
$(T_0)_d=0$ for $d>N_*$ を得る。
第17版の包含 $\mathcal K^{\mathrm{nor}}\subset T_0$ と、
$\mathcal K_n^{\mathrm{nor}}$ の $a$ 次数が $n$ であることから残りが従う。
\end{proof}

\begin{remark}[指数と実用性の監査]
Dub\'eの同じ論文の非斉次版Corollary~8.3では指数が $2^{v-1}$ となる。
これは斉次化のために変数を一つ追加した結果であり、
本件の斉次45変数の適用では $2^{43}$ が正しい。
一方、$F$ を係数体へ移して10変数として適用することはしていない。

式\eqref{eq:v18-effective-small-conjectures}のような小さい次数の結論は、
本定理からは出ない。
\begin{equation}\label{eq:v18-effective-small-conjectures}
 E_{\mathrm{vert}}=t^{10},\qquad
 \mathcal K_{26}^{\mathrm{nor}}=0.
\end{equation}
上界 $N_*$ は実際の計算範囲としては極めて大きい。
本定理の役割は、存在的な有限性を明示的な有限判定範囲に変える点にある。
通常商の全次数計算や、26次の階数証明を実行したことを意味しない。
\end{remark}

\begin{corollary}[全通常商の残余の有限判定範囲]
\label{cor:v18-Evert-finite-decision}
$E_{\mathrm{vert}}$ の全係数は $6\le n\le N_*$ の通常商の厳密な階数で決まる。
具体的には、その各次数で
\[
 [t^n]E_{\mathrm{vert}}
 =\dim(\mathcal C_n/J_n)
 -\#\{(r,s)\in\mathbb Z_{\ge0}^2:4r+6s=n-8\}
\]
を使えばよい。
これは実行可能性に関する見積りを与えるものではなく、
有限個の厳密な線形代数問題への理論的還元である。
\end{corollary}


% END theory/v18_effective_vertex_bound_ja.tex


% BEGIN theory/v18_vertex_hessian_chain_ja.tex
% 第18版：Hessianの頂点socleと、そのq微分が作る頂点torsion。

\section{Hessianの頂点socleとその一段下の類}
\label{sec:v18-vertex-hessian-chain}

この節では、通常商の選択された同変成分
$\mathcal K^{\mathrm{nor}}$ と、全加群
$T_0=H^0_{(a)}(R/G)$ を区別する必要性を、具体的な恒等式で示す。
特に、$\mathcal K^{\mathrm{nor}}=\CC\cdot[H]$ という予想を
$T_0=\CC[F]H$ へ強めることはできない。

独立な変数 $q_0,\ldots,q_4$ を一度使い、
\[
 g_i=\partial_iF(q),\qquad M=(\partial_i\partial_jF(q))_{i,j},
 \qquad C=\operatorname{adj}(M),\qquad H=\det M,
 \qquad H_j=\partial_{q_j}H
\]
と置く。以下の $H$ は従来の $H_{10}$ と同じ非零定数倍の正規化でよい。
恒等式を記した後で $q=q(a)$ を代入する。

\begin{lemma}[Hessianとその微分の乗法恒等式]
\label{lem:v18-hessian-derivative-identity}
$B=\CC[F,q]$、$G_q=(g_0,\ldots,g_4)\subset B$ とする。
各 $j,k$ について
\begin{align}
 q_kH&=2\sum_l C_{kl}g_l,
 \label{eq:v18-qH-socle}\\
 q_kH_j-\delta_{kj}H
 &=2\sum_l (\partial_{q_j}C_{kl})g_l.
 \label{eq:v18-qHj-duality}
\end{align}
従って $B/G_q$ で $qH=0$、$q_kH_j=\delta_{kj}H$ となり、
$(q)^2H_j=0$ である。
\end{lemma}

\begin{proof}
$g_i$ が二次式なのでEulerの恒等式は $Mq=2g$ となる。
$CM=H\operatorname{id}$ を用いると式\eqref{eq:v18-qH-socle}を得る。
これを $q_j$ で微分すれば
\[
 \delta_{kj}H+q_kH_j
 =2\sum_l(\partial_{q_j}C_{kl})g_l
   +2\sum_l C_{kl}M_{lj}
 =2\sum_l(\partial_{q_j}C_{kl})g_l+2\delta_{kj}H,
\]
従って式\eqref{eq:v18-qHj-duality}が成立する。
\end{proof}

\begin{theorem}[全座標に殺されるHessian類]
\label{thm:v18-H-vertex-socle}
Pfaffian代入 $q=q(a)$ の後では、全ての $i,j$ に対して
\begin{equation}\label{eq:v18-aH-explicit}
 a_{ij}H
 =-2\sum_{k,l}a_{ik}(\partial_{q_j}C_{kl})(q(a))g_l(q(a))
 \in G.
\end{equation}
従って $R/G$ における $H$ の類は
\[
 \mathfrak m[H]=0,
 \qquad\mathfrak m=(a_{ij})
\]
を満たす。第17版以前の通常10次非零性により、この類は非零である。
\end{theorem}

\begin{proof}
符号付きPfaffianの恒等式 $aq=0$ を使い、
式\eqref{eq:v18-qHj-duality}に $a_{ik}$ を掛けて $k$ について足す。
左辺の $\sum_k a_{ik}q_kH_j$ が零となるので、
式\eqref{eq:v18-aH-explicit}を得る。
この式は $F$ の局所化や特殊化を一切使わない普遍多項式恒等式である。
\end{proof}

\begin{proposition}[全頂点加群にはHessian以外も存在する]
\label{prop:v18-Hj-vertex-torsion}
各 $j=0,\ldots,4$ について、$H_j(q(a))$ の類は
$T_0=H^0_{\mathfrak m}(R/G)$ に属し、非零である。
その二重次数は $(\deg_F,\deg_a)=(5,8)$ であり、
\begin{equation}\label{eq:v18-Hj-vertex-chain}
 q_k[H_j]=\delta_{kj}[H],\qquad
 \mathfrak m[H]=0.
\end{equation}
特に
\[
 T_0\ne\CC[F]\,[H].
\]
\end{proposition}

\begin{proof}
補題\ref{lem:v18-hessian-derivative-identity}より、
$I=(q_0,\ldots,q_4)$ に対して $I^2[H_j]=0$ である。
第17版の各Schur chart $a_{uv}\ne0$ では、
$R_{a_{uv}}=B'[p_1,p_2,p_3]$、$IR_{a_{uv}}=(p_1,p_2,p_3)$ と書ける。
このとき $q$ は $p$ の一次式で、$F$ の線形座標変換は $p$ に依存しない。
従って $H_j(q(a))$ は $p$ 次数4の斉次式であり、
その類は $(p)$-torsionである。

第17版の局所Čech--Koszul補題では、
$H^0_{(p)}(R_{a_{uv}}/G)$ の $p$ 次数4以上は零である。
従って $[H_j]$ はすべての $a_{uv}\ne0$ で消える。
有限個の座標についてのこの消失から、ある $\mathfrak m$ の冪が $[H_j]$ を殺す。
よって $[H_j]\in T_0$ である。

式\eqref{eq:v18-Hj-vertex-chain}は前補題と前定理による。
特に $q_j[H_j]=[H]\ne0$ なので、各 $[H_j]$ も非零である。
$\CC[F][H]$ は $a$ 次数10にしか成分を持たず、
$[H_j]$ は $a$ 次数8なので、これに属さない。
\end{proof}

\begin{remark}[対角の同変成分との違い]
$\mathcal K_n^{\mathrm{nor}}$ は二重次数 $(n-5,n)$ の
指定されたdeterminant指標の成分である。
上の $[H_j]$ は $(5,8)$ なので、そのまま $\mathcal K_8^{\mathrm{nor}}$ に
入るものではない。一方、$q_jH_j=H$ は $(5,10)$ に入り、
既知の $\mathcal K_{10}^{\mathrm{nor}}$ を与える。
従って全頂点加群のこの非自明な構造は、
$E_{\mathrm{vert}}=t^{10}$ の予想と両立する。
本節は他の対角成分が零であることも、その非零性も主張しない。
\end{remark}

\subsection{独立な$q$変数の完全交叉と、その適用限界}

\begin{proposition}[代入前の普遍勾配は完全交叉である]
\label{prop:v18-universal-q-CI}
$B=\CC[F_{ijk},q_0,\ldots,q_4]$ では $g_0,\ldots,g_4$ は正則列であり、
\begin{equation}\label{eq:v18-q-torsion-bound}
 \bigl(H^0_{(q)}(B/G_q)\bigr)_d=0\qquad(d>5)
\end{equation}
となる。添字は $q$ 次数を表す。
さらにこのtorsionは $F$ 次数5未満には現れない。
\end{proposition}

\begin{proof}
$q\ne0$ では $F\mapsto dF(q)$ は五次元への全射な線形写像である。
例えば $q=e_0$ とすれば、$x_0^3,x_0^2x_1,\ldots,x_0^2x_4$ の係数が
五つの独立な値を与える。
従って $q\ne0$ 上の零点集合は次元 $5+(35-5)=35$ である。
$q=0$ 上も次元35なので、40変数環 $B$ における $G_q$ の高さは5である。
$B$ はCohen--Macaulayであり、五つの生成元は正則列となる。

そのKoszul複体を $L_i=B(-i,-2i)^{\binom5i}$ と書く。
$q$ に関するČech複体との二重複体を考えると、
Koszul方向を先に取れば $H^\bullet_{(q)}(B/G_q)$ を計算する。
Čech方向を先に取ると、自由加群の局所コホモロジーは次数5だけにある。
全次数0の候補は $H^5_{(q)}(L_5)$ であり、
$H^5_{(q)}B$ の最大 $q$ 次数は $-5$ なので、
$L_5$ の $q$ シフト10を加えると最大次数は5となる。
また $L_5$ の $F$ シフトは5であり、$F$ 次数5未満は存在しない。
\end{proof}

\begin{remark}[非平坦な代入を越える際の未解決点]
$\CC[q]\to\CC[a]$、$q\mapsto q(a)$ は平坦ではない。
$q\ne0$ のファイバーは次元5である一方、零ファイバーは
七次元のPl\"ucker錐である。
従って式\eqref{eq:v18-q-torsion-bound}の $q$ 次数を単純に二倍して、
$T_0$ の $a$ 次数が10以下であると結論することはできない。
実際、代入後の五つの $g_i$ の高さは3となり、
Koszulの正次数homology $H_1,H_2$ を無視できない。
この非平坦性を、指定された同変対角成分で制御することが、
巨大な一般上界を小さい構造上界へ置き換える具体的な残課題である。
\end{remark}


% END theory/v18_vertex_hessian_chain_ja.tex


% BEGIN theory/v18_finite_determinacy_ja.tex
% 第18版：normal消滅を満たす普遍形式境界の有限決定性。
\section{normal消滅を満たす境界の明示的有限決定性}
\label{sec:v18-finite-determinacy}

全次数normal消滅と頂点核の有効上界を組み合わせると、
実際の境界の問題を、明示された有限jetの問題へ還元できる。
上界は極めて大きく、実用的な計算範囲を与えるものではない。
また本節の無限合成については、形式的な同値を証明し、収束や代数化は主張しない。

\begin{theorem}[normal消滅部分における形式的有限決定性]
\label{thm:v18-relative-finite-determinacy}
$N\ge10$ が
\[
 \mathcal K_n^{\mathrm{nor}}=0\qquad(n>N)
\]
を満たすとする。先頭項が $P$ の普遍同変形式級数
$d=P+\sum_{n>6}C_n$ が $\mathcal B_d=0$ を満たせば、
一次項が恒等の普遍同変形式座標変換 $\Phi$ によって
\begin{equation}\label{eq:v18-formal-finite-truncation}
 d(F,\Phi(F,a))=d_{\le N}(F,a)
\end{equation}
となる。右辺は変換前の $d$ の $N$ 次以下の切断である。
各形式係数は $F$ の多項式であり、単元変更を使わずに構成できる。
特に二つのそのような境界の $N$-jet が等しければ、普遍形式右同値である。
\end{theorem}

\begin{proof}
$n>N$ とし、既に $n-1$ 次まで $d_{\le N}$ に一致したとする。
それまでの座標変換後の方程式も
定理\ref{thm:v18-normal-contact-law}により $\mathcal B=0$ を満たす。
$d_{\le N}$ も同じ性質を持つので、差の $n$ 次項 $R_n$ は
\[
 \pi_n([R_n])=0
\]
を満たす。仮定 $\mathcal K_n^{\mathrm{nor}}=0$ により $R_n\in J_n$ である。
従って普遍共変場 $Y_{n-5}\in\mathcal G_{n-5}$ を
$dP(Y_{n-5})=-R_n$ となるように選べる。
$a\mapsto a+Y_{n-5}$ を代入すると、六次の先頭項 $P$ の一次微分が
ちょうど $n$ 次で $R_n$ を消し、それより低い次数は変えない。
他の係数への作用と $P$ の二階以上のTaylor項は全て $n$ 次より高い。

この操作を $n=N+1,N+2,\ldots$ に対して行う。
場の次数 $n-5$ は狭義に増加するため、合成は $a$-進位相で収束して
一次項が恒等の形式自己同型を与える。
各 $Y_{n-5}$ は有限次元の複素ベクトル空間における線形方程式を解いて選ぶため、
$F$ 係数の除算は必要ない。
有限jetに寄与する合成は有限個だけであり、各係数は $F$ の多項式のままである。
同じ $N$-jet を持つ二つの方程式をその共通の切断へ移し、
一方の形式自己同型を逆にすれば、最後の主張を得る。
\end{proof}

\begin{corollary}[実際の境界に対する一つの普遍上界]
\label{cor:v18-explicit-finite-boundary}
\begin{equation}\label{eq:v18-explicit-finite-degree}
 N_* =20\cdot18^{\,2^{43}}-10
\end{equation}
と置く。実際のTate境界は、その $N_*$ 次以下の切断に普遍形式右同値である。
また、第17版の偶正規形から出発すれば、普遍形式接触同値の下で
\begin{equation}\label{eq:v18-finite-even-normal-form}
 P+\frac1{96}H+
 \sum_{\substack{26\le n\le N_*\\n\text{ even}}}T_n,
 \qquad [T_n]\in\mathcal K_n^{\mathrm{nor}},
\end{equation}
という有限多項式にできる。各 $T_n$ は、あらかじめ選んだ
$\mathcal K_n^{\mathrm{nor}}$ の代表空間から取れる。
\end{corollary}

\begin{proof}
定理\ref{thm:v18-effective-vertex-cutoff}により
$n>N_*$ では $\mathcal K_n^{\mathrm{nor}}=0$ であり、
定理\ref{thm:v18-all-order-normal-vanishing}により実際の境界のnormal係数は全て零である。
従って前定理を適用できる。

偶正規形では、各偶数次数で通常像 $J_n$ の成分を同じ方法で除き、
核の代表だけを残す。$n$ が偶数なら場の次数 $n-5$ は奇数なので、
$a\mapsto a+Y_{n-5}$ は反転 $a\mapsto-a$ と可換であり、偶性を保つ。
核の代表を選ぶ操作は有限次元の線形代数でできる。
$n>N_*$ では代表空間は零なので、その後の項は全て消せる。
\end{proof}

\begin{remark}[有限決定性の範囲と残る問題]
この有限決定性は $\mathcal B=0$ を満たす普遍同変方程式の範囲内の主張である。
任意の形式摂動についての通常の有限決定性ではない。
自由な直和因子 $\mathcal A[\Theta_8]$ は任意の高次数にも現れ得るので、
その条件を除いた主張は本証明からは導けない。
また $N_*$ は計算済みの最小閾値ではなく、Gr\"obner基底の一般上界から得た安全な上界である。
\[
 E_{\mathrm{vert}}(t)=t^{10}
\]
が証明されれば、式\eqref{eq:v18-finite-even-normal-form}の和は全て消え、
実際の境界が $P+H/96$ と普遍形式接触同値であることが従う。
現在はこのHilbert級数の等式も、$\mathcal K_{26}^{\mathrm{nor}}=0$ も未確定である。
\end{remark}

% END theory/v18_finite_determinacy_ja.tex


% BEGIN characters/v18_character_notes_ja.tex
\section{二十六次の普遍指標と通常商の未確定部分}
\label{sec:v18-characters}

二十六次の普遍空間を、接触写像の階数やKoszul複体の正次数完全性を
仮定せずに計算する。通常商の自由成分には二つの既知の元があるが、
頂点支持成分 $K^{\mathrm{nor}}_{26}$ の消失は本節では確定しない。

\subsection{次数、標的表現、独立指標計算}

$\Delta=\det V$、$A=\Lambda^2V$、$W=A\otimes\Delta^{-1}$ と置く。
$F\in\operatorname{Sym}^3V$、$a\in W$、
$q(a)=\tfrac12a\wedge a\in V^*\otimes\Delta^{-1}$、
$P=F(q(a))\in\Delta^{-3}$ の規約を維持する。
二十六次のKoszul第 $k$ 項は
\[
 K_{k,26}=\operatorname{Hom}_{\mathrm{GL}(V)}
 \left(\operatorname{Sym}^{21-k}(\operatorname{Sym}^3V)
 \otimes\operatorname{Sym}^{26-4k}A,
 \Lambda^kV^*\otimes\Delta^{23-2k}\right),\qquad 0\leq k\leq5.
\]
$K_{0,26}=C_{26}$、$K_{1,26}=D_{22}$、$K_{2,26}=B_{26}$ と記す。
終端 $K_{5,26}$ にも $a$ 次数6が残る。

\begin{proposition}[二十六次の厳密次元]
\label{prop:v18-character-dimensions}
普遍空間の次元は次の通りである。
\[
\begin{array}{c|r|r|c|r}
\text{空間}&\deg_F&\deg_a&\text{標的分割}&\dim\\\hline
C_{26}&21&26&(23,23,23,23,23)&6377\\
G_{21}&20&21&(21,21,20,20,20)&32805\\
U_{20}&20&20&(20,20,20,20,20)&3712\\
D_{22}&20&22&(21,21,21,21,20)&17393\\
B_{26}&19&18&(19,19,19,18,18)&17043\\
K_{3,26}&18&14&(17,17,16,16,16)&7121\\
K_{4,26}&17&10&(15,14,14,14,14)&1151\\
K_{5,26}&16&6&(12,12,12,12,12)&52
\end{array}
\]
補助空間では $\dim C_{17}=31$、$\dim G_{12}=155$、
$\dim U_{11}=9$、$\dim D_{13}=92$ を再確認する。
\end{proposition}

\begin{proof}
三次単項式の35個の重み $\alpha$ からEuler積
$\prod_{|\alpha|=3}(1-tx^\alpha)^{-1}$ を作る。
その $t^m$ 係数に、$\rho=(4,3,2,1,0)$ を用いたWeyl交代和
\[
 \sum_{\sigma\in S_5}\operatorname{sgn}(\sigma)
 [x^{\lambda+\rho-\sigma\rho}]
 \operatorname{ch}\operatorname{Sym}^m(\operatorname{Sym}^3V)
\]
を適用してSchur多重度を得る。
外積二乗側も10個の重みから独立に指標を生成し、
テンソル積の交代和を直接評価する。
Littlewood分解
\[
 \operatorname{Sym}^d(\Lambda^2V)
 =\bigoplus_{j=0}^{\lfloor d/2\rfloor}S_{(d-j,d-j,j,j,0)}V
\]
の各多重度1と、Weyl次元和も照合した。

$C_{26}$ に現れる $F$ 側の分割は
$(23,20-j,20-j,j,j)$、$0\leq j\leq10$ である。
その多重度列は
\[
 (4,6,50,137,532,942,1699,1502,1170,280,55)
\]
となり、和は6377である。
$U_{20}$ の分割 $(20,20-j,20-j,j,j)$ に対する多重度列は
\[
 (2,2,19,54,255,493,983,923,756,187,38)
\]
となり、和は3712である。
$G_{21}$ と $D_{22}$ はそれぞれ外積第3冪・第1冪の垂直帯Pieri計算で、
Koszul各項は外積第 $k$ 冪のPieri計算で、直接のテンソル積指標と一致する。
Koszul源の非零Pieri成分数は $k=2,3,4,5$ の順に33、28、16、4である。

全ての $1\leq m\leq21$ の $F$ 指標と $1\leq d\leq26$ の $a$ 指標を
Newton--Adams再帰 $m h_m=\sum_{k=1}^m p_kh_{m-k}$ により別途生成し、
Euler積の全重み係数辞書と一致することを確認した。
必要な $F$ 次数 $11,12,16,17,18,19,20,21$ の全Schur分解は、
各多重度とWeyl次元の積の和が $\binom{m+34}{m}$ に一致する。
このうち前版にもある次数は、保存JSONの表現形式を正規化してから
全Schur分解を照合した。今回の全重み再計算と、旧証明書との保存値照合は
異なる検証として記録している。
\end{proof}

\subsection{Euler交代和が決める範囲}

今回の交代和は
\[
 6377-17393+17043-7121+1151-52=5.
\]
既証明の $H_i=0$（$i>2$）を適用すると
\[
 \dim H_0-\dim H_1+\dim H_2=5,
\]
および
\[
 \operatorname{rank}\delta_5=52,\qquad
 \operatorname{rank}\delta_4=1151-52=1099,\qquad
 \operatorname{rank}\delta_3=7121-1099=6022
\]
を得る。Pfaffian微分の全射性からは
\[
 \dim\ker(\delta q:G_{21}\to D_{22})=32805-17393=15412
\]
が従う。これらはいずれも $\operatorname{rank}(dP)$ の計算ではない。

$\kappa_{26}=\dim K^{\mathrm{nor}}_{26}$、$b_{26}=\operatorname{rank}\delta_2$ と置く。
既存の分裂
$Q=A_0[\Theta_8]\oplus K^{\mathrm{nor}}$、
$A_0=\mathbf C[\bar U_4,\bar J_6]$ を用いると
\[
 \dim Q_{26}=2+\kappa_{26},\qquad
 \dim J_{26}=6375-\kappa_{26}.
\]
従って
\[
 \dim H_1=11018+\kappa_{26}-b_{26},\qquad
 \dim H_2=11021-b_{26},\qquad
 \dim H_2-\dim H_1=3-\kappa_{26}.
\]
特に、交代和5から $\kappa_{26}=0$ や $H_1=0$ を推論することはできない。
$\kappa_{26}=0$ を別に証明できた場合に限り、
$\dim H_2=\dim H_1+3\geq3$ が従う。

\subsection{単元像、自由成分、必要な通常階数}

多項式環 $\mathbf C[F,a]$ における非零多項式 $P$ の乗法は単射であるから、
$PU_{20}\subset J_{26}$ は3712次元である。
この所属はEuler恒等式 $dP(a)=6P$ からも分かる。
通常商の自由成分の二つの元は
\[
 [\Theta_8U_4^3J_6],\qquad [\Theta_8J_6^3].
\]
これらの独立性は既存のnormal射影により確定している。
全通常商をこの二元が張ることは $\kappa_{26}=0$ と同値であり、
本版の有限rank計算では証明していない。

$P=0$ 上の評価で通常像の2663本の独立列を構成できれば、
\[
 \dim J_{26}\geq3712+2663=6375
\]
と理論上界 $\dim J_{26}\leq6375$ が一致する。
この条件を満たす全minorの新規構成・独立監査は今回は実施していない。
96点でのpilotの非零minorは、その有限の下界だけを与える。
小標本の階数飽和や候補列の枯渇を通常商の上界には用いない。

補助的に、$U_s$ 内の $P$ 可除部分を矩形補分割で計算した。
$P$ の係数空間は
\[
 \operatorname{Hom}_{\mathrm{GL}(V)}
 \left(\operatorname{Sym}^{s-1}(\operatorname{Sym}^3V)
 \otimes\operatorname{Sym}^{s-6}A,\Delta^{s-3}\right)
\]
であり、$F$ 側の分割は
$(s-3,s-3-j,s-3-j,3+j,3+j)$ となる。
保存された完全Schur分解から得る値は
\[
\begin{array}{c|r|r|r}
 &\dim U_s&\dim(P\text{ 可除部分})&\dim(U_s/(P\text{ 可除部分}))\\\hline
U_{13}&47&5&42\\
U_{16}&486&141&345\\
U_{20}&3712&1406&2306
\end{array}
\]
である。この表は単元の源空間の計算であり、候補グラフの基底性を主張しない。

\subsection{制限の標的次元と証明書の範囲}

Pl\"ucker錐へ制限したHom空間の次元は
\[
\begin{array}{c|rrrrrrrr}
 &C_{26}&G_{21}&U_{20}&D_{22}&B_{26}&K_{3,26}&K_{4,26}&K_{5,26}\\\hline
\dim&0&2&2&0&0&22&66&21
\end{array}
\]
である。これはLittlewood分解のCartan成分のみを残したPieri計算であり、
任意の普遍制限写像の全射性を仮定していない。

Fermat固定群の成分別計算による評価階数の上界は
\[
\begin{array}{c|r|r|r}
 &\dim&\text{対角値の階数上界}&\text{対角値と一階法jetの階数上界}\\\hline
C_{26}&6377&1332&4478\\
G_{21}&32805&4854&19814\\
U_{20}&3712&487&1991\\
D_{22}&17393&2834&11092
\end{array}
\]
となる。実際の評価階数との等号は主張しない。

\nolinkurl{v18_character_certificate.py} は全重みの新規整数計算、
\nolinkurl{v18_character_audit.py} は保存データの監査を担う。

\nolinkurl{v18_koszul_pieri_check.py} は別実装による垂直帯列挙を担う。
\nolinkurl{v18_unit_divisibility.py} の補助次数13・15の表には旧版の出典を残す。
通常商のrank pilotと選択された小minorの独立再縮約は
\texttt{rank/} 以下に分けて記録する。

% END characters/v18_character_notes_ja.tex


% BEGIN detector/v18_normal_detector_ja.tex
% 第18版：26次 normal 検出子と独立保存監査。
\section{二十六次normal検出子と厳密な零判定}
\label{sec:v18-normal26-detectors}

normal射影の26次の標的は
\[
 \mathcal A_{18}[\Theta_8]
 =\mathbb C\cdot[\Theta_8U_4^3J_6]\oplus\mathbb C\cdot[\Theta_8J_6^3]
\]
である。以下の二つの検出子は、この二次元成分を完全に読む。
通常商全体の次元は、この標的の次元だけからは決まらない。

$m=-1,-2$ に対して
\[
 F_m=x_0^3+x_1^3+x_2^3+x_3^3+(-2-6m)x_4^3+6mx_2x_3x_4,
 \qquad a_\varepsilon=e_{01}+\varepsilon(e_2-e_4)\wedge(e_3-e_4)
\]
とし、$L_m(C)=[\varepsilon^3]C(F_m,a_\varepsilon)$ と定める。
どちらも三元部分 $f_m$ と $p=(1,1,1)$ について $f_m(p)=0$ を満たす。
三元式の滑らかさの判別条件 $(-2-6m)+8m^3$ はそれぞれ $-4,-54$ で非零である。
第17版のnormal射影定理により、$L_m$ は $\mathcal J_{26}$ を消す。

\begin{proposition}[二つの有理normal検出子]\label{prop:v18-normal-jet-detectors}
基底を $W_1=\Theta_8U_4^3J_6$、$W_2=\Theta_8J_6^3$ とすると
\[
 \begin{array}{c|rrr}
 m&\overline U_4&\overline J_6&h_{f_m}(p)\\\hline
 -1&5&528&-4\\
 -2&36&7128&-54
 \end{array}
\]
であり、評価行列とその行列式は
\[
 M_{26}=
 \begin{pmatrix}
 -264000&-588791808\\
 -17958454272&-19556756702208
 \end{pmatrix},\qquad
 \det M_{26}=-5410806990313291776.
\]
従って
\begin{align*}
 \rho_{26,1}(C)&=
 \frac{531441L_{-1}(C)-16L_{-2}(C)}{147034844352},\\
 \rho_{26,2}(C)&=
 \frac{-8503056L_{-1}(C)+125L_{-2}(C)}{2561935127989248}
\end{align*}
は $\rho_{26,i}(W_j)=\delta_{ij}$ を満たし、normal射影の二係数を与える。
\end{proposition}
\begin{proof}
\texttt{v18\_normal\_detector\_certificate.py} は旧版のtensor graphの定義から
rank二切片上の $U_4,J_6$ を再評価し、三元Hessianを有理数演算で直接求める。
$U_4$ のgraph値は規約通り24で割る。
各行の値は $h_f(p)(\overline U_4^3\overline J_6,\overline J_6^3)$ であり、
逆行列を厳密に求めると上の二つの有理式になる。
出力には全入力テンソル、rank二点、法方向、旧入力とgraph定義のSHA-256を保存した。
\end{proof}

\begin{proposition}[二十六次の新規normal Tate計算と保存監査]
\label{prop:v18-normal26-tate-zero}
二入力で新規計算したraw Tate境界のnormal三次係数は
\[
 L_{-1}(d_{26}^{\mathrm T})=L_{-2}(d_{26}^{\mathrm T})=0.
\]
従って
\[
 \pi_{26}([d_{26}^{\mathrm T}])=0,\qquad
 \pi_{26}([R_{26}^{\mathrm{ev}}])=0.
\]
\end{proposition}
\begin{proof}
\texttt{v18\_tate\_normal\_jet.py} は
$\mathbb Q[\varepsilon]/(\varepsilon^4)$ 上でMaurer--Cartan再帰を最初から行う。
全 $A_1,\ldots,A_{18}$ と、曲率 $R_{19},\ldots,R_{22}$ の$x$次数
$4,3,2,1$ 以下を新規計算し、22個の定数ブロックを得た。
行列式の規約は $d^{\mathrm T}=\tfrac12\det K$ である。
二入力それぞれで、$a$次数0から26、$\varepsilon$次数0から3の
$27\times4=108$ 個の行列式係数は全て零となった。

別の \texttt{v18\_normal\_saved\_audit.py} は研究計算モジュールを一切importせず、
標準ライブラリの \texttt{Fraction} によって、保存ブロックから
置換展開と稠密な部分集合動的計画の二方法で全216係数を再構成した。
各入力の125成分の対称テンソル、$a_0$、法方向、$K_1=-a$ と
$K_2=2F(q,-,-)$ の等式、次数切断のmetadata、計算コードのhashも照合した。
最後に前命題の逆行列を適用して $(\rho_{26,1},\rho_{26,2})=(0,0)$ を確認した。
これは保存監査であり、新規再帰そのものは前記Tateスクリプトによる。
実際の26次係数との等式は系\ref{cor:v18-actual26-normal-equals-raw}による。
\end{proof}

\begin{remark}[結論の範囲]
ここでの零判定は $[R_{26}^{\mathrm{ev}}]\in\mathcal K_{26}^{\mathrm{nor}}$ を意味する。
通常像への所属には $\mathcal K_{26}^{\mathrm{nor}}=0$、通常26次の必要なrank、
またはこの一つの類の直接の消去を別に示す必要がある。
新規normal再帰と独立保存監査のいずれも、その追加主張を行わない。
\end{remark}

% END detector/v18_normal_detector_ja.tex

\clearpage
\section{今後の課題}

第18版では、実際の境界の横断Hessian係数が全次数で零であることを証明した。
従って全次数Hessian正規形の問題は、有限次元の頂点部分
$\mathcal K^{\mathrm{nor}}$ に集中する。
明示的上界 $N_*=20\cdot18^{2^{43}}-10$ は得たが、実用的な閾値ではない。
以下では、この理論的な有限還元を、小さい構造定理または有限証明書へ置き換える。
\begin{enumerate}
 \item 指定された同変頂点成分 $\mathcal K^{\mathrm{nor}}$ を決定する。
 最も重要な未解決式は $E_{\mathrm{vert}}=t^{10}$ である。
 これが成立すれば、今回の全次数normal消滅と合わせて
 実際の境界の普遍形式 $P+H/96$ 正規形が従う。
 全頂点加群 $T_0$ には $H_j=\partial_{q_j}H$ も存在するので、
 $T_0=\mathbb C[F]H$ という強すぎる簡約は使えない。
 独立な $q$ 変数での完全交叉と、非平坦なPfaffian代入の差を
 二重次数とdeterminant指標を保ったまま制御する。

 \item 二十六次の残る頂点類を決める。
 全指標とnormal射影は確定したが、通常26次商の二次元性は未証明である。
 $[R_{26}^{\mathrm{ev}}]\in\mathcal K_{26}^{\mathrm{nor}}$ の直接消去、
 小さい理論的上界、または十分な通常接触minorのいずれかが必要になる。
 96列の有限minorは部分的な下界証拠であり、必要な2663本の独立方向を与えていない。
 24次の核の消失を26次へ外挿せず、Euler交代和を完全性として扱わない。

 \item 全次数normal消滅の幾何学的証明をTate模型内でも説明する。
 split三次式に沿う定数ブロックの構造、Picard形式対数、Hesse pencilの一定な
 抽象曲線との対応を直接表す式を求める。
 今回の証明は全次数を扱うため、normal成分の零判定だけを目的とした
 次数ごとの大型計算は必要ない。
 ただし頂点成分を読むためには、法三次jetを超える計算が引き続き必要である。

 \item 普遍Koszulの $H_1,H_2$ を決める。
 高次 $H_i$ ($i>2$) の消失は既知であるが、二つの個別次元は未確定である。
 26次のEuler和は5であり、通常商の次元を単独で与えない。
 五つの勾配のheightは3なので、代入後も五元正則列であるとは扱わない。

 \item 有限形式正規形の収束・代数化と、表示の縮約を調べる。
 今回の $N_*$-jet への無限合成は形式的である。
 既知の有限次数の代数的・収束的正規化と、そのまま同一視しない。
 また4929補助変数と共役の5890追加未知数を含む既存表示を小さくし、
 $16\times5$ 行列によるincidenceと相対Picardの幾何を直接比較する課題が残る。

 \item 被約特異部分は命題\ref{thm:global-bn-singular}で $C$ と同定したが、
 Jacobian idealが定める非被約な特異スキームの構造を決定する。
 一段階blowup内でincidenceを定める $p^*\OO(3)$ の切断の束方向6次以上の項を求め、
 二段階解消の記述と比較する。
 相対Picardによるblowupの形式記述を大域的な代数モジュライへ延長すること、
 および $G$ の相対対象モジュライの分離性・propernessも別の課題である。
 \item 明示的族の線形束 $L_R$ のPicard座標を、Sklyaninの平行移動点 $p(h)$ と
 直接結ぶ公式を求める。
 現在は固定形式曲線を介した対応と、その一次係数が非零であることが分かっている。
 Abel--Jacobi座標での公式は、gluing対象と非被約点の局所環との関係を精密にする。
 \item 構造対象や次数シフトを指定しない、全体dg圏の同値類を区別する。
 第\ref{sec:five-generator}節の分類は次数付き代数と偏極付きtailsについてのものである。
 抽象point schemeからは四つの次数付き同型類が残る。
 六次境界は元の $X$ を復元するが、その境界を残余圏だけから内在的に指定できるかは
 別に検討する必要がある。gluing対象から非被約点の局所環への普遍的な変形写像も未構成である。
 \item 一般の三次式に対し、二つの中心パラメータを持つ有限表示の
 代数的量子化を構成し、残余圏へ写る二次項を明示する。
 point moduleの計算をline module・理想層へ拡張し、Fano surfaceとの関係を調べる。
 評価核から曲線と線形束を $F_L^{**}/F_L=j_*L$ として回復する構成を、
 他の数値類や非自由支持型へ広げることも考えられる。
\end{enumerate}

\clearpage
\appendix
\section{検証コードと再現方法}

第12版までのPythonコードと第13版の重み空間・Tate再帰は標準ライブラリを用いる。
第13版のテンソルグラフ縮約ではNumPyも用いる。
有限体、Python多倍長整数、有理数の厳密演算を使い、浮動小数点近似で階数を判定しない。

\begin{enumerate}
 \item \texttt{nc\_cubic\_verify\_central.py} は $\mathbb Q[h]\langle x,y,z\rangle$ 内で
 $[G_h,x],[G_h,y],[G_h,z]$ が関係式の両側イデアルに入る有限の証明書を生成する。
 抽出した証明書を別途展開し直して元の多項式と一致することも確認する。
 \item \texttt{nc\_cubic\_verify\_pfaffian.py} は5個の二次式のsyzygy、
 および $a_{01}=1$、他の $a_{ij}=0$ におけるJacobian階数3を厳密計算する。
 \item \texttt{nc\_cubic\_verify\_point\_curve.py} は多重線形化行列の行列式を展開し、
 \eqref{eq:elliptic} を検証する。
 \item \texttt{nc\_cubic\_verify\_hesse.py} はJacobian環の対称性、rotationの数値行列、
 Hesse曲線のWeierstrass変換による $j$-公式、固定曲線の補正を有理数上で検証する。
 \item \texttt{nc\_cubic\_verify\_exceptional.py} は $\mathbb Q(\omega)$ 上で
 Fourier変換・Zhang twist・中心式を検証し、特異性を検出する階数4のMFの
 $d^2=W I$ を全成分で検証する。
 \item \texttt{nc\_cubic\_verify\_points.py} は全非零パラメータでの階数2を保証する
 小行列式、長さ8の局所代数の積表、socleの三次形式、中心三次式の消去、
 および非被約点の $j$-不変量を厳密に検証する。
 \item \texttt{nc\_cubic\_verify\_reconstruction.py} はHessianの表示、
 二つの $j$ の逆数関係、次数2・4の復元式、
 3-isogenyの商曲線とWeierstrass不変量、$X_0(3)$ の二つの $j$-写像、
 Fricke involution、形式展開および一般のSklyanin復元式を検証する。
 \item \texttt{nc\_cubic\_verify\_v5.py} は普遍的な双ベクトルの表示の
 Euler-horizontal性と外積恒等式、写像の階数126、
 traceを含めたFermatの固定接方向、および $j_{\mathrm{fat}}=432$ の
 復元多項式の因数分解と判別式を厳密に検証する。
 \item \texttt{nc\_cubic\_verify\_v6.py} は符号付きPfaffianのkernel恒等式、
 三次式のPfaffian代入の係数行列の階数35、
 $GL(5)$ の29族の生成元についての共変性、および六次Euler恒等式を検証する。
 係数行列は全5005行・35列を対象とし、零でない660行を疎形式で保持する。
 \item \texttt{nc\_cubic\_verify\_v7.py} は35個の三次単項式の重みとWeyl alternantを使い、
 6次から12次までの普遍共変式空間の次元、座標変更とunit変更の次元、
 および10次に関係する $\delta q$ の表現の重複度1を整数で検証する。
 Fermatの8次候補の非零性、有限対称群に対するcharacter、固定軸上の横断式、
 Hesseの $j$-不変量の先頭係数も検証する。
 またFermatの10次Hessian候補について、全2002個の5次単項式から20個の
 characterに適合する候補を選び、Reynolds投影によるJacobian像の階数2、
 Hessian候補を加えた階数3、および非所属を示す一つの単項式の係数を検証する。
 このコードは実際の境界の10次係数を決定するものではない。
 \item \texttt{nc\_cubic\_verify\_v8.py} は、対角三次式のClifford模型で
 Maurer--Cartan方程式を6次まで有理多項式として検証する。
 本文の $5\times5$ 行列式の10次係数 $-648$、正規化比 $1/96$、
 Fermatの別標本、および六次先頭係数2を再計算する。
 32次元外積空間の作用素による独立の積の実装も照合し、
 \texttt{nc\_cubic\_v8\_certificates.json} に行列と厳密な値を保存する。
 \item \texttt{nc\_cubic\_verify\_v9.py} は変更していないv8のClifford演算を用い、
 再帰をblock次数8、行列式を次数12まで延ばす。
 五つの標本で有理多項式としてのMaurer--Cartan方程式と独立な32次元作用素を照合し、
 十二次の混合行列式係数および $dP(\operatorname{Alt}K_7)$ を検証する。
 完全な行列と値を \texttt{nc\_cubic\_v9\_certificates.json} に保存する。
 \item \texttt{nc\_cubic\_v9\_finite\_tate\_dimensions.py} は、
 有限係数空間の5890、homotopyの像の4929、および曲率空間の20330を、
 de Rhamの再帰・hookの次元式・Koszul複体の階数という三つの方法で照合する。
 \item \texttt{v9\_verify\_contact\_bounds.py} はv7の整数character計算を用い、
 $\mathcal G_7,\mathcal D_8,\mathcal U_6,\mathcal C_{12}$ の次元とPieri表を検証する。
 二つの最高重みベクトルの像の係数3と8を厳密展開し、保存された全五標本について
 $K_7$ の交代化、Pfaffian微分、および $dP(\phi_7)$ を独立に再計算する。
 \item \texttt{v10\_rank\_certificate.py} は
 $\Psi_7=\operatorname{Alt}\kappa h(A_1h(A_3^2))$ を有理数で再計算し、
 $\Phi_7$ と合わせた評価行列の行列式 $147288/25$ を検証する。
 再帰の3次までを独立な外積作用素でも照合する。
 \item \texttt{v10\_jet\_contact\_certificate.py} は、四次単元の対角重みから
 五つの星型単項式を列挙し、正規化を錐外の標本で確認する。
 $\rho(\Theta_8U_4)=1$、$\rho(d^{\mathrm T}_{12})=0$、
 および二つの普遍座標変更の係数を厳密計算する。
 \item \texttt{v10\_degree13\_certificate.py} は13次のSchur重複度、
 $\mathcal G_8$、$\mathcal U_7$ の次元を再計算し、
 全五標本で $dP(\operatorname{Alt}K_8)$ を独立に照合する。
 \item \texttt{v10\_jet\_thirteenth.py} は標準点でTate再帰を $A_9$ まで計算し、
 有理多項式でのMC等式、既存の全ブロックとの一致、独立な外積作用素を検査する。
 $[t^{13}]\det K=9200/9$ と普遍係数 $-4600/2189$ を再現する。
 第10版の四つのスクリプトはそれぞれ同名の \texttt{.json} 証明書を生成する。

 \item \texttt{v11\_character\_certificate.py} は14次の四つの次元
 $15,28,7,17$、Schur重複度、六つの最高重みの像と $\delta q$ の全射性を検証する。
 \texttt{v11\_fermat\_scope\_certificate.py} は有限Fermat安定化群の
 指標平均を $\mathbb Q(\zeta_3)$ で計算し、対角制限の核と法方向一階微分の範囲を検証する。
 \item \texttt{v11\_filtered\_koszul\_certificate.py} は三つの自然な縮約を
 整数だけで計算し、独立なsyzygyの小行列式 $-1990656$ を再現する。
 \texttt{v11\_tetrahedral\_u4.py} は四面体縮約、四次補間による方向微分、
 非対角三次式を生む剪断について $U_4$ の公式を検証する。
 \item \texttt{v11\_contact\_rank\_certificate.py} は七つの像の評価行列と
 その非零行列式を再検証する。\texttt{--recompute} を付けると、
 補助モジュール \texttt{v11\_dual\_rank.py} によってTate再帰と
 非対角三次式の一階微分をキャッシュなしで構成し直す。
 有理数でのMC等式、独立な外積作用素、使用ブロックの $x=0$ での作用を検証する。
 \item \texttt{v11\_actual\_fourteenth.py} は標準点の $A_{10}$ までの完全な
 Tate再帰と独立作用素を検証する。
 \texttt{v11\_fourteenth\_coefficients.py} は五点で $A_8$ までを新たに構成し、
 厳密な $x$ 次数切断から $K_9,K_{10}$ と14次係数を求める。
 標準点では完全な再帰の全ブロックと一致することを確認する。
 \item \texttt{v11\_degree14\_complement.py} は六頂点の単元 $J_6$ を直接計算し、
 五点のPfaffian微分、七つの画像基底、$\rho_{14}(\Theta_8J_6)=1$、
 $\rho_{14}(H_{10}U_4)=\rho_{14}(d^{\mathrm T}_{14})=0$ を有理数で検証する。
 \texttt{--recompute-fields} では、階数証明の四点に加えて必要な第五点の
 ベクトル場をTate再帰から新たに作り直す。
 第11版の八つのJSON証明書は、上記スクリプトに対応する入力・行列・結果を保存する。

 \item \texttt{v12\_character\_certificate.py} は15次の次元
 $7,35,1,22$、十二個のPieri成分、八個の最高重みの像、低次の安定化空間、
 Koszul源の次元 $19,3$、Fermat制限の上界を厳密に検証する。
 \item \texttt{v12\_filtered\_contact\_certificate.py} はC型の最高重みの係数
 $v+2$ を $1\le r\le12$ の42成分で整数多項式として検算する。
 また既存の五点で $\Phi_6,\Psi_{10}$ のPfaffian持ち上げ、
 $dH(\Phi_6)=-dP(\Psi_{10})$、$\operatorname{Alt}K_6=-\Phi_6/9$ を検証する。
 これらは本文の全次数の証明に対する有限の補助検算である。
 \item \texttt{v12\_contact\_rank\_certificate.py} は六つの非単元方向の
 全ての交代行列・Pfaffian微分・接触像、および非零な $6\times6$ 行列式を監査する。
 \texttt{--compare-fresh} は新たに計算した三つの通常評価と三つの一階変分との完全一致を確認する。
 \texttt{--recompute} は同じ六行をキャッシュなしで再構成する。
 \item \texttt{v12\_fifteenth\_tate.py} は十一個のブロックから15次の行列式を求める。
 \texttt{--full} は標準点で切断しない $A_1,\ldots,A_{11}$ を作り、
 有理数MC等式・独立な32次元作用素・正確な $x$ 切断との一致を検証する。
 \texttt{--audit} は保存済みの十一個のブロックから係数を再計算する。
 \item \texttt{v12\_fifteenth\_dual.py} は三つの一階変分を新たに構成し、
 両項を含む接触像の方向微分を有理数で検証する。
 使用した全ブロックは独立の外積・縮約作用素とも照合する。
 第12版の六つの証明書は、次元、Pfaffian持ち上げ、六行のコンパクトな評価データ、
 完全な標準再帰、通常評価、一階変分を別々のJSONに保存する。

 \item \texttt{v13\_weight\_space\_certificate.py} は内部重み空間の次元
 $0,50,1002,5890,20330$ と、共役方向の次元951を検証する。
 全6942基底でcontraction、$\delta^2=h^2=0$ を確認し、二つの微分の階数50、951を
 有理数の疎消去でも求める。有限方程式のJacobian同型を使う代数化の証明は本文による。
 \item \texttt{v13\_character\_certificate.py} は16次の四空間の次元
 $40,103,59,18$、Koszul源の次元 $22,2,0$、十六個のPieri成分、
 対角制限の階数上界を整数指標で検証する。
 \item \texttt{v13\_tensor\_contractions.py} は、NumPyの\texttt{object}配列を使い、
 多倍長整数と有理数によるグラフの縮約を実装する。
 \texttt{v13\_tensor\_rank\_certificate.py} は二十個のKoszulベクトル、
 二十一個の非単元接触像、補元、および $H\mathcal U_6$ の所属の正確な証明書を監査する。
 \texttt{--recompute} は全入力で全グラフを再計算し、自由脚から場を復元する方法とも照合する。
 \item \texttt{v13\_tate\_sixteenth.py} は完全な $A_1,\ldots,A_8$ と正確な
 $x$ 次数切断から $K_9,\ldots,K_{12}$ を求める。
 曲率の閉性、Maurer--Cartan等式、独立な外積作用素を照合し、
 16次の行列式と低次正規化から生じるHessian微分を同時に計算する。
 標準点のJSONは既存15次までの係数との一致も含む。
 \item \texttt{v13\_diagonal\_rho\_certificate.py} は十二点の評価行列と
 整数重みを使い、$\rho_{16}(J_{16})=0$ と
 $\rho_{16}(\Theta_8U_4^2)=1$ を有理数で確認する。
 \texttt{--recompute} は二十一個のグラフを全十二点で縮約し直す。
 \item \texttt{v13\_actual\_sixteenth\_certificate.py} は、十二点全てのブロックを
 統合した同名のJSONから、各行列式・Pfaffian微分・Hessian微分を検算し、
 未正規化係数、補正項、正規化後の代表元それぞれに $\rho_{16}$ を適用する。
 \texttt{--merge} は下記の新規再計算による二群と標準点を統合する。
 Tateデータには統合JSONを用い、rho証明書とも照合する。
 中間の二群や非公開キャッシュには依存しない。

\end{enumerate}

実行例（左列、右列の順）：
\begin{quote}\footnotesize\ttfamily
\begin{tabular}{@{}ll@{}}
python3 nc\_cubic\_verify\_central.py & python3 nc\_cubic\_verify\_v7.py \\
python3 nc\_cubic\_verify\_pfaffian.py & python3 nc\_cubic\_verify\_v8.py \\
python3 nc\_cubic\_verify\_point\_curve.py & python3 nc\_cubic\_verify\_v9.py \\
python3 nc\_cubic\_verify\_hesse.py & python3 nc\_cubic\_v9\_finite\_tate\_dimensions.py \\
python3 nc\_cubic\_verify\_exceptional.py & python3 v9\_verify\_contact\_bounds.py \\
python3 nc\_cubic\_verify\_points.py & python3 v10\_rank\_certificate.py \\
python3 nc\_cubic\_verify\_reconstruction.py & python3 v10\_jet\_contact\_certificate.py \\
python3 nc\_cubic\_verify\_v5.py & python3 v10\_degree13\_certificate.py \\
python3 nc\_cubic\_verify\_v6.py & python3 v10\_jet\_thirteenth.py \\
\end{tabular}
\end{quote}

第11版の追加検証は、次の順で実行できる。
\begin{quote}\footnotesize\ttfamily
\begin{tabular}{@{}l@{}}
python3 v11\_character\_certificate.py\\
python3 v11\_fermat\_scope\_certificate.py\\
python3 v11\_filtered\_koszul\_certificate.py\\
python3 v11\_tetrahedral\_u4.py\\
python3 v11\_contact\_rank\_certificate.py --recompute\\
python3 v11\_actual\_fourteenth.py\\
python3 v11\_fourteenth\_coefficients.py\\
python3 v11\_degree14\_complement.py --recompute-fields\\
\end{tabular}
\end{quote}
\texttt{v11\_contact\_rank\_certificate.py} と
\texttt{v11\_degree14\_complement.py} は、オプションを省略すると保存済みの
有理数証明書を監査する。再帰の全再計算には上記オプションを用いる。


第12版の新規計算と監査は次の順で再現できる。
\begin{quote}\footnotesize\ttfamily
\begin{tabular}{@{}l@{}}
python3 v12\_character\_certificate.py\\
python3 v12\_filtered\_contact\_certificate.py\\
python3 v12\_fifteenth\_tate.py --full --indices 0\\
\quad --output v12\_fifteenth\_full.json\\
python3 v12\_fifteenth\_tate.py --indices 1 2\\
\quad --output v12\_fifteenth\_tate.json\\
python3 v12\_fifteenth\_dual.py --indices 0 1 3\\
\quad --output v12\_fifteenth\_dual.json\\
python3 v12\_contact\_rank\_certificate.py --compare-fresh\\
\end{tabular}
\end{quote}
字下げしたオプションは直前のコマンドと同じ行に続ける。
接触階数だけを再構成する場合は
\texttt{python3 v12\_contact\_rank\_certificate.py --recompute} でもよい。
完全再計算は保存済みの私的な再帰キャッシュを使用しない。
\texttt{--audit --output v12\_fifteenth\_full.json} および
\texttt{--audit --output v12\_fifteenth\_tate.json} は行列式の高速監査に使う。

第13版の再現は次の順で行える。テンソル縮約にはNumPyを用いる。
今回の実行環境はPython 3.12.14、NumPy 2.3.5である。
\begin{quote}\footnotesize\ttfamily
\begin{tabular}{@{}l@{}}
python3 v13\_weight\_space\_certificate.py\\
python3 v13\_character\_certificate.py\\
python3 v13\_tensor\_rank\_certificate.py --recompute\\
python3 v13\_diagonal\_rho\_certificate.py --recompute\\
python3 v13\_tate\_sixteenth.py --indices 0\\
\quad --output v13\_tate\_sixteenth.json\\
python3 v13\_tate\_sixteenth.py --manifest\\
\quad v13\_diagonal\_rho\_certificate.json --indices 1 2 3 4 5\\
\quad --no-fields --output v13\_tate\_rho\_group\_a.json\\
python3 v13\_tate\_sixteenth.py --manifest\\
\quad v13\_diagonal\_rho\_certificate.json --indices 6 7 8 9 10 11\\
\quad --no-fields --output v13\_tate\_rho\_group\_b.json\\
python3 v13\_actual\_sixteenth\_certificate.py --merge\\
\end{tabular}
\end{quote}
字下げしたオプションは直前のコマンドと同じ行に続ける。
二群のTate再計算は出力先を分けた別プロセスで並行実行できる。
配布ZIPには統合済みの十二点データを含み、中間の二群は含めない。
高速監査には \texttt{--recompute} や \texttt{--merge} を外し、
Tate標準点には \texttt{--audit --output v13\_tate\_sixteenth.json} を用いる。
旧32個のPythonと21個のJSONは保持し、第13版の7個のPythonと6個のJSONを追加した。

16次の普遍接触階数は、一般三次テンソルによる上下界を別々に証明して決定した。
十二点の対角評価は、その後で一つの商の係数を読み取るために用いる。
全次数の偶化の代数性・収束性は、5890方程式とそのJacobianの同型から証明しており、
有限個のTate係数の一致から推測しているものではない。

Hilbert級数の低次数の数値一致は補助検査である。
全次数の平坦性の証明は、Sklyaninの定理とmonicな三次式による加群分解であり、
有限次数までの計算で代用していない。
PDFの再生成にはLuaLaTeX、LuaTeX-ja、Harano Ajiフォントを用いる。
計算コードが証明するのは有限の代数的恒等式である。
dg圏の同値・生成性・smoothnessは本文の証明によるものであり、数値検査の結論ではない。
第8版以降の行列コードは、対角三次式への特殊化の厳密計算を再現する。
第11版では非対角三次式の一階変分も独立なChevalley記号の外積作用素で検証する。
第10版の12次の消去は、先に商を1次元と証明し、その係数汎関数を構成することで、
有限個の厳密な評価を普遍的な多項式恒等式へ結び付けている。
第11版の14次も、普遍階数を先に決定してから五点の係数汎関数を適用する。
第12版の15次は、七次元の候補空間に対する七個の独立な接触像で全射性を示す。
全次数のPfaffian微分の全射性と境界の偶性は本文の一般的な証明に依拠し、
有限個のTate係数から推測しているわけではない。
対角評価だけによる普遍次元の推測は用いていない。
第9版の有限代数化と収束は全次数の証明であり、有限次数の検算で代用していない。
共変式の次元、Tate族と実際のgluing対象の同定、および全ての滑らかな $F$ での
Jacobian非所属は本文の数学的証明に依拠する。
同様に、第7版の普遍座標の存在、proper imageの性質、Picard族の不変性、
および解消の構成は、有限の表現計算だけからは従わず、本文の幾何的証明を用いる。

\clearpage
\section{第14版の計算証明書と再現範囲}
\label{app:v14-reproducibility}

第14版の追加ファイルは継続用bundleの \texttt{source\_v14/} に置き、
添付由来の69ファイルを含む \texttt{source\_v13/} は変更していない。
新規ファイルは \texttt{characters/}, \texttt{rank/}, \texttt{contact/},
\texttt{detector/}, \texttt{tate/} に分けた。
日本語原稿はLuaLaTeXで組版する。

指標計算は、Euler積による整数の重み計算をNewton--Adams再帰と照合した。
通常接触階数は、一般の整数三次テンソルでの94方向と60方向の独立性に基づく。
両者は別の辺番号付け・縮約方式でも監査し、選択minorの全成分を有理数で再計算した。
五点汎関数は300個の通常接触値と80個の追加方向の値を新規に縮約して作った。
これらの零判定は有理数で行っている。

実際の18次係数は、公開した五点それぞれからTate再帰を初めから作った。
標準点では通常のFraction演算と共通分母整数演算を照合した。
全ての点で、閉性と切断されたMaurer--Cartan等式、
独立な定数ブロック抽出、二方式の行列式係数計算を確認した。
この新規再構成と、保存した行列・値を読む高速監査は別の操作として記録する。

bundleのルートで次を実行すると、入力を一時コピーして保存証明書を監査する。
\begin{verbatim}
python3 -B verification/run_v14_checks.py
\end{verbatim}
指標の全再計算、選択テンソル縮約の独立な再計算、またはTateの新規再構成には、
同じスクリプトの \texttt{--help} にある追加オプションを用いる。
検算にはPython 3.12.14とNumPy 2.3.5を用いた。
計算は整数・有理数・明示された有限体で行い、浮動小数点の近似階数判定を使わない。

原稿の幾何学的な同定と既存定理は引き続き証明の入力である。
有限計算の成功は、原稿全体の全ての幾何学的・圏論的議論を独立に再認証したことを意味しない。
また、文献の新規性や優先権の全面調査は本版の作業範囲に含めていない。


\section{第15版の証明書・保存範囲・再現手順}
\label{app:v15-reproducibility}

第15版の追加研究は継続用bundle v3の \texttt{source\_v15/} に収めた。
添付由来の \texttt{source\_v13/} と、前回の \texttt{source\_v14/} は保存し、
旧原稿も比較と再構成の基準として保持した。
第15版の日本語TeXは単体でLuaLaTeXにより組版でき、
基準となる第14版と追加章からの再組立てにも対応する。

今回の証明の入力は次のように区別する。
指標は $F$ 次数15まで全重みをEuler積とNewton--Adamsで新規に照合した。
接触部分空間の164方向とKoszul核の359方向は、整係数テンソル公式、
公開整数入力、法1009の非零小行列式により特性零の独立性を証明した。
採用した全縮約は別の辺番号付けでも再計算している。
三点の164接触列と二つの補元、14個と1個の二階補正は有理数で再構成した。
最後に、三点のTate係数も保存値の補間によらず初項から再帰で求めた。

\nolinkurl{verification/run_v15_checks.py} は入力を一時コピーして監査し、
JSONとテキストの結果を作る。
\begin{verbatim}
python3 -B verification/run_v15_checks.py
\end{verbatim}
この既定実行は保存証明書の監査を中心とする。
全重み、テンソル縮約、またはTate再帰そのものの新規再計算は、
同スクリプトの \texttt{--help} に表示される追加オプションで区別して指定する。
Pythonのassertionを無効化する最適化オプションは使わない。
依存はPython標準ライブラリとNumPyであり、
検算時の版はPython 3.12.14、NumPy 2.3.5である。

高速監査の成功を、再帰を初めから作り直した結果と同一視しない。
また、第15版は通常20次商の次元を1または2という範囲に留めており、
その完全同定を実際の20次消去の前提にしていない。
幾何学的同定と既存定理は引き続き数学的入力であり、有限計算の成功が
既存原稿の全ての圏論的・幾何学的議論の独立な再認証を意味するわけではない。
文献の新規性・優先権についての全面調査も、今回の計算研究の主張には含めていない。


\section{第16版の証明依存・再現手順・継続点}
\label{app:v16-reproducibility}

第16版の追加研究は継続用bundle v4の \texttt{source\_v16/} に収めた。
第13版、第14版、第15版の研究ソースはバイト単位で保存し、
各版のTeX、PDF、証明書と前版の索引も比較可能な形で保持した。
第16版のTeXは単体でLuaLaTeXにより組版でき、
変更していない第15版と追加章からの再組立ても検査した。

\paragraph{構造証明.}
Pl\"ucker素イデアル、普遍方向の錐上消失、adjugate恒等式、
単元制限環、横断三次項による単射、普遍勾配のheightを順に扱った。
三元三次不変式環の生成定理とCohen--Macaulay環の正則列判定は
明示した一次資料を用いる。
新規性・優先権の全面的調査を実施したという主張は含めない。
記号恒等式の検算例と、証明そのものの論理的入力は本文で区別した。

\paragraph{有限計算の証明書.}
二十二次の指標は、$F$ 次数17、$a$ 次数22までの全重みを
Euler積とNewton--Adams再帰で照合した。
通常像の429方向は公開整数入力上の非零minorで独立性を証明し、
選択した全 $429^2$ 成分を別の縮約実装でも再構成した。
全429方向の二点値、143個の積方向の因数分解、76個の二階補正を
有理数で検算した。
実際のTate値は四入力五回の新規Maurer--Cartan再帰による。
保存ブロックからの行列式再計算と、再帰そのものの再実行は区別する。

\paragraph{再現の入口.}
\begin{verbatim}
python3 -B verification/verify_integrity.py
python3 -B verification/run_v16_checks.py --help
python3 -B verification/run_v16_checks.py
\end{verbatim}
既定の監査と、全指標・全縮約・Tate再帰の再計算はオプションで分けた。
既定の成功が全ての重い再帰をその場で再実行したことを意味しない。
assertionを有効にして実行し、Pythonの最適化オプションは使わない。
入力を一時コピーして検査するため、配布する確定済みソースは変更されない。
依存するPython、NumPy、および日本語LuaLaTeXの条件は同梱の手順に示した。

\paragraph{次回に引き継ぐ境界.}
現在確定しているのは24次開始の偶正規形、20次商2次元、22次商1次元、
二項下位jetに対する全次数の $I_N=J_N$、および自由部分加群の単射である。
全商の同定、24次以降の実際の消滅、残る第2Koszul微分の階数は未決定である。
研究状態、主張台帳、本文の命題・ページ索引、コード索引、
次回の具体的な手順をbundleに収めた。


\section{第17版の全次数構造と二十四次計算の再現範囲}
\label{app:v17-reproducibility}

第17版では、通常商の構造を証明する四つの章と、二十四次を確定する六つの章を加えた。
継続用bundleは第5版であり、原稿の第17版とは版番号を区別する。
\nolinkurl{source_v13/}、\nolinkurl{source_v14/}、
\nolinkurl{source_v15/}、\nolinkurl{source_v16/} は不変の基準として保持した。
今回の新しい脚本、入力、結果、証明原稿は \nolinkurl{source_v17/} に置く。

\subsection{構造理論の証明依存}

行列束とadjugateの恒等式は、対称行列15変数と交代行列10変数を独立変数とした
整数係数多項式の恒等式として検証する。
三元三次式の全共変量加群についてはFisherの
\cite[Theorem~2.7(i)]{FisherHessian2012}を用い、判別式を可逆にしない全係数環上で
normal射影を構成する。
Heisenberg作用のcommutantとHesse族の優勢性は、階数8と行列式118098で確認する。
頂点支持の証明は、\v{C}ech--Koszul双複体の二つの収束するスペクトル系列を比較する。
有限次元性はNoether有限生成性と指定二重次数から従う。
これらの数学的推論を、有限標本の行列の階数から導いたとは扱わない。

\subsection{二十四次の証明書}

全指標計算は $F$ 次数19、$a$ 次数24までの重みをEuler積とNewton再帰で照合し、
Schur--Weyl型の交代和と独立なPieri列挙から各次元を求める。
接触像では1120個の有理点を使い、1365次元の $PU_{18}$ を消した上で
1091本の通常接触方向の非零minorを構成する。
別実装は905個の相異なる因子を新規に縮約し、選択小行列の1190281成分を全て照合する。
保存された因子から全行列を再構成する軽量監査も用意した。
軽量監査と新規テンソル縮約は、実行記録上も区別する。

通常商の上界にはnormal sliceの独立性を使う。
有限体で零になった値、探索の停滞、KoszulのEuler和だけを上界に用いない。
二階補正の所属は全次数のnormal射影零と
$\mathcal K_{24}^{\mathrm{nor}}=0$ を組み合わせて証明する。
四入力での310本の補正の有限評価は追加の照合であり、
そのうち二本が点ごとには非零であることも記録した。

\subsection{Tate計算の実施範囲}

四つの完全な点入力で五回の新規再帰を行った。
このうち一入力は、共通分母整数演算とは別の\texttt{Fraction}演算でも再帰し、
全20定数ブロックと24次以下の全25行列式係数を比較した。
さらに二つの入力について、$\mathbb Q[\varepsilon]/(\varepsilon^4)$ 上で
法三次jetを最初から生成した。
この二回は、点での値から補間した計算ではない。
法三次再帰の保存値監査は独立な\texttt{Fraction}実装で全200係数を再構成し、
入力テンソル・交代行列・実行脚本のhashも確認する。

有理数体上のMaurer--Cartan検査、有限体上の独立な32次元外積作用素検査、
二つの行列式算法による照合を区別して記録する。
過去の切断済み $A_{15},A_{16}$ を全項の再帰として再利用していない。
二十六次へ進む場合は、全 $A_1,\ldots,A_{18}$ と、
曲率 $R_{19},\ldots,R_{22}$ の内部$x$次数 $4,3,2,1$ 以下が新たに必要になる。

\subsection{同梱された入口}

bundleを展開したディレクトリで次を実行する。
Pythonのassertを有効に保ち、\texttt{-O}、\texttt{-OO}を用いない。
\begin{verbatim}
python3 -B verification/verify_integrity.py
python3 -B verification/run_v17_checks.py
python3 -B verification/run_v17_checks.py --fresh-normal-jets
\end{verbatim}
既定の10項目は、指標保存監査、KoszulのPieri計算、通常像の保存監査、
構造恒等式、24次台帳、G12基底の新規厳密計算、normal検出子、
下位補正の保存監査、完全点Tate保存監査、法三次Tate保存監査である。
入口は五つのソース木を一時コピーし、実行前後の元ファイルのhashを比較する。
\begin{verbatim}
python3 -B verification/run_v17_checks.py --full-characters
python3 -B verification/run_v17_checks.py --recompute-graphs
python3 -B verification/run_v17_checks.py --recompute-detectors
\end{verbatim}
これらのオプションは、それぞれ全重み、接触minorの独立新規縮約、
四入力の310補正の新規計算を追加する。
完全点のTate再帰を最初から実行するコマンドは
\nolinkurl{source_v17/README_ja.txt} に記す。

\subsection{到達点と残る課題}

通常商は自由な $\Theta_8$ 成分と有限次元の頂点部分に分裂した。
しかし有限補正多項式 $E_{\mathrm{vert}}$ の全係数と最高次数、
自由部分と一致し始める閾値は未計算である。
二十四次の確定から二十六次の頂点部分の消失を外挿しない。
二十六次の既知の自由方向は
$[\Theta_8U_4^3J_6]$、$[\Theta_8J_6^3]$ である。
実際の係数については、raw Tate項、正規化後の項、normal射影、頂点部分を区別して調べる。
全次数でのHessian正規形と、残る $H_1,H_2$ の個別次元は未決定である。
本版は、新規性と優先権についての網羅的文献調査を行ったとは主張しない。


\section{第18版の証明依存と再現範囲}
\label{app:v18-reproducibility}

本稿の第18版に対応する継続用bundleは第6版である。
第13版から第17版までのソース木、原稿、証明書と索引は基準資料として保持し、
今回の追加を \nolinkurl{source_v18/} にまとめる。
過去の節の「未決定」という記述は当該版の到達点を表し、
現在の到達点は冒頭、今回の七つの節、および今後の課題に記載する。

\subsection{新しい証明の入力}

全次数normal消滅の証明は、既存の固定平面のPicard直線とnormal coneの幾何を入力とする。
その上でHesse pencilの形式一意性を証明し、全多項式共変量へ延長する。
新しい多項式監査はHessian、Hesseパラメータの変更式、$j$ の有理式とその微分、
および優勢性を示す十行十列行列を確認する。
既存の幾何学的証明全体を有限コードで再認証したという意味ではない。

接触変換則はSchur chart、横断次数の重み、三元End共変量のスカラー性から導く。
有限次数上界にはDub\'eの斉次版定理\cite[Theorem~8.2]{Dube1990}を用いる。
45変数全体を保ち、$F$ を係数体に移していない。
Hessianのsocle恒等式はEuler式、adjugate、$aq=0$ の三つの普遍恒等式から導く。
その微分の頂点支持には第17版のSchur chart上のČech--Koszul次数評価を使う。
これらの記述証明と、有限多項式の監査記録は別の役割を持つ。

\subsection{26次計算の実施範囲}

全指標は $F$ 次数21、$a$ 次数26までのEuler積とNewton--Adams再帰を照合する。
Schur重複度、Pieri列挙、全Weyl次元和、および旧版との重複を確認する。
通常26次商を決める完全な接触minorは今回構成していない。
96列の部分minorの探索とその独立監査は、完全な階数証明と区別する。

法三次Tate計算は二入力について最初から行い、
全 $A_1,\ldots,A_{18}$ と必要な切断 $A_{19},\ldots,A_{22}$ を再構成する。
旧版で切断した $A_{17},A_{18}$ を流用していない。
独立な保存監査は二つの行列式算法を使い、二入力の全216係数を再構成する。
通常の点における26次Tate係数の完全再帰は今回の実施範囲に含めない。
準備した完全点ドライバは未実行の補助コードとして明記する。

\subsection{再開用の入口と文書上の訂正}

bundleを展開したディレクトリで次を実行する。
Pythonのassertを有効に保ち、\texttt{-O}、\texttt{-OO}を用いない。
\begin{verbatim}
python3 -B verification/verify_integrity.py
python3 -B verification/run_v18_checks.py
python3 -B verification/run_v18_checks.py --fresh-normal-jets
python3 -B verification/run_v18_checks.py --full-characters
\end{verbatim}
検証入口は六つのソース木を一時コピーし、元のファイルのhashが変わらないことを確認する。
既定の保存監査、新規法三次再帰、全重みの再計算は、実行記録上も区別する。
正確な入力、出力、依存関係と所要時間は
\nolinkurl{source_v18/README_ja.txt} と再現報告にまとめた。

第17版の \nolinkurl{source_v17/README_ja.txt} にある完全点Tateの新規再帰例は、
\texttt{--audit} を付ける点に文書上の誤りがあった。
当該ドライバの \texttt{--audit} は既存JSONの保存ブロック監査を選ぶため、
新規再帰ではこのフラグを外す。
保存されている過去の新規計算記録そのものは変更せず、
現在のREADMEに正しい二種類のコマンドを示す。

\subsection{未解決事項を引き継ぐ}

全次数normal消滅と、normal消滅部分での明示的有限決定性は今回の新しい結論である。
$E_{\mathrm{vert}}=t^{10}$、$\mathcal K_{26}^{\mathrm{nor}}=0$、
実際の26次項の通常像への所属、無限形式変換の収束・代数化は残る。
従って26次開始の既存正規形を28次開始へ更新していない。
新規性や優先権について網羅的な文献調査を行ったとは主張しない。


\begin{thebibliography}{99}
\small
\bibitem{Yek}
A. Yekutieli,
\emph{Twisted deformation quantization of algebraic varieties},
Adv. Math. 268 (2015), 241--305.
Definitions 2.7--2.8, Theorem 3.6, Corollary 11.2.
\url{https://arxiv.org/abs/0905.0488}

\bibitem{May}
E. Mayanskiy,
\emph{Poisson cohomology of two Fano threefolds}.
Lemma 1, Section 2.2, Theorems 1--2.
\url{https://arxiv.org/abs/1303.5916}

\bibitem{Kon}
M. Kontsevich,
\emph{Deformation quantization of Poisson manifolds, I},
Lett. Math. Phys. 66 (2003), 157--216.
Sections 4.6, 6, 8.
\url{https://arxiv.org/abs/q-alg/9709040}

\bibitem{EG}
P. Etingof and V. Ginzburg,
\emph{Noncommutative del Pezzo surfaces and Calabi--Yau algebras},
J. Eur. Math. Soc. 12 (2010), 1371--1416.
Sections 3, 7--8.
\url{https://arxiv.org/abs/0709.3593}

\bibitem{IS}
N. Iyudu and S. Shkarin,
\emph{Sklyanin algebras and a cubic root of 1}.
Theorem 1.1 and Introduction.
\url{https://arxiv.org/abs/2108.06290}

\bibitem{HeYe}
J.-W. He and Y. Ye,
\emph{Clifford deformations of Koszul Frobenius algebras and noncommutative quadrics}.
Lemma 1.2 in the preprint.
\url{https://arxiv.org/abs/1905.04699}

\bibitem{WWY}
C. Walton, X. Wang and M. Yakimov,
\emph{Poisson geometry of PI 3-dimensional Sklyanin algebras}.
Section 2.3, especially Lemma 2.19 (independent of the later PI assumption).
\url{https://arxiv.org/abs/1704.04975}

\bibitem{CCKM}
T. Cassidy, A. Conner, E. Kirkman and W. F. Moore,
\emph{Periodic free resolutions from twisted matrix factorizations}.
Theorems 4.7 and 5.5.
\url{https://arxiv.org/abs/1307.8415}

\bibitem{YZ}
A. Yekutieli and J. J. Zhang,
\emph{Serre duality for noncommutative projective schemes},
Proc. Amer. Math. Soc. 125 (1997), 697--707.
Theorem 4.2 and Corollary 4.3.
\url{https://doi.org/10.1090/S0002-9939-97-03782-9}

\bibitem{IM}
A. Itaba and M. Matsuno,
\emph{AS-regularity of geometric algebras of plane cubic curves},
J. Aust. Math. Soc. 112 (2022), 193--217.
Lemma 2.15 and the proof of Theorem 4.4.
\url{https://doi.org/10.1017/S1446788721000070}

\bibitem{Chen}
X.-W. Chen,
\emph{Generalized Serre duality}.
Theorem 3.10.
\url{https://arxiv.org/abs/math/0610258}

\bibitem{Kuz}
A. Kuznetsov,
\emph{Hochschild homology and semiorthogonal decompositions},
arXiv:0904.4330.
Theorem 7.7, Remark 7.8, Theorem 8.9.
\url{https://arxiv.org/abs/0904.4330}

\bibitem{JLLZ}
A. Jacovskis, X. Lin, Z. Liu and S. Zhang,
\emph{Infinitesimal categorical Torelli theorems for Fano threefolds},
J. Pure Appl. Algebra 227 (2023), 107418.
Theorem 4.4(2).
\url{https://arxiv.org/abs/2203.08187}

\bibitem{KL}
B. Keller and W. Lowen,
\emph{On Hochschild cohomology and Morita deformations},
Int. Math. Res. Not. 2009, no.17, 3221--3235.
\url{https://webusers.imj-prg.fr/~bernhard.keller/publ/KellerLowenHochDefo.pdf}

\bibitem{LowenPre}
W. Lowen,
\emph{Algebroid prestacks and deformations of ringed spaces}.
Theorems 1.4 and 3.26.
\url{https://arxiv.org/abs/math/0511197}

\bibitem{DDL}
O. De Deken and W. Lowen,
\emph{Abelian and derived deformations in the presence of
$\mathbb Z$-generating geometric helices},
J. Noncommut. Geom. 5 (2011), 477--505.
Theorem 4.5, Propositions 5.8--5.11, proof of Theorem 5.13, Lemma 5.14.
\url{https://doi.org/10.4171/JNCG/83}

\bibitem{Orlov}
D. Orlov,
\emph{Derived categories of coherent sheaves and triangulated categories of singularities}.
Theorem 2.5.
\url{https://arxiv.org/abs/math/0503632}

\bibitem{Ueyama}
K. Ueyama,
\emph{Graded maximal Cohen--Macaulay modules over noncommutative graded
Gorenstein isolated singularities},
J. Algebra 383 (2013), 85--103.
\url{https://doi.org/10.1016/j.jalgebra.2013.02.022}
The Serre formula is also stated as Theorem 4 in the author's ICRA XV slides:
\url{https://www.math.uni-bielefeld.de/icra2012/presentations/icra2012_ueyama.pdf}

\bibitem{BKP}
A. Blanc, L. Katzarkov and P. Pandit,
\emph{Generators in formal deformations of categories}.
Definition 3.3 and Section 3.3 (Lurie's comparison, Proposition 5.3.21).
\url{https://arxiv.org/abs/1705.00655}

\bibitem{BS}
J. Burke and G. Stevenson,
\emph{The derived category of a graded Gorenstein ring}.
Definition 5.1, Remark 5.2, Lemma 6.2, Theorem 6.4(3).
\url{https://arxiv.org/abs/1507.00830}

\bibitem{BFK}
M. Ballard, D. Favero and L. Katzarkov,
\emph{A category of kernels for equivariant factorizations and its implications for Hodge theory},
Publ. Math. IHES 120 (2014), 1--111.
Theorem 5.15 and Corollary 5.18.
\url{https://doi.org/10.1007/s10240-013-0059-9}

\bibitem{Elagin}
A. Elagin,
\emph{On equivariant triangulated categories}.
Theorem 4.2 and the dg enhancement construction.
\url{https://arxiv.org/abs/1403.7027}

\bibitem{OrlovGlue}
D. Orlov,
\emph{Smooth and proper noncommutative schemes and gluing of DG categories},
Adv. Math. 302 (2016), 59--105.
Proposition 3.8, Theorems 3.18 and 3.25.
\url{https://arxiv.org/abs/1402.7364}

\bibitem{Zhang}
J. J. Zhang,
\emph{Twisted graded algebras and equivalences of graded categories},
Proc. London Math. Soc. (3) 72 (1996), 281--311.
Theorem 1.1.
\url{https://doi.org/10.1112/plms/s3-72.2.281}

\bibitem{AV}
M. Artin and M. Van den Bergh,
\emph{Twisted homogeneous coordinate rings},
J. Algebra 133 (1990), 249--271.
\url{https://doi.org/10.1016/0021-8693(90)90269-T}

\bibitem{SVdB}
J. T. Stafford and M. Van den Bergh,
\emph{Noncommutative curves and noncommutative surfaces},
Bull. Amer. Math. Soc. 38 (2001), 171--216.
Theorem 3.2 and Section 8 (the quantum triangle).
\url{https://arxiv.org/abs/math/9910082}
\bibitem{FLZ}
S. Feyzbakhsh, Z. Liu and S. Zhang,
\emph{New perspectives on categorical Torelli theorems for del Pezzo threefolds}.
Section 3.2, equations (9)--(12), Proposition 5.5, Remark 7.2.
\url{https://arxiv.org/abs/2304.01321}

\bibitem{ELOI}
A. I. Efimov, V. A. Lunts and D. O. Orlov,
\emph{Deformation theory of objects in homotopy and derived categories I: general theory}.
Proposition 6.1 and Section 7.
\url{https://arxiv.org/abs/math/0702838}

\bibitem{ELOII}
A. I. Efimov, V. A. Lunts and D. O. Orlov,
\emph{Deformation theory of objects in homotopy and derived categories II:
pro-representability of the deformation functor}.
Sections 3, 5--7; Theorem 7.2, Theorem 11.8, Proposition 11.10.
\url{https://arxiv.org/abs/math/0702839}

\bibitem{RZ}
D. Rogalski and J. J. Zhang,
\emph{Canonical Maps to Twisted Rings}.
Introduction and Section 4.
\url{https://arxiv.org/abs/math/0409405}

\bibitem{Getzler}
E. Getzler,
\emph{Lie theory for nilpotent $L_\infty$-algebras},
Ann. of Math. 170 (2009), 271--301.
Proposition 1.1, Section 4, Proposition 4.4.
\url{https://doi.org/10.4007/annals.2009.170.271}

\bibitem{TV}
B. To\"en and M. Vaqui\'e,
\emph{Moduli of objects in dg-categories},
Ann. Sci. \'Ecole Norm. Sup. (4) 40 (2007), 387--444.
Lemma 2.8, Theorem 3.6, Corollaries 3.15, 3.17, 3.21--3.22.
\url{https://arxiv.org/abs/math/0503269}

\bibitem{FP}
S. Feyzbakhsh and L. Pertusi,
\emph{Serre-invariant stability conditions and Ulrich bundles on cubic threefolds},
\'Epijournal de G\'eom\'etrie Alg\'ebrique 7 (2023), Article 1.
Section 6, Proposition 6.3, Corollary 6.5.
\url{https://doi.org/10.46298/epiga.2022.9611}

\bibitem{PY}
L. Pertusi and S. Yang,
\emph{Some remarks on Fano threefolds of index two and stability conditions}.
Theorem 1.2.
\url{https://arxiv.org/abs/2004.02798}

\bibitem{StacksUnramified}
The Stacks Project Authors,
\emph{The Stacks Project},
Lemma 41.17.1, Tag 04HH (\'etale local structure of unramified morphisms).
\url{https://stacks.math.columbia.edu/tag/04HH}

\bibitem{AIK}
J. Alper, A. V. Isaev and N. G. Kruzhilin,
\emph{Associated forms of binary quartics and ternary cubics},
Transformation Groups 21 (2016), no. 3, 593--618.
Section 3.2, equations (3.11), (3.18), Remark 3.1.
\url{https://arxiv.org/abs/1409.8369}

\bibitem{IMRelations}
A. Itaba and M. Matsuno,
\emph{Defining relations of 3-dimensional quadratic AS-regular algebras},
Mathematical Journal of Okayama University 63 (2021), no. 1, 61--86.
Theorems 2.8, 4.12, 4.16 and Lemma 4.17.
\url{https://www.math.okayama-u.ac.jp/mjou/mjou63/_04_Itaba.pdf}

\bibitem{ATT}
A. S. I. Anema, J. Top and A. Tuijp,
\emph{Hesse Pencils and 3-Torsion Structures},
SIGMA 14 (2018), 102, 13 pages.
Theorem 1.1 and Sections 2--4.
\url{https://doi.org/10.3842/SIGMA.2018.102}
\bibitem{KellerHoch}
B. Keller,
\emph{Derived invariance of higher structures on the Hochschild complex}.
Manuscript, 2003, revised May 16, 2018.
Section 4.5, equations (4.5.1)--(4.5.2).
\url{https://webusers.imj-prg.fr/~bernhard.keller/publ/dih.pdf}

\bibitem{ToenMorita}
B. To\"en,
\emph{The homotopy theory of dg-categories and derived Morita theory},
Invent. Math. 167 (2007), 615--667.
Section 8.1 and Theorem 8.9.
\url{https://arxiv.org/abs/math/0408337}

\Needspace{5\baselineskip}
\bibitem{LowenObstruction}
W. Lowen,
\emph{Obstruction Theory for Objects in Abelian and Derived Categories},
Comm. Algebra \textbf{33} (2005), 3195--3223.
Theorem 3.8, Corollary 3.11.
\href{https://doi.org/10.1081/AGB-200066155}{doi:10.1081/AGB-200066155}.
\url{https://arxiv.org/abs/math/0407019}

\bibitem{StacksNilPerfect}
The Stacks Project Authors,
\emph{The Stacks Project}, Lemma 15.80.4, Tag 07LU.
\url{https://stacks.math.columbia.edu/tag/07LU}

\bibitem{StacksPicardSmooth}
The Stacks Project Authors,
\emph{The Stacks Project}, Lemma 108.8.5, Tag 0DPJ.
\url{https://stacks.math.columbia.edu/tag/0DPJ}

\bibitem{PicardFormal}
A. Chatzistamatiou,
\emph{Commutative formal groups arising from schemes},
Appendix A.2, Proposition A.2.2.
\url{https://arxiv.org/abs/1203.4926}

\bibitem{ChowGrass}
W.-L. Chow,
\emph{On the Geometry of Algebraic Homogeneous Spaces},
Ann. of Math. (2) \textbf{50} (1949), 32--67.
\href{https://doi.org/10.2307/1969351}{doi:10.2307/1969351}.

\bibitem{GhorpadeKaipa}
S. R. Ghorpade and K. V. Kaipa,
\emph{Automorphism groups of Grassmann codes},
Finite Fields Appl. \textbf{23} (2013), 80--102.
Theorem 3.6.
\url{https://arxiv.org/abs/1204.5547}

\bibitem{SGA2}
A. Grothendieck, avec un expos\'e de M. Raynaud,
\emph{Cohomologie locale des faisceaux coh\'erents et th\'eor\`emes de Lefschetz
locaux et globaux (SGA 2)},
Y. Laszlo (\'ed.), Documents Math\'ematiques \textbf{4}, SMF, 2005.
Expos\'e XII, Corollaire 3.7.
\url{https://arxiv.org/abs/math/0511279}


\bibitem{LieblichComplexes}
M. Lieblich,
\emph{Moduli of complexes on a proper morphism},
J. Algebraic Geom. 15 (2006), 175--206.
Theorem 3.1.1, Corollary 3.1.2, Theorem 4.2.1, Corollary 4.3.3.
\url{https://arxiv.org/abs/math/0502198}

\bibitem{StacksStrictTransform}
The Stacks Project Authors,
Lemma 31.34.2, Tag 080E (strict transformとblowup).
\url{https://stacks.math.columbia.edu/tag/080E}

\bibitem{StacksProperFinite}
The Stacks Project Authors,
Lemma 37.44.1, Tag 02LS (properかつquasi-finiteな射の有限性).
\url{https://stacks.math.columbia.edu/tag/02LS}

\bibitem{MustataSing}
M. Musta\c{t}\u{a},
\emph{Singularities in birational geometry: minimal log discrepancies and the log canonical threshold},
AIM workshop notes, Section 1, pp.~1--3.
\url{https://www.aimath.org/WWN/singularvariety/sing.pdf}


\bibitem{BeteaWheeler}
D. Betea and M. Wheeler,
\emph{Refined Cauchy and Littlewood identities, plane partitions and
symmetry classes of alternating sign matrices}, arXiv:1402.0229, equation (6).
\url{https://arxiv.org/abs/1402.0229}

\bibitem{DeBoeckPagetWildon}
M. de Boeck, R. Paget and M. Wildon,
\emph{Plethysms of symmetric functions and highest weight representations},
Trans. Amer. Math. Soc. 374 (2021), 8013--8043, Example 1.6.
\url{https://doi.org/10.1090/tran/8481}
\url{https://arxiv.org/abs/1810.03448}


\bibitem{DyckerhoffMF}
T. Dyckerhoff,
\emph{Compact generators in categories of matrix factorizations},
Duke Math. J. \textbf{159} (2011), 223--274.
\url{https://arxiv.org/abs/0904.4713}.

\bibitem{BDSS}
L. Bus\'e, A. Dimca, H. Schenck and G. Sticlaru,
\emph{The Hessian polynomial and the Jacobian ideal of a reduced
hypersurface in $\mathbf P^n$}, arXiv:1910.09195v4 (2021).
\url{https://arxiv.org/abs/1910.09195}.

\bibitem{StacksFormalFunctions}
The Stacks Project Authors,
\emph{Theorem on formal functions}, Theorem 30.20.5, Tag 02OC.
\url{https://stacks.math.columbia.edu/tag/02OC}.


\bibitem{StacksEtaleCriterion}
The Stacks Project Authors,
\emph{\'Etale morphisms}, Proposition 59.26.2, Tag 03PA;
Section 29.37, Tag 02GH.
\url{https://stacks.math.columbia.edu/tag/03PA}

\bibitem{StacksFaithfulCompletion}
The Stacks Project Authors,
Lemma 10.97.3, Tag 00MC (faithful flatness of completion).
\url{https://stacks.math.columbia.edu/tag/00MC}

\bibitem{StacksFlatBaseChange}
The Stacks Project Authors,
\emph{Flat base change}, Lemma 30.5.2, Tag 02KH.
\url{https://stacks.math.columbia.edu/tag/02KH}

% 既存原稿の thebibliography 環境の中へ追加する項目。
\bibitem{Fisher2008}
T. Fisher,
\emph{The invariants of a genus one curve},
Proceedings of the London Mathematical Society (3) \textbf{97} (2008), 753--782.
\url{https://doi.org/10.1112/plms/pdn021}.
著者公開版：\url{https://www.dpmms.cam.ac.uk/~taf1000/papers/g1inv.pdf},
\S3.3, \S4, Lemma~4.3, Theorem~4.4(i).

\bibitem{SwanMKN}
R. G. Swan,
\emph{On a Theorem of Mohan Kumar and Nori},
著者公開原稿、\S3, Corollary~3.8, Lemma~3.9.
\url{https://www.math.uchicago.edu/~swan/MKN.pdf}.

\bibitem{StacksCMRegular}
The Stacks Project Authors,
\emph{The Stacks Project},
Lemmas 10.106.3 and 10.104.2, Tags 00NQ and 02JN.
\url{https://stacks.math.columbia.edu/tag/00NQ},
\url{https://stacks.math.columbia.edu/tag/02JN}.


% 既存のthebibliographyの内部へ追加する。
\bibitem{FisherHessian2012}
T.~Fisher,
\emph{The Hessian of a genus one curve},
Proc.\ London Math.\ Soc.\ (3) \textbf{104} (2012), 613--648,
doi:10.1112/plms/pdr039.
著者版：\url{https://www.dpmms.cam.ac.uk/~taf1000/papers/g1hessians.pdf}.

\bibitem{StacksCechKoszul}
The Stacks Project Authors,
\emph{The Stacks Project},
Tag 012X, \emph{Spectral sequences: double complexes},
\url{https://stacks.math.columbia.edu/tag/012X};
Tag 0G6L,
\url{https://stacks.math.columbia.edu/tag/0G6L}.

% 既存thebibliographyへ追加する第18版理論の一次資料。
\bibitem{Dube1990}
T.~W.~Dub\'e,
\emph{The structure of polynomial ideals and Gr\"obner bases},
SIAM Journal on Computing \textbf{19} (1990), no.~4, 750--773.
\href{https://doi.org/10.1137/0219053}{doi:10.1137/0219053}.
斉次版はTheorem~8.2、非斉次版はCorollary~8.3、いずれもp.~773。
一次論文PDF:
\url{https://cs.uwaterloo.ca/~r5olivei/courses/2025-spring-cs860/D90.pdf}.

\end{thebibliography}
\end{document}
